% Options for packages loaded elsewhere
\PassOptionsToPackage{unicode}{hyperref}
\PassOptionsToPackage{hyphens}{url}
\documentclass[
  11pt,
]{article}
\usepackage{xcolor}
\usepackage[margin=18mm]{geometry}
\usepackage{amsmath,amssymb}
\setcounter{secnumdepth}{5}
\usepackage{iftex}
\ifPDFTeX
  \usepackage[T1]{fontenc}
  \usepackage[utf8]{inputenc}
  \usepackage{textcomp} % provide euro and other symbols
\else % if luatex or xetex
  \usepackage{unicode-math} % this also loads fontspec
  \defaultfontfeatures{Scale=MatchLowercase}
  \defaultfontfeatures[\rmfamily]{Ligatures=TeX,Scale=1}
\fi
\usepackage{lmodern}
\ifPDFTeX\else
  % xetex/luatex font selection
    \setmainfont[]{Noto Serif}
    \setsansfont[]{Noto Sans}
    \setmonofont[]{Noto Sans Mono}
  \setmathfont[]{Noto Sans Math}
\fi
% Use upquote if available, for straight quotes in verbatim environments
\IfFileExists{upquote.sty}{\usepackage{upquote}}{}
\IfFileExists{microtype.sty}{% use microtype if available
  \usepackage[]{microtype}
  \UseMicrotypeSet[protrusion]{basicmath} % disable protrusion for tt fonts
}{}
\makeatletter
\@ifundefined{KOMAClassName}{% if non-KOMA class
  \IfFileExists{parskip.sty}{%
    \usepackage{parskip}
  }{% else
    \setlength{\parindent}{0pt}
    \setlength{\parskip}{6pt plus 2pt minus 1pt}}
}{% if KOMA class
  \KOMAoptions{parskip=half}}
\makeatother
\ifLuaTeX
\usepackage[bidi=basic]{babel}
\else
\usepackage[bidi=default]{babel}
\fi
\babelprovide[main,import]{russian}
\ifPDFTeX
\else
\babelfont{rm}[]{Noto Serif}
\fi
% get rid of language-specific shorthands (see #6817):
\let\LanguageShortHands\languageshorthands
\def\languageshorthands#1{}
\setlength{\emergencystretch}{3em} % prevent overfull lines
\providecommand{\tightlist}{%
  \setlength{\itemsep}{0pt}\setlength{\parskip}{0pt}}
\usepackage{bookmark}
\IfFileExists{xurl.sty}{\usepackage{xurl}}{} % add URL line breaks if available
\urlstyle{same}
\hypersetup{
  pdftitle={Билеты по дискретным случайным процессам},
  pdflang={ru-RU},
  hidelinks,
  pdfcreator={LaTeX via pandoc}}

\title{Билеты по дискретным случайным процессам}
\author{}
\date{}

\begin{document}
\maketitle

{
\setcounter{tocdepth}{2}
\tableofcontents
}
\section{Билет 01. Системы множеств: алгебры, полуалгебры,
сигма-алгебры. Меры Дирака и
Жордана}\label{ux431ux438ux43bux435ux442-01.-ux441ux438ux441ux442ux435ux43cux44b-ux43cux43dux43eux436ux435ux441ux442ux432-ux430ux43bux433ux435ux431ux440ux44b-ux43fux43eux43bux443ux430ux43bux433ux435ux431ux440ux44b-ux441ux438ux433ux43cux430-ux430ux43bux433ux435ux431ux440ux44b.-ux43cux435ux440ux44b-ux434ux438ux440ux430ux43aux430-ux438-ux436ux43eux440ux434ux430ux43dux430}

Источники: Ширяев 1, гл. II; лекции ДСП.

\subsection{1. Короткий
ответ}\label{ux43aux43eux440ux43eux442ux43aux438ux439-ux43eux442ux432ux435ux442}

В теории вероятностей в общем случае нельзя требовать, чтобы хорошая
вероятность была задана на всех подмножествах пространства исходов.
Сначала выбирают систему множеств, устойчивую относительно нужных
операций: алгебру, полуалгебру или \(\sigma\)-алгебру. Алгебра замкнута
относительно конечных операций, \(\sigma\)-алгебра - относительно
счетных операций, а полуалгебра удобна как простая область задания меры
перед ее продолжением. Два базовых примера мер: мера Дирака

\[
\delta_x(A)=\mathbf 1_A(x)
\]

и мера Жордана, которая на прямоугольнике \(I=\prod_{k=1}^d[a_k,b_k)\)
задается формулой

\[
v(I)=\prod_{k=1}^d(b_k-a_k).
\]

\subsection{2. Полный
ответ}\label{ux43fux43eux43bux43dux44bux439-ux43eux442ux432ux435ux442}

\subsubsection{2.1. Пространство элементарных исходов и
события}\label{ux43fux440ux43eux441ux442ux440ux430ux43dux441ux442ux432ux43e-ux44dux43bux435ux43cux435ux43dux442ux430ux440ux43dux44bux445-ux438ux441ux445ux43eux434ux43eux432-ux438-ux441ux43eux431ux44bux442ux438ux44f}

Пусть фиксировано множество \(\Omega\). Его элементы называются
элементарными исходами, а подмножества \(\Omega\) интерпретируются как
события. Однако в общем бесконечном случае нельзя заранее рассчитывать,
что нужная вероятность будет корректно задана на всех подмножествах.
Поэтому выбирают класс допустимых множеств
\(\mathcal F\subset 2^\Omega\) и требуют, чтобы он был устойчив
относительно естественных операций над событиями.

\subsubsection{2.2. Требования к системам
множеств}\label{ux442ux440ux435ux431ux43eux432ux430ux43dux438ux44f-ux43a-ux441ux438ux441ux442ux435ux43cux430ux43c-ux43cux43dux43eux436ux435ux441ux442ux432}

Минимальные требования зависят от того, какие операции мы хотим
выполнять. Если нужно брать только конечные объединения, пересечения и
дополнения, достаточно алгебры. Если нужно работать с пределами
последовательностей событий, например с событиями вида
\(\bigcup_{n=1}^{\infty}A_n\), \(\bigcap_{n=1}^{\infty}A_n\),
\(\limsup A_n\), нужна \(\sigma\)-алгебра. Полуалгебра обычно
используется как ``кирпичики'': на ней легко задать меру, например длину
полуинтервала, а затем продолжить меру на алгебру и дальше на
\(\sigma\)-алгебру.

\subsubsection{2.3. Полуалгебры и алгебры элементарных
множеств}\label{ux43fux43eux43bux443ux430ux43bux433ux435ux431ux440ux44b-ux438-ux430ux43bux433ux435ux431ux440ux44b-ux44dux43bux435ux43cux435ux43dux442ux430ux440ux43dux44bux445-ux43cux43dux43eux436ux435ux441ux442ux432}

Для вероятностей главная система множеств - \(\sigma\)-алгебра событий
\(\mathcal F\). Но перед ней часто стоят более простые классы. Например,
на \(\mathbb R\) естественно начать с полуинтервалов \([a,b)\), на
\(\mathbb R^d\) - с полуоткрытых прямоугольников
\(\prod_{k=1}^d[a_k,b_k)\). Они образуют полуалгебру, потому что
пересечение двух таких прямоугольников снова прямоугольник такого же
вида, а разность прямоугольников раскладывается в конечное дизъюнктное
объединение прямоугольников. Из полуалгебры берут конечные дизъюнктные
объединения и получают алгебру элементарных множеств.

\subsubsection{2.4. Понятие меры и свойство
аддитивности}\label{ux43fux43eux43dux44fux442ux438ux435-ux43cux435ux440ux44b-ux438-ux441ux432ux43eux439ux441ux442ux432ux43e-ux430ux434ux434ux438ux442ux438ux432ux43dux43eux441ux442ux438}

Мера на такой системе множеств - это функция, которая каждому
допустимому множеству ставит в соответствие его ``размер''. В этом
билете достаточно понимать меру как неотрицательную конечно-аддитивную
функцию на выбранной системе множеств. Счетная аддитивность будет
центральной темой следующего билета. Для попарно непересекающихся
допустимых множеств \(A_1,\ldots,A_n\) конечная аддитивность означает

\[
\mu\left(\bigsqcup_{k=1}^n A_k\right)=\sum_{k=1}^n\mu(A_k).
\]

\subsubsection{2.5. Мера
Дирака}\label{ux43cux435ux440ux430-ux434ux438ux440ux430ux43aux430}

Мера Дирака \(\delta_x\) сосредоточивает всю массу в одной точке \(x\).
Она равна \(1\), если множество содержит \(x\), и равна \(0\), если не
содержит. Это простейшая вероятностная мера: вся случайность вырождается
в один исход.

\subsubsection{2.6. Мера Жордана как геометрический
объем}\label{ux43cux435ux440ux430-ux436ux43eux440ux434ux430ux43dux430-ux43aux430ux43a-ux433ux435ux43eux43cux435ux442ux440ux438ux447ux435ux441ux43aux438ux439-ux43eux431ux44aux435ux43c}

Мера Жордана - геометрический объем элементарных множеств. На прямой это
длина, на плоскости площадь, в \(\mathbb R^d\) - \(d\)-мерный объем.
Сначала ее задают на прямоугольниках формулой \(v(I)=\prod(b_k-a_k)\),
затем на элементарном множестве \(A=\bigsqcup_{j=1}^m I_j\) кладут

\[
v(A)=\sum_{j=1}^m v(I_j).
\]

\subsubsection{2.7. Корректность меры
Жордана}\label{ux43aux43eux440ux440ux435ux43aux442ux43dux43eux441ux442ux44c-ux43cux435ux440ux44b-ux436ux43eux440ux434ux430ux43dux430}

Основной вопрос здесь: не зависит ли эта сумма от способа разбиения
\(A\) на прямоугольники? Ответ: не зависит. Это доказывается переходом к
общему измельчению двух разбиений. Поэтому мера Жордана корректно
определена на алгебре элементарных множеств.

\subsection{3. Основные
определения}\label{ux43eux441ux43dux43eux432ux43dux44bux435-ux43eux43fux440ux435ux434ux435ux43bux435ux43dux438ux44f}

\subsubsection{Система
множеств}\label{ux441ux438ux441ux442ux435ux43cux430-ux43cux43dux43eux436ux435ux441ux442ux432}

Любое семейство подмножеств \(\mathcal C\subset 2^\Omega\).

\subsubsection{Алгебра
множеств}\label{ux430ux43bux433ux435ux431ux440ux430-ux43cux43dux43eux436ux435ux441ux442ux432}

Семейство \(\mathcal A\subset 2^\Omega\) называется алгеброй на
\(\Omega\), если:

\begin{enumerate}
\def\labelenumi{\arabic{enumi}.}
\tightlist
\item
  \(\Omega\in\mathcal A\);
\item
  если \(A\in\mathcal A\), то
  \(\overline{A}=\Omega\setminus A\in\mathcal A\);
\item
  если \(A,B\in\mathcal A\), то \(A\cup B\in\mathcal A\).
\end{enumerate}

Из этих условий сразу следует, что \(\varnothing\in\mathcal A\),
\(A\cap B\in\mathcal A\), \(A\setminus B\in\mathcal A\), а также что
алгебра замкнута относительно любых конечных объединений и конечных
пересечений.

\subsubsection{\texorpdfstring{\(\sigma\)-алгебра}{\textbackslash sigma-алгебра}}\label{sigma-ux430ux43bux433ux435ux431ux440ux430}

Семейство \(\mathcal F\subset 2^\Omega\) называется \(\sigma\)-алгеброй,
если:

\begin{enumerate}
\def\labelenumi{\arabic{enumi}.}
\tightlist
\item
  \(\Omega\in\mathcal F\);
\item
  \(A\in\mathcal F\Rightarrow \overline{A}\in\mathcal F\);
\item
  для любой последовательности \(A_1,A_2,\ldots\in\mathcal F\) выполнено
  \(\bigcup_{n=1}^{\infty}A_n\in\mathcal F\).
\end{enumerate}

Любая \(\sigma\)-алгебра является алгеброй, но не всякая алгебра
является \(\sigma\)-алгеброй.

\subsubsection{Полуалгебра}\label{ux43fux43eux43bux443ux430ux43bux433ux435ux431ux440ux430}

Семейство \(\mathcal S\subset 2^\Omega\) называется полуалгеброй, если:

\begin{enumerate}
\def\labelenumi{\arabic{enumi}.}
\tightlist
\item
  \(\varnothing\in\mathcal S\);
\item
  если \(A,B\in\mathcal S\), то \(A\cap B\in\mathcal S\);
\item
  если \(A,B\in\mathcal S\), то разность \(A\setminus B\) представима
  как конечное дизъюнктное объединение множеств из \(\mathcal S\):
\end{enumerate}

\[
A\setminus B=\bigsqcup_{k=1}^n C_k,\qquad C_k\in\mathcal S.
\]

Иногда дополнительно требуют \(\Omega\in\mathcal S\); в вероятностных
приложениях это удобно, но главное свойство полуалгебры - хорошее
поведение пересечений и разностей.

\subsubsection{Мера на
алгебре}\label{ux43cux435ux440ux430-ux43dux430-ux430ux43bux433ux435ux431ux440ux435}

Неотрицательная функция \(\mu:\mathcal A\to[0,\infty]\) называется
конечно-аддитивной мерой, если \(\mu(\varnothing)=0\) и для любых
попарно непересекающихся \(A_1,\ldots,A_n\in\mathcal A\), таких что
\(\bigsqcup_{k=1}^nA_k\in\mathcal A\), выполнено

\[
\mu\left(\bigsqcup_{k=1}^nA_k\right)=\sum_{k=1}^n\mu(A_k).
\]

\subsubsection{Мера
Дирака}\label{ux43cux435ux440ux430-ux434ux438ux440ux430ux43aux430-1}

Пусть \(x\in\Omega\). На любой \(\sigma\)-алгебре
\(\mathcal F\subset 2^\Omega\) мера Дирака в точке \(x\) задается
формулой

\[
\delta_x(A)=
\begin{cases}
1, & x\in A,\\
0, & x\notin A.
\end{cases}
\]

\subsubsection{\texorpdfstring{Полуоткрытый прямоугольник в
\(\mathbb R^d\)}{Полуоткрытый прямоугольник в \textbackslash mathbb R\^{}d}}\label{ux43fux43eux43bux443ux43eux442ux43aux440ux44bux442ux44bux439-ux43fux440ux44fux43cux43eux443ux433ux43eux43bux44cux43dux438ux43a-ux432-mathbb-rd}

Множество вида

\[
I=\prod_{k=1}^d[a_k,b_k),\qquad a_k\le b_k.
\]

Если хотя бы для одного \(k\) выполнено \(a_k=b_k\), прямоугольник
вырожден и имеет объем \(0\).

\subsubsection{Элементарное
множество}\label{ux44dux43bux435ux43cux435ux43dux442ux430ux440ux43dux43eux435-ux43cux43dux43eux436ux435ux441ux442ux432ux43e}

Конечное дизъюнктное объединение полуоткрытых прямоугольников:

\[
A=\bigsqcup_{j=1}^m I_j.
\]

\subsubsection{Мера Жордана на элементарных
множествах}\label{ux43cux435ux440ux430-ux436ux43eux440ux434ux430ux43dux430-ux43dux430-ux44dux43bux435ux43cux435ux43dux442ux430ux440ux43dux44bux445-ux43cux43dux43eux436ux435ux441ux442ux432ux430ux445}

Сначала на прямоугольнике:

\[
v(I)=\prod_{k=1}^d(b_k-a_k),
\]

затем на элементарном множестве:

\[
v(A)=\sum_{j=1}^m v(I_j),\qquad A=\bigsqcup_{j=1}^m I_j.
\]

\subsection{4. Основные теоремы и
свойства}\label{ux43eux441ux43dux43eux432ux43dux44bux435-ux442ux435ux43eux440ux435ux43cux44b-ux438-ux441ux432ux43eux439ux441ux442ux432ux430}

\subsubsection{Свойство 1. Алгебра замкнута относительно конечных
пересечений и
разностей}\label{ux441ux432ux43eux439ux441ux442ux432ux43e-1.-ux430ux43bux433ux435ux431ux440ux430-ux437ux430ux43cux43aux43dux443ux442ux430-ux43eux442ux43dux43eux441ux438ux442ux435ux43bux44cux43dux43e-ux43aux43eux43dux435ux447ux43dux44bux445-ux43fux435ux440ux435ux441ux435ux447ux435ux43dux438ux439-ux438-ux440ux430ux437ux43dux43eux441ux442ux435ux439}

Если \(\mathcal A\) - алгебра и \(A,B\in\mathcal A\), то

\[
A\cap B\in\mathcal A,\qquad A\setminus B\in\mathcal A.
\]

\subsubsection{Доказательство}\label{ux434ux43eux43aux430ux437ux430ux442ux435ux43bux44cux441ux442ux432ux43e}

По законам де Моргана

\[
A\cap B=\overline{\overline{A}\cup \overline{B}}.
\]

Так как \(\overline{A},\overline{B}\in\mathcal A\), то
\(\overline{A}\cup \overline{B}\in\mathcal A\), а значит и его
дополнение лежит в \(\mathcal A\). Далее

\[
A\setminus B=A\cap \overline{B},
\]

а пересечение двух множеств из \(\mathcal A\) уже доказано. \(\square\)

\subsubsection{\texorpdfstring{Свойство 2. \(\sigma\)-алгебра замкнута
относительно счетных
пересечений}{Свойство 2. \textbackslash sigma-алгебра замкнута относительно счетных пересечений}}\label{ux441ux432ux43eux439ux441ux442ux432ux43e-2.-sigma-ux430ux43bux433ux435ux431ux440ux430-ux437ux430ux43cux43aux43dux443ux442ux430-ux43eux442ux43dux43eux441ux438ux442ux435ux43bux44cux43dux43e-ux441ux447ux435ux442ux43dux44bux445-ux43fux435ux440ux435ux441ux435ux447ux435ux43dux438ux439}

Если \(A_n\in\mathcal F\) для всех \(n\), то

\[
\bigcap_{n=1}^{\infty}A_n\in\mathcal F.
\]

\subsubsection{Доказательство}\label{ux434ux43eux43aux430ux437ux430ux442ux435ux43bux44cux441ux442ux432ux43e-1}

Снова используем формулу де Моргана:

\[
\bigcap_{n=1}^{\infty}A_n=\overline{\left(\bigcup_{n=1}^{\infty}\overline{A_n}\right)}.
\]

Каждое \(\overline{A_n}\) лежит в \(\mathcal F\), счетное объединение
лежит в \(\mathcal F\), и дополнение тоже лежит в \(\mathcal F\).
\(\square\)

\subsubsection{Свойство 3. Конечные дизъюнктные объединения элементов
полуалгебры образуют
алгебру}\label{ux441ux432ux43eux439ux441ux442ux432ux43e-3.-ux43aux43eux43dux435ux447ux43dux44bux435-ux434ux438ux437ux44aux44eux43dux43aux442ux43dux44bux435-ux43eux431ux44aux435ux434ux438ux43dux435ux43dux438ux44f-ux44dux43bux435ux43cux435ux43dux442ux43eux432-ux43fux43eux43bux443ux430ux43bux433ux435ux431ux440ux44b-ux43eux431ux440ux430ux437ux443ux44eux442-ux430ux43bux433ux435ux431ux440ux443}

Пусть \(\mathcal S\) - полуалгебра и \(\Omega\in\mathcal S\). Рассмотрим
класс

\[
\mathcal A(\mathcal S)=\left\{\bigsqcup_{k=1}^n S_k:\ S_k\in\mathcal S,\ n<\infty\right\}.
\]

Тогда \(\mathcal A(\mathcal S)\) является алгеброй.

\subsubsection{Доказательство}\label{ux434ux43eux43aux430ux437ux430ux442ux435ux43bux44cux441ux442ux432ux43e-2}

Во-первых, \(\Omega\in\mathcal A(\mathcal S)\), потому что
\(\Omega\in\mathcal S\). Во-вторых, конечное объединение двух множеств
из \(\mathcal A(\mathcal S)\) можно сделать дизъюнктным разбиением,
последовательно вычитая уже учтенные части. Полуалгебра гарантирует, что
каждая разность раскладывается в конечное дизъюнктное объединение
элементов \(\mathcal S\). Значит объединение снова лежит в
\(\mathcal A(\mathcal S)\).

Остается дополнение. Если

\[
A=\bigsqcup_{k=1}^n S_k,
\]

то \(\Omega\setminus A\) получается последовательным вычитанием
\(S_1,\ldots,S_n\) из \(\Omega\). На каждом шаге разность элемента
полуалгебры с элементом полуалгебры раскладывается в конечное
дизъюнктное объединение элементов полуалгебры. После конечного числа
шагов получаем конечное дизъюнктное объединение элементов
\(\mathcal S\). Значит \(\overline{A}\in\mathcal A(\mathcal S)\).
\(\square\)

\subsubsection{Свойство 4. Мера Дирака конечно- и
счетно-аддитивна}\label{ux441ux432ux43eux439ux441ux442ux432ux43e-4.-ux43cux435ux440ux430-ux434ux438ux440ux430ux43aux430-ux43aux43eux43dux435ux447ux43dux43e--ux438-ux441ux447ux435ux442ux43dux43e-ux430ux434ux434ux438ux442ux438ux432ux43dux430}

Если \(A_1,A_2,\ldots\) попарно не пересекаются, то

\[
\delta_x\left(\bigsqcup_{n=1}^{\infty}A_n\right)=\sum_{n=1}^{\infty}\delta_x(A_n).
\]

\subsubsection{Доказательство}\label{ux434ux43eux43aux430ux437ux430ux442ux435ux43bux44cux441ux442ux432ux43e-3}

Если \(x\notin\bigcup_n A_n\), то обе части равны \(0\). Если
\(x\in\bigcup_n A_n\), то из попарной непересекаемости существует ровно
один номер \(n_0\), для которого \(x\in A_{n_0}\). Тогда левая часть
равна \(1\), а в правой сумме ровно одно слагаемое равно \(1\),
остальные равны \(0\). Значит равенство выполнено. \(\square\)

\subsubsection{Свойство 5. Полуоткрытые прямоугольники образуют
полуалгебру}\label{ux441ux432ux43eux439ux441ux442ux432ux43e-5.-ux43fux43eux43bux443ux43eux442ux43aux440ux44bux442ux44bux435-ux43fux440ux44fux43cux43eux443ux433ux43eux43bux44cux43dux438ux43aux438-ux43eux431ux440ux430ux437ux443ux44eux442-ux43fux43eux43bux443ux430ux43bux433ux435ux431ux440ux443}

В \(\mathbb R^d\) класс

\[
\mathcal S_d=\left\{\prod_{k=1}^d[a_k,b_k):\ a_k\le b_k\right\}
\]

является полуалгеброй.

\subsubsection{Доказательство}\label{ux434ux43eux43aux430ux437ux430ux442ux435ux43bux44cux441ux442ux432ux43e-4}

Пустое множество можно получить как вырожденный прямоугольник, например
\([a,a)\times\prod_{k=2}^d[a_k,b_k)\). Пересечение двух прямоугольников

\[
I=\prod_{k=1}^d[a_k,b_k),\qquad J=\prod_{k=1}^d[c_k,d_k)
\]

равно

\[
I\cap J=\prod_{k=1}^d[\max(a_k,c_k),\min(b_k,d_k)),
\]

если все интервалы невырождены; иначе пересечение пусто. Значит
пересечение снова принадлежит \(\mathcal S_d\).

Для разности \(I\setminus J\) нужно разрезать \(I\) гиперплоскостями,
проходящими через координаты \(c_k,d_k\). В каждой координате возникает
конечное разбиение на полуинтервалы. Декартовы произведения этих
полуинтервалов дают конечное разбиение \(I\) на полуоткрытые
прямоугольники; часть из них лежит в \(J\), остальные образуют
\(I\setminus J\). Значит разность есть конечное дизъюнктное объединение
элементов \(\mathcal S_d\). \(\square\)

\subsubsection{Свойство 6. Объем Жордана корректно определен на
элементарных
множествах}\label{ux441ux432ux43eux439ux441ux442ux432ux43e-6.-ux43eux431ux44aux435ux43c-ux436ux43eux440ux434ux430ux43dux430-ux43aux43eux440ux440ux435ux43aux442ux43dux43e-ux43eux43fux440ux435ux434ux435ux43bux435ux43d-ux43dux430-ux44dux43bux435ux43cux435ux43dux442ux430ux440ux43dux44bux445-ux43cux43dux43eux436ux435ux441ux442ux432ux430ux445}

Если элементарное множество \(A\) имеет два дизъюнктных разбиения на
полуоткрытые прямоугольники,

\[
A=\bigsqcup_{i=1}^m I_i=\bigsqcup_{j=1}^n J_j,
\]

то

\[
\sum_{i=1}^m v(I_i)=\sum_{j=1}^n v(J_j).
\]

\subsubsection{Доказательство}\label{ux434ux43eux43aux430ux437ux430ux442ux435ux43bux44cux441ux442ux432ux43e-5}

Рассмотрим все непустые пересечения \(I_i\cap J_j\). Они являются
полуоткрытыми прямоугольниками. Для каждого фиксированного \(i\)
множество \(I_i\) разбивается в дизъюнктное объединение прямоугольников
\(I_i\cap J_j\), поэтому

\[
v(I_i)=\sum_j v(I_i\cap J_j).
\]

Для каждого фиксированного \(j\) аналогично

\[
v(J_j)=\sum_i v(I_i\cap J_j).
\]

Суммируя по \(i\) и \(j\), получаем

\[
\sum_i v(I_i)=\sum_{i,j}v(I_i\cap J_j)=\sum_j v(J_j).
\]

Здесь используется элементарный факт: объем прямоугольника равен сумме
объемов прямоугольников конечного разбиения. Он доказывается
перемножением длин координатных разбиений. \(\square\)

\subsubsection{Свойство 7. Мера Жордана конечно-аддитивна на
элементарных
множествах}\label{ux441ux432ux43eux439ux441ux442ux432ux43e-7.-ux43cux435ux440ux430-ux436ux43eux440ux434ux430ux43dux430-ux43aux43eux43dux435ux447ux43dux43e-ux430ux434ux434ux438ux442ux438ux432ux43dux430-ux43dux430-ux44dux43bux435ux43cux435ux43dux442ux430ux440ux43dux44bux445-ux43cux43dux43eux436ux435ux441ux442ux432ux430ux445}

Если \(A,B\) - непересекающиеся элементарные множества, то

\[
v(A\sqcup B)=v(A)+v(B).
\]

\subsubsection{Доказательство}\label{ux434ux43eux43aux430ux437ux430ux442ux435ux43bux44cux441ux442ux432ux43e-6}

Представим \(A\) и \(B\) как конечные дизъюнктные объединения
полуоткрытых прямоугольников. Тогда \(A\sqcup B\) уже представлено как
конечное дизъюнктное объединение всех этих прямоугольников. По
определению меры Жордана на элементарных множествах

\[
v(A\sqcup B)=\sum_i v(I_i)+\sum_j v(J_j)=v(A)+v(B).
\]

Для конечного числа попарно непересекающихся элементарных множеств
доказательство такое же. \(\square\)

\subsection{5. Примеры}\label{ux43fux440ux438ux43cux435ux440ux44b}

\subsubsection{\texorpdfstring{Пример 1. Минимальная и максимальная
\(\sigma\)-алгебры}{Пример 1. Минимальная и максимальная \textbackslash sigma-алгебры}}\label{ux43fux440ux438ux43cux435ux440-1.-ux43cux438ux43dux438ux43cux430ux43bux44cux43dux430ux44f-ux438-ux43cux430ux43aux441ux438ux43cux430ux43bux44cux43dux430ux44f-sigma-ux430ux43bux433ux435ux431ux440ux44b}

На любом \(\Omega\) класс \(\{\varnothing,\Omega\}\) - минимальная
\(\sigma\)-алгебра. Класс \(2^\Omega\) - максимальная
\(\sigma\)-алгебра.

\subsubsection{\texorpdfstring{Пример 2. Алгебра, которая не является
\(\sigma\)-алгеброй}{Пример 2. Алгебра, которая не является \textbackslash sigma-алгеброй}}\label{ux43fux440ux438ux43cux435ux440-2.-ux430ux43bux433ux435ux431ux440ux430-ux43aux43eux442ux43eux440ux430ux44f-ux43dux435-ux44fux432ux43bux44fux435ux442ux441ux44f-sigma-ux430ux43bux433ux435ux431ux440ux43eux439}

Пусть \(\Omega=\mathbb N\). Рассмотрим класс

\[
\mathcal A=\{A\subset\mathbb N:\ A \text{ конечно или } \overline{A} \text{ конечно}\}.
\]

Это алгебра: дополнение конечного множества ко-конечно, дополнение
ко-конечного множества конечно; конечные объединения конечных множеств
конечны, а объединение с ко-конечным множеством ко-конечно.

Но \(\mathcal A\) не является \(\sigma\)-алгеброй. Одноточечные
множества \(\{2\},\{4\},\{6\},\ldots\) конечны, значит лежат в
\(\mathcal A\). Их счетное объединение - множество четных чисел. Оно
бесконечно, и его дополнение тоже бесконечно. Значит множество четных
чисел не лежит в \(\mathcal A\).

\subsubsection{Пример 3. Полуалгебра полуинтервалов на
прямой}\label{ux43fux440ux438ux43cux435ux440-3.-ux43fux43eux43bux443ux430ux43bux433ux435ux431ux440ux430-ux43fux43eux43bux443ux438ux43dux442ux435ux440ux432ux430ux43bux43eux432-ux43dux430-ux43fux440ux44fux43cux43eux439}

Класс

\[
\mathcal S=\{[a,b):\ -\infty<a\le b<\infty\}
\]

является полуалгеброй на \(\mathbb R\) в локальном смысле. Пересечение
двух полуинтервалов снова полуинтервал:

\[
[a,b)\cap[c,d)=[\max(a,c),\min(b,d)),
\]

если правый конец больше левого; иначе это пустое множество. Разность
\([a,b)\setminus[c,d)\) является объединением не более чем двух
полуинтервалов:

\[
[a,b)\setminus[c,d)= [a,\min(b,c))\sqcup[\max(a,d),b),
\]

после удаления пустых частей.

\subsubsection{Пример 4. Мера
Дирака}\label{ux43fux440ux438ux43cux435ux440-4.-ux43cux435ux440ux430-ux434ux438ux440ux430ux43aux430}

Пусть \(\Omega=\mathbb R\), \(x=0\). Тогда

\[
\delta_0((-1,1))=1,\qquad \delta_0((1,2))=0.
\]

Если случайная величина почти наверное равна \(0\), ее распределение
есть \(\delta_0\).

\subsubsection{Пример 5. Мера Жордана на прямой и
плоскости}\label{ux43fux440ux438ux43cux435ux440-5.-ux43cux435ux440ux430-ux436ux43eux440ux434ux430ux43dux430-ux43dux430-ux43fux440ux44fux43cux43eux439-ux438-ux43fux43bux43eux441ux43aux43eux441ux442ux438}

На прямой

\[
v([a,b))=b-a.
\]

Если

\[
A=[0,1)\sqcup[2,5),
\]

то

\[
v(A)=1+3=4.
\]

В \(\mathbb R^2\) для прямоугольника \([a,b)\times[c,d)\):

\[
v([a,b)\times[c,d))=(b-a)(d-c).
\]

\subsection{6. Доказательства и
обоснования}\label{ux434ux43eux43aux430ux437ux430ux442ux435ux43bux44cux441ux442ux432ux430-ux438-ux43eux431ux43eux441ux43dux43eux432ux430ux43dux438ux44f}

Для устного экзамена удобно держать три доказательные линии.

\textbf{Первая линия: операции над системами множеств.}

Все свойства алгебр и \(\sigma\)-алгебр доказываются через дополнения и
формулы де Моргана:

\[
\overline{\left(\bigcup_n A_n\right)}=\bigcap_n \overline{A_n},\qquad
\overline{\left(\bigcap_n A_n\right)}=\bigcup_n \overline{A_n}.
\]

Если система замкнута относительно дополнения и объединения, она
автоматически замкнута относительно пересечения. Если объединение
разрешено счетное, то и пересечение разрешено счетное.

\textbf{Вторая линия: полуалгебра как удобный класс разбиений.}

Полуалгебра не обязана быть замкнута относительно объединения. Но она
устроена так, что при вычитании одного простого множества из другого
получается конечное число простых множеств. Поэтому из полуалгебры можно
построить алгебру конечными дизъюнктными объединениями.

\textbf{Третья линия: меры Дирака и Жордана.}

Для Дирака доказательство аддитивности сводится к вопросу: где находится
точка \(x\)? В дизъюнктном объединении она может лежать только в одном
слагаемом.

Для Жордана главное - корректность определения. Если одно и то же
элементарное множество разбили на прямоугольники двумя способами, нужно
перейти к общему измельчению всеми попарными пересечениями
прямоугольников. Тогда обе суммы объемов сводятся к одной и той же
двойной сумме.

\subsection{7. Что написать на
доске}\label{ux447ux442ux43e-ux43dux430ux43fux438ux441ux430ux442ux44c-ux43dux430-ux434ux43eux441ux43aux435}

\begin{enumerate}
\def\labelenumi{\arabic{enumi}.}
\tightlist
\item
  Цепочка систем:
\end{enumerate}

\[
\text{полуалгебра}\longrightarrow \text{алгебра конечных дизъюнктных объединений}\longrightarrow \sigma\text{-алгебра}.
\]

\begin{enumerate}
\def\labelenumi{\arabic{enumi}.}
\setcounter{enumi}{1}
\tightlist
\item
  Определения:
\end{enumerate}

\[
A,B\in\mathcal A\Rightarrow \overline{A},A\cup B\in\mathcal A,
\]

\[
A_n\in\mathcal F\Rightarrow \bigcup_{n=1}^{\infty}A_n\in\mathcal F.
\]

\begin{enumerate}
\def\labelenumi{\arabic{enumi}.}
\setcounter{enumi}{2}
\tightlist
\item
  Полуалгебра:
\end{enumerate}

\[
A,B\in\mathcal S\Rightarrow A\cap B\in\mathcal S,\qquad
A\setminus B=\bigsqcup_{k=1}^n C_k,\ C_k\in\mathcal S.
\]

\begin{enumerate}
\def\labelenumi{\arabic{enumi}.}
\setcounter{enumi}{3}
\tightlist
\item
  Дирак:
\end{enumerate}

\[
\delta_x(A)=\mathbf 1_A(x).
\]

\begin{enumerate}
\def\labelenumi{\arabic{enumi}.}
\setcounter{enumi}{4}
\tightlist
\item
  Жордан:
\end{enumerate}

\[
v\left(\prod_{k=1}^d[a_k,b_k)\right)=\prod_{k=1}^d(b_k-a_k),\qquad
v\left(\bigsqcup_{j=1}^m I_j\right)=\sum_{j=1}^m v(I_j).
\]

\subsection{8. Типичные дополнительные
вопросы}\label{ux442ux438ux43fux438ux447ux43dux44bux435-ux434ux43eux43fux43eux43bux43dux438ux442ux435ux43bux44cux43dux44bux435-ux432ux43eux43fux440ux43eux441ux44b}

\textbf{Вопрос. Почему в определении алгебры достаточно требовать
замкнутость относительно дополнения и объединения, но не пересечения?}

Потому что пересечение выражается через дополнение и объединение:

\[
A\cap B=\overline{\overline{A}\cup \overline{B}}.
\]

\textbf{Вопрос. Чем алгебра отличается от \(\sigma\)-алгебры?}

Алгебра замкнута относительно конечных объединений, а \(\sigma\)-алгебра
- относительно счетных объединений. Например, конечные и ко-конечные
подмножества \(\mathbb N\) образуют алгебру, но не \(\sigma\)-алгебру.

\textbf{Вопрос. Почему полуинтервалы \([a,b)\) берут именно
полуоткрытыми?}

Потому что они хорошо дизъюнктно склеиваются без пересечения границ:

\[
[a,c)=[a,b)\sqcup[b,c).
\]

Это удобно для задания длины и дальнейшего продолжения меры.

\textbf{Вопрос. Является ли полуалгебра алгеброй?}

Не обязательно. Полуинтервалы \([a,b)\) не замкнуты относительно
объединения: \([0,1)\cup[2,3)\) не является одним полуинтервалом. Но
конечные дизъюнктные объединения таких полуинтервалов уже образуют
алгебру элементарных множеств.

\textbf{Вопрос. Почему мера Дирака является вероятностной мерой?}

Потому что \(\delta_x(\Omega)=1\), \(\delta_x(A)\ge0\), и она
счетно-аддитивна на дизъюнктных объединениях.

\textbf{Вопрос. Почему мера Жордана не решает всю теорию меры?}

Потому что она хорошо работает для элементарных и жорданово измеримых
множеств, но не образует достаточно богатой \(\sigma\)-аддитивной теории
на всех множествах, которые возникают как пределы. Для этого строят меру
Лебега.

\subsection{9. Типичные
ошибки}\label{ux442ux438ux43fux438ux447ux43dux44bux435-ux43eux448ux438ux431ux43aux438}

\begin{itemize}
\tightlist
\item
  Говорить, что полуалгебра замкнута относительно объединений. Это
  неверно в общем случае.
\item
  Забывать, что алгебра требует только конечные операции, а
  \(\sigma\)-алгебра - счетные.
\item
  Путать \(\varnothing\in\mathcal S\) и \(\Omega\in\mathcal S\): для
  полуалгебры обязательно нужно пустое множество, а наличие \(\Omega\)
  зависит от принятого варианта определения.
\item
  Считать, что мера Жордана задана на всех подмножествах
  \(\mathbb R^d\). Она изначально задана на элементарных множествах, а
  затем на жорданово измеримых множествах.
\item
  Не проверять корректность определения \(v(A)=\sum_jv(I_j)\): одно и то
  же элементарное множество может иметь разные разбиения.
\item
  Называть любое неотрицательное отображение мерой. Нужна хотя бы
  конечная аддитивность, а в вероятностной теории дальше потребуется
  счетная аддитивность.
\end{itemize}

\subsection{10. Связи с другими
билетами}\label{ux441ux432ux44fux437ux438-ux441-ux434ux440ux443ux433ux438ux43cux438-ux431ux438ux43bux435ux442ux430ux43cux438}

\begin{itemize}
\tightlist
\item
  Билет 02: счетная аддитивность и вероятностные пространства используют
  \(\sigma\)-алгебру как систему событий.
\item
  Билет 03: продолжение меры начинается с меры на полуалгебре и строит
  меру на алгебре и \(\sigma\)-алгебре.
\item
  Билет 04: борелевская \(\sigma\)-алгебра - главный пример
  \(\sigma\)-алгебры в метрических пространствах.
\item
  Билет 05: измеримость случайной величины означает, что прообразы
  борелевских множеств лежат в \(\mathcal F\).
\item
  Билет 06: распределение случайной величины - это мера на пространстве
  значений.
\end{itemize}

\subsection{11. Минимум на
удовлетворительно}\label{ux43cux438ux43dux438ux43cux443ux43c-ux43dux430-ux443ux434ux43eux432ux43bux435ux442ux432ux43eux440ux438ux442ux435ux43bux44cux43dux43e}

Нужно уметь строго сформулировать:

\begin{enumerate}
\def\labelenumi{\arabic{enumi}.}
\tightlist
\item
  определение алгебры;
\item
  определение \(\sigma\)-алгебры;
\item
  определение полуалгебры;
\item
  формулу меры Дирака;
\item
  формулу объема прямоугольника для меры Жордана.
\end{enumerate}

Минимальный ответ:

\begin{quote}
Алгебра замкнута относительно конечных операций над множествами,
\(\sigma\)-алгебра - относительно счетных объединений, полуалгебра
замкнута относительно пересечений, а разности в ней раскладываются в
конечные дизъюнктные объединения. Мера Дирака задается формулой
\(\delta_x(A)=\mathbf 1_A(x)\), мера Жордана на прямоугольнике -
формулой \(v(\prod[a_k,b_k))=\prod(b_k-a_k)\).
\end{quote}

\subsection{12. Минимум на
хорошо}\label{ux43cux438ux43dux438ux43cux443ux43c-ux43dux430-ux445ux43eux440ux43eux448ux43e}

Кроме минимума нужно доказать:

\begin{enumerate}
\def\labelenumi{\arabic{enumi}.}
\tightlist
\item
  что алгебра замкнута относительно пересечений и разностей;
\item
  что \(\sigma\)-алгебра замкнута относительно счетных пересечений;
\item
  что полуинтервалы \([a,b)\) образуют полуалгебру;
\item
  что мера Дирака аддитивна.
\end{enumerate}

Также нужно привести пример алгебры, не являющейся \(\sigma\)-алгеброй:
конечные и ко-конечные подмножества \(\mathbb N\).

\subsection{13. Что добавить для
отлично}\label{ux447ux442ux43e-ux434ux43eux431ux430ux432ux438ux442ux44c-ux434ux43bux44f-ux43eux442ux43bux438ux447ux43dux43e}

Для отличного ответа нужно объяснить роль полуалгебры в построении меры:

\[
\mathcal S\longrightarrow \mathcal A(\mathcal S)\longrightarrow \sigma(\mathcal S).
\]

Также стоит доказать корректность меры Жордана через общее измельчение
двух разбиений. Это показывает, что формула

\[
v(A)=\sum_j v(I_j)
\]

не зависит от выбранного представления элементарного множества \(A\).

Полезно добавить предупреждение: мера Жордана - предшественник меры
Лебега, но она недостаточно устойчива для предельных операций, поэтому в
вероятности и анализе переходят к \(\sigma\)-аддитивной мере Лебега.

\subsection{14. Устная версия ответа на 5
минут}\label{ux443ux441ux442ux43dux430ux44f-ux432ux435ux440ux441ux438ux44f-ux43eux442ux432ux435ux442ux430-ux43dux430-5-ux43cux438ux43dux443ux442}

Сначала я фиксирую пространство \(\Omega\) и говорю, что события - это
не обязательно все подмножества \(\Omega\), а элементы некоторой системы
множеств. Для конечных операций достаточно алгебры:
\(\Omega\in\mathcal A\), вместе с \(A\) лежит \(\overline{A}\), вместе с
\(A,B\) лежит \(A\cup B\). Через де Моргана сразу получаю
\(A\cap B\in\mathcal A\) и \(A\setminus B\in\mathcal A\).

Дальше усиливаю определение: \(\sigma\)-алгебра - это алгебра, замкнутая
относительно счетных объединений. Тогда она замкнута и относительно
счетных пересечений, потому что

\[
\bigcap_n A_n=\overline{\left(\bigcup_n \overline{A_n}\right)}.
\]

После этого объясняю полуалгебру. Это система простых множеств, где
пересечение остается простым, а разность раскладывается в конечное
дизъюнктное объединение простых множеств. Главный пример - полуинтервалы
\([a,b)\) на прямой и полуоткрытые прямоугольники в \(\mathbb R^d\).

Затем перехожу к мерам. Мера Дирака в точке \(x\) задается формулой
\(\delta_x(A)=\mathbf 1_A(x)\). Она аддитивна, потому что в дизъюнктном
объединении точка \(x\) может лежать максимум в одном слагаемом.

Мера Жордана - это длина, площадь или объем элементарных множеств. На
прямоугольнике \(I=\prod[a_k,b_k)\) она равна \(v(I)=\prod(b_k-a_k)\).
На конечном дизъюнктном объединении прямоугольников она равна сумме
объемов. Корректность доказывается общим измельчением двух разбиений.

В конце связываю билет со следующими: в билете 2 нужна счетная
аддитивность, в билете 3 - продолжение меры с полуалгебры, в билете 4 -
борелевские \(\sigma\)-алгебры.

\subsection{15. Если осталось 2 минуты перед
экзаменом}\label{ux435ux441ux43bux438-ux43eux441ux442ux430ux43bux43eux441ux44c-2-ux43cux438ux43dux443ux442ux44b-ux43fux435ux440ux435ux434-ux44dux43aux437ux430ux43cux435ux43dux43eux43c}

Запомнить пять формул:

\[
A\cap B=\overline{\overline{A}\cup \overline{B}},
\]

\[
\bigcap_n A_n=\overline{\left(\bigcup_n \overline{A_n}\right)},
\]

\[
A\setminus B=\bigsqcup_{k=1}^n C_k\quad\text{для полуалгебры},
\]

\[
\delta_x(A)=\mathbf 1_A(x),
\]

\[
v\left(\prod_{k=1}^d[a_k,b_k)\right)=\prod_{k=1}^d(b_k-a_k).
\]

Главная ловушка: алгебра - конечные операции, \(\sigma\)-алгебра -
счетные операции, полуалгебра - не алгебра и не обязана быть замкнута по
объединениям.

\newpage

\section{Билет 02. Счетная аддитивность и вероятностные пространства.
Мера, порожденная функцией
распределения}\label{ux431ux438ux43bux435ux442-02.-ux441ux447ux435ux442ux43dux430ux44f-ux430ux434ux434ux438ux442ux438ux432ux43dux43eux441ux442ux44c-ux438-ux432ux435ux440ux43eux44fux442ux43dux43eux441ux442ux43dux44bux435-ux43fux440ux43eux441ux442ux440ux430ux43dux441ux442ux432ux430.-ux43cux435ux440ux430-ux43fux43eux440ux43eux436ux434ux435ux43dux43dux430ux44f-ux444ux443ux43dux43aux446ux438ux435ux439-ux440ux430ux441ux43fux440ux435ux434ux435ux43bux435ux43dux438ux44f}

Источники: Ширяев 1, гл. II; лекции ДСП.

\subsection{1. Короткий
ответ}\label{ux43aux43eux440ux43eux442ux43aux438ux439-ux43eux442ux432ux435ux442-1}

Вероятностное пространство в аксиоматике Колмогорова - это тройка

\[
(\Omega,\mathcal F,\mathbb P),
\]

где \(\Omega\) - пространство элементарных исходов, \(\mathcal F\) -
\(\sigma\)-алгебра событий, а \(\mathbb P\) - счетно-аддитивная мера на
\(\mathcal F\) с полной массой \(1\).

Главная аксиома:

\[
A_i\cap A_j=\varnothing,\ i\ne j
\quad\Longrightarrow\quad
\mathbb P\left(\bigsqcup_{n=1}^{\infty} A_n\right)
=
\sum_{n=1}^{\infty}\mathbb P(A_n).
\]

Из нее следуют \(\mathbb P(\varnothing)=0\), конечная аддитивность,
формула дополнения, монотонность, счетная субаддитивность и
непрерывность вероятности снизу и сверху.

Если \(F:\mathbb R\to[0,1]\) не убывает, непрерывна справа и

\[
\lim_{x\to-\infty}F(x)=0,\qquad
\lim_{x\to+\infty}F(x)=1,
\]

то \(F\) порождает единственную вероятностную меру \(\mu_F\) на
борелевской \(\sigma\)-алгебре \(\mathcal B(\mathbb R)\) такую, что

\[
\mu_F((a,b])=F(b)-F(a),\qquad a<b.
\]

Это мера Лебега-Стилтьеса, соответствующая функции распределения \(F\).
В этом билете существование такой меры обычно формулируется без
доказательства; доказательство опирается на теорему продолжения меры.

\subsection{2. Полный
ответ}\label{ux43fux43eux43bux43dux44bux439-ux43eux442ux432ux435ux442-1}

\subsubsection{2.1. Аксиоматический подход
Колмогорова}\label{ux430ux43aux441ux438ux43eux43cux430ux442ux438ux447ux435ux441ux43aux438ux439-ux43fux43eux434ux445ux43eux434-ux43aux43eux43bux43cux43eux433ux43eux440ux43eux432ux430}

В первом билете выбиралась система множеств, на которой можно говорить о
событиях. Во втором билете на этой системе задается вероятность. Идея
Колмогорова состоит в том, что вероятность является обычной мерой, но
нормированной условием полной массы:

\[
\mathbb P(\Omega)=1.
\]

Пусть \(\Omega\) - множество всех элементарных исходов эксперимента.
Например, для одного броска монеты \(\Omega=\{H,T\}\), для бесконечной
последовательности бросков \(\Omega=\{H,T\}^{\mathbb N}\), для
вещественной случайной величины естественно брать \(\Omega=\mathbb R\).
Не каждое подмножество \(\Omega\) автоматически объявляется событием.
Выбирается \(\sigma\)-алгебра \(\mathcal F\subset 2^\Omega\), то есть
класс подмножеств, замкнутый относительно дополнений и счетных
объединений. Элементы \(\mathcal F\) называются событиями.

\subsubsection{2.2. Определение и аксиомы
вероятности}\label{ux43eux43fux440ux435ux434ux435ux43bux435ux43dux438ux435-ux438-ux430ux43aux441ux438ux43eux43cux44b-ux432ux435ux440ux43eux44fux442ux43dux43eux441ux442ux438}

Вероятность - это функция

\[
\mathbb P:\mathcal F\to[0,1],
\]

которая удовлетворяет трем требованиям:

\begin{enumerate}
\def\labelenumi{\arabic{enumi}.}
\tightlist
\item
  неотрицательность: \(\mathbb P(A)\ge 0\) для всех \(A\in\mathcal F\);
\item
  нормировка: \(\mathbb P(\Omega)=1\);
\item
  счетная аддитивность: если \(A_1,A_2,\ldots\in\mathcal F\) попарно не
  пересекаются, то
\end{enumerate}

\[
\mathbb P\left(\bigcup_{n=1}^{\infty}A_n\right)
=
\sum_{n=1}^{\infty}\mathbb P(A_n).
\]

Так как множества \(A_n\) попарно не пересекаются, объединение часто
обозначают дизъюнктным объединением:

\[
\bigcup_{n=1}^{\infty}A_n=\bigsqcup_{n=1}^{\infty}A_n.
\]

\subsubsection{2.3. Счетная и конечная
аддитивность}\label{ux441ux447ux435ux442ux43dux430ux44f-ux438-ux43aux43eux43dux435ux447ux43dux430ux44f-ux430ux434ux434ux438ux442ux438ux432ux43dux43eux441ux442ux44c}

Счетная аддитивность является принципиальным усилением конечной
аддитивности. Конечная аддитивность говорит только о конечных
разбиениях:

\[
\mathbb P\left(\bigsqcup_{k=1}^m A_k\right)
=
\sum_{k=1}^m\mathbb P(A_k).
\]

Для вероятностной теории этого недостаточно, потому что основные события
часто являются счетными объединениями или пересечениями. Например,
событие ``хотя бы однажды произошел успех'' имеет вид

\[
\bigcup_{n=1}^{\infty}A_n,
\]

событие ``успехи происходят бесконечно часто'' имеет вид

\[
\limsup_{n\to\infty}A_n
=
\bigcap_{m=1}^{\infty}\bigcup_{n=m}^{\infty}A_n,
\]

а предельные теоремы работают именно с пределами последовательностей
событий и случайных величин. Поэтому в аксиоматике Колмогорова требуется
не просто конечная, а счетная аддитивность.

\subsubsection{2.4. Свойства
вероятности}\label{ux441ux432ux43eux439ux441ux442ux432ux430-ux432ux435ux440ux43eux44fux442ux43dux43eux441ux442ux438}

Из аксиом вероятности выводится набор стандартных свойств:
\(\mathbb P(\varnothing)=0\), формула дополнения, монотонность, формула
сложения для объединения, субаддитивность, непрерывность снизу и
непрерывность сверху. Эти свойства не являются дополнительными
аксиомами. Их нужно уметь быстро доказывать из счетной аддитивности.

\subsubsection{2.5. Вероятностная мера и функция
распределения}\label{ux432ux435ux440ux43eux44fux442ux43dux43eux441ux442ux43dux430ux44f-ux43cux435ux440ux430-ux438-ux444ux443ux43dux43aux446ux438ux44f-ux440ux430ux441ux43fux440ux435ux434ux435ux43bux435ux43dux438ux44f}

Вторая часть билета - пример вероятностной меры, заданной не через явное
перечисление исходов, а через функцию распределения. Пусть
\(F:\mathbb R\to[0,1]\) обладает свойствами:

\begin{enumerate}
\def\labelenumi{\arabic{enumi}.}
\tightlist
\item
  \(F\) не убывает;
\item
  \(F\) непрерывна справа;
\item
  \(\lim_{x\to-\infty}F(x)=0\);
\item
  \(\lim_{x\to+\infty}F(x)=1\).
\end{enumerate}

Тогда существует единственная вероятностная мера \(\mu_F\) на
\(\mathcal B(\mathbb R)\), для которой

\[
\mu_F((-\infty,x])=F(x)
\]

и, следовательно,

\[
\mu_F((a,b])=F(b)-F(a).
\]

\subsubsection{2.6. Мера Лебега-Стилтьеса и типы
распределений}\label{ux43cux435ux440ux430-ux43bux435ux431ux435ux433ux430-ux441ux442ux438ux43bux442ux44cux435ux441ux430-ux438-ux442ux438ux43fux44b-ux440ux430ux441ux43fux440ux435ux434ux435ux43bux435ux43dux438ux439}

Эта мера называется мерой Лебега-Стилтьеса, порожденной \(F\).
Интуитивно \(F(x)\) задает массу луча \((-\infty,x]\), а разность
\(F(b)-F(a)\) задает массу полуинтервала \((a,b]\). Если \(F\) имеет
скачок в точке \(a\), то в этой точке у меры есть атом:

\[
\mu_F(\{a\})=F(a)-F(a-),
\]

где

\[
F(a-)=\lim_{x\uparrow a}F(x).
\]

Если \(F\) непрерывна, атомов нет. Если \(F\) абсолютно непрерывна и
имеет плотность \(f\), то

\[
F(x)=\int_{-\infty}^{x} f(t)\,dt,\qquad
\mu_F((a,b])=\int_a^b f(t)\,dt.
\]

Таким образом, дискретные, непрерывные и смешанные распределения входят
в одну общую схему: каждое распределение на прямой задается своей
функцией распределения.

\subsection{3. Основные
определения}\label{ux43eux441ux43dux43eux432ux43dux44bux435-ux43eux43fux440ux435ux434ux435ux43bux435ux43dux438ux44f-1}

\subsubsection{Измеримое
пространство}\label{ux438ux437ux43cux435ux440ux438ux43cux43eux435-ux43fux440ux43eux441ux442ux440ux430ux43dux441ux442ux432ux43e}

Пара \((\Omega,\mathcal F)\), где \(\Omega\) - множество, а
\(\mathcal F\) - \(\sigma\)-алгебра подмножеств \(\Omega\).

\subsubsection{Элементарный
исход}\label{ux44dux43bux435ux43cux435ux43dux442ux430ux440ux43dux44bux439-ux438ux441ux445ux43eux434}

Элемент \(\omega\in\Omega\). Он описывает один возможный исход
эксперимента.

\subsubsection{Событие}\label{ux441ux43eux431ux44bux442ux438ux435}

Множество \(A\in\mathcal F\). Если в эксперименте реализовался исход
\(\omega\in A\), говорят, что событие \(A\) произошло.

\subsubsection{Вероятностная
мера}\label{ux432ux435ux440ux43eux44fux442ux43dux43eux441ux442ux43dux430ux44f-ux43cux435ux440ux430}

Функция \(\mathbb P:\mathcal F\to[0,1]\), удовлетворяющая условиям:

\[
\mathbb P(\Omega)=1,
\]

и для любых попарно непересекающихся \(A_1,A_2,\ldots\in\mathcal F\):

\[
\mathbb P\left(\bigsqcup_{n=1}^{\infty}A_n\right)
=
\sum_{n=1}^{\infty}\mathbb P(A_n).
\]

Неотрицательность обычно включается в определение меры. Так как значения
берутся в \([0,1]\), она уже зафиксирована.

\subsubsection{Вероятностное пространство, или абстрактная
колмогоровская
тройка}\label{ux432ux435ux440ux43eux44fux442ux43dux43eux441ux442ux43dux43eux435-ux43fux440ux43eux441ux442ux440ux430ux43dux441ux442ux432ux43e-ux438ux43bux438-ux430ux431ux441ux442ux440ux430ux43aux442ux43dux430ux44f-ux43aux43eux43bux43cux43eux433ux43eux440ux43eux432ux441ux43aux430ux44f-ux442ux440ux43eux439ux43aux430}

Тройка

\[
(\Omega,\mathcal F,\mathbb P),
\]

где \((\Omega,\mathcal F)\) - измеримое пространство, а \(\mathbb P\) -
вероятностная мера на \(\mathcal F\).

\subsubsection{Попарно несовместные
события}\label{ux43fux43eux43fux430ux440ux43dux43e-ux43dux435ux441ux43eux432ux43cux435ux441ux442ux43dux44bux435-ux441ux43eux431ux44bux442ux438ux44f}

Последовательность \((A_n)_{n\ge1}\) называется попарно несовместной,
если

\[
A_i\cap A_j=\varnothing,\qquad i\ne j.
\]

\subsubsection{Конечная
аддитивность}\label{ux43aux43eux43dux435ux447ux43dux430ux44f-ux430ux434ux434ux438ux442ux438ux432ux43dux43eux441ux442ux44c}

Для попарно несовместных \(A_1,\ldots,A_m\in\mathcal F\):

\[
\mathbb P\left(\bigsqcup_{k=1}^{m}A_k\right)
=
\sum_{k=1}^{m}\mathbb P(A_k).
\]

\subsubsection{Счетная
аддитивность}\label{ux441ux447ux435ux442ux43dux430ux44f-ux430ux434ux434ux438ux442ux438ux432ux43dux43eux441ux442ux44c}

Для попарно несовместных \(A_1,A_2,\ldots\in\mathcal F\):

\[
\mathbb P\left(\bigsqcup_{k=1}^{\infty}A_k\right)
=
\sum_{k=1}^{\infty}\mathbb P(A_k).
\]

\subsubsection{\texorpdfstring{Борелевская \(\sigma\)-алгебра на
\(\mathbb R\)}{Борелевская \textbackslash sigma-алгебра на \textbackslash mathbb R}}\label{ux431ux43eux440ux435ux43bux435ux432ux441ux43aux430ux44f-sigma-ux430ux43bux433ux435ux431ux440ux430-ux43dux430-mathbb-r}

Минимальная \(\sigma\)-алгебра, содержащая все открытые множества
прямой. Обозначается \(\mathcal B(\mathbb R)\).

Эквивалентно, \(\mathcal B(\mathbb R)\) порождена интервалами вида
\((-\infty,x]\):

\[
\mathcal B(\mathbb R)=\sigma\{(-\infty,x]:x\in\mathbb R\}.
\]

\subsubsection{\texorpdfstring{Функция распределения меры \(\mu\) на
\(\mathbb R\)}{Функция распределения меры \textbackslash mu на \textbackslash mathbb R}}\label{ux444ux443ux43dux43aux446ux438ux44f-ux440ux430ux441ux43fux440ux435ux434ux435ux43bux435ux43dux438ux44f-ux43cux435ux440ux44b-mu-ux43dux430-mathbb-r}

\[
F_\mu(x)=\mu((-\infty,x]).
\]

\subsubsection{\texorpdfstring{Функция распределения случайной величины
\(X\)}{Функция распределения случайной величины X}}\label{ux444ux443ux43dux43aux446ux438ux44f-ux440ux430ux441ux43fux440ux435ux434ux435ux43bux435ux43dux438ux44f-ux441ux43bux443ux447ux430ux439ux43dux43eux439-ux432ux435ux43bux438ux447ux438ux43dux44b-x}

\[
F_X(x)=\mathbb P(X\le x)
=
\mathbb P(\{\omega:X(\omega)\le x\}).
\]

\subsubsection{Абстрактная функция
распределения}\label{ux430ux431ux441ux442ux440ux430ux43aux442ux43dux430ux44f-ux444ux443ux43dux43aux446ux438ux44f-ux440ux430ux441ux43fux440ux435ux434ux435ux43bux435ux43dux438ux44f}

Функция \(F:\mathbb R\to[0,1]\), которая не убывает, непрерывна справа и
имеет пределы \(0\) на \(-\infty\) и \(1\) на \(+\infty\).

\subsubsection{Мера, порожденная функцией
распределения}\label{ux43cux435ux440ux430-ux43fux43eux440ux43eux436ux434ux435ux43dux43dux430ux44f-ux444ux443ux43dux43aux446ux438ux435ux439-ux440ux430ux441ux43fux440ux435ux434ux435ux43bux435ux43dux438ux44f}

Единственная вероятностная мера \(\mu_F\) на \(\mathcal B(\mathbb R)\),
для которой

\[
\mu_F((-\infty,x])=F(x)
\]

для всех \(x\in\mathbb R\). Тогда для \(a<b\):

\[
\mu_F((a,b])=F(b)-F(a).
\]

\subsection{4. Основные теоремы и
свойства}\label{ux43eux441ux43dux43eux432ux43dux44bux435-ux442ux435ux43eux440ux435ux43cux44b-ux438-ux441ux432ux43eux439ux441ux442ux432ux430-1}

\subsubsection{Свойство 1. Вероятность пустого множества равна
нулю}\label{ux441ux432ux43eux439ux441ux442ux432ux43e-1.-ux432ux435ux440ux43eux44fux442ux43dux43eux441ux442ux44c-ux43fux443ux441ux442ux43eux433ux43e-ux43cux43dux43eux436ux435ux441ux442ux432ux430-ux440ux430ux432ux43dux430-ux43dux443ux43bux44e}

\[
\mathbb P(\varnothing)=0.
\]

\subsubsection{Доказательство}\label{ux434ux43eux43aux430ux437ux430ux442ux435ux43bux44cux441ux442ux432ux43e-7}

Используем дизъюнктное разложение

\[
\Omega=\Omega\sqcup\varnothing\sqcup\varnothing\sqcup\cdots.
\]

По счетной аддитивности

\[
1=\mathbb P(\Omega)
=
\mathbb P(\Omega)+\sum_{n=2}^{\infty}\mathbb P(\varnothing)
=
1+\sum_{n=2}^{\infty}\mathbb P(\varnothing).
\]

Следовательно,

\[
\sum_{n=2}^{\infty}\mathbb P(\varnothing)=0.
\]

Слагаемые неотрицательны, значит \(\mathbb P(\varnothing)=0\).
\(\square\)

\subsubsection{Свойство 2. Из счетной аддитивности следует конечная
аддитивность}\label{ux441ux432ux43eux439ux441ux442ux432ux43e-2.-ux438ux437-ux441ux447ux435ux442ux43dux43eux439-ux430ux434ux434ux438ux442ux438ux432ux43dux43eux441ux442ux438-ux441ux43bux435ux434ux443ux435ux442-ux43aux43eux43dux435ux447ux43dux430ux44f-ux430ux434ux434ux438ux442ux438ux432ux43dux43eux441ux442ux44c}

Если \(A_1,\ldots,A_m\) попарно не пересекаются, то

\[
\mathbb P\left(\bigsqcup_{k=1}^{m}A_k\right)
=
\sum_{k=1}^{m}\mathbb P(A_k).
\]

\subsubsection{Доказательство}\label{ux434ux43eux43aux430ux437ux430ux442ux435ux43bux44cux441ux442ux432ux43e-8}

Возьмем счетную последовательность

\[
A_1,\ldots,A_m,\varnothing,\varnothing,\ldots
\]

и применим счетную аддитивность:

\[
\mathbb P\left(\bigsqcup_{k=1}^{m}A_k\right)
=
\sum_{k=1}^{m}\mathbb P(A_k)+\sum_{k=m+1}^{\infty}\mathbb P(\varnothing).
\]

По свойству 1 \(\mathbb P(\varnothing)=0\), поэтому

\[
\mathbb P\left(\bigsqcup_{k=1}^{m}A_k\right)
=
\sum_{k=1}^{m}\mathbb P(A_k).
\]

\(\square\)

\subsubsection{Свойство 3. Формула
дополнения}\label{ux441ux432ux43eux439ux441ux442ux432ux43e-3.-ux444ux43eux440ux43cux443ux43bux430-ux434ux43eux43fux43eux43bux43dux435ux43dux438ux44f}

Для любого \(A\in\mathcal F\):

\[
\mathbb P(\overline{A})=1-\mathbb P(A).
\]

\subsubsection{Доказательство}\label{ux434ux43eux43aux430ux437ux430ux442ux435ux43bux44cux441ux442ux432ux43e-9}

Множества \(A\) и \(\overline{A}\) несовместны, причем

\[
A\sqcup \overline{A}=\Omega.
\]

По конечной аддитивности:

\[
1=\mathbb P(\Omega)=\mathbb P(A)+\mathbb P(\overline{A}).
\]

Отсюда

\[
\mathbb P(\overline{A})=1-\mathbb P(A).
\]

\(\square\)

\subsubsection{Свойство 4. Монотонность и формула
разности}\label{ux441ux432ux43eux439ux441ux442ux432ux43e-4.-ux43cux43eux43dux43eux442ux43eux43dux43dux43eux441ux442ux44c-ux438-ux444ux43eux440ux43cux443ux43bux430-ux440ux430ux437ux43dux43eux441ux442ux438}

Если \(A\subset B\), то

\[
\mathbb P(B\setminus A)=\mathbb P(B)-\mathbb P(A),
\qquad
\mathbb P(A)\le \mathbb P(B).
\]

\subsubsection{Доказательство}\label{ux434ux43eux43aux430ux437ux430ux442ux435ux43bux44cux441ux442ux432ux43e-10}

При \(A\subset B\) имеем дизъюнктное разложение

\[
B=A\sqcup(B\setminus A).
\]

Значит

\[
\mathbb P(B)=\mathbb P(A)+\mathbb P(B\setminus A).
\]

Отсюда следует формула разности. Так как
\(\mathbb P(B\setminus A)\ge0\), получаем монотонность:

\[
\mathbb P(A)\le\mathbb P(B).
\]

\(\square\)

\subsubsection{Свойство 5. Формула сложения для двух
событий}\label{ux441ux432ux43eux439ux441ux442ux432ux43e-5.-ux444ux43eux440ux43cux443ux43bux430-ux441ux43bux43eux436ux435ux43dux438ux44f-ux434ux43bux44f-ux434ux432ux443ux445-ux441ux43eux431ux44bux442ux438ux439}

Для любых \(A,B\in\mathcal F\):

\[
\mathbb P(A\cup B)=\mathbb P(A)+\mathbb P(B)-\mathbb P(A\cap B).
\]

\subsubsection{Доказательство}\label{ux434ux43eux43aux430ux437ux430ux442ux435ux43bux44cux441ux442ux432ux43e-11}

Разложим \(A\cup B\) на попарно непересекающиеся части:

\[
A\cup B=(A\setminus B)\sqcup(B\setminus A)\sqcup(A\cap B).
\]

Также

\[
A=(A\setminus B)\sqcup(A\cap B),
\qquad
B=(B\setminus A)\sqcup(A\cap B).
\]

По конечной аддитивности:

\[
\mathbb P(A\cup B)
=
\mathbb P(A\setminus B)+\mathbb P(B\setminus A)+\mathbb P(A\cap B),
\]

а

\[
\mathbb P(A)+\mathbb P(B)
=
\mathbb P(A\setminus B)+\mathbb P(B\setminus A)+2\mathbb P(A\cap B).
\]

Вычитая один раз \(\mathbb P(A\cap B)\), получаем нужную формулу.
\(\square\)

\subsubsection{Свойство 6. Счетная
субаддитивность}\label{ux441ux432ux43eux439ux441ux442ux432ux43e-6.-ux441ux447ux435ux442ux43dux430ux44f-ux441ux443ux431ux430ux434ux434ux438ux442ux438ux432ux43dux43eux441ux442ux44c}

Для любых событий \(A_1,A_2,\ldots\in\mathcal F\):

\[
\mathbb P\left(\bigcup_{n=1}^{\infty}A_n\right)
\le
\sum_{n=1}^{\infty}\mathbb P(A_n).
\]

\subsubsection{Доказательство}\label{ux434ux43eux43aux430ux437ux430ux442ux435ux43bux44cux441ux442ux432ux43e-12}

Преобразуем произвольную последовательность событий в попарно
непересекающуюся последовательность:

\[
B_1=A_1,\qquad
B_n=A_n\setminus\bigcup_{k=1}^{n-1}A_k,\quad n\ge2.
\]

Тогда \(B_n\subset A_n\), множества \(B_n\) попарно не пересекаются и

\[
\bigcup_{n=1}^{\infty}A_n
=
\bigsqcup_{n=1}^{\infty}B_n.
\]

По счетной аддитивности и монотонности:

\[
\mathbb P\left(\bigcup_{n=1}^{\infty}A_n\right)
=
\sum_{n=1}^{\infty}\mathbb P(B_n)
\le
\sum_{n=1}^{\infty}\mathbb P(A_n).
\]

\(\square\)

\subsubsection{Свойство 7. Непрерывность вероятности
снизу}\label{ux441ux432ux43eux439ux441ux442ux432ux43e-7.-ux43dux435ux43fux440ux435ux440ux44bux432ux43dux43eux441ux442ux44c-ux432ux435ux440ux43eux44fux442ux43dux43eux441ux442ux438-ux441ux43dux438ux437ux443}

Если

\[
A_1\subset A_2\subset\cdots,
\qquad
A=\bigcup_{n=1}^{\infty}A_n,
\]

то

\[
\mathbb P(A_n)\uparrow\mathbb P(A),
\qquad
\lim_{n\to\infty}\mathbb P(A_n)=\mathbb P(A).
\]

\subsubsection{Доказательство}\label{ux434ux43eux43aux430ux437ux430ux442ux435ux43bux44cux441ux442ux432ux43e-13}

Представим возрастающую последовательность через непересекающиеся слои:

\[
B_1=A_1,\qquad
B_n=A_n\setminus A_{n-1},\quad n\ge2.
\]

Тогда

\[
A_N=\bigsqcup_{n=1}^{N}B_n,
\qquad
A=\bigsqcup_{n=1}^{\infty}B_n.
\]

Следовательно,

\[
\mathbb P(A_N)=\sum_{n=1}^{N}\mathbb P(B_n),
\qquad
\mathbb P(A)=\sum_{n=1}^{\infty}\mathbb P(B_n).
\]

Переходя к пределу по \(N\), получаем

\[
\lim_{N\to\infty}\mathbb P(A_N)
=
\mathbb P(A).
\]

\(\square\)

\subsubsection{Свойство 8. Непрерывность вероятности
сверху}\label{ux441ux432ux43eux439ux441ux442ux432ux43e-8.-ux43dux435ux43fux440ux435ux440ux44bux432ux43dux43eux441ux442ux44c-ux432ux435ux440ux43eux44fux442ux43dux43eux441ux442ux438-ux441ux432ux435ux440ux445ux443}

Если

\[
A_1\supset A_2\supset\cdots,
\qquad
A=\bigcap_{n=1}^{\infty}A_n,
\]

то

\[
\mathbb P(A_n)\downarrow\mathbb P(A),
\qquad
\lim_{n\to\infty}\mathbb P(A_n)=\mathbb P(A).
\]

\subsubsection{Доказательство}\label{ux434ux43eux43aux430ux437ux430ux442ux435ux43bux44cux441ux442ux432ux43e-14}

Так как \(A_n\downarrow A\), дополнения возрастают:

\[
\overline{A_1}\subset \overline{A_2}\subset\cdots,
\qquad
\bigcup_{n=1}^{\infty}\overline{A_n}=\overline{A}.
\]

По непрерывности снизу:

\[
\mathbb P(\overline{A_n})\uparrow\mathbb P(\overline{A}).
\]

Используя формулу дополнения, получаем

\[
\mathbb P(A_n)=1-\mathbb P(\overline{A_n})
\downarrow
1-\mathbb P(\overline{A})
=
\mathbb P(A).
\]

\(\square\)

Для общих мер в этом утверждении нужно условие конечности хотя бы
\(\mu(A_1)<\infty\). Для вероятностной меры оно выполнено автоматически,
потому что \(\mathbb P(A_1)\le1\).

\subsubsection{Теорема 9. Функция распределения порождает
меру}\label{ux442ux435ux43eux440ux435ux43cux430-9.-ux444ux443ux43dux43aux446ux438ux44f-ux440ux430ux441ux43fux440ux435ux434ux435ux43bux435ux43dux438ux44f-ux43fux43eux440ux43eux436ux434ux430ux435ux442-ux43cux435ux440ux443}

Пусть \(F:\mathbb R\to[0,1]\) не убывает, непрерывна справа и

\[
\lim_{x\to-\infty}F(x)=0,\qquad
\lim_{x\to+\infty}F(x)=1.
\]

Тогда существует единственная вероятностная мера \(\mu_F\) на
\(\mathcal B(\mathbb R)\) такая, что

\[
\mu_F((-\infty,x])=F(x)
\]

для всех \(x\in\mathbb R\). Эквивалентно,

\[
\mu_F((a,b])=F(b)-F(a),\qquad a<b.
\]

\subsubsection{Комментарий к
доказательству}\label{ux43aux43eux43cux43cux435ux43dux442ux430ux440ux438ux439-ux43a-ux434ux43eux43aux430ux437ux430ux442ux435ux43bux44cux441ux442ux432ux443}

По программе здесь доказательство не требуется. Идея такая: сначала
задают массу полуинтервалов формулой

\[
\mu_F((a,b])=F(b)-F(a),
\]

затем задают массу конечных дизъюнктных объединений таких полуинтервалов
как сумму масс. Монотонность \(F\) гарантирует неотрицательность,
телескопические суммы дают конечную аддитивность, а непрерывность справа
нужна для счетной аддитивности на исходной алгебре. После этого
применяется теорема продолжения меры с алгебры или полуалгебры на
порожденную \(\sigma\)-алгебру. Поскольку полуинтервалы порождают
\(\mathcal B(\mathbb R)\), получается мера на борелевских множествах.
Единственность следует из того, что значения на порождающем классе
фиксированы.

\subsection{5. Примеры}\label{ux43fux440ux438ux43cux435ux440ux44b-1}

\subsubsection{Пример 1. Один бросок
монеты}\label{ux43fux440ux438ux43cux435ux440-1.-ux43eux434ux438ux43d-ux431ux440ux43eux441ux43eux43a-ux43cux43eux43dux435ux442ux44b}

Пусть

\[
\Omega=\{H,T\},\qquad
\mathcal F=2^\Omega.
\]

Для симметричной монеты:

\[
\mathbb P(\{H\})=\mathbb P(\{T\})=\frac12.
\]

Тогда

\[
\mathbb P(\Omega)=1,\qquad
\mathbb P(\varnothing)=0.
\]

Счетная аддитивность здесь фактически сводится к конечной, потому что
пространство конечно.

\subsubsection{Пример 2. Счетное пространство
исходов}\label{ux43fux440ux438ux43cux435ux440-2.-ux441ux447ux435ux442ux43dux43eux435-ux43fux440ux43eux441ux442ux440ux430ux43dux441ux442ux432ux43e-ux438ux441ux445ux43eux434ux43eux432}

Пусть

\[
\Omega=\mathbb N,\qquad
\mathcal F=2^{\mathbb N},
\]

и задана последовательность \(p_n\ge0\) такая, что

\[
\sum_{n=1}^{\infty}p_n=1.
\]

Тогда

\[
\mathbb P(A)=\sum_{n\in A}p_n
\]

задает вероятность на \(2^{\mathbb N}\). В частности,

\[
\mathbb P(\{n\})=p_n.
\]

Это типичная дискретная вероятностная мера.

\subsubsection{Пример 3. Вырожденное распределение и мера
Дирака}\label{ux43fux440ux438ux43cux435ux440-3.-ux432ux44bux440ux43eux436ux434ux435ux43dux43dux43eux435-ux440ux430ux441ux43fux440ux435ux434ux435ux43bux435ux43dux438ux435-ux438-ux43cux435ux440ux430-ux434ux438ux440ux430ux43aux430}

Пусть \(a\in\mathbb R\) и

\[
F(x)=
\begin{cases}
0, & x<a,\\
1, & x\ge a.
\end{cases}
\]

Тогда \(\mu_F=\delta_a\), то есть вся масса сосредоточена в точке \(a\):

\[
\mu_F(A)=
\begin{cases}
1, & a\in A,\\
0, & a\notin A.
\end{cases}
\]

Скачок функции распределения равен

\[
F(a)-F(a-)=1-0=1,
\]

поэтому

\[
\mu_F(\{a\})=1.
\]

\subsubsection{\texorpdfstring{Пример 4. Распределение Бернулли на
\(\mathbb R\)}{Пример 4. Распределение Бернулли на \textbackslash mathbb R}}\label{ux43fux440ux438ux43cux435ux440-4.-ux440ux430ux441ux43fux440ux435ux434ux435ux43bux435ux43dux438ux435-ux431ux435ux440ux43dux443ux43bux43bux438-ux43dux430-mathbb-r}

Пусть случайная величина принимает значения \(0\) и \(1\):

\[
\mathbb P(X=1)=p,\qquad
\mathbb P(X=0)=1-p.
\]

Ее функция распределения:

\[
F(x)=
\begin{cases}
0, & x<0,\\
1-p, & 0\le x<1,\\
1, & x\ge1.
\end{cases}
\]

У меры два атома:

\[
\mu_F(\{0\})=1-p,\qquad
\mu_F(\{1\})=p.
\]

\subsubsection{\texorpdfstring{Пример 5. Равномерное распределение на
\([0,1]\)}{Пример 5. Равномерное распределение на {[}0,1{]}}}\label{ux43fux440ux438ux43cux435ux440-5.-ux440ux430ux432ux43dux43eux43cux435ux440ux43dux43eux435-ux440ux430ux441ux43fux440ux435ux434ux435ux43bux435ux43dux438ux435-ux43dux430-01}

Функция распределения:

\[
F(x)=
\begin{cases}
0, & x<0,\\
x, & 0\le x<1,\\
1, & x\ge1.
\end{cases}
\]

Тогда для \(0\le a<b\le1\):

\[
\mu_F((a,b])=F(b)-F(a)=b-a.
\]

В этом случае мера совпадает с длиной на отрезке \([0,1]\).

\subsubsection{Пример 6. Экспоненциальное
распределение}\label{ux43fux440ux438ux43cux435ux440-6.-ux44dux43aux441ux43fux43eux43dux435ux43dux446ux438ux430ux43bux44cux43dux43eux435-ux440ux430ux441ux43fux440ux435ux434ux435ux43bux435ux43dux438ux435}

При параметре \(\lambda>0\):

\[
F(x)=
\begin{cases}
0, & x<0,\\
1-e^{-\lambda x}, & x\ge0.
\end{cases}
\]

Тогда для \(0\le a<b\):

\[
\mu_F((a,b])=e^{-\lambda a}-e^{-\lambda b}
=
\int_a^b \lambda e^{-\lambda t}\,dt.
\]

Плотность равна

\[
f(t)=\lambda e^{-\lambda t}\mathbf 1_{\{t\ge0\}}.
\]

\subsubsection{Пример 7. Почему конечной аддитивности
мало}\label{ux43fux440ux438ux43cux435ux440-7.-ux43fux43eux447ux435ux43cux443-ux43aux43eux43dux435ux447ux43dux43eux439-ux430ux434ux434ux438ux442ux438ux432ux43dux43eux441ux442ux438-ux43cux430ux43bux43e}

На алгебре конечных и ко-конечных подмножеств \(\mathbb N\) можно задать
функцию

\[
\nu(A)=
\begin{cases}
0, & A\text{ конечно},\\
1, & A\text{ ко-конечно}.
\end{cases}
\]

Она конечно-аддитивна и нормирована: \(\nu(\mathbb N)=1\). Но она не
является счетно-аддитивной, потому что

\[
\mathbb N=\bigsqcup_{n=1}^{\infty}\{n\},
\]

и тогда счетная аддитивность потребовала бы

\[
1=\nu(\mathbb N)=\sum_{n=1}^{\infty}\nu(\{n\})=0.
\]

Это показывает, что счетная аддитивность не является формальностью. Она
запрещает такие патологические вероятности и делает возможным предельный
анализ.

\subsection{6. Доказательства и
обоснования}\label{ux434ux43eux43aux430ux437ux430ux442ux435ux43bux44cux441ux442ux432ux430-ux438-ux43eux431ux43eux441ux43dux43eux432ux430ux43dux438ux44f-1}

\textbf{Почему \(\sigma\)-алгебра нужна вместе со счетной
аддитивностью.}

Счетная аддитивность говорит о вероятности счетного объединения:

\[
\mathbb P\left(\bigsqcup_{n=1}^{\infty}A_n\right).
\]

Чтобы левая часть имела смысл, само объединение должно быть событием.
Поэтому область определения вероятности должна быть замкнута
относительно счетных объединений. Именно это и обеспечивает
\(\sigma\)-алгебра. Если бы \(\mathcal F\) была только алгеброй, счетное
объединение событий могло бы не принадлежать \(\mathcal F\), и аксиома
счетной аддитивности не была бы применима.

\textbf{Почему непрерывность снизу эквивалентна счетной аддитивности при
конечной аддитивности.}

Если мера конечно-аддитивна и непрерывна снизу, то для попарно
непересекающихся \(A_n\) положим

\[
C_N=\bigsqcup_{n=1}^{N}A_n,
\qquad
C=\bigsqcup_{n=1}^{\infty}A_n.
\]

Тогда

\[
C_1\subset C_2\subset\cdots,
\qquad
C=\bigcup_{N=1}^{\infty}C_N.
\]

По непрерывности снизу:

\[
\mathbb P(C)=\lim_{N\to\infty}\mathbb P(C_N).
\]

По конечной аддитивности:

\[
\mathbb P(C_N)=\sum_{n=1}^{N}\mathbb P(A_n).
\]

Следовательно,

\[
\mathbb P(C)
=
\lim_{N\to\infty}\sum_{n=1}^{N}\mathbb P(A_n)
=
\sum_{n=1}^{\infty}\mathbb P(A_n).
\]

Значит получаем счетную аддитивность. Обратное направление уже доказано
в свойстве о непрерывности снизу. Поэтому в вероятностной мере счетная
аддитивность и непрерывность снизу тесно связаны.

\textbf{Почему правая непрерывность важна для функции распределения.}

Если \(\mu\) - вероятностная мера и

\[
F(x)=\mu((-\infty,x]),
\]

то \(F\) автоматически непрерывна справа. Действительно, если
\(x_n\downarrow x\), то множества

\[
(-\infty,x_n]\downarrow(-\infty,x].
\]

По непрерывности меры сверху:

\[
\mu((-\infty,x_n])\downarrow\mu((-\infty,x]).
\]

То есть

\[
F(x_n)\downarrow F(x).
\]

Это и есть правая непрерывность. Поэтому в обратной теореме о построении
меры по \(F\) правая непрерывность не случайное техническое условие, а
необходимое свойство любой настоящей функции распределения.

\textbf{Почему \(\mu_F(\{a\})=F(a)-F(a-)\).}

Рассмотрим интервалы

\[
\left(a-\frac1n,a\right],\qquad n=1,2,\ldots
\]

Они убывают к одноточечному множеству:

\[
\bigcap_{n=1}^{\infty}\left(a-\frac1n,a\right]=\{a\}.
\]

По непрерывности сверху:

\[
\mu_F(\{a\})
=
\lim_{n\to\infty}\mu_F\left(\left(a-\frac1n,a\right]\right).
\]

По определению меры на полуинтервалах:

\[
\mu_F\left(\left(a-\frac1n,a\right]\right)
=
F(a)-F\left(a-\frac1n\right).
\]

Переходя к пределу, получаем

\[
\mu_F(\{a\})=F(a)-F(a-).
\]

\textbf{Почему мера по функции распределения единственна.}

Полуинтервалы \((-\infty,x]\) порождают борелевскую \(\sigma\)-алгебру:

\[
\mathcal B(\mathbb R)=\sigma\{(-\infty,x]:x\in\mathbb R\}.
\]

Если две вероятностные меры совпадают на всех множествах
\((-\infty,x]\), то они совпадают на порожденной ими \(\sigma\)-алгебре.
Строго это доказывается через \(\pi\)-\(\lambda\)-теорему или через
теорему единственности продолжения меры. Поэтому функция распределения
однозначно задает борелевскую вероятность на прямой.

\subsection{7. Что написать на
доске}\label{ux447ux442ux43e-ux43dux430ux43fux438ux441ux430ux442ux44c-ux43dux430-ux434ux43eux441ux43aux435-1}

\begin{enumerate}
\def\labelenumi{\arabic{enumi}.}
\tightlist
\item
  Определение колмогоровской тройки:
\end{enumerate}

\[
(\Omega,\mathcal F,\mathbb P),\qquad
\mathbb P:\mathcal F\to[0,1],\qquad
\mathbb P(\Omega)=1.
\]

\begin{enumerate}
\def\labelenumi{\arabic{enumi}.}
\setcounter{enumi}{1}
\tightlist
\item
  Счетная аддитивность:
\end{enumerate}

\[
A_i\cap A_j=\varnothing,\ i\ne j
\quad\Rightarrow\quad
\mathbb P\left(\bigsqcup_{n=1}^{\infty}A_n\right)
=
\sum_{n=1}^{\infty}\mathbb P(A_n).
\]

\begin{enumerate}
\def\labelenumi{\arabic{enumi}.}
\setcounter{enumi}{2}
\tightlist
\item
  Базовые следствия:
\end{enumerate}

\[
\mathbb P(\varnothing)=0,\qquad
\mathbb P(\overline{A})=1-\mathbb P(A),
\]

\[
A\subset B\Rightarrow \mathbb P(A)\le\mathbb P(B),
\]

\[
\mathbb P(A\cup B)=\mathbb P(A)+\mathbb P(B)-\mathbb P(A\cap B),
\]

\[
\mathbb P\left(\bigcup_{n=1}^{\infty}A_n\right)
\le
\sum_{n=1}^{\infty}\mathbb P(A_n).
\]

\begin{enumerate}
\def\labelenumi{\arabic{enumi}.}
\setcounter{enumi}{3}
\tightlist
\item
  Непрерывность вероятности:
\end{enumerate}

\[
A_n\uparrow A\Rightarrow \mathbb P(A_n)\uparrow\mathbb P(A),
\]

\[
A_n\downarrow A\Rightarrow \mathbb P(A_n)\downarrow\mathbb P(A).
\]

\begin{enumerate}
\def\labelenumi{\arabic{enumi}.}
\setcounter{enumi}{4}
\tightlist
\item
  Функция распределения и порожденная мера:
\end{enumerate}

\[
F(x)=\mu_F((-\infty,x]),
\]

\[
\mu_F((a,b])=F(b)-F(a),
\]

\[
\mu_F(\{a\})=F(a)-F(a-).
\]

\begin{enumerate}
\def\labelenumi{\arabic{enumi}.}
\setcounter{enumi}{5}
\tightlist
\item
  Условия на \(F\):
\end{enumerate}

\[
F\text{ не убывает},\qquad
F\text{ непрерывна справа},
\]

\[
\lim_{x\to-\infty}F(x)=0,\qquad
\lim_{x\to+\infty}F(x)=1.
\]

\subsection{8. Типичные дополнительные
вопросы}\label{ux442ux438ux43fux438ux447ux43dux44bux435-ux434ux43eux43fux43eux43bux43dux438ux442ux435ux43bux44cux43dux44bux435-ux432ux43eux43fux440ux43eux441ux44b-1}

\textbf{Вопрос. Чем вероятностная мера отличается от обычной меры?}

Вероятностная мера - это мера полной массы \(1\). То есть она
счетно-аддитивна и неотрицательна, как обычная мера, но дополнительно
нормирована:

\[
\mathbb P(\Omega)=1.
\]

\textbf{Вопрос. Почему нельзя ограничиться конечной аддитивностью?}

Потому что в теории вероятностей постоянно встречаются предельные
события: счетные объединения, счетные пересечения, \(\limsup A_n\),
\(\liminf A_n\). Конечная аддитивность не дает непрерывности вероятности
по монотонным последовательностям событий. Без нее нельзя корректно
доказывать предельные теоремы.

\textbf{Вопрос. Счетная аддитивность применяется к любым \(A_n\)?}

Нет. Равенство с суммой применяется к попарно непересекающимся событиям.
Для произвольных событий верна только субаддитивность:

\[
\mathbb P\left(\bigcup_n A_n\right)\le \sum_n\mathbb P(A_n).
\]

\textbf{Вопрос. Почему \(\mathbb P(A)\le1\)?}

Потому что \(A\subset\Omega\), а по монотонности:

\[
\mathbb P(A)\le\mathbb P(\Omega)=1.
\]

\textbf{Вопрос. Что означает \(A_n\uparrow A\)?}

Это означает, что

\[
A_1\subset A_2\subset\cdots,
\qquad
A=\bigcup_{n=1}^{\infty}A_n.
\]

Тогда вероятности также возрастают к вероятности предела:

\[
\mathbb P(A_n)\uparrow\mathbb P(A).
\]

\textbf{Вопрос. Что означает \(A_n\downarrow A\)?}

Это означает, что

\[
A_1\supset A_2\supset\cdots,
\qquad
A=\bigcap_{n=1}^{\infty}A_n.
\]

Для вероятностной меры:

\[
\mathbb P(A_n)\downarrow\mathbb P(A).
\]

\textbf{Вопрос. Почему функция распределения должна быть непрерывна
справа?}

Потому что для любой меры

\[
F(x)=\mu((-\infty,x])
\]

и при \(x_n\downarrow x\) множества \((-\infty,x_n]\) убывают к
\((-\infty,x]\). Непрерывность меры сверху дает
\(F(x_n)\downarrow F(x)\).

\textbf{Вопрос. Что такое атом распределения?}

Атом в точке \(a\) - это положительная масса одноточечного множества:

\[
\mu(\{a\})>0.
\]

Для функции распределения:

\[
\mu(\{a\})=F(a)-F(a-).
\]

Значит атомы соответствуют скачкам функции распределения.

\textbf{Вопрос. Если функция распределения непрерывна, значит ли это,
что есть плотность?}

Нет. Непрерывность означает отсутствие атомов, но не гарантирует
существование плотности. Бывают непрерывные сингулярные распределения,
например распределение Кантора.

\textbf{Вопрос. Почему мера \(\mu_F\) единственна?}

Потому что значения на лучах \((-\infty,x]\) задают значения функции
распределения, а эти лучи порождают всю борелевскую \(\sigma\)-алгебру
\(\mathcal B(\mathbb R)\). Две вероятностные меры, совпадающие на таком
порождающем классе, совпадают на всех борелевских множествах.

\subsection{9. Типичные
ошибки}\label{ux442ux438ux43fux438ux447ux43dux44bux435-ux43eux448ux438ux431ux43aux438-1}

\begin{enumerate}
\def\labelenumi{\arabic{enumi}.}
\item
  Писать счетную аддитивность для непересекающихся событий не указав
  условие \(A_i\cap A_j=\varnothing\).
\item
  Путать аддитивность и субаддитивность. Для произвольных событий обычно
  верно только
\end{enumerate}

\[
\mathbb P\left(\bigcup_n A_n\right)\le\sum_n\mathbb P(A_n).
\]

\begin{enumerate}
\def\labelenumi{\arabic{enumi}.}
\setcounter{enumi}{2}
\item
  Называть любое подмножество \(\Omega\) событием. В абстрактной модели
  событие - это только элемент \(\mathcal F\).
\item
  Забывать нормировку \(\mathbb P(\Omega)=1\). Без нее была бы просто
  конечная мера, а не вероятностная мера.
\item
  Считать непрерывность сверху верной для любой меры без условий. Для
  общей меры нужно \(\mu(A_1)<\infty\). В вероятностном пространстве это
  условие автоматическое.
\item
  Писать для функции распределения меру интервала как \(F(b)-F(a)\) без
  уточнения типа интервала. Формула
\end{enumerate}

\[
\mu_F((a,b])=F(b)-F(a)
\]

соответствует правонепрерывной функции распределения
\(F(x)=\mu((-\infty,x])\).

\begin{enumerate}
\def\labelenumi{\arabic{enumi}.}
\setcounter{enumi}{6}
\item
  Думать, что непрерывная функция распределения обязательно имеет
  плотность. Это неверно: отсутствие скачков не равно абсолютной
  непрерывности.
\item
  Забывать, что теорема о мере, порожденной \(F\), строит меру на
  \(\mathcal B(\mathbb R)\), а не сразу на всех подмножествах
  \(\mathbb R\).
\end{enumerate}

\subsection{10. Связи с другими
билетами}\label{ux441ux432ux44fux437ux438-ux441-ux434ux440ux443ux433ux438ux43cux438-ux431ux438ux43bux435ux442ux430ux43cux438-1}

\textbf{Билет 01.} Здесь используются \(\sigma\)-алгебры и меры. Билет 2
добавляет счетную аддитивность и нормировку \(\mathbb P(\Omega)=1\).

\textbf{Билет 03.} Теорема о мере, порожденной функцией распределения,
доказывается через продолжение меры с полуалгебры или алгебры на
\(\sigma\)-алгебру.

\textbf{Билет 04.} Мера \(\mu_F\) строится на борелевской
\(\sigma\)-алгебре \(\mathcal B(\mathbb R)\), поэтому нужны борелевские
множества и пополнение меры.

\textbf{Билет 05.} Случайная величина - это измеримое отображение. Ее
распределение является образом вероятностной меры, а функция
распределения задает это распределение на прямой.

\textbf{Билет 06.} Распределения случайных величин и векторов опираются
на вероятностные меры и функции распределения.

\textbf{Билет 07.} Интеграл Лебега строится по мере. Математическое
ожидание - это интеграл по вероятностной мере.

\textbf{Билет 15 и Билет 16.} Теорема Колмогорова о согласованных
конечномерных распределениях и случайные процессы используют абстрактные
вероятностные пространства как базовую модель.

\subsection{11. Минимум на
удовлетворительно}\label{ux43cux438ux43dux438ux43cux443ux43c-ux43dux430-ux443ux434ux43eux432ux43bux435ux442ux432ux43eux440ux438ux442ux435ux43bux44cux43dux43e-1}

Нужно уметь сказать:

\begin{enumerate}
\def\labelenumi{\arabic{enumi}.}
\tightlist
\item
  Вероятностное пространство - это тройка
  \((\Omega,\mathcal F,\mathbb P)\).
\item
  \(\Omega\) - исходы, \(\mathcal F\) - \(\sigma\)-алгебра событий,
  \(\mathbb P\) - вероятность.
\item
  Главные условия:
\end{enumerate}

\[
\mathbb P(\Omega)=1,
\qquad
\mathbb P\left(\bigsqcup_{n=1}^{\infty}A_n\right)
=
\sum_{n=1}^{\infty}\mathbb P(A_n).
\]

\begin{enumerate}
\def\labelenumi{\arabic{enumi}.}
\setcounter{enumi}{3}
\tightlist
\item
  Из аксиом следуют \(\mathbb P(\varnothing)=0\),
  \(\mathbb P(\overline{A})=1-\mathbb P(A)\), монотонность.
\item
  Функция распределения \(F\) порождает меру \(\mu_F\) на
  \(\mathcal B(\mathbb R)\) по формуле
\end{enumerate}

\[
\mu_F((a,b])=F(b)-F(a).
\]

\begin{enumerate}
\def\labelenumi{\arabic{enumi}.}
\setcounter{enumi}{5}
\tightlist
\item
  Условия на \(F\): неубывание, правая непрерывность, пределы \(0\) и
  \(1\) на бесконечностях.
\end{enumerate}

\subsection{12. Минимум на
хорошо}\label{ux43cux438ux43dux438ux43cux443ux43c-ux43dux430-ux445ux43eux440ux43eux448ux43e-1}

Кроме минимума на удовлетворительно, нужно:

\begin{enumerate}
\def\labelenumi{\arabic{enumi}.}
\tightlist
\item
  Доказать \(\mathbb P(\varnothing)=0\) из счетной аддитивности.
\item
  Доказать формулу дополнения и монотонность.
\item
  Доказать счетную субаддитивность через дизъюнктизацию:
\end{enumerate}

\[
B_1=A_1,\qquad
B_n=A_n\setminus\bigcup_{k=1}^{n-1}A_k.
\]

\begin{enumerate}
\def\labelenumi{\arabic{enumi}.}
\setcounter{enumi}{3}
\tightlist
\item
  Сформулировать непрерывность вероятности снизу и сверху:
\end{enumerate}

\[
A_n\uparrow A\Rightarrow \mathbb P(A_n)\uparrow\mathbb P(A),
\]

\[
A_n\downarrow A\Rightarrow \mathbb P(A_n)\downarrow\mathbb P(A).
\]

\begin{enumerate}
\def\labelenumi{\arabic{enumi}.}
\setcounter{enumi}{4}
\tightlist
\item
  Объяснить, почему счетная аддитивность нужна для предельных событий.
\item
  Привести 2-3 примера мер, порожденных функцией распределения: Дирак,
  Бернулли, равномерное распределение.
\end{enumerate}

\subsection{13. Что добавить для
отлично}\label{ux447ux442ux43e-ux434ux43eux431ux430ux432ux438ux442ux44c-ux434ux43bux44f-ux43eux442ux43bux438ux447ux43dux43e-1}

Для сильного ответа стоит добавить:

\begin{enumerate}
\def\labelenumi{\arabic{enumi}.}
\tightlist
\item
  Связь счетной аддитивности с непрерывностью меры снизу: при конечной
  аддитивности непрерывность снизу фактически дает счетную аддитивность.
\item
  Контрпример конечно-аддитивной, но не счетно-аддитивной функции на
  конечных и ко-конечных подмножествах \(\mathbb N\).
\item
  Пояснение, что правая непрерывность функции распределения следует из
  непрерывности меры сверху.
\item
  Формулу атома:
\end{enumerate}

\[
\mu_F(\{a\})=F(a)-F(a-).
\]

\begin{enumerate}
\def\labelenumi{\arabic{enumi}.}
\setcounter{enumi}{4}
\tightlist
\item
  Объяснение, что непрерывная функция распределения не обязана иметь
  плотность.
\item
  Идею построения \(\mu_F\): полуинтервалы \(\to\) алгебра конечных
  объединений \(\to\) продолжение меры на \(\mathcal B(\mathbb R)\).
\item
  Упоминание, что доказательство существования \(\mu_F\) в этом билете
  не требуется, но относится к теореме продолжения меры.
\end{enumerate}

\subsection{14. Устная версия ответа на 5
минут}\label{ux443ux441ux442ux43dux430ux44f-ux432ux435ux440ux441ux438ux44f-ux43eux442ux432ux435ux442ux430-ux43dux430-5-ux43cux438ux43dux443ux442-1}

Вероятностное пространство в аксиоматике Колмогорова - это тройка
\((\Omega,\mathcal F,\mathbb P)\). Здесь \(\Omega\) - множество
элементарных исходов, \(\mathcal F\) - \(\sigma\)-алгебра событий, а
\(\mathbb P\) - вероятностная мера. Событиями являются не все
подмножества \(\Omega\), а только элементы \(\mathcal F\).

Вероятность удовлетворяет трем основным условиям: неотрицательность,
нормировка \(\mathbb P(\Omega)=1\) и счетная аддитивность. Счетная
аддитивность означает, что для попарно непересекающихся событий
\(A_1,A_2,\ldots\) выполнено

\[
\mathbb P\left(\bigsqcup_{n=1}^{\infty}A_n\right)
=
\sum_{n=1}^{\infty}\mathbb P(A_n).
\]

Это сильнее конечной аддитивности и нужно для работы с предельными
событиями. Например, события вида \(\bigcup_n A_n\), \(\bigcap_n A_n\) и
\(\limsup A_n\) естественно возникают в теории случайных процессов и
предельных теоремах.

Из счетной аддитивности выводятся основные свойства. Во-первых,
\(\mathbb P(\varnothing)=0\): применяем счетную аддитивность к
разложению \(\Omega=\Omega\sqcup\varnothing\sqcup\cdots\). Во-вторых, из
разложения \(\Omega=A\sqcup \overline{A}\) следует
\(\mathbb P(\overline{A})=1-\mathbb P(A)\). В-третьих, если
\(A\subset B\), то \(B=A\sqcup(B\setminus A)\), поэтому
\(\mathbb P(B)=\mathbb P(A)+\mathbb P(B\setminus A)\) и
\(\mathbb P(A)\le\mathbb P(B)\).

Также важна счетная субаддитивность:

\[
\mathbb P\left(\bigcup_n A_n\right)\le\sum_n\mathbb P(A_n).
\]

Она доказывается заменой \(A_n\) на непересекающиеся множества
\(B_1=A_1\), \(B_n=A_n\setminus\bigcup_{k<n}A_k\). Еще одно важное
следствие - непрерывность вероятности. Если \(A_n\uparrow A\), то
\(\mathbb P(A_n)\uparrow\mathbb P(A)\); если \(A_n\downarrow A\), то
\(\mathbb P(A_n)\downarrow\mathbb P(A)\).

Теперь пример меры, порожденной функцией распределения. Пусть
\(F:\mathbb R\to[0,1]\) не убывает, непрерывна справа и имеет пределы
\(0\) при \(x\to-\infty\) и \(1\) при \(x\to+\infty\). Тогда существует
единственная вероятностная мера \(\mu_F\) на \(\mathcal B(\mathbb R)\)
такая, что

\[
\mu_F((-\infty,x])=F(x).
\]

Отсюда для полуинтервала:

\[
\mu_F((a,b])=F(b)-F(a).
\]

Доказательство существования здесь обычно не требуется. Идея: задать
меру на полуинтервалах, затем на конечных объединениях полуинтервалов,
после чего применить теорему продолжения меры. Скачки \(F\) дают атомы:

\[
\mu_F(\{a\})=F(a)-F(a-).
\]

Например, ступенчатая функция распределения дает дискретную меру,
непрерывная линейная на \([0,1]\) функция дает равномерное
распределение, а функция \(1-e^{-\lambda x}\) при \(x\ge0\) дает
экспоненциальное распределение.

\subsection{15. Если осталось 2 минуты перед
экзаменом}\label{ux435ux441ux43bux438-ux43eux441ux442ux430ux43bux43eux441ux44c-2-ux43cux438ux43dux443ux442ux44b-ux43fux435ux440ux435ux434-ux44dux43aux437ux430ux43cux435ux43dux43eux43c-1}

Запомнить главное:

\[
(\Omega,\mathcal F,\mathbb P)
\]

\begin{itemize}
\tightlist
\item
  это пространство исходов, \(\sigma\)-алгебра событий и вероятность.
\end{itemize}

Главная аксиома:

\[
\mathbb P\left(\bigsqcup_{n=1}^{\infty}A_n\right)
=
\sum_{n=1}^{\infty}\mathbb P(A_n)
\]

для попарно непересекающихся \(A_n\).

Основные следствия:

\[
\mathbb P(\varnothing)=0,\qquad
\mathbb P(\overline{A})=1-\mathbb P(A),
\]

\[
A\subset B\Rightarrow\mathbb P(A)\le\mathbb P(B),
\]

\[
A_n\uparrow A\Rightarrow\mathbb P(A_n)\uparrow\mathbb P(A),
\]

\[
A_n\downarrow A\Rightarrow\mathbb P(A_n)\downarrow\mathbb P(A).
\]

Функция распределения \(F\) должна быть неубывающей, непрерывной справа,
с пределами \(0\) и \(1\) на бесконечностях. Она задает меру на
\(\mathcal B(\mathbb R)\):

\[
\mu_F((-\infty,x])=F(x),
\qquad
\mu_F((a,b])=F(b)-F(a).
\]

Скачок функции распределения равен массе точки:

\[
\mu_F(\{a\})=F(a)-F(a-).
\]

\newpage

\section{\texorpdfstring{Билет 03. Продолжение меры с полуалгебры на
алгебру и \(\sigma\)-алгебру. Мера
Лебега}{Билет 03. Продолжение меры с полуалгебры на алгебру и \textbackslash sigma-алгебру. Мера Лебега}}\label{ux431ux438ux43bux435ux442-03.-ux43fux440ux43eux434ux43eux43bux436ux435ux43dux438ux435-ux43cux435ux440ux44b-ux441-ux43fux43eux43bux443ux430ux43bux433ux435ux431ux440ux44b-ux43dux430-ux430ux43bux433ux435ux431ux440ux443-ux438-sigma-ux430ux43bux433ux435ux431ux440ux443.-ux43cux435ux440ux430-ux43bux435ux431ux435ux433ux430}

Источники: Ширяев 1, гл. II, §3; лекции ДСП.

\subsection{1. Короткий
ответ}\label{ux43aux43eux440ux43eux442ux43aux438ux439-ux43eux442ux432ux435ux442-2}

Меру часто задают сначала на простых множествах, где ее значение
очевидно, а затем продолжают на более богатые классы:

\[
\mathcal S\longrightarrow\mathcal A(\mathcal S)\longrightarrow\sigma(\mathcal S).
\]

Здесь \(\mathcal S\) - полуалгебра, \(\mathcal A(\mathcal S)\) - алгебра
конечных дизъюнктных объединений элементов \(\mathcal S\), а
\(\sigma(\mathcal S)\) - порожденная \(\sigma\)-алгебра.

Если

\[
A=\bigsqcup_{k=1}^n S_k,\qquad S_k\in\mathcal S,
\]

то продолжение на алгебру задается формулой

\[
\mu(A)=\sum_{k=1}^n\mu(S_k).
\]

Главный вопрос на этом шаге - корректность: сумма не должна зависеть от
выбранного разбиения \(A\).

Продолжение на \(\sigma\)-алгебру строится через внешнюю меру:

\[
\mu^*(E)=
\inf\left\{
\sum_{n=1}^{\infty}\mu(A_n):
E\subset\bigcup_{n=1}^{\infty}A_n,\ A_n\in\mathcal A
\right\}.
\]

Из внешней меры выделяют измеримые множества по критерию Каратеодори:

\[
\mu^*(B)=\mu^*(B\cap E)+\mu^*(B\setminus E)
\]

для всех \(B\subset\Omega\). Эти множества образуют \(\sigma\)-алгебру,
а \(\mu^*\) на ней становится счетно-аддитивной мерой. При
\(\sigma\)-конечности продолжение на \(\sigma(\mathcal A)\) единственно.

Мера Лебега получается как продолжение объема полуоткрытых
прямоугольников:

\[
\lambda_d\left(\prod_{j=1}^d[a_j,b_j)\right)
=
\prod_{j=1}^d(b_j-a_j).
\]

На прямой это просто длина:

\[
\lambda([a,b))=b-a.
\]

\subsection{2. Полный
ответ}\label{ux43fux43eux43bux43dux44bux439-ux43eux442ux432ux435ux442-2}

\subsubsection{3.1. Введение и необходимость продолжения
меры}\label{ux432ux432ux435ux434ux435ux43dux438ux435-ux438-ux43dux435ux43eux431ux445ux43eux434ux438ux43cux43eux441ux442ux44c-ux43fux440ux43eux434ux43eux43bux436ux435ux43dux438ux44f-ux43cux435ux440ux44b}

В билете 1 были введены системы множеств: полуалгебры, алгебры и
\(\sigma\)-алгебры. В билете 2 вероятность была определена как
счетно-аддитивная мера на \(\sigma\)-алгебре. Теперь нужно объяснить,
откуда такие меры берутся.

Обычно меру невозможно удобно задать сразу на всей \(\sigma\)-алгебре.
Например, для меры Лебега естественно начать не со всех борелевских
множеств, а с полуинтервалов на прямой:

\[
\lambda([a,b))=b-a.
\]

В \(\mathbb R^d\) естественно начать с полуоткрытых прямоугольников:

\[
I=\prod_{j=1}^d[a_j,b_j),
\qquad
\lambda_d(I)=\prod_{j=1}^d(b_j-a_j).
\]

Эти множества простые: их легко пересекать, разрезать и считать их
длину, площадь или объем. Но для вероятности, интеграла Лебега и
случайных процессов нужны более сложные множества: счетные объединения,
счетные пересечения, пределы последовательностей множеств. Поэтому нужна
процедура продолжения меры.

\subsubsection{3.2. Продолжение с полуалгебры на
алгебру}\label{ux43fux440ux43eux434ux43eux43bux436ux435ux43dux438ux435-ux441-ux43fux43eux43bux443ux430ux43bux433ux435ux431ux440ux44b-ux43dux430-ux430ux43bux433ux435ux431ux440ux443}

Первый шаг - продолжение с полуалгебры \(\mathcal S\) на алгебру
\(\mathcal A(\mathcal S)\). Алгебра \(\mathcal A(\mathcal S)\) состоит
из конечных дизъюнктных объединений элементов \(\mathcal S\):

\[
\mathcal A(\mathcal S)=
\left\{
\bigsqcup_{k=1}^n S_k:\ S_k\in\mathcal S,\ n<\infty
\right\}.
\]

Если \(A\in\mathcal A(\mathcal S)\), то \(A\) представимо в виде

\[
A=\bigsqcup_{k=1}^n S_k.
\]

Естественное определение:

\[
\mu(A)=\sum_{k=1}^n\mu(S_k).
\]

Нужно доказать, что если у \(A\) есть другое разбиение, то сумма
получится та же. Это делается через общее измельчение двух разбиений.
Если

\[
A=\bigsqcup_{i=1}^m S_i=\bigsqcup_{j=1}^n T_j,
\]

то элементы \(S_i\cap T_j\) дают общее разбиение. Так как полуалгебра
замкнута относительно пересечений, эти множества снова лежат в
\(\mathcal S\). Аддитивность на \(\mathcal S\) дает равенство сумм.

\subsubsection{\texorpdfstring{3.3. Продолжение на \(\sigma\)-алгебру:
внешняя
мера}{3.3. Продолжение на \textbackslash sigma-алгебру: внешняя мера}}\label{ux43fux440ux43eux434ux43eux43bux436ux435ux43dux438ux435-ux43dux430-sigma-ux430ux43bux433ux435ux431ux440ux443-ux432ux43dux435ux448ux43dux44fux44f-ux43cux435ux440ux430}

Второй шаг - продолжение с алгебры на \(\sigma\)-алгебру. Пусть \(\mu\)
- предмера на алгебре \(\mathcal A\), то есть счетно-аддитивная мера на
\(\mathcal A\) в том смысле, что счетная аддитивность выполняется для
тех счетных дизъюнктных объединений, которые снова лежат в
\(\mathcal A\). Тогда строят внешнюю меру:

\[
\mu^*(E)=
\inf\left\{
\sum_{n=1}^{\infty}\mu(A_n):
E\subset\bigcup_{n=1}^{\infty}A_n,\ A_n\in\mathcal A
\right\}.
\]

Смысл: множество \(E\) покрывается счетным набором уже известных
измеримых множеств \(A_n\), а \(\mu^*(E)\) - наименьшая возможная цена
такого покрытия.

Внешняя мера определена на всех подмножествах \(\Omega\), но на всех
подмножествах она обычно не является мерой. Поэтому выделяют класс
хороших множеств.

\subsubsection{3.4. Условие измеримости по
Каратеодори}\label{ux443ux441ux43bux43eux432ux438ux435-ux438ux437ux43cux435ux440ux438ux43cux43eux441ux442ux438-ux43fux43e-ux43aux430ux440ux430ux442ux435ux43eux434ux43eux440ux438}

Множество \(E\) называется измеримым по Каратеодори, если для любого
\(B\subset\Omega\):

\[
\mu^*(B)=\mu^*(B\cap E)+\mu^*(B\setminus E).
\]

Это означает, что \(E\) корректно разрезает любое множество \(B\) на две
части без потери и без лишней массы. Теорема Каратеодори утверждает, что
все такие \(E\) образуют \(\sigma\)-алгебру, а \(\mu^*\) на ней является
счетно-аддитивной мерой. Кроме того, эта \(\sigma\)-алгебра содержит
исходную алгебру \(\mathcal A\), и новая мера совпадает с исходной
\(\mu\) на \(\mathcal A\).

\subsubsection{\texorpdfstring{3.5. Единственность продолжения и
\(\sigma\)-конечность}{3.5. Единственность продолжения и \textbackslash sigma-конечность}}\label{ux435ux434ux438ux43dux441ux442ux432ux435ux43dux43dux43eux441ux442ux44c-ux43fux440ux43eux434ux43eux43bux436ux435ux43dux438ux44f-ux438-sigma-ux43aux43eux43dux435ux447ux43dux43eux441ux442ux44c}

Если исходная мера \(\sigma\)-конечна, то продолжение на
\(\sigma(\mathcal A)\) единственно. \(\sigma\)-конечность означает, что
все пространство можно покрыть счетным числом множеств конечной меры:

\[
\Omega=\bigcup_{n=1}^{\infty}E_n,
\qquad
\mu(E_n)<\infty.
\]

Для меры Лебега это условие выполнено:

\[
\mathbb R^d=\bigcup_{n=1}^{\infty}[-n,n)^d,
\qquad
\lambda_d([-n,n)^d)=(2n)^d<\infty.
\]

\subsubsection{3.6. Резюме: построение меры
Лебега}\label{ux440ux435ux437ux44eux43cux435-ux43fux43eux441ux442ux440ux43eux435ux43dux438ux435-ux43cux435ux440ux44b-ux43bux435ux431ux435ux433ux430}

Мера Лебега строится именно по этой схеме. Сначала задают объем
полуоткрытого прямоугольника. Затем получают меру на алгебре
элементарных множеств, то есть конечных объединений прямоугольников.
После этого через внешнюю меру и критерий Каратеодори получают меру на
\(\sigma\)-алгебре. На борелевских множествах получается борелевская
мера Лебега, а после добавления всех подмножеств множеств меры ноль
получается полная мера Лебега.

\subsection{3. Основные
определения}\label{ux43eux441ux43dux43eux432ux43dux44bux435-ux43eux43fux440ux435ux434ux435ux43bux435ux43dux438ux44f-2}

\subsubsection{Продолжение
меры}\label{ux43fux440ux43eux434ux43eux43bux436ux435ux43dux438ux435-ux43cux435ux440ux44b}

Пусть \(\mathcal C\subset\mathcal D\). Мера \(\nu\) на \(\mathcal D\)
называется продолжением меры \(\mu\) с \(\mathcal C\), если

\[
\nu(A)=\mu(A),\qquad A\in\mathcal C.
\]

\subsubsection{Полуалгебра}\label{ux43fux43eux43bux443ux430ux43bux433ux435ux431ux440ux430-1}

Семейство \(\mathcal S\subset 2^\Omega\) такое, что
\(\varnothing\in\mathcal S\), пересечение двух множеств из
\(\mathcal S\) снова лежит в \(\mathcal S\), а разность двух множеств из
\(\mathcal S\) раскладывается в конечное дизъюнктное объединение
множеств из \(\mathcal S\):

\[
S\setminus T=\bigsqcup_{k=1}^n S_k,
\qquad S_k\in\mathcal S.
\]

\subsubsection{Алгебра, порожденная
полуалгеброй}\label{ux430ux43bux433ux435ux431ux440ux430-ux43fux43eux440ux43eux436ux434ux435ux43dux43dux430ux44f-ux43fux43eux43bux443ux430ux43bux433ux435ux431ux440ux43eux439}

\[
\mathcal A(\mathcal S)=
\left\{
\bigsqcup_{k=1}^n S_k:\ S_k\in\mathcal S,\ n<\infty
\right\}.
\]

Если \(\Omega\in\mathcal S\), то \(\mathcal A(\mathcal S)\) является
алгеброй.

\subsubsection{\texorpdfstring{\(\sigma\)-алгебра, порожденная системой
множеств}{\textbackslash sigma-алгебра, порожденная системой множеств}}\label{sigma-ux430ux43bux433ux435ux431ux440ux430-ux43fux43eux440ux43eux436ux434ux435ux43dux43dux430ux44f-ux441ux438ux441ux442ux435ux43cux43eux439-ux43cux43dux43eux436ux435ux441ux442ux432}

\[
\sigma(\mathcal C)=
\bigcap\{\mathcal F:\ \mathcal F\text{ - }\sigma\text{-алгебра},\ \mathcal C\subset\mathcal F\}.
\]

Это наименьшая \(\sigma\)-алгебра, содержащая \(\mathcal C\).

\subsubsection{Предмера на
алгебре}\label{ux43fux440ux435ux434ux43cux435ux440ux430-ux43dux430-ux430ux43bux433ux435ux431ux440ux435}

Функция \(\mu:\mathcal A\to[0,\infty]\) называется предмерой, если
\(\mu(\varnothing)=0\) и для любых попарно непересекающихся
\(A_1,A_2,\ldots\in\mathcal A\), таких что
\(\bigsqcup_{n=1}^{\infty}A_n\in\mathcal A\), выполнено

\[
\mu\left(\bigsqcup_{n=1}^{\infty}A_n\right)
=
\sum_{n=1}^{\infty}\mu(A_n).
\]

\subsubsection{\texorpdfstring{\(\sigma\)-конечность}{\textbackslash sigma-конечность}}\label{sigma-ux43aux43eux43dux435ux447ux43dux43eux441ux442ux44c}

Мера \(\mu\) называется \(\sigma\)-конечной, если существуют множества
\(E_n\) конечной меры такие, что

\[
\Omega=\bigcup_{n=1}^{\infty}E_n,
\qquad
\mu(E_n)<\infty.
\]

\subsubsection{Внешняя
мера}\label{ux432ux43dux435ux448ux43dux44fux44f-ux43cux435ux440ux430}

Функция \(\mu^*:2^\Omega\to[0,\infty]\), для которой:

\begin{enumerate}
\def\labelenumi{\arabic{enumi}.}
\tightlist
\item
  \(\mu^*(\varnothing)=0\);
\item
  \(A\subset B\Rightarrow\mu^*(A)\le\mu^*(B)\);
\item
  для любых \(A_1,A_2,\ldots\subset\Omega\):
\end{enumerate}

\[
\mu^*\left(\bigcup_{n=1}^{\infty}A_n\right)
\le
\sum_{n=1}^{\infty}\mu^*(A_n).
\]

\subsubsection{Измеримость по
Каратеодори}\label{ux438ux437ux43cux435ux440ux438ux43cux43eux441ux442ux44c-ux43fux43e-ux43aux430ux440ux430ux442ux435ux43eux434ux43eux440ux438}

Множество \(E\subset\Omega\) называется \(\mu^*\)-измеримым, если для
любого \(B\subset\Omega\):

\[
\mu^*(B)=\mu^*(B\cap E)+\mu^*(B\setminus E).
\]

\subsubsection{Мера
Лебега}\label{ux43cux435ux440ux430-ux43bux435ux431ux435ux433ux430}

Полная мера \(\lambda_d\) на \(\mathbb R^d\), полученная продолжением
объема полуоткрытых прямоугольников:

\[
\lambda_d\left(\prod_{j=1}^d[a_j,b_j)\right)
=
\prod_{j=1}^d(b_j-a_j).
\]

\subsection{4. Основные теоремы и
свойства}\label{ux43eux441ux43dux43eux432ux43dux44bux435-ux442ux435ux43eux440ux435ux43cux44b-ux438-ux441ux432ux43eux439ux441ux442ux432ux430-2}

\subsubsection{Теорема 1. Продолжение с полуалгебры на
алгебру}\label{ux442ux435ux43eux440ux435ux43cux430-1.-ux43fux440ux43eux434ux43eux43bux436ux435ux43dux438ux435-ux441-ux43fux43eux43bux443ux430ux43bux433ux435ux431ux440ux44b-ux43dux430-ux430ux43bux433ux435ux431ux440ux443}

Пусть \(\mathcal S\) - полуалгебра, \(\Omega\in\mathcal S\), и на
\(\mathcal S\) задана неотрицательная функция \(\mu\) с
\(\mu(\varnothing)=0\), согласованная с конечными дизъюнктными
разложениями: если

\[
S=\bigsqcup_{k=1}^n S_k,\qquad S,S_k\in\mathcal S,
\]

то

\[
\mu(S)=\sum_{k=1}^n\mu(S_k).
\]

Тогда на \(\mathcal A(\mathcal S)\) можно определить

\[
\mu\left(\bigsqcup_{k=1}^n S_k\right)
=
\sum_{k=1}^n\mu(S_k).
\]

Это определение корректно и задает аддитивное продолжение на алгебру.

\subsubsection{Доказательство}\label{ux434ux43eux43aux430ux437ux430ux442ux435ux43bux44cux441ux442ux432ux43e-15}

Пусть

\[
A=\bigsqcup_{i=1}^m S_i
=
\bigsqcup_{j=1}^n T_j,
\qquad S_i,T_j\in\mathcal S.
\]

Тогда для каждого \(i\):

\[
S_i=S_i\cap A
=
\bigsqcup_{j=1}^n(S_i\cap T_j).
\]

Так как \(\mathcal S\) замкнута относительно пересечений, все
\(S_i\cap T_j\) лежат в \(\mathcal S\). По аддитивности:

\[
\mu(S_i)=\sum_{j=1}^n\mu(S_i\cap T_j).
\]

Суммируем по \(i\):

\[
\sum_{i=1}^m\mu(S_i)
=
\sum_{i=1}^m\sum_{j=1}^n\mu(S_i\cap T_j).
\]

С другой стороны, для каждого \(j\):

\[
T_j=\bigsqcup_{i=1}^m(S_i\cap T_j),
\]

поэтому

\[
\sum_{j=1}^n\mu(T_j)
=
\sum_{j=1}^n\sum_{i=1}^m\mu(S_i\cap T_j).
\]

Двойные суммы одинаковы, значит значения по двум разбиениям совпадают.
\(\square\)

\subsubsection{Теорема 2. Теорема продолжения
меры}\label{ux442ux435ux43eux440ux435ux43cux430-2.-ux442ux435ux43eux440ux435ux43cux430-ux43fux440ux43eux434ux43eux43bux436ux435ux43dux438ux44f-ux43cux435ux440ux44b}

Пусть \(\mu\) - предмера на алгебре \(\mathcal A\). Тогда существует
мера \(\nu\) на \(\sigma\)-алгебре \(\mathcal M\), содержащей
\(\sigma(\mathcal A)\), такая что

\[
\nu(A)=\mu(A),\qquad A\in\mathcal A.
\]

Если \(\mu\) \(\sigma\)-конечна, то продолжение на
\(\sigma(\mathcal A)\) единственно.

\subsubsection{Схема
доказательства}\label{ux441ux445ux435ux43cux430-ux434ux43eux43aux430ux437ux430ux442ux435ux43bux44cux441ux442ux432ux430}

Строится внешняя мера

\[
\mu^*(E)=
\inf\left\{
\sum_{n=1}^{\infty}\mu(A_n):
E\subset\bigcup_{n=1}^{\infty}A_n,\ A_n\in\mathcal A
\right\}.
\]

Затем берутся все множества \(E\), удовлетворяющие критерию Каратеодори:

\[
\mu^*(B)=\mu^*(B\cap E)+\mu^*(B\setminus E),
\qquad B\subset\Omega.
\]

Доказывается, что они образуют \(\sigma\)-алгебру, что \(\mu^*\) на ней
счетно-аддитивна, что все множества из \(\mathcal A\) измеримы и что
\(\mu^*=\mu\) на \(\mathcal A\). Единственность при
\(\sigma\)-конечности доказывается разбиением пространства на счетное
объединение множеств конечной меры и применением теоремы единственности
на каждом таком куске.

\subsubsection{Теорема 3. Построение меры
Лебега}\label{ux442ux435ux43eux440ux435ux43cux430-3.-ux43fux43eux441ux442ux440ux43eux435ux43dux438ux435-ux43cux435ux440ux44b-ux43bux435ux431ux435ux433ux430}

На полуалгебре

\[
\mathcal S_d=
\left\{
\prod_{j=1}^d[a_j,b_j):\ a_j\le b_j
\right\}
\]

формула

\[
\lambda_d\left(\prod_{j=1}^d[a_j,b_j)\right)
=
\prod_{j=1}^d(b_j-a_j)
\]

задает объем. Он продолжается на алгебру элементарных множеств, затем на
\(\mathcal B(\mathbb R^d)\), а после пополнения дает полную меру Лебега.

\subsubsection{Обоснование}\label{ux43eux431ux43eux441ux43dux43eux432ux430ux43dux438ux435}

Полуоткрытые прямоугольники образуют полуалгебру. Объем аддитивен при
конечном разрезании прямоугольников. Значит его можно продолжить на
алгебру элементарных множеств. Эта мера \(\sigma\)-конечна, потому что

\[
\mathbb R^d=\bigcup_{n=1}^{\infty}[-n,n)^d,
\qquad
\lambda_d([-n,n)^d)<\infty.
\]

По теореме продолжения получается единственная борелевская мера,
совпадающая с объемом на прямоугольниках. Ее пополнение и называется
мерой Лебега.

\subsubsection{Свойство 4. Внешняя мера
Лебега}\label{ux441ux432ux43eux439ux441ux442ux432ux43e-4.-ux432ux43dux435ux448ux43dux44fux44f-ux43cux435ux440ux430-ux43bux435ux431ux435ux433ux430}

Для \(E\subset\mathbb R^d\):

\[
\lambda_d^*(E)
=
\inf\left\{
\sum_{n=1}^{\infty}\lambda_d(I_n):
E\subset\bigcup_{n=1}^{\infty}I_n,\ I_n\in\mathcal S_d
\right\}.
\]

Если \(E\) лебегово измеримо, то

\[
\lambda_d(E)=\lambda_d^*(E).
\]

\subsubsection{Свойство 5. Инвариантность меры Лебега относительно
сдвига}\label{ux441ux432ux43eux439ux441ux442ux432ux43e-5.-ux438ux43dux432ux430ux440ux438ux430ux43dux442ux43dux43eux441ux442ux44c-ux43cux435ux440ux44b-ux43bux435ux431ux435ux433ux430-ux43eux442ux43dux43eux441ux438ux442ux435ux43bux44cux43dux43e-ux441ux434ux432ux438ux433ux430}

Если \(E\) лебегово измеримо и \(h\in\mathbb R^d\), то

\[
\lambda_d(E+h)=\lambda_d(E),
\qquad
E+h=\{x+h:\ x\in E\}.
\]

Идея доказательства: сдвиг прямоугольника не меняет его объем, поэтому
не меняет суммарные объемы покрытий. Значит внешняя мера сохраняется, а
критерий Каратеодори переносит измеримость на сдвинутое множество.

\subsubsection{Свойство 6. Точки и счетные множества имеют меру
ноль}\label{ux441ux432ux43eux439ux441ux442ux432ux43e-6.-ux442ux43eux447ux43aux438-ux438-ux441ux447ux435ux442ux43dux44bux435-ux43cux43dux43eux436ux435ux441ux442ux432ux430-ux438ux43cux435ux44eux442-ux43cux435ux440ux443-ux43dux43eux43bux44c}

Для любой точки \(x\in\mathbb R^d\):

\[
\lambda_d(\{x\})=0.
\]

\subsubsection{Доказательство}\label{ux434ux43eux43aux430ux437ux430ux442ux435ux43bux44cux441ux442ux432ux43e-16}

Точку можно покрыть кубом со стороной \(\varepsilon>0\), поэтому

\[
\lambda_d^*(\{x\})\le\varepsilon^d.
\]

Так как \(\varepsilon\) можно брать сколь угодно малым, получаем
\(\lambda_d(\{x\})=0\).

Если \(E=\{x_1,x_2,\ldots\}\) счетно, то по счетной субаддитивности:

\[
\lambda_d(E)
\le
\sum_{n=1}^{\infty}\lambda_d(\{x_n\})
=0.
\]

Значит \(\lambda_d(E)=0\). \(\square\)

\subsection{5. Примеры}\label{ux43fux440ux438ux43cux435ux440ux44b-2}

\subsubsection{Пример 1. Длина конечного объединения
полуинтервалов}\label{ux43fux440ux438ux43cux435ux440-1.-ux434ux43bux438ux43dux430-ux43aux43eux43dux435ux447ux43dux43eux433ux43e-ux43eux431ux44aux435ux434ux438ux43dux435ux43dux438ux44f-ux43fux43eux43bux443ux438ux43dux442ux435ux440ux432ux430ux43bux43eux432}

Пусть

\[
A=[0,1)\sqcup[2,5).
\]

Тогда

\[
\lambda(A)=1+3=4.
\]

Если разбить иначе:

\[
A=[0,\tfrac12)\sqcup[\tfrac12,1)\sqcup[2,3)\sqcup[3,5),
\]

то

\[
\lambda(A)=\frac12+\frac12+1+2=4.
\]

Это иллюстрирует независимость от разбиения.

\subsubsection{Пример 2. Открытые и закрытые
интервалы}\label{ux43fux440ux438ux43cux435ux440-2.-ux43eux442ux43aux440ux44bux442ux44bux435-ux438-ux437ux430ux43aux440ux44bux442ux44bux435-ux438ux43dux442ux435ux440ux432ux430ux43bux44b}

После продолжения меры:

\[
\lambda((a,b))=\lambda([a,b))=\lambda((a,b])=\lambda([a,b])=b-a.
\]

Концы не влияют на меру, потому что точки имеют меру \(0\).

\subsubsection{Пример 3. Рациональные числа на
отрезке}\label{ux43fux440ux438ux43cux435ux440-3.-ux440ux430ux446ux438ux43eux43dux430ux43bux44cux43dux44bux435-ux447ux438ux441ux43bux430-ux43dux430-ux43eux442ux440ux435ux437ux43aux435}

Множество \(\mathbb Q\cap[0,1]\) счетно, поэтому

\[
\lambda(\mathbb Q\cap[0,1])=0.
\]

Но оно всюду плотно в \([0,1]\). Значит мера Лебега измеряет длину, а не
топологическую плотность.

\subsubsection{Пример 4. Иррациональные числа на
отрезке}\label{ux43fux440ux438ux43cux435ux440-4.-ux438ux440ux440ux430ux446ux438ux43eux43dux430ux43bux44cux43dux44bux435-ux447ux438ux441ux43bux430-ux43dux430-ux43eux442ux440ux435ux437ux43aux435}

Так как

\[
[0,1]=(\mathbb Q\cap[0,1])\sqcup([0,1]\setminus\mathbb Q),
\]

то

\[
\lambda([0,1]\setminus\mathbb Q)=1.
\]

\subsubsection{Пример 5. Канторово
множество}\label{ux43fux440ux438ux43cux435ux440-5.-ux43aux430ux43dux442ux43eux440ux43eux432ux43e-ux43cux43dux43eux436ux435ux441ux442ux432ux43e}

Классическое канторово множество несчетно, но имеет меру \(0\).
Суммарная длина удаленных интервалов равна

\[
\sum_{n=1}^{\infty}\frac{2^{n-1}}{3^n}
=
\frac13\sum_{n=0}^{\infty}\left(\frac23\right)^n
=1.
\]

Из отрезка длины \(1\) удалена суммарная длина \(1\), поэтому оставшееся
канторово множество имеет меру \(0\).

\subsection{6. Доказательства и
обоснования}\label{ux434ux43eux43aux430ux437ux430ux442ux435ux43bux44cux441ux442ux432ux430-ux438-ux43eux431ux43eux441ux43dux43eux432ux430ux43dux438ux44f-2}

\textbf{Почему продолжение на алгебру задается суммой.}

Если \(A=\bigsqcup_{k=1}^n S_k\), то любая аддитивная мера обязана
удовлетворять

\[
\mu(A)=\sum_{k=1}^n\mu(S_k).
\]

Иначе она нарушала бы аддитивность. Поэтому формула единственно
возможная. Сложность только в том, чтобы доказать независимость от
разбиения.

\textbf{Почему нужна внешняя мера.}

Алгебра закрыта только относительно конечных операций. Но
\(\sigma\)-алгебра содержит результаты счетных операций. Внешняя мера
назначает размер любому множеству через счетные покрытия уже известными
множествами:

\[
E\subset\bigcup_{n=1}^{\infty}A_n.
\]

Инфимум сумм \(\sum_n\mu(A_n)\) дает наилучшую оценку размера \(E\)
сверху.

\textbf{Почему не все множества сразу измеримы.}

Внешняя мера субаддитивна на всех подмножествах, но аддитивность может
ломаться. Поэтому оставляют только те множества \(E\), которые правильно
разрезают любое \(B\):

\[
\mu^*(B)=\mu^*(B\cap E)+\mu^*(B\setminus E).
\]

Это и есть критерий Каратеодори.

\textbf{Почему \(\sigma\)-конечность важна для единственности.}

Если пространство покрыто счетным числом множеств конечной меры, то
совпадение двух продолжений можно проверять на каждом конечном куске, а
затем склеивать по счетному покрытию. Без \(\sigma\)-конечности
продолжение может быть не единственным.

\textbf{Почему мера Лебега продолжает геометрический объем.}

На прямоугольниках объем известен. При конечном разрезании
прямоугольника сумма объемов частей равна объему целого. Это дает меру
на полуалгебре и алгебре элементарных множеств. Теорема продолжения
переносит эту меру на борелевские множества, а пополнение дает полную
меру Лебега.

\subsection{7. Что написать на
доске}\label{ux447ux442ux43e-ux43dux430ux43fux438ux441ux430ux442ux44c-ux43dux430-ux434ux43eux441ux43aux435-2}

\begin{enumerate}
\def\labelenumi{\arabic{enumi}.}
\tightlist
\item
  Схема:
\end{enumerate}

\[
\mathcal S\to\mathcal A(\mathcal S)\to\sigma(\mathcal S).
\]

\begin{enumerate}
\def\labelenumi{\arabic{enumi}.}
\setcounter{enumi}{1}
\tightlist
\item
  Алгебра из полуалгебры:
\end{enumerate}

\[
\mathcal A(\mathcal S)=
\left\{
\bigsqcup_{k=1}^n S_k:\ S_k\in\mathcal S
\right\}.
\]

\begin{enumerate}
\def\labelenumi{\arabic{enumi}.}
\setcounter{enumi}{2}
\tightlist
\item
  Продолжение на алгебру:
\end{enumerate}

\[
\mu\left(\bigsqcup_{k=1}^n S_k\right)
=
\sum_{k=1}^n\mu(S_k).
\]

\begin{enumerate}
\def\labelenumi{\arabic{enumi}.}
\setcounter{enumi}{3}
\tightlist
\item
  Внешняя мера:
\end{enumerate}

\[
\mu^*(E)=
\inf\left\{
\sum_{n=1}^{\infty}\mu(A_n):
E\subset\bigcup_{n=1}^{\infty}A_n,\ A_n\in\mathcal A
\right\}.
\]

\begin{enumerate}
\def\labelenumi{\arabic{enumi}.}
\setcounter{enumi}{4}
\tightlist
\item
  Критерий Каратеодори:
\end{enumerate}

\[
\mu^*(B)=\mu^*(B\cap E)+\mu^*(B\setminus E).
\]

\begin{enumerate}
\def\labelenumi{\arabic{enumi}.}
\setcounter{enumi}{5}
\tightlist
\item
  Мера Лебега:
\end{enumerate}

\[
\lambda_d\left(\prod_{j=1}^d[a_j,b_j)\right)
=
\prod_{j=1}^d(b_j-a_j).
\]

\begin{enumerate}
\def\labelenumi{\arabic{enumi}.}
\setcounter{enumi}{6}
\tightlist
\item
  \(\sigma\)-конечность:
\end{enumerate}

\[
\mathbb R^d=\bigcup_{n=1}^{\infty}[-n,n)^d.
\]

\subsection{8. Типичные дополнительные
вопросы}\label{ux442ux438ux43fux438ux447ux43dux44bux435-ux434ux43eux43fux43eux43bux43dux438ux442ux435ux43bux44cux43dux44bux435-ux432ux43eux43fux440ux43eux441ux44b-2}

\textbf{Вопрос. Почему нельзя сразу пользоваться формулой конечной суммы
на \(\sigma(\mathcal S)\)?}

Потому что элементы \(\sigma(\mathcal S)\) могут не иметь конечного
разбиения на элементы \(\mathcal S\). Они возникают через счетные
объединения, пересечения и пределы.

\textbf{Вопрос. Что такое предмера?}

Это счетно-аддитивная мера на алгебре, еще не продолженная на
\(\sigma\)-алгебру.

\textbf{Вопрос. Что дает \(\sigma\)-конечность?}

Единственность продолжения на порожденную \(\sigma\)-алгебру.

\textbf{Вопрос. Внешняя мера определена на всех подмножествах. Значит,
все подмножества измеримы?}

Нет. Внешняя мера есть на всех подмножествах, но счетно-аддитивной мерой
она становится только на \(\sigma\)-алгебре измеримых по Каратеодори
множеств.

\textbf{Вопрос. Чем борелевская мера Лебега отличается от полной меры
Лебега?}

Борелевская мера задана на \(\mathcal B(\mathbb R^d)\). Полная мера
Лебега дополнительно содержит все подмножества борелевских множеств меры
ноль.

\textbf{Вопрос. Почему \(\lambda([a,b])=b-a\), хотя исходно задавали
\([a,b)\)?}

Потому что

\[
[a,b]=[a,b)\cup\{b\},
\]

а точка имеет меру ноль.

\textbf{Вопрос. Может ли несчетное множество иметь меру ноль?}

Да. Пример - канторово множество.

\subsection{9. Типичные
ошибки}\label{ux442ux438ux43fux438ux447ux43dux44bux435-ux43eux448ux438ux431ux43aux438-2}

\begin{enumerate}
\def\labelenumi{\arabic{enumi}.}
\tightlist
\item
  Путать классы множеств:
\end{enumerate}

\[
\mathcal S\subset\mathcal A(\mathcal S)\subset\sigma(\mathcal S).
\]

\begin{enumerate}
\def\labelenumi{\arabic{enumi}.}
\setcounter{enumi}{1}
\tightlist
\item
  Забывать, что формула
\end{enumerate}

\[
\mu\left(\bigsqcup_{k=1}^nS_k\right)=\sum_{k=1}^n\mu(S_k)
\]

требует попарной непересекаемости.

\begin{enumerate}
\def\labelenumi{\arabic{enumi}.}
\setcounter{enumi}{2}
\item
  Не проверять независимость от разбиения.
\item
  Считать внешнюю меру полноценной мерой на всех подмножествах.
\item
  Забывать \(\sigma\)-конечность в утверждении об единственности
  продолжения.
\item
  Путать борелевскую меру Лебега и полную меру Лебега.
\item
  Думать, что множество меры ноль обязательно счетно.
\item
  Использовать старое обозначение дополнения через верхний индекс. В
  этих конспектах дополнение пишется чертой сверху:
\end{enumerate}

\[
\overline{A}=\Omega\setminus A.
\]

\subsection{10. Связи с другими
билетами}\label{ux441ux432ux44fux437ux438-ux441-ux434ux440ux443ux433ux438ux43cux438-ux431ux438ux43bux435ux442ux430ux43cux438-2}

\textbf{Билет 01.} Полуалгебры, алгебры и \(\sigma\)-алгебры нужны для
построения мер через продолжение.

\textbf{Билет 02.} Вероятностная мера - частный случай меры. Мера по
функции распределения строится той же идеей продолжения.

\textbf{Билет 04.} Борелевские множества и пополнение меры продолжают
тему меры Лебега.

\textbf{Билет 05.} Измеримые функции определяются относительно
\(\sigma\)-алгебр, полученных такими построениями.

\textbf{Билет 07.} Интеграл Лебега строится по мере.

\textbf{Билет 11.} Замена переменной в интеграле использует свойства
меры Лебега.

\textbf{Билет 15 и Билет 16.} Построение мер на пространствах
последовательностей является более сложной версией идеи продолжения.

\subsection{11. Минимум на
удовлетворительно}\label{ux43cux438ux43dux438ux43cux443ux43c-ux43dux430-ux443ux434ux43eux432ux43bux435ux442ux432ux43eux440ux438ux442ux435ux43bux44cux43dux43e-2}

Нужно уметь сказать:

\begin{enumerate}
\def\labelenumi{\arabic{enumi}.}
\tightlist
\item
  Меру задают на простом классе множеств и продолжают:
\end{enumerate}

\[
\mathcal S\to\mathcal A(\mathcal S)\to\sigma(\mathcal S).
\]

\begin{enumerate}
\def\labelenumi{\arabic{enumi}.}
\setcounter{enumi}{1}
\tightlist
\item
  На алгебре конечных дизъюнктных объединений:
\end{enumerate}

\[
\mu\left(\bigsqcup_{k=1}^nS_k\right)=\sum_{k=1}^n\mu(S_k).
\]

\begin{enumerate}
\def\labelenumi{\arabic{enumi}.}
\setcounter{enumi}{2}
\item
  На \(\sigma\)-алгебру продолжение строится через внешнюю меру и
  критерий Каратеодори.
\item
  Мера Лебега продолжает длину:
\end{enumerate}

\[
\lambda([a,b))=b-a.
\]

\begin{enumerate}
\def\labelenumi{\arabic{enumi}.}
\setcounter{enumi}{4}
\tightlist
\item
  В \(\mathbb R^d\) она продолжает объем прямоугольника:
\end{enumerate}

\[
\lambda_d\left(\prod_{j=1}^d[a_j,b_j)\right)
=
\prod_{j=1}^d(b_j-a_j).
\]

\subsection{12. Минимум на
хорошо}\label{ux43cux438ux43dux438ux43cux443ux43c-ux43dux430-ux445ux43eux440ux43eux448ux43e-2}

Кроме минимума на удовлетворительно, нужно:

\begin{enumerate}
\def\labelenumi{\arabic{enumi}.}
\tightlist
\item
  Объяснить корректность продолжения на алгебру через общее измельчение
  разбиений.
\item
  Сформулировать внешнюю меру:
\end{enumerate}

\[
\mu^*(E)=
\inf\left\{
\sum_{n=1}^{\infty}\mu(A_n):
E\subset\bigcup_{n=1}^{\infty}A_n
\right\}.
\]

\begin{enumerate}
\def\labelenumi{\arabic{enumi}.}
\setcounter{enumi}{2}
\tightlist
\item
  Сформулировать критерий Каратеодори:
\end{enumerate}

\[
\mu^*(B)=\mu^*(B\cap E)+\mu^*(B\setminus E).
\]

\begin{enumerate}
\def\labelenumi{\arabic{enumi}.}
\setcounter{enumi}{3}
\tightlist
\item
  Сказать, что при \(\sigma\)-конечности продолжение единственно.
\item
  Привести примеры: длина интервала, мера точки, мера
  \(\mathbb Q\cap[0,1]\).
\end{enumerate}

\subsection{13. Что добавить для
отлично}\label{ux447ux442ux43e-ux434ux43eux431ux430ux432ux438ux442ux44c-ux434ux43bux44f-ux43eux442ux43bux438ux447ux43dux43e-2}

Для сильного ответа стоит добавить:

\begin{enumerate}
\def\labelenumi{\arabic{enumi}.}
\tightlist
\item
  Различие между мерой на полуалгебре, предмерой на алгебре и мерой на
  \(\sigma\)-алгебре.
\item
  Подробную схему доказательства теоремы продолжения через внешнюю меру.
\item
  Объяснение, почему внешняя мера не является мерой на всех
  подмножествах.
\item
  Различие между борелевской мерой Лебега и полной мерой Лебега.
\item
  Доказательство, что точка и любое счетное множество имеют меру ноль.
\item
  Пример канторова множества как несчетного множества меры ноль.
\item
  Упоминание, что без \(\sigma\)-конечности единственность продолжения
  может нарушаться.
\end{enumerate}

\subsection{14. Устная версия ответа на 5
минут}\label{ux443ux441ux442ux43dux430ux44f-ux432ux435ux440ux441ux438ux44f-ux43eux442ux432ux435ux442ux430-ux43dux430-5-ux43cux438ux43dux443ux442-2}

Тема билета - как построить меру на сложной \(\sigma\)-алгебре, если
сначала она задана только на простых множествах. Схема:

\[
\mathcal S\to\mathcal A(\mathcal S)\to\sigma(\mathcal S).
\]

Сначала есть полуалгебра \(\mathcal S\), например полуоткрытые интервалы
или прямоугольники. Из нее строится алгебра \(\mathcal A(\mathcal S)\)
конечных дизъюнктных объединений. Если

\[
A=\bigsqcup_{k=1}^nS_k,
\]

то кладем

\[
\mu(A)=\sum_{k=1}^n\mu(S_k).
\]

Главное доказать, что это не зависит от разбиения. Если есть два
разбиения \(A=\bigsqcup_iS_i=\bigsqcup_jT_j\), берем пересечения
\(S_i\cap T_j\). Они дают общее измельчение, и по аддитивности суммы
совпадают.

Дальше строится внешняя мера:

\[
\mu^*(E)=
\inf\left\{
\sum_{n=1}^{\infty}\mu(A_n):
E\subset\bigcup_{n=1}^{\infty}A_n,\ A_n\in\mathcal A
\right\}.
\]

Она задает размер любого множества через наилучшие счетные покрытия.
Затем по критерию Каратеодори выбираются измеримые множества:

\[
\mu^*(B)=\mu^*(B\cap E)+\mu^*(B\setminus E).
\]

Они образуют \(\sigma\)-алгебру, и на ней \(\mu^*\) становится мерой.
При \(\sigma\)-конечности продолжение на порожденную \(\sigma\)-алгебру
единственно.

Мера Лебега - главный пример. На полуоткрытых прямоугольниках задаем

\[
\lambda_d\left(\prod_{j=1}^d[a_j,b_j)\right)
=
\prod_{j=1}^d(b_j-a_j).
\]

Затем продолжаем на алгебру элементарных множеств, потом на борелевскую
\(\sigma\)-алгебру и после пополнения получаем полную меру Лебега. На
прямой:

\[
\lambda([a,b))=b-a.
\]

Точки и счетные множества имеют меру ноль, например
\(\lambda(\mathbb Q\cap[0,1])=0\), а \(\lambda([0,1])=1\).

\subsection{15. Если осталось 2 минуты перед
экзаменом}\label{ux435ux441ux43bux438-ux43eux441ux442ux430ux43bux43eux441ux44c-2-ux43cux438ux43dux443ux442ux44b-ux43fux435ux440ux435ux434-ux44dux43aux437ux430ux43cux435ux43dux43eux43c-2}

Запомнить:

\[
\mathcal S\to\mathcal A(\mathcal S)\to\sigma(\mathcal S).
\]

На алгебре:

\[
\mu\left(\bigsqcup_{k=1}^nS_k\right)=\sum_{k=1}^n\mu(S_k).
\]

Внешняя мера:

\[
\mu^*(E)=
\inf\left\{
\sum_n\mu(A_n):
E\subset\bigcup_n A_n
\right\}.
\]

Каратеодори:

\[
\mu^*(B)=\mu^*(B\cap E)+\mu^*(B\setminus E).
\]

Мера Лебега:

\[
\lambda([a,b))=b-a,
\qquad
\lambda_d\left(\prod_{j=1}^d[a_j,b_j)\right)
=
\prod_{j=1}^d(b_j-a_j).
\]

Главные ловушки: забыть независимость от разбиения, спутать внешнюю меру
с мерой на всех подмножествах и забыть \(\sigma\)-конечность для
единственности.

\newpage

\section{\texorpdfstring{Билет 04. Борелевские множества, борелевские
\(\sigma\)-алгебры, пополнение меры и борелевские меры
Лебега}{Билет 04. Борелевские множества, борелевские \textbackslash sigma-алгебры, пополнение меры и борелевские меры Лебега}}\label{ux431ux438ux43bux435ux442-04.-ux431ux43eux440ux435ux43bux435ux432ux441ux43aux438ux435-ux43cux43dux43eux436ux435ux441ux442ux432ux430-ux431ux43eux440ux435ux43bux435ux432ux441ux43aux438ux435-sigma-ux430ux43bux433ux435ux431ux440ux44b-ux43fux43eux43fux43eux43bux43dux435ux43dux438ux435-ux43cux435ux440ux44b-ux438-ux431ux43eux440ux435ux43bux435ux432ux441ux43aux438ux435-ux43cux435ux440ux44b-ux43bux435ux431ux435ux433ux430}

Источники: Ширяев 1, гл. II; лекции ДСП.

\subsection{1. Короткий
ответ}\label{ux43aux43eux440ux43eux442ux43aux438ux439-ux43eux442ux432ux435ux442-3}

Если \(X\) - топологическое пространство, то борелевская
\(\sigma\)-алгебра на \(X\) - это минимальная \(\sigma\)-алгебра,
содержащая все открытые множества:

\[
\mathcal B(X)=\sigma\{G\subset X:\ G\text{ открыто}\}.
\]

Ее элементы называются борелевскими множествами. В \(\mathbb R^d\)
борелевскую \(\sigma\)-алгебру можно порождать открытыми шарами,
открытыми множествами, замкнутыми множествами, прямоугольниками или
полуинтервалами:

\[
\mathcal B(\mathbb R^d)
=
\sigma\left\{\prod_{j=1}^d(a_j,b_j):a_j<b_j\right\}.
\]

Борелевская мера - это мера, заданная на \(\mathcal B(X)\). Борелевская
мера Лебега на \(\mathbb R^d\) - это мера на
\(\mathcal B(\mathbb R^d)\), которая на прямоугольниках равна обычному
объему:

\[
\lambda_d^B\left(\prod_{j=1}^d[a_j,b_j)\right)
=
\prod_{j=1}^d(b_j-a_j).
\]

Мера называется полной, если любое подмножество множества меры ноль
измеримо:

\[
N\subset Z,\quad Z\in\mathcal F,\quad \mu(Z)=0
\quad\Longrightarrow\quad
N\in\mathcal F.
\]

Пополнение меры добавляет к \(\mathcal F\) все подмножества нулевых
множеств. Полная мера Лебега - это пополнение борелевской меры Лебега.
Поэтому каждое борелевское множество лебегово измеримо, но не каждое
лебегово измеримое множество является борелевским.

\subsection{2. Полный
ответ}\label{ux43fux43eux43bux43dux44bux439-ux43eux442ux432ux435ux442-3}

\subsubsection{2.1. Введение: Соотношение борелевских и лебеговых
множеств}\label{ux432ux432ux435ux434ux435ux43dux438ux435-ux441ux43eux43eux442ux43dux43eux448ux435ux43dux438ux435-ux431ux43eux440ux435ux43bux435ux432ux441ux43aux438ux445-ux438-ux43bux435ux431ux435ux433ux43eux432ux44bux445-ux43cux43dux43eux436ux435ux441ux442ux432}

В предыдущем билете мера Лебега строилась как продолжение
геометрического объема. Теперь нужно различить две близкие, но не
одинаковые системы множеств: борелевские множества и лебегово измеримые
множества.

\subsubsection{\texorpdfstring{2.2. Определение борелевской
\(\sigma\)-алгебры}{2.2. Определение борелевской \textbackslash sigma-алгебры}}\label{ux43eux43fux440ux435ux434ux435ux43bux435ux43dux438ux435-ux431ux43eux440ux435ux43bux435ux432ux441ux43aux43eux439-sigma-ux430ux43bux433ux435ux431ux440ux44b}

Борелевская \(\sigma\)-алгебра появляется из топологии. Если на
множестве \(X\) задана топология, то известны открытые множества. Но для
меры и вероятности открытых множеств недостаточно: нужна
\(\sigma\)-алгебра, чтобы брать дополнения, счетные объединения и
счетные пересечения. Поэтому берут минимальную \(\sigma\)-алгебру,
содержащую все открытые множества:

\[
\mathcal B(X)=\sigma(\text{открытые множества}).
\]

Такую \(\sigma\)-алгебру называют борелевской, а ее элементы -
борелевскими множествами.

На прямой и в \(\mathbb R^d\) это основной класс множеств для
распределений. Например, если \(X\) - случайная величина, то ее
распределение является вероятностной мерой на \(\mathcal B(\mathbb R)\).
Если \(\xi\) - случайный вектор в \(\mathbb R^d\), то его распределение
является мерой на \(\mathcal B(\mathbb R^d)\).

\subsubsection{\texorpdfstring{2.3. Порождающие классы в
\(\mathbb R^d\)}{2.3. Порождающие классы в \textbackslash mathbb R\^{}d}}\label{ux43fux43eux440ux43eux436ux434ux430ux44eux449ux438ux435-ux43aux43bux430ux441ux441ux44b-ux432-mathbb-rd}

В \(\mathbb R^d\) борелевская \(\sigma\)-алгебра может быть порождена
разными удобными классами:

\[
\mathcal B(\mathbb R^d)
=
\sigma(\text{открытые шары})
=
\sigma(\text{замкнутые множества})
=
\sigma(\text{полуоткрытые прямоугольники}).
\]

Это важно практически: функцию распределения удобно задавать через лучи
\((-\infty,x]\), меру Лебега - через прямоугольники, а топологические
рассуждения - через открытые множества. Все эти подходы ведут к одной и
той же борелевской \(\sigma\)-алгебре.

\subsubsection{2.4. Борелевская мера и проблема
полноты}\label{ux431ux43eux440ux435ux43bux435ux432ux441ux43aux430ux44f-ux43cux435ux440ux430-ux438-ux43fux440ux43eux431ux43bux435ux43cux430-ux43fux43eux43bux43dux43eux442ux44b}

Борелевская мера - это просто мера на \(\mathcal B(X)\). Например,
борелевская мера Лебега \(\lambda_d^B\) на \(\mathbb R^d\) задается как
продолжение объема прямоугольников на \(\mathcal B(\mathbb R^d)\):

\[
\lambda_d^B\left(\prod_{j=1}^d[a_j,b_j)\right)
=
\prod_{j=1}^d(b_j-a_j).
\]

Но борелевская мера Лебега не полна. Это значит, что может существовать
борелевское множество \(Z\) меры ноль и подмножество \(N\subset Z\),
которое не является борелевским. Тогда \(N\) должно иметь меру ноль с
геометрической точки зрения, но оно не лежит в области определения
борелевской меры.

\subsubsection{2.5. Процедура пополнения
меры}\label{ux43fux440ux43eux446ux435ux434ux443ux440ux430-ux43fux43eux43fux43eux43bux43dux435ux43dux438ux44f-ux43cux435ux440ux44b}

Чтобы исправить это, меру пополняют. Пусть есть пространство с мерой
\((\Omega,\mathcal F,\mu)\). Полное пополнение имеет \(\sigma\)-алгебру

\[
\overline{\mathcal F}^{\,\mu}
=
\{A\cup N:\ A\in\mathcal F,\ N\subset Z,\ Z\in\mathcal F,\ \mu(Z)=0\}.
\]

Идея: берем старые измеримые множества \(A\) и разрешаем добавлять к ним
произвольные подмножества старых нулевых множеств. Мера на таком
множестве определяется по формуле

\[
\overline{\mu}(A\cup N)=\mu(A).
\]

Это корректно: добавление или удаление подмножеств нулевых множеств не
должно менять меру.

\subsubsection{2.6. Мера Лебега и иерархия
измеримости}\label{ux43cux435ux440ux430-ux43bux435ux431ux435ux433ux430-ux438-ux438ux435ux440ux430ux440ux445ux438ux44f-ux438ux437ux43cux435ux440ux438ux43cux43eux441ux442ux438}

Полная мера Лебега \(\lambda_d\) - это пополнение борелевской меры
Лебега \(\lambda_d^B\). Поэтому есть строгая цепочка:

\[
\mathcal B(\mathbb R^d)
\subset
\mathcal L(\mathbb R^d)
\subset
2^{\mathbb R^d},
\]

где \(\mathcal L(\mathbb R^d)\) - \(\sigma\)-алгебра лебегово измеримых
множеств. Первое включение строго: существуют лебегово измеримые
множества, которые не являются борелевскими. Второе включение тоже
строго: существуют неизмеримые по Лебегу множества, например множество
Витали.

\subsection{3. Основные
определения}\label{ux43eux441ux43dux43eux432ux43dux44bux435-ux43eux43fux440ux435ux434ux435ux43bux435ux43dux438ux44f-3}

\subsubsection{Топологическое
пространство}\label{ux442ux43eux43fux43eux43bux43eux433ux438ux447ux435ux441ux43aux43eux435-ux43fux440ux43eux441ux442ux440ux430ux43dux441ux442ux432ux43e}

Пара \((X,\tau)\), где \(\tau\) - семейство открытых подмножеств \(X\),
замкнутое относительно произвольных объединений и конечных пересечений.

\subsubsection{\texorpdfstring{Борелевская
\(\sigma\)-алгебра}{Борелевская \textbackslash sigma-алгебра}}\label{ux431ux43eux440ux435ux43bux435ux432ux441ux43aux430ux44f-sigma-ux430ux43bux433ux435ux431ux440ux430}

\[
\mathcal B(X)=\sigma(\tau)
=
\sigma\{G\subset X:\ G\text{ открыто}\}.
\]

Это минимальная \(\sigma\)-алгебра, содержащая все открытые множества.

\subsubsection{Борелевское
множество}\label{ux431ux43eux440ux435ux43bux435ux432ux441ux43aux43eux435-ux43cux43dux43eux436ux435ux441ux442ux432ux43e}

Элемент \(\mathcal B(X)\).

\subsubsection{Борелевская
мера}\label{ux431ux43eux440ux435ux43bux435ux432ux441ux43aux430ux44f-ux43cux435ux440ux430}

Мера \(\mu\) на измеримом пространстве \((X,\mathcal B(X))\).

\subsubsection{Борелевская вероятностная
мера}\label{ux431ux43eux440ux435ux43bux435ux432ux441ux43aux430ux44f-ux432ux435ux440ux43eux44fux442ux43dux43eux441ux442ux43dux430ux44f-ux43cux435ux440ux430}

Борелевская мера \(\mu\) такая, что

\[
\mu(X)=1.
\]

\subsubsection{Борелевская мера
Лебега}\label{ux431ux43eux440ux435ux43bux435ux432ux441ux43aux430ux44f-ux43cux435ux440ux430-ux43bux435ux431ux435ux433ux430}

Мера \(\lambda_d^B\) на \(\mathcal B(\mathbb R^d)\), совпадающая с
объемом на полуоткрытых прямоугольниках:

\[
\lambda_d^B\left(\prod_{j=1}^d[a_j,b_j)\right)
=
\prod_{j=1}^d(b_j-a_j).
\]

\subsubsection{Полная
мера}\label{ux43fux43eux43bux43dux430ux44f-ux43cux435ux440ux430}

Мера \(\mu\) на \((\Omega,\mathcal F)\) называется полной, если

\[
N\subset Z,\quad Z\in\mathcal F,\quad \mu(Z)=0
\quad\Longrightarrow\quad
N\in\mathcal F.
\]

\subsubsection{\texorpdfstring{Пополнение \(\sigma\)-алгебры по
мере}{Пополнение \textbackslash sigma-алгебры по мере}}\label{ux43fux43eux43fux43eux43bux43dux435ux43dux438ux435-sigma-ux430ux43bux433ux435ux431ux440ux44b-ux43fux43e-ux43cux435ux440ux435}

\[
\overline{\mathcal F}^{\,\mu}
=
\{A\cup N:\ A\in\mathcal F,\ N\subset Z,\ Z\in\mathcal F,\ \mu(Z)=0\}.
\]

\subsubsection{Пополненная
мера}\label{ux43fux43eux43fux43eux43bux43dux435ux43dux43dux430ux44f-ux43cux435ux440ux430}

Если \(E=A\cup N\in\overline{\mathcal F}^{\,\mu}\), где
\(A\in\mathcal F\) и \(N\) лежит в измеримом нулевом множестве, то

\[
\overline{\mu}(E)=\mu(A).
\]

\subsubsection{\texorpdfstring{Лебегова
\(\sigma\)-алгебра}{Лебегова \textbackslash sigma-алгебра}}\label{ux43bux435ux431ux435ux433ux43eux432ux430-sigma-ux430ux43bux433ux435ux431ux440ux430}

Пополнение борелевской \(\sigma\)-алгебры \(\mathcal B(\mathbb R^d)\)
относительно борелевской меры Лебега:

\[
\mathcal L(\mathbb R^d)
=
\overline{\mathcal B(\mathbb R^d)}^{\,\lambda_d^B}.
\]

\subsubsection{Мера
Лебега}\label{ux43cux435ux440ux430-ux43bux435ux431ux435ux433ux430-1}

Пополнение борелевской меры Лебега. Обозначается \(\lambda_d\).

\subsection{4. Основные теоремы и
свойства}\label{ux43eux441ux43dux43eux432ux43dux44bux435-ux442ux435ux43eux440ux435ux43cux44b-ux438-ux441ux432ux43eux439ux441ux442ux432ux430-3}

\subsubsection{\texorpdfstring{Свойство 1. Борелевская
\(\sigma\)-алгебра
существует}{Свойство 1. Борелевская \textbackslash sigma-алгебра существует}}\label{ux441ux432ux43eux439ux441ux442ux432ux43e-1.-ux431ux43eux440ux435ux43bux435ux432ux441ux43aux430ux44f-sigma-ux430ux43bux433ux435ux431ux440ux430-ux441ux443ux449ux435ux441ux442ux432ux443ux435ux442}

Для любого топологического пространства \(X\) существует минимальная
\(\sigma\)-алгебра, содержащая все открытые множества.

\subsubsection{Доказательство}\label{ux434ux43eux43aux430ux437ux430ux442ux435ux43bux44cux441ux442ux432ux43e-17}

Рассмотрим все \(\sigma\)-алгебры на \(X\), содержащие все открытые
множества. Такой класс не пуст, потому что \(2^X\) - \(\sigma\)-алгебра
и содержит все подмножества \(X\). Пересечение любого семейства
\(\sigma\)-алгебр снова является \(\sigma\)-алгеброй. Поэтому

\[
\mathcal B(X)=
\bigcap\{\mathcal F:\ \mathcal F\text{ - }\sigma\text{-алгебра},\ \tau\subset\mathcal F\}
\]

является \(\sigma\)-алгеброй, содержит все открытые множества и
содержится в любой другой \(\sigma\)-алгебре с этим свойством. Значит
она минимальна. \(\square\)

\subsubsection{Свойство 2. Замкнутые множества
борелевские}\label{ux441ux432ux43eux439ux441ux442ux432ux43e-2.-ux437ux430ux43cux43aux43dux443ux442ux44bux435-ux43cux43dux43eux436ux435ux441ux442ux432ux430-ux431ux43eux440ux435ux43bux435ux432ux441ux43aux438ux435}

Если \(F\subset X\) замкнуто, то \(F\in\mathcal B(X)\).

\subsubsection{Доказательство}\label{ux434ux43eux43aux430ux437ux430ux442ux435ux43bux44cux441ux442ux432ux43e-18}

Если \(F\) замкнуто, то \(\overline{F}=X\setminus F\) открыто. Открытое
множество лежит в \(\mathcal B(X)\), а \(\sigma\)-алгебра замкнута
относительно дополнений. Значит \(F\in\mathcal B(X)\). \(\square\)

\subsubsection{\texorpdfstring{Свойство 3. В пространстве со счетной
базой борелевская \(\sigma\)-алгебра порождена
базой}{Свойство 3. В пространстве со счетной базой борелевская \textbackslash sigma-алгебра порождена базой}}\label{ux441ux432ux43eux439ux441ux442ux432ux43e-3.-ux432-ux43fux440ux43eux441ux442ux440ux430ux43dux441ux442ux432ux435-ux441ux43e-ux441ux447ux435ux442ux43dux43eux439-ux431ux430ux437ux43eux439-ux431ux43eux440ux435ux43bux435ux432ux441ux43aux430ux44f-sigma-ux430ux43bux433ux435ux431ux440ux430-ux43fux43eux440ux43eux436ux434ux435ux43dux430-ux431ux430ux437ux43eux439}

Пусть \(\{U_n\}_{n\ge1}\) - счетная база топологии \(X\). Тогда

\[
\mathcal B(X)=\sigma(U_1,U_2,\ldots).
\]

\subsubsection{Доказательство}\label{ux434ux43eux43aux430ux437ux430ux442ux435ux43bux44cux441ux442ux432ux43e-19}

Так как все \(U_n\) открыты, то

\[
\sigma(U_1,U_2,\ldots)\subset\mathcal B(X).
\]

Обратно, любое открытое множество \(G\) представляется как объединение
базовых элементов:

\[
G=\bigcup_{n:\ U_n\subset G}U_n.
\]

Так как база счетная, это счетное объединение. Значит
\(G\in\sigma(U_1,U_2,\ldots)\) для любого открытого \(G\).
Следовательно, \(\mathcal B(X)\subset\sigma(U_1,U_2,\ldots)\). Получаем
равенство. \(\square\)

\subsubsection{\texorpdfstring{Свойство 4. Порождающие классы для
\(\mathcal B(\mathbb R^d)\)}{Свойство 4. Порождающие классы для \textbackslash mathcal B(\textbackslash mathbb R\^{}d)}}\label{ux441ux432ux43eux439ux441ux442ux432ux43e-4.-ux43fux43eux440ux43eux436ux434ux430ux44eux449ux438ux435-ux43aux43bux430ux441ux441ux44b-ux434ux43bux44f-mathcal-bmathbb-rd}

В \(\mathbb R^d\):

\[
\mathcal B(\mathbb R^d)
=
\sigma(\text{открытые шары})
=
\sigma(\text{замкнутые множества})
=
\sigma(\text{полуоткрытые прямоугольники}).
\]

Можно также брать только шары или прямоугольники с рациональными
координатами.

\subsubsection{Доказательство-идея}\label{ux434ux43eux43aux430ux437ux430ux442ux435ux43bux44cux441ux442ux432ux43e-ux438ux434ux435ux44f}

Открытые шары порождают топологию. Замкнутые множества порождают ту же
\(\sigma\)-алгебру, потому что дополнение замкнутого множества открыто,
а дополнение открытого замкнуто. Прямоугольники с рациональными
координатами образуют счетную базу обычной топологии в \(\mathbb R^d\).
По свойству 3 они порождают борелевскую \(\sigma\)-алгебру.

\subsubsection{\texorpdfstring{Свойство 5. Борелевская
\(\sigma\)-алгебра на прямой порождается
лучами}{Свойство 5. Борелевская \textbackslash sigma-алгебра на прямой порождается лучами}}\label{ux441ux432ux43eux439ux441ux442ux432ux43e-5.-ux431ux43eux440ux435ux43bux435ux432ux441ux43aux430ux44f-sigma-ux430ux43bux433ux435ux431ux440ux430-ux43dux430-ux43fux440ux44fux43cux43eux439-ux43fux43eux440ux43eux436ux434ux430ux435ux442ux441ux44f-ux43bux443ux447ux430ux43cux438}

На \(\mathbb R\):

\[
\mathcal B(\mathbb R)=\sigma\{(-\infty,x]:x\in\mathbb R\}.
\]

\subsubsection{Доказательство}\label{ux434ux43eux43aux430ux437ux430ux442ux435ux43bux44cux441ux442ux432ux43e-20}

Каждый луч \((-\infty,x]\) замкнут, значит борелевский. Поэтому

\[
\sigma\{(-\infty,x]\}\subset\mathcal B(\mathbb R).
\]

Обратно, интервалы вида \((a,b]\) выражаются через лучи:

\[
(a,b]=(-\infty,b]\setminus(-\infty,a].
\]

Открытый интервал можно представить как счетное объединение:

\[
(a,b)=\bigcup_{n=1}^{\infty}\left(a,b-\frac1n\right]
\]

для достаточно больших \(n\); точнее можно брать все \(n\), для которых
\(b-\frac1n>a\). Значит открытые интервалы лежат в
\(\sigma\{(-\infty,x]\}\). Любое открытое множество на прямой является
счетным объединением открытых интервалов с рациональными концами,
поэтому оно также лежит в этой \(\sigma\)-алгебре. Следовательно,

\[
\mathcal B(\mathbb R)\subset\sigma\{(-\infty,x]\}.
\]

Получаем равенство. \(\square\)

\subsubsection{Теорема 6. Пополнение
меры}\label{ux442ux435ux43eux440ux435ux43cux430-6.-ux43fux43eux43fux43eux43bux43dux435ux43dux438ux435-ux43cux435ux440ux44b}

Пусть \((\Omega,\mathcal F,\mu)\) - пространство с мерой. Тогда

\[
\overline{\mathcal F}^{\,\mu}
=
\{A\cup N:\ A\in\mathcal F,\ N\subset Z,\ Z\in\mathcal F,\ \mu(Z)=0\}
\]

является \(\sigma\)-алгеброй, а формула

\[
\overline{\mu}(A\cup N)=\mu(A)
\]

задает на ней полную меру. Кроме того, \(\overline{\mu}\) продолжает
\(\mu\).

\subsubsection{Доказательство-идея}\label{ux434ux43eux43aux430ux437ux430ux442ux435ux43bux44cux441ux442ux432ux43e-ux438ux434ux435ux44f-1}

Если множество отличается от измеримого множества только на подмножество
нулевого множества, то его можно считать измеримым после пополнения.
Удобная эквивалентная форма:

\[
E\in\overline{\mathcal F}^{\,\mu}
\quad\Longleftrightarrow\quad
\exists A\in\mathcal F:\ E\triangle A\subset Z,\ Z\in\mathcal F,\ \mu(Z)=0,
\]

где \(E\triangle A\) - симметрическая разность. Из этой формы легко
проверить замкнутость относительно дополнений и счетных объединений.
Если

\[
E_n\triangle A_n\subset Z_n,\qquad \mu(Z_n)=0,
\]

то

\[
\left(\bigcup_n E_n\right)\triangle\left(\bigcup_n A_n\right)
\subset
\bigcup_n Z_n,
\]

а \(\bigcup_n Z_n\) имеет меру \(0\) по счетной субаддитивности. Значит
счетное объединение снова лежит в пополнении.

Корректность \(\overline{\mu}\) следует из того, что два возможных
измеримых представителя отличаются только на множестве меры \(0\),
поэтому имеют одинаковую меру. Полнота очевидна из построения: все
подмножества старых и новых нулевых множеств добавлены.

\subsubsection{Свойство 7. Борелевская мера Лебега не
полна}\label{ux441ux432ux43eux439ux441ux442ux432ux43e-7.-ux431ux43eux440ux435ux43bux435ux432ux441ux43aux430ux44f-ux43cux435ux440ux430-ux43bux435ux431ux435ux433ux430-ux43dux435-ux43fux43eux43bux43dux430}

Борелевская мера Лебега \(\lambda_d^B\) на \(\mathcal B(\mathbb R^d)\)
не является полной.

\subsubsection{Обоснование}\label{ux43eux431ux43eux441ux43dux43eux432ux430ux43dux438ux435-1}

Достаточно рассмотреть \(d=1\). Канторово множество \(C\) замкнуто,
значит борелевское, и имеет меру Лебега \(0\). Все подмножества \(C\)
имеют внешнюю меру \(0\) и входят в лебегово пополнение. Но подмножеств
у \(C\) больше, чем борелевских множеств на прямой: борелевских множеств
имеет мощность континуума, а подмножеств канторова множества -
\(2^{\mathfrak c}\). Поэтому существуют подмножества \(C\), которые не
являются борелевскими. Они лежат в лебеговом пополнении, но не лежат в
\(\mathcal B(\mathbb R)\). Значит борелевская мера Лебега не полна.

\subsubsection{\texorpdfstring{Свойство 8. Лебегова \(\sigma\)-алгебра
строго шире
борелевской}{Свойство 8. Лебегова \textbackslash sigma-алгебра строго шире борелевской}}\label{ux441ux432ux43eux439ux441ux442ux432ux43e-8.-ux43bux435ux431ux435ux433ux43eux432ux430-sigma-ux430ux43bux433ux435ux431ux440ux430-ux441ux442ux440ux43eux433ux43e-ux448ux438ux440ux435-ux431ux43eux440ux435ux43bux435ux432ux441ux43aux43eux439}

На \(\mathbb R^d\):

\[
\mathcal B(\mathbb R^d)\subsetneq\mathcal L(\mathbb R^d).
\]

При этом

\[
\mathcal L(\mathbb R^d)\ne 2^{\mathbb R^d}.
\]

Первое строгое включение следует из предыдущего свойства. Второе
означает существование неизмеримых по Лебегу множеств, например множеств
Витали.

\subsection{5. Примеры}\label{ux43fux440ux438ux43cux435ux440ux44b-3}

\subsubsection{Пример 1. Интервалы
борелевские}\label{ux43fux440ux438ux43cux435ux440-1.-ux438ux43dux442ux435ux440ux432ux430ux43bux44b-ux431ux43eux440ux435ux43bux435ux432ux441ux43aux438ux435}

На прямой открытые интервалы борелевские по определению, замкнутые
интервалы борелевские как замкнутые множества, а полуинтервалы
борелевские, потому что

\[
(a,b]=(-\infty,b]\setminus(-\infty,a].
\]

\subsubsection{Пример 2. Счетные множества
борелевские}\label{ux43fux440ux438ux43cux435ux440-2.-ux441ux447ux435ux442ux43dux44bux435-ux43cux43dux43eux436ux435ux441ux442ux432ux430-ux431ux43eux440ux435ux43bux435ux432ux441ux43aux438ux435}

Каждая точка \(\{x\}\) замкнута, значит борелевская. Если
\(A=\{x_1,x_2,\ldots\}\) счетно, то

\[
A=\bigcup_{n=1}^{\infty}\{x_n\},
\]

поэтому \(A\) борелевско. В частности, \(\mathbb Q\) и
\(\mathbb Q\cap[0,1]\) борелевские.

\subsubsection{Пример 3. Счетные множества имеют меру Лебега
ноль}\label{ux43fux440ux438ux43cux435ux440-3.-ux441ux447ux435ux442ux43dux44bux435-ux43cux43dux43eux436ux435ux441ux442ux432ux430-ux438ux43cux435ux44eux442-ux43cux435ux440ux443-ux43bux435ux431ux435ux433ux430-ux43dux43eux43bux44c}

Так как \(\lambda(\{x\})=0\), то

\[
\lambda(\mathbb Q\cap[0,1])
\le
\sum_{q\in\mathbb Q\cap[0,1]}\lambda(\{q\})
=0.
\]

Значит

\[
\lambda(\mathbb Q\cap[0,1])=0.
\]

\subsubsection{Пример 4. Подмножества нулевых множеств после
пополнения}\label{ux43fux440ux438ux43cux435ux440-4.-ux43fux43eux434ux43cux43dux43eux436ux435ux441ux442ux432ux430-ux43dux443ux43bux435ux432ux44bux445-ux43cux43dux43eux436ux435ux441ux442ux432-ux43fux43eux441ux43bux435-ux43fux43eux43fux43eux43bux43dux435ux43dux438ux44f}

Пусть \(C\) - канторово множество. Оно замкнуто, значит борелевское, и

\[
\lambda(C)=0.
\]

Любое \(N\subset C\) входит в лебегово пополнение и имеет меру \(0\). Но
не каждое такое \(N\) борелевское.

\subsubsection{Пример 5. Неизмеримое
множество}\label{ux43fux440ux438ux43cux435ux440-5.-ux43dux435ux438ux437ux43cux435ux440ux438ux43cux43eux435-ux43cux43dux43eux436ux435ux441ux442ux432ux43e}

Множество Витали \(V\subset[0,1]\) не является лебегово измеримым.
Значит даже после пополнения меры Лебега область определения не
становится равной \(2^{\mathbb R}\).

\subsection{6. Доказательства и
обоснования}\label{ux434ux43eux43aux430ux437ux430ux442ux435ux43bux44cux441ux442ux432ux430-ux438-ux43eux431ux43eux441ux43dux43eux432ux430ux43dux438ux44f-3}

\textbf{Почему борелевская \(\sigma\)-алгебра зависит от топологии.}

Определение

\[
\mathcal B(X)=\sigma(\text{открытые множества})
\]

использует открытые множества. Если на одном и том же \(X\) взять другую
топологию, то изменится класс открытых множеств, а значит может
измениться и борелевская \(\sigma\)-алгебра.

\textbf{Почему в \(\mathbb R^d\) достаточно рациональных
прямоугольников.}

Обычная топология в \(\mathbb R^d\) имеет счетную базу из
прямоугольников с рациональными координатами. Любое открытое множество
является объединением всех таких базовых прямоугольников, которые в нем
лежат. Так как база счетна, это счетное объединение. Поэтому
\(\sigma\)-алгебра, порожденная рациональными прямоугольниками, уже
содержит все открытые множества.

\textbf{Почему пополнение не меняет старую меру.}

Если \(A\in\mathcal F\), то можно представить \(A=A\cup\varnothing\),
где \(\varnothing\) лежит в любом нулевом множестве. Тогда

\[
\overline{\mu}(A)=\mu(A).
\]

Поэтому пополнение является именно продолжением меры, а не новой мерой с
другими значениями на старых множествах.

\textbf{Почему добавляются только подмножества нулевых множеств.}

Если \(N\subset Z\) и \(\mu(Z)=0\), то по монотонности любая разумная
мера должна давать

\[
0\le\mu(N)\le\mu(Z)=0.
\]

Значит единственно возможное значение для \(N\) - ноль. Поэтому такие
множества можно добавить без противоречия. Для произвольных подмножеств
положительных множеств такой аргумент не работает.

\textbf{Почему не все лебегово измеримые множества борелевские.}

Пополнение добавляет все подмножества борелевских нулевых множеств. У
канторова множества мера \(0\), но подмножеств у него больше, чем всех
борелевских множеств на прямой. Поэтому среди этих подмножеств есть
неборелевские, но после пополнения они становятся лебегово измеримыми.

\subsection{7. Что написать на
доске}\label{ux447ux442ux43e-ux43dux430ux43fux438ux441ux430ux442ux44c-ux43dux430-ux434ux43eux441ux43aux435-3}

\begin{enumerate}
\def\labelenumi{\arabic{enumi}.}
\tightlist
\item
  Борелевская \(\sigma\)-алгебра:
\end{enumerate}

\[
\mathcal B(X)=\sigma\{G\subset X:\ G\text{ открыто}\}.
\]

\begin{enumerate}
\def\labelenumi{\arabic{enumi}.}
\setcounter{enumi}{1}
\tightlist
\item
  В \(\mathbb R^d\):
\end{enumerate}

\[
\mathcal B(\mathbb R^d)
=
\sigma(\text{открытые шары})
=
\sigma(\text{полуоткрытые прямоугольники}).
\]

\begin{enumerate}
\def\labelenumi{\arabic{enumi}.}
\setcounter{enumi}{2}
\tightlist
\item
  Борелевская мера Лебега:
\end{enumerate}

\[
\lambda_d^B\left(\prod_{j=1}^d[a_j,b_j)\right)
=
\prod_{j=1}^d(b_j-a_j).
\]

\begin{enumerate}
\def\labelenumi{\arabic{enumi}.}
\setcounter{enumi}{3}
\tightlist
\item
  Полная мера:
\end{enumerate}

\[
N\subset Z,\quad Z\in\mathcal F,\quad \mu(Z)=0
\Rightarrow
N\in\mathcal F.
\]

\begin{enumerate}
\def\labelenumi{\arabic{enumi}.}
\setcounter{enumi}{4}
\tightlist
\item
  Пополнение:
\end{enumerate}

\[
\overline{\mathcal F}^{\,\mu}
=
\{A\cup N:\ A\in\mathcal F,\ N\subset Z,\ Z\in\mathcal F,\ \mu(Z)=0\}.
\]

\begin{enumerate}
\def\labelenumi{\arabic{enumi}.}
\setcounter{enumi}{5}
\tightlist
\item
  Борелевская и лебегова \(\sigma\)-алгебры:
\end{enumerate}

\[
\mathcal B(\mathbb R^d)\subsetneq\mathcal L(\mathbb R^d)\subsetneq 2^{\mathbb R^d}.
\]

\subsection{8. Типичные дополнительные
вопросы}\label{ux442ux438ux43fux438ux447ux43dux44bux435-ux434ux43eux43fux43eux43bux43dux438ux442ux435ux43bux44cux43dux44bux435-ux432ux43eux43fux440ux43eux441ux44b-3}

\textbf{Вопрос. Что такое борелевское множество?}

Это элемент минимальной \(\sigma\)-алгебры, содержащей все открытые
множества:

\[
A\in\mathcal B(X).
\]

\textbf{Вопрос. Почему замкнутые множества борелевские?}

Потому что дополнение замкнутого множества открыто, а \(\sigma\)-алгебра
замкнута относительно дополнений.

\textbf{Вопрос. Почему в \(\mathbb R\) борелевская \(\sigma\)-алгебра
порождается лучами \((-\infty,x]\)?}

Потому что из лучей получаются полуинтервалы \((a,b]\), из них -
открытые интервалы, а открытые интервалы с рациональными концами
порождают обычную топологию прямой.

\textbf{Вопрос. Что значит, что мера полная?}

Это значит, что все подмножества измеримых множеств меры ноль тоже
измеримы.

\textbf{Вопрос. Чем борелевская мера Лебега отличается от меры Лебега?}

Борелевская мера Лебега задана только на \(\mathcal B(\mathbb R^d)\).
Полная мера Лебега задана на большей \(\sigma\)-алгебре
\(\mathcal L(\mathbb R^d)\), полученной пополнением.

\textbf{Вопрос. Может ли лебегово измеримое множество быть
неборелевским?}

Да. Например, некоторые подмножества канторова множества лебегово
измеримы, потому что лежат в множестве меры ноль, но не являются
борелевскими.

\textbf{Вопрос. Все ли подмножества \(\mathbb R\) становятся измеримыми
после пополнения?}

Нет. Существуют неизмеримые по Лебегу множества, например множества
Витали.

\subsection{9. Типичные
ошибки}\label{ux442ux438ux43fux438ux447ux43dux44bux435-ux43eux448ux438ux431ux43aux438-3}

\begin{enumerate}
\def\labelenumi{\arabic{enumi}.}
\tightlist
\item
  Считать борелевские и лебегово измеримые множества одним и тем же
  классом. Правильно:
\end{enumerate}

\[
\mathcal B(\mathbb R^d)\subsetneq\mathcal L(\mathbb R^d).
\]

\begin{enumerate}
\def\labelenumi{\arabic{enumi}.}
\setcounter{enumi}{1}
\item
  Не уточнять топологию в определении \(\mathcal B(X)\).
\item
  Думать, что пополнение добавляет все подмножества пространства. Оно
  добавляет только подмножества множеств меры ноль и их объединения со
  старыми измеримыми множествами.
\item
  Путать внешнюю меру и пополненную меру. Внешняя мера определена на
  всех подмножествах, но не является мерой на всех подмножествах.
\item
  Забывать, что борелевская мера Лебега не полна.
\item
  Думать, что всякое множество меры ноль счетно. Канторово множество
  несчетно и имеет меру ноль.
\item
  Использовать дополнение через верхний индекс. В этих конспектах:
\end{enumerate}

\[
\overline{A}=X\setminus A.
\]

\subsection{10. Связи с другими
билетами}\label{ux441ux432ux44fux437ux438-ux441-ux434ux440ux443ux433ux438ux43cux438-ux431ux438ux43bux435ux442ux430ux43cux438-3}

\textbf{Билет 01.} Борелевская \(\sigma\)-алгебра - важнейший пример
\(\sigma\)-алгебры.

\textbf{Билет 02.} Распределения случайных величин являются
вероятностными мерами на борелевских \(\sigma\)-алгебрах.

\textbf{Билет 03.} Борелевская мера Лебега строится как продолжение меры
с полуалгебры прямоугольников, а полная мера Лебега получается
пополнением.

\textbf{Билет 05.} Измеримость случайной величины означает, что
прообразы борелевских множеств измеримы.

\textbf{Билет 06.} Распределения случайных величин и векторов задаются
на \(\mathcal B(\mathbb R)\) и \(\mathcal B(\mathbb R^d)\).

\textbf{Билет 07.} Интеграл Лебега обычно строится по полной мере Лебега
или по абстрактной полной мере.

\subsection{11. Минимум на
удовлетворительно}\label{ux43cux438ux43dux438ux43cux443ux43c-ux43dux430-ux443ux434ux43eux432ux43bux435ux442ux432ux43eux440ux438ux442ux435ux43bux44cux43dux43e-3}

Нужно уметь сказать:

\begin{enumerate}
\def\labelenumi{\arabic{enumi}.}
\tightlist
\item
  Борелевская \(\sigma\)-алгебра:
\end{enumerate}

\[
\mathcal B(X)=\sigma(\text{открытые множества}).
\]

\begin{enumerate}
\def\labelenumi{\arabic{enumi}.}
\setcounter{enumi}{1}
\tightlist
\item
  Борелевская мера - мера на \(\mathcal B(X)\).
\item
  Борелевская мера Лебега на \(\mathbb R^d\) совпадает с объемом на
  прямоугольниках.
\item
  Полная мера содержит все подмножества множеств меры ноль.
\item
  Полная мера Лебега - пополнение борелевской меры Лебега.
\end{enumerate}

\subsection{12. Минимум на
хорошо}\label{ux43cux438ux43dux438ux43cux443ux43c-ux43dux430-ux445ux43eux440ux43eux448ux43e-3}

Кроме минимума на удовлетворительно, нужно:

\begin{enumerate}
\def\labelenumi{\arabic{enumi}.}
\tightlist
\item
  Доказать существование \(\mathcal B(X)\) как пересечения всех
  \(\sigma\)-алгебр, содержащих открытые множества.
\item
  Объяснить, почему замкнутые множества борелевские.
\item
  Назвать порождающие классы для \(\mathcal B(\mathbb R^d)\): шары,
  рациональные прямоугольники, полуоткрытые прямоугольники.
\item
  Сформулировать пополнение:
\end{enumerate}

\[
\overline{\mathcal F}^{\,\mu}
=
\{A\cup N:\ A\in\mathcal F,\ N\subset Z,\ \mu(Z)=0\}.
\]

\begin{enumerate}
\def\labelenumi{\arabic{enumi}.}
\setcounter{enumi}{4}
\tightlist
\item
  Объяснить различие между борелевскими и лебегово измеримыми
  множествами.
\end{enumerate}

\subsection{13. Что добавить для
отлично}\label{ux447ux442ux43e-ux434ux43eux431ux430ux432ux438ux442ux44c-ux434ux43bux44f-ux43eux442ux43bux438ux447ux43dux43e-3}

Для сильного ответа стоит добавить:

\begin{enumerate}
\def\labelenumi{\arabic{enumi}.}
\tightlist
\item
  Доказательство, что в пространстве со счетной базой борелевская
  \(\sigma\)-алгебра порождена базой.
\item
  Доказательство, что \(\mathcal B(\mathbb R)=\sigma\{(-\infty,x]\}\).
\item
  Эквивалентное описание пополнения через симметрическую разность:
\end{enumerate}

\[
E\triangle A\subset Z,\qquad \mu(Z)=0.
\]

\begin{enumerate}
\def\labelenumi{\arabic{enumi}.}
\setcounter{enumi}{3}
\tightlist
\item
  Обоснование, почему борелевская мера Лебега не полна, через
  подмножества канторова множества.
\item
  Цепочку:
\end{enumerate}

\[
\mathcal B(\mathbb R^d)\subsetneq\mathcal L(\mathbb R^d)\subsetneq 2^{\mathbb R^d}.
\]

\subsection{14. Устная версия ответа на 5
минут}\label{ux443ux441ux442ux43dux430ux44f-ux432ux435ux440ux441ux438ux44f-ux43eux442ux432ux435ux442ux430-ux43dux430-5-ux43cux438ux43dux443ux442-3}

Борелевская \(\sigma\)-алгебра строится из топологии. Если \(X\) -
топологическое пространство, то

\[
\mathcal B(X)=\sigma\{G\subset X:\ G\text{ открыто}\}.
\]

Это минимальная \(\sigma\)-алгебра, содержащая все открытые множества.
Ее элементы называются борелевскими множествами. Замкнутые множества
тоже борелевские, потому что они являются дополнениями открытых
множеств.

В \(\mathbb R^d\) борелевскую \(\sigma\)-алгебру можно порождать разными
классами: открытыми шарами, замкнутыми множествами, прямоугольниками,
полуоткрытыми прямоугольниками. Если взять прямоугольники с
рациональными координатами, они образуют счетную базу топологии, поэтому
тоже порождают \(\mathcal B(\mathbb R^d)\).

Борелевская мера - это мера на \(\mathcal B(X)\). Борелевская мера
Лебега на \(\mathbb R^d\) - это мера, которая на полуоткрытых
прямоугольниках равна объему:

\[
\lambda_d^B\left(\prod_{j=1}^d[a_j,b_j)\right)
=
\prod_{j=1}^d(b_j-a_j).
\]

Но борелевская мера Лебега не полна: могут быть подмножества борелевских
множеств меры ноль, которые сами не являются борелевскими.

Поэтому вводят пополнение меры. Мера \(\mu\) на \((\Omega,\mathcal F)\)
называется полной, если

\[
N\subset Z,\quad Z\in\mathcal F,\quad \mu(Z)=0
\Rightarrow N\in\mathcal F.
\]

Пополнение добавляет к \(\mathcal F\) все множества вида

\[
A\cup N,
\]

где \(A\in\mathcal F\), а \(N\) лежит в измеримом множестве меры ноль.
Формально:

\[
\overline{\mathcal F}^{\,\mu}
=
\{A\cup N:\ A\in\mathcal F,\ N\subset Z,\ Z\in\mathcal F,\ \mu(Z)=0\}.
\]

Полная мера Лебега - это пополнение борелевской меры Лебега. Поэтому

\[
\mathcal B(\mathbb R^d)\subsetneq\mathcal L(\mathbb R^d)\subsetneq 2^{\mathbb R^d}.
\]

Первое включение строго, потому что есть лебегово измеримые
неборелевские множества, например некоторые подмножества канторова
множества. Второе включение строго, потому что существуют неизмеримые
множества, например множества Витали.

\subsection{15. Если осталось 2 минуты перед
экзаменом}\label{ux435ux441ux43bux438-ux43eux441ux442ux430ux43bux43eux441ux44c-2-ux43cux438ux43dux443ux442ux44b-ux43fux435ux440ux435ux434-ux44dux43aux437ux430ux43cux435ux43dux43eux43c-3}

Запомнить:

\[
\mathcal B(X)=\sigma(\text{открытые множества}).
\]

Борелевская мера - мера на \(\mathcal B(X)\).

Борелевская мера Лебега:

\[
\lambda_d^B\left(\prod_{j=1}^d[a_j,b_j)\right)
=
\prod_{j=1}^d(b_j-a_j).
\]

Полная мера:

\[
N\subset Z,\quad \mu(Z)=0
\Rightarrow
N\text{ измеримо}.
\]

Пополнение:

\[
\overline{\mathcal F}^{\,\mu}
=
\{A\cup N:\ A\in\mathcal F,\ N\subset Z,\ Z\in\mathcal F,\ \mu(Z)=0\}.
\]

Главная мысль:

\[
\mathcal B(\mathbb R^d)\subsetneq\mathcal L(\mathbb R^d)\subsetneq 2^{\mathbb R^d}.
\]

Борелевские множества - не то же самое, что все лебегово измеримые
множества.

\newpage

\section{Билет 05. Измеримые функции, отображения и случайные
элементы}\label{ux431ux438ux43bux435ux442-05.-ux438ux437ux43cux435ux440ux438ux43cux44bux435-ux444ux443ux43dux43aux446ux438ux438-ux43eux442ux43eux431ux440ux430ux436ux435ux43dux438ux44f-ux438-ux441ux43bux443ux447ux430ux439ux43dux44bux435-ux44dux43bux435ux43cux435ux43dux442ux44b}

Источники: Ширяев 1, гл. II; Сердобольская, лекция 1; лекции ДСП.

\subsection{1. Короткий
ответ}\label{ux43aux43eux440ux43eux442ux43aux438ux439-ux43eux442ux432ux435ux442-4}

Измеримое отображение - это отображение между измеримыми пространствами,
у которого прообраз любого измеримого множества снова измерим. Если

\[
X:(\Omega,\mathcal F)\to(E,\mathcal E),
\]

то измеримость означает

\[
X^{-1}(B)=\{\omega\in\Omega:X(\omega)\in B\}\in\mathcal F
\qquad \forall B\in\mathcal E.
\]

Случайная величина - это измеримое отображение из вероятностного
пространства в \(\mathbb R\) с борелевской \(\sigma\)-алгеброй:

\[
X:(\Omega,\mathcal F,P)\to(\mathbb R,\mathcal B(\mathbb R)).
\]

Случайный вектор - измеримое отображение в \(\mathbb R^d\), случайный
элемент - измеримое отображение в произвольное измеримое, обычно
метрическое, пространство. Измеримость нужна для того, чтобы события
вида \(\{X\in B\}\) имели вероятность.

\subsection{2. Полный
ответ}\label{ux43fux43eux43bux43dux44bux439-ux43eux442ux432ux435ux442-4}

\subsubsection{2.1. Мотивация и связь с наблюдаемыми
величинами}\label{ux43cux43eux442ux438ux432ux430ux446ux438ux44f-ux438-ux441ux432ux44fux437ux44c-ux441-ux43dux430ux431ux43bux44eux434ux430ux435ux43cux44bux43cux438-ux432ux435ux43bux438ux447ux438ux43dux430ux43cux438}

В предыдущих билетах вводились \(\sigma\)-алгебры, меры, вероятностные
пространства и борелевские множества. Теперь нужно связать пространство
элементарных исходов с наблюдаемыми величинами. В вероятности мы редко
работаем напрямую с исходом \(\omega\in\Omega\). Обычно нас интересует
численная или более сложная характеристика исхода: результат измерения,
вектор наблюдений, траектория процесса, последовательность значений.

\subsubsection{2.2. Определение измеримого
отображения}\label{ux43eux43fux440ux435ux434ux435ux43bux435ux43dux438ux435-ux438ux437ux43cux435ux440ux438ux43cux43eux433ux43e-ux43eux442ux43eux431ux440ux430ux436ux435ux43dux438ux44f}

Пусть есть измеримое пространство \((\Omega,\mathcal F)\) и другое
измеримое пространство \((E,\mathcal E)\). Отображение

\[
X:\Omega\to E
\]

называется измеримым, если для любого измеримого множества
\(B\in\mathcal E\) его прообраз измерим:

\[
X^{-1}(B)\in\mathcal F.
\]

Интуитивно это означает: если в пространстве значений \(E\) можно задать
наблюдаемое событие \(B\), то событие ``значение \(X\) попало в \(B\)''
должно быть событием в исходном вероятностном пространстве. Иначе
выражение \(P\{X\in B\}\) просто не будет определено.

\subsubsection{2.3. Вещественная случайная
величина}\label{ux432ux435ux449ux435ux441ux442ux432ux435ux43dux43dux430ux44f-ux441ux43bux443ux447ux430ux439ux43dux430ux44f-ux432ux435ux43bux438ux447ux438ux43dux430}

Главный частный случай - вещественная случайная величина. Это измеримое
отображение

\[
X:(\Omega,\mathcal F,P)\to(\mathbb R,\mathcal B(\mathbb R)).
\]

По определению для любого борелевского множества \(B\subset\mathbb R\)
должно быть

\[
\{X\in B\}\in\mathcal F.
\]

Проверять все борелевские множества напрямую неудобно. Поэтому
используют порождающие классы. Так как

\[
\mathcal B(\mathbb R)=\sigma\{(-\infty,x]:x\in\mathbb R\},
\]

то для вещественной функции достаточно проверить, что

\[
\{X\le x\}\in\mathcal F\qquad \forall x\in\mathbb R.
\]

Эквивалентно можно проверять множества \(\{X<x\}\), \(\{X>x\}\) или
\(\{X\ge x\}\) для всех \(x\).

\subsubsection{2.4. Случайный вектор и критерий
измеримости}\label{ux441ux43bux443ux447ux430ux439ux43dux44bux439-ux432ux435ux43aux442ux43eux440-ux438-ux43aux440ux438ux442ux435ux440ux438ux439-ux438ux437ux43cux435ux440ux438ux43cux43eux441ux442ux438}

Случайный вектор

\[
X=(X_1,\ldots,X_d)
\]

со значениями в \(\mathbb R^d\) - это измеримое отображение

\[
X:(\Omega,\mathcal F,P)\to(\mathbb R^d,\mathcal B(\mathbb R^d)).
\]

Важный критерий: \(X\) измерим тогда и только тогда, когда все его
координаты \(X_1,\ldots,X_d\) измеримы как вещественные случайные
величины. Это удобно, потому что на практике вектор обычно задается
координатами.

\subsubsection{2.5. Случайный элемент в метрическом
пространстве}\label{ux441ux43bux443ux447ux430ux439ux43dux44bux439-ux44dux43bux435ux43cux435ux43dux442-ux432-ux43cux435ux442ux440ux438ux447ux435ux441ux43aux43eux43c-ux43fux440ux43eux441ux442ux440ux430ux43dux441ux442ux432ux435}

Еще более общий объект - случайный элемент. Если \((E,\rho)\) -
метрическое пространство, то случайным элементом со значениями в \(E\)
называют измеримое отображение

\[
X:(\Omega,\mathcal F,P)\to(E,\mathcal B(E)),
\]

где \(\mathcal B(E)\) - борелевская \(\sigma\)-алгебра метрического
пространства. Примеры: случайная точка в \(\mathbb R^d\), случайная
последовательность в \(\mathbb R^{\mathbb N}\), случайная функция в
пространстве траекторий.

\subsubsection{2.6. Измеримость случайных
последовательностей}\label{ux438ux437ux43cux435ux440ux438ux43cux43eux441ux442ux44c-ux441ux43bux443ux447ux430ux439ux43dux44bux445-ux43fux43eux441ux43bux435ux434ux43eux432ux430ux442ux435ux43bux44cux43dux43eux441ux442ux435ux439}

Для дискретных случайных процессов это особенно важно.
Последовательность случайных величин \((X_1,X_2,\ldots)\) можно
рассматривать двояко:

\begin{enumerate}
\def\labelenumi{\arabic{enumi}.}
\tightlist
\item
  как семейство измеримых отображений \(X_n:\Omega\to\mathbb R\);
\item
  как одно отображение
\end{enumerate}

\[
X:\Omega\to\mathbb R^{\mathbb N},\qquad
X(\omega)=(X_1(\omega),X_2(\omega),\ldots).
\]

Если на \(\mathbb R^{\mathbb N}\) взять цилиндрическую, то есть
произведенную, \(\sigma\)-алгебру, то такое отображение измеримо тогда и
только тогда, когда измеримы все координаты \(X_n\). Это будет
использоваться в билетах про конечномерные распределения, теорему
Колмогорова и дискретные случайные процессы.

\subsubsection{2.7. Роль условия измеримости в теории
вероятностей}\label{ux440ux43eux43bux44c-ux443ux441ux43bux43eux432ux438ux44f-ux438ux437ux43cux435ux440ux438ux43cux43eux441ux442ux438-ux432-ux442ux435ux43eux440ux438ux438-ux432ux435ux440ux43eux44fux442ux43dux43eux441ux442ux435ux439}

Итак, роль билета такая: измеримость - минимальное техническое условие,
которое превращает функцию от исхода в корректный вероятностный объект.
После этого можно задавать распределение \(P_X(B)=P\{X\in B\}\),
интегрировать функции от \(X\), говорить о моментах, независимости,
условном ожидании и процессах.

\subsection{3. Основные
определения}\label{ux43eux441ux43dux43eux432ux43dux44bux435-ux43eux43fux440ux435ux434ux435ux43bux435ux43dux438ux44f-4}

\subsubsection{Измеримое
пространство}\label{ux438ux437ux43cux435ux440ux438ux43cux43eux435-ux43fux440ux43eux441ux442ux440ux430ux43dux441ux442ux432ux43e-1}

Пара \((\Omega,\mathcal F)\), где \(\Omega\) - множество, а
\(\mathcal F\) - \(\sigma\)-алгебра его подмножеств.

\subsubsection{Прообраз
множества}\label{ux43fux440ux43eux43eux431ux440ux430ux437-ux43cux43dux43eux436ux435ux441ux442ux432ux430}

Если \(X:\Omega\to E\) и \(B\subset E\), то

\[
X^{-1}(B)=\{\omega\in\Omega:X(\omega)\in B\}.
\]

Важно: прообраз всегда является подмножеством \(\Omega\), а не
подмножеством \(E\).

\subsubsection{Измеримое
отображение}\label{ux438ux437ux43cux435ux440ux438ux43cux43eux435-ux43eux442ux43eux431ux440ux430ux436ux435ux43dux438ux435}

Пусть \((\Omega,\mathcal F)\) и \((E,\mathcal E)\) - измеримые
пространства. Отображение

\[
X:\Omega\to E
\]

называется \(\mathcal F/\mathcal E\)-измеримым, если

\[
X^{-1}(B)\in\mathcal F
\qquad \forall B\in\mathcal E.
\]

\subsubsection{Борелевски измеримая
функция}\label{ux431ux43eux440ux435ux43bux435ux432ux441ux43aux438-ux438ux437ux43cux435ux440ux438ux43cux430ux44f-ux444ux443ux43dux43aux446ux438ux44f}

Если \(E\) - топологическое или метрическое пространство и
\(\mathcal E=\mathcal B(E)\), то измеримое отображение в
\((E,\mathcal B(E))\) называют борелевски измеримым.

\subsubsection{Вещественная случайная
величина}\label{ux432ux435ux449ux435ux441ux442ux432ux435ux43dux43dux430ux44f-ux441ux43bux443ux447ux430ux439ux43dux430ux44f-ux432ux435ux43bux438ux447ux438ux43dux430-1}

Измеримое отображение

\[
X:(\Omega,\mathcal F,P)\to(\mathbb R,\mathcal B(\mathbb R)).
\]

То есть для любого \(B\in\mathcal B(\mathbb R)\) событие \(\{X\in B\}\)
принадлежит \(\mathcal F\).

\subsubsection{Случайный
вектор}\label{ux441ux43bux443ux447ux430ux439ux43dux44bux439-ux432ux435ux43aux442ux43eux440}

Измеримое отображение

\[
X=(X_1,\ldots,X_d):(\Omega,\mathcal F,P)\to(\mathbb R^d,\mathcal B(\mathbb R^d)).
\]

\subsubsection{Случайный элемент метрического
пространства}\label{ux441ux43bux443ux447ux430ux439ux43dux44bux439-ux44dux43bux435ux43cux435ux43dux442-ux43cux435ux442ux440ux438ux447ux435ux441ux43aux43eux433ux43e-ux43fux440ux43eux441ux442ux440ux430ux43dux441ux442ux432ux430}

Если \((E,\rho)\) - метрическое пространство, то случайный элемент со
значениями в \(E\) - это измеримое отображение

\[
X:(\Omega,\mathcal F,P)\to(E,\mathcal B(E)).
\]

\subsubsection{Случайная
последовательность}\label{ux441ux43bux443ux447ux430ux439ux43dux430ux44f-ux43fux43eux441ux43bux435ux434ux43eux432ux430ux442ux435ux43bux44cux43dux43eux441ux442ux44c}

Случайный элемент пространства последовательностей:

\[
X:\Omega\to\mathbb R^{\mathbb N},\qquad
X(\omega)=(X_1(\omega),X_2(\omega),\ldots),
\]

обычно с произведенной \(\sigma\)-алгеброй

\[
\mathcal B(\mathbb R)^{\otimes\mathbb N}
=
\sigma\{\pi_n^{-1}(B):n\in\mathbb N,\ B\in\mathcal B(\mathbb R)\},
\]

где \(\pi_n(x_1,x_2,\ldots)=x_n\) - координатная проекция.

\subsubsection{Сечение
процесса}\label{ux441ux435ux447ux435ux43dux438ux435-ux43fux440ux43eux446ux435ux441ux441ux430}

Для случайного процесса \(\{X_t:t\in T\}\) фиксированное значение
\(X_t(\omega)\) как функция \(\omega\) называется сечением процесса в
момент \(t\).

\subsubsection{Траектория
процесса}\label{ux442ux440ux430ux435ux43aux442ux43eux440ux438ux44f-ux43fux440ux43eux446ux435ux441ux441ux430}

Для фиксированного \(\omega\) функция

\[
t\mapsto X_t(\omega)
\]

называется траекторией процесса.

\subsection{4. Основные теоремы и
свойства}\label{ux43eux441ux43dux43eux432ux43dux44bux435-ux442ux435ux43eux440ux435ux43cux44b-ux438-ux441ux432ux43eux439ux441ux442ux432ux430-4}

\subsubsection{Свойство 1. Достаточно проверять порождающий
класс}\label{ux441ux432ux43eux439ux441ux442ux432ux43e-1.-ux434ux43eux441ux442ux430ux442ux43eux447ux43dux43e-ux43fux440ux43eux432ux435ux440ux44fux442ux44c-ux43fux43eux440ux43eux436ux434ux430ux44eux449ux438ux439-ux43aux43bux430ux441ux441}

Пусть \((\Omega,\mathcal F)\) и \((E,\mathcal E)\) - измеримые
пространства, а \(\mathcal E=\sigma(\mathcal C)\). Тогда отображение
\(X:\Omega\to E\) измеримо тогда и только тогда, когда

\[
X^{-1}(C)\in\mathcal F
\qquad \forall C\in\mathcal C.
\]

\subsubsection{Доказательство}\label{ux434ux43eux43aux430ux437ux430ux442ux435ux43bux44cux441ux442ux432ux43e-21}

Рассмотрим класс

\[
\mathcal D=\{B\subset E:X^{-1}(B)\in\mathcal F\}.
\]

Этот класс является \(\sigma\)-алгеброй. Действительно,

\[
X^{-1}(E)=\Omega\in\mathcal F,
\]

\[
X^{-1}(E\setminus B)=\Omega\setminus X^{-1}(B),
\]

и

\[
X^{-1}\left(\bigcup_{n=1}^{\infty}B_n\right)
=
\bigcup_{n=1}^{\infty}X^{-1}(B_n).
\]

Если \(X^{-1}(C)\in\mathcal F\) для всех \(C\in\mathcal C\), то
\(\mathcal C\subset\mathcal D\). Так как \(\mathcal D\) -
\(\sigma\)-алгебра, она содержит \(\sigma(\mathcal C)=\mathcal E\).
Значит \(X\) измеримо. Обратное направление следует прямо из
определения. \(\square\)

\subsubsection{Свойство 2. Критерии измеримости вещественной
функции}\label{ux441ux432ux43eux439ux441ux442ux432ux43e-2.-ux43aux440ux438ux442ux435ux440ux438ux438-ux438ux437ux43cux435ux440ux438ux43cux43eux441ux442ux438-ux432ux435ux449ux435ux441ux442ux432ux435ux43dux43dux43eux439-ux444ux443ux43dux43aux446ux438ux438}

Для функции \(X:\Omega\to\mathbb R\) следующие условия эквивалентны:

\begin{enumerate}
\def\labelenumi{\arabic{enumi}.}
\tightlist
\item
  \(X\) измерима как отображение в
  \((\mathbb R,\mathcal B(\mathbb R))\);
\item
  \(\{X\le x\}\in\mathcal F\) для всех \(x\in\mathbb R\);
\item
  \(\{X<x\}\in\mathcal F\) для всех \(x\in\mathbb R\);
\item
  \(\{X>x\}\in\mathcal F\) для всех \(x\in\mathbb R\);
\item
  \(\{X\ge x\}\in\mathcal F\) для всех \(x\in\mathbb R\).
\end{enumerate}

\subsubsection{Доказательство-идея}\label{ux434ux43eux43aux430ux437ux430ux442ux435ux43bux44cux441ux442ux432ux43e-ux438ux434ux435ux44f-2}

Борелевская \(\sigma\)-алгебра на прямой порождается лучами:

\[
\mathcal B(\mathbb R)=\sigma\{(-\infty,x]:x\in\mathbb R\}.
\]

По свойству 1 достаточно проверять прообразы таких лучей:

\[
X^{-1}((-\infty,x])=\{X\le x\}.
\]

Остальные варианты получаются через дополнения и счетные пересечения:

\[
\{X<x\}=\bigcup_{n=1}^{\infty}\{X\le x-\tfrac1n\},
\]

\[
\{X\ge x\}=\Omega\setminus\{X<x\}.
\]

\subsubsection{Свойство 3. Для борелевской измеримости достаточно
открытых
множеств}\label{ux441ux432ux43eux439ux441ux442ux432ux43e-3.-ux434ux43bux44f-ux431ux43eux440ux435ux43bux435ux432ux441ux43aux43eux439-ux438ux437ux43cux435ux440ux438ux43cux43eux441ux442ux438-ux434ux43eux441ux442ux430ux442ux43eux447ux43dux43e-ux43eux442ux43aux440ux44bux442ux44bux445-ux43cux43dux43eux436ux435ux441ux442ux432}

Пусть \(E\) - топологическое пространство. Отображение

\[
X:(\Omega,\mathcal F)\to(E,\mathcal B(E))
\]

измеримо тогда и только тогда, когда

\[
X^{-1}(G)\in\mathcal F
\]

для любого открытого множества \(G\subset E\).

\subsubsection{Доказательство}\label{ux434ux43eux43aux430ux437ux430ux442ux435ux43bux44cux441ux442ux432ux43e-22}

По определению

\[
\mathcal B(E)=\sigma\{G\subset E:G\text{ открыто}\}.
\]

Остается применить свойство 1.

\subsubsection{Свойство 4. Непрерывные отображения борелевски
измеримы}\label{ux441ux432ux43eux439ux441ux442ux432ux43e-4.-ux43dux435ux43fux440ux435ux440ux44bux432ux43dux44bux435-ux43eux442ux43eux431ux440ux430ux436ux435ux43dux438ux44f-ux431ux43eux440ux435ux43bux435ux432ux441ux43aux438-ux438ux437ux43cux435ux440ux438ux43cux44b}

Пусть \(E\) и \(S\) - топологические пространства, а

\[
f:E\to S
\]

непрерывно. Тогда \(f\) измеримо как отображение

\[
(E,\mathcal B(E))\to(S,\mathcal B(S)).
\]

\subsubsection{Доказательство}\label{ux434ux43eux43aux430ux437ux430ux442ux435ux43bux44cux441ux442ux432ux43e-23}

Для любого открытого \(G\subset S\) прообраз \(f^{-1}(G)\) открыт в
\(E\), значит принадлежит \(\mathcal B(E)\). По свойству 3 отображение
\(f\) борелевски измеримо. \(\square\)

\subsubsection{Свойство 5. Композиция измеримых отображений
измерима}\label{ux441ux432ux43eux439ux441ux442ux432ux43e-5.-ux43aux43eux43cux43fux43eux437ux438ux446ux438ux44f-ux438ux437ux43cux435ux440ux438ux43cux44bux445-ux43eux442ux43eux431ux440ux430ux436ux435ux43dux438ux439-ux438ux437ux43cux435ux440ux438ux43cux430}

Если

\[
X:(\Omega,\mathcal F)\to(E,\mathcal E)
\]

измеримо и

\[
f:(E,\mathcal E)\to(S,\mathcal S)
\]

измеримо, то композиция

\[
f\circ X:(\Omega,\mathcal F)\to(S,\mathcal S)
\]

измерима.

\subsubsection{Доказательство}\label{ux434ux43eux43aux430ux437ux430ux442ux435ux43bux44cux441ux442ux432ux43e-24}

Для любого \(C\in\mathcal S\)

\[
(f\circ X)^{-1}(C)=X^{-1}(f^{-1}(C)).
\]

Так как \(f\) измеримо, \(f^{-1}(C)\in\mathcal E\). Так как \(X\)
измеримо, \(X^{-1}(f^{-1}(C))\in\mathcal F\). \(\square\)

\subsubsection{Свойство 6. Критерий измеримости случайного
вектора}\label{ux441ux432ux43eux439ux441ux442ux432ux43e-6.-ux43aux440ux438ux442ux435ux440ux438ux439-ux438ux437ux43cux435ux440ux438ux43cux43eux441ux442ux438-ux441ux43bux443ux447ux430ux439ux43dux43eux433ux43e-ux432ux435ux43aux442ux43eux440ux430}

Отображение

\[
X=(X_1,\ldots,X_d):\Omega\to\mathbb R^d
\]

измеримо относительно \(\mathcal B(\mathbb R^d)\) тогда и только тогда,
когда все координаты

\[
X_1,\ldots,X_d
\]

измеримы как вещественные функции.

\subsubsection{Доказательство}\label{ux434ux43eux43aux430ux437ux430ux442ux435ux43bux44cux441ux442ux432ux43e-25}

Если \(X\) измеримо, то координата

\[
X_k=\pi_k\circ X
\]

измерима, потому что проекция \(\pi_k:\mathbb R^d\to\mathbb R\)
непрерывна, а композиция измеримых отображений измерима.

Обратно, пусть все \(X_k\) измеримы. Достаточно проверить прообразы
прямоугольников

\[
(-\infty,x_1]\times\cdots\times(-\infty,x_d],
\]

потому что такие множества порождают \(\mathcal B(\mathbb R^d)\). Но

\[
X^{-1}\left((-\infty,x_1]\times\cdots\times(-\infty,x_d]\right)
=
\bigcap_{k=1}^d\{X_k\le x_k\}.
\]

Каждое множество справа лежит в \(\mathcal F\), а \(\mathcal F\)
замкнута относительно конечных пересечений. Значит \(X\) измерим.
\(\square\)

\subsubsection{Свойство 7. Борелевская функция от случайного вектора
дает случайную
величину}\label{ux441ux432ux43eux439ux441ux442ux432ux43e-7.-ux431ux43eux440ux435ux43bux435ux432ux441ux43aux430ux44f-ux444ux443ux43dux43aux446ux438ux44f-ux43eux442-ux441ux43bux443ux447ux430ux439ux43dux43eux433ux43e-ux432ux435ux43aux442ux43eux440ux430-ux434ux430ux435ux442-ux441ux43bux443ux447ux430ux439ux43dux443ux44e-ux432ux435ux43bux438ux447ux438ux43dux443}

Если \(X:\Omega\to\mathbb R^d\) - случайный вектор, а
\(g:\mathbb R^d\to\mathbb R^m\) - борелевски измеримое отображение, то

\[
g(X)
\]

является случайным вектором в \(\mathbb R^m\).

В частности, если \(X,Y\) - вещественные случайные величины, то
измеримы:

\[
X+Y,\qquad XY,\qquad \max(X,Y),\qquad \min(X,Y),
\]

если \(Y\ne0\) всюду, то также \(X/Y\). Если \(Y=0\) возможно, нужно
определить значение на множестве \(\{Y=0\}\), например положить там
\(0\), и рассматривать такую измеримую модификацию.

\subsubsection{Обоснование}\label{ux43eux431ux43eux441ux43dux43eux432ux430ux43dux438ux435-2}

Вектор \((X,Y)\) измерим по свойству 6, а операции сложения, умножения,
максимума и минимума являются непрерывными функциями
\(\mathbb R^2\to\mathbb R\).

\subsubsection{Свойство 8. Предел измеримых вещественных функций
измерим}\label{ux441ux432ux43eux439ux441ux442ux432ux43e-8.-ux43fux440ux435ux434ux435ux43b-ux438ux437ux43cux435ux440ux438ux43cux44bux445-ux432ux435ux449ux435ux441ux442ux432ux435ux43dux43dux44bux445-ux444ux443ux43dux43aux446ux438ux439-ux438ux437ux43cux435ux440ux438ux43c}

Если \(X_n:\Omega\to\mathbb R\) измеримы, то функции

\[
\sup_n X_n,\qquad \inf_n X_n,\qquad \limsup_{n\to\infty}X_n,\qquad \liminf_{n\to\infty}X_n
\]

измеримы как расширенные вещественные функции. Если
\(X_n(\omega)\to X(\omega)\) для всех \(\omega\), то \(X\) измерима.

\subsubsection{Доказательство-идея}\label{ux434ux43eux43aux430ux437ux430ux442ux435ux43bux44cux441ux442ux432ux43e-ux438ux434ux435ux44f-3}

Например,

\[
\left\{\sup_n X_n>a\right\}
=
\bigcup_{n=1}^{\infty}\{X_n>a\}\in\mathcal F.
\]

Значит \(\sup_nX_n\) измерима. Инфимум выражается через супремум:

\[
\inf_nX_n=-\sup_n(-X_n).
\]

Далее

\[
\limsup_{n\to\infty}X_n=\inf_{m\ge1}\sup_{n\ge m}X_n,
\]

\[
\liminf_{n\to\infty}X_n=\sup_{m\ge1}\inf_{n\ge m}X_n.
\]

Если обычный предел существует, то

\[
X=\limsup_{n\to\infty}X_n=\liminf_{n\to\infty}X_n.
\]

\subsubsection{Свойство 9. Измеримость случайной
последовательности}\label{ux441ux432ux43eux439ux441ux442ux432ux43e-9.-ux438ux437ux43cux435ux440ux438ux43cux43eux441ux442ux44c-ux441ux43bux443ux447ux430ux439ux43dux43eux439-ux43fux43eux441ux43bux435ux434ux43eux432ux430ux442ux435ux43bux44cux43dux43eux441ux442ux438}

Пусть

\[
X(\omega)=(X_1(\omega),X_2(\omega),\ldots)
\]

и на \(\mathbb R^{\mathbb N}\) взята произведенная \(\sigma\)-алгебра
\(\mathcal B(\mathbb R)^{\otimes\mathbb N}\). Тогда \(X\) измеримо тогда
и только тогда, когда каждая координата \(X_n\) измерима.

\subsubsection{Доказательство-идея}\label{ux434ux43eux43aux430ux437ux430ux442ux435ux43bux44cux441ux442ux432ux43e-ux438ux434ux435ux44f-4}

Произведенная \(\sigma\)-алгебра порождается цилиндрами

\[
\pi_{i_1}^{-1}(B_1)\cap\cdots\cap\pi_{i_k}^{-1}(B_k),
\]

где \(B_j\in\mathcal B(\mathbb R)\). Прообраз такого цилиндра равен

\[
\{X_{i_1}\in B_1\}\cap\cdots\cap\{X_{i_k}\in B_k\},
\]

что измеримо при измеримости всех координат.

\subsection{5. Примеры}\label{ux43fux440ux438ux43cux435ux440ux44b-4}

\subsubsection{Пример 1. Индикатор
события}\label{ux43fux440ux438ux43cux435ux440-1.-ux438ux43dux434ux438ux43aux430ux442ux43eux440-ux441ux43eux431ux44bux442ux438ux44f}

Пусть \(A\in\mathcal F\). Тогда

\[
\mathbf 1_A(\omega)=
\begin{cases}
1, & \omega\in A,\\
0, & \omega\notin A
\end{cases}
\]

является случайной величиной. Действительно, для любого \(x\) множество
\(\{\mathbf 1_A\le x\}\) равно одному из множеств

\[
\varnothing,\qquad \Omega\setminus A,\qquad \Omega,
\]

и все они принадлежат \(\mathcal F\).

\subsubsection{Пример 2. Простая случайная
величина}\label{ux43fux440ux438ux43cux435ux440-2.-ux43fux440ux43eux441ux442ux430ux44f-ux441ux43bux443ux447ux430ux439ux43dux430ux44f-ux432ux435ux43bux438ux447ux438ux43dux430}

Если \(A_1,\ldots,A_n\in\mathcal F\) и \(a_1,\ldots,a_n\in\mathbb R\),
то

\[
X(\omega)=\sum_{k=1}^n a_k\mathbf 1_{A_k}(\omega)
\]

измерима. Она принимает конечное число значений. Чтобы увидеть
измеримость строго, можно разбить \(\Omega\) на конечное число атомов
вида

\[
\bigcap_{k=1}^n C_k,\qquad C_k\in\{A_k,\Omega\setminus A_k\}.
\]

На каждом таком атоме \(X\) постоянна, поэтому \(\{X\in B\}\) является
конечным объединением измеримых атомов. Такие функции будут основой
конструкции интеграла Лебега в Билете 07.

\subsubsection{Пример 3. Случайный
вектор}\label{ux43fux440ux438ux43cux435ux440-3.-ux441ux43bux443ux447ux430ux439ux43dux44bux439-ux432ux435ux43aux442ux43eux440}

Если \(X\) и \(Y\) - случайные величины, то

\[
Z=(X,Y):\Omega\to\mathbb R^2
\]

является случайным вектором. Например, событие попадания в прямоугольник
имеет вид

\[
\{Z\in(-\infty,x]\times(-\infty,y]\}
=
\{X\le x\}\cap\{Y\le y\}.
\]

Это событие измеримо.

\subsubsection{Пример 4. Случайная
последовательность}\label{ux43fux440ux438ux43cux435ux440-4.-ux441ux43bux443ux447ux430ux439ux43dux430ux44f-ux43fux43eux441ux43bux435ux434ux43eux432ux430ux442ux435ux43bux44cux43dux43eux441ux442ux44c}

Пусть \(X_1,X_2,\ldots\) - случайные величины. Тогда

\[
X(\omega)=(X_1(\omega),X_2(\omega),\ldots)
\]

является случайным элементом пространства \(\mathbb R^{\mathbb N}\) с
произведенной \(\sigma\)-алгеброй. События, зависящие от конечного числа
координат, например

\[
\{X_1\le a_1,\ X_2\le a_2,\ X_5\in B\},
\]

измеримы. Именно такие события задают конечномерные распределения.

\subsubsection{Пример 5. Непрерывная функция от случайной
величины}\label{ux43fux440ux438ux43cux435ux440-5.-ux43dux435ux43fux440ux435ux440ux44bux432ux43dux430ux44f-ux444ux443ux43dux43aux446ux438ux44f-ux43eux442-ux441ux43bux443ux447ux430ux439ux43dux43eux439-ux432ux435ux43bux438ux447ux438ux43dux44b}

Если \(X\) - случайная величина, то

\[
e^X,\qquad X^2,\qquad \sin X,\qquad |X|
\]

тоже случайные величины, потому что соответствующие функции
\(\mathbb R\to\mathbb R\) непрерывны и, следовательно, борелевски
измеримы.

\subsubsection{Пример 6. Не всякая функция является случайной
величиной}\label{ux43fux440ux438ux43cux435ux440-6.-ux43dux435-ux432ux441ux44fux43aux430ux44f-ux444ux443ux43dux43aux446ux438ux44f-ux44fux432ux43bux44fux435ux442ux441ux44f-ux441ux43bux443ux447ux430ux439ux43dux43eux439-ux432ux435ux43bux438ux447ux438ux43dux43eux439}

Пусть \(\Omega=\mathbb R\) и \(\mathcal F=\{\varnothing,\mathbb R\}\) -
тривиальная \(\sigma\)-алгебра. Функция

\[
X(\omega)=\omega
\]

не является случайной величиной относительно этой \(\mathcal F\), потому
что

\[
\{X\le0\}=(-\infty,0]\notin\mathcal F.
\]

Проблема не в самой формуле \(X(\omega)=\omega\), а в слишком бедной
\(\sigma\)-алгебре событий.

\subsection{6. Доказательства и
обоснования}\label{ux434ux43eux43aux430ux437ux430ux442ux435ux43bux44cux441ux442ux432ux430-ux438-ux43eux431ux43eux441ux43dux43eux432ux430ux43dux438ux44f-4}

\textbf{Почему в определении стоят прообразы, а не образы.}

Вероятность задана на подмножествах \(\Omega\), то есть на событиях из
\(\mathcal F\). Если в пространстве значений есть множество
\(B\subset E\), то событие ``значение \(X\) лежит в \(B\)'' находится в
исходном пространстве:

\[
\{\omega:X(\omega)\in B\}=X^{-1}(B).
\]

Именно его вероятность можно пытаться вычислить. Образ множества
\(X(A)\) лежит в \(E\) и сам по себе не является событием в \(\Omega\).

\textbf{Почему нельзя проверять только одно множество.}

Если проверить только \(\{X\le0\}\), это ничего не говорит о событиях
\(\{X\le x\}\) при других \(x\). Для распределения и интеграла нужно,
чтобы были измеримы все события вида \(\{X\in B\}\), где \(B\)
борелевское. Поэтому проверяют либо все борелевские множества, либо
порождающий класс, например все лучи \((-\infty,x]\).

\textbf{Почему координатный критерий работает только для правильной
\(\sigma\)-алгебры.}

Для \(\mathbb R^d\) с борелевской \(\sigma\)-алгеброй прямоугольники
порождают все борелевские множества, поэтому измеримость координат
достаточна. Если на пространстве значений взять произвольную более
богатую \(\sigma\)-алгебру, то координат может быть недостаточно: нужно
проверять именно \(\sigma\)-алгебру, выбранную в пространстве значений.

\textbf{Почему в метрическом пространстве часто используют счетную
базу.}

Если метрическое пространство сепарабельно, то у него есть счетная база
открытых шаров. Тогда для проверки борелевской измеримости достаточно
проверять прообразы шаров из этой базы. Это удобно для пространств
последовательностей и функций. Без сепарабельности некоторые привычные
упрощения могут не работать.

\textbf{Предупреждение про неограниченные и расширенные функции.}

Обычная вещественная случайная величина принимает значения в
\(\mathbb R\). Иногда разрешают значения \(\pm\infty\) и говорят о
расширенной случайной величине со значениями в

\[
\overline{\mathbb R}=[-\infty,\infty].
\]

Тогда нужно явно понимать, какая \(\sigma\)-алгебра стоит на
\(\overline{\mathbb R}\), и отдельно следить за корректностью операций
вроде \(+\infty-\infty\).

\textbf{Предупреждение про процессы с несчетным временем.}

Для счетного набора времен измеримость всех координат хорошо согласуется
с произведенной \(\sigma\)-алгеброй. Для процессов с несчетным
множеством времен, например \(t\in[0,1]\), утверждения о траекториях как
элементах пространства функций требуют дополнительных условий и выбора
конкретной \(\sigma\)-алгебры или топологии на пространстве траекторий.

\subsection{7. Что написать на
доске}\label{ux447ux442ux43e-ux43dux430ux43fux438ux441ux430ux442ux44c-ux43dux430-ux434ux43eux441ux43aux435-4}

\begin{enumerate}
\def\labelenumi{\arabic{enumi}.}
\tightlist
\item
  Измеримое отображение:
\end{enumerate}

\[
X:(\Omega,\mathcal F)\to(E,\mathcal E),
\qquad
X^{-1}(B)\in\mathcal F\quad \forall B\in\mathcal E.
\]

\begin{enumerate}
\def\labelenumi{\arabic{enumi}.}
\setcounter{enumi}{1}
\tightlist
\item
  Случайная величина:
\end{enumerate}

\[
X:(\Omega,\mathcal F,P)\to(\mathbb R,\mathcal B(\mathbb R)).
\]

\begin{enumerate}
\def\labelenumi{\arabic{enumi}.}
\setcounter{enumi}{2}
\tightlist
\item
  Критерий для \(\mathbb R\):
\end{enumerate}

\[
X\text{ измерима}
\quad\Longleftrightarrow\quad
\{X\le x\}\in\mathcal F\quad \forall x\in\mathbb R.
\]

\begin{enumerate}
\def\labelenumi{\arabic{enumi}.}
\setcounter{enumi}{3}
\tightlist
\item
  Проверка по порождающему классу:
\end{enumerate}

\[
\mathcal E=\sigma(\mathcal C),\quad
X^{-1}(C)\in\mathcal F\ \forall C\in\mathcal C
\Rightarrow
X\text{ измеримо}.
\]

\begin{enumerate}
\def\labelenumi{\arabic{enumi}.}
\setcounter{enumi}{4}
\tightlist
\item
  Векторный критерий:
\end{enumerate}

\[
X=(X_1,\ldots,X_d)\text{ измерим}
\quad\Longleftrightarrow\quad
X_1,\ldots,X_d\text{ измеримы}.
\]

\begin{enumerate}
\def\labelenumi{\arabic{enumi}.}
\setcounter{enumi}{5}
\tightlist
\item
  Композиция:
\end{enumerate}

\[
X\text{ измеримо},\quad g\text{ борелевски измерима}
\Rightarrow
g(X)\text{ измеримо}.
\]

\begin{enumerate}
\def\labelenumi{\arabic{enumi}.}
\setcounter{enumi}{6}
\tightlist
\item
  Случайная последовательность:
\end{enumerate}

\[
X:\Omega\to\mathbb R^{\mathbb N},
\qquad
X(\omega)=(X_1(\omega),X_2(\omega),\ldots).
\]

\subsection{8. Типичные дополнительные
вопросы}\label{ux442ux438ux43fux438ux447ux43dux44bux435-ux434ux43eux43fux43eux43bux43dux438ux442ux435ux43bux44cux43dux44bux435-ux432ux43eux43fux440ux43eux441ux44b-4}

\textbf{Вопрос. Зачем нужна измеримость случайной величины?}

Чтобы события вида \(\{X\in B\}\) принадлежали \(\mathcal F\) и имели
вероятность. Без измеримости нельзя корректно определить распределение
\(P_X(B)=P\{X\in B\}\).

\textbf{Вопрос. Почему в определении измеримости используется прообраз?}

Потому что вероятность задана на событиях в \(\Omega\). Множество \(B\)
лежит в пространстве значений, а событие в исходном пространстве равно
\(X^{-1}(B)\).

\textbf{Вопрос. Нужно ли проверять все борелевские множества на прямой?}

По определению да, но практически нет. Достаточно проверить порождающий
класс, например все лучи \((-\infty,x]\):

\[
\{X\le x\}\in\mathcal F\quad \forall x.
\]

\textbf{Вопрос. Почему непрерывная функция от случайной величины
измерима?}

Непрерывная функция борелевски измерима, а композиция измеримых
отображений измерима.

\textbf{Вопрос. Когда вектор \((X_1,\ldots,X_d)\) является случайным?}

Тогда и только тогда, когда все координаты \(X_1,\ldots,X_d\) являются
случайными величинами.

\textbf{Вопрос. Что такое случайный элемент метрического пространства?}

Это измеримое отображение из вероятностного пространства в метрическое
пространство с его борелевской \(\sigma\)-алгеброй:

\[
X:(\Omega,\mathcal F,P)\to(E,\mathcal B(E)).
\]

\textbf{Вопрос. Чем случайный элемент отличается от случайной величины?}

Случайная величина принимает значения в \(\mathbb R\). Случайный элемент
может принимать значения в более общем пространстве: \(\mathbb R^d\),
\(\mathbb R^{\mathbb N}\), пространстве функций, множестве состояний
цепи.

\textbf{Вопрос. Что такое траектория случайного процесса?}

При фиксированном \(\omega\) это функция времени:

\[
t\mapsto X_t(\omega).
\]

А при фиксированном \(t\) функция \(\omega\mapsto X_t(\omega)\)
называется сечением процесса.

\textbf{Вопрос. Может ли функция быть неслучайной из-за выбора
\(\sigma\)-алгебры?}

Да. Если \(\mathcal F\) слишком бедная, прообразы борелевских множеств
могут не лежать в \(\mathcal F\). Например, тождественная функция на
\(\mathbb R\) не измерима относительно тривиальной \(\sigma\)-алгебры
\(\{\varnothing,\mathbb R\}\).

\subsection{9. Типичные
ошибки}\label{ux442ux438ux43fux438ux447ux43dux44bux435-ux43eux448ux438ux431ux43aux438-4}

\begin{enumerate}
\def\labelenumi{\arabic{enumi}.}
\item
  Говорить ``случайная величина - любая функция на \(\Omega\)''.
  Правильно: случайная величина - измеримая функция.
\item
  Путать \(X(\omega)\) и распределение \(P_X\). Первое - значение
  функции при исходе, второе - мера на пространстве значений.
\item
  Проверять образ \(X(A)\) вместо прообраза \(X^{-1}(B)\). В определении
  измеримости всегда стоят прообразы.
\item
  Проверять одно событие, например \(\{X>0\}\), и считать, что
  измеримость доказана. Нужно проверить порождающий класс или всю
  \(\sigma\)-алгебру.
\item
  Забывать, какая \(\sigma\)-алгебра стоит в пространстве значений. Для
  вещественных величин обычно это \(\mathcal B(\mathbb R)\).
\item
  Считать, что координатная измеримость автоматически решает любую
  задачу со случайными функциями. Для пространств траекторий нужно
  учитывать выбранную \(\sigma\)-алгебру и топологию.
\item
  Путать сечение и траекторию процесса: \(X_t(\omega)\) как функция
  \(\omega\) - сечение, как функция \(t\) при фиксированном \(\omega\) -
  траектория.
\end{enumerate}

\subsection{10. Связи с другими
билетами}\label{ux441ux432ux44fux437ux438-ux441-ux434ux440ux443ux433ux438ux43cux438-ux431ux438ux43bux435ux442ux430ux43cux438-4}

\textbf{Билет 01.} Измеримость использует \(\sigma\)-алгебры как классы
допустимых событий.

\textbf{Билет 02.} Вероятностное пространство \((\Omega,\mathcal F,P)\)
- исходная область, на которой определяются случайные величины.

\textbf{Билет 04.} Для вещественных и векторных случайных величин
пространство значений снабжается борелевской \(\sigma\)-алгеброй.

\textbf{Билет 06.} Распределение случайного элемента определяется через
измеримость:

\[
P_X(B)=P(X^{-1}(B)).
\]

\textbf{Билет 07.} Интеграл Лебега строится сначала для измеримых
функций.

\textbf{Билет 08.} Характеристическая функция \(\mathbb E e^{itX}\)
имеет смысл, потому что \(e^{itX}\) измерима.

\textbf{Билет 10.} Условное математическое ожидание является измеримой
функцией относительно заданной \(\sigma\)-алгебры.

\textbf{Билет 15.} Теорема Колмогорова строит случайный процесс как
случайный элемент пространства последовательностей по согласованным
конечномерным распределениям.

\textbf{Билет 16.} Дискретный случайный процесс можно понимать как
последовательность случайных величин и как один случайный элемент
пространства последовательностей.

\subsection{11. Минимум на
удовлетворительно}\label{ux43cux438ux43dux438ux43cux443ux43c-ux43dux430-ux443ux434ux43eux432ux43bux435ux442ux432ux43eux440ux438ux442ux435ux43bux44cux43dux43e-4}

Нужно уметь сказать:

\begin{enumerate}
\def\labelenumi{\arabic{enumi}.}
\tightlist
\item
  Измеримое отображение:
\end{enumerate}

\[
X^{-1}(B)\in\mathcal F\quad \forall B\in\mathcal E.
\]

\begin{enumerate}
\def\labelenumi{\arabic{enumi}.}
\setcounter{enumi}{1}
\tightlist
\item
  Случайная величина - измеримое отображение
\end{enumerate}

\[
X:(\Omega,\mathcal F,P)\to(\mathbb R,\mathcal B(\mathbb R)).
\]

\begin{enumerate}
\def\labelenumi{\arabic{enumi}.}
\setcounter{enumi}{2}
\tightlist
\item
  Для вещественной случайной величины достаточно проверять
\end{enumerate}

\[
\{X\le x\}\in\mathcal F\quad \forall x.
\]

\begin{enumerate}
\def\labelenumi{\arabic{enumi}.}
\setcounter{enumi}{3}
\item
  Случайный вектор измерим тогда и только тогда, когда измеримы его
  координаты.
\item
  Один пример: \(\mathbf 1_A\) является случайной величиной, если
  \(A\in\mathcal F\).
\end{enumerate}

\subsection{12. Минимум на
хорошо}\label{ux43cux438ux43dux438ux43cux443ux43c-ux43dux430-ux445ux43eux440ux43eux448ux43e-4}

Кроме минимума на удовлетворительно, нужно:

\begin{enumerate}
\def\labelenumi{\arabic{enumi}.}
\tightlist
\item
  Объяснить, почему достаточно проверять порождающий класс.
\item
  Доказать или аккуратно проговорить критерий для \(\{X\le x\}\).
\item
  Сформулировать измеримость композиции.
\item
  Объяснить, почему непрерывная функция от случайной величины измерима.
\item
  Разобрать случайный вектор и случайную последовательность.
\end{enumerate}

\subsection{13. Что добавить для
отлично}\label{ux447ux442ux43e-ux434ux43eux431ux430ux432ux438ux442ux44c-ux434ux43bux44f-ux43eux442ux43bux438ux447ux43dux43e-4}

Для сильного ответа стоит добавить:

\begin{enumerate}
\def\labelenumi{\arabic{enumi}.}
\tightlist
\item
  Доказательство через класс
\end{enumerate}

\[
\mathcal D=\{B:X^{-1}(B)\in\mathcal F\}.
\]

\begin{enumerate}
\def\labelenumi{\arabic{enumi}.}
\setcounter{enumi}{1}
\tightlist
\item
  Координатный критерий для \(\mathbb R^d\) с доказательством через
  прямоугольники.
\item
  Связь случайной последовательности с произведенной \(\sigma\)-алгеброй
  и цилиндрами.
\item
  Различие между сечением и траекторией случайного процесса.
\item
  Предупреждение: для несчетного времени и пространств траекторий нужны
  дополнительные условия.
\end{enumerate}

\subsection{14. Устная версия ответа на 5
минут}\label{ux443ux441ux442ux43dux430ux44f-ux432ux435ux440ux441ux438ux44f-ux43eux442ux432ux435ux442ux430-ux43dux430-5-ux43cux438ux43dux443ux442-4}

Измеримые функции нужны для того, чтобы функции от исхода были
корректными вероятностными объектами. Пусть есть измеримые пространства
\((\Omega,\mathcal F)\) и \((E,\mathcal E)\). Отображение

\[
X:\Omega\to E
\]

называется измеримым, если

\[
X^{-1}(B)\in\mathcal F
\qquad \forall B\in\mathcal E.
\]

Это означает, что для любого наблюдаемого множества значений \(B\)
событие \(\{X\in B\}\) действительно является событием, которому можно
приписать вероятность.

Вещественная случайная величина - это измеримое отображение

\[
X:(\Omega,\mathcal F,P)\to(\mathbb R,\mathcal B(\mathbb R)).
\]

По определению нужно проверять все борелевские множества, но практически
достаточно проверять лучи:

\[
\{X\le x\}\in\mathcal F
\qquad \forall x\in\mathbb R.
\]

Это верно потому, что лучи \((-\infty,x]\) порождают
\(\mathcal B(\mathbb R)\), а класс множеств \(B\), для которых
\(X^{-1}(B)\) измерим, является \(\sigma\)-алгеброй.

Случайный вектор

\[
X=(X_1,\ldots,X_d)
\]

измерим тогда и только тогда, когда измеримы все координаты.
Доказательство идет через прямоугольники:

\[
\{X_1\le x_1,\ldots,X_d\le x_d\}
=
\bigcap_{k=1}^d\{X_k\le x_k\}.
\]

Если \(g\) - борелевски измеримая функция, а \(X\) - случайный элемент,
то \(g(X)\) снова измерим. В частности, непрерывные функции от случайных
величин, суммы, произведения, максимумы и минимумы случайных величин
снова являются случайными величинами.

Случайный элемент метрического пространства - это измеримое отображение
в \((E,\mathcal B(E))\). Например, случайная последовательность может
рассматриваться как одно отображение

\[
X:\Omega\to\mathbb R^{\mathbb N},
\qquad
X(\omega)=(X_1(\omega),X_2(\omega),\ldots),
\]

если на пространстве последовательностей взять произведенную
\(\sigma\)-алгебру. Это связывает билет с конечномерными
распределениями, теоремой Колмогорова и дискретными случайными
процессами.

\subsection{15. Если осталось 2 минуты перед
экзаменом}\label{ux435ux441ux43bux438-ux43eux441ux442ux430ux43bux43eux441ux44c-2-ux43cux438ux43dux443ux442ux44b-ux43fux435ux440ux435ux434-ux44dux43aux437ux430ux43cux435ux43dux43eux43c-4}

Запомнить:

\[
X:(\Omega,\mathcal F)\to(E,\mathcal E)
\text{ измеримо}
\quad\Longleftrightarrow\quad
X^{-1}(B)\in\mathcal F\ \forall B\in\mathcal E.
\]

Случайная величина:

\[
X:(\Omega,\mathcal F,P)\to(\mathbb R,\mathcal B(\mathbb R)).
\]

Для \(\mathbb R\):

\[
X\text{ измерима}
\quad\Longleftrightarrow\quad
\{X\le x\}\in\mathcal F\ \forall x.
\]

Для вектора:

\[
(X_1,\ldots,X_d)\text{ измерим}
\quad\Longleftrightarrow\quad
X_1,\ldots,X_d\text{ измеримы}.
\]

Главная мысль: измеримость нужна, чтобы \(\{X\in B\}\) было событием и
имело вероятность. Главная ловушка: случайная величина - не любая
функция, а измеримая функция.

\newpage

\section{Билет 06. Распределения случайных величин, векторов и
элементов. Конечномерные распределения и
плотности}\label{ux431ux438ux43bux435ux442-06.-ux440ux430ux441ux43fux440ux435ux434ux435ux43bux435ux43dux438ux44f-ux441ux43bux443ux447ux430ux439ux43dux44bux445-ux432ux435ux43bux438ux447ux438ux43d-ux432ux435ux43aux442ux43eux440ux43eux432-ux438-ux44dux43bux435ux43cux435ux43dux442ux43eux432.-ux43aux43eux43dux435ux447ux43dux43eux43cux435ux440ux43dux44bux435-ux440ux430ux441ux43fux440ux435ux434ux435ux43bux435ux43dux438ux44f-ux438-ux43fux43bux43eux442ux43dux43eux441ux442ux438}

Источники: Ширяев 1, гл. II; Сердобольская, лекция 1; лекции ДСП.

\subsection{1. Короткий
ответ}\label{ux43aux43eux440ux43eux442ux43aux438ux439-ux43eux442ux432ux435ux442-5}

Распределение случайного элемента - это образ вероятностной меры при
измеримом отображении. Если

\[
X:(\Omega,\mathcal F,P)\to(E,\mathcal E)
\]

измерим, то его распределение задается формулой

\[
P_X(B)=P\{X\in B\}=P(X^{-1}(B)),\qquad B\in\mathcal E.
\]

Для вещественной случайной величины распределение можно задавать
функцией распределения

\[
F_X(x)=P\{X\le x\}.
\]

Для случайного вектора \(X=(X_1,\ldots,X_d)\):

\[
F_X(x_1,\ldots,x_d)=P\{X_1\le x_1,\ldots,X_d\le x_d\}.
\]

Для процесса \(\{X_t:t\in T\}\) конечномерные распределения - это
распределения векторов

\[
(X_{t_1},\ldots,X_{t_n}).
\]

Если распределение имеет плотность \(p\), то вероятность множества
считается интегралом:

\[
P_X(B)=\int_B p(x)\,dx.
\]

\subsection{2. Полный
ответ}\label{ux43fux43eux43bux43dux44bux439-ux43eux442ux432ux435ux442-5}

\subsubsection{6.1. Введение: перенос вероятностной
меры}\label{ux432ux432ux435ux434ux435ux43dux438ux435-ux43fux435ux440ux435ux43dux43eux441-ux432ux435ux440ux43eux44fux442ux43dux43eux441ux442ux43dux43eux439-ux43cux435ux440ux44b}

В Билете 05 случайная величина была определена как измеримое
отображение. Теперь главный вопрос: какую вероятность она задает на
пространстве своих значений. Исходная вероятность \(P\) живет на
\((\Omega,\mathcal F)\), но наблюдаем мы обычно не сам исход \(\omega\),
а значение \(X(\omega)\). Поэтому вероятность нужно перенести с
\(\Omega\) на пространство значений.

\subsubsection{6.2. Определение распределения
(закона)}\label{ux43eux43fux440ux435ux434ux435ux43bux435ux43dux438ux435-ux440ux430ux441ux43fux440ux435ux434ux435ux43bux435ux43dux438ux44f-ux437ux430ux43aux43eux43dux430}

Пусть

\[
X:(\Omega,\mathcal F,P)\to(E,\mathcal E)
\]

\begin{itemize}
\tightlist
\item
  случайный элемент. Его распределением, или законом, называется мера
  \(P_X\) на \((E,\mathcal E)\):
\end{itemize}

\[
P_X(B)=P\{\omega:X(\omega)\in B\}=P(X^{-1}(B)),\qquad B\in\mathcal E.
\]

Эта формула корректна именно потому, что \(X\) измерим: для каждого
\(B\in\mathcal E\) прообраз \(X^{-1}(B)\) лежит в \(\mathcal F\), значит
имеет вероятность.

\subsubsection{6.3. Свойства распределения как вероятностной
меры}\label{ux441ux432ux43eux439ux441ux442ux432ux430-ux440ux430ux441ux43fux440ux435ux434ux435ux43bux435ux43dux438ux44f-ux43aux430ux43a-ux432ux435ux440ux43eux44fux442ux43dux43eux441ux442ux43dux43eux439-ux43cux435ux440ux44b}

Распределение является вероятностной мерой. Действительно,

\[
P_X(E)=P(X^{-1}(E))=P(\Omega)=1,
\]

а счетная аддитивность следует из того, что прообраз сохраняет
дизъюнктные объединения:

\[
X^{-1}\left(\bigsqcup_{n=1}^{\infty}B_n\right)
=
\bigsqcup_{n=1}^{\infty}X^{-1}(B_n).
\]

\subsubsection{6.4. Функция распределения вещественной случайной
величины}\label{ux444ux443ux43dux43aux446ux438ux44f-ux440ux430ux441ux43fux440ux435ux434ux435ux43bux435ux43dux438ux44f-ux432ux435ux449ux435ux441ux442ux432ux435ux43dux43dux43eux439-ux441ux43bux443ux447ux430ux439ux43dux43eux439-ux432ux435ux43bux438ux447ux438ux43dux44b}

Для вещественной случайной величины \(X\) распределение \(P_X\) - это
вероятностная мера на \(\mathcal B(\mathbb R)\). Часто ее задают не
самой мерой, а функцией распределения

\[
F_X(x)=P\{X\le x\}=P_X((-\infty,x]).
\]

Функция распределения не убывает, непрерывна справа и имеет пределы

\[
\lim_{x\to-\infty}F_X(x)=0,\qquad
\lim_{x\to+\infty}F_X(x)=1.
\]

Наоборот, всякая функция с этими свойствами задает вероятностную меру на
\(\mathcal B(\mathbb R)\); это уже использовалось в Билете 02.

\subsubsection{6.5. Распределение и функция распределения случайного
вектора}\label{ux440ux430ux441ux43fux440ux435ux434ux435ux43bux435ux43dux438ux435-ux438-ux444ux443ux43dux43aux446ux438ux44f-ux440ux430ux441ux43fux440ux435ux434ux435ux43bux435ux43dux438ux44f-ux441ux43bux443ux447ux430ux439ux43dux43eux433ux43e-ux432ux435ux43aux442ux43eux440ux430}

Для случайного вектора

\[
X=(X_1,\ldots,X_d)
\]

распределение \(P_X\) - это вероятностная мера на
\(\mathcal B(\mathbb R^d)\). Его функция распределения:

\[
F_X(x_1,\ldots,x_d)
=
P\{X_1\le x_1,\ldots,X_d\le x_d\}.
\]

Она задает вероятности прямоугольников через разности значений \(F_X\).
Например, в размерности \(2\):

\[
P\{a<X\le b,\ c<Y\le d\}
=
F(b,d)-F(a,d)-F(b,c)+F(a,c).
\]

В размерности \(d\) используется аналогичная формула
включений-исключений по вершинам прямоугольника.

\subsubsection{6.6. Комплексные случайные
величины}\label{ux43aux43eux43cux43fux43bux435ux43aux441ux43dux44bux435-ux441ux43bux443ux447ux430ux439ux43dux44bux435-ux432ux435ux43bux438ux447ux438ux43dux44b}

Комплексная случайная величина рассматривается как двумерный
вещественный вектор:

\[
Z=X+iY
\quad\Longleftrightarrow\quad
(\operatorname{Re}Z,\operatorname{Im}Z)=(X,Y).
\]

Поэтому ее распределение - мера на \(\mathcal B(\mathbb C)\), что
естественно отождествляется с \(\mathcal B(\mathbb R^2)\). Для
комплексного случайного вектора из \(\mathbb C^m\) получают
распределение в \(\mathbb R^{2m}\). Важно: у комплексных чисел нет
естественного линейного порядка, поэтому обычно говорят не о функции
\(P\{Z\le z\}\), а о распределении как мере или о совместной функции
распределения вещественных и мнимых частей.

\subsubsection{6.7. Распределение в метрических и функциональных
пространствах}\label{ux440ux430ux441ux43fux440ux435ux434ux435ux43bux435ux43dux438ux435-ux432-ux43cux435ux442ux440ux438ux447ux435ux441ux43aux438ux445-ux438-ux444ux443ux43dux43aux446ux438ux43eux43dux430ux43bux44cux43dux44bux445-ux43fux440ux43eux441ux442ux440ux430ux43dux441ux442ux432ux430ux445}

Если \(X\) - случайный элемент метрического пространства \((E,\rho)\),
то его распределение - вероятностная мера на борелевской
\(\sigma\)-алгебре \(\mathcal B(E)\). Например, случайная
последовательность имеет распределение на пространстве
\(\mathbb R^{\mathbb N}\) с произведенной \(\sigma\)-алгеброй, а
случайная траектория - на выбранном пространстве функций, если это
пространство снабжено подходящей \(\sigma\)-алгеброй.

\subsubsection{6.8. Конечномерные распределения случайного
процесса}\label{ux43aux43eux43dux435ux447ux43dux43eux43cux435ux440ux43dux44bux435-ux440ux430ux441ux43fux440ux435ux434ux435ux43bux435ux43dux438ux44f-ux441ux43bux443ux447ux430ux439ux43dux43eux433ux43e-ux43fux440ux43eux446ux435ux441ux441ux430}

Для случайного процесса \(\{X_t:t\in T\}\) центральный объект -
конечномерные распределения. Для любого набора моментов
\(t_1,\ldots,t_n\) рассматривают случайный вектор

\[
(X_{t_1},\ldots,X_{t_n})
\]

и его распределение:

\[
P_{t_1,\ldots,t_n}(B)
=
P\{(X_{t_1},\ldots,X_{t_n})\in B\},
\qquad B\in\mathcal B(\mathbb R^n).
\]

Если процесс вещественный, его конечномерная функция распределения:

\[
F_{t_1,\ldots,t_n}(x_1,\ldots,x_n)
=
P\{X_{t_1}\le x_1,\ldots,X_{t_n}\le x_n\}.
\]

Эти распределения описывают совместное поведение процесса в конечном
числе моментов времени. Для полного задания закона процесса на
пространстве последовательностей нужно согласованное семейство таких
распределений; это будет темой Билета 15 и Билета 16.

\subsubsection{6.9. Распределения с
плотностью}\label{ux440ux430ux441ux43fux440ux435ux434ux435ux43bux435ux43dux438ux44f-ux441-ux43fux43bux43eux442ux43dux43eux441ux442ux44cux44e}

Отдельный важный случай - распределения с плотностью. Если \(X\)
принимает значения в \(\mathbb R^d\) и существует неотрицательная
функция \(p\) такая, что

\[
\int_{\mathbb R^d}p(x)\,dx=1
\]

и

\[
P_X(B)=\int_B p(x)\,dx,
\]

то \(p\) называется плотностью распределения \(X\) относительно меры
Лебега. В программе этого билета отдельно выделен случай непрерывной
плотности: если \(p\) непрерывна, то вероятности прямоугольников и
достаточно хороших областей можно понимать через обычный несобственный
интеграл Римана. Например,

\[
F_X(x)=\int_{-\infty}^{x}p(u)\,du
\]

в одномерном случае и

\[
F_X(x_1,\ldots,x_d)
=
\int_{-\infty}^{x_1}\cdots\int_{-\infty}^{x_d}
p(u_1,\ldots,u_d)\,du_d\cdots du_1
\]

в \(d\)-мерном случае.

\subsection{3. Основные
определения}\label{ux43eux441ux43dux43eux432ux43dux44bux435-ux43eux43fux440ux435ux434ux435ux43bux435ux43dux438ux44f-5}

\subsubsection{Распределение случайного
элемента}\label{ux440ux430ux441ux43fux440ux435ux434ux435ux43bux435ux43dux438ux435-ux441ux43bux443ux447ux430ux439ux43dux43eux433ux43e-ux44dux43bux435ux43cux435ux43dux442ux430}

Пусть \(X:(\Omega,\mathcal F,P)\to(E,\mathcal E)\) измерим.
Распределение \(X\) - это мера

\[
P_X=P\circ X^{-1}
\]

на \((E,\mathcal E)\), заданная формулой

\[
P_X(B)=P(X^{-1}(B)),\qquad B\in\mathcal E.
\]

\subsubsection{Закон случайной
величины}\label{ux437ux430ux43aux43eux43d-ux441ux43bux443ux447ux430ux439ux43dux43eux439-ux432ux435ux43bux438ux447ux438ux43dux44b}

То же самое, что распределение. Записи

\[
X\sim\mu,\qquad P_X=\mu
\]

означают, что случайная величина \(X\) имеет распределение \(\mu\).

\subsubsection{Функция распределения вещественной случайной
величины}\label{ux444ux443ux43dux43aux446ux438ux44f-ux440ux430ux441ux43fux440ux435ux434ux435ux43bux435ux43dux438ux44f-ux432ux435ux449ux435ux441ux442ux432ux435ux43dux43dux43eux439-ux441ux43bux443ux447ux430ux439ux43dux43eux439-ux432ux435ux43bux438ux447ux438ux43dux44b-1}

\[
F_X(x)=P\{X\le x\}=P_X((-\infty,x]).
\]

\subsubsection{Атом
распределения}\label{ux430ux442ux43eux43c-ux440ux430ux441ux43fux440ux435ux434ux435ux43bux435ux43dux438ux44f}

Точка \(a\), для которой

\[
P\{X=a\}>0.
\]

Скачок функции распределения в точке \(a\) равен массе атома:

\[
P\{X=a\}=F_X(a)-F_X(a-),
\]

где

\[
F_X(a-)=\lim_{x\uparrow a}F_X(x).
\]

\subsubsection{Дискретное
распределение}\label{ux434ux438ux441ux43aux440ux435ux442ux43dux43eux435-ux440ux430ux441ux43fux440ux435ux434ux435ux43bux435ux43dux438ux435}

Распределение, сосредоточенное на конечном или счетном множестве
\(\{x_k\}\):

\[
P\{X=x_k\}=p_k,\qquad p_k\ge0,\qquad \sum_k p_k=1.
\]

\subsubsection{Распределение случайного
вектора}\label{ux440ux430ux441ux43fux440ux435ux434ux435ux43bux435ux43dux438ux435-ux441ux43bux443ux447ux430ux439ux43dux43eux433ux43e-ux432ux435ux43aux442ux43eux440ux430}

Для

\[
X=(X_1,\ldots,X_d)
\]

это мера \(P_X\) на \(\mathcal B(\mathbb R^d)\):

\[
P_X(B)=P\{(X_1,\ldots,X_d)\in B\}.
\]

\subsubsection{Функция распределения случайного
вектора}\label{ux444ux443ux43dux43aux446ux438ux44f-ux440ux430ux441ux43fux440ux435ux434ux435ux43bux435ux43dux438ux44f-ux441ux43bux443ux447ux430ux439ux43dux43eux433ux43e-ux432ux435ux43aux442ux43eux440ux430}

\[
F_X(x_1,\ldots,x_d)
=
P\{X_1\le x_1,\ldots,X_d\le x_d\}.
\]

\subsubsection{Комплексная случайная
величина}\label{ux43aux43eux43cux43fux43bux435ux43aux441ux43dux430ux44f-ux441ux43bux443ux447ux430ux439ux43dux430ux44f-ux432ux435ux43bux438ux447ux438ux43dux430}

Случайная величина со значениями в \(\mathbb C\), понимаемая как
случайный вектор

\[
Z=\operatorname{Re}Z+i\operatorname{Im}Z
\quad\leftrightarrow\quad
(\operatorname{Re}Z,\operatorname{Im}Z)\in\mathbb R^2.
\]

Ее распределение - мера на
\(\mathcal B(\mathbb C)\simeq\mathcal B(\mathbb R^2)\).

\subsubsection{Случайный элемент метрического
пространства}\label{ux441ux43bux443ux447ux430ux439ux43dux44bux439-ux44dux43bux435ux43cux435ux43dux442-ux43cux435ux442ux440ux438ux447ux435ux441ux43aux43eux433ux43e-ux43fux440ux43eux441ux442ux440ux430ux43dux441ux442ux432ux430-1}

Измеримое отображение

\[
X:(\Omega,\mathcal F,P)\to(E,\mathcal B(E)),
\]

где \((E,\rho)\) - метрическое пространство. Его распределение -
вероятностная мера на \(\mathcal B(E)\).

\subsubsection{Конечномерное распределение
процесса}\label{ux43aux43eux43dux435ux447ux43dux43eux43cux435ux440ux43dux43eux435-ux440ux430ux441ux43fux440ux435ux434ux435ux43bux435ux43dux438ux435-ux43fux440ux43eux446ux435ux441ux441ux430}

Для процесса \(\{X_t:t\in T\}\) и набора \(t_1,\ldots,t_n\) это
распределение случайного вектора

\[
(X_{t_1},\ldots,X_{t_n}).
\]

\subsubsection{Конечномерная функция
распределения}\label{ux43aux43eux43dux435ux447ux43dux43eux43cux435ux440ux43dux430ux44f-ux444ux443ux43dux43aux446ux438ux44f-ux440ux430ux441ux43fux440ux435ux434ux435ux43bux435ux43dux438ux44f}

\[
F_{t_1,\ldots,t_n}(x_1,\ldots,x_n)
=
P\{X_{t_1}\le x_1,\ldots,X_{t_n}\le x_n\}.
\]

\subsubsection{Плотность
распределения}\label{ux43fux43bux43eux442ux43dux43eux441ux442ux44c-ux440ux430ux441ux43fux440ux435ux434ux435ux43bux435ux43dux438ux44f}

Измеримая функция \(p:\mathbb R^d\to\mathbb R\), неотрицательная почти
всюду, такая, что

\[
\int_{\mathbb R^d}p(x)\,dx=1
\]

и для допустимых множеств \(B\) выполнено

\[
P_X(B)=\int_B p(x)\,dx.
\]

В языке меры это означает, что \(P_X\) абсолютно непрерывно относительно
меры Лебега, а \(p\) является производной Радона-Никодима; формально это
будет связано с Билетом 10.

\subsubsection{Маргинальное
распределение}\label{ux43cux430ux440ux433ux438ux43dux430ux43bux44cux43dux43eux435-ux440ux430ux441ux43fux440ux435ux434ux435ux43bux435ux43dux438ux435}

Распределение части координат случайного вектора. Если \(X=(X_1,X_2)\),
то распределение \(X_1\) получается из совместного распределения:

\[
P_{X_1}(B)=P_{(X_1,X_2)}(B\times\mathbb R).
\]

\subsection{4. Основные теоремы и
свойства}\label{ux43eux441ux43dux43eux432ux43dux44bux435-ux442ux435ux43eux440ux435ux43cux44b-ux438-ux441ux432ux43eux439ux441ux442ux432ux430-5}

\subsubsection{Свойство 1. Распределение случайного элемента является
вероятностной
мерой}\label{ux441ux432ux43eux439ux441ux442ux432ux43e-1.-ux440ux430ux441ux43fux440ux435ux434ux435ux43bux435ux43dux438ux435-ux441ux43bux443ux447ux430ux439ux43dux43eux433ux43e-ux44dux43bux435ux43cux435ux43dux442ux430-ux44fux432ux43bux44fux435ux442ux441ux44f-ux432ux435ux440ux43eux44fux442ux43dux43eux441ux442ux43dux43eux439-ux43cux435ux440ux43eux439}

Пусть \(X:(\Omega,\mathcal F,P)\to(E,\mathcal E)\) измерим. Тогда
\(P_X\) - вероятностная мера на \((E,\mathcal E)\).

\subsubsection{Доказательство}\label{ux434ux43eux43aux430ux437ux430ux442ux435ux43bux44cux441ux442ux432ux43e-26}

Во-первых,

\[
P_X(E)=P(X^{-1}(E))=P(\Omega)=1.
\]

Во-вторых, если \(B_1,B_2,\ldots\in\mathcal E\) попарно не пересекаются,
то их прообразы тоже попарно не пересекаются, и

\[
X^{-1}\left(\bigsqcup_{n=1}^{\infty}B_n\right)
=
\bigsqcup_{n=1}^{\infty}X^{-1}(B_n).
\]

По счетной аддитивности \(P\):

\[
P_X\left(\bigsqcup_{n=1}^{\infty}B_n\right)
=
\sum_{n=1}^{\infty}P_X(B_n).
\]

Значит \(P_X\) - вероятностная мера. \(\square\)

\subsubsection{Свойство 2. Распределение вещественной случайной величины
задается функцией
распределения}\label{ux441ux432ux43eux439ux441ux442ux432ux43e-2.-ux440ux430ux441ux43fux440ux435ux434ux435ux43bux435ux43dux438ux435-ux432ux435ux449ux435ux441ux442ux432ux435ux43dux43dux43eux439-ux441ux43bux443ux447ux430ux439ux43dux43eux439-ux432ux435ux43bux438ux447ux438ux43dux44b-ux437ux430ux434ux430ux435ux442ux441ux44f-ux444ux443ux43dux43aux446ux438ux435ux439-ux440ux430ux441ux43fux440ux435ux434ux435ux43bux435ux43dux438ux44f}

Если

\[
F_X(x)=P\{X\le x\},
\]

то значения \(F_X\) однозначно задают меру \(P_X\) на
\(\mathcal B(\mathbb R)\).

\subsubsection{Обоснование}\label{ux43eux431ux43eux441ux43dux43eux432ux430ux43dux438ux435-3}

Для каждого \(x\)

\[
P_X((-\infty,x])=F_X(x).
\]

Лучами \((-\infty,x]\) порождается \(\mathcal B(\mathbb R)\), поэтому
вероятностная мера на \(\mathcal B(\mathbb R)\) однозначно определяется
значениями на этих лучах. Теорема о существовании меры по функции
распределения формулировалась в Билете 02 без доказательства.

\subsubsection{Свойство 3. Свойства функции распределения на
прямой}\label{ux441ux432ux43eux439ux441ux442ux432ux43e-3.-ux441ux432ux43eux439ux441ux442ux432ux430-ux444ux443ux43dux43aux446ux438ux438-ux440ux430ux441ux43fux440ux435ux434ux435ux43bux435ux43dux438ux44f-ux43dux430-ux43fux440ux44fux43cux43eux439}

Функция распределения \(F\) любой вещественной случайной величины
обладает свойствами:

\begin{enumerate}
\def\labelenumi{\arabic{enumi}.}
\tightlist
\item
  \(F\) не убывает;
\item
  \(F\) непрерывна справа;
\item
  \(\lim_{x\to-\infty}F(x)=0\);
\item
  \(\lim_{x\to+\infty}F(x)=1\).
\end{enumerate}

\subsubsection{Доказательство-идея}\label{ux434ux43eux43aux430ux437ux430ux442ux435ux43bux44cux441ux442ux432ux43e-ux438ux434ux435ux44f-5}

Монотонность следует из вложения лучей:

\[
x\le y\Rightarrow (-\infty,x]\subset(-\infty,y].
\]

Правые пределы и пределы на бесконечностях следуют из непрерывности
вероятностной меры сверху и снизу для монотонных последовательностей
событий. Например, если \(x_n\downarrow x\), то

\[
(-\infty,x_n]\downarrow(-\infty,x],
\]

поэтому

\[
F(x_n)\downarrow F(x).
\]

\subsubsection{Свойство 4. Масса точки равна скачку функции
распределения}\label{ux441ux432ux43eux439ux441ux442ux432ux43e-4.-ux43cux430ux441ux441ux430-ux442ux43eux447ux43aux438-ux440ux430ux432ux43dux430-ux441ux43aux430ux447ux43aux443-ux444ux443ux43dux43aux446ux438ux438-ux440ux430ux441ux43fux440ux435ux434ux435ux43bux435ux43dux438ux44f}

Для вещественной случайной величины

\[
P\{X=a\}=F(a)-F(a-).
\]

\subsubsection{Доказательство}\label{ux434ux43eux43aux430ux437ux430ux442ux435ux43bux44cux441ux442ux432ux43e-27}

События \(\{X\le x\}\) при \(x\uparrow a\) возрастают к событию
\(\{X<a\}\). Поэтому

\[
F(a-)=P\{X<a\}.
\]

Тогда

\[
P\{X=a\}=P\{X\le a\}-P\{X<a\}=F(a)-F(a-).
\]

\subsubsection{Свойство 5. Вероятность промежутка через функцию
распределения}\label{ux441ux432ux43eux439ux441ux442ux432ux43e-5.-ux432ux435ux440ux43eux44fux442ux43dux43eux441ux442ux44c-ux43fux440ux43eux43cux435ux436ux443ux442ux43aux430-ux447ux435ux440ux435ux437-ux444ux443ux43dux43aux446ux438ux44e-ux440ux430ux441ux43fux440ux435ux434ux435ux43bux435ux43dux438ux44f}

Для \(a<b\):

\[
P\{a<X\le b\}=F(b)-F(a).
\]

Также

\[
P\{X\le b\}=F(b),\qquad
P\{X>b\}=1-F(b).
\]

Остальные варианты интервалов требуют учета атомов на концах.

\subsubsection{Свойство 6. Совместная функция распределения задает
распределение
вектора}\label{ux441ux432ux43eux439ux441ux442ux432ux43e-6.-ux441ux43eux432ux43cux435ux441ux442ux43dux430ux44f-ux444ux443ux43dux43aux446ux438ux44f-ux440ux430ux441ux43fux440ux435ux434ux435ux43bux435ux43dux438ux44f-ux437ux430ux434ux430ux435ux442-ux440ux430ux441ux43fux440ux435ux434ux435ux43bux435ux43dux438ux435-ux432ux435ux43aux442ux43eux440ux430}

Для случайного вектора \(X\in\mathbb R^d\) функция

\[
F_X(x)=P\{X_1\le x_1,\ldots,X_d\le x_d\}
\]

задает значения распределения на прямоугольниках вида

\[
(-\infty,x_1]\times\cdots\times(-\infty,x_d],
\]

а такие прямоугольники порождают \(\mathcal B(\mathbb R^d)\). Поэтому
\(F_X\) однозначно задает распределение \(P_X\).

\subsubsection{Свойство 7. Маргинальные распределения получаются из
совместного}\label{ux441ux432ux43eux439ux441ux442ux432ux43e-7.-ux43cux430ux440ux433ux438ux43dux430ux43bux44cux43dux44bux435-ux440ux430ux441ux43fux440ux435ux434ux435ux43bux435ux43dux438ux44f-ux43fux43eux43bux443ux447ux430ux44eux442ux441ux44f-ux438ux437-ux441ux43eux432ux43cux435ux441ux442ux43dux43eux433ux43e}

Если \(X=(X_1,\ldots,X_d)\) и нужно распределение координат с номерами
\(i_1,\ldots,i_k\), то

\[
P_{(X_{i_1},\ldots,X_{i_k})}(B)
=
P_X(\pi^{-1}(B)),
\]

где \(\pi:\mathbb R^d\to\mathbb R^k\) - соответствующая координатная
проекция.

Для функций распределения это соответствует отправлению лишних
аргументов в \(+\infty\). Например:

\[
F_{X_1}(x)=\lim_{y\to+\infty}F_{(X_1,X_2)}(x,y).
\]

\subsubsection{Свойство 8. Плотность неотрицательна и имеет интеграл
1}\label{ux441ux432ux43eux439ux441ux442ux432ux43e-8.-ux43fux43bux43eux442ux43dux43eux441ux442ux44c-ux43dux435ux43eux442ux440ux438ux446ux430ux442ux435ux43bux44cux43dux430-ux438-ux438ux43cux435ux435ux442-ux438ux43dux442ux435ux433ux440ux430ux43b-1}

Если \(p\) - плотность распределения на \(\mathbb R^d\), то

\[
p(x)\ge0\quad\text{почти всюду}
\]

и

\[
\int_{\mathbb R^d}p(x)\,dx=1.
\]

Вероятность области \(B\) равна интегралу плотности по этой области:

\[
P_X(B)=\int_Bp(x)\,dx.
\]

В одномерном непрерывном случае:

\[
F(x)=\int_{-\infty}^{x}p(u)\,du.
\]

\subsubsection{Свойство 9. Плотность совместного распределения дает
маргинальные плотности
интегрированием}\label{ux441ux432ux43eux439ux441ux442ux432ux43e-9.-ux43fux43bux43eux442ux43dux43eux441ux442ux44c-ux441ux43eux432ux43cux435ux441ux442ux43dux43eux433ux43e-ux440ux430ux441ux43fux440ux435ux434ux435ux43bux435ux43dux438ux44f-ux434ux430ux435ux442-ux43cux430ux440ux433ux438ux43dux430ux43bux44cux43dux44bux435-ux43fux43bux43eux442ux43dux43eux441ux442ux438-ux438ux43dux442ux435ux433ux440ux438ux440ux43eux432ux430ux43dux438ux435ux43c}

Если \((X,Y)\) имеет совместную плотность \(p_{X,Y}(x,y)\), то плотность
\(X\) равна

\[
p_X(x)=\int_{-\infty}^{+\infty}p_{X,Y}(x,y)\,dy,
\]

если этот интеграл корректно определен.

В многомерном случае маргинальная плотность получается интегрированием
совместной плотности по лишним координатам.

\subsubsection{Свойство 10. Конечномерные распределения процесса должны
быть
согласованы}\label{ux441ux432ux43eux439ux441ux442ux432ux43e-10.-ux43aux43eux43dux435ux447ux43dux43eux43cux435ux440ux43dux44bux435-ux440ux430ux441ux43fux440ux435ux434ux435ux43bux435ux43dux438ux44f-ux43fux440ux43eux446ux435ux441ux441ux430-ux434ux43eux43bux436ux43dux44b-ux431ux44bux442ux44c-ux441ux43eux433ux43bux430ux441ux43eux432ux430ux43dux44b}

Если конечномерные распределения действительно получены из одного
процесса, то они согласованы при перестановке координат и при
вычеркивании координат. Например,

\[
P_{t_1,t_2}(B_1\times\mathbb R)=P_{t_1}(B_1),
\]

и

\[
P_{t_2,t_1}(B_2\times B_1)=P_{t_1,t_2}(B_1\times B_2).
\]

Полная теорема о том, когда согласованное семейство задает процесс,
относится к Билету 15.

\subsection{5. Примеры}\label{ux43fux440ux438ux43cux435ux440ux44b-5}

\subsubsection{Пример 1. Распределение
Бернулли}\label{ux43fux440ux438ux43cux435ux440-1.-ux440ux430ux441ux43fux440ux435ux434ux435ux43bux435ux43dux438ux435-ux431ux435ux440ux43dux443ux43bux43bux438}

Пусть

\[
P\{X=1\}=p,\qquad P\{X=0\}=1-p.
\]

Тогда распределение \(X\) сосредоточено на \(\{0,1\}\):

\[
P_X=(1-p)\delta_0+p\delta_1.
\]

Функция распределения:

\[
F_X(x)=
\begin{cases}
0, & x<0,\\
1-p, & 0\le x<1,\\
1, & x\ge1.
\end{cases}
\]

У этой функции скачки в точках \(0\) и \(1\), равные массам
соответствующих атомов.

\subsubsection{\texorpdfstring{Пример 2. Равномерное распределение на
\([0,1]\)}{Пример 2. Равномерное распределение на {[}0,1{]}}}\label{ux43fux440ux438ux43cux435ux440-2.-ux440ux430ux432ux43dux43eux43cux435ux440ux43dux43eux435-ux440ux430ux441ux43fux440ux435ux434ux435ux43bux435ux43dux438ux435-ux43dux430-01}

Если \(X\) равномерна на \([0,1]\), то

\[
p(x)=
\begin{cases}
1, & x\in[0,1],\\
0, & x\notin[0,1].
\end{cases}
\]

Функция распределения:

\[
F_X(x)=
\begin{cases}
0, & x<0,\\
x, & 0\le x\le1,\\
1, & x>1.
\end{cases}
\]

\subsubsection{Пример 3. Случайный вектор с
плотностью}\label{ux43fux440ux438ux43cux435ux440-3.-ux441ux43bux443ux447ux430ux439ux43dux44bux439-ux432ux435ux43aux442ux43eux440-ux441-ux43fux43bux43eux442ux43dux43eux441ux442ux44cux44e}

Пусть \((X,Y)\) имеет плотность \(p(x,y)\). Тогда

\[
P\{(X,Y)\in B\}=\iint_B p(x,y)\,dx\,dy.
\]

Например,

\[
P\{a<X\le b,\ c<Y\le d\}
=
\int_a^b\int_c^d p(x,y)\,dy\,dx.
\]

Маргинальная плотность \(X\):

\[
p_X(x)=\int_{-\infty}^{+\infty}p(x,y)\,dy.
\]

\subsubsection{Пример 4. Комплексная случайная
величина}\label{ux43fux440ux438ux43cux435ux440-4.-ux43aux43eux43cux43fux43bux435ux43aux441ux43dux430ux44f-ux441ux43bux443ux447ux430ux439ux43dux430ux44f-ux432ux435ux43bux438ux447ux438ux43dux430}

Пусть

\[
Z=X+iY.
\]

Распределение \(Z\) - это распределение вектора \((X,Y)\) на плоскости.
Например,

\[
P\{Z\in B\}=P\{(X,Y)\in \widetilde B\},
\]

где \(\widetilde B\subset\mathbb R^2\) - множество, соответствующее
\(B\subset\mathbb C\) при отождествлении \(\mathbb C\) и
\(\mathbb R^2\).

\subsubsection{Пример 5. Конечномерное распределение случайной
последовательности}\label{ux43fux440ux438ux43cux435ux440-5.-ux43aux43eux43dux435ux447ux43dux43eux43cux435ux440ux43dux43eux435-ux440ux430ux441ux43fux440ux435ux434ux435ux43bux435ux43dux438ux435-ux441ux43bux443ux447ux430ux439ux43dux43eux439-ux43fux43eux441ux43bux435ux434ux43eux432ux430ux442ux435ux43bux44cux43dux43eux441ux442ux438}

Пусть \(X_1,X_2,\ldots\) - случайная последовательность. Для индексов
\(1,3,5\) конечномерное распределение:

\[
P_{1,3,5}(B)=P\{(X_1,X_3,X_5)\in B\},
\qquad B\in\mathcal B(\mathbb R^3).
\]

Его функция распределения:

\[
F_{1,3,5}(x_1,x_3,x_5)
=
P\{X_1\le x_1,\ X_3\le x_3,\ X_5\le x_5\}.
\]

\subsubsection{Пример 6. Распределение случайного
элемента}\label{ux43fux440ux438ux43cux435ux440-6.-ux440ux430ux441ux43fux440ux435ux434ux435ux43bux435ux43dux438ux435-ux441ux43bux443ux447ux430ux439ux43dux43eux433ux43e-ux44dux43bux435ux43cux435ux43dux442ux430}

Пусть

\[
X(\omega)=(X_1(\omega),X_2(\omega),\ldots)
\]

\begin{itemize}
\tightlist
\item
  случайная последовательность как элемент \(\mathbb R^{\mathbb N}\).
  Тогда ее распределение - мера на пространстве последовательностей.
  Цилиндрическое событие
\end{itemize}

\[
C=\{x\in\mathbb R^{\mathbb N}:x_1\le a,\ x_2\in B,\ x_5\le c\}
\]

имеет вероятность

\[
P_X(C)=P\{X_1\le a,\ X_2\in B,\ X_5\le c\}.
\]

\subsection{6. Доказательства и
обоснования}\label{ux434ux43eux43aux430ux437ux430ux442ux435ux43bux44cux441ux442ux432ux430-ux438-ux43eux431ux43eux441ux43dux43eux432ux430ux43dux438ux44f-5}

\textbf{Почему распределение называют образом меры.}

Формула

\[
P_X(B)=P(X^{-1}(B))
\]

переносит меру \(P\) с исходного пространства на пространство значений.
Поэтому \(P_X\) называют образом меры \(P\) при отображении \(X\) и
пишут

\[
P_X=P\circ X^{-1}.
\]

\textbf{Почему одной функции распределения хватает на прямой.}

На прямой борелевская \(\sigma\)-алгебра порождается лучами
\((-\infty,x]\). Функция \(F_X\) задает вероятности всех таких лучей.
Так как вероятностная мера на порождающем классе при стандартных
условиях продолжается единственным образом, \(F_X\) задает все
распределение.

\textbf{Почему для вектора нужны совместные распределения, а не только
маргинальные.}

Маргинальные распределения \(X\) и \(Y\) не определяют совместное
распределение \((X,Y)\). Например, если \(X\) симметрична и принимает
значения \(\pm1\), то пары

\[
(X,X)
\quad\text{и}\quad
(X,-X)
\]

имеют одинаковые одномерные маргинальные распределения, но разные
совместные распределения.

\textbf{Почему плотность не является вероятностью точки.}

Для непрерывного распределения с плотностью

\[
P\{X=a\}=\int_{\{a\}}p(x)\,dx=0.
\]

Значение \(p(a)\) может быть больше \(1\) и не является вероятностью.
Вероятности получаются только после интегрирования плотности по
множеству положительной длины или объема.

\textbf{Почему комплексный случай сводится к вещественному векторному.}

Пространство \(\mathbb C^m\) естественно отождествляется с
\(\mathbb R^{2m}\):

\[
(z_1,\ldots,z_m)
\leftrightarrow
(\operatorname{Re}z_1,\operatorname{Im}z_1,\ldots,\operatorname{Re}z_m,\operatorname{Im}z_m).
\]

Поэтому распределения комплексных величин - это обычные распределения
вещественных векторов четной размерности.

\textbf{Почему конечномерные распределения процесса должны быть
согласованы.}

Если распределение вектора \((X_{t_1},X_{t_2})\) известно, то
распределение одной координаты \(X_{t_1}\) получается из него
автоматически:

\[
P\{X_{t_1}\in B\}=P\{(X_{t_1},X_{t_2})\in B\times\mathbb R\}.
\]

Поэтому нельзя произвольно назначать распределения для разных наборов
времен: они должны давать одинаковые ответы на пересекающихся наборах
координат.

\subsection{7. Что написать на
доске}\label{ux447ux442ux43e-ux43dux430ux43fux438ux441ux430ux442ux44c-ux43dux430-ux434ux43eux441ux43aux435-5}

\begin{enumerate}
\def\labelenumi{\arabic{enumi}.}
\tightlist
\item
  Распределение как образ меры:
\end{enumerate}

\[
P_X=P\circ X^{-1},
\qquad
P_X(B)=P\{X\in B\}.
\]

\begin{enumerate}
\def\labelenumi{\arabic{enumi}.}
\setcounter{enumi}{1}
\tightlist
\item
  Функция распределения:
\end{enumerate}

\[
F_X(x)=P\{X\le x\}.
\]

\begin{enumerate}
\def\labelenumi{\arabic{enumi}.}
\setcounter{enumi}{2}
\tightlist
\item
  Случайный вектор:
\end{enumerate}

\[
F_X(x_1,\ldots,x_d)
=
P\{X_1\le x_1,\ldots,X_d\le x_d\}.
\]

\begin{enumerate}
\def\labelenumi{\arabic{enumi}.}
\setcounter{enumi}{3}
\tightlist
\item
  Конечномерные распределения процесса:
\end{enumerate}

\[
P_{t_1,\ldots,t_n}(B)
=
P\{(X_{t_1},\ldots,X_{t_n})\in B\}.
\]

\begin{enumerate}
\def\labelenumi{\arabic{enumi}.}
\setcounter{enumi}{4}
\tightlist
\item
  Конечномерная функция распределения:
\end{enumerate}

\[
F_{t_1,\ldots,t_n}(x_1,\ldots,x_n)
=
P\{X_{t_1}\le x_1,\ldots,X_{t_n}\le x_n\}.
\]

\begin{enumerate}
\def\labelenumi{\arabic{enumi}.}
\setcounter{enumi}{5}
\tightlist
\item
  Плотность:
\end{enumerate}

\[
p\ge0,\qquad
\int_{\mathbb R^d}p(x)\,dx=1,\qquad
P_X(B)=\int_Bp(x)\,dx.
\]

\begin{enumerate}
\def\labelenumi{\arabic{enumi}.}
\setcounter{enumi}{6}
\tightlist
\item
  Маргинализация:
\end{enumerate}

\[
P_{X_1}(B)=P_{(X_1,X_2)}(B\times\mathbb R),
\qquad
p_X(x)=\int p_{X,Y}(x,y)\,dy.
\]

\subsection{8. Типичные дополнительные
вопросы}\label{ux442ux438ux43fux438ux447ux43dux44bux435-ux434ux43eux43fux43eux43bux43dux438ux442ux435ux43bux44cux43dux44bux435-ux432ux43eux43fux440ux43eux441ux44b-5}

\textbf{Вопрос. Что такое распределение случайной величины?}

Это вероятностная мера на пространстве значений:

\[
P_X(B)=P\{X\in B\}.
\]

Для вещественной случайной величины это мера на
\(\mathcal B(\mathbb R)\).

\textbf{Вопрос. Почему формула \(P_X(B)=P\{X\in B\}\) корректна?}

Потому что \(X\) измерима, значит \(\{X\in B\}=X^{-1}(B)\in\mathcal F\)
для каждого \(B\) из \(\sigma\)-алгебры пространства значений.

\textbf{Вопрос. Чем распределение отличается от функции распределения?}

Распределение - это мера \(P_X\). Функция распределения - числовая
функция \(F_X(x)=P\{X\le x\}\), которая в вещественном случае задает эту
меру.

\textbf{Вопрос. Чем распределение отличается от плотности?}

Распределение - мера. Плотность - функция, которую нужно интегрировать,
чтобы получить вероятности:

\[
P_X(B)=\int_Bp(x)\,dx.
\]

Плотность существует не у всякого распределения.

\textbf{Вопрос. Может ли плотность быть больше единицы?}

Да. Плотность не является вероятностью. Например, равномерное
распределение на \([0,\tfrac12]\) имеет плотность \(2\) на этом
интервале.

\textbf{Вопрос. Как получить распределение координаты из совместного
распределения?}

Нужно взять проекцию или отправить остальные координаты в \(\mathbb R\):

\[
P_{X_1}(B)=P_{(X_1,X_2)}(B\times\mathbb R).
\]

\textbf{Вопрос. Задают ли маргинальные распределения совместное?}

Нет. Они не содержат информации о зависимости между координатами. Для
совместного закона нужна совместная мера или совместная функция
распределения.

\textbf{Вопрос. Как понимать распределение комплексной случайной
величины?}

Как распределение двумерного вещественного вектора
\((\operatorname{Re}Z,\operatorname{Im}Z)\) на плоскости.

\textbf{Вопрос. Что такое конечномерные распределения процесса?}

Это распределения всех конечных наборов координат:

\[
(X_{t_1},\ldots,X_{t_n}).
\]

\textbf{Вопрос. Почему конечномерные распределения должны быть
согласованы?}

Потому что распределение меньшего набора координат должно получаться из
распределения большего набора маргинализацией.

\subsection{9. Типичные
ошибки}\label{ux442ux438ux43fux438ux447ux43dux44bux435-ux43eux448ux438ux431ux43aux438-5}

\begin{enumerate}
\def\labelenumi{\arabic{enumi}.}
\item
  Путать случайную величину \(X\), ее распределение \(P_X\), функцию
  распределения \(F_X\) и плотность \(p\).
\item
  Говорить о плотности у любого распределения. У дискретного
  распределения относительно меры Лебега плотности нет.
\item
  Считать значение плотности вероятностью: \(p(x)\) - не вероятность
  точки.
\item
  Забывать, что комплексная случайная величина задается двумя
  вещественными координатами.
\item
  Думать, что одномерные маргиналы задают совместное распределение.
\item
  Забывать согласованность конечномерных распределений процесса.
\item
  Использовать формулу \(F_X(x)=P\{X\le x\}\) для объектов без порядка,
  например для общего метрического пространства. Там основной объект -
  мера \(P_X\).
\end{enumerate}

\subsection{10. Связи с другими
билетами}\label{ux441ux432ux44fux437ux438-ux441-ux434ux440ux443ux433ux438ux43cux438-ux431ux438ux43bux435ux442ux430ux43cux438-5}

\textbf{Билет 02.} Функция распределения на прямой задает вероятностную
меру; это основа перехода от \(F\) к \(P_X\).

\textbf{Билет 04.} Распределения вещественных и векторных случайных
величин живут на борелевских \(\sigma\)-алгебрах.

\textbf{Билет 05.} Измеримость случайного элемента нужна, чтобы формула
\(P_X(B)=P(X^{-1}(B))\) была корректной.

\textbf{Билет 07.} Интегралы по плотности и математические ожидания
требуют аппарата интеграла Лебега.

\textbf{Билет 08.} Характеристическая функция является преобразованием
Фурье распределения.

\textbf{Билет 09.} Независимость случайных величин выражается через
факторизацию совместного распределения.

\textbf{Билет 10.} Плотность в общем виде описывается теоремой
Радона-Никодима.

\textbf{Билет 11.} Замена переменной связывает интегралы по
вероятностному пространству и по распределению.

\textbf{Билет 15.} Согласованные конечномерные распределения задают
случайный процесс.

\textbf{Билет 16.} Дискретный случайный процесс можно рассматривать как
случайный элемент пространства последовательностей; его закон задается
распределением на этом пространстве.

\subsection{11. Минимум на
удовлетворительно}\label{ux43cux438ux43dux438ux43cux443ux43c-ux43dux430-ux443ux434ux43eux432ux43bux435ux442ux432ux43eux440ux438ux442ux435ux43bux44cux43dux43e-5}

Нужно уметь сказать:

\begin{enumerate}
\def\labelenumi{\arabic{enumi}.}
\tightlist
\item
  Распределение случайного элемента:
\end{enumerate}

\[
P_X(B)=P\{X\in B\}.
\]

\begin{enumerate}
\def\labelenumi{\arabic{enumi}.}
\setcounter{enumi}{1}
\tightlist
\item
  Для вещественной случайной величины:
\end{enumerate}

\[
F_X(x)=P\{X\le x\}.
\]

\begin{enumerate}
\def\labelenumi{\arabic{enumi}.}
\setcounter{enumi}{2}
\tightlist
\item
  Для случайного вектора:
\end{enumerate}

\[
F_X(x_1,\ldots,x_d)
=
P\{X_1\le x_1,\ldots,X_d\le x_d\}.
\]

\begin{enumerate}
\def\labelenumi{\arabic{enumi}.}
\setcounter{enumi}{3}
\item
  Для процесса конечномерные распределения - это распределения векторов
  \((X_{t_1},\ldots,X_{t_n})\).
\item
  При плотности:
\end{enumerate}

\[
P_X(B)=\int_Bp(x)\,dx.
\]

\subsection{12. Минимум на
хорошо}\label{ux43cux438ux43dux438ux43cux443ux43c-ux43dux430-ux445ux43eux440ux43eux448ux43e-5}

Кроме минимума на удовлетворительно, нужно:

\begin{enumerate}
\def\labelenumi{\arabic{enumi}.}
\tightlist
\item
  Объяснить, почему \(P_X\) является вероятностной мерой.
\item
  Назвать свойства функции распределения: монотонность, правую
  непрерывность, пределы \(0\) и \(1\).
\item
  Показать, как вероятности интервалов выражаются через \(F\).
\item
  Объяснить разницу между распределением, функцией распределения и
  плотностью.
\item
  Рассказать, как получить маргинальное распределение из совместного.
\end{enumerate}

\subsection{13. Что добавить для
отлично}\label{ux447ux442ux43e-ux434ux43eux431ux430ux432ux438ux442ux44c-ux434ux43bux44f-ux43eux442ux43bux438ux447ux43dux43e-5}

Для сильного ответа стоит добавить:

\begin{enumerate}
\def\labelenumi{\arabic{enumi}.}
\tightlist
\item
  Доказательство переноса меры через прообраз.
\item
  Формулу скачка:
\end{enumerate}

\[
P\{X=a\}=F(a)-F(a-).
\]

\begin{enumerate}
\def\labelenumi{\arabic{enumi}.}
\setcounter{enumi}{2}
\tightlist
\item
  Комплексный случай как распределение вещественного вектора из
  действительных и мнимых частей.
\item
  Согласованность конечномерных распределений процесса.
\item
  Предупреждение, что маргинальные распределения не определяют
  совместное.
\item
  Уточнение: непрерывная плотность позволяет считать вероятности через
  интеграл Римана на прямоугольниках и хороших областях, но общий язык
  плотностей - мера Лебега.
\end{enumerate}

\subsection{14. Устная версия ответа на 5
минут}\label{ux443ux441ux442ux43dux430ux44f-ux432ux435ux440ux441ux438ux44f-ux43eux442ux432ux435ux442ux430-ux43dux430-5-ux43cux438ux43dux443ux442-5}

Распределение случайного элемента - это способ перенести вероятность с
пространства исходов на пространство значений. Пусть

\[
X:(\Omega,\mathcal F,P)\to(E,\mathcal E)
\]

измерим. Тогда его распределение:

\[
P_X(B)=P\{X\in B\}=P(X^{-1}(B)),\qquad B\in\mathcal E.
\]

Это вероятностная мера на \((E,\mathcal E)\). Измеримость нужна, чтобы
\(X^{-1}(B)\) было событием.

Для вещественной случайной величины распределение можно задавать
функцией распределения:

\[
F_X(x)=P\{X\le x\}.
\]

Она не убывает, непрерывна справа, имеет предел \(0\) при
\(x\to-\infty\) и предел \(1\) при \(x\to+\infty\). Вероятность
полуинтервала выражается как

\[
P\{a<X\le b\}=F(b)-F(a).
\]

Масса точки равна скачку:

\[
P\{X=a\}=F(a)-F(a-).
\]

Для случайного вектора \(X=(X_1,\ldots,X_d)\) распределение - мера на
\(\mathcal B(\mathbb R^d)\), а совместная функция распределения:

\[
F_X(x_1,\ldots,x_d)
=
P\{X_1\le x_1,\ldots,X_d\le x_d\}.
\]

Комплексная случайная величина рассматривается как двумерный
вещественный вектор \((\operatorname{Re}Z,\operatorname{Im}Z)\).

Для процесса \(\{X_t:t\in T\}\) конечномерные распределения - это
распределения векторов

\[
(X_{t_1},\ldots,X_{t_n})
\]

для всех конечных наборов времен. Они должны быть согласованы при
перестановке и вычеркивании координат.

Если распределение имеет плотность \(p\), то

\[
p\ge0\text{ почти всюду},\qquad
\int p(x)\,dx=1,\qquad
P_X(B)=\int_Bp(x)\,dx.
\]

В одномерном случае

\[
F(x)=\int_{-\infty}^{x}p(u)\,du.
\]

Важно не путать распределение, функцию распределения и плотность:
распределение - мера, функция распределения - способ задать меру на
прямой, плотность - функция под интегралом и существует не всегда.

\subsection{15. Если осталось 2 минуты перед
экзаменом}\label{ux435ux441ux43bux438-ux43eux441ux442ux430ux43bux43eux441ux44c-2-ux43cux438ux43dux443ux442ux44b-ux43fux435ux440ux435ux434-ux44dux43aux437ux430ux43cux435ux43dux43eux43c-5}

Запомнить:

\[
P_X=P\circ X^{-1},
\qquad
P_X(B)=P\{X\in B\}.
\]

Для вещественной случайной величины:

\[
F_X(x)=P\{X\le x\}.
\]

Для вектора:

\[
F_X(x_1,\ldots,x_d)
=
P\{X_1\le x_1,\ldots,X_d\le x_d\}.
\]

Для процесса:

\[
F_{t_1,\ldots,t_n}(x_1,\ldots,x_n)
=
P\{X_{t_1}\le x_1,\ldots,X_{t_n}\le x_n\}.
\]

При плотности:

\[
P_X(B)=\int_Bp(x)\,dx,\qquad
p\ge0,\qquad
\int p=1.
\]

Главная ловушка: \(P_X\), \(F_X\) и \(p\) - разные объекты. Плотность
есть не у всякого распределения.

\newpage

\section{Билет 07. Простые функции, интеграл Лебега и теоремы Леви и
Лебега о предельном
переходе}\label{ux431ux438ux43bux435ux442-07.-ux43fux440ux43eux441ux442ux44bux435-ux444ux443ux43dux43aux446ux438ux438-ux438ux43dux442ux435ux433ux440ux430ux43b-ux43bux435ux431ux435ux433ux430-ux438-ux442ux435ux43eux440ux435ux43cux44b-ux43bux435ux432ux438-ux438-ux43bux435ux431ux435ux433ux430-ux43e-ux43fux440ux435ux434ux435ux43bux44cux43dux43eux43c-ux43fux435ux440ux435ux445ux43eux434ux435}

Источники: Ширяев 1, гл. II, §6; лекции ДСП.

\subsection{1. Короткий
ответ}\label{ux43aux43eux440ux43eux442ux43aux438ux439-ux43eux442ux432ux435ux442-6}

Интеграл Лебега по неотрицательной мере строится в три шага: сначала для
индикаторов, затем для простых неотрицательных функций, затем для
произвольных неотрицательных измеримых функций через супремум по простым
минорантам.

Если

\[
s=\sum_{k=1}^n a_k\mathbf 1_{A_k},
\qquad
a_k\ge0,\quad A_k\in\mathcal F,
\]

и множества \(A_k\) попарно не пересекаются, то

\[
\int s\,d\mu=\sum_{k=1}^n a_k\mu(A_k).
\]

Для неотрицательной измеримой функции \(f\):

\[
\int f\,d\mu
=
\sup\left\{\int s\,d\mu:\ 0\le s\le f,\ s\text{ простая}\right\}.
\]

Теорема Беппо Леви: если \(0\le f_n\uparrow f\), то

\[
\int f_n\,d\mu\uparrow\int f\,d\mu.
\]

Мажорантная теорема Лебега: если \(f_n\to f\) почти всюду и
\(|f_n|\le g\), где \(g\in L^1\), то

\[
\int f_n\,d\mu\to\int f\,d\mu.
\]

\subsection{2. Полный
ответ}\label{ux43fux43eux43bux43dux44bux439-ux43eux442ux432ux435ux442-6}

\subsubsection{2.1. Введение и роль интеграла
Лебега}\label{ux432ux432ux435ux434ux435ux43dux438ux435-ux438-ux440ux43eux43bux44c-ux438ux43dux442ux435ux433ux440ux430ux43bux430-ux43bux435ux431ux435ux433ux430}

В предыдущих билетах были введены измеримые функции и распределения.
Теперь нужен аппарат, который позволяет интегрировать измеримые функции
по мере. В вероятности этот интеграл становится математическим
ожиданием, а в дальнейшем через него выражаются моменты,
характеристические функции, условные ожидания, корреляции и спектральные
характеристики.

\subsubsection{2.2. Измеримое пространство и интеграл от индикаторной
функции}\label{ux438ux437ux43cux435ux440ux438ux43cux43eux435-ux43fux440ux43eux441ux442ux440ux430ux43dux441ux442ux432ux43e-ux438-ux438ux43dux442ux435ux433ux440ux430ux43b-ux43eux442-ux438ux43dux434ux438ux43aux430ux442ux43eux440ux43dux43eux439-ux444ux443ux43dux43aux446ux438ux438}

Пусть задано пространство с неотрицательной мерой

\[
(\Omega,\mathcal F,\mu).
\]

Неотрицательность меры означает, что

\[
\mu(A)\in[0,\infty]
\qquad
A\in\mathcal F.
\]

Сначала интеграл определяют для самой простой измеримой функции -
индикатора множества:

\[
\int \mathbf 1_A\,d\mu=\mu(A).
\]

Это естественно: индикатор равен \(1\) на \(A\) и \(0\) вне \(A\),
значит площадь под ним должна равняться мере множества \(A\).

\subsubsection{2.3. Интеграл от простых
функций}\label{ux438ux43dux442ux435ux433ux440ux430ux43b-ux43eux442-ux43fux440ux43eux441ux442ux44bux445-ux444ux443ux43dux43aux446ux438ux439}

Дальше рассматривают простые функции. Простая функция - измеримая
функция, принимающая конечное число значений. Если \(s\ge0\) и

\[
s=\sum_{k=1}^n a_k\mathbf 1_{A_k},
\]

где \(a_k\ge0\), \(A_k\in\mathcal F\) и \(A_k\) попарно не пересекаются,
то интеграл задают формулой

\[
\int s\,d\mu=\sum_{k=1}^n a_k\mu(A_k).
\]

Эта формула не зависит от выбранного представления простой функции: если
разложение измельчить до общего конечного разбиения, сумма останется той
же. Поэтому интеграл простой функции определен корректно.

\subsubsection{2.4. Построение интеграла для неотрицательных измеримых
функций}\label{ux43fux43eux441ux442ux440ux43eux435ux43dux438ux435-ux438ux43dux442ux435ux433ux440ux430ux43bux430-ux434ux43bux44f-ux43dux435ux43eux442ux440ux438ux446ux430ux442ux435ux43bux44cux43dux44bux445-ux438ux437ux43cux435ux440ux438ux43cux44bux445-ux444ux443ux43dux43aux446ux438ux439}

Для произвольной неотрицательной измеримой функции \(f\) интеграл строят
приближением снизу простыми функциями:

\[
\int f\,d\mu
=
\sup\left\{\int s\,d\mu:\ 0\le s\le f,\ s\text{ простая измеримая}\right\}.
\]

Интуиция такая: сложную функцию можно снизу аппроксимировать
ступенчатыми измеримыми функциями. Интеграл \(f\) - это наибольший
предел площадей таких нижних ступенчатых приближений. Значение интеграла
может быть бесконечным.

\subsubsection{2.5. Функции произвольного знака и условие
интегрируемости}\label{ux444ux443ux43dux43aux446ux438ux438-ux43fux440ux43eux438ux437ux432ux43eux43bux44cux43dux43eux433ux43e-ux437ux43dux430ux43aux430-ux438-ux443ux441ux43bux43eux432ux438ux435-ux438ux43dux442ux435ux433ux440ux438ux440ux443ux435ux43cux43eux441ux442ux438}

Если функция \(f\) может принимать отрицательные значения, ее
раскладывают на положительную и отрицательную части:

\[
f^+=\max(f,0),\qquad
f^-=\max(-f,0),
\]

так что

\[
f=f^+-f^-,
\qquad
|f|=f^++f^-.
\]

Если хотя бы один из интегралов \(\int f^+\,d\mu\), \(\int f^-\,d\mu\)
конечен, можно определить расширенный интеграл

\[
\int f\,d\mu=\int f^+\,d\mu-\int f^-\,d\mu,
\]

не допуская неопределенности \(\infty-\infty\). Если

\[
\int |f|\,d\mu<\infty,
\]

то \(f\) называется интегрируемой, или суммируемой, и пишут
\(f\in L^1(\mu)\).

\subsubsection{2.6. Теорема Беппо Леви о монотонной
сходимости}\label{ux442ux435ux43eux440ux435ux43cux430-ux431ux435ux43fux43fux43e-ux43bux435ux432ux438-ux43e-ux43cux43eux43dux43eux442ux43eux43dux43dux43eux439-ux441ux445ux43eux434ux438ux43cux43eux441ux442ux438}

Главная сила интеграла Лебега - хорошие теоремы о предельном переходе. В
программе эти теоремы идут без доказательства, но условия нужно знать
точно.

Первая - теорема Беппо Леви, или теорема о монотонной сходимости. Если

\[
0\le f_1\le f_2\le\cdots
\]

и

\[
f_n(\omega)\uparrow f(\omega)
\]

для всех \(\omega\) или почти всюду, то

\[
\int f_n\,d\mu\uparrow\int f\,d\mu.
\]

Здесь не нужен интегрируемый мажорант; зато нужна неотрицательность и
монотонный рост.

\subsubsection{2.7. Мажорантная теорема Лебега о доминированной
сходимости}\label{ux43cux430ux436ux43eux440ux430ux43dux442ux43dux430ux44f-ux442ux435ux43eux440ux435ux43cux430-ux43bux435ux431ux435ux433ux430-ux43e-ux434ux43eux43cux438ux43dux438ux440ux43eux432ux430ux43dux43dux43eux439-ux441ux445ux43eux434ux438ux43cux43eux441ux442ux438}

Вторая - мажорантная теорема Лебега, или теорема Лебега о доминированной
сходимости. Если

\[
f_n\to f
\]

почти всюду и существует интегрируемая функция \(g\in L^1(\mu)\) такая,
что

\[
|f_n|\le g
\qquad \text{почти всюду для всех }n,
\]

то \(f\) интегрируема и

\[
\int f_n\,d\mu\to\int f\,d\mu.
\]

Эта теорема позволяет менять предел и интеграл при наличии общего
интегрируемого контроля. Без такого контроля утверждение может быть
неверным.

\subsection{3. Основные
определения}\label{ux43eux441ux43dux43eux432ux43dux44bux435-ux43eux43fux440ux435ux434ux435ux43bux435ux43dux438ux44f-6}

\subsubsection{Пространство с неотрицательной
мерой}\label{ux43fux440ux43eux441ux442ux440ux430ux43dux441ux442ux432ux43e-ux441-ux43dux435ux43eux442ux440ux438ux446ux430ux442ux435ux43bux44cux43dux43eux439-ux43cux435ux440ux43eux439}

Тройка \((\Omega,\mathcal F,\mu)\), где \(\mathcal F\) -
\(\sigma\)-алгебра, а \(\mu:\mathcal F\to[0,\infty]\) -
счетно-аддитивная неотрицательная мера.

\subsubsection{Индикатор
множества}\label{ux438ux43dux434ux438ux43aux430ux442ux43eux440-ux43cux43dux43eux436ux435ux441ux442ux432ux430}

Для \(A\in\mathcal F\):

\[
\mathbf 1_A(\omega)=
\begin{cases}
1, & \omega\in A,\\
0, & \omega\notin A.
\end{cases}
\]

\subsubsection{Простая
функция}\label{ux43fux440ux43eux441ux442ux430ux44f-ux444ux443ux43dux43aux446ux438ux44f}

Измеримая функция \(s:\Omega\to\mathbb R\) называется простой, если она
принимает конечное число значений. Ее можно записать в виде

\[
s=\sum_{k=1}^n a_k\mathbf 1_{A_k},
\]

где \(A_k\in\mathcal F\). В каноническом представлении значения \(a_k\)
различны, а множества

\[
A_k=\{\omega:s(\omega)=a_k\}
\]

попарно не пересекаются и образуют разбиение области, где \(s\)
принимает эти значения.

\subsubsection{Неотрицательная простая
функция}\label{ux43dux435ux43eux442ux440ux438ux446ux430ux442ux435ux43bux44cux43dux430ux44f-ux43fux440ux43eux441ux442ux430ux44f-ux444ux443ux43dux43aux446ux438ux44f}

Простая функция \(s\), для которой \(s(\omega)\ge0\) для всех
\(\omega\). Если неравенство выполнено только почти всюду, функцию можно
заменить равной ей почти всюду неотрицательной модификацией.

\subsubsection{Интеграл
индикатора}\label{ux438ux43dux442ux435ux433ux440ux430ux43b-ux438ux43dux434ux438ux43aux430ux442ux43eux440ux430}

\[
\int \mathbf 1_A\,d\mu=\mu(A).
\]

\subsubsection{Интеграл неотрицательной простой
функции}\label{ux438ux43dux442ux435ux433ux440ux430ux43b-ux43dux435ux43eux442ux440ux438ux446ux430ux442ux435ux43bux44cux43dux43eux439-ux43fux440ux43eux441ux442ux43eux439-ux444ux443ux43dux43aux446ux438ux438}

Если

\[
s=\sum_{k=1}^n a_k\mathbf 1_{A_k},
\qquad
a_k\ge0,
\]

и \(A_k\) попарно не пересекаются, то

\[
\int s\,d\mu=\sum_{k=1}^n a_k\mu(A_k).
\]

\subsubsection{Интеграл неотрицательной измеримой
функции}\label{ux438ux43dux442ux435ux433ux440ux430ux43b-ux43dux435ux43eux442ux440ux438ux446ux430ux442ux435ux43bux44cux43dux43eux439-ux438ux437ux43cux435ux440ux438ux43cux43eux439-ux444ux443ux43dux43aux446ux438ux438}

Для \(f:\Omega\to[0,\infty]\):

\[
\int f\,d\mu
=
\sup\left\{\int s\,d\mu:\ 0\le s\le f,\ s\text{ простая измеримая}\right\}.
\]

\subsubsection{Положительная и отрицательная части
функции}\label{ux43fux43eux43bux43eux436ux438ux442ux435ux43bux44cux43dux430ux44f-ux438-ux43eux442ux440ux438ux446ux430ux442ux435ux43bux44cux43dux430ux44f-ux447ux430ux441ux442ux438-ux444ux443ux43dux43aux446ux438ux438}

\[
f^+=\max(f,0),\qquad
f^-=\max(-f,0).
\]

Тогда

\[
f=f^+-f^-,
\qquad
|f|=f^++f^-.
\]

\subsubsection{Интегрируемая
функция}\label{ux438ux43dux442ux435ux433ux440ux438ux440ux443ux435ux43cux430ux44f-ux444ux443ux43dux43aux446ux438ux44f}

Измеримая функция \(f\) называется интегрируемой, если

\[
\int |f|\,d\mu<\infty.
\]

В этом случае

\[
\int f\,d\mu=\int f^+\,d\mu-\int f^-\,d\mu
\]

и оба интеграла справа конечны.

\subsubsection{Почти
всюду}\label{ux43fux43eux447ux442ux438-ux432ux441ux44eux434ux443}

Свойство выполнено почти всюду, если оно нарушается только на множестве
меры ноль.

\subsubsection{Мажоранта}\label{ux43cux430ux436ux43eux440ux430ux43dux442ux430}

Функция \(g\) называется интегрируемой мажорантой для последовательности
\(f_n\), если

\[
|f_n|\le g\quad \text{почти всюду для всех }n,
\qquad
g\in L^1(\mu).
\]

\subsection{4. Основные теоремы и
свойства}\label{ux43eux441ux43dux43eux432ux43dux44bux435-ux442ux435ux43eux440ux435ux43cux44b-ux438-ux441ux432ux43eux439ux441ux442ux432ux430-6}

\subsubsection{Свойство 1. Корректность интеграла простой
функции}\label{ux441ux432ux43eux439ux441ux442ux432ux43e-1.-ux43aux43eux440ux440ux435ux43aux442ux43dux43eux441ux442ux44c-ux438ux43dux442ux435ux433ux440ux430ux43bux430-ux43fux440ux43eux441ux442ux43eux439-ux444ux443ux43dux43aux446ux438ux438}

Если неотрицательная простая функция имеет два представления

\[
s=\sum_{i=1}^m a_i\mathbf 1_{A_i}
=
\sum_{j=1}^n b_j\mathbf 1_{B_j},
\]

где множества в каждом представлении попарно не пересекаются, то

\[
\sum_{i=1}^m a_i\mu(A_i)
=
\sum_{j=1}^n b_j\mu(B_j).
\]

\subsubsection{Доказательство-идея}\label{ux434ux43eux43aux430ux437ux430ux442ux435ux43bux44cux441ux442ux432ux43e-ux438ux434ux435ux44f-6}

Перейдем к общему измельчению

\[
C_{ij}=A_i\cap B_j.
\]

На каждом непустом \(C_{ij}\) значения двух представлений совпадают, то
есть \(a_i=b_j\). Поэтому обе суммы после разбиения на \(C_{ij}\)
становятся одной и той же суммой.

\subsubsection{Свойство 2. Интеграл индикатора равен
мере}\label{ux441ux432ux43eux439ux441ux442ux432ux43e-2.-ux438ux43dux442ux435ux433ux440ux430ux43b-ux438ux43dux434ux438ux43aux430ux442ux43eux440ux430-ux440ux430ux432ux435ux43d-ux43cux435ux440ux435}

Для любого \(A\in\mathcal F\):

\[
\int \mathbf 1_A\,d\mu=\mu(A).
\]

Это базовая нормировка интеграла Лебега.

\subsubsection{Свойство 3. Монотонность
интеграла}\label{ux441ux432ux43eux439ux441ux442ux432ux43e-3.-ux43cux43eux43dux43eux442ux43eux43dux43dux43eux441ux442ux44c-ux438ux43dux442ux435ux433ux440ux430ux43bux430}

Если \(0\le f\le g\), то

\[
\int f\,d\mu\le\int g\,d\mu.
\]

\subsubsection{Доказательство}\label{ux434ux43eux43aux430ux437ux430ux442ux435ux43bux44cux441ux442ux432ux43e-28}

Любая простая функция \(s\), удовлетворяющая \(0\le s\le f\),
автоматически удовлетворяет \(0\le s\le g\). Поэтому супремум по простым
минорантам для \(f\) не больше супремума для \(g\).

\subsubsection{Свойство 4. Линейность для неотрицательных
функций}\label{ux441ux432ux43eux439ux441ux442ux432ux43e-4.-ux43bux438ux43dux435ux439ux43dux43eux441ux442ux44c-ux434ux43bux44f-ux43dux435ux43eux442ux440ux438ux446ux430ux442ux435ux43bux44cux43dux44bux445-ux444ux443ux43dux43aux446ux438ux439}

Если \(f,g\ge0\) измеримы и \(a,b\ge0\), то

\[
\int(af+bg)\,d\mu
=
a\int f\,d\mu+b\int g\,d\mu.
\]

Здесь допускаются бесконечные значения, если не возникает противоречивых
операций.

\subsubsection{Доказательство-идея}\label{ux434ux43eux43aux430ux437ux430ux442ux435ux43bux44cux441ux442ux432ux43e-ux438ux434ux435ux44f-7}

Для простых функций это следует из формулы суммы по разбиению. Для общих
неотрицательных функций берут возрастающие последовательности простых
функций, приближающие \(f\) и \(g\) снизу, и применяют теорему Беппо
Леви.

\subsubsection{\texorpdfstring{Свойство 5. Линейность на
\(L^1\)}{Свойство 5. Линейность на L\^{}1}}\label{ux441ux432ux43eux439ux441ux442ux432ux43e-5.-ux43bux438ux43dux435ux439ux43dux43eux441ux442ux44c-ux43dux430-l1}

Если \(f,g\in L^1(\mu)\) и \(a,b\in\mathbb R\), то

\[
\int(af+bg)\,d\mu
=
a\int f\,d\mu+b\int g\,d\mu.
\]

\subsubsection{Обоснование}\label{ux43eux431ux43eux441ux43dux43eux432ux430ux43dux438ux435-4}

Интегрируемость гарантирует конечность интегралов положительных и
отрицательных частей, поэтому не возникает неопределенности
\(\infty-\infty\).

\subsubsection{Свойство 6. Если функция равна нулю почти всюду, то ее
интеграл равен
нулю}\label{ux441ux432ux43eux439ux441ux442ux432ux43e-6.-ux435ux441ux43bux438-ux444ux443ux43dux43aux446ux438ux44f-ux440ux430ux432ux43dux430-ux43dux443ux43bux44e-ux43fux43eux447ux442ux438-ux432ux441ux44eux434ux443-ux442ux43e-ux435ux435-ux438ux43dux442ux435ux433ux440ux430ux43b-ux440ux430ux432ux435ux43d-ux43dux443ux43bux44e}

Если \(f\ge0\) и \(f=0\) почти всюду, то

\[
\int f\,d\mu=0.
\]

Если \(f=g\) почти всюду и обе функции интегрируемы, то

\[
\int f\,d\mu=\int g\,d\mu.
\]

\subsubsection{Свойство 7. Аппроксимация неотрицательной измеримой
функции
простыми}\label{ux441ux432ux43eux439ux441ux442ux432ux43e-7.-ux430ux43fux43fux440ux43eux43aux441ux438ux43cux430ux446ux438ux44f-ux43dux435ux43eux442ux440ux438ux446ux430ux442ux435ux43bux44cux43dux43eux439-ux438ux437ux43cux435ux440ux438ux43cux43eux439-ux444ux443ux43dux43aux446ux438ux438-ux43fux440ux43eux441ux442ux44bux43cux438}

Для любой неотрицательной измеримой функции \(f\) существует
последовательность неотрицательных простых функций \(s_n\) такая, что

\[
0\le s_1\le s_2\le\cdots\le f,
\qquad
s_n(\omega)\uparrow f(\omega).
\]

\subsubsection{Доказательство-идея}\label{ux434ux43eux43aux430ux437ux430ux442ux435ux43bux44cux441ux442ux432ux43e-ux438ux434ux435ux44f-8}

Можно взять ступенчатое округление вниз. Например, на уровне \(n\)
разбить значения \(f\) на интервалы длины \(2^{-n}\) до высоты \(n\) и
положить

\[
s_n(\omega)=\frac{k}{2^n}
\]

на множестве

\[
\left\{\frac{k}{2^n}\le f(\omega)<\frac{k+1}{2^n}\right\},
\]

а для \(f(\omega)\ge n\) ограничить значение сверху числом \(n\). Эти
множества измеримы, потому что \(f\) измерима.

\subsubsection{Теорема 8. Беппо Леви, или монотонная
сходимость}\label{ux442ux435ux43eux440ux435ux43cux430-8.-ux431ux435ux43fux43fux43e-ux43bux435ux432ux438-ux438ux43bux438-ux43cux43eux43dux43eux442ux43eux43dux43dux430ux44f-ux441ux445ux43eux434ux438ux43cux43eux441ux442ux44c}

Пусть \(f_n:\Omega\to[0,\infty]\) измеримы,

\[
0\le f_1\le f_2\le\cdots,
\]

и

\[
f_n(\omega)\uparrow f(\omega)
\]

почти всюду. Тогда

\[
\int f_n\,d\mu\uparrow\int f\,d\mu.
\]

В программе эта теорема идет без доказательства.

\subsubsection{Идея
доказательства}\label{ux438ux434ux435ux44f-ux434ux43eux43aux430ux437ux430ux442ux435ux43bux44cux441ux442ux432ux430}

Неравенство

\[
\lim_n\int f_n\,d\mu\le\int f\,d\mu
\]

следует из монотонности. Обратное неравенство получают так: берут
простую функцию \(s\le f\), немного уменьшают ее до \(\alpha s\), где
\(0<\alpha<1\), и используют множества, на которых \(f_n\) уже превысила
\(\alpha s\). Потом переходят к пределу по \(n\) и по
\(\alpha\uparrow1\).

\subsubsection{Теорема 9. Мажорантная теорема Лебега, или доминированная
сходимость}\label{ux442ux435ux43eux440ux435ux43cux430-9.-ux43cux430ux436ux43eux440ux430ux43dux442ux43dux430ux44f-ux442ux435ux43eux440ux435ux43cux430-ux43bux435ux431ux435ux433ux430-ux438ux43bux438-ux434ux43eux43cux438ux43dux438ux440ux43eux432ux430ux43dux43dux430ux44f-ux441ux445ux43eux434ux438ux43cux43eux441ux442ux44c}

Пусть \(f_n\) измеримы,

\[
f_n\to f
\]

почти всюду, и существует \(g\in L^1(\mu)\) такая, что

\[
|f_n|\le g
\qquad \text{почти всюду для всех }n.
\]

Тогда \(f\in L^1(\mu)\) и

\[
\int f_n\,d\mu\to\int f\,d\mu.
\]

Более сильная форма вывода:

\[
\int |f_n-f|\,d\mu\to0.
\]

Из нее уже следует сходимость интегралов.

В программе эта теорема также идет без полного доказательства.

\subsubsection{Идея
доказательства}\label{ux438ux434ux435ux44f-ux434ux43eux43aux430ux437ux430ux442ux435ux43bux44cux441ux442ux432ux430-1}

Из \(|f_n|\le g\) и \(f_n\to f\) следует \(|f|\le g\), значит \(f\)
интегрируема. Далее к неотрицательным функциям

\[
2g-|f_n-f|
\]

или к верхним и нижним оценкам применяют лемму Фату. Смысл:
интегрируемая мажоранта запрещает массе убегать в бесконечность или
концентрироваться на все меньших множествах с большой высотой.

\subsubsection{Свойство 10. Нужны именно условия предельных
теорем}\label{ux441ux432ux43eux439ux441ux442ux432ux43e-10.-ux43dux443ux436ux43dux44b-ux438ux43cux435ux43dux43dux43e-ux443ux441ux43bux43eux432ux438ux44f-ux43fux440ux435ux434ux435ux43bux44cux43dux44bux445-ux442ux435ux43eux440ux435ux43c}

Если \(f_n\to f\) почти всюду, но нет ни монотонности, ни интегрируемой
мажоранты, то из этого не следует

\[
\int f_n\,d\mu\to\int f\,d\mu.
\]

Например, на \([0,1]\) с мерой Лебега функции

\[
f_n(x)=n\mathbf 1_{(0,1/n)}(x)
\]

сходятся к \(0\) почти всюду, но

\[
\int_0^1 f_n(x)\,dx=1
\]

для всех \(n\).

\subsection{5. Примеры}\label{ux43fux440ux438ux43cux435ux440ux44b-6}

\subsubsection{Пример 1.
Индикатор}\label{ux43fux440ux438ux43cux435ux440-1.-ux438ux43dux434ux438ux43aux430ux442ux43eux440}

Если \(A\in\mathcal F\), то

\[
\int \mathbf 1_A\,d\mu=\mu(A).
\]

В вероятностном пространстве \((\Omega,\mathcal F,P)\):

\[
\mathbb E\mathbf 1_A=P(A).
\]

\subsubsection{Пример 2. Простая
функция}\label{ux43fux440ux438ux43cux435ux440-2.-ux43fux440ux43eux441ux442ux430ux44f-ux444ux443ux43dux43aux446ux438ux44f}

Пусть \(A,B\in\mathcal F\), \(A\cap B=\varnothing\), и

\[
s=2\mathbf 1_A+5\mathbf 1_B.
\]

Тогда

\[
\int s\,d\mu=2\mu(A)+5\mu(B).
\]

Если \(\mu(A)=3\), \(\mu(B)=1\), то

\[
\int s\,d\mu=11.
\]

\subsubsection{Пример 3. Дискретное математическое
ожидание}\label{ux43fux440ux438ux43cux435ux440-3.-ux434ux438ux441ux43aux440ux435ux442ux43dux43eux435-ux43cux430ux442ux435ux43cux430ux442ux438ux447ux435ux441ux43aux43eux435-ux43eux436ux438ux434ux430ux43dux438ux435}

Пусть \(X\) принимает значения \(x_1,\ldots,x_n\) с вероятностями
\(p_1,\ldots,p_n\). Тогда

\[
X=\sum_{k=1}^n x_k\mathbf 1_{\{X=x_k\}},
\]

и

\[
\mathbb E X=\int X\,dP=\sum_{k=1}^n x_kp_k
\]

при условии, что \(X\) неотрицательна или интегрируема. Это обычная
формула ожидания для дискретной величины.

\subsubsection{Пример 4. Интеграл по
плотности}\label{ux43fux440ux438ux43cux435ux440-4.-ux438ux43dux442ux435ux433ux440ux430ux43b-ux43fux43e-ux43fux43bux43eux442ux43dux43eux441ux442ux438}

Если \(X\) имеет плотность \(p\) на \(\mathbb R\), то вероятность
события \(X\in B\) выражается как интеграл индикатора по распределению:

\[
P\{X\in B\}=\int \mathbf 1_B(x)p(x)\,dx.
\]

Это связывает билет с распределениями из Билета 06.

\subsubsection{Пример 5. Теорема Беппо
Леви}\label{ux43fux440ux438ux43cux435ux440-5.-ux442ux435ux43eux440ux435ux43cux430-ux431ux435ux43fux43fux43e-ux43bux435ux432ux438}

На \([0,1]\) пусть

\[
f_n(x)=\mathbf 1_{[0,1-1/n]}(x).
\]

Тогда

\[
0\le f_n\uparrow \mathbf 1_{[0,1)}
\]

и

\[
\int_0^1 f_n(x)\,dx=1-\frac1n
\uparrow
1
=
\int_0^1 \mathbf 1_{[0,1)}(x)\,dx.
\]

\subsubsection{Пример 6. Мажорантная теорема
Лебега}\label{ux43fux440ux438ux43cux435ux440-6.-ux43cux430ux436ux43eux440ux430ux43dux442ux43dux430ux44f-ux442ux435ux43eux440ux435ux43cux430-ux43bux435ux431ux435ux433ux430}

На \([0,1]\) пусть

\[
f_n(x)=x^n.
\]

Тогда \(f_n(x)\to0\) для \(x\in[0,1)\) и \(f_n(1)=1\). Значит
\(f_n\to0\) почти всюду. Кроме того,

\[
0\le f_n(x)\le1,
\]

а \(g(x)=1\) интегрируема на \([0,1]\). Поэтому

\[
\int_0^1 x^n\,dx\to0.
\]

И действительно,

\[
\int_0^1 x^n\,dx=\frac1{n+1}.
\]

\subsubsection{Пример 7. Почему мажоранта
нужна}\label{ux43fux440ux438ux43cux435ux440-7.-ux43fux43eux447ux435ux43cux443-ux43cux430ux436ux43eux440ux430ux43dux442ux430-ux43dux443ux436ux43dux430}

На \([0,1]\) функции

\[
f_n(x)=n\mathbf 1_{(0,1/n)}(x)
\]

сходятся к \(0\) почти всюду, но

\[
\int_0^1 f_n\,dx=1.
\]

Значит без интегрируемой мажоранты нельзя просто переносить предел под
интеграл.

\subsection{6. Доказательства и
обоснования}\label{ux434ux43eux43aux430ux437ux430ux442ux435ux43bux44cux441ux442ux432ux430-ux438-ux43eux431ux43eux441ux43dux43eux432ux430ux43dux438ux44f-6}

\textbf{Почему простые функции являются первым шагом.}

Простая функция принимает конечное число значений, поэтому ее интеграл
сводится к конечной сумме мер уровней. Это ровно тот случай, где мера
множеств напрямую превращается в интеграл функции.

\textbf{Почему нужно приближение снизу.}

Для неотрицательной функции \(f\) нельзя сразу сложить значения по всем
уровням, потому что значений может быть бесконечно много. Вместо этого
берут все простые функции \(s\), которые лежат ниже \(f\), и выбирают
супремум их интегралов:

\[
\int f\,d\mu=\sup_{0\le s\le f}\int s\,d\mu.
\]

Так определение сохраняет монотонность и согласуется с интегралом
простых функций.

\textbf{Почему интеграл может быть бесконечным.}

Для неотрицательных функций значение \(\int f\,d\mu\) допускается в
\([0,\infty]\). Например, на \(\mathbb R\):

\[
\int_{\mathbb R}1\,dx=\infty.
\]

Это не ошибка, а часть конструкции. Конечность требуется только для
интегрируемости в смысле \(L^1\).

\textbf{Почему для знакопеременных функций нужна абсолютная
интегрируемость.}

Если обе части \(f^+\) и \(f^-\) имеют бесконечный интеграл, выражение

\[
\int f^+\,d\mu-\int f^-\,d\mu
\]

имеет вид \(\infty-\infty\) и не определено. Условие

\[
\int |f|\,d\mu<\infty
\]

гарантирует, что обе части конечны.

\textbf{Почему теорема Леви не требует мажоранты.}

В теореме Леви функции растут монотонно снизу. Интегралы не могут
колебаться или терять массу: они тоже образуют возрастающую
последовательность. Поэтому предел интегралов совпадает с интегралом
предельной функции, возможно бесконечным.

\textbf{Почему теорема Лебега требует мажоранту.}

При обычной почти всюду сходимости графики могут становиться очень
высокими на очень маленьких множествах. Точечный предел этого не видит,
а интеграл видит площадь. Интегрируемая мажоранта \(g\) запрещает такие
пики и дает контроль:

\[
|f_n|\le g,\qquad \int g\,d\mu<\infty.
\]

\textbf{Что значит ``без доказательства'' в программе.}

Для теорем Беппо Леви и Лебега на экзамене достаточно дать точные
условия и вывод, а также объяснить идею. Полные доказательства можно не
разворачивать, если преподаватель специально не попросит.

\subsection{7. Что написать на
доске}\label{ux447ux442ux43e-ux43dux430ux43fux438ux441ux430ux442ux44c-ux43dux430-ux434ux43eux441ux43aux435-6}

\begin{enumerate}
\def\labelenumi{\arabic{enumi}.}
\tightlist
\item
  Простая функция:
\end{enumerate}

\[
s=\sum_{k=1}^n a_k\mathbf 1_{A_k}.
\]

\begin{enumerate}
\def\labelenumi{\arabic{enumi}.}
\setcounter{enumi}{1}
\tightlist
\item
  Интеграл простой неотрицательной функции:
\end{enumerate}

\[
\int s\,d\mu=\sum_{k=1}^n a_k\mu(A_k).
\]

\begin{enumerate}
\def\labelenumi{\arabic{enumi}.}
\setcounter{enumi}{2}
\tightlist
\item
  Интеграл неотрицательной функции:
\end{enumerate}

\[
\int f\,d\mu
=
\sup_{0\le s\le f,\ s\text{ простая}}
\int s\,d\mu.
\]

\begin{enumerate}
\def\labelenumi{\arabic{enumi}.}
\setcounter{enumi}{3}
\tightlist
\item
  Положительная и отрицательная части:
\end{enumerate}

\[
f=f^+-f^-,
\qquad
|f|=f^++f^-.
\]

\begin{enumerate}
\def\labelenumi{\arabic{enumi}.}
\setcounter{enumi}{4}
\tightlist
\item
  Интегрируемость:
\end{enumerate}

\[
f\in L^1(\mu)
\quad\Longleftrightarrow\quad
\int |f|\,d\mu<\infty.
\]

\begin{enumerate}
\def\labelenumi{\arabic{enumi}.}
\setcounter{enumi}{5}
\tightlist
\item
  Беппо Леви:
\end{enumerate}

\[
0\le f_n\uparrow f
\quad\Longrightarrow\quad
\int f_n\,d\mu\uparrow\int f\,d\mu.
\]

\begin{enumerate}
\def\labelenumi{\arabic{enumi}.}
\setcounter{enumi}{6}
\tightlist
\item
  Лебег:
\end{enumerate}

\[
f_n\to f,\quad |f_n|\le g,\quad g\in L^1
\quad\Longrightarrow\quad
\int f_n\,d\mu\to\int f\,d\mu.
\]

\subsection{8. Типичные дополнительные
вопросы}\label{ux442ux438ux43fux438ux447ux43dux44bux435-ux434ux43eux43fux43eux43bux43dux438ux442ux435ux43bux44cux43dux44bux435-ux432ux43eux43fux440ux43eux441ux44b-6}

\textbf{Вопрос. Что такое простая функция?}

Это измеримая функция, принимающая конечное число значений. Ее можно
записать как

\[
s=\sum_{k=1}^n a_k\mathbf 1_{A_k},
\qquad A_k\in\mathcal F.
\]

\textbf{Вопрос. Почему интеграл индикатора равен мере множества?}

Индикатор \(\mathbf 1_A\) равен \(1\) на \(A\) и \(0\) вне \(A\),
поэтому площадь под ним равна размеру множества \(A\):

\[
\int\mathbf 1_A\,d\mu=\mu(A).
\]

\textbf{Вопрос. Как определяется интеграл неотрицательной измеримой
функции?}

Через супремум интегралов простых минорант:

\[
\int f\,d\mu=\sup_{0\le s\le f}\int s\,d\mu.
\]

\textbf{Вопрос. Что говорит теорема Беппо Леви?}

Если \(0\le f_n\uparrow f\), то

\[
\int f_n\,d\mu\uparrow\int f\,d\mu.
\]

\textbf{Вопрос. Чем теорема Леви отличается от теоремы Лебега?}

В теореме Леви нужна монотонность и неотрицательность, но не нужна
мажоранта. В теореме Лебега монотонность не нужна, зато нужна
интегрируемая мажоранта \(g\).

\textbf{Вопрос. Что такое интегрируемая мажоранта?}

Это функция \(g\in L^1(\mu)\) такая, что

\[
|f_n|\le g
\]

почти всюду для всех \(n\).

\textbf{Вопрос. Почему нельзя применять теорему Лебега без мажоранты?}

Почти всюду сходимость не контролирует интегралы. Пример:

\[
f_n=n\mathbf 1_{(0,1/n)}\to0
\]

почти всюду, но

\[
\int_0^1 f_n\,dx=1.
\]

\textbf{Вопрос. Что значит \(f\in L^1\)?}

Это значит, что \(f\) измерима и

\[
\int |f|\,d\mu<\infty.
\]

\textbf{Вопрос. Почему для знакопеременной функции нельзя просто
интегрировать положительную и отрицательную части без условий?}

Потому что может возникнуть неопределенность \(\infty-\infty\).
Абсолютная интегрируемость исключает эту проблему.

\subsection{9. Типичные
ошибки}\label{ux442ux438ux43fux438ux447ux43dux44bux435-ux43eux448ux438ux431ux43aux438-6}

\begin{enumerate}
\def\labelenumi{\arabic{enumi}.}
\item
  Применять мажорантную теорему Лебега, не предъявив \(g\in L^1\).
\item
  Путать теорему Беппо Леви и теорему Лебега: в Леви нужна монотонность,
  в Лебеге - мажоранта.
\item
  Забывать, что для неотрицательной функции интеграл может быть равен
  \(\infty\).
\item
  Для знакопеременной функции писать \(\int f=\int f^+-\int f^-\), не
  проверив отсутствие неопределенности \(\infty-\infty\).
\item
  Считать, что почти всюду сходимость сама по себе гарантирует
  сходимость интегралов.
\item
  Записывать интеграл простой функции без требования измеримости
  множеств \(A_k\).
\item
  Забывать, что представление простой функции можно измельчать, и
  интеграл не зависит от конкретной записи.
\end{enumerate}

\subsection{10. Связи с другими
билетами}\label{ux441ux432ux44fux437ux438-ux441-ux434ux440ux443ux433ux438ux43cux438-ux431ux438ux43bux435ux442ux430ux43cux438-6}

\textbf{Билет 03.} Интеграл Лебега строится по мере; мера Лебега
получается как продолжение меры с полуалгебры.

\textbf{Билет 04.} Для интеграла по \(\mathbb R^d\) важны борелевские и
лебегово измеримые множества.

\textbf{Билет 05.} Интегрировать можно измеримые функции; измеримость -
входное условие конструкции.

\textbf{Билет 06.} Вероятности через плотности записываются как
интегралы по мере Лебега.

\textbf{Билет 08.} Характеристические функции и интегрируемость
ограниченных измеримых функций опираются на интеграл Лебега.

\textbf{Билет 09.} Теоремы Фубини и независимость используют интегралы
по произведению мер.

\textbf{Билет 10.} Теорема Радона-Никодима и условное ожидание
формулируются через интегралы.

\textbf{Билет 11.} Замена переменной, моменты и дисперсии выражаются
через интеграл по вероятностному пространству и по распределению.

\subsection{11. Минимум на
удовлетворительно}\label{ux43cux438ux43dux438ux43cux443ux43c-ux43dux430-ux443ux434ux43eux432ux43bux435ux442ux432ux43eux440ux438ux442ux435ux43bux44cux43dux43e-6}

Нужно уметь сказать:

\begin{enumerate}
\def\labelenumi{\arabic{enumi}.}
\tightlist
\item
  Простая функция:
\end{enumerate}

\[
s=\sum_{k=1}^n a_k\mathbf 1_{A_k}.
\]

\begin{enumerate}
\def\labelenumi{\arabic{enumi}.}
\setcounter{enumi}{1}
\tightlist
\item
  Интеграл простой функции:
\end{enumerate}

\[
\int s\,d\mu=\sum_{k=1}^n a_k\mu(A_k).
\]

\begin{enumerate}
\def\labelenumi{\arabic{enumi}.}
\setcounter{enumi}{2}
\tightlist
\item
  Интеграл неотрицательной функции:
\end{enumerate}

\[
\int f\,d\mu=\sup_{0\le s\le f}\int s\,d\mu.
\]

\begin{enumerate}
\def\labelenumi{\arabic{enumi}.}
\setcounter{enumi}{3}
\tightlist
\item
  Теорема Беппо Леви:
\end{enumerate}

\[
0\le f_n\uparrow f
\Rightarrow
\int f_n\,d\mu\uparrow\int f\,d\mu.
\]

\begin{enumerate}
\def\labelenumi{\arabic{enumi}.}
\setcounter{enumi}{4}
\tightlist
\item
  Теорема Лебега:
\end{enumerate}

\[
f_n\to f,\quad |f_n|\le g\in L^1
\Rightarrow
\int f_n\,d\mu\to\int f\,d\mu.
\]

\subsection{12. Минимум на
хорошо}\label{ux43cux438ux43dux438ux43cux443ux43c-ux43dux430-ux445ux43eux440ux43eux448ux43e-6}

Кроме минимума на удовлетворительно, нужно:

\begin{enumerate}
\def\labelenumi{\arabic{enumi}.}
\tightlist
\item
  Объяснить, почему интеграл простой функции не зависит от
  представления.
\item
  Назвать монотонность и линейность интеграла.
\item
  Рассказать про положительную и отрицательную части функции.
\item
  Объяснить, что такое \(L^1\) и абсолютная интегрируемость.
\item
  Привести пример, где без мажоранты сходимость интегралов ломается.
\end{enumerate}

\subsection{13. Что добавить для
отлично}\label{ux447ux442ux43e-ux434ux43eux431ux430ux432ux438ux442ux44c-ux434ux43bux44f-ux43eux442ux43bux438ux447ux43dux43e-6}

Для сильного ответа стоит добавить:

\begin{enumerate}
\def\labelenumi{\arabic{enumi}.}
\tightlist
\item
  Схему аппроксимации неотрицательной измеримой функции простыми
  функциями снизу.
\item
  Идею доказательства теоремы Беппо Леви через простую функцию
  \(s\le f\) и множества \(\{f_n\ge\alpha s\}\).
\item
  Идею мажорантной теоремы Лебега: мажоранта не дает массе уходить в
  пики.
\item
  Различие между конечным интегралом, бесконечным интегралом
  неотрицательной функции и интегрируемостью в \(L^1\).
\item
  Связь с математическим ожиданием:
\end{enumerate}

\[
\mathbb E X=\int X\,dP.
\]

\subsection{14. Устная версия ответа на 5
минут}\label{ux443ux441ux442ux43dux430ux44f-ux432ux435ux440ux441ux438ux44f-ux43eux442ux432ux435ux442ux430-ux43dux430-5-ux43cux438ux43dux443ux442-6}

Интеграл Лебега строится по неотрицательной мере
\((\Omega,\mathcal F,\mu)\). Начинают с индикатора:

\[
\int\mathbf 1_A\,d\mu=\mu(A).
\]

Потом берут простые функции. Простая функция - это измеримая функция,
принимающая конечное число значений. Если

\[
s=\sum_{k=1}^n a_k\mathbf 1_{A_k},
\]

где \(a_k\ge0\), \(A_k\in\mathcal F\) и \(A_k\) попарно не пересекаются,
то

\[
\int s\,d\mu=\sum_{k=1}^n a_k\mu(A_k).
\]

Корректность следует из перехода к общему измельчению двух разбиений.

Для произвольной неотрицательной измеримой функции \(f\) интеграл
определяется как супремум по простым функциям снизу:

\[
\int f\,d\mu
=
\sup_{0\le s\le f,\ s\text{ простая}}
\int s\,d\mu.
\]

Значение может быть бесконечным.

Для знакопеременной функции вводят

\[
f^+=\max(f,0),\qquad f^-=\max(-f,0).
\]

Если

\[
\int |f|\,d\mu<\infty,
\]

то \(f\in L^1\) и

\[
\int f\,d\mu=\int f^+\,d\mu-\int f^-\,d\mu.
\]

Главные теоремы о предельном переходе идут в программе без
доказательства. Теорема Беппо Леви: если

\[
0\le f_n\uparrow f,
\]

то

\[
\int f_n\,d\mu\uparrow\int f\,d\mu.
\]

Мажорантная теорема Лебега: если \(f_n\to f\) почти всюду и есть
\(g\in L^1\) такая, что

\[
|f_n|\le g,
\]

то

\[
\int f_n\,d\mu\to\int f\,d\mu.
\]

Важно: почти всюду сходимости самой по себе недостаточно. Например,
\(f_n=n\mathbf 1_{(0,1/n)}\) сходится к нулю почти всюду на \([0,1]\),
но интегралы равны \(1\).

\subsection{15. Если осталось 2 минуты перед
экзаменом}\label{ux435ux441ux43bux438-ux43eux441ux442ux430ux43bux43eux441ux44c-2-ux43cux438ux43dux443ux442ux44b-ux43fux435ux440ux435ux434-ux44dux43aux437ux430ux43cux435ux43dux43eux43c-6}

Запомнить:

\[
s=\sum_{k=1}^n a_k\mathbf 1_{A_k},
\qquad
\int s\,d\mu=\sum_{k=1}^n a_k\mu(A_k).
\]

Для \(f\ge0\):

\[
\int f\,d\mu
=
\sup_{0\le s\le f,\ s\text{ простая}}
\int s\,d\mu.
\]

Беппо Леви:

\[
0\le f_n\uparrow f
\Rightarrow
\int f_n\,d\mu\uparrow\int f\,d\mu.
\]

Лебег:

\[
f_n\to f,\quad |f_n|\le g,\quad g\in L^1
\Rightarrow
\int f_n\,d\mu\to\int f\,d\mu.
\]

Главная ловушка: теорема Лебега требует интегрируемую мажоранту, а
просто почти всюду сходимости недостаточно.

\newpage

\section{Билет 08. Интегрируемость ограниченных измеримых функций.
Характеристические функции конечномерных
распределений}\label{ux431ux438ux43bux435ux442-08.-ux438ux43dux442ux435ux433ux440ux438ux440ux443ux435ux43cux43eux441ux442ux44c-ux43eux433ux440ux430ux43dux438ux447ux435ux43dux43dux44bux445-ux438ux437ux43cux435ux440ux438ux43cux44bux445-ux444ux443ux43dux43aux446ux438ux439.-ux445ux430ux440ux430ux43aux442ux435ux440ux438ux441ux442ux438ux447ux435ux441ux43aux438ux435-ux444ux443ux43dux43aux446ux438ux438-ux43aux43eux43dux435ux447ux43dux43eux43cux435ux440ux43dux44bux445-ux440ux430ux441ux43fux440ux435ux434ux435ux43bux435ux43dux438ux439}

Источники: Ширяев 1, гл. II, §§6,12; лекции ДСП.

\subsection{1. Короткий
ответ}\label{ux43aux43eux440ux43eux442ux43aux438ux439-ux43eux442ux432ux435ux442-7}

Если \((\Omega,\mathcal F,\mu)\) - пространство конечной меры, то всякая
ограниченная измеримая функция интегрируема:

\[
|f|\le M,\qquad \mu(\Omega)<\infty
\quad\Longrightarrow\quad
\int_\Omega |f|\,d\mu\le M\mu(\Omega)<\infty.
\]

В вероятностном пространстве это означает: всякая ограниченная случайная
величина имеет конечное математическое ожидание.

Характеристическая функция распределения - это его преобразование Фурье.
Для случайного вектора \(X=(X_1,\ldots,X_n)\):

\[
\varphi_X(t)
=\mathbb E e^{i\langle t,X\rangle}
=\int_{\mathbb R^n} e^{i\langle t,x\rangle}\,P_X(dx),
\qquad t\in\mathbb R^n.
\]

Она всегда существует, потому что

\[
\left|e^{i\langle t,X\rangle}\right|=1.
\]

Для случайного процесса конечномерные характеристические функции - это
характеристические функции векторов

\[
(X_{t_1},\ldots,X_{t_n}).
\]

Они однозначно задают соответствующие конечномерные распределения.

\subsection{2. Полный
ответ}\label{ux43fux43eux43bux43dux44bux439-ux43eux442ux432ux435ux442-7}

\subsubsection{2.1. Введение: основные идеи
билета}\label{ux432ux432ux435ux434ux435ux43dux438ux435-ux43eux441ux43dux43eux432ux43dux44bux435-ux438ux434ux435ux438-ux431ux438ux43bux435ux442ux430}

В билете соединяются две идеи. Первая относится к интегралу Лебега: на
пространстве конечной меры ограниченность и измеримость автоматически
дают интегрируемость. Вторая относится к распределениям:
характеристическая функция кодирует распределение как преобразование
Фурье и поэтому удобна для работы с суммами, предельными теоремами и
конечномерными распределениями случайных процессов.

\subsubsection{2.2. Интегрируемость ограниченных измеримых
функций}\label{ux438ux43dux442ux435ux433ux440ux438ux440ux443ux435ux43cux43eux441ux442ux44c-ux43eux433ux440ux430ux43dux438ux447ux435ux43dux43dux44bux445-ux438ux437ux43cux435ux440ux438ux43cux44bux445-ux444ux443ux43dux43aux446ux438ux439}

Пусть задано пространство с мерой

\[
(\Omega,\mathcal F,\mu).
\]

Функция \(f:\Omega\to\mathbb R\) или \(f:\Omega\to\mathbb C\) называется
интегрируемой, если

\[
\int_\Omega |f|\,d\mu<\infty.
\]

Для комплексной функции \(f=u+iv\) условие \(\int |f|\,d\mu<\infty\)
эквивалентно интегрируемости действительной и мнимой частей:
\(u,v\in L^1(\mu)\).

Если \(f\) измерима и ограничена, то существует \(M<\infty\) такое, что

\[
|f(\omega)|\le M
\]

для всех \(\omega\) или хотя бы почти всюду. Тогда

\[
\int_\Omega |f|\,d\mu
\le
\int_\Omega M\,d\mu
=M\mu(\Omega).
\]

Если \(\mu(\Omega)<\infty\), правая часть конечна, значит
\(f\in L^1(\mu)\). В вероятности \(\mu=P\) и \(P(\Omega)=1\), поэтому
оценка принимает особенно простой вид:

\[
\mathbb E|X|\le M.
\]

Главное ограничение этой теоремы: мера должна быть конечной. На
бесконечной мере утверждение неверно. Например, функция \(f(x)\equiv1\)
ограничена и измерима на \((\mathbb R,\mathcal B(\mathbb R),\lambda)\),
но

\[
\int_\mathbb R 1\,d\lambda=\infty.
\]

\subsubsection{2.3. Определение характеристической функции случайной
величины и
вектора}\label{ux43eux43fux440ux435ux434ux435ux43bux435ux43dux438ux435-ux445ux430ux440ux430ux43aux442ux435ux440ux438ux441ux442ux438ux447ux435ux441ux43aux43eux439-ux444ux443ux43dux43aux446ux438ux438-ux441ux43bux443ux447ux430ux439ux43dux43eux439-ux432ux435ux43bux438ux447ux438ux43dux44b-ux438-ux432ux435ux43aux442ux43eux440ux430}

Теперь перейдем к характеристическим функциям. Пусть \(X\) - случайная
величина с распределением \(P_X\) на \(\mathbb R\). Ее
характеристическая функция:

\[
\varphi_X(t)=\mathbb E e^{itX}
=\int_\mathbb R e^{itx}\,P_X(dx),
\qquad t\in\mathbb R.
\]

Если распределение задается функцией распределения \(F_X\), то пишут
также

\[
\varphi_X(t)=\int_\mathbb R e^{itx}\,dF_X(x).
\]

Если есть плотность \(p_X\), то

\[
\varphi_X(t)=\int_\mathbb R e^{itx}p_X(x)\,dx.
\]

Для случайного вектора \(X=(X_1,\ldots,X_n)\) характеристическая функция
задается формулой

\[
\varphi_X(t)
=\mathbb E e^{i\langle t,X\rangle}
=\int_{\mathbb R^n} e^{i\langle t,x\rangle}\,P_X(dx),
\qquad t=(t_1,\ldots,t_n)\in\mathbb R^n,
\]

где

\[
\langle t,x\rangle=t_1x_1+\cdots+t_nx_n.
\]

Эта величина всегда корректно определена: случайная величина
\(e^{i\langle t,X\rangle}\) комплекснозначна, измерима и ограничена по
модулю единицей. Поэтому она интегрируема на вероятностном пространстве.
Это важное преимущество характеристических функций перед производящими
функциями моментов: моментная производящая функция \(\mathbb E e^{tX}\)
может не существовать, а характеристическая функция существует всегда.

\subsubsection{2.4. Характеристические функции конечномерных
распределений
процесса}\label{ux445ux430ux440ux430ux43aux442ux435ux440ux438ux441ux442ux438ux447ux435ux441ux43aux438ux435-ux444ux443ux43dux43aux446ux438ux438-ux43aux43eux43dux435ux447ux43dux43eux43cux435ux440ux43dux44bux445-ux440ux430ux441ux43fux440ux435ux434ux435ux43bux435ux43dux438ux439-ux43fux440ux43eux446ux435ux441ux441ux430}

Для случайного процесса \((X_s)_{s\in T}\) конечномерное распределение
задается выбором конечного набора моментов времени

\[
s_1,\ldots,s_n\in T
\]

и законом случайного вектора

\[
(X_{s_1},\ldots,X_{s_n}).
\]

Его характеристическая функция:

\[
\varphi_{s_1,\ldots,s_n}(u_1,\ldots,u_n)
=
\mathbb E\exp\left(i\sum_{k=1}^n u_kX_{s_k}\right).
\]

Именно такие функции называются характеристическими функциями
конечномерных распределений процесса. Если знать их для всех конечных
наборов индексов, то знать все конечномерные распределения. В дальнейшем
это нужно для задания случайных процессов, гауссовских
последовательностей, стационарности и предельных теорем.

\subsubsection{2.5. Основные свойства и
непрерывность}\label{ux43eux441ux43dux43eux432ux43dux44bux435-ux441ux432ux43eux439ux441ux442ux432ux430-ux438-ux43dux435ux43fux440ux435ux440ux44bux432ux43dux43eux441ux442ux44c}

Основные свойства характеристической функции следуют из простых свойств
комплексной экспоненты и интеграла Лебега.

Во-первых,

\[
\varphi_X(0)=\mathbb E1=1.
\]

Во-вторых,

\[
|\varphi_X(t)|
=|\mathbb E e^{i\langle t,X\rangle}|
\le \mathbb E |e^{i\langle t,X\rangle}|
=1.
\]

В-третьих,

\[
\varphi_X(-t)=\overline{\varphi_X(t)}.
\]

В частности, для одномерной случайной величины с симметричным
относительно нуля распределением характеристическая функция
действительнозначна:

\[
\varphi_X(t)=\mathbb E\cos(tX).
\]

В-четвертых, характеристическая функция равномерно непрерывна. Для
одномерного случая:

\[
|\varphi(t+h)-\varphi(t)|
=
\left|\mathbb E\left(e^{itX}(e^{ihX}-1)\right)\right|
\le
\mathbb E|e^{ihX}-1|.
\]

При \(h\to0\) подынтегральное выражение стремится к нулю и ограничено
числом \(2\), поэтому по теореме Лебега о доминированной сходимости
правая часть стремится к нулю. Оценка не зависит от \(t\), значит
непрерывность равномерная.

\subsubsection{2.6. Положительная определенность и теорема
Бохнера}\label{ux43fux43eux43bux43eux436ux438ux442ux435ux43bux44cux43dux430ux44f-ux43eux43fux440ux435ux434ux435ux43bux435ux43dux43dux43eux441ux442ux44c-ux438-ux442ux435ux43eux440ux435ux43cux430-ux431ux43eux445ux43dux435ux440ux430}

В-пятых, характеристическая функция положительно определена:

\[
\sum_{j,k=1}^m c_j\overline{c_k}\,
\varphi_X(t_j-t_k)\ge0
\]

для любых \(t_1,\ldots,t_m\in\mathbb R^n\) и
\(c_1,\ldots,c_m\in\mathbb C\). Действительно,

\[
\sum_{j,k=1}^m c_j\overline{c_k}\,
\varphi_X(t_j-t_k)
=
\mathbb E\left|\sum_{j=1}^m c_je^{i\langle t_j,X\rangle}\right|^2
\ge0.
\]

Это свойство важно в обратную сторону: по теореме Бохнера непрерывная в
нуле положительно определенная функция \(\varphi\) с \(\varphi(0)=1\)
является характеристической функцией некоторого вероятностного
распределения. Для экзамена обычно достаточно знать необходимость этих
условий и смысл положительной определенности.

\subsubsection{2.7. Однозначность распределения и формула
обращения}\label{ux43eux434ux43dux43eux437ux43dux430ux447ux43dux43eux441ux442ux44c-ux440ux430ux441ux43fux440ux435ux434ux435ux43bux435ux43dux438ux44f-ux438-ux444ux43eux440ux43cux443ux43bux430-ux43eux431ux440ux430ux449ux435ux43dux438ux44f}

В-шестых, характеристическая функция однозначно определяет
распределение. Если две случайные величины или два случайных вектора
имеют одну и ту же характеристическую функцию, то их распределения
совпадают. В одномерном случае это можно усилить формулой обращения:
если \(a<b\) - точки непрерывности функции распределения \(F\), то

\[
F(b)-F(a)
=
\lim_{T\to\infty}
\frac1{2\pi}
\int_{-T}^{T}
\frac{e^{-ita}-e^{-itb}}{it}\,\varphi(t)\,dt.
\]

Если дополнительно \(\varphi\in L^1(\mathbb R)\), то распределение имеет
непрерывную плотность

\[
f(x)=\frac1{2\pi}\int_\mathbb R e^{-itx}\varphi(t)\,dt.
\]

Смысл: характеристическая функция - не просто вспомогательная величина,
а полная запись закона распределения.

\subsubsection{2.8. Характеристические функции независимых величин и
преобразований}\label{ux445ux430ux440ux430ux43aux442ux435ux440ux438ux441ux442ux438ux447ux435ux441ux43aux438ux435-ux444ux443ux43dux43aux446ux438ux438-ux43dux435ux437ux430ux432ux438ux441ux438ux43cux44bux445-ux432ux435ux43bux438ux447ux438ux43d-ux438-ux43fux440ux435ux43eux431ux440ux430ux437ux43eux432ux430ux43dux438ux439}

В-седьмых, характеристические функции особенно удобны для сумм
независимых величин. Если \(X\) и \(Y\) независимы, то

\[
\varphi_{X+Y}(t)
=\mathbb E e^{it(X+Y)}
=\mathbb E(e^{itX}e^{itY})
=\mathbb E e^{itX}\,\mathbb E e^{itY}
=\varphi_X(t)\varphi_Y(t).
\]

Для независимых \(X_1,\ldots,X_n\):

\[
\varphi_{X_1+\cdots+X_n}(t)
=
\prod_{k=1}^n\varphi_{X_k}(t).
\]

Именно поэтому характеристические функции используются в центральной
предельной теореме: свертка распределений превращается в произведение
функций.

Для векторов действует аналогичная формула. Если \(X\) и \(Y\)
независимые случайные векторы в \(\mathbb R^n\), то

\[
\varphi_{X+Y}(t)=\varphi_X(t)\varphi_Y(t).
\]

А если компоненты \(X_1,\ldots,X_n\) независимы, то

\[
\varphi_{(X_1,\ldots,X_n)}(t_1,\ldots,t_n)
=
\prod_{k=1}^n\varphi_{X_k}(t_k).
\]

Обратное тоже верно в стандартной форме: если совместная
характеристическая функция вектора факторизуется в произведение
одномерных характеристических функций координат, то координаты
независимы, потому что характеристическая функция однозначно задает
совместное распределение.

Еще одно полезное свойство - поведение при линейных преобразованиях.
Если \(Y=AX+b\), где \(A\) - матрица, \(b\) - вектор, то

\[
\varphi_Y(t)
=e^{i\langle t,b\rangle}\varphi_X(A^Tt).
\]

В одномерном случае для \(Y=aX+b\):

\[
\varphi_Y(t)=e^{itb}\varphi_X(at).
\]

Для конечномерных распределений процесса это означает, что линейные
комбинации координат также удобно описываются через одну и ту же
характеристическую функцию:

\[
\mathbb E e^{i(u_1X_{s_1}+\cdots+u_nX_{s_n})}.
\]

\subsubsection{2.9.
Заключение}\label{ux437ux430ux43aux43bux44eux447ux435ux43dux438ux435}

Итак, в этом билете главная логика такая: ограниченные измеримые функции
на конечной мере можно интегрировать; комплексная экспонента от
случайного вектора всегда ограничена; поэтому характеристическая функция
всегда существует и является надежным способом кодировать конечномерные
распределения.

\subsection{3. Основные
определения}\label{ux43eux441ux43dux43eux432ux43dux44bux435-ux43eux43fux440ux435ux434ux435ux43bux435ux43dux438ux44f-7}

\subsubsection{Интегрируемая
функция}\label{ux438ux43dux442ux435ux433ux440ux438ux440ux443ux435ux43cux430ux44f-ux444ux443ux43dux43aux446ux438ux44f-1}

Измеримая функция \(f\) называется интегрируемой относительно меры
\(\mu\), если

\[
\int |f|\,d\mu<\infty.
\]

В этом случае пишут \(f\in L^1(\mu)\).

\subsubsection{Ограниченная
функция}\label{ux43eux433ux440ux430ux43dux438ux447ux435ux43dux43dux430ux44f-ux444ux443ux43dux43aux446ux438ux44f}

Функция \(f\) ограничена, если существует \(M<\infty\) такое, что

\[
|f(\omega)|\le M
\]

для всех \(\omega\) или почти всюду. Для интеграла Лебега достаточно
ограничения почти всюду.

\subsubsection{Конечная
мера}\label{ux43aux43eux43dux435ux447ux43dux430ux44f-ux43cux435ux440ux430}

Мера \(\mu\) называется конечной, если

\[
\mu(\Omega)<\infty.
\]

Вероятностная мера всегда конечна, потому что \(P(\Omega)=1\).

\subsubsection{Комплекснозначная случайная
величина}\label{ux43aux43eux43cux43fux43bux435ux43aux441ux43dux43eux437ux43dux430ux447ux43dux430ux44f-ux441ux43bux443ux447ux430ux439ux43dux430ux44f-ux432ux435ux43bux438ux447ux438ux43dux430}

Случайная величина \(Z=U+iV\), где \(U\) и \(V\) - действительные
случайные величины. Ее математическое ожидание задается как

\[
\mathbb EZ=\mathbb EU+i\mathbb EV,
\]

если соответствующие ожидания существуют. Если \(\mathbb E|Z|<\infty\),
то ожидание существует.

\subsubsection{Распределение случайного
вектора}\label{ux440ux430ux441ux43fux440ux435ux434ux435ux43bux435ux43dux438ux435-ux441ux43bux443ux447ux430ux439ux43dux43eux433ux43e-ux432ux435ux43aux442ux43eux440ux430-1}

Для \(X:\Omega\to\mathbb R^n\) распределение \(P_X\) задается формулой

\[
P_X(B)=P\{X\in B\},
\qquad B\in\mathcal B(\mathbb R^n).
\]

\subsubsection{\texorpdfstring{Характеристическая функция распределения
на
\(\mathbb R^n\)}{Характеристическая функция распределения на \textbackslash mathbb R\^{}n}}\label{ux445ux430ux440ux430ux43aux442ux435ux440ux438ux441ux442ux438ux447ux435ux441ux43aux430ux44f-ux444ux443ux43dux43aux446ux438ux44f-ux440ux430ux441ux43fux440ux435ux434ux435ux43bux435ux43dux438ux44f-ux43dux430-mathbb-rn}

Если \(\mu\) - вероятностная мера на \(\mathbb R^n\), то

\[
\varphi_\mu(t)=\int_{\mathbb R^n} e^{i\langle t,x\rangle}\,\mu(dx),
\qquad t\in\mathbb R^n.
\]

\subsubsection{Характеристическая функция случайного
вектора}\label{ux445ux430ux440ux430ux43aux442ux435ux440ux438ux441ux442ux438ux447ux435ux441ux43aux430ux44f-ux444ux443ux43dux43aux446ux438ux44f-ux441ux43bux443ux447ux430ux439ux43dux43eux433ux43e-ux432ux435ux43aux442ux43eux440ux430}

Для \(X=(X_1,\ldots,X_n)\):

\[
\varphi_X(t)=\mathbb E e^{i\langle t,X\rangle}.
\]

\subsubsection{Характеристическая функция случайной
величины}\label{ux445ux430ux440ux430ux43aux442ux435ux440ux438ux441ux442ux438ux447ux435ux441ux43aux430ux44f-ux444ux443ux43dux43aux446ux438ux44f-ux441ux43bux443ux447ux430ux439ux43dux43eux439-ux432ux435ux43bux438ux447ux438ux43dux44b}

Частный случай при \(n=1\):

\[
\varphi_X(t)=\mathbb E e^{itX}.
\]

\subsubsection{Конечномерное распределение
процесса}\label{ux43aux43eux43dux435ux447ux43dux43eux43cux435ux440ux43dux43eux435-ux440ux430ux441ux43fux440ux435ux434ux435ux43bux435ux43dux438ux435-ux43fux440ux43eux446ux435ux441ux441ux430-1}

Для процесса \((X_s)_{s\in T}\) и индексов \(s_1,\ldots,s_n\) это
распределение вектора

\[
(X_{s_1},\ldots,X_{s_n}).
\]

\subsubsection{Характеристическая функция конечномерного
распределения}\label{ux445ux430ux440ux430ux43aux442ux435ux440ux438ux441ux442ux438ux447ux435ux441ux43aux430ux44f-ux444ux443ux43dux43aux446ux438ux44f-ux43aux43eux43dux435ux447ux43dux43eux43cux435ux440ux43dux43eux433ux43e-ux440ux430ux441ux43fux440ux435ux434ux435ux43bux435ux43dux438ux44f}

\[
\varphi_{s_1,\ldots,s_n}(u_1,\ldots,u_n)
=
\mathbb E\exp\left(i\sum_{k=1}^n u_kX_{s_k}\right).
\]

\subsubsection{Положительно определенная
функция}\label{ux43fux43eux43bux43eux436ux438ux442ux435ux43bux44cux43dux43e-ux43eux43fux440ux435ux434ux435ux43bux435ux43dux43dux430ux44f-ux444ux443ux43dux43aux446ux438ux44f}

Функция \(\varphi:\mathbb R^n\to\mathbb C\) положительно определена,
если для любых \(t_1,\ldots,t_m\in\mathbb R^n\) и
\(c_1,\ldots,c_m\in\mathbb C\)

\[
\sum_{j,k=1}^m c_j\overline{c_k}\varphi(t_j-t_k)\ge0.
\]

\subsection{4. Основные теоремы и
свойства}\label{ux43eux441ux43dux43eux432ux43dux44bux435-ux442ux435ux43eux440ux435ux43cux44b-ux438-ux441ux432ux43eux439ux441ux442ux432ux430-7}

\subsubsection{Теорема 1. Ограниченная измеримая функция интегрируема на
конечной
мере}\label{ux442ux435ux43eux440ux435ux43cux430-1.-ux43eux433ux440ux430ux43dux438ux447ux435ux43dux43dux430ux44f-ux438ux437ux43cux435ux440ux438ux43cux430ux44f-ux444ux443ux43dux43aux446ux438ux44f-ux438ux43dux442ux435ux433ux440ux438ux440ux443ux435ux43cux430-ux43dux430-ux43aux43eux43dux435ux447ux43dux43eux439-ux43cux435ux440ux435}

Пусть \((\Omega,\mathcal F,\mu)\) - пространство конечной меры, \(f\)
измерима и \(|f|\le M\) почти всюду. Тогда \(f\in L^1(\mu)\) и

\[
\int |f|\,d\mu\le M\mu(\Omega).
\]

\subsubsection{Доказательство}\label{ux434ux43eux43aux430ux437ux430ux442ux435ux43bux44cux441ux442ux432ux43e-29}

Так как \(|f|\le M\) почти всюду, по монотонности интеграла

\[
\int |f|\,d\mu\le\int M\,d\mu=M\mu(\Omega)<\infty.
\]

Значит \(f\) интегрируема.

\subsubsection{Следствие для
вероятности}\label{ux441ux43bux435ux434ux441ux442ux432ux438ux435-ux434ux43bux44f-ux432ux435ux440ux43eux44fux442ux43dux43eux441ux442ux438}

Если \(X\) - ограниченная случайная величина, то

\[
\mathbb E|X|\le M.
\]

Поэтому \(X\) имеет математическое ожидание.

\subsubsection{Следствие для характеристической
функции}\label{ux441ux43bux435ux434ux441ux442ux432ux438ux435-ux434ux43bux44f-ux445ux430ux440ux430ux43aux442ux435ux440ux438ux441ux442ux438ux447ux435ux441ux43aux43eux439-ux444ux443ux43dux43aux446ux438ux438}

Для любого случайного вектора \(X\) и любого \(t\in\mathbb R^n\)
случайная величина

\[
e^{i\langle t,X\rangle}
\]

ограничена по модулю единицей, значит интегрируема. Поэтому
\(\varphi_X(t)\) существует для всех \(t\).

\subsubsection{Свойство 1.
Нормировка}\label{ux441ux432ux43eux439ux441ux442ux432ux43e-1.-ux43dux43eux440ux43cux438ux440ux43eux432ux43aux430}

\[
\varphi_X(0)=1.
\]

\subsubsection{Свойство 2.
Ограниченность}\label{ux441ux432ux43eux439ux441ux442ux432ux43e-2.-ux43eux433ux440ux430ux43dux438ux447ux435ux43dux43dux43eux441ux442ux44c}

\[
|\varphi_X(t)|\le1.
\]

\subsubsection{Свойство 3. Сопряженная
симметрия}\label{ux441ux432ux43eux439ux441ux442ux432ux43e-3.-ux441ux43eux43fux440ux44fux436ux435ux43dux43dux430ux44f-ux441ux438ux43cux43cux435ux442ux440ux438ux44f}

\[
\varphi_X(-t)=\overline{\varphi_X(t)}.
\]

\subsubsection{Свойство 4. Равномерная
непрерывность}\label{ux441ux432ux43eux439ux441ux442ux432ux43e-4.-ux440ux430ux432ux43dux43eux43cux435ux440ux43dux430ux44f-ux43dux435ux43fux440ux435ux440ux44bux432ux43dux43eux441ux442ux44c}

Характеристическая функция равномерно непрерывна на \(\mathbb R^n\).

\subsubsection{Идея
доказательства}\label{ux438ux434ux435ux44f-ux434ux43eux43aux430ux437ux430ux442ux435ux43bux44cux441ux442ux432ux430-2}

Для \(h\to0\):

\[
|\varphi(t+h)-\varphi(t)|
\le
\mathbb E|e^{i\langle h,X\rangle}-1|
\to0
\]

по мажорантной теореме Лебега, так как модуль подынтегрального выражения
не больше \(2\).

\subsubsection{Свойство 5. Положительная
определенность}\label{ux441ux432ux43eux439ux441ux442ux432ux43e-5.-ux43fux43eux43bux43eux436ux438ux442ux435ux43bux44cux43dux430ux44f-ux43eux43fux440ux435ux434ux435ux43bux435ux43dux43dux43eux441ux442ux44c}

Характеристическая функция положительно определена:

\[
\sum_{j,k=1}^m c_j\overline{c_k}\varphi(t_j-t_k)\ge0.
\]

\subsubsection{Доказательство}\label{ux434ux43eux43aux430ux437ux430ux442ux435ux43bux44cux441ux442ux432ux43e-30}

\[
\sum_{j,k=1}^m c_j\overline{c_k}\varphi(t_j-t_k)
=
\mathbb E\left|\sum_{j=1}^m c_je^{i\langle t_j,X\rangle}\right|^2
\ge0.
\]

\subsubsection{Свойство 6.
Единственность}\label{ux441ux432ux43eux439ux441ux442ux432ux43e-6.-ux435ux434ux438ux43dux441ux442ux432ux435ux43dux43dux43eux441ux442ux44c}

Если

\[
\varphi_X(t)=\varphi_Y(t)
\qquad \text{для всех }t,
\]

то

\[
P_X=P_Y.
\]

То есть характеристическая функция однозначно определяет распределение.

\subsubsection{Свойство 7. Формула обращения в одномерном
случае}\label{ux441ux432ux43eux439ux441ux442ux432ux43e-7.-ux444ux43eux440ux43cux443ux43bux430-ux43eux431ux440ux430ux449ux435ux43dux438ux44f-ux432-ux43eux434ux43dux43eux43cux435ux440ux43dux43eux43c-ux441ux43bux443ux447ux430ux435}

Если \(a<b\) - точки непрерывности функции распределения \(F\), то

\[
F(b)-F(a)
=
\lim_{T\to\infty}
\frac1{2\pi}
\int_{-T}^{T}
\frac{e^{-ita}-e^{-itb}}{it}\varphi(t)\,dt.
\]

Если \(\varphi\in L^1(\mathbb R)\), то есть плотность

\[
f(x)=\frac1{2\pi}\int_\mathbb R e^{-itx}\varphi(t)\,dt.
\]

\subsubsection{Свойство 8. Линейное
преобразование}\label{ux441ux432ux43eux439ux441ux442ux432ux43e-8.-ux43bux438ux43dux435ux439ux43dux43eux435-ux43fux440ux435ux43eux431ux440ux430ux437ux43eux432ux430ux43dux438ux435}

Если \(Y=AX+b\), то

\[
\varphi_Y(t)=e^{i\langle t,b\rangle}\varphi_X(A^Tt).
\]

В частности, если \(Y=aX+b\), то

\[
\varphi_Y(t)=e^{itb}\varphi_X(at).
\]

\subsubsection{Свойство 9. Сумма независимых
величин}\label{ux441ux432ux43eux439ux441ux442ux432ux43e-9.-ux441ux443ux43cux43cux430-ux43dux435ux437ux430ux432ux438ux441ux438ux43cux44bux445-ux432ux435ux43bux438ux447ux438ux43d}

Если \(X\) и \(Y\) независимы, то

\[
\varphi_{X+Y}(t)=\varphi_X(t)\varphi_Y(t).
\]

Для независимых \(X_1,\ldots,X_n\):

\[
\varphi_{X_1+\cdots+X_n}(t)=\prod_{k=1}^n\varphi_{X_k}(t).
\]

\subsubsection{Свойство 10. Независимость координат через
факторизацию}\label{ux441ux432ux43eux439ux441ux442ux432ux43e-10.-ux43dux435ux437ux430ux432ux438ux441ux438ux43cux43eux441ux442ux44c-ux43aux43eux43eux440ux434ux438ux43dux430ux442-ux447ux435ux440ux435ux437-ux444ux430ux43aux442ux43eux440ux438ux437ux430ux446ux438ux44e}

Координаты \(X_1,\ldots,X_n\) независимы тогда и только тогда, когда

\[
\varphi_{(X_1,\ldots,X_n)}(t_1,\ldots,t_n)
=
\prod_{k=1}^n\varphi_{X_k}(t_k).
\]

\subsubsection{Свойство 11. Моменты через
производные}\label{ux441ux432ux43eux439ux441ux442ux432ux43e-11.-ux43cux43eux43cux435ux43dux442ux44b-ux447ux435ux440ux435ux437-ux43fux440ux43eux438ux437ux432ux43eux434ux43dux44bux435}

Если \(\mathbb E|X|^m<\infty\), то характеристическая функция имеет
производные до порядка \(m\), и

\[
\varphi_X^{(k)}(0)=i^k\mathbb E X^k,
\qquad k\le m.
\]

Эта связь подробно разбирается в Билете 13. В этом билете важно не
перепутать: характеристическая функция существует всегда, а моменты и
производные нужного порядка могут не существовать.

\subsubsection{Свойство 12. Согласованность конечномерных
характеристических
функций}\label{ux441ux432ux43eux439ux441ux442ux432ux43e-12.-ux441ux43eux433ux43bux430ux441ux43eux432ux430ux43dux43dux43eux441ux442ux44c-ux43aux43eux43dux435ux447ux43dux43eux43cux435ux440ux43dux44bux445-ux445ux430ux440ux430ux43aux442ux435ux440ux438ux441ux442ux438ux447ux435ux441ux43aux438ux445-ux444ux443ux43dux43aux446ux438ux439}

Если конечномерные распределения процесса согласованы, то их
характеристические функции тоже согласованы. Например,

\[
\varphi_{s_1,\ldots,s_n,s_{n+1}}(u_1,\ldots,u_n,0)
=
\varphi_{s_1,\ldots,s_n}(u_1,\ldots,u_n).
\]

Перестановка индексов \(s_k\) соответствует такой же перестановке
аргументов \(u_k\). Это характеристическая версия согласованности из
Билета 15.

\subsection{5. Примеры}\label{ux43fux440ux438ux43cux435ux440ux44b-7}

\subsubsection{Пример 1. Ограниченная случайная
величина}\label{ux43fux440ux438ux43cux435ux440-1.-ux43eux433ux440ux430ux43dux438ux447ux435ux43dux43dux430ux44f-ux441ux43bux443ux447ux430ux439ux43dux430ux44f-ux432ux435ux43bux438ux447ux438ux43dux430}

Пусть \(X\) принимает значения в отрезке \([-3,3]\). Тогда

\[
|X|\le3
\]

и

\[
\mathbb E|X|\le3.
\]

Значит \(X\) интегрируема, даже если мы не знаем ее точного
распределения.

\subsubsection{Пример 2. Почему нужна конечность
меры}\label{ux43fux440ux438ux43cux435ux440-2.-ux43fux43eux447ux435ux43cux443-ux43dux443ux436ux43dux430-ux43aux43eux43dux435ux447ux43dux43eux441ux442ux44c-ux43cux435ux440ux44b}

На \((\mathbb R,\mathcal B(\mathbb R),\lambda)\) функция

\[
f(x)\equiv1
\]

ограничена и измерима, но

\[
\int_\mathbb R |f(x)|\,dx=\infty.
\]

Значит ограниченность сама по себе не гарантирует интегрируемость на
бесконечной мере.

\subsubsection{Пример 3. Вырожденное
распределение}\label{ux43fux440ux438ux43cux435ux440-3.-ux432ux44bux440ux43eux436ux434ux435ux43dux43dux43eux435-ux440ux430ux441ux43fux440ux435ux434ux435ux43bux435ux43dux438ux435}

Если \(X=c\) почти наверное, то

\[
\varphi_X(t)=\mathbb E e^{itc}=e^{itc}.
\]

Это характеристическая функция меры Дирака \(\delta_c\).

\subsubsection{Пример 4. Бернуллиевская случайная
величина}\label{ux43fux440ux438ux43cux435ux440-4.-ux431ux435ux440ux43dux443ux43bux43bux438ux435ux432ux441ux43aux430ux44f-ux441ux43bux443ux447ux430ux439ux43dux430ux44f-ux432ux435ux43bux438ux447ux438ux43dux430}

Пусть

\[
P\{X=1\}=p,\qquad P\{X=0\}=q,\qquad p+q=1.
\]

Тогда

\[
\varphi_X(t)=\mathbb E e^{itX}=q+pe^{it}.
\]

Если \(X_1,\ldots,X_n\) независимы и имеют такое же распределение, то
для \(S_n=X_1+\cdots+X_n\)

\[
\varphi_{S_n}(t)=(q+pe^{it})^n.
\]

Это характеристическая функция биномиального распределения.

\subsubsection{Пример 5. Дискретный
вектор}\label{ux43fux440ux438ux43cux435ux440-5.-ux434ux438ux441ux43aux440ux435ux442ux43dux44bux439-ux432ux435ux43aux442ux43eux440}

Пусть случайный вектор \(X\) принимает значения
\(x^{(1)},\ldots,x^{(m)}\in\mathbb R^n\) с вероятностями
\(p_1,\ldots,p_m\). Тогда

\[
\varphi_X(t)=\sum_{r=1}^m p_r e^{i\langle t,x^{(r)}\rangle}.
\]

Например, если \(X=(X_1,X_2)\) принимает значения \((0,0),(1,0),(0,1)\)
с вероятностями \(p_0,p_1,p_2\), то

\[
\varphi_X(t_1,t_2)=p_0+p_1e^{it_1}+p_2e^{it_2}.
\]

\subsubsection{Пример 6. Нормальное
распределение}\label{ux43fux440ux438ux43cux435ux440-6.-ux43dux43eux440ux43cux430ux43bux44cux43dux43eux435-ux440ux430ux441ux43fux440ux435ux434ux435ux43bux435ux43dux438ux435}

Если

\[
X\sim N(m,\sigma^2),
\]

то

\[
\varphi_X(t)=\exp\left(imt-\frac{\sigma^2t^2}{2}\right).
\]

Для гауссовского вектора \(X\sim N(m,\Sigma)\):

\[
\varphi_X(t)
=
\exp\left(i\langle t,m\rangle-\frac12 t^T\Sigma t\right).
\]

Эта формула будет основной в Билете 18.

\subsubsection{Пример 7. Независимые
координаты}\label{ux43fux440ux438ux43cux435ux440-7.-ux43dux435ux437ux430ux432ux438ux441ux438ux43cux44bux435-ux43aux43eux43eux440ux434ux438ux43dux430ux442ux44b}

Пусть \(X_1\) и \(X_2\) независимы. Тогда

\[
\varphi_{(X_1,X_2)}(t_1,t_2)
=
\mathbb E e^{i(t_1X_1+t_2X_2)}
=
\varphi_{X_1}(t_1)\varphi_{X_2}(t_2).
\]

Если такая факторизация нарушается, координаты не являются независимыми.

\subsubsection{Пример 8. Характеристическая функция существует без
математического
ожидания}\label{ux43fux440ux438ux43cux435ux440-8.-ux445ux430ux440ux430ux43aux442ux435ux440ux438ux441ux442ux438ux447ux435ux441ux43aux430ux44f-ux444ux443ux43dux43aux446ux438ux44f-ux441ux443ux449ux435ux441ux442ux432ux443ux435ux442-ux431ux435ux437-ux43cux430ux442ux435ux43cux430ux442ux438ux447ux435ux441ux43aux43eux433ux43e-ux43eux436ux438ux434ux430ux43dux438ux44f}

У распределения Коши с плотностью

\[
p(x)=\frac1{\pi(1+x^2)}
\]

не существует \(\mathbb E|X|\), но характеристическая функция существует
и равна

\[
\varphi_X(t)=e^{-|t|}.
\]

Это хороший пример против ошибки ``если нет матожидания, то нет
характеристической функции''.

\subsubsection{Пример 9. Конечномерная характеристическая функция
процесса}\label{ux43fux440ux438ux43cux435ux440-9.-ux43aux43eux43dux435ux447ux43dux43eux43cux435ux440ux43dux430ux44f-ux445ux430ux440ux430ux43aux442ux435ux440ux438ux441ux442ux438ux447ux435ux441ux43aux430ux44f-ux444ux443ux43dux43aux446ux438ux44f-ux43fux440ux43eux446ux435ux441ux441ux430}

Пусть \((X_k)_{k\ge1}\) - случайная последовательность. Для трех
моментов времени \(1,2,5\):

\[
\varphi_{1,2,5}(u_1,u_2,u_3)
=
\mathbb E e^{i(u_1X_1+u_2X_2+u_3X_5)}.
\]

Если \(X_1,X_2,X_5\) независимы, то

\[
\varphi_{1,2,5}(u_1,u_2,u_3)
=
\varphi_{X_1}(u_1)\varphi_{X_2}(u_2)\varphi_{X_5}(u_3).
\]

Если последовательность стационарна в узком смысле, то такие функции не
меняются при общем сдвиге индексов:

\[
\varphi_{1,2,5}=\varphi_{1+h,2+h,5+h}.
\]

\subsection{6. Схемы доказательств /
обоснований}\label{ux441ux445ux435ux43cux44b-ux434ux43eux43aux430ux437ux430ux442ux435ux43bux44cux441ux442ux432-ux43eux431ux43eux441ux43dux43eux432ux430ux43dux438ux439}

\textbf{Почему ограниченная измеримая функция интегрируема на конечной
мере.}

Достаточно сравнить \(|f|\) с константой \(M\):

\[
0\le |f|\le M.
\]

По монотонности интеграла

\[
\int |f|\,d\mu\le \int M\,d\mu=M\mu(\Omega).
\]

Конечность \(\mu(\Omega)\) завершает доказательство.

\textbf{Почему характеристическая функция всегда существует.}

Для любого \(t\):

\[
\left|e^{i\langle t,X\rangle}\right|=1.
\]

Это ограниченная измеримая комплекснозначная случайная величина. На
вероятностном пространстве любая ограниченная случайная величина
интегрируема, значит

\[
\mathbb E e^{i\langle t,X\rangle}
\]

корректно определено.

\textbf{Почему \(|\varphi(t)|\le1\).}

Используется неравенство

\[
|\mathbb EZ|\le\mathbb E|Z|
\]

для интегрируемой комплексной случайной величины \(Z\). При
\(Z=e^{i\langle t,X\rangle}\) получаем

\[
|\varphi(t)|\le \mathbb E1=1.
\]

\textbf{Почему характеристическая функция непрерывна.}

Разность:

\[
\varphi(t+h)-\varphi(t)
=
\mathbb E\left[e^{i\langle t,X\rangle}(e^{i\langle h,X\rangle}-1)\right].
\]

Подынтегральное выражение по модулю не больше \(2\) и стремится к нулю
при \(h\to0\). Поэтому можно применить теорему Лебега о доминированной
сходимости.

\textbf{Почему характеристическая функция суммы независимых величин
равна произведению.}

Если \(X\) и \(Y\) независимы, то \(e^{itX}\) и \(e^{itY}\) тоже
независимы как измеримые функции от независимых случайных величин. Тогда

\[
\mathbb E(e^{itX}e^{itY})
=
\mathbb E e^{itX}\,\mathbb E e^{itY}.
\]

\textbf{Почему характеристическая функция задает распределение.}

Это теорема единственности. Идея доказательства: по \(\varphi(t)\) можно
восстановить интегралы от достаточно хороших функций через обратное
преобразование Фурье. Затем индикаторы интервалов аппроксимируются
такими функциями, и восстанавливаются значения

\[
P\{a<X\le b\}=F(b)-F(a)
\]

в точках непрерывности \(F\). Следовательно, восстанавливается вся
функция распределения.

\textbf{Почему конечномерные характеристические функции задают
конечномерные распределения процесса.}

Для каждого набора \(s_1,\ldots,s_n\) функция

\[
\varphi_{s_1,\ldots,s_n}
\]

является характеристической функцией случайного вектора
\((X_{s_1},\ldots,X_{s_n})\). По теореме единственности она задает его
распределение. Если это известно для всех конечных наборов индексов, то
известны все конечномерные распределения процесса.

\subsection{7. Что написать на
доске}\label{ux447ux442ux43e-ux43dux430ux43fux438ux441ux430ux442ux44c-ux43dux430-ux434ux43eux441ux43aux435-7}

\begin{enumerate}
\def\labelenumi{\arabic{enumi}.}
\tightlist
\item
  Интегрируемость ограниченной функции:
\end{enumerate}

\[
|f|\le M,\quad \mu(\Omega)<\infty
\quad\Rightarrow\quad
\int |f|\,d\mu\le M\mu(\Omega)<\infty.
\]

\begin{enumerate}
\def\labelenumi{\arabic{enumi}.}
\setcounter{enumi}{1}
\tightlist
\item
  Характеристическая функция случайной величины:
\end{enumerate}

\[
\varphi_X(t)=\mathbb E e^{itX}
=\int_\mathbb R e^{itx}\,P_X(dx).
\]

\begin{enumerate}
\def\labelenumi{\arabic{enumi}.}
\setcounter{enumi}{2}
\tightlist
\item
  Характеристическая функция случайного вектора:
\end{enumerate}

\[
\varphi_X(t)=\mathbb E e^{i\langle t,X\rangle}.
\]

\begin{enumerate}
\def\labelenumi{\arabic{enumi}.}
\setcounter{enumi}{3}
\tightlist
\item
  Конечномерная характеристическая функция процесса:
\end{enumerate}

\[
\varphi_{s_1,\ldots,s_n}(u_1,\ldots,u_n)
=
\mathbb E\exp\left(i\sum_{k=1}^n u_kX_{s_k}\right).
\]

\begin{enumerate}
\def\labelenumi{\arabic{enumi}.}
\setcounter{enumi}{4}
\tightlist
\item
  Базовые свойства:
\end{enumerate}

\[
\varphi(0)=1,\qquad |\varphi(t)|\le1,\qquad
\varphi(-t)=\overline{\varphi(t)}.
\]

\begin{enumerate}
\def\labelenumi{\arabic{enumi}.}
\setcounter{enumi}{5}
\tightlist
\item
  Независимые суммы:
\end{enumerate}

\[
X\perp Y
\quad\Rightarrow\quad
\varphi_{X+Y}(t)=\varphi_X(t)\varphi_Y(t).
\]

\begin{enumerate}
\def\labelenumi{\arabic{enumi}.}
\setcounter{enumi}{6}
\tightlist
\item
  Формула обращения:
\end{enumerate}

\[
F(b)-F(a)
=
\lim_{T\to\infty}
\frac1{2\pi}
\int_{-T}^{T}
\frac{e^{-ita}-e^{-itb}}{it}\varphi(t)\,dt.
\]

\subsection{8. Типичные дополнительные
вопросы}\label{ux442ux438ux43fux438ux447ux43dux44bux435-ux434ux43eux43fux43eux43bux43dux438ux442ux435ux43bux44cux43dux44bux435-ux432ux43eux43fux440ux43eux441ux44b-7}

\textbf{Вопрос. Почему ограниченная измеримая функция интегрируема
именно на конечной мере?}

Потому что

\[
\int |f|\,d\mu\le M\mu(\Omega).
\]

Если \(\mu(\Omega)<\infty\), правая часть конечна. Если мера бесконечна,
оценка не дает конечности; пример \(f\equiv1\) на \(\mathbb R\)
показывает, что утверждение вообще неверно.

\textbf{Вопрос. Почему характеристическая функция всегда существует?}

Потому что комплексная экспонента имеет модуль \(1\):

\[
|e^{i\langle t,X\rangle}|=1.
\]

Значит это ограниченная случайная величина, а на вероятностном
пространстве ограниченные случайные величины интегрируемы.

\textbf{Вопрос. Чем характеристическая функция отличается от моментной
производящей функции?}

Характеристическая функция:

\[
\varphi_X(t)=\mathbb E e^{itX}
\]

существует для всех \(t\in\mathbb R\). Моментная производящая функция:

\[
M_X(t)=\mathbb E e^{tX}
\]

может не существовать ни в какой окрестности нуля. Например, у
распределения Коши характеристическая функция есть, а моменты и
\(M_X(t)\) не существуют.

\textbf{Вопрос. Что значит, что характеристическая функция -
преобразование Фурье распределения?}

Это значит, что вместо плотности или меры \(P_X\) рассматривается
интеграл

\[
\int e^{itx}\,P_X(dx).
\]

Если есть плотность \(p\), то это обычное преобразование Фурье:

\[
\varphi(t)=\int e^{itx}p(x)\,dx.
\]

\textbf{Вопрос. Можно ли из характеристической функции восстановить
распределение?}

Да. По теореме единственности характеристическая функция однозначно
задает распределение. В одномерном случае есть формула обращения для
\(F(b)-F(a)\) в точках непрерывности \(F\).

\textbf{Вопрос. Почему для суммы независимых величин характеристические
функции перемножаются?}

Потому что

\[
e^{it(X+Y)}=e^{itX}e^{itY},
\]

а для независимых случайных величин ожидание произведения равно
произведению ожиданий.

\textbf{Вопрос. Верно ли произведение характеристических функций без
независимости?}

Нет. Формула

\[
\varphi_{X+Y}(t)=\varphi_X(t)\varphi_Y(t)
\]

требует независимости. Без независимости распределение суммы зависит от
совместного распределения, а не только от маргинальных законов.

\textbf{Вопрос. Как записать характеристическую функцию конечномерного
распределения процесса?}

Для индексов \(s_1,\ldots,s_n\):

\[
\varphi_{s_1,\ldots,s_n}(u_1,\ldots,u_n)
=
\mathbb E\exp\left(i\sum_{k=1}^n u_kX_{s_k}\right).
\]

Это характеристическая функция вектора \((X_{s_1},\ldots,X_{s_n})\).

\textbf{Вопрос. Что дает положительная определенность характеристической
функции?}

Она является необходимым свойством любой характеристической функции:

\[
\sum_{j,k}c_j\overline{c_k}\varphi(t_j-t_k)\ge0.
\]

Интуитивно это следует из неотрицательности квадрата модуля. В обратную
сторону теорема Бохнера говорит, что непрерывная в нуле положительно
определенная функция с \(\varphi(0)=1\) является характеристической
функцией некоторого распределения.

\textbf{Вопрос. Как связаны характеристические функции и моменты?}

Если существует момент нужного порядка, то производные в нуле выражают
моменты:

\[
\varphi_X^{(k)}(0)=i^k\mathbb E X^k.
\]

Но существование \(\varphi\) само по себе не означает существование
моментов. Подробно это отдельная тема Билета 13.

\textbf{Вопрос. Как проверить независимость координат через
характеристическую функцию?}

Нужно проверить факторизацию:

\[
\varphi_{(X_1,\ldots,X_n)}(t_1,\ldots,t_n)
=
\prod_{k=1}^n\varphi_{X_k}(t_k).
\]

Если она выполнена для всех аргументов, координаты независимы.

\subsection{9. Типичные
ошибки}\label{ux442ux438ux43fux438ux447ux43dux44bux435-ux43eux448ux438ux431ux43aux438-7}

\begin{enumerate}
\def\labelenumi{\arabic{enumi}.}
\item
  Говорить, что ограниченная функция всегда интегрируема, не упоминая
  конечность меры.
\item
  Забывать измеримость: ограниченность без измеримости не дает интеграла
  Лебега как объекта на данном измеримом пространстве.
\item
  Путать интегрируемость \(X\) и существование характеристической
  функции. Для \(\varphi_X\) нужна интегрируемость \(e^{itX}\), а она
  есть всегда, потому что модуль равен \(1\).
\item
  Путать характеристическую функцию с моментной производящей функцией:
\end{enumerate}

\[
\mathbb E e^{itX}\ne \mathbb E e^{tX}.
\]

\begin{enumerate}
\def\labelenumi{\arabic{enumi}.}
\setcounter{enumi}{4}
\item
  Считать, что если нет математического ожидания, то нет
  характеристической функции.
\item
  Применять формулу
\end{enumerate}

\[
\varphi_{X+Y}=\varphi_X\varphi_Y
\]

без независимости.

\begin{enumerate}
\def\labelenumi{\arabic{enumi}.}
\setcounter{enumi}{6}
\tightlist
\item
  Считать, что плотность существует у любого распределения. Формула
\end{enumerate}

\[
\varphi(t)=\int e^{itx}p(x)\,dx
\]

верна только при наличии плотности.

\begin{enumerate}
\def\labelenumi{\arabic{enumi}.}
\setcounter{enumi}{7}
\tightlist
\item
  Забывать в векторном случае скалярное произведение:
\end{enumerate}

\[
e^{i\langle t,X\rangle},
\]

а не произведение \(e^{itX}\) в одномерном смысле.

\begin{enumerate}
\def\labelenumi{\arabic{enumi}.}
\setcounter{enumi}{8}
\item
  Писать формулу обращения без условия о точках непрерывности функции
  распределения.
\item
  Смешивать конечномерные распределения процесса с одномерными. Для
  процесса одномерные законы \(X_t\) сами по себе обычно не задают
  совместную зависимость между разными моментами времени.
\end{enumerate}

\subsection{10. Связи с другими
билетами}\label{ux441ux432ux44fux437ux438-ux441-ux434ux440ux443ux433ux438ux43cux438-ux431ux438ux43bux435ux442ux430ux43cux438-7}

\textbf{Билет 05.} Характеристическая функция определена для случайного
вектора как измеримого отображения. Измеримость нужна, чтобы
\(e^{i\langle t,X\rangle}\) была случайной величиной.

\textbf{Билет 06.} Распределение случайного вектора \(P_X\) - мера, по
которой интегрируется \(e^{i\langle t,x\rangle}\). Конечномерные
распределения процесса в этом билете получают характеристические
функции.

\textbf{Билет 07.} Интегрируемость и теорема Лебега о доминированной
сходимости нужны для существования и непрерывности характеристической
функции.

\textbf{Билет 09.} Независимость дает факторизацию совместного
распределения и произведение характеристических функций.

\textbf{Билет 11.} Формула замены переменной объясняет равенство

\[
\mathbb E e^{i\langle t,X\rangle}
=
\int e^{i\langle t,x\rangle}\,P_X(dx).
\]

\textbf{Билет 13.} Производные характеристической функции в нуле связаны
с математическим ожиданием, вторыми моментами и матрицей корреляции.

\textbf{Билет 15.} Согласованные конечномерные распределения можно
описывать через согласованные конечномерные характеристические функции.

\textbf{Билет 17.} Узкая стационарность означает неизменность всех
конечномерных распределений при сдвиге, а значит неизменность их
характеристических функций.

\textbf{Билет 18.} Гауссовские векторы и последовательности удобно
задаются характеристической функцией

\[
\exp\left(i\langle t,m\rangle-\frac12t^T\Sigma t\right).
\]

\textbf{Билет 20.} Характеристическая функция - Фурье-образ
распределения; это связывает вероятностные меры, дельта-функции и
обобщенные функции.

\textbf{Билет 21.} Метод характеристических функций используется в
доказательстве центральной предельной теоремы и закона больших чисел.

\subsection{11. Минимум на
удовлетворительно}\label{ux43cux438ux43dux438ux43cux443ux43c-ux43dux430-ux443ux434ux43eux432ux43bux435ux442ux432ux43eux440ux438ux442ux435ux43bux44cux43dux43e-7}

Нужно сказать три вещи.

\begin{enumerate}
\def\labelenumi{\arabic{enumi}.}
\tightlist
\item
  Если \(|f|\le M\) и \(\mu(\Omega)<\infty\), то
\end{enumerate}

\[
\int |f|\,d\mu\le M\mu(\Omega)<\infty.
\]

\begin{enumerate}
\def\labelenumi{\arabic{enumi}.}
\setcounter{enumi}{1}
\tightlist
\item
  Характеристическая функция случайной величины:
\end{enumerate}

\[
\varphi_X(t)=\mathbb E e^{itX}.
\]

Для вектора:

\[
\varphi_X(t)=\mathbb E e^{i\langle t,X\rangle}.
\]

\begin{enumerate}
\def\labelenumi{\arabic{enumi}.}
\setcounter{enumi}{2}
\tightlist
\item
  Она всегда существует, потому что \(|e^{i\langle t,X\rangle}|=1\), и
  однозначно задает распределение.
\end{enumerate}

Минимальный пример: для \(X\sim \operatorname{Bernoulli}(p)\)

\[
\varphi_X(t)=q+pe^{it}.
\]

\subsection{12. Минимум на
хорошо}\label{ux43cux438ux43dux438ux43cux443ux43c-ux43dux430-ux445ux43eux440ux43eux448ux43e-7}

Кроме минимума нужно назвать свойства:

\[
\varphi(0)=1,\qquad |\varphi(t)|\le1,\qquad
\varphi(-t)=\overline{\varphi(t)}.
\]

Также нужно объяснить:

\begin{itemize}
\tightlist
\item
  почему \(\varphi\) равномерно непрерывна;
\item
  почему для независимых сумм функции перемножаются;
\item
  как записывается характеристическая функция конечномерного
  распределения процесса;
\item
  что без конечности меры ограниченность не гарантирует интегрируемость.
\end{itemize}

\subsection{13. Что добавить для
отлично}\label{ux447ux442ux43e-ux434ux43eux431ux430ux432ux438ux442ux44c-ux434ux43bux44f-ux43eux442ux43bux438ux447ux43dux43e-7}

Для отличного ответа стоит добавить:

\begin{enumerate}
\def\labelenumi{\arabic{enumi}.}
\item
  Доказательство интегрируемости ограниченной функции через оценку
  \(M\mu(\Omega)\).
\item
  Доказательство существования характеристической функции через
  ограниченность комплексной экспоненты.
\item
  Положительную определенность:
\end{enumerate}

\[
\sum_{j,k} c_j\overline{c_k}\varphi(t_j-t_k)\ge0.
\]

\begin{enumerate}
\def\labelenumi{\arabic{enumi}.}
\setcounter{enumi}{3}
\item
  Теорему единственности и идею формулы обращения.
\item
  Факторизацию совместной характеристической функции как критерий
  независимости координат.
\item
  Согласованность конечномерных характеристических функций:
\end{enumerate}

\[
\varphi_{s_1,\ldots,s_n,s_{n+1}}(u_1,\ldots,u_n,0)
=
\varphi_{s_1,\ldots,s_n}(u_1,\ldots,u_n).
\]

\begin{enumerate}
\def\labelenumi{\arabic{enumi}.}
\setcounter{enumi}{6}
\tightlist
\item
  Предупреждение: характеристическая функция существует всегда, но
  моменты и производные могут не существовать.
\end{enumerate}

\subsection{14. Устная версия ответа на 5
минут}\label{ux443ux441ux442ux43dux430ux44f-ux432ux435ux440ux441ux438ux44f-ux43eux442ux432ux435ux442ux430-ux43dux430-5-ux43cux438ux43dux443ux442-7}

Сначала я сформулирую базовый факт об интегрируемости. Пусть
\((\Omega,\mathcal F,\mu)\) - пространство конечной меры, \(f\) измерима
и ограничена: \(|f|\le M\). Тогда

\[
\int |f|\,d\mu\le\int M\,d\mu=M\mu(\Omega)<\infty.
\]

Значит \(f\) интегрируема. В вероятности \(P(\Omega)=1\), поэтому всякая
ограниченная случайная величина интегрируема.

Теперь характеристические функции. Для случайной величины

\[
\varphi_X(t)=\mathbb E e^{itX}
=\int_\mathbb R e^{itx}\,P_X(dx).
\]

Для случайного вектора \(X\in\mathbb R^n\):

\[
\varphi_X(t)=\mathbb E e^{i\langle t,X\rangle}
=\int_{\mathbb R^n}e^{i\langle t,x\rangle}\,P_X(dx).
\]

Она всегда существует, потому что \(|e^{i\langle t,X\rangle}|=1\). Это
как раз применение первой части билета.

Основные свойства:

\[
\varphi(0)=1,\qquad |\varphi(t)|\le1,\qquad
\varphi(-t)=\overline{\varphi(t)}.
\]

Также \(\varphi\) равномерно непрерывна: разность оценивается через

\[
\mathbb E|e^{i\langle h,X\rangle}-1|,
\]

и это стремится к нулю по теореме Лебега о доминированной сходимости.

Если \(X\) и \(Y\) независимы, то

\[
\varphi_{X+Y}(t)=\mathbb E e^{itX}e^{itY}
=\varphi_X(t)\varphi_Y(t).
\]

Это одно из главных применений характеристических функций.

Для процесса \((X_s)\) конечномерная характеристическая функция имеет
вид

\[
\varphi_{s_1,\ldots,s_n}(u_1,\ldots,u_n)
=
\mathbb E\exp\left(i\sum_{k=1}^n u_kX_{s_k}\right).
\]

Она является характеристической функцией вектора
\((X_{s_1},\ldots,X_{s_n})\) и поэтому однозначно задает соответствующее
конечномерное распределение.

В конце можно привести пример. Для Бернулли:

\[
\varphi_X(t)=q+pe^{it}.
\]

Для нормального распределения:

\[
\varphi_X(t)=\exp\left(imt-\frac{\sigma^2t^2}{2}\right).
\]

Главная ловушка: характеристическая функция существует всегда, но это не
значит, что существуют математическое ожидание, дисперсия или моментная
производящая функция.

\subsection{15. Если осталось 2 минуты перед
экзаменом}\label{ux435ux441ux43bux438-ux43eux441ux442ux430ux43bux43eux441ux44c-2-ux43cux438ux43dux443ux442ux44b-ux43fux435ux440ux435ux434-ux44dux43aux437ux430ux43cux435ux43dux43eux43c-7}

Запомнить нужно следующий каркас.

\begin{enumerate}
\def\labelenumi{\arabic{enumi}.}
\tightlist
\item
  Ограниченная измеримая функция на конечной мере интегрируема:
\end{enumerate}

\[
|f|\le M,\quad \mu(\Omega)<\infty
\Rightarrow
\int |f|\,d\mu\le M\mu(\Omega).
\]

\begin{enumerate}
\def\labelenumi{\arabic{enumi}.}
\setcounter{enumi}{1}
\tightlist
\item
  Характеристическая функция:
\end{enumerate}

\[
\varphi_X(t)=\mathbb E e^{itX},
\qquad
\varphi_X(t)=\mathbb E e^{i\langle t,X\rangle}
\]

для вектора.

\begin{enumerate}
\def\labelenumi{\arabic{enumi}.}
\setcounter{enumi}{2}
\item
  Она всегда существует, потому что модуль экспоненты равен \(1\).
\item
  Главные свойства:
\end{enumerate}

\[
\varphi(0)=1,\qquad |\varphi(t)|\le1,\qquad
\varphi(-t)=\overline{\varphi(t)}.
\]

\begin{enumerate}
\def\labelenumi{\arabic{enumi}.}
\setcounter{enumi}{4}
\tightlist
\item
  Независимые суммы:
\end{enumerate}

\[
\varphi_{X+Y}=\varphi_X\varphi_Y.
\]

\begin{enumerate}
\def\labelenumi{\arabic{enumi}.}
\setcounter{enumi}{5}
\tightlist
\item
  Конечномерное распределение процесса:
\end{enumerate}

\[
\varphi_{s_1,\ldots,s_n}(u)
=
\mathbb E\exp\left(i\sum u_kX_{s_k}\right).
\]

\begin{enumerate}
\def\labelenumi{\arabic{enumi}.}
\setcounter{enumi}{6}
\tightlist
\item
  Главная ошибка: не забыть условие конечности меры для интегрируемости
  и условие независимости для произведения характеристических функций.
\end{enumerate}

\newpage

\section{Билет 09. Теоремы Фубини и независимые случайные
величины}\label{ux431ux438ux43bux435ux442-09.-ux442ux435ux43eux440ux435ux43cux44b-ux444ux443ux431ux438ux43dux438-ux438-ux43dux435ux437ux430ux432ux438ux441ux438ux43cux44bux435-ux441ux43bux443ux447ux430ux439ux43dux44bux435-ux432ux435ux43bux438ux447ux438ux43dux44b}

Источники: Ширяев 1, гл. I, §3; гл. II, §§5-6; лекции ДСП.

\subsection{1. Короткий
ответ}\label{ux43aux43eux440ux43eux442ux43aux438ux439-ux43eux442ux432ux435ux442-8}

В этом билете две связанные идеи.

Первая: если на произведении измеримых пространств задана произведенная
мера, то интеграл по произведению можно считать как повторный интеграл.
Для неотрицательной измеримой функции работает теорема Тонелли:

\[
\int_{\Omega_1\times\Omega_2} f(\omega_1,\omega_2)\,(\mu_1\otimes\mu_2)(d\omega_1,d\omega_2)
=
\int_{\Omega_1}
\left(
\int_{\Omega_2} f(\omega_1,\omega_2)\,\mu_2(d\omega_2)
\right)
\mu_1(d\omega_1),
\]

и аналогично с другим порядком интегрирования. Значения могут быть
бесконечными.

Вторая: независимость означает, что совместное распределение
раскладывается в произведение маргинальных распределений. Для случайных
элементов \(X\) и \(Y\):

\[
X\ \text{и}\ Y\ \text{независимы}
\quad\Longleftrightarrow\quad
P_{(X,Y)}=P_X\otimes P_Y.
\]

Отсюда через Фубини получается ключевая формула:

\[
\mathbb E[g(X)h(Y)]
=
\mathbb E g(X)\,\mathbb E h(Y),
\]

если соответствующие интегралы корректно определены. Для независимых
случайных величин с плотностями совместная плотность равна произведению:

\[
p_{X,Y}(x,y)=p_X(x)p_Y(y)
\]

почти всюду.

\subsection{2. Полный
ответ}\label{ux43fux43eux43bux43dux44bux439-ux43eux442ux432ux435ux442-8}

\subsubsection{2.1. Введение и область
применения}\label{ux432ux432ux435ux434ux435ux43dux438ux435-ux438-ux43eux431ux43bux430ux441ux442ux44c-ux43fux440ux438ux43cux435ux43dux435ux43dux438ux44f}

Теоремы Фубини нужны для того, чтобы корректно переходить от интеграла
по произведению пространств к повторным интегралам. В вероятности это
появляется постоянно: совместные распределения случайных векторов,
математические ожидания функций нескольких случайных величин, свертки
распределений, марковские переходы, корреляционные функции и
спектральные формулы.

\subsubsection{2.2. Конструкция произведения пространств и
мер}\label{ux43aux43eux43dux441ux442ux440ux443ux43aux446ux438ux44f-ux43fux440ux43eux438ux437ux432ux435ux434ux435ux43dux438ux44f-ux43fux440ux43eux441ux442ux440ux430ux43dux441ux442ux432-ux438-ux43cux435ux440}

Пусть заданы два измеримых пространства с мерами:

\[
(\Omega_1,\mathcal F_1,\mu_1),
\qquad
(\Omega_2,\mathcal F_2,\mu_2).
\]

На декартовом произведении \(\Omega_1\times\Omega_2\) рассматривают
произведенную \(\sigma\)-алгебру

\[
\mathcal F_1\otimes\mathcal F_2
=
\sigma\{A\times B:\ A\in\mathcal F_1,\ B\in\mathcal F_2\}.
\]

Если меры \(\mu_1\) и \(\mu_2\) \(\sigma\)-конечны, то существует
единственная мера \(\mu_1\otimes\mu_2\) на
\(\mathcal F_1\otimes\mathcal F_2\), такая что

\[
(\mu_1\otimes\mu_2)(A\times B)=\mu_1(A)\mu_2(B).
\]

Для вероятностных мер \(\sigma\)-конечность автоматически выполнена,
потому что \(P(\Omega)=1\).

\subsubsection{2.3. Сечения функций и теорема
Тонелли}\label{ux441ux435ux447ux435ux43dux438ux44f-ux444ux443ux43dux43aux446ux438ux439-ux438-ux442ux435ux43eux440ux435ux43cux430-ux442ux43eux43dux435ux43bux43bux438}

Пусть

\[
f:\Omega_1\times\Omega_2\to[0,\infty]
\]

измерима относительно \(\mathcal F_1\otimes\mathcal F_2\). Тогда для
фиксированного \(\omega_1\) функция

\[
\omega_2\mapsto f(\omega_1,\omega_2)
\]

измерима по \(\mathcal F_2\), а для фиксированного \(\omega_2\) функция

\[
\omega_1\mapsto f(\omega_1,\omega_2)
\]

измерима по \(\mathcal F_1\). Эти функции называются сечениями функции
\(f\).

Теорема Тонелли утверждает, что для неотрицательной измеримой функции
повторные интегралы существуют как расширенные неотрицательные числа и
равны интегралу по произведенной мере:

\[
\int f\,d(\mu_1\otimes\mu_2)
=
\int_{\Omega_1}
\left(\int_{\Omega_2} f(\omega_1,\omega_2)\,\mu_2(d\omega_2)\right)
\mu_1(d\omega_1)
\]

\[
=
\int_{\Omega_2}
\left(\int_{\Omega_1} f(\omega_1,\omega_2)\,\mu_1(d\omega_1)\right)
\mu_2(d\omega_2).
\]

Здесь не требуется заранее знать конечность интеграла. Можно получить
значение \(+\infty\). Это важно: для неотрицательных функций порядок
интегрирования можно менять без проверки абсолютной интегрируемости.

\subsubsection{2.4. Теорема Фубини для интегрируемых
функций}\label{ux442ux435ux43eux440ux435ux43cux430-ux444ux443ux431ux438ux43dux438-ux434ux43bux44f-ux438ux43dux442ux435ux433ux440ux438ux440ux443ux435ux43cux44bux445-ux444ux443ux43dux43aux446ux438ux439}

Теорема Фубини - это версия для знакопеременных или комплекснозначных
функций. Пусть

\[
f\in L^1(\mu_1\otimes\mu_2),
\qquad
\int_{\Omega_1\times\Omega_2}|f|\,d(\mu_1\otimes\mu_2)<\infty.
\]

Тогда:

\begin{enumerate}
\def\labelenumi{\arabic{enumi}.}
\tightlist
\item
  для \(\mu_1\)-почти всех \(\omega_1\) функция \(f(\omega_1,\cdot)\)
  интегрируема по \(\mu_2\);
\item
  для \(\mu_2\)-почти всех \(\omega_2\) функция \(f(\cdot,\omega_2)\)
  интегрируема по \(\mu_1\);
\item
  функции повторных интегралов измеримы и интегрируемы;
\item
  выполнено равенство
\end{enumerate}

\[
\int_{\Omega_1\times\Omega_2} f\,d(\mu_1\otimes\mu_2)
=
\int_{\Omega_1}
\left(\int_{\Omega_2} f(\omega_1,\omega_2)\,\mu_2(d\omega_2)\right)
\mu_1(d\omega_1)
\]

\[
=
\int_{\Omega_2}
\left(\int_{\Omega_1} f(\omega_1,\omega_2)\,\mu_1(d\omega_1)\right)
\mu_2(d\omega_2).
\]

Главное условие Фубини - интегрируемость \(|f|\) по произведенной мере.
Практически часто проверяют достаточное условие:

\[
\int_{\Omega_1}
\left(\int_{\Omega_2}|f(\omega_1,\omega_2)|\,\mu_2(d\omega_2)\right)
\mu_1(d\omega_1)
<\infty.
\]

Тогда по Тонелли

\[
\int |f|\,d(\mu_1\otimes\mu_2)<\infty,
\]

значит применима теорема Фубини.

\subsubsection{2.5. Независимость событий и
сигма-алгебр}\label{ux43dux435ux437ux430ux432ux438ux441ux438ux43cux43eux441ux442ux44c-ux441ux43eux431ux44bux442ux438ux439-ux438-ux441ux438ux433ux43cux430-ux430ux43bux433ux435ux431ux440}

Теперь перейдем к независимости. Для двух событий \(A,B\in\mathcal F\)
независимость определяется формулой

\[
P(A\cap B)=P(A)P(B).
\]

Интуиция: знание того, что произошло \(A\), не меняет вероятность \(B\).
Если \(P(A)>0\), это равносильно

\[
P(B\mid A)=P(B).
\]

Для нескольких событий \(A_1,\ldots,A_n\) независимость в совокупности
означает, что для любого непустого набора индексов \(i_1,\ldots,i_k\)

\[
P(A_{i_1}\cap\cdots\cap A_{i_k})
=
P(A_{i_1})\cdots P(A_{i_k}).
\]

Важно: попарная независимость не равна независимости в совокупности.

Для \(\sigma\)-алгебр
\(\mathcal G_1,\ldots,\mathcal G_n\subset\mathcal F\) независимость
означает:

\[
P(G_1\cap\cdots\cap G_n)
=
P(G_1)\cdots P(G_n)
\]

для любых \(G_i\in\mathcal G_i\). Случайные элементы \(X_1,\ldots,X_n\)
независимы, если независимы порожденные ими \(\sigma\)-алгебры:

\[
\sigma(X_1),\ldots,\sigma(X_n).
\]

Эквивалентно, для любых борелевских множеств \(B_i\):

\[
P\{X_1\in B_1,\ldots,X_n\in B_n\}
=
\prod_{i=1}^n P\{X_i\in B_i\}.
\]

\subsubsection{2.6. Независимость случайных элементов и совместное
распределение}\label{ux43dux435ux437ux430ux432ux438ux441ux438ux43cux43eux441ux442ux44c-ux441ux43bux443ux447ux430ux439ux43dux44bux445-ux44dux43bux435ux43cux435ux43dux442ux43eux432-ux438-ux441ux43eux432ux43cux435ux441ux442ux43dux43eux435-ux440ux430ux441ux43fux440ux435ux434ux435ux43bux435ux43dux438ux435}

В терминах распределений это записывается особенно компактно. Для двух
случайных элементов

\[
X:\Omega\to E,
\qquad
Y:\Omega\to F
\]

с распределениями \(P_X\) и \(P_Y\) совместное распределение вектора
\((X,Y)\) есть мера

\[
P_{(X,Y)}(C)=P\{(X,Y)\in C\},
\qquad
C\in\mathcal E\otimes\mathcal F.
\]

Тогда

\[
X\ \text{и}\ Y\ \text{независимы}
\quad\Longleftrightarrow\quad
P_{(X,Y)}=P_X\otimes P_Y.
\]

Действительно, на прямоугольниках \(A\times B\) равенство дает

\[
P\{X\in A,\ Y\in B\}
=
P_X(A)P_Y(B),
\]

то есть ровно определение независимости. А поскольку прямоугольники
порождают произведенную \(\sigma\)-алгебру, равенство переносится на все
измеримые множества стандартным аргументом через монотонный класс.

\subsubsection{2.7. Функции распределения и
плотности}\label{ux444ux443ux43dux43aux446ux438ux438-ux440ux430ux441ux43fux440ux435ux434ux435ux43bux435ux43dux438ux44f-ux438-ux43fux43bux43eux442ux43dux43eux441ux442ux438}

Если \(X\) и \(Y\) - вещественные случайные величины, то независимость
эквивалентна факторизации совместной функции распределения:

\[
F_{X,Y}(x,y)
=
P\{X\le x,\ Y\le y\}
=
F_X(x)F_Y(y)
\]

для всех \(x,y\in\mathbb R\).

Если у совместного распределения есть плотность \(p_{X,Y}\) относительно
меры Лебега на \(\mathbb R^2\), то маргинальные плотности выражаются
через повторные интегралы:

\[
p_X(x)=\int_{\mathbb R}p_{X,Y}(x,y)\,dy,
\qquad
p_Y(y)=\int_{\mathbb R}p_{X,Y}(x,y)\,dx.
\]

Здесь используется Фубини/Тонелли. В этом случае

\[
X\ \text{и}\ Y\ \text{независимы}
\quad\Longleftrightarrow\quad
p_{X,Y}(x,y)=p_X(x)p_Y(y)
\]

почти всюду относительно двумерной меры Лебега.

\subsubsection{2.8. Математическое ожидание произведения
функций}\label{ux43cux430ux442ux435ux43cux430ux442ux438ux447ux435ux441ux43aux43eux435-ux43eux436ux438ux434ux430ux43dux438ux435-ux43fux440ux43eux438ux437ux432ux435ux434ux435ux43dux438ux44f-ux444ux443ux43dux43aux446ux438ux439}

Фубини дает основную вычислительную формулу для независимых величин.
Пусть \(X\) и \(Y\) независимы, а \(g\) и \(h\) - измеримые функции.
Если \(g(X)\) и \(h(Y)\) неотрицательны или если все нужные величины
интегрируемы, то

\[
\mathbb E[g(X)h(Y)]
=
\int_{E\times F} g(x)h(y)\,P_{(X,Y)}(dx,dy).
\]

По независимости

\[
P_{(X,Y)}=P_X\otimes P_Y.
\]

По Фубини:

\[
\mathbb E[g(X)h(Y)]
=
\int_E\int_F g(x)h(y)\,P_Y(dy)\,P_X(dx)
\]

\[
=
\left(\int_E g(x)\,P_X(dx)\right)
\left(\int_F h(y)\,P_Y(dy)\right)
=
\mathbb E g(X)\,\mathbb E h(Y).
\]

В частности, если \(X,Y\in L^1\) и независимы, то

\[
\mathbb E[XY]=\mathbb EX\,\mathbb EY.
\]

Если дополнительно \(X,Y\in L^2\), то

\[
\operatorname{Cov}(X,Y)=0.
\]

Но обратное неверно: нулевая ковариация вообще не означает
независимости. Исключение - специальные классы, например совместно
гауссовские величины, что относится к Билет 18.

\subsubsection{2.9. Распределение суммы независимых величин
(свертка)}\label{ux440ux430ux441ux43fux440ux435ux434ux435ux43bux435ux43dux438ux435-ux441ux443ux43cux43cux44b-ux43dux435ux437ux430ux432ux438ux441ux438ux43cux44bux445-ux432ux435ux43bux438ux447ux438ux43d-ux441ux432ux435ux440ux442ux43aux430}

Еще одно важное следствие - распределение суммы независимых величин.
Если \(X\) и \(Y\) независимы, то закон \(X+Y\) равен свертке законов:

\[
P_{X+Y}=P_X*P_Y.
\]

Для борелевского множества \(B\subset\mathbb R\):

\[
P\{X+Y\in B\}
=
\int_{\mathbb R}\int_{\mathbb R}\mathbf 1_B(x+y)\,P_X(dx)P_Y(dy).
\]

В дискретном случае:

\[
P\{X+Y=z\}
=
\sum_x P\{X=x\}P\{Y=z-x\}.
\]

Если есть плотности \(p_X\) и \(p_Y\), то плотность суммы:

\[
p_{X+Y}(z)
=
\int_{\mathbb R}p_X(x)p_Y(z-x)\,dx.
\]

Через характеристические функции это превращается в произведение:

\[
\varphi_{X+Y}(t)=\varphi_X(t)\varphi_Y(t),
\]

что связывает билет с Билет 08 и предельными теоремами.

\subsubsection{2.10. Резюме}\label{ux440ux435ux437ux44eux43cux435}

Итак, логика билета такая: произведенная мера описывает независимое
совместное распределение; Фубини и Тонелли позволяют корректно
интегрировать по этому произведению; поэтому из независимости следуют
факторизация ожиданий, произведение плотностей, свертка распределений и
произведение характеристических функций.

\subsection{3. Основные
определения}\label{ux43eux441ux43dux43eux432ux43dux44bux435-ux43eux43fux440ux435ux434ux435ux43bux435ux43dux438ux44f-8}

\subsubsection{\texorpdfstring{Произведенная
\(\sigma\)-алгебра}{Произведенная \textbackslash sigma-алгебра}}\label{ux43fux440ux43eux438ux437ux432ux435ux434ux435ux43dux43dux430ux44f-sigma-ux430ux43bux433ux435ux431ux440ux430}

Для измеримых пространств \((\Omega_1,\mathcal F_1)\) и
\((\Omega_2,\mathcal F_2)\):

\[
\mathcal F_1\otimes\mathcal F_2
=
\sigma\{A\times B:\ A\in\mathcal F_1,\ B\in\mathcal F_2\}.
\]

\subsubsection{Произведение
мер}\label{ux43fux440ux43eux438ux437ux432ux435ux434ux435ux43dux438ux435-ux43cux435ux440}

Если \(\mu_1\) и \(\mu_2\) \(\sigma\)-конечны, то \(\mu_1\otimes\mu_2\)
- мера на \(\mathcal F_1\otimes\mathcal F_2\), задаваемая на
прямоугольниках формулой

\[
(\mu_1\otimes\mu_2)(A\times B)=\mu_1(A)\mu_2(B).
\]

\subsubsection{Сечение
множества}\label{ux441ux435ux447ux435ux43dux438ux435-ux43cux43dux43eux436ux435ux441ux442ux432ux430}

Для \(C\subset\Omega_1\times\Omega_2\):

\[
C_{\omega_1}=\{\omega_2:\ (\omega_1,\omega_2)\in C\},
\qquad
C^{\omega_2}=\{\omega_1:\ (\omega_1,\omega_2)\in C\}.
\]

\subsubsection{Сечение
функции}\label{ux441ux435ux447ux435ux43dux438ux435-ux444ux443ux43dux43aux446ux438ux438}

Для \(f(\omega_1,\omega_2)\):

\[
f_{\omega_1}(\omega_2)=f(\omega_1,\omega_2),
\qquad
f^{\omega_2}(\omega_1)=f(\omega_1,\omega_2).
\]

\subsubsection{Интегрируемость на
произведении}\label{ux438ux43dux442ux435ux433ux440ux438ux440ux443ux435ux43cux43eux441ux442ux44c-ux43dux430-ux43fux440ux43eux438ux437ux432ux435ux434ux435ux43dux438ux438}

Функция \(f\) интегрируема относительно \(\mu_1\otimes\mu_2\), если

\[
\int_{\Omega_1\times\Omega_2}|f|\,d(\mu_1\otimes\mu_2)<\infty.
\]

\subsubsection{Независимые
события}\label{ux43dux435ux437ux430ux432ux438ux441ux438ux43cux44bux435-ux441ux43eux431ux44bux442ux438ux44f}

События \(A\) и \(B\) независимы, если

\[
P(A\cap B)=P(A)P(B).
\]

\subsubsection{Независимость событий в
совокупности}\label{ux43dux435ux437ux430ux432ux438ux441ux438ux43cux43eux441ux442ux44c-ux441ux43eux431ux44bux442ux438ux439-ux432-ux441ux43eux432ux43eux43aux443ux43fux43dux43eux441ux442ux438}

События \(A_1,\ldots,A_n\) независимы в совокупности, если для любого
непустого набора индексов \(i_1,\ldots,i_k\):

\[
P(A_{i_1}\cap\cdots\cap A_{i_k})
=
P(A_{i_1})\cdots P(A_{i_k}).
\]

\subsubsection{\texorpdfstring{Независимые
\(\sigma\)-алгебры}{Независимые \textbackslash sigma-алгебры}}\label{ux43dux435ux437ux430ux432ux438ux441ux438ux43cux44bux435-sigma-ux430ux43bux433ux435ux431ux440ux44b}

\(\sigma\)-алгебры \(\mathcal G_1,\ldots,\mathcal G_n\) независимы, если
для любых \(G_i\in\mathcal G_i\):

\[
P(G_1\cap\cdots\cap G_n)
=
\prod_{i=1}^nP(G_i).
\]

\subsubsection{Независимые случайные
элементы}\label{ux43dux435ux437ux430ux432ux438ux441ux438ux43cux44bux435-ux441ux43bux443ux447ux430ux439ux43dux44bux435-ux44dux43bux435ux43cux435ux43dux442ux44b}

Случайные элементы \(X_1,\ldots,X_n\) независимы, если независимы
\(\sigma(X_1),\ldots,\sigma(X_n)\).

\subsubsection{Распределительный критерий
независимости}\label{ux440ux430ux441ux43fux440ux435ux434ux435ux43bux438ux442ux435ux43bux44cux43dux44bux439-ux43aux440ux438ux442ux435ux440ux438ux439-ux43dux435ux437ux430ux432ux438ux441ux438ux43cux43eux441ux442ux438}

\[
X_1,\ldots,X_n\ \text{независимы}
\quad\Longleftrightarrow\quad
P_{(X_1,\ldots,X_n)}
=
P_{X_1}\otimes\cdots\otimes P_{X_n}.
\]

\subsubsection{Свертка
распределений}\label{ux441ux432ux435ux440ux442ux43aux430-ux440ux430ux441ux43fux440ux435ux434ux435ux43bux435ux43dux438ux439}

Для вероятностных мер \(\mu\) и \(\nu\) на \(\mathbb R\) свертка
\(\mu*\nu\) задается формулой

\[
(\mu*\nu)(B)
=
\int_{\mathbb R}\int_{\mathbb R}\mathbf 1_B(x+y)\,\mu(dx)\nu(dy).
\]

Если \(X,Y\) независимы, то

\[
P_{X+Y}=P_X*P_Y.
\]

\subsection{4. Основные теоремы и
свойства}\label{ux43eux441ux43dux43eux432ux43dux44bux435-ux442ux435ux43eux440ux435ux43cux44b-ux438-ux441ux432ux43eux439ux441ux442ux432ux430-8}

\subsubsection{Теорема 1. Существование произведения
мер}\label{ux442ux435ux43eux440ux435ux43cux430-1.-ux441ux443ux449ux435ux441ux442ux432ux43eux432ux430ux43dux438ux435-ux43fux440ux43eux438ux437ux432ux435ux434ux435ux43dux438ux44f-ux43cux435ux440}

Если \(\mu_1\) и \(\mu_2\) - \(\sigma\)-конечные меры, то существует
единственная мера \(\mu_1\otimes\mu_2\) на
\(\mathcal F_1\otimes\mathcal F_2\), такая что

\[
(\mu_1\otimes\mu_2)(A\times B)=\mu_1(A)\mu_2(B).
\]

Для вероятностных пространств это условие всегда выполнено.

\subsubsection{Теорема 2.
Тонелли}\label{ux442ux435ux43eux440ux435ux43cux430-2.-ux442ux43eux43dux435ux43bux43bux438}

Пусть \(f\ge0\) и \(f\) измерима на \(\Omega_1\times\Omega_2\). Тогда
функции

\[
\omega_1\mapsto\int_{\Omega_2}f(\omega_1,\omega_2)\,\mu_2(d\omega_2),
\]

\[
\omega_2\mapsto\int_{\Omega_1}f(\omega_1,\omega_2)\,\mu_1(d\omega_1)
\]

измеримы, и

\[
\int f\,d(\mu_1\otimes\mu_2)
=
\int_{\Omega_1}\int_{\Omega_2}f(\omega_1,\omega_2)\,\mu_2(d\omega_2)\mu_1(d\omega_1)
\]

\[
=
\int_{\Omega_2}\int_{\Omega_1}f(\omega_1,\omega_2)\,\mu_1(d\omega_1)\mu_2(d\omega_2).
\]

Значения могут быть равны \(+\infty\).

\subsubsection{Теорема 3.
Фубини}\label{ux442ux435ux43eux440ux435ux43cux430-3.-ux444ux443ux431ux438ux43dux438}

Если \(f\in L^1(\mu_1\otimes\mu_2)\), то повторные интегралы существуют
почти всюду, являются интегрируемыми функциями и

\[
\int f\,d(\mu_1\otimes\mu_2)
=
\int_{\Omega_1}\int_{\Omega_2}f\,d\mu_2\,d\mu_1
=
\int_{\Omega_2}\int_{\Omega_1}f\,d\mu_1\,d\mu_2.
\]

\subsubsection{Следствие. Проверка условия
Фубини}\label{ux441ux43bux435ux434ux441ux442ux432ux438ux435.-ux43fux440ux43eux432ux435ux440ux43aux430-ux443ux441ux43bux43eux432ux438ux44f-ux444ux443ux431ux438ux43dux438}

Если

\[
\int_{\Omega_1}\int_{\Omega_2}|f(\omega_1,\omega_2)|\,\mu_2(d\omega_2)\mu_1(d\omega_1)<\infty,
\]

то \(f\in L^1(\mu_1\otimes\mu_2)\) и порядок интегрирования можно
менять.

\subsubsection{Теорема 4. Независимость и произведение
распределений}\label{ux442ux435ux43eux440ux435ux43cux430-4.-ux43dux435ux437ux430ux432ux438ux441ux438ux43cux43eux441ux442ux44c-ux438-ux43fux440ux43eux438ux437ux432ux435ux434ux435ux43dux438ux435-ux440ux430ux441ux43fux440ux435ux434ux435ux43bux435ux43dux438ux439}

Случайные элементы \(X\) и \(Y\) независимы тогда и только тогда, когда

\[
P_{(X,Y)}=P_X\otimes P_Y.
\]

Для \(n\) случайных элементов:

\[
P_{(X_1,\ldots,X_n)}
=
P_{X_1}\otimes\cdots\otimes P_{X_n}.
\]

\subsubsection{Теорема 5. Факторизация функций
распределения}\label{ux442ux435ux43eux440ux435ux43cux430-5.-ux444ux430ux43aux442ux43eux440ux438ux437ux430ux446ux438ux44f-ux444ux443ux43dux43aux446ux438ux439-ux440ux430ux441ux43fux440ux435ux434ux435ux43bux435ux43dux438ux44f}

Для вещественных случайных величин \(X_1,\ldots,X_n\) независимость
эквивалентна равенству

\[
F_{X_1,\ldots,X_n}(x_1,\ldots,x_n)
=
\prod_{k=1}^n F_{X_k}(x_k)
\]

для всех \(x_1,\ldots,x_n\).

\subsubsection{Теорема 6. Факторизация
плотности}\label{ux442ux435ux43eux440ux435ux43cux430-6.-ux444ux430ux43aux442ux43eux440ux438ux437ux430ux446ux438ux44f-ux43fux43bux43eux442ux43dux43eux441ux442ux438}

Если у \((X,Y)\) есть совместная плотность \(p_{X,Y}\), то \(X\) и \(Y\)
независимы тогда и только тогда, когда

\[
p_{X,Y}(x,y)=p_X(x)p_Y(y)
\]

почти всюду.

\subsubsection{Теорема 7. Ожидание произведения независимых
функций}\label{ux442ux435ux43eux440ux435ux43cux430-7.-ux43eux436ux438ux434ux430ux43dux438ux435-ux43fux440ux43eux438ux437ux432ux435ux434ux435ux43dux438ux44f-ux43dux435ux437ux430ux432ux438ux441ux438ux43cux44bux445-ux444ux443ux43dux43aux446ux438ux439}

Если \(X\) и \(Y\) независимы, то для неотрицательных измеримых \(g,h\):

\[
\mathbb E[g(X)h(Y)]
=
\mathbb E g(X)\,\mathbb E h(Y),
\]

где равенство понимается в расширенном смысле \([0,+\infty]\) через
теорему Тонелли. Случай \(0\cdot\infty\) понимается естественно: если,
например, \(\mathbb E g(X)=0\), то \(g(X)=0\) почти наверное, поэтому
обе части равны нулю. Если \(g,h\) знакопеременные, достаточно условий
\(\mathbb E|g(X)|<\infty\) и \(\mathbb E|h(Y)|<\infty\); при
независимости тогда \(g(X)h(Y)\) интегрируема и формула верна без
расширенных значений.

В частности:

\[
\mathbb E[XY]=\mathbb EX\,\mathbb EY
\]

для независимых \(X,Y\in L^1\).

\subsubsection{Теорема 8. Независимые
преобразования}\label{ux442ux435ux43eux440ux435ux43cux430-8.-ux43dux435ux437ux430ux432ux438ux441ux438ux43cux44bux435-ux43fux440ux435ux43eux431ux440ux430ux437ux43eux432ux430ux43dux438ux44f}

Если \(X\) и \(Y\) независимы, а \(g,h\) измеримы, то \(g(X)\) и
\(h(Y)\) независимы. Более общо, измеримые функции от независимых групп
случайных величин независимы.

\subsubsection{Свойство 9. Сумма независимых
величин}\label{ux441ux432ux43eux439ux441ux442ux432ux43e-9.-ux441ux443ux43cux43cux430-ux43dux435ux437ux430ux432ux438ux441ux438ux43cux44bux445-ux432ux435ux43bux438ux447ux438ux43d-1}

Если \(X\) и \(Y\) независимы, то

\[
P_{X+Y}=P_X*P_Y.
\]

Если есть плотности:

\[
p_{X+Y}(z)=\int_{\mathbb R}p_X(x)p_Y(z-x)\,dx.
\]

Если есть характеристические функции:

\[
\varphi_{X+Y}(t)=\varphi_X(t)\varphi_Y(t).
\]

\subsection{5. Примеры}\label{ux43fux440ux438ux43cux435ux440ux44b-8}

\subsubsection{Пример 1. Повторный интеграл по
прямоугольнику}\label{ux43fux440ux438ux43cux435ux440-1.-ux43fux43eux432ux442ux43eux440ux43dux44bux439-ux438ux43dux442ux435ux433ux440ux430ux43b-ux43fux43e-ux43fux440ux44fux43cux43eux443ux433ux43eux43bux44cux43dux438ux43aux443}

Пусть \(f(x,y)=x+y\) на \([0,1]\times[0,2]\) с мерой Лебега. Тогда

\[
\int_0^1\int_0^2(x+y)\,dy\,dx
=
\int_0^1(2x+2)\,dx
=
3.
\]

В другом порядке:

\[
\int_0^2\int_0^1(x+y)\,dx\,dy
=
\int_0^2\left(\frac12+y\right)\,dy
=
3.
\]

Равенство обеспечивается Фубини, так как функция ограничена на множестве
конечной меры.

\subsubsection{Пример 2. Независимые бернуллиевские
величины}\label{ux43fux440ux438ux43cux435ux440-2.-ux43dux435ux437ux430ux432ux438ux441ux438ux43cux44bux435-ux431ux435ux440ux43dux443ux43bux43bux438ux435ux432ux441ux43aux438ux435-ux432ux435ux43bux438ux447ux438ux43dux44b}

Пусть \(X,Y\) независимы и

\[
P\{X=1\}=p,\qquad P\{Y=1\}=q.
\]

Тогда

\[
P\{X=1,Y=1\}=pq,
\]

а совместное распределение задается произведениями вероятностей:

\[
P\{X=x,Y=y\}=P\{X=x\}P\{Y=y\}.
\]

Для суммы \(S=X+Y\):

\[
P\{S=2\}=pq,
\]

\[
P\{S=1\}=p(1-q)+(1-p)q,
\]

\[
P\{S=0\}=(1-p)(1-q).
\]

Это дискретная свертка.

\subsubsection{Пример 3. Независимые величины с
плотностями}\label{ux43fux440ux438ux43cux435ux440-3.-ux43dux435ux437ux430ux432ux438ux441ux438ux43cux44bux435-ux432ux435ux43bux438ux447ux438ux43dux44b-ux441-ux43fux43bux43eux442ux43dux43eux441ux442ux44fux43cux438}

Пусть \(X\) и \(Y\) независимы, причем \(X\) имеет плотность \(p_X\), а
\(Y\) - плотность \(p_Y\). Тогда

\[
p_{X,Y}(x,y)=p_X(x)p_Y(y).
\]

Например, если \(X,Y\) независимы и равномерны на \([0,1]\), то

\[
p_{X,Y}(x,y)=\mathbf 1_{[0,1]}(x)\mathbf 1_{[0,1]}(y),
\]

то есть совместная плотность постоянна на единичном квадрате.

\subsubsection{Пример 4. Плотность суммы независимых равномерных
величин}\label{ux43fux440ux438ux43cux435ux440-4.-ux43fux43bux43eux442ux43dux43eux441ux442ux44c-ux441ux443ux43cux43cux44b-ux43dux435ux437ux430ux432ux438ux441ux438ux43cux44bux445-ux440ux430ux432ux43dux43eux43cux435ux440ux43dux44bux445-ux432ux435ux43bux438ux447ux438ux43d}

Пусть \(X,Y\) независимы и равномерны на \([0,1]\). Тогда

\[
p_{X+Y}(z)
=
\int_{\mathbb R}\mathbf 1_{[0,1]}(x)\mathbf 1_{[0,1]}(z-x)\,dx.
\]

Отсюда

\[
p_{X+Y}(z)=
\begin{cases}
z, & 0\le z\le1,\\
2-z, & 1<z\le2,\\
0, & \text{иначе}.
\end{cases}
\]

Это стандартный пример свертки.

\subsubsection{Пример 5. Нулевая ковариация без
независимости}\label{ux43fux440ux438ux43cux435ux440-5.-ux43dux443ux43bux435ux432ux430ux44f-ux43aux43eux432ux430ux440ux438ux430ux446ux438ux44f-ux431ux435ux437-ux43dux435ux437ux430ux432ux438ux441ux438ux43cux43eux441ux442ux438}

Пусть \(X\) принимает значения \(-1,0,1\) с вероятностями \(1/3\), а

\[
Y=X^2.
\]

Тогда \(Y\) полностью определяется по \(X\), значит \(X\) и \(Y\)
зависимы. Но

\[
\mathbb EX=0,\qquad
\mathbb E[XY]=\mathbb E[X^3]=0,
\]

поэтому

\[
\operatorname{Cov}(X,Y)=0.
\]

Следовательно, некоррелированность не равна независимости.

\subsubsection{Пример 6. Попарная независимость не дает независимости в
совокупности}\label{ux43fux440ux438ux43cux435ux440-6.-ux43fux43eux43fux430ux440ux43dux430ux44f-ux43dux435ux437ux430ux432ux438ux441ux438ux43cux43eux441ux442ux44c-ux43dux435-ux434ux430ux435ux442-ux43dux435ux437ux430ux432ux438ux441ux438ux43cux43eux441ux442ux438-ux432-ux441ux43eux432ux43eux43aux443ux43fux43dux43eux441ux442ux438}

Пусть бросаются две независимые честные монеты, а события:

\[
A=\{\text{первая монета - орел}\},
\]

\[
B=\{\text{вторая монета - орел}\},
\]

\[
C=\{\text{результаты монет совпали}\}.
\]

Тогда \(A,B,C\) попарно независимы, но

\[
P(A\cap B\cap C)=\frac14
\]

и

\[
P(A)P(B)P(C)=\frac18.
\]

Значит независимости в совокупности нет.

\subsection{6. Схемы доказательств /
обоснований}\label{ux441ux445ux435ux43cux44b-ux434ux43eux43aux430ux437ux430ux442ux435ux43bux44cux441ux442ux432-ux43eux431ux43eux441ux43dux43eux432ux430ux43dux438ux439-1}

\textbf{Схема доказательства существования произведения мер.}

Сначала на прямоугольниках задают

\[
\mu(A\times B)=\mu_1(A)\mu_2(B).
\]

Затем продолжают эту функцию с полукольца или алгебры прямоугольников на
порожденную \(\sigma\)-алгебру. Единственность получается из теоремы
продолжения меры при \(\sigma\)-конечности. Для вероятностных мер это
особенно просто, потому что меры конечны.

\textbf{Схема доказательства Тонелли.}

\begin{enumerate}
\def\labelenumi{\arabic{enumi}.}
\tightlist
\item
  Сначала проверяют формулу для индикатора прямоугольника:
\end{enumerate}

\[
f=\mathbf 1_{A\times B}.
\]

Тогда обе стороны равны \(\mu_1(A)\mu_2(B)\).

\begin{enumerate}
\def\labelenumi{\arabic{enumi}.}
\setcounter{enumi}{1}
\item
  Затем переходят к индикаторам множеств из
  \(\mathcal F_1\otimes\mathcal F_2\) через монотонный класс.
\item
  Потом доказывают формулу для неотрицательных простых функций по
  линейности.
\item
  Для произвольной \(f\ge0\) берут простые функции \(f_n\uparrow f\) и
  применяют теорему Беппо Леви о монотонной сходимости из Билет 07.
\end{enumerate}

\textbf{Схема доказательства Фубини.}

Для интегрируемой функции \(f\) раскладывают

\[
f=f^+-f^-.
\]

Так как

\[
\int |f|\,d(\mu_1\otimes\mu_2)<\infty,
\]

то интегралы от \(f^+\) и \(f^-\) конечны. К ним применяют Тонелли,
затем вычитают равенства. Это дает существование повторных интегралов
почти всюду и равенство всех трех интегралов.

\textbf{Схема доказательства критерия независимости через произведение
распределений.}

Если \(X\) и \(Y\) независимы, то для прямоугольников \(A\times B\):

\[
P_{(X,Y)}(A\times B)
=
P\{X\in A,Y\in B\}
=
P_X(A)P_Y(B)
=
(P_X\otimes P_Y)(A\times B).
\]

Так как прямоугольники порождают произведенную \(\sigma\)-алгебру,
равенство мер на прямоугольниках переносится на всю \(\sigma\)-алгебру.
Обратное следует сразу из той же формулы на прямоугольниках.

\textbf{Схема доказательства формулы
\(\mathbb E[g(X)h(Y)]=\mathbb E g(X)\mathbb E h(Y)\).}

По теореме о замене переменной:

\[
\mathbb E[g(X)h(Y)]
=
\int g(x)h(y)\,P_{(X,Y)}(dx,dy).
\]

По независимости:

\[
P_{(X,Y)}=P_X\otimes P_Y.
\]

По Фубини:

\[
\int g(x)h(y)\,P_X(dx)P_Y(dy)
=
\left(\int g\,dP_X\right)\left(\int h\,dP_Y\right).
\]

\subsection{7. Что написать на
доске}\label{ux447ux442ux43e-ux43dux430ux43fux438ux441ux430ux442ux44c-ux43dux430-ux434ux43eux441ux43aux435-8}

\begin{enumerate}
\def\labelenumi{\arabic{enumi}.}
\tightlist
\item
  Произведенная мера:
\end{enumerate}

\[
(\mu_1\otimes\mu_2)(A\times B)=\mu_1(A)\mu_2(B).
\]

\begin{enumerate}
\def\labelenumi{\arabic{enumi}.}
\setcounter{enumi}{1}
\tightlist
\item
  Тонелли:
\end{enumerate}

\[
f\ge0
\quad\Rightarrow\quad
\int f\,d(\mu_1\otimes\mu_2)
=
\int\int f\,d\mu_2\,d\mu_1
=
\int\int f\,d\mu_1\,d\mu_2.
\]

\begin{enumerate}
\def\labelenumi{\arabic{enumi}.}
\setcounter{enumi}{2}
\tightlist
\item
  Фубини:
\end{enumerate}

\[
f\in L^1(\mu_1\otimes\mu_2)
\quad\Rightarrow\quad
\int f\,d(\mu_1\otimes\mu_2)
=
\int\int f\,d\mu_2\,d\mu_1.
\]

\begin{enumerate}
\def\labelenumi{\arabic{enumi}.}
\setcounter{enumi}{3}
\tightlist
\item
  Независимость:
\end{enumerate}

\[
P(A\cap B)=P(A)P(B).
\]

\begin{enumerate}
\def\labelenumi{\arabic{enumi}.}
\setcounter{enumi}{4}
\tightlist
\item
  Независимость случайных элементов:
\end{enumerate}

\[
P_{(X,Y)}=P_X\otimes P_Y.
\]

\begin{enumerate}
\def\labelenumi{\arabic{enumi}.}
\setcounter{enumi}{5}
\tightlist
\item
  Ожидание произведения:
\end{enumerate}

\[
\mathbb E[g(X)h(Y)]
=
\mathbb E g(X)\,\mathbb E h(Y).
\]

\begin{enumerate}
\def\labelenumi{\arabic{enumi}.}
\setcounter{enumi}{6}
\tightlist
\item
  Плотности и свертка:
\end{enumerate}

\[
p_{X,Y}(x,y)=p_X(x)p_Y(y),
\qquad
p_{X+Y}(z)=\int p_X(x)p_Y(z-x)\,dx.
\]

\subsection{8. Типичные дополнительные
вопросы}\label{ux442ux438ux43fux438ux447ux43dux44bux435-ux434ux43eux43fux43eux43bux43dux438ux442ux435ux43bux44cux43dux44bux435-ux432ux43eux43fux440ux43eux441ux44b-8}

\textbf{Вопрос. Чем Тонелли отличается от Фубини?}

Тонелли работает для \(f\ge0\) и допускает бесконечные значения
интегралов. Фубини работает для знакопеременных функций, но требует
\(f\in L^1\), то есть \(\int |f|<\infty\).

\textbf{Вопрос. Можно ли менять порядок интегрирования без условий?}

Нет. Для неотрицательных функций можно по Тонелли. Для знакопеременных
нужно проверять абсолютную интегрируемость или другое условие,
гарантирующее применимость Фубини.

\textbf{Вопрос. Почему в Фубини повторные интегралы определены только
почти всюду?}

Потому что для отдельных сечений интеграл может быть не определен или
бесконечен на множестве внешней переменной меры ноль. Интеграл Лебега не
различает изменения на нулевых множествах.

\textbf{Вопрос. Что значит независимость случайных величин через
\(\sigma\)-алгебры?}

Это независимость всей информации, содержащейся в величинах. Формально
независимы \(\sigma(X)\) и \(\sigma(Y)\), то есть все события вида
\(\{X\in A\}\) и \(\{Y\in B\}\).

\textbf{Вопрос. Как связаны независимость и совместное распределение?}

Независимость означает, что совместный закон равен произведению
маргинальных законов:

\[
P_{(X,Y)}=P_X\otimes P_Y.
\]

\textbf{Вопрос. Как проверить независимость по функции распределения?}

Для вещественных \(X,Y\) нужно проверить

\[
F_{X,Y}(x,y)=F_X(x)F_Y(y)
\]

для всех \(x,y\).

\textbf{Вопрос. Как проверить независимость по плотности?}

Если есть совместная плотность, то нужно проверить

\[
p_{X,Y}(x,y)=p_X(x)p_Y(y)
\]

почти всюду.

\textbf{Вопрос. Почему из независимости следует
\(\mathbb E[XY]=\mathbb EX\,\mathbb EY\)?}

Потому что совместное распределение равно \(P_X\otimes P_Y\), а по
Фубини интеграл от \(xy\) по произведению мер раскладывается в
произведение интегралов.

\textbf{Вопрос. Верно ли обратное: если
\(\mathbb E[XY]=\mathbb EX\,\mathbb EY\), то \(X,Y\) независимы?}

Нет. Это означает только нулевую ковариацию при наличии вторых моментов.
Независимость сильнее.

\textbf{Вопрос. Почему попарная независимость не равна независимости в
совокупности?}

Потому что попарная независимость проверяет только пересечения двух
событий, а независимость в совокупности требует факторизации всех
пересечений любых поднаборов.

\textbf{Вопрос. Что происходит с функциями от независимых величин?}

Если \(X\) и \(Y\) независимы, а \(g,h\) измеримы, то \(g(X)\) и
\(h(Y)\) тоже независимы.

\textbf{Вопрос. Как распределение суммы связано с независимостью?}

При независимости закон суммы есть свертка:

\[
P_{X+Y}=P_X*P_Y.
\]

Без независимости маргинальных распределений недостаточно: нужен
совместный закон.

\subsection{9. Типичные
ошибки}\label{ux442ux438ux43fux438ux447ux43dux44bux435-ux43eux448ux438ux431ux43aux438-8}

\begin{enumerate}
\def\labelenumi{\arabic{enumi}.}
\tightlist
\item
  Менять порядок интегрирования для знакопеременной функции без проверки
  \(\int |f|<\infty\).
\item
  Называть любую перестановку интегралов ``Фубини'', хотя для \(f\ge0\)
  точнее говорить ``Тонелли''.
\item
  Забывать, что повторные интегралы в Фубини определены почти всюду, а
  не обязательно при каждом значении внешней переменной.
\item
  Путать попарную независимость и независимость в совокупности.
\item
  Считать, что некоррелированность означает независимость.
\item
  Проверять независимость только на одном значении функции
  распределения, а не для всех \(x,y\).
\item
  Забывать ``почти всюду'' в равенстве плотностей
\end{enumerate}

\[
p_{X,Y}=p_Xp_Y.
\]

\begin{enumerate}
\def\labelenumi{\arabic{enumi}.}
\setcounter{enumi}{7}
\tightlist
\item
  Использовать формулу
\end{enumerate}

\[
\mathbb E[g(X)h(Y)]=\mathbb E g(X)\mathbb E h(Y)
\]

без независимости или без условий существования ожиданий. 9. Считать,
что распределение суммы всегда является сверткой маргинальных
распределений. Это верно только при независимости. 10. Путать
произведение случайных величин \(XY\) и произведение мер
\(P_X\otimes P_Y\).

\subsection{10. Связи с другими
билетами}\label{ux441ux432ux44fux437ux438-ux441-ux434ux440ux443ux433ux438ux43cux438-ux431ux438ux43bux435ux442ux430ux43cux438-8}

\textbf{Билет 06.} Совместные распределения, конечномерные распределения
и плотности дают язык, в котором формулируется независимость.

\textbf{Билет 07.} Тонелли и Фубини опираются на интеграл Лебега и
теоремы о монотонной и мажорированной сходимости.

\textbf{Билет 08.} Для независимых сумм характеристические функции
перемножаются:

\[
\varphi_{X+Y}=\varphi_X\varphi_Y.
\]

\textbf{Билет 10.} Условное ожидание использует независимость от
\(\sigma\)-алгебры: если \(X\) независима от \(\mathcal G\), то

\[
\mathbb E[X\mid\mathcal G]=\mathbb EX.
\]

\textbf{Билет 11.} Замена переменной в интеграле используется вместе с
Фубини для перехода от \(\mathbb E g(X,Y)\) к интегралу по
\(P_{(X,Y)}\).

\textbf{Билет 14.} Независимость дает нулевую ковариацию, но нулевая
ковариация не дает независимости.

\textbf{Билет 18.} Для совместно гауссовских величин некоррелированность
эквивалентна независимости.

\textbf{Билет 19.} Независимые приращения и блуждания используют
факторизацию совместных законов и свертки.

\textbf{Билет 21.} Закон больших чисел и центральная предельная теорема
требуют независимости или ее ослабленных вариантов.

\subsection{11. Минимум на
удовлетворительно}\label{ux43cux438ux43dux438ux43cux443ux43c-ux43dux430-ux443ux434ux43eux432ux43bux435ux442ux432ux43eux440ux438ux442ux435ux43bux44cux43dux43e-8}

Нужно уметь сказать:

\begin{enumerate}
\def\labelenumi{\arabic{enumi}.}
\tightlist
\item
  Теорема Фубини позволяет считать интеграл по произведению мер как
  повторный интеграл.
\item
  Для \(f\ge0\) работает Тонелли, для знакопеременной \(f\) нужно
  \(f\in L^1\).
\item
  Независимость событий:
\end{enumerate}

\[
P(A\cap B)=P(A)P(B).
\]

\begin{enumerate}
\def\labelenumi{\arabic{enumi}.}
\setcounter{enumi}{3}
\tightlist
\item
  Независимость случайных величин означает
\end{enumerate}

\[
P_{(X,Y)}=P_X\otimes P_Y.
\]

\begin{enumerate}
\def\labelenumi{\arabic{enumi}.}
\setcounter{enumi}{4}
\tightlist
\item
  При независимости
\end{enumerate}

\[
\mathbb E[XY]=\mathbb EX\,\mathbb EY.
\]

\subsection{12. Минимум на
хорошо}\label{ux43cux438ux43dux438ux43cux443ux43c-ux43dux430-ux445ux43eux440ux43eux448ux43e-8}

Помимо минимума, нужно:

\begin{enumerate}
\def\labelenumi{\arabic{enumi}.}
\tightlist
\item
  различать Тонелли и Фубини;
\item
  записать условие абсолютной интегрируемости;
\item
  сформулировать независимость \(\sigma\)-алгебр и случайных величин;
\item
  объяснить факторизацию функции распределения:
\end{enumerate}

\[
F_{X,Y}(x,y)=F_X(x)F_Y(y);
\]

\begin{enumerate}
\def\labelenumi{\arabic{enumi}.}
\setcounter{enumi}{4}
\tightlist
\item
  привести пример свертки для суммы независимых величин;
\item
  сказать, что попарная независимость не равна независимости в
  совокупности.
\end{enumerate}

\subsection{13. Что добавить для
отлично}\label{ux447ux442ux43e-ux434ux43eux431ux430ux432ux438ux442ux44c-ux434ux43bux44f-ux43eux442ux43bux438ux447ux43dux43e-8}

Для сильного ответа стоит добавить:

\begin{enumerate}
\def\labelenumi{\arabic{enumi}.}
\tightlist
\item
  произведенную \(\sigma\)-алгебру \(\mathcal F_1\otimes\mathcal F_2\);
\item
  существование и единственность произведенной меры при
  \(\sigma\)-конечности;
\item
  доказательную схему через индикаторы прямоугольников, простые функции
  и монотонную сходимость;
\item
  равенство плотностей почти всюду:
\end{enumerate}

\[
p_{X,Y}=p_Xp_Y;
\]

\begin{enumerate}
\def\labelenumi{\arabic{enumi}.}
\setcounter{enumi}{4}
\tightlist
\item
  формулу свертки плотностей;
\item
  контрпример ``некоррелированность не означает независимость'';
\item
  контрпример ``попарная независимость не означает независимость в
  совокупности''.
\end{enumerate}

\subsection{14. Устная версия ответа на 5
минут}\label{ux443ux441ux442ux43dux430ux44f-ux432ux435ux440ux441ux438ux44f-ux43eux442ux432ux435ux442ux430-ux43dux430-5-ux43cux438ux43dux443ux442-8}

Сначала я ввожу произведение измеримых пространств. Если есть
\((\Omega_1,\mathcal F_1,\mu_1)\) и \((\Omega_2,\mathcal F_2,\mu_2)\),
то на произведении берут

\[
\mathcal F_1\otimes\mathcal F_2
=
\sigma\{A\times B\}.
\]

При \(\sigma\)-конечности существует произведенная мера:

\[
(\mu_1\otimes\mu_2)(A\times B)=\mu_1(A)\mu_2(B).
\]

Дальше формулирую Тонелли: если \(f\ge0\) измерима, то

\[
\int f\,d(\mu_1\otimes\mu_2)
=
\int\int f\,d\mu_2\,d\mu_1
=
\int\int f\,d\mu_1\,d\mu_2.
\]

Значение может быть бесконечным. Затем формулирую Фубини: если

\[
\int |f|\,d(\mu_1\otimes\mu_2)<\infty,
\]

то повторные интегралы существуют почти всюду, конечны после внешнего
интегрирования и равны интегралу по произведению. Отличие такое: Тонелли
- для неотрицательных функций без проверки конечности, Фубини - для
интегрируемых функций.

После этого перехожу к независимости. События \(A\) и \(B\) независимы,
если

\[
P(A\cap B)=P(A)P(B).
\]

\(\sigma\)-алгебры независимы, если такая факторизация верна для всех
событий из них. Случайные величины независимы, если независимы
порожденные ими \(\sigma\)-алгебры. Эквивалентная и самая полезная
запись:

\[
P_{(X,Y)}=P_X\otimes P_Y.
\]

Для вещественных величин это значит

\[
F_{X,Y}(x,y)=F_X(x)F_Y(y).
\]

Если есть плотности, то

\[
p_{X,Y}(x,y)=p_X(x)p_Y(y)
\]

почти всюду.

Главное вычислительное следствие: если \(X,Y\) независимы, то

\[
\mathbb E[g(X)h(Y)]
=
\mathbb E g(X)\mathbb E h(Y),
\]

потому что интеграл идет по произведенной мере и раскладывается по
Фубини. В частности, для независимых интегрируемых \(X,Y\):

\[
\mathbb E[XY]=\mathbb EX\,\mathbb EY.
\]

Для суммы независимых величин получаем свертку:

\[
p_{X+Y}(z)=\int p_X(x)p_Y(z-x)\,dx,
\]

а для характеристических функций:

\[
\varphi_{X+Y}(t)=\varphi_X(t)\varphi_Y(t).
\]

В конце называю ловушки: нельзя менять порядок интегрирования без
условий, попарная независимость не равна независимости в совокупности, а
нулевая ковариация не означает независимости.

\subsection{15. Если осталось 2 минуты перед
экзаменом}\label{ux435ux441ux43bux438-ux43eux441ux442ux430ux43bux43eux441ux44c-2-ux43cux438ux43dux443ux442ux44b-ux43fux435ux440ux435ux434-ux44dux43aux437ux430ux43cux435ux43dux43eux43c-8}

Запомнить три блока.

\textbf{Фубини/Тонелли:}

\[
f\ge0
\Rightarrow
\int f\,d(\mu\otimes\nu)
=
\int\int f(x,y)\,\nu(dy)\mu(dx)
=
\int\int f(x,y)\,\mu(dx)\nu(dy).
\]

Для знакопеременной функции нужно

\[
\int |f|\,d(\mu\otimes\nu)<\infty.
\]

\textbf{Независимость:}

\[
P(A\cap B)=P(A)P(B),
\qquad
P_{(X,Y)}=P_X\otimes P_Y.
\]

\textbf{Главные следствия:}

\[
\mathbb E[g(X)h(Y)]
=
\mathbb E g(X)\mathbb E h(Y),
\]

\[
p_{X,Y}=p_Xp_Y,
\qquad
P_{X+Y}=P_X*P_Y.
\]

Главная ловушка: Фубини требует условий, а независимость намного сильнее
некоррелированности.

\newpage

\section{Билет 10. Теорема Радона-Никодима, плотности распределений и
условные математические
ожидания}\label{ux431ux438ux43bux435ux442-10.-ux442ux435ux43eux440ux435ux43cux430-ux440ux430ux434ux43eux43dux430-ux43dux438ux43aux43eux434ux438ux43cux430-ux43fux43bux43eux442ux43dux43eux441ux442ux438-ux440ux430ux441ux43fux440ux435ux434ux435ux43bux435ux43dux438ux439-ux438-ux443ux441ux43bux43eux432ux43dux44bux435-ux43cux430ux442ux435ux43cux430ux442ux438ux447ux435ux441ux43aux438ux435-ux43eux436ux438ux434ux430ux43dux438ux44f}

Источники: Ширяев 1, гл. II, §§6-7; лекции ДСП.

\subsection{1. Короткий
ответ}\label{ux43aux43eux440ux43eux442ux43aux438ux439-ux43eux442ux432ux435ux442-9}

Теорема Радона-Никодима говорит: если одна мера абсолютно непрерывна
относительно другой, то первая выражается через интеграл некоторой
функции по второй мере. Эта функция называется плотностью, или
производной Радона-Никодима.

Главная формула:

\[
\nu\ll\mu
\quad\Longrightarrow\quad
\nu(A)=\int_A f\,d\mu,
\qquad
f=\frac{d\nu}{d\mu}.
\]

В вероятности плотность распределения - это частный случай этой
производной:

\[
P_X\ll\lambda_n
\quad\Longrightarrow\quad
P_X(B)=\int_B p_X(x)\,dx,
\qquad
p_X=\frac{dP_X}{d\lambda_n}.
\]

Условное математическое ожидание - это среднее значение случайной
величины при заданной информации. Если \(X\in L^1\) и
\(\mathcal G\subset\mathcal F\), то

\[
Y=\mathbb E[X\mid\mathcal G]
\]

означает:

\begin{enumerate}
\def\labelenumi{\arabic{enumi}.}
\tightlist
\item
  \(Y\) является \(\mathcal G\)-измеримой;
\item
  для всех \(G\in\mathcal G\)
\end{enumerate}

\[
\int_G Y\,dP=\int_G X\,dP.
\]

Существование условного ожидания получается из теоремы Радона-Никодима:
на \(\mathcal G\) рассматривают меру со знаком

\[
Q(G)=\int_G X\,dP,
\]

она абсолютно непрерывна относительно \(P|_{\mathcal G}\), а ее
производная Радона-Никодима и есть \(\mathbb E[X\mid\mathcal G]\).

\subsection{2. Полный
ответ}\label{ux43fux43eux43bux43dux44bux439-ux43eux442ux432ux435ux442-9}

\subsubsection{2.1. Три уровня производной
Радона-Никодима}\label{ux442ux440ux438-ux443ux440ux43eux432ux43dux44f-ux43fux440ux43eux438ux437ux432ux43eux434ux43dux43eux439-ux440ux430ux434ux43eux43dux430-ux43dux438ux43aux43eux434ux438ux43cux430}

В этом билете соединяются три уровня одной идеи.

Первый уровень - теорема Радона-Никодима. Она отвечает на вопрос: когда
меру \(\nu\) можно записать как интеграл по другой мере \(\mu\)? Ответ:
когда \(\nu\) абсолютно непрерывна относительно \(\mu\), то есть когда
\(\mu\)-нулевые множества являются и \(\nu\)-нулевыми.

Второй уровень - плотности распределений. Если распределение случайного
вектора абсолютно непрерывно относительно меры Лебега, то у него есть
плотность. То есть плотность - это не отдельный объект, который
существует всегда, а производная Радона-Никодима распределения по мере
Лебега.

Третий уровень - условное математическое ожидание. Оно строится как
производная Радона-Никодима меры

\[
G\mapsto\int_G X\,dP
\]

относительно вероятности, но только на меньшей \(\sigma\)-алгебре
\(\mathcal G\). Поэтому условное ожидание - это не ``значение при
фиксированном исходе'', а \(\mathcal G\)-измеримая случайная величина,
которая сохраняет интегралы по всем событиям из \(\mathcal G\).

\subsubsection{2.2. Абсолютная непрерывность
мер}\label{ux430ux431ux441ux43eux43bux44eux442ux43dux430ux44f-ux43dux435ux43fux440ux435ux440ux44bux432ux43dux43eux441ux442ux44c-ux43cux435ux440}

Пусть \((E,\mathcal E)\) - измеримое пространство, а \(\mu\) и \(\nu\) -
меры на нем. Говорят, что \(\nu\) абсолютно непрерывна относительно
\(\mu\), и пишут

\[
\nu\ll\mu,
\]

если

\[
\mu(A)=0\quad\Longrightarrow\quad \nu(A)=0
\]

для всех \(A\in\mathcal E\). Интуитивно: мера \(\nu\) не видит ничего
вне того, что видит \(\mu\); у \(\nu\) нет массы на \(\mu\)-нулевых
множествах.

\subsubsection{2.3. Теорема
Радона-Никодима}\label{ux442ux435ux43eux440ux435ux43cux430-ux440ux430ux434ux43eux43dux430-ux43dux438ux43aux43eux434ux438ux43cux430}

Теорема Радона-Никодима в вероятностно нужной форме звучит так. Пусть
\(\mu\) - \(\sigma\)-конечная мера, а \(\nu\) - \(\sigma\)-конечная
неотрицательная мера на \((E,\mathcal E)\), причем

\[
\nu\ll\mu.
\]

Тогда существует \(\mathcal E\)-измеримая функция

\[
f:E\to[0,\infty]
\]

такая, что для всех \(A\in\mathcal E\)

\[
\nu(A)=\int_A f\,d\mu.
\]

Функция \(f\) называется производной Радона-Никодима меры \(\nu\)
относительно \(\mu\) и обозначается

\[
f=\frac{d\nu}{d\mu}.
\]

Эта функция единственна \(\mu\)-почти всюду. Это значит, что если \(f\)
и \(g\) дают одну и ту же меру:

\[
\nu(A)=\int_A f\,d\mu=\int_A g\,d\mu
\]

для всех \(A\), то

\[
f=g
\qquad \mu\text{-почти всюду}.
\]

Для мер со знаком формулировка аналогична: если \(\eta\) - мера со
знаком конечной вариации или более общо допустимая в стандартной
формулировке, \(\eta\ll\mu\), то существует измеримая функция \(f\),
возможно знакопеременная, такая что

\[
\eta(A)=\int_A f\,d\mu.
\]

Именно этот вариант нужен для условного ожидания, потому что

\[
Q(G)=\int_G X\,dP
\]

может быть мерой со знаком.

\subsubsection{2.4. Плотность распределения как производная
Радона-Никодима}\label{ux43fux43bux43eux442ux43dux43eux441ux442ux44c-ux440ux430ux441ux43fux440ux435ux434ux435ux43bux435ux43dux438ux44f-ux43aux430ux43a-ux43fux440ux43eux438ux437ux432ux43eux434ux43dux430ux44f-ux440ux430ux434ux43eux43dux430-ux43dux438ux43aux43eux434ux438ux43cux430}

Теперь связь с плотностями распределений. Пусть
\(X:\Omega\to\mathbb R^n\) - случайный вектор, а \(P_X\) - его
распределение:

\[
P_X(B)=P\{X\in B\},
\qquad
B\in\mathcal B(\mathbb R^n).
\]

Если

\[
P_X\ll\lambda_n,
\]

где \(\lambda_n\) - \(n\)-мерная мера Лебега, то по теореме
Радона-Никодима существует борелевская функция \(p_X\ge0\), такая что

\[
P_X(B)=\int_B p_X(x)\,\lambda_n(dx)
=
\int_B p_X(x)\,dx.
\]

Эта функция \(p_X\) называется плотностью распределения \(X\)
относительно меры Лебега:

\[
p_X=\frac{dP_X}{d\lambda_n}.
\]

Она определена не поточечно единственным образом, а только почти всюду
относительно \(\lambda_n\). Поэтому значения плотности в отдельных
точках не имеют вероятностного смысла. Вероятности получаются только
после интегрирования по множествам:

\[
P\{X\in B\}=\int_B p_X(x)\,dx.
\]

Для конечномерных распределений случайного процесса \((X_t)_{t\in T}\)
все то же самое применяется к вектору

\[
(X_{t_1},\ldots,X_{t_n}).
\]

Если его распределение абсолютно непрерывно относительно \(\lambda_n\),
то есть плотность конечномерного распределения:

\[
p_{t_1,\ldots,t_n}(x_1,\ldots,x_n)
=
\frac{dP_{(X_{t_1},\ldots,X_{t_n})}}{d\lambda_n}.
\]

Тогда

\[
P\{(X_{t_1},\ldots,X_{t_n})\in B\}
=
\int_B p_{t_1,\ldots,t_n}(x)\,dx.
\]

Если есть совместная плотность \(p_{X,Y}(x,y)\), то маргинальные
плотности находятся интегрированием:

\[
p_X(x)=\int_{\mathbb R^m}p_{X,Y}(x,y)\,dy,
\qquad
p_Y(y)=\int_{\mathbb R^n}p_{X,Y}(x,y)\,dx,
\]

что связано с Билетом 09 о Фубини.

Для условных плотностей используется та же логика. Если \((X,Y)\) имеет
совместную плотность \(p_{X,Y}(x,y)\), то для \(P_Y\)-почти всех \(y\),
для которых

\[
p_Y(y)>0,
\]

можно взять версию условной плотности \(X\) при \(Y=y\) по формуле

\[
p_{X\mid Y}(x\mid y)
=
\frac{p_{X,Y}(x,y)}{p_Y(y)}.
\]

При \(p_Y(y)=0\) значение версии условной плотности задают произвольно.

Тогда

\[
\mathbb E[g(X)\mid Y=y]
=
\int g(x)p_{X\mid Y}(x\mid y)\,dx,
\]

если интеграл определен. А условное ожидание относительно случайной
величины \(Y\) записывается как

\[
\mathbb E[g(X)\mid Y]
=
m(Y),
\]

где

\[
m(y)=\int g(x)p_{X\mid Y}(x\mid y)\,dx.
\]

Важно: формула с условной плотностью - удобный вычислительный случай.
Общий объект \(\mathbb E[X\mid\mathcal G]\) существует и без плотностей.

\subsubsection{2.5. Определение условного математического
ожидания}\label{ux43eux43fux440ux435ux434ux435ux43bux435ux43dux438ux435-ux443ux441ux43bux43eux432ux43dux43eux433ux43e-ux43cux430ux442ux435ux43cux430ux442ux438ux447ux435ux441ux43aux43eux433ux43e-ux43eux436ux438ux434ux430ux43dux438ux44f}

Перейдем к строгому определению условного математического ожидания.
Пусть

\[
(\Omega,\mathcal F,P)
\]

вероятностное пространство, \(\mathcal G\subset\mathcal F\) -
под-\(\sigma\)-алгебра, а \(X\) - интегрируемая случайная величина:

\[
\mathbb E|X|<\infty.
\]

Случайная величина \(Y\) называется условным математическим ожиданием
\(X\) относительно \(\mathcal G\), если:

\begin{enumerate}
\def\labelenumi{\arabic{enumi}.}
\tightlist
\item
  \(Y\) является \(\mathcal G\)-измеримой;
\item
  \(Y\) интегрируема;
\item
  для каждого события \(G\in\mathcal G\)
\end{enumerate}

\[
\int_G Y\,dP=\int_G X\,dP.
\]

Пишут

\[
Y=\mathbb E[X\mid\mathcal G].
\]

Смысл условий такой. \(\mathcal G\) - это доступная информация.
\(\mathcal G\)-измеримость означает, что \(Y\) зависит только от этой
информации. Интегральное равенство означает, что на любом событии,
различимом с помощью \(\mathcal G\), средняя масса \(Y\) такая же, как у
исходной величины \(X\).

\subsubsection{2.6. Существование и единственность условного
ожидания}\label{ux441ux443ux449ux435ux441ux442ux432ux43eux432ux430ux43dux438ux435-ux438-ux435ux434ux438ux43dux441ux442ux432ux435ux43dux43dux43eux441ux442ux44c-ux443ux441ux43bux43eux432ux43dux43eux433ux43e-ux43eux436ux438ux434ux430ux43dux438ux44f}

Существование доказывается через Радона-Никодима. На \(\mathcal G\)
определим функцию множеств

\[
Q(G)=\int_G X\,dP,
\qquad
G\in\mathcal G.
\]

Так как \(X\in L^1\), \(Q\) является конечной мерой со знаком на
\((\Omega,\mathcal G)\). Если

\[
P(G)=0,
\]

то

\[
Q(G)=\int_G X\,dP=0,
\]

значит

\[
Q\ll P|_{\mathcal G}.
\]

По теореме Радона-Никодима существует \(\mathcal G\)-измеримая функция
\(Y\), такая что

\[
Q(G)=\int_G Y\,dP.
\]

То есть

\[
\int_G X\,dP=\int_G Y\,dP
\]

для всех \(G\in\mathcal G\). Именно эта \(Y\) и есть
\(\mathbb E[X\mid\mathcal G]\).

Единственность условного ожидания тоже понимается почти наверное. Если
\(Y_1\) и \(Y_2\) удовлетворяют определению, то

\[
Y_1=Y_2
\qquad P\text{-почти наверное}.
\]

Поэтому условное ожидание - это класс эквивалентности случайных величин
с точностью до изменения на множестве вероятности ноль. Конкретно
выбранную функцию называют версией условного ожидания.

\subsubsection{2.7. Условная вероятность и ожидание относительно
случайной
величины}\label{ux443ux441ux43bux43eux432ux43dux430ux44f-ux432ux435ux440ux43eux44fux442ux43dux43eux441ux442ux44c-ux438-ux43eux436ux438ux434ux430ux43dux438ux435-ux43eux442ux43dux43eux441ux438ux442ux435ux43bux44cux43dux43e-ux441ux43bux443ux447ux430ux439ux43dux43eux439-ux432ux435ux43bux438ux447ux438ux43dux44b}

Условная вероятность - частный случай условного ожидания. Для события
\(A\in\mathcal F\):

\[
P(A\mid\mathcal G)=\mathbb E[\mathbf 1_A\mid\mathcal G].
\]

Это \(\mathcal G\)-измеримая случайная величина, для которой

\[
\int_G P(A\mid\mathcal G)\,dP
=
P(A\cap G),
\qquad
G\in\mathcal G.
\]

Если \(\mathcal G=\sigma(Y)\), то пишут

\[
\mathbb E[X\mid Y]
\]

вместо

\[
\mathbb E[X\mid\sigma(Y)].
\]

Так как \(\mathbb E[X\mid Y]\) является \(\sigma(Y)\)-измеримой, при
стандартных условиях она представляется как функция от \(Y\):

\[
\mathbb E[X\mid Y]=m(Y)
\]

для некоторой борелевской функции \(m\). Эту функцию часто обозначают

\[
m(y)=\mathbb E[X\mid Y=y].
\]

Но нужно помнить: \(\mathbb E[X\mid Y=y]\) как поточечная функция
определена только с точностью до \(P_Y\)-нулевых множеств. При
непрерывном \(Y\) событие \(\{Y=y\}\) обычно имеет вероятность ноль,
поэтому нельзя определять условное ожидание простой формулой
\(\mathbb E[X\mathbf 1_{\{Y=y\}}]/P\{Y=y\}\).

\subsubsection{2.8. Свойства условного
ожидания}\label{ux441ux432ux43eux439ux441ux442ux432ux430-ux443ux441ux43bux43eux432ux43dux43eux433ux43e-ux43eux436ux438ux434ux430ux43dux438ux44f}

Основные свойства условного ожидания следуют из определения и
единственности.

Линейность:

\[
\mathbb E[aX+bZ\mid\mathcal G]
=
a\mathbb E[X\mid\mathcal G]
+b\mathbb E[Z\mid\mathcal G],
\]

если все ожидания корректно определены.

Монотонность:

\[
X\le Z\ \text{п.н.}
\quad\Longrightarrow\quad
\mathbb E[X\mid\mathcal G]\le\mathbb E[Z\mid\mathcal G]\ \text{п.н.}
\]

Сохранение среднего:

\[
\mathbb E\,\mathbb E[X\mid\mathcal G]=\mathbb EX.
\]

Если \(X\) уже \(\mathcal G\)-измерима, то

\[
\mathbb E[X\mid\mathcal G]=X.
\]

Если \(X\) независима от \(\mathcal G\), то

\[
\mathbb E[X\mid\mathcal G]=\mathbb EX.
\]

Башенное свойство: если
\(\mathcal H\subset\mathcal G\subset\mathcal F\), то

\[
\mathbb E[\mathbb E[X\mid\mathcal G]\mid\mathcal H]
=
\mathbb E[X\mid\mathcal H].
\]

В частности,

\[
\mathbb E[\mathbb E[X\mid\mathcal G]]
=
\mathbb EX.
\]

Вынос известного множителя: если \(Z\) ограничена и
\(\mathcal G\)-измерима, то

\[
\mathbb E[ZX\mid\mathcal G]
=
Z\,\mathbb E[X\mid\mathcal G].
\]

Более общо формула верна при интегрируемости \(ZX\) и
\(Z\,\mathbb E[X\mid\mathcal G]\).

Это свойство часто формулируют как ``известное относительно
\(\mathcal G\) можно выносить за знак условного ожидания''.

\subsubsection{2.9. Геометрический смысл и
итоги}\label{ux433ux435ux43eux43cux435ux442ux440ux438ux447ux435ux441ux43aux438ux439-ux441ux43cux44bux441ux43b-ux438-ux438ux442ux43eux433ux438}

Есть и геометрический смысл: если \(X\in L^2\), то
\(\mathbb E[X\mid\mathcal G]\) - ортогональная проекция \(X\) на
подпространство \(\mathcal G\)-измеримых квадратично интегрируемых
случайных величин. Это связывает билет с Билетом 12.

Итак, Радон-Никодим дает общий язык плотностей: плотность меры
относительно меры, плотность распределения относительно Лебега, условная
плотность как отношение совместной и маргинальной плотностей. Условное
ожидание - тот же механизм, примененный к мере \(G\mapsto\int_G X\,dP\)
на под-\(\sigma\)-алгебре.

\subsection{3. Основные
определения}\label{ux43eux441ux43dux43eux432ux43dux44bux435-ux43eux43fux440ux435ux434ux435ux43bux435ux43dux438ux44f-9}

\subsubsection{Абсолютная непрерывность
мер}\label{ux430ux431ux441ux43eux43bux44eux442ux43dux430ux44f-ux43dux435ux43fux440ux435ux440ux44bux432ux43dux43eux441ux442ux44c-ux43cux435ux440-1}

Мера \(\nu\) абсолютно непрерывна относительно меры \(\mu\), если

\[
\mu(A)=0\quad\Rightarrow\quad \nu(A)=0
\]

для всех измеримых \(A\). Обозначение:

\[
\nu\ll\mu.
\]

\subsubsection{Сингулярность
мер}\label{ux441ux438ux43dux433ux443ux43bux44fux440ux43dux43eux441ux442ux44c-ux43cux435ux440}

Меры \(\nu\) и \(\mu\) сингулярны, если существует измеримое множество
\(A\), такое что

\[
\mu(A)=0,\qquad \nu(E\setminus A)=0.
\]

Иными словами, меры сосредоточены на непересекающихся с точностью до
нулевых множеств частях пространства.

\subsubsection{Производная
Радона-Никодима}\label{ux43fux440ux43eux438ux437ux432ux43eux434ux43dux430ux44f-ux440ux430ux434ux43eux43dux430-ux43dux438ux43aux43eux434ux438ux43cux430}

Если \(\nu\ll\mu\) и

\[
\nu(A)=\int_A f\,d\mu
\]

для всех \(A\), то

\[
f=\frac{d\nu}{d\mu}
\]

называется производной Радона-Никодима, или плотностью \(\nu\)
относительно \(\mu\).

\subsubsection{Плотность
распределения}\label{ux43fux43bux43eux442ux43dux43eux441ux442ux44c-ux440ux430ux441ux43fux440ux435ux434ux435ux43bux435ux43dux438ux44f-1}

Если \(P_X\ll\lambda_n\), то

\[
p_X=\frac{dP_X}{d\lambda_n}
\]

называется плотностью распределения \(X\) относительно меры Лебега.

\subsubsection{Плотность конечномерного
распределения}\label{ux43fux43bux43eux442ux43dux43eux441ux442ux44c-ux43aux43eux43dux435ux447ux43dux43eux43cux435ux440ux43dux43eux433ux43e-ux440ux430ux441ux43fux440ux435ux434ux435ux43bux435ux43dux438ux44f}

Для процесса \((X_t)\) плотность вектора \((X_{t_1},\ldots,X_{t_n})\):

\[
p_{t_1,\ldots,t_n}
=
\frac{dP_{(X_{t_1},\ldots,X_{t_n})}}{d\lambda_n},
\]

если соответствующее распределение абсолютно непрерывно относительно
\(\lambda_n\).

\subsubsection{\texorpdfstring{Под-\(\sigma\)-алгебра}{Под-\textbackslash sigma-алгебра}}\label{ux43fux43eux434-sigma-ux430ux43bux433ux435ux431ux440ux430}

\(\mathcal G\subset\mathcal F\) - \(\sigma\)-алгебра, содержащая
доступную информацию. Чем больше \(\mathcal G\), тем больше событий мы
различаем.

\subsubsection{\texorpdfstring{Условное математическое ожидание
относительно
\(\mathcal G\)}{Условное математическое ожидание относительно \textbackslash mathcal G}}\label{ux443ux441ux43bux43eux432ux43dux43eux435-ux43cux430ux442ux435ux43cux430ux442ux438ux447ux435ux441ux43aux43eux435-ux43eux436ux438ux434ux430ux43dux438ux435-ux43eux442ux43dux43eux441ux438ux442ux435ux43bux44cux43dux43e-mathcal-g}

Для \(X\in L^1\):

\[
Y=\mathbb E[X\mid\mathcal G]
\]

если \(Y\) \(\mathcal G\)-измерима, интегрируема и

\[
\int_GY\,dP=\int_GX\,dP
\]

для всех \(G\in\mathcal G\).

\subsubsection{Версия условного
ожидания}\label{ux432ux435ux440ux441ux438ux44f-ux443ux441ux43bux43eux432ux43dux43eux433ux43e-ux43eux436ux438ux434ux430ux43dux438ux44f}

Конкретная \(\mathcal G\)-измеримая функция \(Y\), удовлетворяющая
определению. Две версии совпадают \(P\)-почти наверное.

\subsubsection{\texorpdfstring{Условная вероятность относительно
\(\mathcal G\)}{Условная вероятность относительно \textbackslash mathcal G}}\label{ux443ux441ux43bux43eux432ux43dux430ux44f-ux432ux435ux440ux43eux44fux442ux43dux43eux441ux442ux44c-ux43eux442ux43dux43eux441ux438ux442ux435ux43bux44cux43dux43e-mathcal-g}

\[
P(A\mid\mathcal G)=\mathbb E[\mathbf 1_A\mid\mathcal G].
\]

\subsubsection{Условное ожидание относительно случайной
величины}\label{ux443ux441ux43bux43eux432ux43dux43eux435-ux43eux436ux438ux434ux430ux43dux438ux435-ux43eux442ux43dux43eux441ux438ux442ux435ux43bux44cux43dux43e-ux441ux43bux443ux447ux430ux439ux43dux43eux439-ux432ux435ux43bux438ux447ux438ux43dux44b}

\[
\mathbb E[X\mid Y]
=
\mathbb E[X\mid\sigma(Y)].
\]

Обычно

\[
\mathbb E[X\mid Y]=m(Y)
\]

для некоторой борелевской функции \(m\).

\subsubsection{Условная
плотность}\label{ux443ux441ux43bux43eux432ux43dux430ux44f-ux43fux43bux43eux442ux43dux43eux441ux442ux44c}

Если \((X,Y)\) имеет совместную плотность \(p_{X,Y}\) и \(Y\) имеет
плотность \(p_Y\), то при \(p_Y(y)>0\)

\[
p_{X\mid Y}(x\mid y)
=
\frac{p_{X,Y}(x,y)}{p_Y(y)}.
\]

\subsection{4. Основные теоремы и
свойства}\label{ux43eux441ux43dux43eux432ux43dux44bux435-ux442ux435ux43eux440ux435ux43cux44b-ux438-ux441ux432ux43eux439ux441ux442ux432ux430-9}

\subsubsection{Теорема 1. Радона-Никодима для неотрицательных
мер}\label{ux442ux435ux43eux440ux435ux43cux430-1.-ux440ux430ux434ux43eux43dux430-ux43dux438ux43aux43eux434ux438ux43cux430-ux434ux43bux44f-ux43dux435ux43eux442ux440ux438ux446ux430ux442ux435ux43bux44cux43dux44bux445-ux43cux435ux440}

Пусть \(\mu\) - \(\sigma\)-конечная мера, \(\nu\) - \(\sigma\)-конечная
неотрицательная мера и

\[
\nu\ll\mu.
\]

Тогда существует измеримая функция \(f\ge0\), такая что

\[
\nu(A)=\int_A f\,d\mu
\]

для всех \(A\). Функция \(f\) единственна \(\mu\)-почти всюду и
обозначается

\[
f=\frac{d\nu}{d\mu}.
\]

\subsubsection{Теорема 2. Радона-Никодима для меры со
знаком}\label{ux442ux435ux43eux440ux435ux43cux430-2.-ux440ux430ux434ux43eux43dux430-ux43dux438ux43aux43eux434ux438ux43cux430-ux434ux43bux44f-ux43cux435ux440ux44b-ux441ux43e-ux437ux43dux430ux43aux43eux43c}

Если \(\eta\) - конечная мера со знаком, \(\eta\ll\mu\), а \(\mu\)
\(\sigma\)-конечна, то существует измеримая функция \(f\), интегрируемая
по \(\mu\), такая что

\[
\eta(A)=\int_A f\,d\mu.
\]

Этот вариант нужен для построения \(\mathbb E[X\mid\mathcal G]\) при
знакопеременной \(X\).

\subsubsection{Теорема 3. Плотность распределения как производная
Радона-Никодима}\label{ux442ux435ux43eux440ux435ux43cux430-3.-ux43fux43bux43eux442ux43dux43eux441ux442ux44c-ux440ux430ux441ux43fux440ux435ux434ux435ux43bux435ux43dux438ux44f-ux43aux430ux43a-ux43fux440ux43eux438ux437ux432ux43eux434ux43dux430ux44f-ux440ux430ux434ux43eux43dux430-ux43dux438ux43aux43eux434ux438ux43cux430}

Для случайного вектора \(X\) в \(\mathbb R^n\) плотность существует
тогда и только тогда, когда

\[
P_X\ll\lambda_n.
\]

Тогда

\[
p_X=\frac{dP_X}{d\lambda_n},
\qquad
P\{X\in B\}=\int_B p_X(x)\,dx.
\]

\subsubsection{Свойство 4. Плотность определена почти
всюду}\label{ux441ux432ux43eux439ux441ux442ux432ux43e-4.-ux43fux43bux43eux442ux43dux43eux441ux442ux44c-ux43eux43fux440ux435ux434ux435ux43bux435ux43dux430-ux43fux43eux447ux442ux438-ux432ux441ux44eux434ux443}

Если \(p\) и \(\tilde p\) - две плотности одного распределения, то

\[
p=\tilde p
\qquad \lambda_n\text{-почти всюду}.
\]

Значения плотности в отдельных точках можно менять без изменения
распределения.

\subsubsection{Теорема 5. Существование условного математического
ожидания}\label{ux442ux435ux43eux440ux435ux43cux430-5.-ux441ux443ux449ux435ux441ux442ux432ux43eux432ux430ux43dux438ux435-ux443ux441ux43bux43eux432ux43dux43eux433ux43e-ux43cux430ux442ux435ux43cux430ux442ux438ux447ux435ux441ux43aux43eux433ux43e-ux43eux436ux438ux434ux430ux43dux438ux44f}

Если \(X\in L^1\) и \(\mathcal G\subset\mathcal F\), то существует

\[
\mathbb E[X\mid\mathcal G],
\]

причем оно единственно \(P\)-почти наверное.

\subsubsection{Доказательство}\label{ux434ux43eux43aux430ux437ux430ux442ux435ux43bux44cux441ux442ux432ux43e-31}

На \((\Omega,\mathcal G)\) задаем меру со знаком

\[
Q(G)=\int_GX\,dP.
\]

Она абсолютно непрерывна относительно \(P|_{\mathcal G}\). По
Радону-Никодиму существует \(\mathcal G\)-измеримая \(Y\), такая что

\[
Q(G)=\int_GY\,dP.
\]

Это и есть определение условного ожидания.

\subsubsection{Свойство 6.
Линейность}\label{ux441ux432ux43eux439ux441ux442ux432ux43e-6.-ux43bux438ux43dux435ux439ux43dux43eux441ux442ux44c}

\[
\mathbb E[aX+bZ\mid\mathcal G]
=
a\mathbb E[X\mid\mathcal G]
+b\mathbb E[Z\mid\mathcal G].
\]

\subsubsection{Свойство 7.
Монотонность}\label{ux441ux432ux43eux439ux441ux442ux432ux43e-7.-ux43cux43eux43dux43eux442ux43eux43dux43dux43eux441ux442ux44c}

\[
X\le Z\ \text{п.н.}
\quad\Rightarrow\quad
\mathbb E[X\mid\mathcal G]\le \mathbb E[Z\mid\mathcal G]\ \text{п.н.}
\]

\subsubsection{Свойство 8. Сохранение
среднего}\label{ux441ux432ux43eux439ux441ux442ux432ux43e-8.-ux441ux43eux445ux440ux430ux43dux435ux43dux438ux435-ux441ux440ux435ux434ux43dux435ux433ux43e}

\[
\mathbb E\,\mathbb E[X\mid\mathcal G]=\mathbb EX.
\]

\subsubsection{Свойство 9. Если величина уже
известна}\label{ux441ux432ux43eux439ux441ux442ux432ux43e-9.-ux435ux441ux43bux438-ux432ux435ux43bux438ux447ux438ux43dux430-ux443ux436ux435-ux438ux437ux432ux435ux441ux442ux43dux430}

Если \(X\) \(\mathcal G\)-измерима, то

\[
\mathbb E[X\mid\mathcal G]=X.
\]

\subsubsection{Свойство 10.
Независимость}\label{ux441ux432ux43eux439ux441ux442ux432ux43e-10.-ux43dux435ux437ux430ux432ux438ux441ux438ux43cux43eux441ux442ux44c}

Если \(X\) независима от \(\mathcal G\), то

\[
\mathbb E[X\mid\mathcal G]=\mathbb EX.
\]

\subsubsection{Свойство 11. Башенное
свойство}\label{ux441ux432ux43eux439ux441ux442ux432ux43e-11.-ux431ux430ux448ux435ux43dux43dux43eux435-ux441ux432ux43eux439ux441ux442ux432ux43e}

Если \(\mathcal H\subset\mathcal G\), то

\[
\mathbb E[\mathbb E[X\mid\mathcal G]\mid\mathcal H]
=
\mathbb E[X\mid\mathcal H].
\]

\subsubsection{Свойство 12. Вынос известного
множителя}\label{ux441ux432ux43eux439ux441ux442ux432ux43e-12.-ux432ux44bux43dux43eux441-ux438ux437ux432ux435ux441ux442ux43dux43eux433ux43e-ux43cux43dux43eux436ux438ux442ux435ux43bux44f}

Если \(Z\) \(\mathcal G\)-измерима и \(ZX\in L^1\), то

\[
\mathbb E[ZX\mid\mathcal G]
=
Z\mathbb E[X\mid\mathcal G].
\]

\subsubsection{Свойство 13. Условная
вероятность}\label{ux441ux432ux43eux439ux441ux442ux432ux43e-13.-ux443ux441ux43bux43eux432ux43dux430ux44f-ux432ux435ux440ux43eux44fux442ux43dux43eux441ux442ux44c}

Для \(A\in\mathcal F\):

\[
P(A\mid\mathcal G)=\mathbb E[\mathbf 1_A\mid\mathcal G],
\]

и для всех \(G\in\mathcal G\)

\[
\int_G P(A\mid\mathcal G)\,dP=P(A\cap G).
\]

\subsubsection{Свойство 14. Условное ожидание при совместной
плотности}\label{ux441ux432ux43eux439ux441ux442ux432ux43e-14.-ux443ux441ux43bux43eux432ux43dux43eux435-ux43eux436ux438ux434ux430ux43dux438ux435-ux43fux440ux438-ux441ux43eux432ux43cux435ux441ux442ux43dux43eux439-ux43fux43bux43eux442ux43dux43eux441ux442ux438}

Если \((X,Y)\) имеет совместную плотность \(p_{X,Y}\), то

\[
\mathbb E[g(X)\mid Y=y]
=
\int g(x)\frac{p_{X,Y}(x,y)}{p_Y(y)}\,dx
\]

при \(p_Y(y)>0\). Тогда

\[
\mathbb E[g(X)\mid Y]=m(Y),
\]

где

\[
m(y)=\int g(x)p_{X\mid Y}(x\mid y)\,dx.
\]

\subsection{5. Примеры}\label{ux43fux440ux438ux43cux435ux440ux44b-9}

\subsubsection{Пример 1. Обычная плотность как
Радон-Никодим}\label{ux43fux440ux438ux43cux435ux440-1.-ux43eux431ux44bux447ux43dux430ux44f-ux43fux43bux43eux442ux43dux43eux441ux442ux44c-ux43aux430ux43a-ux440ux430ux434ux43eux43d-ux43dux438ux43aux43eux434ux438ux43c}

Пусть \(X\) равномерна на \([0,1]\). Тогда

\[
P_X(B)=\lambda(B\cap[0,1]).
\]

Значит

\[
P_X\ll\lambda,
\qquad
\frac{dP_X}{d\lambda}(x)=\mathbf 1_{[0,1]}(x)
\]

почти всюду.

\subsubsection{Пример 2. Нет плотности относительно
Лебега}\label{ux43fux440ux438ux43cux435ux440-2.-ux43dux435ux442-ux43fux43bux43eux442ux43dux43eux441ux442ux438-ux43eux442ux43dux43eux441ux438ux442ux435ux43bux44cux43dux43e-ux43bux435ux431ux435ux433ux430}

Если \(X\equiv a\), то

\[
P_X=\delta_a.
\]

Мера \(\delta_a\) не абсолютно непрерывна относительно меры Лебега,
потому что

\[
\lambda(\{a\})=0,
\qquad
\delta_a(\{a\})=1.
\]

Значит плотности относительно Лебега нет. Это важный контрпример:
распределение есть всегда, плотность - не всегда.

\subsubsection{Пример 3. Условное ожидание по конечному
разбиению}\label{ux43fux440ux438ux43cux435ux440-3.-ux443ux441ux43bux43eux432ux43dux43eux435-ux43eux436ux438ux434ux430ux43dux438ux435-ux43fux43e-ux43aux43eux43dux435ux447ux43dux43eux43cux443-ux440ux430ux437ux431ux438ux435ux43dux438ux44e}

Пусть

\[
\mathcal G=\sigma(D_1,\ldots,D_n),
\]

где \(D_i\) попарно не пересекаются, \(\bigcup_iD_i=\Omega\),
\(P(D_i)>0\). Тогда

\[
\mathbb E[X\mid\mathcal G]
=
\sum_{i=1}^n
\frac{\int_{D_i}X\,dP}{P(D_i)}\mathbf 1_{D_i}.
\]

То есть на каждом атоме разбиения условное ожидание равно среднему \(X\)
по этому атому.

\subsubsection{Пример 4. Условное ожидание при
независимости}\label{ux43fux440ux438ux43cux435ux440-4.-ux443ux441ux43bux43eux432ux43dux43eux435-ux43eux436ux438ux434ux430ux43dux438ux435-ux43fux440ux438-ux43dux435ux437ux430ux432ux438ux441ux438ux43cux43eux441ux442ux438}

Если \(X\) независима от \(\mathcal G\), то

\[
\mathbb E[X\mid\mathcal G]=\mathbb EX.
\]

Например, если \(X\) и \(Y\) независимы, то

\[
\mathbb E[X\mid Y]=\mathbb EX.
\]

\subsubsection{Пример 5. Если величина уже
известна}\label{ux43fux440ux438ux43cux435ux440-5.-ux435ux441ux43bux438-ux432ux435ux43bux438ux447ux438ux43dux430-ux443ux436ux435-ux438ux437ux432ux435ux441ux442ux43dux430}

Если \(X\) является функцией от \(Y\), например

\[
X=Y^2,
\]

то \(X\) \(\sigma(Y)\)-измерима, поэтому

\[
\mathbb E[X\mid Y]=X=Y^2.
\]

\subsubsection{Пример 6. Условное ожидание через совместную
плотность}\label{ux43fux440ux438ux43cux435ux440-6.-ux443ux441ux43bux43eux432ux43dux43eux435-ux43eux436ux438ux434ux430ux43dux438ux435-ux447ux435ux440ux435ux437-ux441ux43eux432ux43cux435ux441ux442ux43dux443ux44e-ux43fux43bux43eux442ux43dux43eux441ux442ux44c}

Пусть \((X,Y)\) имеет совместную плотность \(p_{X,Y}\). Тогда

\[
p_Y(y)=\int p_{X,Y}(x,y)\,dx.
\]

Если \(p_Y(y)>0\), то

\[
\mathbb E[X\mid Y=y]
=
\int x\frac{p_{X,Y}(x,y)}{p_Y(y)}\,dx.
\]

В частности,

\[
\mathbb E[X\mid Y]=m(Y),
\]

где

\[
m(y)=\int x\,p_{X\mid Y}(x\mid y)\,dx.
\]

\subsubsection{Пример 7. Условное ожидание как
прогноз}\label{ux43fux440ux438ux43cux435ux440-7.-ux443ux441ux43bux43eux432ux43dux43eux435-ux43eux436ux438ux434ux430ux43dux438ux435-ux43aux430ux43a-ux43fux440ux43eux433ux43dux43eux437}

Пусть

\[
X=Y+\varepsilon,
\]

где \(\varepsilon\) независима от \(Y\) и \(\mathbb E\varepsilon=0\).
Тогда

\[
\mathbb E[X\mid Y]
=
\mathbb E[Y+\varepsilon\mid Y]
=
Y+\mathbb E[\varepsilon\mid Y]
=
Y.
\]

Здесь \(Y\) - известная часть, а независимый шум усредняется до нуля.

\subsection{6. Схемы доказательств /
обоснований}\label{ux441ux445ux435ux43cux44b-ux434ux43eux43aux430ux437ux430ux442ux435ux43bux44cux441ux442ux432-ux43eux431ux43eux441ux43dux43eux432ux430ux43dux438ux439-2}

\textbf{Идея теоремы Радона-Никодима.}

Полное доказательство - теорема теории меры. Интуитивная идея такая:
если \(\nu\ll\mu\), то \(\nu\) не имеет массы вне \(\mu\)-массы, значит
ее можно описать локальной ``плотностью'' относительно \(\mu\). В
простом случае, когда пространство разбито на конечные атомы \(A_i\) с
\(\mu(A_i)>0\), плотность равна

\[
f=\sum_i \frac{\nu(A_i)}{\mu(A_i)}\mathbf 1_{A_i}.
\]

Общая теорема строит аналог этой локальной плотности для произвольной
\(\sigma\)-алгебры.

\textbf{Доказательство существования условного ожидания.}

\begin{enumerate}
\def\labelenumi{\arabic{enumi}.}
\tightlist
\item
  Берем \(X\in L^1\) и под-\(\sigma\)-алгебру \(\mathcal G\).
\item
  На \(\mathcal G\) задаем
\end{enumerate}

\[
Q(G)=\int_GX\,dP.
\]

\begin{enumerate}
\def\labelenumi{\arabic{enumi}.}
\setcounter{enumi}{2}
\tightlist
\item
  \(Q\) - конечная мера со знаком.
\item
  Если \(P(G)=0\), то \(Q(G)=0\), значит \(Q\ll P|_{\mathcal G}\).
\item
  По Радону-Никодиму существует \(Y=dQ/dP|_{\mathcal G}\).
\item
  Получаем
\end{enumerate}

\[
\int_GY\,dP=Q(G)=\int_GX\,dP.
\]

Значит \(Y=\mathbb E[X\mid\mathcal G]\).

\textbf{Доказательство единственности условного ожидания.}

Пусть \(Y_1\) и \(Y_2\) удовлетворяют определению. Тогда для всех
\(G\in\mathcal G\)

\[
\int_G(Y_1-Y_2)\,dP=0.
\]

Положим

\[
G_+=\{Y_1>Y_2\}.
\]

Так как \(Y_1,Y_2\) \(\mathcal G\)-измеримы, то \(G_+\in\mathcal G\).
Тогда

\[
\int_{G_+}(Y_1-Y_2)\,dP=0.
\]

Но подынтегральная функция положительна на \(G_+\), значит \(P(G_+)=0\).
Аналогично для \(\{Y_2>Y_1\}\). Поэтому \(Y_1=Y_2\) почти наверное.

\textbf{Доказательство башенного свойства.}

Пусть \(\mathcal H\subset\mathcal G\). Обозначим

\[
Z=\mathbb E[\mathbb E[X\mid\mathcal G]\mid\mathcal H].
\]

Для любого \(H\in\mathcal H\) имеем \(H\in\mathcal G\), поэтому

\[
\int_H Z\,dP
=
\int_H \mathbb E[X\mid\mathcal G]\,dP
=
\int_H X\,dP.
\]

Кроме того, \(Z\) \(\mathcal H\)-измерима. Значит по определению

\[
Z=\mathbb E[X\mid\mathcal H].
\]

\textbf{Доказательство свойства независимости.}

Если \(X\) независима от \(\mathcal G\), то постоянная \(\mathbb EX\)
является \(\mathcal G\)-измеримой. Для любого \(G\in\mathcal G\):

\[
\int_G X\,dP
=
\mathbb E[X\mathbf 1_G]
=
\mathbb EX\,P(G)
=
\int_G \mathbb EX\,dP.
\]

По единственности

\[
\mathbb E[X\mid\mathcal G]=\mathbb EX.
\]

\textbf{Доказательство формулы для условной плотности.}

Для борелевского \(A\):

\[
P\{X\in A,\ Y\in dy\}
\]

имеет плотность по \(y\), равную

\[
\int_A p_{X,Y}(x,y)\,dx.
\]

Делим на маргинальную плотность

\[
p_Y(y)=\int p_{X,Y}(x,y)\,dx
\]

и получаем условную плотность

\[
p_{X\mid Y}(x\mid y)=\frac{p_{X,Y}(x,y)}{p_Y(y)}.
\]

\subsection{7. Что написать на
доске}\label{ux447ux442ux43e-ux43dux430ux43fux438ux441ux430ux442ux44c-ux43dux430-ux434ux43eux441ux43aux435-9}

\begin{enumerate}
\def\labelenumi{\arabic{enumi}.}
\tightlist
\item
  Абсолютная непрерывность:
\end{enumerate}

\[
\nu\ll\mu
\quad\Longleftrightarrow\quad
\mu(A)=0\Rightarrow\nu(A)=0.
\]

\begin{enumerate}
\def\labelenumi{\arabic{enumi}.}
\setcounter{enumi}{1}
\tightlist
\item
  Радон-Никодим:
\end{enumerate}

\[
\nu(A)=\int_A \frac{d\nu}{d\mu}\,d\mu.
\]

\begin{enumerate}
\def\labelenumi{\arabic{enumi}.}
\setcounter{enumi}{2}
\tightlist
\item
  Плотность распределения:
\end{enumerate}

\[
P_X\ll\lambda_n,
\qquad
p_X=\frac{dP_X}{d\lambda_n},
\qquad
P_X(B)=\int_Bp_X(x)\,dx.
\]

\begin{enumerate}
\def\labelenumi{\arabic{enumi}.}
\setcounter{enumi}{3}
\tightlist
\item
  Условное ожидание:
\end{enumerate}

\[
Y=\mathbb E[X\mid\mathcal G]
\iff
\begin{cases}
Y\ \mathcal G\text{-измерима},\\
\int_GY\,dP=\int_GX\,dP,\quad G\in\mathcal G.
\end{cases}
\]

\begin{enumerate}
\def\labelenumi{\arabic{enumi}.}
\setcounter{enumi}{4}
\tightlist
\item
  Построение через Радона-Никодима:
\end{enumerate}

\[
Q(G)=\int_GX\,dP,
\qquad
\mathbb E[X\mid\mathcal G]=\frac{dQ}{dP|_{\mathcal G}}.
\]

\begin{enumerate}
\def\labelenumi{\arabic{enumi}.}
\setcounter{enumi}{5}
\tightlist
\item
  Главные свойства:
\end{enumerate}

\[
\mathbb E\,\mathbb E[X\mid\mathcal G]=\mathbb EX,
\]

\[
\mathcal H\subset\mathcal G
\Rightarrow
\mathbb E[\mathbb E[X\mid\mathcal G]\mid\mathcal H]
=
\mathbb E[X\mid\mathcal H],
\]

\[
X\perp\mathcal G
\Rightarrow
\mathbb E[X\mid\mathcal G]=\mathbb EX.
\]

\begin{enumerate}
\def\labelenumi{\arabic{enumi}.}
\setcounter{enumi}{6}
\tightlist
\item
  Условная плотность:
\end{enumerate}

\[
p_{X\mid Y}(x\mid y)=\frac{p_{X,Y}(x,y)}{p_Y(y)}.
\]

\subsection{8. Типичные дополнительные
вопросы}\label{ux442ux438ux43fux438ux447ux43dux44bux435-ux434ux43eux43fux43eux43bux43dux438ux442ux435ux43bux44cux43dux44bux435-ux432ux43eux43fux440ux43eux441ux44b-9}

\textbf{Вопрос. Что означает \(\nu\ll\mu\)?}

Это значит, что всякое \(\mu\)-нулевое множество является
\(\nu\)-нулевым:

\[
\mu(A)=0\Rightarrow\nu(A)=0.
\]

\textbf{Вопрос. Почему плотность распределения - это производная
Радона-Никодима?}

Потому что плотность \(p_X\) задает распределение через интеграл по мере
Лебега:

\[
P_X(B)=\int_Bp_X\,d\lambda.
\]

Это ровно представление Радона-Никодима для \(P_X\ll\lambda\).

\textbf{Вопрос. Любое ли распределение имеет плотность?}

Нет. Распределение Дирака \(\delta_a\) не имеет плотности относительно
меры Лебега, потому что оно дает массу точке, а мера Лебега этой точки
равна нулю.

\textbf{Вопрос. Плотность единственна?}

Только почти всюду относительно меры, по которой берется плотность.
Значения на нулевых множествах можно менять.

\textbf{Вопрос. Что такое условное ожидание строго?}

Это \(\mathcal G\)-измеримая случайная величина \(Y\), такая что

\[
\int_GY\,dP=\int_GX\,dP
\]

для всех \(G\in\mathcal G\).

\textbf{Вопрос. Почему условное ожидание существует?}

Потому что мера со знаком

\[
Q(G)=\int_GX\,dP
\]

абсолютно непрерывна относительно \(P\) на \(\mathcal G\), и можно
применить теорему Радона-Никодима.

\textbf{Вопрос. Почему условное ожидание определено только почти
наверное?}

Потому что производная Радона-Никодима единственна только с точностью до
изменения на множестве меры ноль.

\textbf{Вопрос. Что будет, если \(\mathcal G\) тривиальна?}

Если

\[
\mathcal G=\{\varnothing,\Omega\},
\]

то

\[
\mathbb E[X\mid\mathcal G]=\mathbb EX.
\]

\textbf{Вопрос. Что будет, если \(X\) \(\mathcal G\)-измерима?}

Тогда

\[
\mathbb E[X\mid\mathcal G]=X.
\]

\textbf{Вопрос. Что будет, если \(X\) независима от \(\mathcal G\)?}

Тогда

\[
\mathbb E[X\mid\mathcal G]=\mathbb EX.
\]

\textbf{Вопрос. Что такое башенное свойство?}

Если \(\mathcal H\subset\mathcal G\), то

\[
\mathbb E[\mathbb E[X\mid\mathcal G]\mid\mathcal H]
=
\mathbb E[X\mid\mathcal H].
\]

\textbf{Вопрос. Чем условная вероятность отличается от условного
ожидания?}

Условная вероятность - частный случай:

\[
P(A\mid\mathcal G)=\mathbb E[\mathbf 1_A\mid\mathcal G].
\]

\textbf{Вопрос. Почему нельзя всегда писать
\(P(A\mid Y=y)=P(A\cap\{Y=y\})/P(Y=y)\)?}

Если \(Y\) непрерывна, то обычно

\[
P(Y=y)=0.
\]

Нужно использовать условное ожидание относительно \(\sigma(Y)\) или
регулярные условные распределения.

\textbf{Вопрос. Как найти \(\mathbb E[X\mid Y]\) при совместной
плотности?}

Сначала найти

\[
p_Y(y)=\int p_{X,Y}(x,y)\,dx,
\]

затем

\[
p_{X\mid Y}(x\mid y)=\frac{p_{X,Y}(x,y)}{p_Y(y)},
\]

и после этого

\[
\mathbb E[X\mid Y=y]=\int x p_{X\mid Y}(x\mid y)\,dx.
\]

\subsection{9. Типичные
ошибки}\label{ux442ux438ux43fux438ux447ux43dux44bux435-ux43eux448ux438ux431ux43aux438-9}

\begin{enumerate}
\def\labelenumi{\arabic{enumi}.}
\tightlist
\item
  Говорить о плотности распределения, не проверив абсолютную
  непрерывность относительно меры Лебега.
\item
  Считать значение плотности \(p(x)\) вероятностью точки.
\item
  Забывать, что производная Радона-Никодима единственна только почти
  всюду.
\item
  Путать абсолютную непрерывность меры \(\nu\ll\mu\) с абсолютной
  непрерывностью функции распределения.
\item
  Путать условное ожидание и условную вероятность: условная вероятность
  - это условное ожидание индикатора.
\item
  Определять \(\mathbb E[X\mid Y=y]\) через деление на \(P(Y=y)\) в
  непрерывном случае.
\item
  Забывать условие \(X\in L^1\) для обычного условного ожидания.
\item
  Думать, что \(\mathbb E[X\mid\mathcal G]\) - число. Обычно это
  случайная величина.
\item
  Забывать \(\mathcal G\)-измеримость в определении условного ожидания.
\item
  Применять формулу условной плотности при \(p_Y(y)=0\).
\item
  Путать равенство почти наверное с поточечным равенством.
\item
  Забывать, что при независимости условное ожидание становится
  константой \(\mathbb EX\).
\end{enumerate}

\subsection{10. Связи с другими
билетами}\label{ux441ux432ux44fux437ux438-ux441-ux434ux440ux443ux433ux438ux43cux438-ux431ux438ux43bux435ux442ux430ux43cux438-9}

\textbf{Билет 06.} Плотности распределений в общем виде определяются как
производные Радона-Никодима.

\textbf{Билет 07.} Все определения используют интеграл Лебега и свойства
интегрируемости.

\textbf{Билет 09.} Совместные и маргинальные плотности, условные
плотности и независимость используют Фубини.

\textbf{Билет 11.} Замена переменной позволяет записывать

\[
\mathbb E g(X)=\int g\,dP_X,
\]

а при плотности

\[
\mathbb E g(X)=\int g(x)p_X(x)\,dx.
\]

\textbf{Билет 12.} В \(L^2\) условное ожидание является ортогональным
проектором на подпространство \(\mathcal G\)-измеримых величин.

\textbf{Билет 14.} Условное ожидание связано с регрессией и оптимальным
среднеквадратичным прогнозом.

\textbf{Билет 22.} Марковские переходные функции можно понимать как
условные распределения будущего при известном текущем состоянии.

\textbf{Билет 23.} Марковское свойство формулируется через условные
вероятности и условные распределения.

\subsection{11. Минимум на
удовлетворительно}\label{ux43cux438ux43dux438ux43cux443ux43c-ux43dux430-ux443ux434ux43eux432ux43bux435ux442ux432ux43eux440ux438ux442ux435ux43bux44cux43dux43e-9}

Нужно уметь сказать:

\begin{enumerate}
\def\labelenumi{\arabic{enumi}.}
\tightlist
\item
  Абсолютная непрерывность:
\end{enumerate}

\[
\nu\ll\mu
\quad\Longleftrightarrow\quad
\mu(A)=0\Rightarrow\nu(A)=0.
\]

\begin{enumerate}
\def\labelenumi{\arabic{enumi}.}
\setcounter{enumi}{1}
\tightlist
\item
  Радон-Никодим:
\end{enumerate}

\[
\nu(A)=\int_A f\,d\mu,
\qquad
f=\frac{d\nu}{d\mu}.
\]

\begin{enumerate}
\def\labelenumi{\arabic{enumi}.}
\setcounter{enumi}{2}
\tightlist
\item
  Плотность распределения:
\end{enumerate}

\[
P_X(B)=\int_Bp_X(x)\,dx.
\]

\begin{enumerate}
\def\labelenumi{\arabic{enumi}.}
\setcounter{enumi}{3}
\tightlist
\item
  Условное ожидание:
\end{enumerate}

\[
Y=\mathbb E[X\mid\mathcal G]
\]

если \(Y\) \(\mathcal G\)-измерима и

\[
\int_GY\,dP=\int_GX\,dP
\]

для всех \(G\in\mathcal G\).

\begin{enumerate}
\def\labelenumi{\arabic{enumi}.}
\setcounter{enumi}{4}
\tightlist
\item
  Условная вероятность:
\end{enumerate}

\[
P(A\mid\mathcal G)=\mathbb E[\mathbf 1_A\mid\mathcal G].
\]

\subsection{12. Минимум на
хорошо}\label{ux43cux438ux43dux438ux43cux443ux43c-ux43dux430-ux445ux43eux440ux43eux448ux43e-9}

Помимо минимума, нужно:

\begin{enumerate}
\def\labelenumi{\arabic{enumi}.}
\tightlist
\item
  объяснить, что плотность существует не всегда;
\item
  сказать, что плотность и условное ожидание определены почти всюду;
\item
  вывести условное ожидание через Радона-Никодима:
\end{enumerate}

\[
Q(G)=\int_GX\,dP,
\qquad
\mathbb E[X\mid\mathcal G]=\frac{dQ}{dP|_{\mathcal G}};
\]

\begin{enumerate}
\def\labelenumi{\arabic{enumi}.}
\setcounter{enumi}{3}
\tightlist
\item
  назвать свойства:
\end{enumerate}

\[
\mathbb E\,\mathbb E[X\mid\mathcal G]=\mathbb EX,
\]

\[
\mathbb E[X\mid\mathcal G]=X
\quad\text{если }X\text{ }\mathcal G\text{-измерима},
\]

\[
\mathbb E[X\mid\mathcal G]=\mathbb EX
\quad\text{если }X\perp\mathcal G;
\]

\begin{enumerate}
\def\labelenumi{\arabic{enumi}.}
\setcounter{enumi}{4}
\tightlist
\item
  привести пример условного ожидания по разбиению.
\end{enumerate}

\subsection{13. Что добавить для
отлично}\label{ux447ux442ux43e-ux434ux43eux431ux430ux432ux438ux442ux44c-ux434ux43bux44f-ux43eux442ux43bux438ux447ux43dux43e-9}

Для сильного ответа стоит добавить:

\begin{enumerate}
\def\labelenumi{\arabic{enumi}.}
\tightlist
\item
  формулировку Радона-Никодима для мер со знаком;
\item
  доказательство существования условного ожидания через меру
  \(Q(G)=\int_GX\,dP\);
\item
  доказательство единственности почти наверное;
\item
  башенное свойство;
\item
  вынос \(\mathcal G\)-измеримого множителя;
\item
  условную плотность
\end{enumerate}

\[
p_{X\mid Y}(x\mid y)=\frac{p_{X,Y}(x,y)}{p_Y(y)};
\]

\begin{enumerate}
\def\labelenumi{\arabic{enumi}.}
\setcounter{enumi}{6}
\tightlist
\item
  предупреждение про непрерывное условие \(Y=y\), где \(P(Y=y)=0\).
\end{enumerate}

\subsection{14. Устная версия ответа на 5
минут}\label{ux443ux441ux442ux43dux430ux44f-ux432ux435ux440ux441ux438ux44f-ux43eux442ux432ux435ux442ux430-ux43dux430-5-ux43cux438ux43dux443ux442-9}

Сначала формулирую абсолютную непрерывность мер:

\[
\nu\ll\mu
\quad\Longleftrightarrow\quad
\mu(A)=0\Rightarrow\nu(A)=0.
\]

Затем формулирую теорему Радона-Никодима: если \(\mu\)
\(\sigma\)-конечна, \(\nu\ll\mu\), то существует измеримая функция
\(f\), такая что

\[
\nu(A)=\int_A f\,d\mu.
\]

Эта функция называется производной Радона-Никодима:

\[
f=\frac{d\nu}{d\mu},
\]

и она единственна \(\mu\)-почти всюду.

Дальше объясняю плотности. Если \(X\) - случайный вектор в
\(\mathbb R^n\) и

\[
P_X\ll\lambda_n,
\]

то

\[
p_X=\frac{dP_X}{d\lambda_n}
\]

и

\[
P\{X\in B\}=\int_Bp_X(x)\,dx.
\]

Если абсолютной непрерывности нет, плотности относительно Лебега нет.
Например, у константы \(X=a\) распределение \(\delta_a\), и оно не имеет
лебеговой плотности.

Потом перехожу к условному ожиданию. Пусть \(X\in L^1\),
\(\mathcal G\subset\mathcal F\). Случайная величина \(Y\) называется
\(\mathbb E[X\mid\mathcal G]\), если \(Y\) \(\mathcal G\)-измерима и для
любого \(G\in\mathcal G\)

\[
\int_GY\,dP=\int_GX\,dP.
\]

Существование строится через Радона-Никодима: задаем на \(\mathcal G\)

\[
Q(G)=\int_GX\,dP.
\]

Тогда \(Q\ll P|_{\mathcal G}\), поэтому

\[
\mathbb E[X\mid\mathcal G]
=
\frac{dQ}{dP|_{\mathcal G}}.
\]

Условная вероятность - частный случай:

\[
P(A\mid\mathcal G)=\mathbb E[\mathbf 1_A\mid\mathcal G].
\]

Называю свойства:

\[
\mathbb E\,\mathbb E[X\mid\mathcal G]=\mathbb EX,
\]

если \(X\) \(\mathcal G\)-измерима, то

\[
\mathbb E[X\mid\mathcal G]=X,
\]

если \(X\) независима от \(\mathcal G\), то

\[
\mathbb E[X\mid\mathcal G]=\mathbb EX,
\]

и башня:

\[
\mathcal H\subset\mathcal G
\Rightarrow
\mathbb E[\mathbb E[X\mid\mathcal G]\mid\mathcal H]
=
\mathbb E[X\mid\mathcal H].
\]

В конце даю вычислительный пример: если есть совместная плотность
\(p_{X,Y}\), то

\[
p_{X\mid Y}(x\mid y)=\frac{p_{X,Y}(x,y)}{p_Y(y)}
\]

и

\[
\mathbb E[X\mid Y=y]=\int x p_{X\mid Y}(x\mid y)\,dx.
\]

Главные ловушки: плотность существует не всегда; условное ожидание -
случайная величина, а не число; все версии определены почти наверное.

\subsection{15. Если осталось 2 минуты перед
экзаменом}\label{ux435ux441ux43bux438-ux43eux441ux442ux430ux43bux43eux441ux44c-2-ux43cux438ux43dux443ux442ux44b-ux43fux435ux440ux435ux434-ux44dux43aux437ux430ux43cux435ux43dux43eux43c-9}

Запомнить четыре формулы.

\textbf{Радон-Никодим:}

\[
\nu\ll\mu
\Rightarrow
\nu(A)=\int_A\frac{d\nu}{d\mu}\,d\mu.
\]

\textbf{Плотность распределения:}

\[
P_X\ll\lambda_n,
\qquad
p_X=\frac{dP_X}{d\lambda_n},
\qquad
P_X(B)=\int_Bp_X\,dx.
\]

\textbf{Условное ожидание:}

\[
Y=\mathbb E[X\mid\mathcal G]
\iff
Y\text{ }\mathcal G\text{-измерима и }
\int_GY\,dP=\int_GX\,dP.
\]

\textbf{Построение через Радона-Никодима:}

\[
Q(G)=\int_GX\,dP,
\qquad
\mathbb E[X\mid\mathcal G]=\frac{dQ}{dP|_{\mathcal G}}.
\]

Главная ловушка: условное ожидание и плотность определяются почти всюду;
поточечные значения на нулевых множествах не важны.

\newpage

\section{Билет 11. Замена переменной в интеграле Лебега. Моменты через
интегралы по пространству и по
распределению}\label{ux431ux438ux43bux435ux442-11.-ux437ux430ux43cux435ux43dux430-ux43fux435ux440ux435ux43cux435ux43dux43dux43eux439-ux432-ux438ux43dux442ux435ux433ux440ux430ux43bux435-ux43bux435ux431ux435ux433ux430.-ux43cux43eux43cux435ux43dux442ux44b-ux447ux435ux440ux435ux437-ux438ux43dux442ux435ux433ux440ux430ux43bux44b-ux43fux43e-ux43fux440ux43eux441ux442ux440ux430ux43dux441ux442ux432ux443-ux438-ux43fux43e-ux440ux430ux441ux43fux440ux435ux434ux435ux43bux435ux43dux438ux44e}

Источники: Ширяев 1, гл. II, §6; лекции ДСП.

\subsection{1. Короткий
ответ}\label{ux43aux43eux440ux43eux442ux43aux438ux439-ux43eux442ux432ux435ux442-10}

Главная идея билета: интеграл функции от случайного элемента можно
считать не по исходному вероятностному пространству, а по распределению
этого элемента.

Если

\[
X:(\Omega,\mathcal F,P)\to(E,\mathcal E)
\]

\begin{itemize}
\tightlist
\item
  случайный элемент, а \(P_X\) - его распределение:
\end{itemize}

\[
P_X(B)=P\{X\in B\},
\qquad B\in\mathcal E,
\]

то для измеримой функции \(g:E\to\mathbb R\) или \(\mathbb C\):

\[
\mathbb E g(X)
=
\int_\Omega g(X(\omega))\,P(d\omega)
=
\int_E g(x)\,P_X(dx),
\]

если \(g\ge0\) или если \(g(X)\) интегрируема.

Если \(X\) имеет плотность \(p_X\) относительно меры Лебега на
\(\mathbb R^n\), то

\[
\mathbb E g(X)=\int_{\mathbb R^n} g(x)p_X(x)\,dx.
\]

Отсюда моменты считаются через распределение:

\[
\mathbb EX^k=\int_\mathbb R x^k\,P_X(dx),
\qquad
\mathbb E|X|^k=\int_\mathbb R |x|^k\,P_X(dx),
\]

а при плотности:

\[
\mathbb EX^k=\int_\mathbb R x^k p_X(x)\,dx.
\]

Для случайного вектора \(X\in\mathbb R^n\):

\[
\mathbb EX=\int_{\mathbb R^n}x\,P_X(dx),
\qquad
\operatorname{Cov}(X)=\int_{\mathbb R^n}(x-m)(x-m)^T\,P_X(dx),
\]

где \(m=\mathbb EX\).

\subsection{2. Полный
ответ}\label{ux43fux43eux43bux43dux44bux439-ux43eux442ux432ux435ux442-10}

В предыдущих билетах были введены распределения, плотности и интеграл
Лебега. Теперь нужно связать математическое ожидание на вероятностном
пространстве с интегралом по распределению. Это и есть вероятностная
форма замены переменной.

\subsubsection{2.1. Вероятностная форма замены
переменной}\label{ux432ux435ux440ux43eux44fux442ux43dux43eux441ux442ux43dux430ux44f-ux444ux43eux440ux43cux430-ux437ux430ux43cux435ux43dux44b-ux43fux435ux440ux435ux43cux435ux43dux43dux43eux439}

Пусть

\[
(\Omega,\mathcal F,P)
\]

\begin{itemize}
\tightlist
\item
  вероятностное пространство, \((E,\mathcal E)\) - измеримое
  пространство, а
\end{itemize}

\[
X:\Omega\to E
\]

\begin{itemize}
\tightlist
\item
  случайный элемент, то есть \(\mathcal F/\mathcal E\)-измеримое
  отображение. Его распределение, или образ меры \(P\) при отображении
  \(X\), задается формулой
\end{itemize}

\[
P_X(B)=P(X^{-1}(B))=P\{\omega:X(\omega)\in B\},
\qquad B\in\mathcal E.
\]

Распределение \(P_X\) - это вероятностная мера на пространстве значений
\(E\). Оно содержит всю информацию, необходимую для вычисления
вероятностей событий вида \(\{X\in B\}\) и средних величин вида
\(g(X)\).

Теорема о замене переменной в интеграле Лебега для образа меры
формулируется так. Пусть \(g:E\to\mathbb R\) или \(\mathbb C\) -
\(\mathcal E\)-измеримая функция. Тогда

\[
\int_\Omega g(X(\omega))\,P(d\omega)
=
\int_E g(x)\,P_X(dx),
\]

если \(g\ge0\); в этом случае обе части понимаются как расширенные
неотрицательные интегралы. Если \(g\) знакопеременная или
комплекснозначная, то равенство верно при интегрируемости:

\[
\int_E |g(x)|\,P_X(dx)<\infty,
\]

что равносильно

\[
\mathbb E|g(X)|<\infty.
\]

В вероятностной записи:

\[
\mathbb E g(X)=\int_E g(x)\,P_X(dx).
\]

Это не ``формальная замена буквы \(\omega\) на \(x\)''. Смысл в том, что
распределение \(P_X\) уже учитывает, с какими вероятностями значения
\(X\) попадают в разные множества.

Для вещественной случайной величины \(X\) можно писать

\[
\mathbb E g(X)=\int_\mathbb R g(x)\,P_X(dx).
\]

Если \(F_X\) - функция распределения, то \(P_X\) является мерой
Лебега-Стилтьеса, порожденной \(F_X\), и часто пишут

\[
\mathbb E g(X)=\int_\mathbb R g(x)\,dF_X(x).
\]

Если \(P_X\) абсолютно непрерывно относительно меры Лебега и имеет
плотность \(p_X\), то по Билету 10

\[
P_X(dx)=p_X(x)\,dx,
\]

и формула принимает вид

\[
\mathbb E g(X)=\int_\mathbb R g(x)p_X(x)\,dx.
\]

\subsubsection{2.2. Случайные векторы и
процессы}\label{ux441ux43bux443ux447ux430ux439ux43dux44bux435-ux432ux435ux43aux442ux43eux440ux44b-ux438-ux43fux440ux43eux446ux435ux441ux441ux44b}

Для случайного вектора \(X=(X_1,\ldots,X_n)\):

\[
\mathbb E g(X)
=
\int_{\mathbb R^n}g(x_1,\ldots,x_n)\,P_X(dx_1,\ldots,dx_n).
\]

Если есть плотность \(p_X\) на \(\mathbb R^n\), то

\[
\mathbb E g(X)
=
\int_{\mathbb R^n}g(x)p_X(x)\,dx.
\]

Это особенно важно для конечномерных распределений случайного процесса.
Пусть \((X_t)_{t\in T}\) - процесс и выбраны моменты времени

\[
t_1,\ldots,t_n.
\]

Тогда вектор

\[
X^{(t)}=(X_{t_1},\ldots,X_{t_n})
\]

имеет конечномерное распределение \(P_{t_1,\ldots,t_n}\). Для любой
измеримой функции \(g:\mathbb R^n\to\mathbb R\):

\[
\mathbb E g(X_{t_1},\ldots,X_{t_n})
=
\int_{\mathbb R^n}g(x_1,\ldots,x_n)\,P_{t_1,\ldots,t_n}(dx).
\]

Если есть плотность конечномерного распределения \(p_{t_1,\ldots,t_n}\),
то

\[
\mathbb E g(X_{t_1},\ldots,X_{t_n})
=
\int_{\mathbb R^n}g(x_1,\ldots,x_n)
p_{t_1,\ldots,t_n}(x_1,\ldots,x_n)\,dx_1\cdots dx_n.
\]

\subsubsection{2.3. Замена переменной и
диффеоморфизмы}\label{ux437ux430ux43cux435ux43dux430-ux43fux435ux440ux435ux43cux435ux43dux43dux43eux439-ux438-ux434ux438ux444ux444ux435ux43eux43cux43eux440ux444ux438ux437ux43cux44b}

Теперь классическая замена переменной в интеграле Лебега. Пусть
\(U,V\subset\mathbb R^n\) - открытые множества, а

\[
\Phi:U\to V
\]

\begin{itemize}
\tightlist
\item
  взаимно однозначное непрерывно дифференцируемое отображение с
  непрерывно дифференцируемым обратным, то есть \(C^1\)-диффеоморфизм.
  Тогда для неотрицательной или интегрируемой функции \(h\):
\end{itemize}

\[
\int_V h(y)\,dy
=
\int_U h(\Phi(x))\,|\det D\Phi(x)|\,dx.
\]

Эквивалентно, если \(y=\Phi(x)\), то элемент объема меняется по правилу

\[
dy=|\det D\Phi(x)|\,dx.
\]

Эта формула нужна, например, при переходе от декартовых координат к
полярным, при линейных преобразованиях векторов и при нахождении
плотности преобразованной случайной величины.

Если \(Y=\Phi(X)\), где \(X\) имеет плотность \(p_X\) на \(U\), а
\(\Phi:U\to V\) - диффеоморфизм, то плотность \(Y\) равна

\[
p_Y(y)
=
p_X(\Phi^{-1}(y))
\left|\det D\Phi^{-1}(y)\right|.
\]

В одномерном случае, если \(Y=\phi(X)\) и \(\phi\) строго монотонна,
дифференцируема и \(\phi'(x)\ne0\), то

\[
p_Y(y)=p_X(\phi^{-1}(y))
\left|(\phi^{-1})'(y)\right|.
\]

Если отображение не взаимно однозначно, надо суммировать по всем
прообразам:

\[
p_Y(y)=\sum_{x:\phi(x)=y}\frac{p_X(x)}{|\phi'(x)|},
\]

когда такая формула корректна.

\subsubsection{2.4. Моменты и числовые
характеристики}\label{ux43cux43eux43cux435ux43dux442ux44b-ux438-ux447ux438ux441ux43bux43eux432ux44bux435-ux445ux430ux440ux430ux43aux442ux435ux440ux438ux441ux442ux438ux43aux438}

Теперь моменты. Момент порядка \(k\) вещественной случайной величины
\(X\) - это математическое ожидание

\[
\mathbb EX^k,
\]

если оно существует. Абсолютный момент порядка \(k\):

\[
\mathbb E|X|^k.
\]

Существование абсолютного момента означает

\[
\mathbb E|X|^k<\infty.
\]

Через распределение:

\[
\mathbb E|X|^k
=
\int_\mathbb R |x|^k\,P_X(dx),
\]

а если этот интеграл конечен, то момент

\[
\mathbb EX^k
=
\int_\mathbb R x^k\,P_X(dx)
\]

корректно определен. При плотности:

\[
\mathbb E|X|^k=\int_\mathbb R |x|^k p_X(x)\,dx,
\qquad
\mathbb EX^k=\int_\mathbb R x^k p_X(x)\,dx.
\]

Математическое ожидание:

\[
\mathbb EX=\int_\Omega X(\omega)\,P(d\omega)
=
\int_\mathbb R x\,P_X(dx).
\]

При плотности:

\[
\mathbb EX=\int_\mathbb R x p_X(x)\,dx.
\]

Для дискретной случайной величины, принимающей значения \(x_i\) с
вероятностями \(p_i\):

\[
\mathbb E g(X)=\sum_i g(x_i)p_i.
\]

Значит

\[
\mathbb EX=\sum_i x_i p_i,
\qquad
\mathbb EX^k=\sum_i x_i^k p_i.
\]

Дисперсия:

\[
\operatorname{Var}X
=
\mathbb E[(X-\mathbb EX)^2].
\]

Через распределение:

\[
\operatorname{Var}X
=
\int_\mathbb R (x-m)^2\,P_X(dx),
\qquad
m=\mathbb EX.
\]

Эквивалентная вычислительная формула:

\[
\operatorname{Var}X=\mathbb EX^2-(\mathbb EX)^2,
\]

если \(\mathbb EX^2<\infty\). При плотности:

\[
\operatorname{Var}X
=
\int_\mathbb R (x-m)^2p_X(x)\,dx
=
\int_\mathbb R x^2p_X(x)\,dx
-
\left(\int_\mathbb R xp_X(x)\,dx\right)^2.
\]

\subsubsection{2.5. Ковариация и характеристики случайных
векторов}\label{ux43aux43eux432ux430ux440ux438ux430ux446ux438ux44f-ux438-ux445ux430ux440ux430ux43aux442ux435ux440ux438ux441ux442ux438ux43aux438-ux441ux43bux443ux447ux430ux439ux43dux44bux445-ux432ux435ux43aux442ux43eux440ux43eux432}

Для двух случайных величин \(X,Y\) с конечными вторыми моментами
ковариация:

\[
\operatorname{Cov}(X,Y)
=
\mathbb E[(X-\mathbb EX)(Y-\mathbb EY)].
\]

Если \((X,Y)\) имеет совместное распределение \(P_{X,Y}\), то

\[
\operatorname{Cov}(X,Y)
=
\int_{\mathbb R^2}(x-m_X)(y-m_Y)\,P_{X,Y}(dx,dy).
\]

При совместной плотности \(p_{X,Y}\):

\[
\operatorname{Cov}(X,Y)
=
\int_{\mathbb R^2}(x-m_X)(y-m_Y)p_{X,Y}(x,y)\,dx\,dy.
\]

Также

\[
\operatorname{Cov}(X,Y)=\mathbb E[XY]-\mathbb EX\,\mathbb EY,
\]

если \(\mathbb E X^2<\infty\) и \(\mathbb E Y^2<\infty\).

Для случайного вектора \(X=(X_1,\ldots,X_n)^T\) с конечным вторым
моментом:

\[
m=\mathbb EX
=
\int_{\mathbb R^n}x\,P_X(dx),
\]

а матрица вторых моментов:

\[
M=\mathbb E[XX^T]
=
\int_{\mathbb R^n}xx^T\,P_X(dx).
\]

Ковариационная матрица:

\[
\Sigma
=
\operatorname{Cov}(X)
=
\mathbb E[(X-m)(X-m)^T]
=
\int_{\mathbb R^n}(x-m)(x-m)^T\,P_X(dx).
\]

Эквивалентно:

\[
\Sigma=M-mm^T.
\]

При плотности \(p_X\) все интегралы по \(P_X\) заменяются на интегралы
по мере Лебега:

\[
m=\int_{\mathbb R^n}x p_X(x)\,dx,
\]

\[
M=\int_{\mathbb R^n}xx^T p_X(x)\,dx,
\]

\[
\Sigma=\int_{\mathbb R^n}(x-m)(x-m)^T p_X(x)\,dx.
\]

Элементы матрицы:

\[
\Sigma_{ij}
=
\operatorname{Cov}(X_i,X_j)
=
\int_{\mathbb R^n}(x_i-m_i)(x_j-m_j)\,P_X(dx).
\]

Если есть плотность:

\[
\Sigma_{ij}
=
\int_{\mathbb R^n}(x_i-m_i)(x_j-m_j)p_X(x)\,dx.
\]

Ковариационная матрица всегда симметрична и неотрицательно определена:

\[
a^T\Sigma a
=
\operatorname{Var}(a^TX)
\ge0
\]

для любого \(a\in\mathbb R^n\). Это связывает билет с Билетом 14 и
Билетом 18.

Итак, основная логика билета: случайная величина живет на \(\Omega\), но
ее распределение живет на пространстве значений. Теорема о замене
переменной позволяет считать все средние функции от случайной величины
через распределение. Если есть плотность, интеграл по распределению
превращается в обычный лебегов интеграл с плотностью.

\subsection{3. Основные
определения}\label{ux43eux441ux43dux43eux432ux43dux44bux435-ux43eux43fux440ux435ux434ux435ux43bux435ux43dux438ux44f-10}

\subsubsection{Образ
меры}\label{ux43eux431ux440ux430ux437-ux43cux435ux440ux44b}

Если \(X:(\Omega,\mathcal F)\to(E,\mathcal E)\) измеримо, а \(P\) - мера
на \(\Omega\), то образ меры \(P\) при \(X\) задается формулой

\[
P_X(B)=P(X^{-1}(B)),
\qquad B\in\mathcal E.
\]

В вероятности \(P_X\) называется распределением случайного элемента
\(X\).

\subsubsection{Замена переменной через образ
меры}\label{ux437ux430ux43cux435ux43dux430-ux43fux435ux440ux435ux43cux435ux43dux43dux43eux439-ux447ux435ux440ux435ux437-ux43eux431ux440ux430ux437-ux43cux435ux440ux44b}

Для измеримой \(g\):

\[
\int_\Omega g(X(\omega))\,P(d\omega)
=
\int_E g(x)\,P_X(dx).
\]

\subsubsection{Интеграл Лебега-Стилтьеса по функции
распределения}\label{ux438ux43dux442ux435ux433ux440ux430ux43b-ux43bux435ux431ux435ux433ux430-ux441ux442ux438ux43bux442ux44cux435ux441ux430-ux43fux43e-ux444ux443ux43dux43aux446ux438ux438-ux440ux430ux441ux43fux440ux435ux434ux435ux43bux435ux43dux438ux44f}

Если \(X\) вещественная, \(F_X\) - ее функция распределения, а \(P_X\) -
соответствующая мера, то

\[
\int_\mathbb R g(x)\,P_X(dx)
\]

часто записывают как

\[
\int_\mathbb R g(x)\,dF_X(x).
\]

\subsubsection{Плотность
распределения}\label{ux43fux43bux43eux442ux43dux43eux441ux442ux44c-ux440ux430ux441ux43fux440ux435ux434ux435ux43bux435ux43dux438ux44f-2}

Если \(P_X\ll\lambda_n\), то

\[
p_X=\frac{dP_X}{d\lambda_n}
\]

и

\[
P_X(B)=\int_B p_X(x)\,dx.
\]

\subsubsection{Классическая замена переменной с
якобианом}\label{ux43aux43bux430ux441ux441ux438ux447ux435ux441ux43aux430ux44f-ux437ux430ux43cux435ux43dux430-ux43fux435ux440ux435ux43cux435ux43dux43dux43eux439-ux441-ux44fux43aux43eux431ux438ux430ux43dux43eux43c}

Если \(\Phi:U\to V\) - \(C^1\)-диффеоморфизм, то

\[
\int_V h(y)\,dy
=
\int_U h(\Phi(x))|\det D\Phi(x)|\,dx.
\]

\subsubsection{\texorpdfstring{Момент порядка
\(k\)}{Момент порядка k}}\label{ux43cux43eux43cux435ux43dux442-ux43fux43eux440ux44fux434ux43aux430-k}

Для вещественной случайной величины:

\[
\mathbb EX^k,
\]

если математическое ожидание существует.

\subsubsection{\texorpdfstring{Абсолютный момент порядка
\(k\)}{Абсолютный момент порядка k}}\label{ux430ux431ux441ux43eux43bux44eux442ux43dux44bux439-ux43cux43eux43cux435ux43dux442-ux43fux43eux440ux44fux434ux43aux430-k}

\[
\mathbb E|X|^k.
\]

Если \(\mathbb E|X|^k<\infty\), то момент \(\mathbb EX^k\) существует.

\subsubsection{Дисперсия}\label{ux434ux438ux441ux43fux435ux440ux441ux438ux44f}

\[
\operatorname{Var}X=\mathbb E[(X-\mathbb EX)^2],
\]

если \(\mathbb EX^2<\infty\).

\subsubsection{Ковариация}\label{ux43aux43eux432ux430ux440ux438ux430ux446ux438ux44f}

\[
\operatorname{Cov}(X,Y)
=
\mathbb E[(X-\mathbb EX)(Y-\mathbb EY)],
\]

если \(X,Y\in L^2\).

\subsubsection{Матрица вторых
моментов}\label{ux43cux430ux442ux440ux438ux446ux430-ux432ux442ux43eux440ux44bux445-ux43cux43eux43cux435ux43dux442ux43eux432}

\[
M=\mathbb E[XX^T].
\]

\subsubsection{Ковариационная
матрица}\label{ux43aux43eux432ux430ux440ux438ux430ux446ux438ux43eux43dux43dux430ux44f-ux43cux430ux442ux440ux438ux446ux430}

\[
\Sigma=\mathbb E[(X-\mathbb EX)(X-\mathbb EX)^T].
\]

\subsection{4. Основные теоремы и
свойства}\label{ux43eux441ux43dux43eux432ux43dux44bux435-ux442ux435ux43eux440ux435ux43cux44b-ux438-ux441ux432ux43eux439ux441ux442ux432ux430-10}

\subsubsection{Теорема 1. Замена переменной для образа
меры}\label{ux442ux435ux43eux440ux435ux43cux430-1.-ux437ux430ux43cux435ux43dux430-ux43fux435ux440ux435ux43cux435ux43dux43dux43eux439-ux434ux43bux44f-ux43eux431ux440ux430ux437ux430-ux43cux435ux440ux44b}

Пусть \(X:(\Omega,\mathcal F,P)\to(E,\mathcal E)\) - случайный элемент,
\(P_X\) - его распределение, \(g\) - измеримая функция. Тогда

\[
\int_\Omega g(X(\omega))\,P(d\omega)
=
\int_E g(x)\,P_X(dx),
\]

если \(g\ge0\) или если один из интегралов от \(|g|\) конечен.

\subsubsection{Теорема 2. Доказательная схема замены
переменной}\label{ux442ux435ux43eux440ux435ux43cux430-2.-ux434ux43eux43aux430ux437ux430ux442ux435ux43bux44cux43dux430ux44f-ux441ux445ux435ux43cux430-ux437ux430ux43cux435ux43dux44b-ux43fux435ux440ux435ux43cux435ux43dux43dux43eux439}

Сначала формула проверяется для индикатора:

\[
g=\mathbf 1_B.
\]

Тогда

\[
\int_\Omega \mathbf 1_B(X(\omega))\,dP
=
P\{X\in B\}
=
P_X(B)
=
\int_E\mathbf 1_B(x)\,P_X(dx).
\]

Затем формула переносится на простые функции по линейности, на
неотрицательные измеримые функции по теореме Беппо Леви, а на
интегрируемые знакопеременные функции через положительную и
отрицательную части.

\subsubsection{Теорема 3. Случай с
плотностью}\label{ux442ux435ux43eux440ux435ux43cux430-3.-ux441ux43bux443ux447ux430ux439-ux441-ux43fux43bux43eux442ux43dux43eux441ux442ux44cux44e}

Если \(X\) имеет плотность \(p_X\) относительно \(\lambda_n\), то

\[
\mathbb E g(X)=\int_{\mathbb R^n}g(x)p_X(x)\,dx,
\]

для \(g\ge0\) или интегрируемой \(g(X)\).

\subsubsection{\texorpdfstring{Теорема 4. Классическая замена переменной
в
\(\mathbb R^n\)}{Теорема 4. Классическая замена переменной в \textbackslash mathbb R\^{}n}}\label{ux442ux435ux43eux440ux435ux43cux430-4.-ux43aux43bux430ux441ux441ux438ux447ux435ux441ux43aux430ux44f-ux437ux430ux43cux435ux43dux430-ux43fux435ux440ux435ux43cux435ux43dux43dux43eux439-ux432-mathbb-rn}

Если \(\Phi:U\to V\) - \(C^1\)-диффеоморфизм, то

\[
\int_V h(y)\,dy
=
\int_U h(\Phi(x))|\det D\Phi(x)|\,dx.
\]

\subsubsection{Следствие 5. Плотность
преобразования}\label{ux441ux43bux435ux434ux441ux442ux432ux438ux435-5.-ux43fux43bux43eux442ux43dux43eux441ux442ux44c-ux43fux440ux435ux43eux431ux440ux430ux437ux43eux432ux430ux43dux438ux44f}

Если \(Y=\Phi(X)\), \(X\) имеет плотность \(p_X\), а \(\Phi\) -
диффеоморфизм, то

\[
p_Y(y)=p_X(\Phi^{-1}(y))
\left|\det D\Phi^{-1}(y)\right|.
\]

\subsubsection{Свойство 6. Моменты через
распределение}\label{ux441ux432ux43eux439ux441ux442ux432ux43e-6.-ux43cux43eux43cux435ux43dux442ux44b-ux447ux435ux440ux435ux437-ux440ux430ux441ux43fux440ux435ux434ux435ux43bux435ux43dux438ux435}

Для вещественной \(X\):

\[
\mathbb E|X|^k=\int_\mathbb R |x|^k\,P_X(dx),
\]

\[
\mathbb EX^k=\int_\mathbb R x^k\,P_X(dx),
\]

если соответствующие интегралы существуют.

\subsubsection{Свойство 7. Моменты при
плотности}\label{ux441ux432ux43eux439ux441ux442ux432ux43e-7.-ux43cux43eux43cux435ux43dux442ux44b-ux43fux440ux438-ux43fux43bux43eux442ux43dux43eux441ux442ux438}

Если \(X\) имеет плотность \(p_X\), то

\[
\mathbb E|X|^k=\int_\mathbb R |x|^k p_X(x)\,dx,
\]

\[
\mathbb EX^k=\int_\mathbb R x^k p_X(x)\,dx.
\]

\subsubsection{Свойство 8. Дисперсия через
распределение}\label{ux441ux432ux43eux439ux441ux442ux432ux43e-8.-ux434ux438ux441ux43fux435ux440ux441ux438ux44f-ux447ux435ux440ux435ux437-ux440ux430ux441ux43fux440ux435ux434ux435ux43bux435ux43dux438ux435}

Если \(m=\mathbb EX\) и \(\mathbb EX^2<\infty\), то

\[
\operatorname{Var}X=\int_\mathbb R (x-m)^2\,P_X(dx)
=
\mathbb EX^2-m^2.
\]

При плотности:

\[
\operatorname{Var}X=\int_\mathbb R (x-m)^2p_X(x)\,dx.
\]

\subsubsection{Свойство 9. Ковариация через совместное
распределение}\label{ux441ux432ux43eux439ux441ux442ux432ux43e-9.-ux43aux43eux432ux430ux440ux438ux430ux446ux438ux44f-ux447ux435ux440ux435ux437-ux441ux43eux432ux43cux435ux441ux442ux43dux43eux435-ux440ux430ux441ux43fux440ux435ux434ux435ux43bux435ux43dux438ux435}

Если \(X,Y\in L^2\), то

\[
\operatorname{Cov}(X,Y)
=
\int_{\mathbb R^2}(x-m_X)(y-m_Y)\,P_{X,Y}(dx,dy).
\]

При совместной плотности:

\[
\operatorname{Cov}(X,Y)
=
\int_{\mathbb R^2}(x-m_X)(y-m_Y)p_{X,Y}(x,y)\,dx\,dy.
\]

\subsubsection{Свойство 10. Ковариационная
матрица}\label{ux441ux432ux43eux439ux441ux442ux432ux43e-10.-ux43aux43eux432ux430ux440ux438ux430ux446ux438ux43eux43dux43dux430ux44f-ux43cux430ux442ux440ux438ux446ux430}

Для \(X\in\mathbb R^n\) с конечным вторым моментом:

\[
\Sigma=\int_{\mathbb R^n}(x-m)(x-m)^T\,P_X(dx).
\]

При плотности:

\[
\Sigma=\int_{\mathbb R^n}(x-m)(x-m)^T p_X(x)\,dx.
\]

\subsubsection{Свойство 11. Неотрицательная определенность
ковариационной
матрицы}\label{ux441ux432ux43eux439ux441ux442ux432ux43e-11.-ux43dux435ux43eux442ux440ux438ux446ux430ux442ux435ux43bux44cux43dux430ux44f-ux43eux43fux440ux435ux434ux435ux43bux435ux43dux43dux43eux441ux442ux44c-ux43aux43eux432ux430ux440ux438ux430ux446ux438ux43eux43dux43dux43eux439-ux43cux430ux442ux440ux438ux446ux44b}

Для любого \(a\in\mathbb R^n\):

\[
a^T\Sigma a=\operatorname{Var}(a^TX)\ge0.
\]

\subsection{5. Примеры}\label{ux43fux440ux438ux43cux435ux440ux44b-10}

\subsubsection{Пример 1. Дискретная случайная
величина}\label{ux43fux440ux438ux43cux435ux440-1.-ux434ux438ux441ux43aux440ux435ux442ux43dux430ux44f-ux441ux43bux443ux447ux430ux439ux43dux430ux44f-ux432ux435ux43bux438ux447ux438ux43dux430}

Пусть \(X\) принимает значения \(x_i\) с вероятностями \(p_i\). Тогда
распределение

\[
P_X=\sum_i p_i\delta_{x_i}.
\]

Поэтому

\[
\mathbb E g(X)=\int g\,dP_X=\sum_i g(x_i)p_i.
\]

В частности:

\[
\mathbb EX=\sum_i x_i p_i,
\qquad
\operatorname{Var}X=\sum_i(x_i-m)^2p_i.
\]

\subsubsection{Пример 2. Непрерывная случайная величина с
плотностью}\label{ux43fux440ux438ux43cux435ux440-2.-ux43dux435ux43fux440ux435ux440ux44bux432ux43dux430ux44f-ux441ux43bux443ux447ux430ux439ux43dux430ux44f-ux432ux435ux43bux438ux447ux438ux43dux430-ux441-ux43fux43bux43eux442ux43dux43eux441ux442ux44cux44e}

Пусть \(X\) имеет плотность

\[
p_X(x)=\mathbf 1_{[0,1]}(x).
\]

Тогда

\[
\mathbb EX=\int_0^1 x\,dx=\frac12,
\]

\[
\mathbb EX^2=\int_0^1 x^2\,dx=\frac13,
\]

\[
\operatorname{Var}X=\frac13-\frac14=\frac1{12}.
\]

\subsubsection{Пример 3. Вектор с
плотностью}\label{ux43fux440ux438ux43cux435ux440-3.-ux432ux435ux43aux442ux43eux440-ux441-ux43fux43bux43eux442ux43dux43eux441ux442ux44cux44e}

Пусть \((X,Y)\) имеет совместную плотность \(p(x,y)\). Тогда

\[
\mathbb E[XY]
=
\int_{\mathbb R^2}xy\,p(x,y)\,dx\,dy.
\]

Если \(m_X=\mathbb EX\), \(m_Y=\mathbb EY\), то

\[
\operatorname{Cov}(X,Y)
=
\int_{\mathbb R^2}(x-m_X)(y-m_Y)p(x,y)\,dx\,dy.
\]

\subsubsection{Пример 4. Полярные
координаты}\label{ux43fux440ux438ux43cux435ux440-4.-ux43fux43eux43bux44fux440ux43dux44bux435-ux43aux43eux43eux440ux434ux438ux43dux430ux442ux44b}

Для перехода

\[
x=r\cos\theta,\qquad y=r\sin\theta
\]

якобиан равен

\[
\left|\det\frac{\partial(x,y)}{\partial(r,\theta)}\right|=r.
\]

Поэтому

\[
\int_{\mathbb R^2}h(x,y)\,dx\,dy
=
\int_0^\infty\int_0^{2\pi}
h(r\cos\theta,r\sin\theta)\,r\,d\theta\,dr.
\]

\subsubsection{Пример 5. Линейное преобразование
плотности}\label{ux43fux440ux438ux43cux435ux440-5.-ux43bux438ux43dux435ux439ux43dux43eux435-ux43fux440ux435ux43eux431ux440ux430ux437ux43eux432ux430ux43dux438ux435-ux43fux43bux43eux442ux43dux43eux441ux442ux438}

Пусть \(X\in\mathbb R^n\) имеет плотность \(p_X\), а

\[
Y=AX+b,
\]

где \(A\) - невырожденная матрица. Тогда

\[
p_Y(y)
=
p_X(A^{-1}(y-b))|\det A^{-1}|
=
\frac{p_X(A^{-1}(y-b))}{|\det A|}.
\]

Для моментов:

\[
\mathbb EY=A\mathbb EX+b,
\]

\[
\operatorname{Cov}(Y)=A\operatorname{Cov}(X)A^T.
\]

\subsubsection{Пример 6. Момент может не
существовать}\label{ux43fux440ux438ux43cux435ux440-6.-ux43cux43eux43cux435ux43dux442-ux43cux43eux436ux435ux442-ux43dux435-ux441ux443ux449ux435ux441ux442ux432ux43eux432ux430ux442ux44c}

У распределения Коши с плотностью

\[
p(x)=\frac1{\pi(1+x^2)}
\]

математическое ожидание не существует как конечный лебегов интеграл,
потому что

\[
\int_\mathbb R |x|p(x)\,dx=\infty.
\]

Поэтому нельзя автоматически писать \(\mathbb EX=\int xp(x)\,dx\) как
конечное число.

\subsection{6. Схемы доказательств /
обоснований}\label{ux441ux445ux435ux43cux44b-ux434ux43eux43aux430ux437ux430ux442ux435ux43bux44cux441ux442ux432-ux43eux431ux43eux441ux43dux43eux432ux430ux43dux438ux439-3}

\textbf{Доказательство формулы \(\mathbb E g(X)=\int g\,dP_X\).}

\begin{enumerate}
\def\labelenumi{\arabic{enumi}.}
\tightlist
\item
  Для \(g=\mathbf 1_B\):
\end{enumerate}

\[
\mathbb E\mathbf 1_B(X)=P\{X\in B\}=P_X(B).
\]

\begin{enumerate}
\def\labelenumi{\arabic{enumi}.}
\setcounter{enumi}{1}
\tightlist
\item
  Для простой функции
\end{enumerate}

\[
g=\sum_{k=1}^n a_k\mathbf 1_{B_k}
\]

используем линейность интеграла.

\begin{enumerate}
\def\labelenumi{\arabic{enumi}.}
\setcounter{enumi}{2}
\item
  Для \(g\ge0\) берем простые функции \(g_n\uparrow g\) и применяем
  теорему Беппо Леви.
\item
  Для знакопеременной \(g\) пишем
\end{enumerate}

\[
g=g^+-g^-,
\]

и требуем интегрируемость, чтобы не получить неопределенность
\(\infty-\infty\).

\textbf{Обоснование формулы с плотностью.}

Если

\[
P_X(dx)=p_X(x)\,dx,
\]

то

\[
\int g\,dP_X
=
\int g(x)p_X(x)\,dx
\]

по определению плотности как производной Радона-Никодима.

\textbf{Обоснование дисперсии.}

Если \(m=\mathbb EX\), то

\[
\operatorname{Var}X
=
\mathbb E[(X-m)^2]
=
\mathbb E[X^2-2mX+m^2]
\]

\[
=
\mathbb EX^2-2m\mathbb EX+m^2
=
\mathbb EX^2-m^2.
\]

\textbf{Обоснование неотрицательной определенности ковариационной
матрицы.}

Для любого \(a\in\mathbb R^n\):

\[
a^T\Sigma a
=
\mathbb E[a^T(X-m)(X-m)^Ta]
=
\mathbb E[(a^T(X-m))^2]
\ge0.
\]

\textbf{Обоснование плотности преобразования.}

Если \(Y=\Phi(X)\), то для борелевского \(B\subset V\):

\[
P\{Y\in B\}
=
\int_{\Phi^{-1}(B)}p_X(x)\,dx.
\]

По классической замене переменной \(x=\Phi^{-1}(y)\):

\[
\int_{\Phi^{-1}(B)}p_X(x)\,dx
=
\int_B p_X(\Phi^{-1}(y))
\left|\det D\Phi^{-1}(y)\right|\,dy.
\]

Значит это и есть плотность \(Y\).

\subsection{7. Что написать на
доске}\label{ux447ux442ux43e-ux43dux430ux43fux438ux441ux430ux442ux44c-ux43dux430-ux434ux43eux441ux43aux435-10}

\begin{enumerate}
\def\labelenumi{\arabic{enumi}.}
\tightlist
\item
  Образ меры:
\end{enumerate}

\[
P_X(B)=P\{X\in B\}.
\]

\begin{enumerate}
\def\labelenumi{\arabic{enumi}.}
\setcounter{enumi}{1}
\tightlist
\item
  Замена переменной:
\end{enumerate}

\[
\mathbb E g(X)
=
\int_\Omega g(X(\omega))\,dP
=
\int_E g(x)\,P_X(dx).
\]

\begin{enumerate}
\def\labelenumi{\arabic{enumi}.}
\setcounter{enumi}{2}
\tightlist
\item
  При плотности:
\end{enumerate}

\[
\mathbb E g(X)
=
\int_{\mathbb R^n}g(x)p_X(x)\,dx.
\]

\begin{enumerate}
\def\labelenumi{\arabic{enumi}.}
\setcounter{enumi}{3}
\tightlist
\item
  Классическая замена:
\end{enumerate}

\[
\int_V h(y)\,dy
=
\int_U h(\Phi(x))|\det D\Phi(x)|\,dx.
\]

\begin{enumerate}
\def\labelenumi{\arabic{enumi}.}
\setcounter{enumi}{4}
\tightlist
\item
  Моменты:
\end{enumerate}

\[
\mathbb EX^k=\int x^k\,P_X(dx),
\qquad
\mathbb E|X|^k=\int |x|^k\,P_X(dx).
\]

\begin{enumerate}
\def\labelenumi{\arabic{enumi}.}
\setcounter{enumi}{5}
\tightlist
\item
  Дисперсия:
\end{enumerate}

\[
\operatorname{Var}X
=
\int (x-m)^2\,P_X(dx)
=
\mathbb EX^2-m^2.
\]

\begin{enumerate}
\def\labelenumi{\arabic{enumi}.}
\setcounter{enumi}{6}
\tightlist
\item
  Ковариационная матрица:
\end{enumerate}

\[
\Sigma=\int (x-m)(x-m)^T\,P_X(dx),
\qquad
a^T\Sigma a\ge0.
\]

\subsection{8. Типичные дополнительные
вопросы}\label{ux442ux438ux43fux438ux447ux43dux44bux435-ux434ux43eux43fux43eux43bux43dux438ux442ux435ux43bux44cux43dux44bux435-ux432ux43eux43fux440ux43eux441ux44b-10}

\textbf{Вопрос. Что такое образ меры?}

Это мера на пространстве значений:

\[
P_X(B)=P(X^{-1}(B)).
\]

В вероятности она называется распределением \(X\).

\textbf{Вопрос. В чем смысл формулы \(\mathbb E g(X)=\int g\,dP_X\)?}

Она говорит, что для вычисления среднего функции от \(X\) не нужно знать
все вероятностное пространство. Достаточно знать распределение \(X\).

\textbf{Вопрос. При каких условиях верна формула замены переменной?}

Для \(g\ge0\) она верна в расширенном смысле. Для знакопеременной
функции нужна интегрируемость:

\[
\int |g|\,dP_X<\infty.
\]

\textbf{Вопрос. Чем интеграл по распределению отличается от интеграла по
плотности?}

Интеграл по распределению

\[
\int g\,dP_X
\]

имеет смысл всегда. Интеграл

\[
\int g(x)p_X(x)\,dx
\]

требует существования плотности \(p_X\) относительно меры Лебега.

\textbf{Вопрос. Что делать, если плотности нет?}

Интегрировать по мере распределения \(P_X\) или использовать сумму в
дискретном случае. Нельзя искусственно писать лебегову плотность.

\textbf{Вопрос. Как посчитать момент порядка \(k\)?}

Через замену переменной с \(g(x)=x^k\):

\[
\mathbb EX^k=\int x^k\,P_X(dx),
\]

если интеграл корректно определен.

\textbf{Вопрос. Чем момент отличается от абсолютного момента?}

Абсолютный момент:

\[
\mathbb E|X|^k.
\]

Если он конечен, то обычный момент \(\mathbb EX^k\) точно существует.
Обратное в смысле условной сходимости для Лебега не используется:
математическое ожидание требует интегрируемости положительной и
отрицательной частей без неопределенности.

\textbf{Вопрос. Как выразить дисперсию через распределение?}

Если \(m=\mathbb EX\), то

\[
\operatorname{Var}X=\int (x-m)^2\,P_X(dx).
\]

При плотности:

\[
\operatorname{Var}X=\int (x-m)^2p_X(x)\,dx.
\]

\textbf{Вопрос. Как выразить ковариацию через совместное распределение?}

\[
\operatorname{Cov}(X,Y)
=
\int (x-m_X)(y-m_Y)\,P_{X,Y}(dx,dy).
\]

При совместной плотности заменить \(P_{X,Y}(dx,dy)\) на
\(p_{X,Y}(x,y)\,dx\,dy\).

\textbf{Вопрос. Почему ковариационная матрица неотрицательно
определена?}

Потому что

\[
a^T\Sigma a=\operatorname{Var}(a^TX)\ge0.
\]

\textbf{Вопрос. Как меняется плотность при диффеоморфизме?}

Если \(Y=\Phi(X)\), то

\[
p_Y(y)=p_X(\Phi^{-1}(y))
\left|\det D\Phi^{-1}(y)\right|.
\]

\textbf{Вопрос. Почему нельзя забывать якобиан?}

Потому что мера Лебега объема меняется при растяжении координат. Якобиан
учитывает локальное изменение объема.

\subsection{9. Типичные
ошибки}\label{ux442ux438ux43fux438ux447ux43dux44bux435-ux43eux448ux438ux431ux43aux438-10}

\begin{enumerate}
\def\labelenumi{\arabic{enumi}.}
\tightlist
\item
  Писать интеграл по плотности, когда плотности нет.
\item
  Путать \(P\) на \(\Omega\) и \(P_X\) на пространстве значений.
\item
  Забывать условия интегрируемости для моментов.
\item
  Считать, что \(\int xp(x)\,dx\) существует как число без проверки
  \(\int |x|p(x)\,dx<\infty\).
\item
  Путать момент \(\mathbb EX^2\) и дисперсию \(\operatorname{Var}X\).
\item
  Забывать центрирование в ковариации.
\item
  Считать ковариационную матрицу матрицей вторых моментов без вычитания
  \(mm^T\).
\item
  Забывать якобиан в классической замене переменной.
\item
  В формуле плотности преобразования использовать \(|\det D\Phi|\)
  вместо \(|\det D\Phi^{-1}|\) при записи через \(y\).
\item
  Интегрировать по маргинальному распределению вместо совместного при
  вычислении \(\mathbb E g(X,Y)\).
\end{enumerate}

\subsection{10. Связи с другими
билетами}\label{ux441ux432ux44fux437ux438-ux441-ux434ux440ux443ux433ux438ux43cux438-ux431ux438ux43bux435ux442ux430ux43cux438-10}

\textbf{Билет 06.} Распределения случайных величин, векторов и
конечномерные распределения - это меры, по которым интегрируют в этом
билете.

\textbf{Билет 07.} Все формулы основаны на интеграле Лебега и предельных
теоремах для перехода от простых функций к измеримым.

\textbf{Билет 08.} Характеристическая функция - это частный случай:

\[
\varphi_X(t)=\mathbb E e^{itX}
=
\int e^{itx}\,P_X(dx).
\]

\textbf{Билет 09.} Для совместных распределений и плотностей
используются Фубини и повторные интегралы.

\textbf{Билет 10.} Плотность \(p_X\) - это производная Радона-Никодима
\(dP_X/d\lambda\).

\textbf{Билет 13.} Моменты можно выражать через производные
характеристической функции, а здесь - через интегралы по распределению.

\textbf{Билет 14.} Ковариация и корреляция строятся из вторых моментов.

\textbf{Билет 18.} Гауссовский закон задается средним вектором и
ковариационной матрицей.

\textbf{Билет 24.} Корреляционная функция стационарной
последовательности выражается через моменты второго порядка.

\subsection{11. Минимум на
удовлетворительно}\label{ux43cux438ux43dux438ux43cux443ux43c-ux43dux430-ux443ux434ux43eux432ux43bux435ux442ux432ux43eux440ux438ux442ux435ux43bux44cux43dux43e-10}

Нужно уметь сказать:

\begin{enumerate}
\def\labelenumi{\arabic{enumi}.}
\tightlist
\item
  Распределение:
\end{enumerate}

\[
P_X(B)=P\{X\in B\}.
\]

\begin{enumerate}
\def\labelenumi{\arabic{enumi}.}
\setcounter{enumi}{1}
\tightlist
\item
  Замена переменной:
\end{enumerate}

\[
\mathbb E g(X)
=
\int_\Omega g(X)\,dP
=
\int_E g(x)\,P_X(dx).
\]

\begin{enumerate}
\def\labelenumi{\arabic{enumi}.}
\setcounter{enumi}{2}
\tightlist
\item
  При плотности:
\end{enumerate}

\[
\mathbb E g(X)=\int g(x)p_X(x)\,dx.
\]

\begin{enumerate}
\def\labelenumi{\arabic{enumi}.}
\setcounter{enumi}{3}
\tightlist
\item
  Момент:
\end{enumerate}

\[
\mathbb EX^k=\int x^k\,P_X(dx).
\]

\begin{enumerate}
\def\labelenumi{\arabic{enumi}.}
\setcounter{enumi}{4}
\tightlist
\item
  Дисперсия:
\end{enumerate}

\[
\operatorname{Var}X=\mathbb EX^2-(\mathbb EX)^2.
\]

\subsection{12. Минимум на
хорошо}\label{ux43cux438ux43dux438ux43cux443ux43c-ux43dux430-ux445ux43eux440ux43eux448ux43e-10}

Помимо минимума, нужно:

\begin{enumerate}
\def\labelenumi{\arabic{enumi}.}
\tightlist
\item
  объяснить, что формула верна для \(g\ge0\) или интегрируемой \(g\);
\item
  различать интеграл по \(P_X\) и интеграл по плотности;
\item
  записать ковариацию через совместное распределение;
\item
  записать матрицу
\end{enumerate}

\[
\Sigma=\int (x-m)(x-m)^T\,P_X(dx);
\]

\begin{enumerate}
\def\labelenumi{\arabic{enumi}.}
\setcounter{enumi}{4}
\tightlist
\item
  привести дискретный и непрерывный пример;
\item
  назвать классическую замену переменной с якобианом.
\end{enumerate}

\subsection{13. Что добавить для
отлично}\label{ux447ux442ux43e-ux434ux43eux431ux430ux432ux438ux442ux44c-ux434ux43bux44f-ux43eux442ux43bux438ux447ux43dux43e-10}

Для сильного ответа стоит добавить:

\begin{enumerate}
\def\labelenumi{\arabic{enumi}.}
\tightlist
\item
  доказательство через индикаторы, простые функции и теорему Беппо Леви;
\item
  интеграл Лебега-Стилтьеса
\end{enumerate}

\[
\int g\,dF_X;
\]

\begin{enumerate}
\def\labelenumi{\arabic{enumi}.}
\setcounter{enumi}{2}
\tightlist
\item
  формулу плотности преобразования:
\end{enumerate}

\[
p_Y(y)=p_X(\Phi^{-1}(y))|\det D\Phi^{-1}(y)|;
\]

\begin{enumerate}
\def\labelenumi{\arabic{enumi}.}
\setcounter{enumi}{3}
\tightlist
\item
  неотрицательную определенность ковариационной матрицы;
\item
  предупреждение про распределение Коши или другой пример, где момент не
  существует;
\item
  связь с характеристическими функциями и гауссовскими законами.
\end{enumerate}

\subsection{14. Устная версия ответа на 5
минут}\label{ux443ux441ux442ux43dux430ux44f-ux432ux435ux440ux441ux438ux44f-ux43eux442ux432ux435ux442ux430-ux43dux430-5-ux43cux438ux43dux443ux442-10}

Сначала я формулирую главную теорему. Пусть \(X:\Omega\to E\) -
случайный элемент, а \(P_X\) - его распределение:

\[
P_X(B)=P\{X\in B\}.
\]

Тогда для измеримой функции \(g\):

\[
\mathbb E g(X)
=
\int_\Omega g(X(\omega))\,dP
=
\int_E g(x)\,P_X(dx),
\]

если \(g\ge0\) или если \(g(X)\) интегрируема. Это называется заменой
переменной в интеграле Лебега через образ меры.

Дальше говорю: если \(X\) вещественная, можно писать

\[
\mathbb E g(X)=\int_\mathbb R g(x)\,dF_X(x).
\]

Если есть плотность \(p_X\), то

\[
\mathbb E g(X)=\int_\mathbb R g(x)p_X(x)\,dx.
\]

Важно: плотность существует не всегда. Если плотности нет, надо
интегрировать по \(P_X\), а не по \(p_Xdx\).

Затем формулирую классическую замену переменной. Если \(\Phi:U\to V\) -
\(C^1\)-диффеоморфизм, то

\[
\int_V h(y)\,dy
=
\int_U h(\Phi(x))|\det D\Phi(x)|\,dx.
\]

Отсюда для \(Y=\Phi(X)\):

\[
p_Y(y)=p_X(\Phi^{-1}(y))|\det D\Phi^{-1}(y)|.
\]

После этого перехожу к моментам. Для вещественной случайной величины:

\[
\mathbb EX^k=\int x^k\,P_X(dx),
\qquad
\mathbb E|X|^k=\int |x|^k\,P_X(dx).
\]

При плотности:

\[
\mathbb EX^k=\int x^kp_X(x)\,dx.
\]

Дисперсия:

\[
\operatorname{Var}X
=
\int (x-m)^2\,P_X(dx)
=
\mathbb EX^2-m^2.
\]

Для пары:

\[
\operatorname{Cov}(X,Y)
=
\int (x-m_X)(y-m_Y)\,P_{X,Y}(dx,dy).
\]

Для вектора:

\[
\Sigma
=
\int (x-m)(x-m)^T\,P_X(dx).
\]

И завершаю свойством:

\[
a^T\Sigma a=\operatorname{Var}(a^TX)\ge0.
\]

Главные ловушки: не писать плотность без абсолютной непрерывности, не
забывать существование моментов и якобиан при классической замене
переменной.

\subsection{15. Если осталось 2 минуты перед
экзаменом}\label{ux435ux441ux43bux438-ux43eux441ux442ux430ux43bux43eux441ux44c-2-ux43cux438ux43dux443ux442ux44b-ux43fux435ux440ux435ux434-ux44dux43aux437ux430ux43cux435ux43dux43eux43c-10}

Запомнить четыре блока.

\textbf{Замена переменной через распределение:}

\[
\mathbb E g(X)
=
\int_\Omega g(X)\,dP
=
\int_E g(x)\,P_X(dx).
\]

\textbf{Если есть плотность:}

\[
\mathbb E g(X)
=
\int g(x)p_X(x)\,dx.
\]

\textbf{Моменты и дисперсия:}

\[
\mathbb EX^k=\int x^k\,P_X(dx),
\qquad
\operatorname{Var}X=\int (x-m)^2\,P_X(dx).
\]

\textbf{Векторный случай:}

\[
m=\int x\,P_X(dx),
\qquad
\Sigma=\int (x-m)(x-m)^T\,P_X(dx).
\]

Главная ловушка: интеграл по плотности допустим только при наличии
плотности; иначе нужно интегрировать по мере распределения \(P_X\).

\newpage

\section{Билет 12. Упорядоченные и алгебраические структуры; гильбертовы
пространства и
проекторы}\label{ux431ux438ux43bux435ux442-12.-ux443ux43fux43eux440ux44fux434ux43eux447ux435ux43dux43dux44bux435-ux438-ux430ux43bux433ux435ux431ux440ux430ux438ux447ux435ux441ux43aux438ux435-ux441ux442ux440ux443ux43aux442ux443ux440ux44b-ux433ux438ux43bux44cux431ux435ux440ux442ux43eux432ux44b-ux43fux440ux43eux441ux442ux440ux430ux43dux441ux442ux432ux430-ux438-ux43fux440ux43eux435ux43aux442ux43eux440ux44b}

Источники: Ширяев 1, гл. II, §11; Сердобольская, лекция 6; лекции ДСП.

\subsection{1. Короткий
ответ}\label{ux43aux43eux440ux43eux442ux43aux438ux439-ux43eux442ux432ux435ux442-11}

В этом билете нужно связать общий язык структур с вероятностной
геометрией второго порядка.

\textbf{Упорядоченная структура} - это множество с отношением порядка.
Например, для случайных величин можно писать

\[
X\le Y \quad \text{почти наверное}.
\]

\textbf{Алгебраическая структура} - это множество с операциями. В
вероятности главные примеры: линейные пространства случайных величин,
пространства \(L^p\), алгебры событий, \(\sigma\)-алгебры.

Главная часть билета: случайные величины с конечным вторым моментом
образуют гильбертово пространство

\[
L^2(\Omega,\mathcal F,P),
\]

со скалярным произведением

\[
\langle X,Y\rangle=\mathbb E[XY]
\]

в вещественном случае и нормой

\[
\|X\|_2=(\mathbb EX^2)^{1/2}.
\]

Для центрированных величин скалярное произведение становится
ковариацией:

\[
\langle X-\mathbb EX,\ Y-\mathbb EY\rangle
=
\operatorname{Cov}(X,Y).
\]

Если \(H\) - замкнутое подпространство гильбертова пространства, то для
любого \(x\) существует единственная ортогональная проекция
\(P_Hx\in H\):

\[
x-P_Hx\perp H.
\]

То есть \(P_Hx\) - ближайший к \(x\) элемент из \(H\):

\[
\|x-P_Hx\|=\min_{h\in H}\|x-h\|.
\]

В вероятности это дает геометрический смысл условного ожидания:

\[
\mathbb E[X\mid\mathcal G]
\]

\begin{itemize}
\tightlist
\item
  ортогональная проекция \(X\in L^2\) на замкнутое подпространство
  \(\mathcal G\)-измеримых величин из \(L^2\).
\end{itemize}

\subsection{2. Полный
ответ}\label{ux43fux43eux43bux43dux44bux439-ux43eux442ux432ux435ux442-11}

\subsubsection{2.1. Переход к геометрии случайных
величин}\label{ux43fux435ux440ux435ux445ux43eux434-ux43a-ux433ux435ux43eux43cux435ux442ux440ux438ux438-ux441ux43bux443ux447ux430ux439ux43dux44bux445-ux432ux435ux43bux438ux447ux438ux43d}

В программе этот билет стоит как переход от меры и интеграла к линейной
геометрии случайных величин. До этого случайные величины рассматривались
как измеримые функции, у которых есть распределения и моменты. Теперь
добавляется идея: если у случайных величин есть вторые моменты, то с
ними можно работать как с векторами в гильбертовом пространстве.

\subsubsection{2.2. Упорядоченные
структуры}\label{ux443ux43fux43eux440ux44fux434ux43eux447ux435ux43dux43dux44bux435-ux441ux442ux440ux443ux43aux442ux443ux440ux44b}

Сначала нужен общий язык структур.

\textbf{Упорядоченная структура} - это множество \(M\) с отношением
\(\le\). Отношение называется частичным порядком, если для любых
\(a,b,c\in M\) выполнены:

\begin{enumerate}
\def\labelenumi{\arabic{enumi}.}
\tightlist
\item
  рефлексивность:
\end{enumerate}

\[
a\le a;
\]

\begin{enumerate}
\def\labelenumi{\arabic{enumi}.}
\setcounter{enumi}{1}
\tightlist
\item
  антисимметричность:
\end{enumerate}

\[
a\le b,\ b\le a\quad\Rightarrow\quad a=b;
\]

\begin{enumerate}
\def\labelenumi{\arabic{enumi}.}
\setcounter{enumi}{2}
\tightlist
\item
  транзитивность:
\end{enumerate}

\[
a\le b,\ b\le c\quad\Rightarrow\quad a\le c.
\]

Если любые два элемента сравнимы, порядок называется линейным, или
полным. Например, на \(\mathbb R\) обычный порядок полный. На множестве
подмножеств \(\Omega\) порядок по включению

\[
A\subset B
\]

частичный: не любые два множества сравнимы.

В вероятности важны несколько порядков:

\[
A\subset B
\]

для событий,

\[
X\le Y
\]

для случайных величин, обычно в смысле почти наверное. Строго говоря,
это частичный порядок на классах эквивалентности по равенству почти
наверное; на конкретных версиях случайных величин это только
предпорядок. Также используют порядок

\[
f\le g
\]

для измеримых функций. Этот порядок нужен в теореме Беппо Леви,
монотонности интеграла и монотонности условного ожидания.

Есть также порядок на неотрицательно определенных матрицах:

\[
A\le B
\quad\Longleftrightarrow\quad
B-A\ge0
\]

в смысле квадратичных форм:

\[
x^T(B-A)x\ge0
\]

для всех \(x\). Такой порядок используется для ковариационных матриц.

\subsubsection{2.3. Алгебраические
структуры}\label{ux430ux43bux433ux435ux431ux440ux430ux438ux447ux435ux441ux43aux438ux435-ux441ux442ux440ux443ux43aux442ux443ux440ux44b}

\textbf{Алгебраическая структура} - это множество с одной или
несколькими операциями, удовлетворяющими аксиомам. Примеры:

\begin{enumerate}
\def\labelenumi{\arabic{enumi}.}
\tightlist
\item
  полугруппа - множество с ассоциативной операцией;
\item
  моноид - полугруппа с нейтральным элементом;
\item
  группа - моноид, где у каждого элемента есть обратный;
\item
  кольцо - множество с операциями сложения и умножения;
\item
  поле - кольцо, где можно делить на ненулевые элементы;
\item
  линейное пространство - множество, где можно складывать векторы и
  умножать их на числа из поля.
\end{enumerate}

\subsubsection{2.4. Линейное пространство случайных
величин}\label{ux43bux438ux43dux435ux439ux43dux43eux435-ux43fux440ux43eux441ux442ux440ux430ux43dux441ux442ux432ux43e-ux441ux43bux443ux447ux430ux439ux43dux44bux445-ux432ux435ux43bux438ux447ux438ux43d}

В этом билете главный пример - линейное пространство случайных величин.
Если \(X\) и \(Y\) - случайные величины, а \(a,b\in\mathbb R\), то

\[
aX+bY
\]

снова случайная величина. Если \(X,Y\in L^2\), то по неравенству

\[
(aX+bY)^2\le 2a^2X^2+2b^2Y^2
\]

получаем \(aX+bY\in L^2\). Значит \(L^2\) - линейное пространство.

\subsubsection{2.5. Гильбертовы
пространства}\label{ux433ux438ux43bux44cux431ux435ux440ux442ux43eux432ux44b-ux43fux440ux43eux441ux442ux440ux430ux43dux441ux442ux432ux430}

Теперь переходим к гильбертовым пространствам. Сначала определения.

Линейное пространство \(V\) над \(\mathbb R\) называется пространством
со скалярным произведением, если задана функция

\[
\langle x,y\rangle
\]

такая, что:

\begin{enumerate}
\def\labelenumi{\arabic{enumi}.}
\tightlist
\item
  линейность по первому аргументу:
\end{enumerate}

\[
\langle ax_1+bx_2,y\rangle
=
a\langle x_1,y\rangle+b\langle x_2,y\rangle;
\]

\begin{enumerate}
\def\labelenumi{\arabic{enumi}.}
\setcounter{enumi}{1}
\tightlist
\item
  симметрия:
\end{enumerate}

\[
\langle x,y\rangle=\langle y,x\rangle;
\]

\begin{enumerate}
\def\labelenumi{\arabic{enumi}.}
\setcounter{enumi}{2}
\tightlist
\item
  положительная определенность:
\end{enumerate}

\[
\langle x,x\rangle\ge0,
\qquad
\langle x,x\rangle=0\Longleftrightarrow x=0.
\]

В комплексном случае вместо симметрии берут эрмитову симметрию:

\[
\langle x,y\rangle=\overline{\langle y,x\rangle},
\]

и одно из направлений линейно, другое сопряженно-линейно.

Скалярное произведение задает норму:

\[
\|x\|=\sqrt{\langle x,x\rangle}.
\]

Пространство со скалярным произведением называется гильбертовым, если
оно полно по этой норме: всякая фундаментальная последовательность
сходится к элементу того же пространства.

\subsubsection{\texorpdfstring{2.6. Пространство
\(L^2\)}{2.6. Пространство L\^{}2}}\label{ux43fux440ux43eux441ux442ux440ux430ux43dux441ux442ux432ux43e-l2}

Главный пример для вероятности:

\[
L^2(\Omega,\mathcal F,P)
=
\{X:\mathbb E X^2<\infty\}/\sim,
\]

где

\[
X\sim Y
\quad\Longleftrightarrow\quad
X=Y\ \text{почти наверное}.
\]

Важно: строго элементы \(L^2\) - не отдельные функции, а классы
эквивалентности по равенству почти наверное. В устном ответе обычно
говорят ``случайная величина из \(L^2\)'', понимая этот класс.

Скалярное произведение:

\[
\langle X,Y\rangle=\mathbb E[XY].
\]

Оно корректно определено, потому что по неравенству Коши-Буняковского

\[
\mathbb E|XY|
\le
(\mathbb EX^2)^{1/2}(\mathbb EY^2)^{1/2}<\infty.
\]

Норма:

\[
\|X\|_2=(\mathbb EX^2)^{1/2}.
\]

Расстояние:

\[
d(X,Y)=\|X-Y\|_2
=
\left(\mathbb E[(X-Y)^2]\right)^{1/2}.
\]

То есть близость в \(L^2\) - это малая среднеквадратичная ошибка.

Пространство \(L^2\) полно, поэтому это гильбертово пространство. В
вероятности это фундаментально: задачи прогноза, регрессии и фильтрации
можно формулировать как задачи ортогонального проектирования.

\subsubsection{2.7. Ортогональность и
ковариация}\label{ux43eux440ux442ux43eux433ux43eux43dux430ux43bux44cux43dux43eux441ux442ux44c-ux438-ux43aux43eux432ux430ux440ux438ux430ux446ux438ux44f}

Ортогональность в \(L^2\):

\[
X\perp Y
\quad\Longleftrightarrow\quad
\langle X,Y\rangle=\mathbb E[XY]=0.
\]

Если \(X\) и \(Y\) центрированы, то есть

\[
\mathbb EX=\mathbb EY=0,
\]

то

\[
\operatorname{Cov}(X,Y)=\mathbb E[XY]=\langle X,Y\rangle.
\]

В общем случае удобнее работать с центрированными величинами

\[
X^0=X-\mathbb EX,
\qquad
Y^0=Y-\mathbb EY.
\]

Тогда

\[
\operatorname{Cov}(X,Y)
=
\mathbb E[X^0Y^0]
=
\langle X^0,Y^0\rangle.
\]

Отсюда корреляция в Билете 14 имеет геометрический смысл косинуса угла:

\[
\rho(X,Y)
=
\frac{\langle X^0,Y^0\rangle}{\|X^0\|_2\|Y^0\|_2}.
\]

\subsubsection{2.8. Ортогональные
проекторы}\label{ux43eux440ux442ux43eux433ux43eux43dux430ux43bux44cux43dux44bux435-ux43fux440ux43eux435ux43aux442ux43eux440ux44b}

Теперь проекторы. Пусть \(H\) - линейное подпространство гильбертова
пространства \(\mathcal H\). Элемент \(y\in H\) называется ортогональной
проекцией \(x\) на \(H\), если

\[
x-y\perp H,
\]

то есть

\[
\langle x-y,h\rangle=0
\]

для всех \(h\in H\).

Главная теорема о проекции: если \(H\) - замкнутое линейное
подпространство гильбертова пространства \(\mathcal H\), то для любого
\(x\in\mathcal H\) существует единственный элемент \(P_Hx\in H\), такой
что

\[
x-P_Hx\perp H.
\]

Этот элемент является ближайшим к \(x\) элементом из \(H\):

\[
\|x-P_Hx\|
=
\inf_{h\in H}\|x-h\|.
\]

Более того, для любого \(h\in H\):

\[
\|x-h\|^2
=
\|x-P_Hx\|^2+\|P_Hx-h\|^2.
\]

Это просто теорема Пифагора, потому что

\[
x-h=(x-P_Hx)+(P_Hx-h),
\]

а две части ортогональны.

Отображение

\[
P_H:\mathcal H\to H
\]

называется ортогональным проектором. Оно обладает свойствами:

\begin{enumerate}
\def\labelenumi{\arabic{enumi}.}
\tightlist
\item
  линейность:
\end{enumerate}

\[
P_H(ax+by)=aP_Hx+bP_Hy;
\]

\begin{enumerate}
\def\labelenumi{\arabic{enumi}.}
\setcounter{enumi}{1}
\tightlist
\item
  идемпотентность:
\end{enumerate}

\[
P_H^2=P_H;
\]

\begin{enumerate}
\def\labelenumi{\arabic{enumi}.}
\setcounter{enumi}{2}
\tightlist
\item
  самосопряженность:
\end{enumerate}

\[
\langle P_Hx,y\rangle=\langle x,P_Hy\rangle;
\]

\begin{enumerate}
\def\labelenumi{\arabic{enumi}.}
\setcounter{enumi}{3}
\tightlist
\item
  сжатие нормы:
\end{enumerate}

\[
\|P_Hx\|\le\|x\|;
\]

\begin{enumerate}
\def\labelenumi{\arabic{enumi}.}
\setcounter{enumi}{4}
\tightlist
\item
  разложение:
\end{enumerate}

\[
x=P_Hx+(x-P_Hx),
\qquad
P_Hx\in H,
\quad
x-P_Hx\in H^\perp.
\]

Условие замкнутости \(H\) важно. Если подпространство не замкнуто,
ближайший элемент может не существовать внутри самого \(H\).

\subsubsection{2.9. Матрица Грама и линейная
регрессия}\label{ux43cux430ux442ux440ux438ux446ux430-ux433ux440ux430ux43cux430-ux438-ux43bux438ux43dux435ux439ux43dux430ux44f-ux440ux435ux433ux440ux435ux441ux441ux438ux44f}

В конечномерном случае проекция вычисляется через матрицу Грама. Пусть
\(h_1,\ldots,h_n\) - линейно независимые элементы гильбертова
пространства, а

\[
H=\operatorname{span}(h_1,\ldots,h_n).
\]

Ищем проекцию в виде

\[
P_Hx=\sum_{j=1}^n a_jh_j.
\]

Условие ортогональности ошибки:

\[
x-\sum_{j=1}^n a_jh_j\perp h_i,
\qquad i=1,\ldots,n.
\]

Значит коэффициенты удовлетворяют нормальным уравнениям:

\[
\sum_{j=1}^n a_j\langle h_j,h_i\rangle
=
\langle x,h_i\rangle,
\qquad i=1,\ldots,n.
\]

Матрица

\[
G=(\langle h_i,h_j\rangle)_{i,j=1}^n
\]

называется матрицей Грама. Если \(h_1,\ldots,h_n\) линейно независимы,
то \(G\) положительно определена и обратима. В ортонормированном случае

\[
\langle h_i,h_j\rangle=\delta_{ij},
\]

и формула упрощается:

\[
P_Hx=\sum_{j=1}^n \langle x,h_j\rangle h_j.
\]

В вероятности это дает линейную регрессию. Пусть
\(X,Y_1,\ldots,Y_n\in L^2\), и мы хотим приблизить \(X\) линейной
комбинацией \(Y_1,\ldots,Y_n\):

\[
\widehat X=\sum_{j=1}^n a_jY_j.
\]

Оптимальная в среднеквадратичном смысле оценка минимизирует

\[
\mathbb E\left(X-\sum_{j=1}^n a_jY_j\right)^2.
\]

Она является ортогональной проекцией \(X\) на

\[
H=\operatorname{span}(Y_1,\ldots,Y_n).
\]

Нормальные уравнения:

\[
\mathbb E\left[\left(X-\sum_{j=1}^n a_jY_j\right)Y_i\right]=0,
\qquad i=1,\ldots,n.
\]

То есть ошибка оптимального прогноза ортогональна всем используемым
регрессорам.

Если работают с центрированными величинами, то матрица Грама становится
ковариационной матрицей:

\[
G_{ij}
=
\langle Y_i^0,Y_j^0\rangle
=
\operatorname{Cov}(Y_i,Y_j).
\]

Поэтому ковариационные матрицы всегда неотрицательно определены:

\[
a^TGa
=
\left\|\sum_{i=1}^n a_iY_i^0\right\|_2^2
\ge0.
\]

\subsubsection{2.10. Условное ожидание как ортогональная
проекция}\label{ux443ux441ux43bux43eux432ux43dux43eux435-ux43eux436ux438ux434ux430ux43dux438ux435-ux43aux430ux43a-ux43eux440ux442ux43eux433ux43eux43dux430ux43bux44cux43dux430ux44f-ux43fux440ux43eux435ux43aux446ux438ux44f}

Теперь условное ожидание как проектор. Пусть \(X\in L^2\) и
\(\mathcal G\subset\mathcal F\) - под-\(\sigma\)-алгебра. Рассмотрим

\[
L^2(\mathcal G)
=
\{Y\in L^2(\Omega,\mathcal F,P):Y\ \mathcal G\text{-измерима}\}.
\]

Это замкнутое линейное подпространство \(L^2(\mathcal F)\). Тогда

\[
\mathbb E[X\mid\mathcal G]
\]

является ортогональной проекцией \(X\) на \(L^2(\mathcal G)\):

\[
P_{L^2(\mathcal G)}X=\mathbb E[X\mid\mathcal G].
\]

Почему? По определению условного ожидания для любого \(G\in\mathcal G\):

\[
\int_G(X-\mathbb E[X\mid\mathcal G])\,dP=0.
\]

Это означает ортогональность ошибки индикаторам \(\mathbf 1_G\). А через
линейность и предельный переход ошибка ортогональна всем ограниченным
\(\mathcal G\)-измеримым величинам, а затем всем величинам из
\(L^2(\mathcal G)\). Значит

\[
X-\mathbb E[X\mid\mathcal G]\perp L^2(\mathcal G).
\]

Отсюда следует оптимальное свойство условного ожидания:

\[
\mathbb E[(X-\mathbb E[X\mid\mathcal G])^2]
=
\min_{Y\in L^2(\mathcal G)}
\mathbb E[(X-Y)^2].
\]

То есть условное ожидание - лучший среднеквадратичный прогноз \(X\) по
информации \(\mathcal G\).

\subsubsection{2.11.
Заключение}\label{ux437ux430ux43aux43bux44eux447ux435ux43dux438ux435-1}

Итак, билет строит мост: порядок и алгебраические операции задают общий
язык; \(L^2\) превращает случайные величины в векторы; ковариации
становятся скалярными произведениями; условные ожидания и линейные
оценки становятся ортогональными проекциями.

\subsection{3. Основные
определения}\label{ux43eux441ux43dux43eux432ux43dux44bux435-ux43eux43fux440ux435ux434ux435ux43bux435ux43dux438ux44f-11}

\subsubsection{Частично упорядоченное
множество}\label{ux447ux430ux441ux442ux438ux447ux43dux43e-ux443ux43fux43eux440ux44fux434ux43eux447ux435ux43dux43dux43eux435-ux43cux43dux43eux436ux435ux441ux442ux432ux43e}

Множество \(M\) с отношением \(\le\), которое рефлексивно,
антисимметрично и транзитивно.

\subsubsection{Линейный
порядок}\label{ux43bux438ux43dux435ux439ux43dux44bux439-ux43fux43eux440ux44fux434ux43eux43a}

Частичный порядок, в котором любые два элемента сравнимы.

\subsubsection{Верхняя и нижняя
грань}\label{ux432ux435ux440ux445ux43dux44fux44f-ux438-ux43dux438ux436ux43dux44fux44f-ux433ux440ux430ux43dux44c}

Для \(A\subset M\) элемент \(u\) - верхняя грань, если \(a\le u\) для
всех \(a\in A\). Нижняя грань определяется аналогично.

\subsubsection{Супремум и
инфимум}\label{ux441ux443ux43fux440ux435ux43cux443ux43c-ux438-ux438ux43dux444ux438ux43cux443ux43c}

Наименьшая верхняя грань называется супремумом, наибольшая нижняя грань
- инфимумом.

\subsubsection{Алгебраическая
структура}\label{ux430ux43bux433ux435ux431ux440ux430ux438ux447ux435ux441ux43aux430ux44f-ux441ux442ux440ux443ux43aux442ux443ux440ux430}

Множество с заданными операциями, удовлетворяющими аксиомам. Примеры:
группа, кольцо, поле, линейное пространство.

\subsubsection{Линейное
пространство}\label{ux43bux438ux43dux435ux439ux43dux43eux435-ux43fux440ux43eux441ux442ux440ux430ux43dux441ux442ux432ux43e}

Множество \(V\), где определены сложение векторов и умножение на
скаляры, причем выполнены стандартные аксиомы линейности.

\subsubsection{Линейное
подпространство}\label{ux43bux438ux43dux435ux439ux43dux43eux435-ux43fux43eux434ux43fux440ux43eux441ux442ux440ux430ux43dux441ux442ux432ux43e}

Подмножество \(H\subset V\), замкнутое относительно линейных комбинаций:

\[
x,y\in H,\ a,b\in\mathbb R
\quad\Rightarrow\quad
ax+by\in H.
\]

\subsubsection{Скалярное
произведение}\label{ux441ux43aux430ux43bux44fux440ux43dux43eux435-ux43fux440ux43eux438ux437ux432ux435ux434ux435ux43dux438ux435}

Функция \(\langle x,y\rangle\), линейная по одному аргументу,
симметричная в вещественном случае и положительно определенная.

\subsubsection{Норма, порожденная скалярным
произведением}\label{ux43dux43eux440ux43cux430-ux43fux43eux440ux43eux436ux434ux435ux43dux43dux430ux44f-ux441ux43aux430ux43bux44fux440ux43dux44bux43c-ux43fux440ux43eux438ux437ux432ux435ux434ux435ux43dux438ux435ux43c}

\[
\|x\|=\sqrt{\langle x,x\rangle}.
\]

\subsubsection{Гильбертово
пространство}\label{ux433ux438ux43bux44cux431ux435ux440ux442ux43eux432ux43e-ux43fux440ux43eux441ux442ux440ux430ux43dux441ux442ux432ux43e}

Полное линейное пространство со скалярным произведением.

\subsubsection{\texorpdfstring{Пространство
\(L^2\)}{Пространство L\^{}2}}\label{ux43fux440ux43eux441ux442ux440ux430ux43dux441ux442ux432ux43e-l2-1}

\[
L^2(\Omega,\mathcal F,P)=\{X:\mathbb EX^2<\infty\}/\sim,
\]

где \(X\sim Y\), если \(X=Y\) почти наверное.

\subsubsection{\texorpdfstring{Скалярное произведение в
\(L^2\)}{Скалярное произведение в L\^{}2}}\label{ux441ux43aux430ux43bux44fux440ux43dux43eux435-ux43fux440ux43eux438ux437ux432ux435ux434ux435ux43dux438ux435-ux432-l2}

\[
\langle X,Y\rangle=\mathbb E[XY].
\]

Для комплексных величин обычно берут

\[
\langle X,Y\rangle=\mathbb E[X\overline Y]
\]

или сопряжение в другом аргументе, в зависимости от конвенции.

\subsubsection{Ортогональность}\label{ux43eux440ux442ux43eux433ux43eux43dux430ux43bux44cux43dux43eux441ux442ux44c}

\[
x\perp y
\quad\Longleftrightarrow\quad
\langle x,y\rangle=0.
\]

В \(L^2\):

\[
X\perp Y
\quad\Longleftrightarrow\quad
\mathbb E[XY]=0.
\]

\subsubsection{Ортогональное
дополнение}\label{ux43eux440ux442ux43eux433ux43eux43dux430ux43bux44cux43dux43eux435-ux434ux43eux43fux43eux43bux43dux435ux43dux438ux435}

\[
H^\perp=\{x:\langle x,h\rangle=0\ \text{для всех }h\in H\}.
\]

\subsubsection{Ортогональная
проекция}\label{ux43eux440ux442ux43eux433ux43eux43dux430ux43bux44cux43dux430ux44f-ux43fux440ux43eux435ux43aux446ux438ux44f}

Элемент \(P_Hx\in H\), такой что

\[
x-P_Hx\perp H.
\]

\subsubsection{Ортогональный
проектор}\label{ux43eux440ux442ux43eux433ux43eux43dux430ux43bux44cux43dux44bux439-ux43fux440ux43eux435ux43aux442ux43eux440}

Отображение \(P_H:x\mapsto P_Hx\) на замкнутое подпространство \(H\).

\subsubsection{Матрица
Грама}\label{ux43cux430ux442ux440ux438ux446ux430-ux433ux440ux430ux43cux430}

Для \(h_1,\ldots,h_n\):

\[
G=(\langle h_i,h_j\rangle)_{i,j=1}^n.
\]

\subsubsection{Ортонормированная
система}\label{ux43eux440ux442ux43eux43dux43eux440ux43cux438ux440ux43eux432ux430ux43dux43dux430ux44f-ux441ux438ux441ux442ux435ux43cux430}

Система \((e_i)\), для которой

\[
\langle e_i,e_j\rangle=\delta_{ij}.
\]

\subsection{4. Основные теоремы и
свойства}\label{ux43eux441ux43dux43eux432ux43dux44bux435-ux442ux435ux43eux440ux435ux43cux44b-ux438-ux441ux432ux43eux439ux441ux442ux432ux430-11}

\subsubsection{\texorpdfstring{Теорема 1. \(L^2\) - гильбертово
пространство}{Теорема 1. L\^{}2 - гильбертово пространство}}\label{ux442ux435ux43eux440ux435ux43cux430-1.-l2---ux433ux438ux43bux44cux431ux435ux440ux442ux43eux432ux43e-ux43fux440ux43eux441ux442ux440ux430ux43dux441ux442ux432ux43e}

Пространство случайных величин с конечным вторым моментом и скалярным
произведением

\[
\langle X,Y\rangle=\mathbb E[XY]
\]

является гильбертовым пространством после отождествления случайных
величин, равных почти наверное.

\subsubsection{Условие
корректности}\label{ux443ux441ux43bux43eux432ux438ux435-ux43aux43eux440ux440ux435ux43aux442ux43dux43eux441ux442ux438}

Если \(X,Y\in L^2\), то \(\mathbb E|XY|<\infty\) по Коши-Буняковскому.

\subsubsection{Теорема 2.
Коши-Буняковский}\label{ux442ux435ux43eux440ux435ux43cux430-2.-ux43aux43eux448ux438-ux431ux443ux43dux44fux43aux43eux432ux441ux43aux438ux439}

В любом пространстве со скалярным произведением:

\[
|\langle x,y\rangle|
\le
\|x\|\|y\|.
\]

В \(L^2\):

\[
|\mathbb E[XY]|
\le
(\mathbb EX^2)^{1/2}(\mathbb EY^2)^{1/2}.
\]

\subsubsection{Теорема 3. Проекция на замкнутое
подпространство}\label{ux442ux435ux43eux440ux435ux43cux430-3.-ux43fux440ux43eux435ux43aux446ux438ux44f-ux43dux430-ux437ux430ux43cux43aux43dux443ux442ux43eux435-ux43fux43eux434ux43fux440ux43eux441ux442ux440ux430ux43dux441ux442ux432ux43e}

Если \(H\) - замкнутое линейное подпространство гильбертова пространства
\(\mathcal H\), то для любого \(x\in\mathcal H\) существует единственный
\(P_Hx\in H\), такой что

\[
x-P_Hx\perp H.
\]

При этом

\[
\|x-P_Hx\|
=
\min_{h\in H}\|x-h\|.
\]

\subsubsection{Свойство 4. Пифагорова формула для
проекции}\label{ux441ux432ux43eux439ux441ux442ux432ux43e-4.-ux43fux438ux444ux430ux433ux43eux440ux43eux432ux430-ux444ux43eux440ux43cux443ux43bux430-ux434ux43bux44f-ux43fux440ux43eux435ux43aux446ux438ux438}

Для любого \(h\in H\):

\[
\|x-h\|^2
=
\|x-P_Hx\|^2+\|P_Hx-h\|^2.
\]

\subsubsection{Свойство 5. Свойства ортогонального
проектора}\label{ux441ux432ux43eux439ux441ux442ux432ux43e-5.-ux441ux432ux43eux439ux441ux442ux432ux430-ux43eux440ux442ux43eux433ux43eux43dux430ux43bux44cux43dux43eux433ux43e-ux43fux440ux43eux435ux43aux442ux43eux440ux430}

Ортогональный проектор \(P_H\):

\[
P_H^2=P_H,
\]

линеен,

\[
\langle P_Hx,y\rangle=\langle x,P_Hy\rangle,
\]

и

\[
\|P_Hx\|\le\|x\|.
\]

\subsubsection{Теорема 6. Нормальные уравнения конечномерной
проекции}\label{ux442ux435ux43eux440ux435ux43cux430-6.-ux43dux43eux440ux43cux430ux43bux44cux43dux44bux435-ux443ux440ux430ux432ux43dux435ux43dux438ux44f-ux43aux43eux43dux435ux447ux43dux43eux43cux435ux440ux43dux43eux439-ux43fux440ux43eux435ux43aux446ux438ux438}

Если \(H=\operatorname{span}(h_1,\ldots,h_n)\) и

\[
P_Hx=\sum_{j=1}^n a_jh_j,
\]

то коэффициенты \(a_j\) определяются системой

\[
\sum_{j=1}^n a_j\langle h_j,h_i\rangle
=
\langle x,h_i\rangle,
\qquad i=1,\ldots,n.
\]

\subsubsection{Свойство 7. Проекция на ортонормированную
систему}\label{ux441ux432ux43eux439ux441ux442ux432ux43e-7.-ux43fux440ux43eux435ux43aux446ux438ux44f-ux43dux430-ux43eux440ux442ux43eux43dux43eux440ux43cux438ux440ux43eux432ux430ux43dux43dux443ux44e-ux441ux438ux441ux442ux435ux43cux443}

Если \(e_1,\ldots,e_n\) ортонормированы, то

\[
P_Hx=\sum_{j=1}^n\langle x,e_j\rangle e_j.
\]

\subsubsection{Свойство 8. Матрица Грама неотрицательно
определена}\label{ux441ux432ux43eux439ux441ux442ux432ux43e-8.-ux43cux430ux442ux440ux438ux446ux430-ux433ux440ux430ux43cux430-ux43dux435ux43eux442ux440ux438ux446ux430ux442ux435ux43bux44cux43dux43e-ux43eux43fux440ux435ux434ux435ux43bux435ux43dux430}

Для любых \(a_1,\ldots,a_n\):

\[
\sum_{i,j=1}^n a_i a_j\langle h_i,h_j\rangle
=
\left\|\sum_{i=1}^n a_ih_i\right\|^2
\ge0.
\]

Если \(h_1,\ldots,h_n\) линейно независимы, то матрица Грама
положительно определена.

\subsubsection{Свойство 9. Ковариационная матрица как матрица
Грама}\label{ux441ux432ux43eux439ux441ux442ux432ux43e-9.-ux43aux43eux432ux430ux440ux438ux430ux446ux438ux43eux43dux43dux430ux44f-ux43cux430ux442ux440ux438ux446ux430-ux43aux430ux43a-ux43cux430ux442ux440ux438ux446ux430-ux433ux440ux430ux43cux430}

Для центрированных величин \(X_1,\ldots,X_n\in L^2\):

\[
G_{ij}=\mathbb E[X_iX_j]
\]

\begin{itemize}
\tightlist
\item
  матрица Грама. Поэтому она неотрицательно определена.
\end{itemize}

Для нецентрированных величин ковариационная матрица является матрицей
Грама центрированных величин:

\[
\Sigma_{ij}
=
\operatorname{Cov}(X_i,X_j)
=
\langle X_i-\mathbb EX_i,\ X_j-\mathbb EX_j\rangle.
\]

\subsubsection{Теорема 10. Условное ожидание как ортогональная
проекция}\label{ux442ux435ux43eux440ux435ux43cux430-10.-ux443ux441ux43bux43eux432ux43dux43eux435-ux43eux436ux438ux434ux430ux43dux438ux435-ux43aux430ux43a-ux43eux440ux442ux43eux433ux43eux43dux430ux43bux44cux43dux430ux44f-ux43fux440ux43eux435ux43aux446ux438ux44f}

Если \(X\in L^2\) и \(\mathcal G\subset\mathcal F\), то

\[
\mathbb E[X\mid\mathcal G]
=
P_{L^2(\mathcal G)}X.
\]

Эквивалентно:

\[
X-\mathbb E[X\mid\mathcal G]\perp L^2(\mathcal G).
\]

И оптимальное свойство:

\[
\mathbb E[(X-\mathbb E[X\mid\mathcal G])^2]
=
\min_{Y\in L^2(\mathcal G)}
\mathbb E[(X-Y)^2].
\]

\subsection{5. Примеры}\label{ux43fux440ux438ux43cux435ux440ux44b-11}

\subsubsection{Пример 1. Евклидово
пространство}\label{ux43fux440ux438ux43cux435ux440-1.-ux435ux432ux43aux43bux438ux434ux43eux432ux43e-ux43fux440ux43eux441ux442ux440ux430ux43dux441ux442ux432ux43e}

В \(\mathbb R^n\):

\[
\langle x,y\rangle=x^Ty,
\qquad
\|x\|=\sqrt{x_1^2+\cdots+x_n^2}.
\]

Проекция точки на прямую или плоскость - обычная ортогональная проекция
из геометрии.

\subsubsection{\texorpdfstring{Пример 2. Пространство
\(L^2\)}{Пример 2. Пространство L\^{}2}}\label{ux43fux440ux438ux43cux435ux440-2.-ux43fux440ux43eux441ux442ux440ux430ux43dux441ux442ux432ux43e-l2}

Пусть \(X,Y\in L^2\). Тогда

\[
\langle X,Y\rangle=\mathbb E[XY].
\]

Если \(\mathbb EX=\mathbb EY=0\), то

\[
X\perp Y
\quad\Longleftrightarrow\quad
\operatorname{Cov}(X,Y)=0.
\]

То есть ортогональность центрированных величин - это
некоррелированность.

\subsubsection{Пример 3. Проекция на
константы}\label{ux43fux440ux438ux43cux435ux440-3.-ux43fux440ux43eux435ux43aux446ux438ux44f-ux43dux430-ux43aux43eux43dux441ux442ux430ux43dux442ux44b}

Пусть \(H\) - подпространство констант:

\[
H=\{c:\ c\in\mathbb R\}.
\]

Проекция \(X\in L^2\) на \(H\) равна

\[
P_HX=\mathbb EX.
\]

Действительно,

\[
\mathbb E[X-\mathbb EX]=0,
\]

то есть ошибка ортогональна всем константам.

\subsubsection{Пример 4. Условное ожидание по
разбиению}\label{ux43fux440ux438ux43cux435ux440-4.-ux443ux441ux43bux43eux432ux43dux43eux435-ux43eux436ux438ux434ux430ux43dux438ux435-ux43fux43e-ux440ux430ux437ux431ux438ux435ux43dux438ux44e}

Пусть \(\mathcal G=\sigma(D_1,\ldots,D_n)\), где \(D_i\) - атомы
разбиения и \(P(D_i)>0\). Тогда

\[
\mathbb E[X\mid\mathcal G]
=
\sum_{i=1}^n
\frac{\int_{D_i}X\,dP}{P(D_i)}\mathbf 1_{D_i}.
\]

Это проекция \(X\) на подпространство функций, постоянных на каждом
\(D_i\).

\subsubsection{Пример 5. Линейная регрессия на одну
величину}\label{ux43fux440ux438ux43cux435ux440-5.-ux43bux438ux43dux435ux439ux43dux430ux44f-ux440ux435ux433ux440ux435ux441ux441ux438ux44f-ux43dux430-ux43eux434ux43dux443-ux432ux435ux43bux438ux447ux438ux43dux443}

Пусть \(X,Y\in L^2\), \(\mathbb EY=0\), \(\mathbb EY^2>0\). Ищем лучшую
оценку \(X\) вида \(aY\). Минимизируем

\[
\mathbb E(X-aY)^2.
\]

Условие ортогональности:

\[
\mathbb E[(X-aY)Y]=0.
\]

Отсюда

\[
a=\frac{\mathbb E[XY]}{\mathbb E[Y^2]}.
\]

Если \(X,Y\) центрированы, то

\[
a=\frac{\operatorname{Cov}(X,Y)}{\operatorname{Var}Y}.
\]

\subsubsection{Пример 6. Матрица Грама и ковариационная
матрица}\label{ux43fux440ux438ux43cux435ux440-6.-ux43cux430ux442ux440ux438ux446ux430-ux433ux440ux430ux43cux430-ux438-ux43aux43eux432ux430ux440ux438ux430ux446ux438ux43eux43dux43dux430ux44f-ux43cux430ux442ux440ux438ux446ux430}

Пусть \(X=(X_1,\ldots,X_n)\), где все \(X_i\) центрированы и имеют
вторые моменты. Тогда

\[
G_{ij}=\mathbb E[X_iX_j]
\]

является матрицей Грама. Для любого \(a\in\mathbb R^n\):

\[
a^TGa
=
\mathbb E\left(\sum_{i=1}^n a_iX_i\right)^2
\ge0.
\]

Поэтому ковариационная матрица неотрицательно определена.

\subsubsection{Пример 7. Ортогональность не равна
независимости}\label{ux43fux440ux438ux43cux435ux440-7.-ux43eux440ux442ux43eux433ux43eux43dux430ux43bux44cux43dux43eux441ux442ux44c-ux43dux435-ux440ux430ux432ux43dux430-ux43dux435ux437ux430ux432ux438ux441ux438ux43cux43eux441ux442ux438}

Пусть \(X\) равномерна на \([-1,1]\), а

\[
Y=X^2.
\]

Тогда \(Y\) полностью определяется по \(X\), значит независимости нет.
Но

\[
\mathbb E[X]=0,
\qquad
\mathbb E[XY]=\mathbb E[X^3]=0,
\]

поэтому центрированные величины могут быть ортогональны или
некоррелированы без независимости.

\subsection{6. Схемы доказательств /
обоснований}\label{ux441ux445ux435ux43cux44b-ux434ux43eux43aux430ux437ux430ux442ux435ux43bux44cux441ux442ux432-ux43eux431ux43eux441ux43dux43eux432ux430ux43dux438ux439-4}

\textbf{Почему \(L^2\) - пространство со скалярным произведением.}

Для \(X,Y\in L^2\) произведение \(XY\) интегрируемо по
Коши-Буняковскому:

\[
\mathbb E|XY|
\le
(\mathbb EX^2)^{1/2}(\mathbb EY^2)^{1/2}.
\]

Линейность и симметрия следуют из линейности математического ожидания.
Положительность:

\[
\langle X,X\rangle=\mathbb EX^2\ge0.
\]

Если \(\mathbb EX^2=0\), то \(X=0\) почти наверное, поэтому в \(L^2\)
это нулевой элемент.

\textbf{Идея полноты \(L^2\).}

Если \((X_n)\) фундаментальна по норме \(\|\cdot\|_2\), то существует
\(X\in L^2\), к которому \(X_n\) сходится в среднем квадратичном:

\[
\mathbb E[(X_n-X)^2]\to0.
\]

Полное доказательство относится к теории \(L^p\), но для экзамена важно
знать результат: \(L^2\) полно, значит это гильбертово пространство.

\textbf{Доказательство теоремы о проекции: идея.}

Нужно минимизировать функцию

\[
h\mapsto\|x-h\|^2
\]

по \(h\in H\). В гильбертовом пространстве замкнутость \(H\)
обеспечивает существование ближайшего элемента. Если \(y=P_Hx\) -
минимум, то для любого направления \(u\in H\) функция

\[
\varphi(t)=\|x-y-tu\|^2
\]

имеет минимум при \(t=0\). Поэтому

\[
\varphi'(0)=-2\langle x-y,u\rangle=0.
\]

Значит

\[
x-y\perp H.
\]

Обратно, если \(x-y\perp H\), то для любого \(h\in H\)

\[
x-h=(x-y)+(y-h)
\]

и по Пифагору

\[
\|x-h\|^2=\|x-y\|^2+\|y-h\|^2\ge\|x-y\|^2.
\]

Значит \(y\) - ближайший элемент.

\textbf{Доказательство нормальных уравнений.}

Пусть

\[
\widehat x=\sum_{j=1}^n a_jh_j.
\]

Условие проекции:

\[
x-\widehat x\perp h_i
\]

для каждого \(i\). Поэтому

\[
\left\langle x-\sum_{j=1}^n a_jh_j,\ h_i\right\rangle=0.
\]

Раскрывая линейность, получаем

\[
\sum_{j=1}^n a_j\langle h_j,h_i\rangle=\langle x,h_i\rangle.
\]

\textbf{Почему условное ожидание - проекция.}

Пусть \(Y=\mathbb E[X\mid\mathcal G]\). По определению условного
ожидания

\[
\int_G(X-Y)\,dP=0
\]

для всех \(G\in\mathcal G\). Это значит

\[
\mathbb E[(X-Y)\mathbf 1_G]=0.
\]

Линейные комбинации индикаторов дают простые \(\mathcal G\)-измеримые
функции. Предельным переходом получаем

\[
\mathbb E[(X-Y)Z]=0
\]

для всех \(Z\in L^2(\mathcal G)\). Значит

\[
X-Y\perp L^2(\mathcal G),
\]

и \(Y\) является ортогональной проекцией.

\subsection{7. Что написать на
доске}\label{ux447ux442ux43e-ux43dux430ux43fux438ux441ux430ux442ux44c-ux43dux430-ux434ux43eux441ux43aux435-11}

\begin{enumerate}
\def\labelenumi{\arabic{enumi}.}
\tightlist
\item
  Частичный порядок:
\end{enumerate}

\[
a\le a,\quad
a\le b,\ b\le a\Rightarrow a=b,\quad
a\le b,\ b\le c\Rightarrow a\le c.
\]

\begin{enumerate}
\def\labelenumi{\arabic{enumi}.}
\setcounter{enumi}{1}
\tightlist
\item
  \(L^2\):
\end{enumerate}

\[
L^2=\{X:\mathbb EX^2<\infty\}/\sim.
\]

\begin{enumerate}
\def\labelenumi{\arabic{enumi}.}
\setcounter{enumi}{2}
\tightlist
\item
  Скалярное произведение:
\end{enumerate}

\[
\langle X,Y\rangle=\mathbb E[XY],
\qquad
\|X\|_2=(\mathbb EX^2)^{1/2}.
\]

\begin{enumerate}
\def\labelenumi{\arabic{enumi}.}
\setcounter{enumi}{3}
\tightlist
\item
  Центрированные величины:
\end{enumerate}

\[
\operatorname{Cov}(X,Y)
=
\langle X-\mathbb EX,\ Y-\mathbb EY\rangle.
\]

\begin{enumerate}
\def\labelenumi{\arabic{enumi}.}
\setcounter{enumi}{4}
\tightlist
\item
  Проектор:
\end{enumerate}

\[
P_Hx\in H,
\qquad
x-P_Hx\perp H.
\]

\begin{enumerate}
\def\labelenumi{\arabic{enumi}.}
\setcounter{enumi}{5}
\tightlist
\item
  Нормальные уравнения:
\end{enumerate}

\[
\sum_j a_j\langle h_j,h_i\rangle=\langle x,h_i\rangle.
\]

\begin{enumerate}
\def\labelenumi{\arabic{enumi}.}
\setcounter{enumi}{6}
\tightlist
\item
  Условное ожидание как проекция:
\end{enumerate}

\[
\mathbb E[X\mid\mathcal G]
=
P_{L^2(\mathcal G)}X.
\]

\begin{enumerate}
\def\labelenumi{\arabic{enumi}.}
\setcounter{enumi}{7}
\tightlist
\item
  Матрица Грама:
\end{enumerate}

\[
G_{ij}=\langle h_i,h_j\rangle,
\qquad
a^TGa=\left\|\sum_i a_ih_i\right\|^2\ge0.
\]

\subsection{8. Типичные дополнительные
вопросы}\label{ux442ux438ux43fux438ux447ux43dux44bux435-ux434ux43eux43fux43eux43bux43dux438ux442ux435ux43bux44cux43dux44bux435-ux432ux43eux43fux440ux43eux441ux44b-11}

\textbf{Вопрос. Что такое частичный порядок?}

Это отношение \(\le\), которое рефлексивно, антисимметрично и
транзитивно. Пример: включение множеств \(A\subset B\).

\textbf{Вопрос. Чем частичный порядок отличается от линейного?}

В линейном порядке любые два элемента сравнимы. В частичном - не
обязательно. Например, два события могут не быть вложены друг в друга.

\textbf{Вопрос. Какие алгебраические структуры важны в вероятности?}

Линейные пространства случайных величин, пространства \(L^p\), алгебры и
\(\sigma\)-алгебры событий, матрицы с операциями сложения и умножения.

\textbf{Вопрос. Почему \(L^2\) - линейное пространство?}

Если \(X,Y\in L^2\), то \(aX+bY\in L^2\), потому что квадрат линейной
комбинации оценивается через \(X^2\) и \(Y^2\).

\textbf{Вопрос. Почему \(\mathbb E[XY]\) корректно определено для
\(X,Y\in L^2\)?}

По неравенству Коши-Буняковского:

\[
\mathbb E|XY|\le(\mathbb EX^2)^{1/2}(\mathbb EY^2)^{1/2}.
\]

\textbf{Вопрос. Что значит полнота \(L^2\)?}

Всякая фундаментальная по среднеквадратичной норме последовательность
случайных величин сходится в \(L^2\) к случайной величине из \(L^2\).

\textbf{Вопрос. Что такое ортогональность в \(L^2\)?}

\[
X\perp Y
\quad\Longleftrightarrow\quad
\mathbb E[XY]=0.
\]

Для центрированных величин это то же самое, что некоррелированность.

\textbf{Вопрос. Ортогональность означает независимость?}

Нет. Ортогональность означает нулевое скалярное произведение, то есть
для центрированных величин нулевую ковариацию. Независимость сильнее.

\textbf{Вопрос. Зачем нужна замкнутость подпространства в теореме о
проекции?}

Чтобы ближайший элемент действительно лежал в \(H\). Без замкнутости
минимум может достигаться только в замыкании \(H\).

\textbf{Вопрос. Как найти проекцию на конечномерное подпространство?}

Записать

\[
P_Hx=\sum_j a_jh_j
\]

и решить нормальные уравнения:

\[
\sum_j a_j\langle h_j,h_i\rangle=\langle x,h_i\rangle.
\]

\textbf{Вопрос. Что такое матрица Грама?}

\[
G_{ij}=\langle h_i,h_j\rangle.
\]

Она неотрицательно определена, потому что

\[
a^TGa=\left\|\sum_i a_ih_i\right\|^2\ge0.
\]

\textbf{Вопрос. Почему ковариационная матрица неотрицательно
определена?}

Она является матрицей Грама центрированных случайных величин:

\[
\Sigma_{ij}
=
\langle X_i-\mathbb EX_i,\ X_j-\mathbb EX_j\rangle.
\]

\textbf{Вопрос. Почему условное ожидание - проекция?}

Потому что ошибка

\[
X-\mathbb E[X\mid\mathcal G]
\]

ортогональна всем \(\mathcal G\)-измеримым величинам из \(L^2\).

\textbf{Вопрос. Какой оптимальный смысл условного ожидания?}

Это лучший среднеквадратичный прогноз:

\[
\mathbb E[(X-\mathbb E[X\mid\mathcal G])^2]
=
\min_{Y\in L^2(\mathcal G)}
\mathbb E[(X-Y)^2].
\]

\subsection{9. Типичные
ошибки}\label{ux442ux438ux43fux438ux447ux43dux44bux435-ux43eux448ux438ux431ux43aux438-11}

\begin{enumerate}
\def\labelenumi{\arabic{enumi}.}
\tightlist
\item
  Путать ортогональность и независимость.
\item
  Забывать, что в \(L^2\) элементы считаются с точностью до равенства
  почти наверное.
\item
  Использовать скалярное произведение \(\mathbb E[XY]\) без проверки
  вторых моментов.
\item
  Называть любое линейное подпространство подходящим для проекции,
  забывая замкнутость.
\item
  Путать матрицу вторых моментов \(\mathbb E[XX^T]\) и ковариационную
  матрицу \(\mathbb E[(X-m)(X-m)^T]\).
\item
  Забывать центрирование при связи ковариации со скалярным
  произведением.
\item
  Считать, что условное ожидание как проекция верно только для конечных
  разбиений. На самом деле оно верно для любой под-\(\sigma\)-алгебры
  при \(X\in L^2\).
\item
  В нормальных уравнениях ставить коэффициенты без матрицы Грама, как
  будто система всегда ортонормирована.
\item
  Путать проектор \(P_H\) и произвольный линейный оператор с \(P^2=P\):
  ортогональный проектор еще самосопряжен.
\item
  Забывать, что \(L^2\)-норма измеряет среднеквадратичную ошибку, а не
  поточечное расстояние.
\end{enumerate}

\subsection{10. Связи с другими
билетами}\label{ux441ux432ux44fux437ux438-ux441-ux434ux440ux443ux433ux438ux43cux438-ux431ux438ux43bux435ux442ux430ux43cux438-11}

\textbf{Билет 10.} Условное ожидание определяется через интегральное
свойство, а в \(L^2\) получает геометрический смысл ортогональной
проекции.

\textbf{Билет 11.} Вторые моменты, дисперсии и ковариации вычисляются
через интегралы по распределению; здесь они становятся нормами и
скалярными произведениями.

\textbf{Билет 13.} Матрица вторых моментов и ковариационная матрица
связаны с производными характеристической функции.

\textbf{Билет 14.} Коэффициент корреляции - косинус угла между
центрированными величинами в \(L^2\).

\textbf{Билет 18.} Гауссовские векторы задаются средним и ковариационной
матрицей; ортогонализация и проекции особенно удобны в гауссовском
случае.

\textbf{Билет 24.} Корреляционная функция и эргодичность по среднему
используют \(L^2\)-оценки и дисперсии средних.

\textbf{Билет 25.} Спектральная теория стационарных последовательностей
строится в гильбертовом пространстве второго порядка.

\subsection{11. Минимум на
удовлетворительно}\label{ux43cux438ux43dux438ux43cux443ux43c-ux43dux430-ux443ux434ux43eux432ux43bux435ux442ux432ux43eux440ux438ux442ux435ux43bux44cux43dux43e-11}

Нужно уметь сказать:

\begin{enumerate}
\def\labelenumi{\arabic{enumi}.}
\tightlist
\item
  Частичный порядок - рефлексивность, антисимметричность,
  транзитивность.
\item
  \(L^2\) - пространство случайных величин с конечным вторым моментом:
\end{enumerate}

\[
L^2=\{X:\mathbb EX^2<\infty\}.
\]

\begin{enumerate}
\def\labelenumi{\arabic{enumi}.}
\setcounter{enumi}{2}
\tightlist
\item
  Скалярное произведение:
\end{enumerate}

\[
\langle X,Y\rangle=\mathbb E[XY].
\]

\begin{enumerate}
\def\labelenumi{\arabic{enumi}.}
\setcounter{enumi}{3}
\tightlist
\item
  Для центрированных величин:
\end{enumerate}

\[
\operatorname{Cov}(X,Y)=\langle X,Y\rangle.
\]

\begin{enumerate}
\def\labelenumi{\arabic{enumi}.}
\setcounter{enumi}{4}
\tightlist
\item
  Проекция на замкнутое подпространство:
\end{enumerate}

\[
x-P_Hx\perp H.
\]

\begin{enumerate}
\def\labelenumi{\arabic{enumi}.}
\setcounter{enumi}{5}
\tightlist
\item
  Условное ожидание в \(L^2\):
\end{enumerate}

\[
\mathbb E[X\mid\mathcal G]
=
P_{L^2(\mathcal G)}X.
\]

\subsection{12. Минимум на
хорошо}\label{ux43cux438ux43dux438ux43cux443ux43c-ux43dux430-ux445ux43eux440ux43eux448ux43e-11}

Помимо минимума, нужно:

\begin{enumerate}
\def\labelenumi{\arabic{enumi}.}
\tightlist
\item
  объяснить, что \(L^2\) является гильбертовым пространством;
\item
  записать норму
\end{enumerate}

\[
\|X\|_2=(\mathbb EX^2)^{1/2};
\]

\begin{enumerate}
\def\labelenumi{\arabic{enumi}.}
\setcounter{enumi}{2}
\tightlist
\item
  сформулировать теорему о проекции на замкнутое подпространство;
\item
  записать нормальные уравнения для конечномерной проекции;
\item
  объяснить матрицу Грама и ее неотрицательную определенность;
\item
  привести пример проекции на константы:
\end{enumerate}

\[
P_HX=\mathbb EX.
\]

\subsection{13. Что добавить для
отлично}\label{ux447ux442ux43e-ux434ux43eux431ux430ux432ux438ux442ux44c-ux434ux43bux44f-ux43eux442ux43bux438ux447ux43dux43e-11}

Для сильного ответа стоит добавить:

\begin{enumerate}
\def\labelenumi{\arabic{enumi}.}
\tightlist
\item
  различие между частичным и линейным порядком;
\item
  примеры алгебраических структур в вероятности;
\item
  доказательную идею проекторной теоремы через минимизацию
  \(\|x-h\|^2\);
\item
  Пифагорову формулу для проекции;
\item
  свойства проектора: линейность, идемпотентность, самосопряженность;
\item
  условное ожидание как лучший среднеквадратичный прогноз;
\item
  предупреждение, что ортогональность не равна независимости.
\end{enumerate}

\subsection{14. Устная версия ответа на 5
минут}\label{ux443ux441ux442ux43dux430ux44f-ux432ux435ux440ux441ux438ux44f-ux43eux442ux432ux435ux442ux430-ux43dux430-5-ux43cux438ux43dux443ux442-11}

Сначала я называю общий язык. Упорядоченная структура - это множество с
отношением порядка. Частичный порядок рефлексивен, антисимметричен и
транзитивен. Примеры: включение событий \(A\subset B\), порядок
случайных величин \(X\le Y\) почти наверное, порядок на неотрицательно
определенных матрицах.

Алгебраическая структура - это множество с операциями. В этом билете
главный пример - линейное пространство случайных величин. Если
\(X,Y\in L^2\), то \(aX+bY\in L^2\).

Дальше перехожу к гильбертову пространству. Рассматриваем

\[
L^2(\Omega,\mathcal F,P)=\{X:\mathbb EX^2<\infty\}/\sim.
\]

Скалярное произведение:

\[
\langle X,Y\rangle=\mathbb E[XY],
\]

норма:

\[
\|X\|_2=(\mathbb EX^2)^{1/2}.
\]

По Коши-Буняковскому \(\mathbb E[XY]\) корректно определено.
Пространство \(L^2\) полно, поэтому это гильбертово пространство.

Для центрированных величин

\[
X^0=X-\mathbb EX,\qquad Y^0=Y-\mathbb EY
\]

получаем

\[
\operatorname{Cov}(X,Y)=\langle X^0,Y^0\rangle.
\]

Значит ковариация - это скалярное произведение центрированных случайных
величин.

Потом формулирую теорему о проекции. Если \(H\) - замкнутое
подпространство гильбертова пространства, то для любого \(x\) существует
единственный \(P_Hx\in H\), такой что

\[
x-P_Hx\perp H.
\]

Это ближайший элемент из \(H\):

\[
\|x-P_Hx\|=\min_{h\in H}\|x-h\|.
\]

Проектор линеен, идемпотентен и самосопряжен.

В конечномерном случае, если

\[
P_Hx=\sum_j a_jh_j,
\]

то коэффициенты находятся из нормальных уравнений:

\[
\sum_j a_j\langle h_j,h_i\rangle=\langle x,h_i\rangle.
\]

Матрица \(\langle h_i,h_j\rangle\) - матрица Грама. Она неотрицательно
определена.

Завершаю вероятностным смыслом. Если \(X\in L^2\) и
\(\mathcal G\subset\mathcal F\), то

\[
\mathbb E[X\mid\mathcal G]
=
P_{L^2(\mathcal G)}X.
\]

То есть условное ожидание - лучший среднеквадратичный прогноз \(X\) по
информации \(\mathcal G\):

\[
\mathbb E[(X-\mathbb E[X\mid\mathcal G])^2]
=
\min_{Y\in L^2(\mathcal G)}\mathbb E[(X-Y)^2].
\]

Главные ловушки: не путать ортогональность и независимость, не забывать
вторые моменты и замкнутость подпространства.

\subsection{15. Если осталось 2 минуты перед
экзаменом}\label{ux435ux441ux43bux438-ux43eux441ux442ux430ux43bux43eux441ux44c-2-ux43cux438ux43dux443ux442ux44b-ux43fux435ux440ux435ux434-ux44dux43aux437ux430ux43cux435ux43dux43eux43c-11}

Запомнить четыре формулы.

\textbf{\(L^2\) как гильбертово пространство:}

\[
\langle X,Y\rangle=\mathbb E[XY],
\qquad
\|X\|_2=(\mathbb EX^2)^{1/2}.
\]

\textbf{Ковариация как скалярное произведение:}

\[
\operatorname{Cov}(X,Y)
=
\langle X-\mathbb EX,\ Y-\mathbb EY\rangle.
\]

\textbf{Проектор:}

\[
P_Hx\in H,
\qquad
x-P_Hx\perp H.
\]

\textbf{Условное ожидание как проекция:}

\[
\mathbb E[X\mid\mathcal G]
=
P_{L^2(\mathcal G)}X.
\]

Главная ловушка: ортогональность/некоррелированность не означает
независимость; проекция гарантируется на замкнутое подпространство.

\newpage

\section{Билет 13. Связи математического ожидания и матрицы корреляции с
характеристической
функцией}\label{ux431ux438ux43bux435ux442-13.-ux441ux432ux44fux437ux438-ux43cux430ux442ux435ux43cux430ux442ux438ux447ux435ux441ux43aux43eux433ux43e-ux43eux436ux438ux434ux430ux43dux438ux44f-ux438-ux43cux430ux442ux440ux438ux446ux44b-ux43aux43eux440ux440ux435ux43bux44fux446ux438ux438-ux441-ux445ux430ux440ux430ux43aux442ux435ux440ux438ux441ux442ux438ux447ux435ux441ux43aux43eux439-ux444ux443ux43dux43aux446ux438ux435ux439}

Источники: Ширяев 1, гл. II, §12; лекции ДСП.

\subsection{1. Короткий
ответ}\label{ux43aux43eux440ux43eux442ux43aux438ux439-ux43eux442ux432ux435ux442-12}

Характеристическая функция конечномерного распределения кодирует не
только само распределение, но и моменты, если они существуют.

Для случайной величины \(X\):

\[
\varphi_X(t)=\mathbb E e^{itX}.
\]

Если \(\mathbb E|X|^k<\infty\), то \(\varphi_X\) имеет производные до
порядка \(k\), и

\[
\varphi_X^{(r)}(0)=i^r\mathbb E X^r,
\qquad r=1,\ldots,k.
\]

В частности,

\[
\mathbb EX=\frac{\varphi_X'(0)}{i}=-i\varphi_X'(0),
\]

а если существует второй момент, то

\[
\mathbb EX^2=-\varphi_X''(0),
\qquad
\operatorname{Var}X=-\varphi_X''(0)-(\mathbb EX)^2.
\]

Для случайного вектора \(X=(X_1,\ldots,X_n)^T\):

\[
\varphi_X(t)=\mathbb E e^{i\langle t,X\rangle},
\qquad t\in\mathbb R^n.
\]

Если \(\mathbb E|X_j|<\infty\), то

\[
\frac{\partial\varphi_X}{\partial t_j}(0)=i\mathbb EX_j.
\]

Если существуют вторые моменты, то

\[
\frac{\partial^2\varphi_X}{\partial t_i\partial t_j}(0)
=
-\mathbb E[X_iX_j].
\]

Значит средний вектор и матрица вторых моментов восстанавливаются из
производных характеристической функции в нуле:

\[
m_j=\mathbb EX_j=-i\partial_j\varphi_X(0),
\]

\[
R_{ij}=\mathbb E[X_iX_j]=-\partial_{ij}\varphi_X(0).
\]

Ковариационная матрица:

\[
\Sigma_{ij}
=
\operatorname{Cov}(X_i,X_j)
=
R_{ij}-m_im_j.
\]

Терминологически в курсах случайных процессов матрицу
\(R=(\mathbb E X_iX_j)\) часто называют корреляционной матрицей, а
\(\Sigma\) - ковариационной. Если величины центрированы, то
\(R=\Sigma\).

\subsection{2. Полный
ответ}\label{ux43fux43eux43bux43dux44bux439-ux43eux442ux432ux435ux442-12}

\subsubsection{2.1. Введение и связь с предыдущими
билетами}\label{ux432ux432ux435ux434ux435ux43dux438ux435-ux438-ux441ux432ux44fux437ux44c-ux441-ux43fux440ux435ux434ux44bux434ux443ux449ux438ux43cux438-ux431ux438ux43bux435ux442ux430ux43cux438}

Этот билет продолжает Билет 08 о характеристических функциях и Билет 11
о моментах через интегралы по распределению. В билете 8
характеристическая функция была введена как преобразование Фурье
конечномерного распределения:

\[
\varphi_X(t)
=
\mathbb E e^{i\langle t,X\rangle}
=
\int_{\mathbb R^n} e^{i\langle t,x\rangle}\,P_X(dx).
\]

Здесь нужно объяснить, как из этой функции получить математическое
ожидание, вторые моменты и корреляционную или ковариационную матрицу.

\subsubsection{2.2. Одномерный случай: моменты случайной
величины}\label{ux43eux434ux43dux43eux43cux435ux440ux43dux44bux439-ux441ux43bux443ux447ux430ux439-ux43cux43eux43cux435ux43dux442ux44b-ux441ux43bux443ux447ux430ux439ux43dux43eux439-ux432ux435ux43bux438ux447ux438ux43dux44b}

Начнем с одномерного случая. Пусть \(X\) - вещественная случайная
величина. Ее характеристическая функция:

\[
\varphi(t)=\mathbb E e^{itX}.
\]

Она существует всегда, потому что

\[
|e^{itX}|=1.
\]

Но производные характеристической функции и моменты требуют
дополнительных условий. Если \(\mathbb E|X|<\infty\), то можно
дифференцировать под знаком математического ожидания:

\[
\varphi'(t)=\mathbb E\left[iXe^{itX}\right].
\]

При \(t=0\):

\[
\varphi'(0)=i\mathbb EX.
\]

Поэтому

\[
\mathbb EX=\frac{\varphi'(0)}{i}=-i\varphi'(0).
\]

Если \(\mathbb EX^2<\infty\), то

\[
\varphi''(t)=\mathbb E\left[(iX)^2e^{itX}\right]
=
-\mathbb E\left[X^2e^{itX}\right],
\]

и при \(t=0\):

\[
\varphi''(0)=-\mathbb EX^2.
\]

Значит

\[
\mathbb EX^2=-\varphi''(0).
\]

Дисперсия:

\[
\operatorname{Var}X
=
\mathbb EX^2-(\mathbb EX)^2
=
-\varphi''(0)-(-i\varphi'(0))^2.
\]

Можно записать и проще:

\[
\operatorname{Var}X=-\varphi''(0)-(\mathbb EX)^2.
\]

Общая одномерная формула такая. Если

\[
\mathbb E|X|^k<\infty,
\]

то для \(r=1,\ldots,k\):

\[
\varphi^{(r)}(t)
=
\mathbb E\left[(iX)^r e^{itX}\right],
\]

а в нуле:

\[
\varphi^{(r)}(0)
=
i^r\mathbb EX^r.
\]

Отсюда

\[
\mathbb EX^r
=
\frac{\varphi^{(r)}(0)}{i^r}.
\]

\subsubsection{2.3. Случайный вектор: математическое ожидание и матрица
вторых
моментов}\label{ux441ux43bux443ux447ux430ux439ux43dux44bux439-ux432ux435ux43aux442ux43eux440-ux43cux430ux442ux435ux43cux430ux442ux438ux447ux435ux441ux43aux43eux435-ux43eux436ux438ux434ux430ux43dux438ux435-ux438-ux43cux430ux442ux440ux438ux446ux430-ux432ux442ux43eux440ux44bux445-ux43cux43eux43cux435ux43dux442ux43eux432}

Теперь случайный вектор. Пусть

\[
X=(X_1,\ldots,X_n)^T
\]

и \(P_X\) - его конечномерное распределение на \(\mathbb R^n\).
Характеристическая функция:

\[
\varphi(t)
=
\varphi_X(t_1,\ldots,t_n)
=
\mathbb E\exp\left(i\sum_{j=1}^n t_jX_j\right).
\]

Если компоненты интегрируемы, то частные производные первого порядка в
нуле дают математические ожидания:

\[
\partial_j\varphi(t)
=
\mathbb E\left[iX_j e^{i\langle t,X\rangle}\right],
\]

и

\[
\partial_j\varphi(0)=i\mathbb EX_j.
\]

Поэтому средний вектор

\[
m=\mathbb EX=(m_1,\ldots,m_n)^T
\]

получается по формуле

\[
m_j=-i\partial_j\varphi(0).
\]

Если существуют вторые моменты, то вторые частные производные дают
смешанные вторые моменты:

\[
\partial_{ij}\varphi(t)
=
\frac{\partial^2\varphi(t)}{\partial t_i\partial t_j}
=
\mathbb E\left[(iX_i)(iX_j)e^{i\langle t,X\rangle}\right]
=
-\mathbb E\left[X_iX_j e^{i\langle t,X\rangle}\right].
\]

При \(t=0\):

\[
\partial_{ij}\varphi(0)
=
-\mathbb E[X_iX_j].
\]

Обозначим матрицу вторых моментов:

\[
R=\mathbb E[XX^T],
\qquad
R_{ij}=\mathbb E[X_iX_j].
\]

Тогда

\[
R_{ij}=-\partial_{ij}\varphi(0).
\]

В терминах матрицы Гессе:

\[
R=-D^2\varphi(0),
\]

если \(D^2\varphi(0)\) понимается как матрица вторых частных
производных.

\subsubsection{2.4. Ковариационная матрица и
терминология}\label{ux43aux43eux432ux430ux440ux438ux430ux446ux438ux43eux43dux43dux430ux44f-ux43cux430ux442ux440ux438ux446ux430-ux438-ux442ux435ux440ux43cux438ux43dux43eux43bux43eux433ux438ux44f}

Ковариационная матрица:

\[
\Sigma
=
\operatorname{Cov}(X)
=
\mathbb E[(X-m)(X-m)^T].
\]

По формуле из Билета 11:

\[
\Sigma=R-mm^T.
\]

Следовательно, через характеристическую функцию:

\[
\Sigma_{ij}
=
-\partial_{ij}\varphi(0)-m_im_j.
\]

Если выразить \(m_i\) и \(m_j\) тоже через производные \(\varphi\), то

\[
\Sigma_{ij}
=
-\partial_{ij}\varphi(0)
+\partial_i\varphi(0)\partial_j\varphi(0).
\]

Проверка знака: так как \(\partial_i\varphi(0)=im_i\), то

\[
\partial_i\varphi(0)\partial_j\varphi(0)
=
-m_im_j.
\]

Поэтому последняя формула действительно равна

\[
-\partial_{ij}\varphi(0)-m_im_j.
\]

Отдельно полезно сказать про термин ``матрица корреляции''. В разных
курсах используются три близких объекта:

\begin{enumerate}
\def\labelenumi{\arabic{enumi}.}
\tightlist
\item
  матрица вторых моментов, или корреляционная матрица в теории случайных
  процессов:
\end{enumerate}

\[
R_{ij}=\mathbb E[X_iX_j];
\]

\begin{enumerate}
\def\labelenumi{\arabic{enumi}.}
\setcounter{enumi}{1}
\tightlist
\item
  ковариационная матрица:
\end{enumerate}

\[
\Sigma_{ij}=\operatorname{Cov}(X_i,X_j);
\]

\begin{enumerate}
\def\labelenumi{\arabic{enumi}.}
\setcounter{enumi}{2}
\tightlist
\item
  нормированная матрица коэффициентов корреляции:
\end{enumerate}

\[
\rho_{ij}
=
\frac{\operatorname{Cov}(X_i,X_j)}
{\sqrt{\operatorname{Var}X_i\,\operatorname{Var}X_j}},
\]

если дисперсии положительны.

В билете 13 обычно нужна связь характеристической функции именно с
математическим ожиданием и матрицей вторых моментов \(R\), а затем с
ковариационной матрицей \(\Sigma\). Если \(X\) центрирован, то

\[
m=0,
\qquad
\Sigma=R,
\]

и корреляционная матрица в смысле вторых моментов совпадает с
ковариационной.

\subsubsection{2.5. Приложение к конечномерным распределениям случайных
процессов}\label{ux43fux440ux438ux43bux43eux436ux435ux43dux438ux435-ux43a-ux43aux43eux43dux435ux447ux43dux43eux43cux435ux440ux43dux44bux43c-ux440ux430ux441ux43fux440ux435ux434ux435ux43bux435ux43dux438ux44fux43c-ux441ux43bux443ux447ux430ux439ux43dux44bux445-ux43fux440ux43eux446ux435ux441ux441ux43eux432}

Для конечномерного распределения случайного процесса запись такая же.
Если выбран набор моментов времени

\[
t_1,\ldots,t_n,
\]

и

\[
X=(X_{t_1},\ldots,X_{t_n})^T,
\]

то его конечномерная характеристическая функция:

\[
\varphi_{t_1,\ldots,t_n}(u_1,\ldots,u_n)
=
\mathbb E\exp\left(i\sum_{j=1}^n u_jX_{t_j}\right).
\]

Тогда

\[
\mathbb EX_{t_j}
=
-i\frac{\partial}{\partial u_j}
\varphi_{t_1,\ldots,t_n}(0),
\]

а корреляционная функция на выбранном конечном наборе индексов:

\[
R(t_i,t_j)
=
\mathbb E[X_{t_i}X_{t_j}]
=
-\frac{\partial^2}{\partial u_i\partial u_j}
\varphi_{t_1,\ldots,t_n}(0).
\]

Ковариационная функция:

\[
K(t_i,t_j)
=
\operatorname{Cov}(X_{t_i},X_{t_j})
=
R(t_i,t_j)-m(t_i)m(t_j).
\]

В стационарных задачах второго порядка это переходит в билеты Билет 17,
Билет 24 и Билет 25: среднее постоянно, а корреляция или ковариация
зависит от лага.

\subsubsection{2.6. Логарифм характеристической функции и
кумулянты}\label{ux43bux43eux433ux430ux440ux438ux444ux43c-ux445ux430ux440ux430ux43aux442ux435ux440ux438ux441ux442ux438ux447ux435ux441ux43aux43eux439-ux444ux443ux43dux43aux446ux438ux438-ux438-ux43aux443ux43cux443ux43bux44fux43dux442ux44b}

Иногда удобнее использовать логарифм характеристической функции. Так как
\(\varphi(0)=1\), в некоторой окрестности нуля при достаточной гладкости
можно рассматривать

\[
\psi(t)=\log\varphi(t).
\]

Тогда первые производные \(\psi\) дают среднее, а вторые - ковариации:

\[
\partial_j\psi(0)=i m_j,
\]

\[
\partial_{ij}\psi(0)=-\Sigma_{ij}.
\]

Это форма через кумулянты: первая кумулянта - среднее, вторая кумулянта
- ковариация. Для экзамена достаточно знать основную формулу через
производные \(\varphi\); логарифм можно добавить как сильное замечание.

\subsubsection{2.7. Главная идея
билета}\label{ux433ux43bux430ux432ux43dux430ux44f-ux438ux434ux435ux44f-ux431ux438ux43bux435ux442ux430}

Главная идея всего билета: характеристическая функция является
преобразованием Фурье распределения. Дифференцирование по параметрам
\(t_j\) выносит из экспоненты множители \(iX_j\). Поэтому производные в
нуле превращаются в моменты.

\subsection{3. Основные
определения}\label{ux43eux441ux43dux43eux432ux43dux44bux435-ux43eux43fux440ux435ux434ux435ux43bux435ux43dux438ux44f-12}

\subsubsection{Характеристическая функция случайной
величины}\label{ux445ux430ux440ux430ux43aux442ux435ux440ux438ux441ux442ux438ux447ux435ux441ux43aux430ux44f-ux444ux443ux43dux43aux446ux438ux44f-ux441ux43bux443ux447ux430ux439ux43dux43eux439-ux432ux435ux43bux438ux447ux438ux43dux44b-1}

\[
\varphi_X(t)=\mathbb E e^{itX},
\qquad t\in\mathbb R.
\]

\subsubsection{Характеристическая функция случайного
вектора}\label{ux445ux430ux440ux430ux43aux442ux435ux440ux438ux441ux442ux438ux447ux435ux441ux43aux430ux44f-ux444ux443ux43dux43aux446ux438ux44f-ux441ux43bux443ux447ux430ux439ux43dux43eux433ux43e-ux432ux435ux43aux442ux43eux440ux430-1}

\[
\varphi_X(t)=\mathbb E e^{i\langle t,X\rangle},
\qquad t\in\mathbb R^n.
\]

\subsubsection{Характеристическая функция конечномерного распределения
процесса}\label{ux445ux430ux440ux430ux43aux442ux435ux440ux438ux441ux442ux438ux447ux435ux441ux43aux430ux44f-ux444ux443ux43dux43aux446ux438ux44f-ux43aux43eux43dux435ux447ux43dux43eux43cux435ux440ux43dux43eux433ux43e-ux440ux430ux441ux43fux440ux435ux434ux435ux43bux435ux43dux438ux44f-ux43fux440ux43eux446ux435ux441ux441ux430}

Для процесса \((X_s)_{s\in T}\) и индексов \(s_1,\ldots,s_n\):

\[
\varphi_{s_1,\ldots,s_n}(u)
=
\mathbb E\exp\left(i\sum_{j=1}^n u_jX_{s_j}\right).
\]

\subsubsection{Математическое ожидание
вектора}\label{ux43cux430ux442ux435ux43cux430ux442ux438ux447ux435ux441ux43aux43eux435-ux43eux436ux438ux434ux430ux43dux438ux435-ux432ux435ux43aux442ux43eux440ux430}

\[
m=\mathbb EX,
\qquad
m_j=\mathbb EX_j.
\]

Оно существует, если все компоненты \(X_j\) интегрируемы.

\subsubsection{Матрица вторых
моментов}\label{ux43cux430ux442ux440ux438ux446ux430-ux432ux442ux43eux440ux44bux445-ux43cux43eux43cux435ux43dux442ux43eux432-1}

\[
R=\mathbb E[XX^T],
\qquad
R_{ij}=\mathbb E[X_iX_j].
\]

В теории случайных процессов ее часто называют корреляционной матрицей.

\subsubsection{Ковариационная
матрица}\label{ux43aux43eux432ux430ux440ux438ux430ux446ux438ux43eux43dux43dux430ux44f-ux43cux430ux442ux440ux438ux446ux430-1}

\[
\Sigma=\operatorname{Cov}(X)=\mathbb E[(X-m)(X-m)^T],
\]

то есть

\[
\Sigma_{ij}
=
\operatorname{Cov}(X_i,X_j)
=
\mathbb E[(X_i-m_i)(X_j-m_j)].
\]

\subsubsection{\texorpdfstring{Связь матриц \(R\) и
\(\Sigma\)}{Связь матриц R и \textbackslash Sigma}}\label{ux441ux432ux44fux437ux44c-ux43cux430ux442ux440ux438ux446-r-ux438-sigma}

\[
\Sigma=R-mm^T.
\]

Если \(m=0\), то

\[
\Sigma=R.
\]

\subsubsection{Нормированная матрица коэффициентов
корреляции}\label{ux43dux43eux440ux43cux438ux440ux43eux432ux430ux43dux43dux430ux44f-ux43cux430ux442ux440ux438ux446ux430-ux43aux43eux44dux444ux444ux438ux446ux438ux435ux43dux442ux43eux432-ux43aux43eux440ux440ux435ux43bux44fux446ux438ux438}

Если \(\operatorname{Var}X_i>0\), то

\[
\rho_{ij}
=
\frac{\Sigma_{ij}}
{\sqrt{\Sigma_{ii}\Sigma_{jj}}}.
\]

Диагональные элементы равны \(1\), а \(|\rho_{ij}|\le1\).

\subsubsection{Мультииндексная форма
момента}\label{ux43cux443ux43bux44cux442ux438ux438ux43dux434ux435ux43aux441ux43dux430ux44f-ux444ux43eux440ux43cux430-ux43cux43eux43cux435ux43dux442ux430}

Для \(\alpha=(\alpha_1,\ldots,\alpha_n)\):

\[
|\alpha|=\alpha_1+\cdots+\alpha_n,
\]

\[
\partial^\alpha
=
\frac{\partial^{|\alpha|}}
{\partial t_1^{\alpha_1}\cdots\partial t_n^{\alpha_n}},
\]

и при существовании соответствующего абсолютного момента:

\[
\partial^\alpha\varphi(0)
=
i^{|\alpha|}
\mathbb E\left[X_1^{\alpha_1}\cdots X_n^{\alpha_n}\right].
\]

\subsection{4. Основные теоремы и
свойства}\label{ux43eux441ux43dux43eux432ux43dux44bux435-ux442ux435ux43eux440ux435ux43cux44b-ux438-ux441ux432ux43eux439ux441ux442ux432ux430-12}

\subsubsection{Теорема 1. Моменты одномерной величины через производные
характеристической
функции}\label{ux442ux435ux43eux440ux435ux43cux430-1.-ux43cux43eux43cux435ux43dux442ux44b-ux43eux434ux43dux43eux43cux435ux440ux43dux43eux439-ux432ux435ux43bux438ux447ux438ux43dux44b-ux447ux435ux440ux435ux437-ux43fux440ux43eux438ux437ux432ux43eux434ux43dux44bux435-ux445ux430ux440ux430ux43aux442ux435ux440ux438ux441ux442ux438ux447ux435ux441ux43aux43eux439-ux444ux443ux43dux43aux446ux438ux438}

Если

\[
\mathbb E|X|^k<\infty,
\]

то \(\varphi_X\) имеет непрерывные производные до порядка \(k\), и

\[
\varphi_X^{(r)}(t)
=
\mathbb E[(iX)^r e^{itX}],
\qquad r=1,\ldots,k.
\]

В частности,

\[
\varphi_X^{(r)}(0)=i^r\mathbb EX^r.
\]

\subsubsection{Следствие. Среднее и
дисперсия}\label{ux441ux43bux435ux434ux441ux442ux432ux438ux435.-ux441ux440ux435ux434ux43dux435ux435-ux438-ux434ux438ux441ux43fux435ux440ux441ux438ux44f}

Если \(\mathbb E|X|<\infty\), то

\[
\mathbb EX=-i\varphi_X'(0).
\]

Если \(\mathbb EX^2<\infty\), то

\[
\mathbb EX^2=-\varphi_X''(0),
\]

\[
\operatorname{Var}X=-\varphi_X''(0)-(-i\varphi_X'(0))^2.
\]

\subsubsection{Теорема 2. Моменты вектора через частные
производные}\label{ux442ux435ux43eux440ux435ux43cux430-2.-ux43cux43eux43cux435ux43dux442ux44b-ux432ux435ux43aux442ux43eux440ux430-ux447ux435ux440ux435ux437-ux447ux430ux441ux442ux43dux44bux435-ux43fux440ux43eux438ux437ux432ux43eux434ux43dux44bux435}

Пусть \(X=(X_1,\ldots,X_n)^T\) и существует момент

\[
\mathbb E|X_1|^{\alpha_1}\cdots |X_n|^{\alpha_n}<\infty.
\]

Тогда

\[
\partial^\alpha\varphi_X(0)
=
i^{|\alpha|}
\mathbb E\left[X_1^{\alpha_1}\cdots X_n^{\alpha_n}\right].
\]

Для первого и второго порядков:

\[
\partial_j\varphi(0)=i\mathbb EX_j,
\]

\[
\partial_{ij}\varphi(0)=-\mathbb E[X_iX_j].
\]

\subsubsection{Свойство 3. Матрица вторых моментов через матрицу Гессе
характеристической
функции}\label{ux441ux432ux43eux439ux441ux442ux432ux43e-3.-ux43cux430ux442ux440ux438ux446ux430-ux432ux442ux43eux440ux44bux445-ux43cux43eux43cux435ux43dux442ux43eux432-ux447ux435ux440ux435ux437-ux43cux430ux442ux440ux438ux446ux443-ux433ux435ux441ux441ux435-ux445ux430ux440ux430ux43aux442ux435ux440ux438ux441ux442ux438ux447ux435ux441ux43aux43eux439-ux444ux443ux43dux43aux446ux438ux438}

Если все вторые моменты существуют, то

\[
R=-D^2\varphi(0),
\]

то есть

\[
R_{ij}=-\partial_{ij}\varphi(0).
\]

\subsubsection{Свойство 4. Ковариационная матрица через производные
характеристической
функции}\label{ux441ux432ux43eux439ux441ux442ux432ux43e-4.-ux43aux43eux432ux430ux440ux438ux430ux446ux438ux43eux43dux43dux430ux44f-ux43cux430ux442ux440ux438ux446ux430-ux447ux435ux440ux435ux437-ux43fux440ux43eux438ux437ux432ux43eux434ux43dux44bux435-ux445ux430ux440ux430ux43aux442ux435ux440ux438ux441ux442ux438ux447ux435ux441ux43aux43eux439-ux444ux443ux43dux43aux446ux438ux438}

Если все вторые моменты существуют, то

\[
\Sigma_{ij}
=
-\partial_{ij}\varphi(0)
+\partial_i\varphi(0)\partial_j\varphi(0).
\]

Эквивалентно:

\[
\Sigma=-D^2\varphi(0)-mm^T,
\qquad
m_j=-i\partial_j\varphi(0).
\]

\subsubsection{Свойство 5. Через логарифм характеристической
функции}\label{ux441ux432ux43eux439ux441ux442ux432ux43e-5.-ux447ux435ux440ux435ux437-ux43bux43eux433ux430ux440ux438ux444ux43c-ux445ux430ux440ux430ux43aux442ux435ux440ux438ux441ux442ux438ux447ux435ux441ux43aux43eux439-ux444ux443ux43dux43aux446ux438ux438}

Если \(\psi(t)=\log\varphi(t)\) определена и дважды дифференцируема в
окрестности нуля, то

\[
\partial_j\psi(0)=im_j,
\]

\[
\partial_{ij}\psi(0)=-\Sigma_{ij}.
\]

\subsubsection{Свойство 6. Неотрицательная определенность матриц второго
порядка}\label{ux441ux432ux43eux439ux441ux442ux432ux43e-6.-ux43dux435ux43eux442ux440ux438ux446ux430ux442ux435ux43bux44cux43dux430ux44f-ux43eux43fux440ux435ux434ux435ux43bux435ux43dux43dux43eux441ux442ux44c-ux43cux430ux442ux440ux438ux446-ux432ux442ux43eux440ux43eux433ux43e-ux43fux43eux440ux44fux434ux43aux430}

Матрица вторых моментов \(R\) неотрицательно определена:

\[
a^TRa
=
\mathbb E(a^TX)^2
\ge0.
\]

Ковариационная матрица \(\Sigma\) тоже неотрицательно определена:

\[
a^T\Sigma a
=
\operatorname{Var}(a^TX)
\ge0.
\]

Это согласуется с Билетом 12: \(R\) и \(\Sigma\) являются матрицами
Грама соответствующих случайных величин в \(L^2\).

\subsubsection{Свойство 7. Независимые
координаты}\label{ux441ux432ux43eux439ux441ux442ux432ux43e-7.-ux43dux435ux437ux430ux432ux438ux441ux438ux43cux44bux435-ux43aux43eux43eux440ux434ux438ux43dux430ux442ux44b}

Если \(X_1,\ldots,X_n\) независимы, то

\[
\varphi_X(t_1,\ldots,t_n)
=
\prod_{j=1}^n\varphi_{X_j}(t_j).
\]

Тогда смешанные моменты, если они существуют, раскладываются. Например:

\[
\mathbb E[X_iX_j]=\mathbb EX_i\,\mathbb EX_j,
\qquad i\ne j.
\]

Поэтому

\[
\operatorname{Cov}(X_i,X_j)=0
\]

для \(i\ne j\), если вторые моменты конечны.

\subsubsection{Свойство 8. Для гауссовского
вектора}\label{ux441ux432ux43eux439ux441ux442ux432ux43e-8.-ux434ux43bux44f-ux433ux430ux443ux441ux441ux43eux432ux441ux43aux43eux433ux43e-ux432ux435ux43aux442ux43eux440ux430}

Если

\[
X\sim N(m,\Sigma),
\]

то

\[
\varphi_X(t)
=
\exp\left(i\langle t,m\rangle-\frac12 t^T\Sigma t\right).
\]

Дифференцирование в нуле дает тот же средний вектор \(m\) и
ковариационную матрицу \(\Sigma\). Это связь с Билетом 18.

\subsection{5. Примеры}\label{ux43fux440ux438ux43cux435ux440ux44b-12}

\subsubsection{Пример 1. Вырожденная случайная
величина}\label{ux43fux440ux438ux43cux435ux440-1.-ux432ux44bux440ux43eux436ux434ux435ux43dux43dux430ux44f-ux441ux43bux443ux447ux430ux439ux43dux430ux44f-ux432ux435ux43bux438ux447ux438ux43dux430}

Пусть \(X=c\) почти наверное. Тогда

\[
\varphi(t)=e^{itc}.
\]

Производная:

\[
\varphi'(t)=ice^{itc},
\qquad
\varphi'(0)=ic.
\]

Значит

\[
\mathbb EX=-i\varphi'(0)=c.
\]

Вторая производная:

\[
\varphi''(0)=-c^2,
\]

поэтому

\[
\mathbb EX^2=-\varphi''(0)=c^2,
\]

и

\[
\operatorname{Var}X=c^2-c^2=0.
\]

\subsubsection{Пример 2. Бернуллиевская случайная
величина}\label{ux43fux440ux438ux43cux435ux440-2.-ux431ux435ux440ux43dux443ux43bux43bux438ux435ux432ux441ux43aux430ux44f-ux441ux43bux443ux447ux430ux439ux43dux430ux44f-ux432ux435ux43bux438ux447ux438ux43dux430}

Пусть

\[
P\{X=1\}=p,
\qquad
P\{X=0\}=q,
\qquad
p+q=1.
\]

Тогда

\[
\varphi(t)=q+pe^{it}.
\]

Первая производная:

\[
\varphi'(t)=ipe^{it},
\qquad
\varphi'(0)=ip.
\]

Значит

\[
\mathbb EX=p.
\]

Вторая производная:

\[
\varphi''(t)=-pe^{it},
\qquad
\varphi''(0)=-p.
\]

Поэтому

\[
\mathbb EX^2=p.
\]

Дисперсия:

\[
\operatorname{Var}X=p-p^2=pq.
\]

\subsubsection{Пример 3. Нормальная случайная
величина}\label{ux43fux440ux438ux43cux435ux440-3.-ux43dux43eux440ux43cux430ux43bux44cux43dux430ux44f-ux441ux43bux443ux447ux430ux439ux43dux430ux44f-ux432ux435ux43bux438ux447ux438ux43dux430}

Если

\[
X\sim N(m,\sigma^2),
\]

то

\[
\varphi(t)=\exp\left(imt-\frac12\sigma^2t^2\right).
\]

Логарифм:

\[
\psi(t)=\log\varphi(t)=imt-\frac12\sigma^2t^2.
\]

Тогда

\[
\psi'(0)=im,
\qquad
\psi''(0)=-\sigma^2.
\]

Значит

\[
\mathbb EX=m,
\qquad
\operatorname{Var}X=\sigma^2.
\]

\subsubsection{Пример 4. Дискретный двумерный
вектор}\label{ux43fux440ux438ux43cux435ux440-4.-ux434ux438ux441ux43aux440ux435ux442ux43dux44bux439-ux434ux432ux443ux43cux435ux440ux43dux44bux439-ux432ux435ux43aux442ux43eux440}

Пусть \(X=(X_1,X_2)\) принимает значения

\[
(0,0),\ (1,0),\ (0,1)
\]

с вероятностями \(p_0,p_1,p_2\), где \(p_0+p_1+p_2=1\). Тогда

\[
\varphi(t_1,t_2)
=
p_0+p_1e^{it_1}+p_2e^{it_2}.
\]

Средние:

\[
\partial_1\varphi(0)=ip_1,
\qquad
\partial_2\varphi(0)=ip_2,
\]

поэтому

\[
\mathbb EX_1=p_1,
\qquad
\mathbb EX_2=p_2.
\]

Смешанный второй момент:

\[
\partial_{12}\varphi(0)=0,
\]

значит

\[
\mathbb E[X_1X_2]=0.
\]

Ковариация:

\[
\operatorname{Cov}(X_1,X_2)
=
0-p_1p_2
=
-p_1p_2.
\]

Это хороший пример того, что смешанный второй момент и ковариация - не
одно и то же.

\subsubsection{Пример 5. Независимые
координаты}\label{ux43fux440ux438ux43cux435ux440-5.-ux43dux435ux437ux430ux432ux438ux441ux438ux43cux44bux435-ux43aux43eux43eux440ux434ux438ux43dux430ux442ux44b}

Пусть \(X_1\) и \(X_2\) независимы. Тогда

\[
\varphi_{(X_1,X_2)}(t_1,t_2)
=
\varphi_{X_1}(t_1)\varphi_{X_2}(t_2).
\]

Если существуют вторые моменты, то

\[
\mathbb E[X_1X_2]
=
\mathbb EX_1\,\mathbb EX_2,
\]

и поэтому

\[
\operatorname{Cov}(X_1,X_2)=0.
\]

Но обратное неверно: нулевая ковариация не означает независимость.

\subsubsection{Пример 6. Конечномерное распределение
процесса}\label{ux43fux440ux438ux43cux435ux440-6.-ux43aux43eux43dux435ux447ux43dux43eux43cux435ux440ux43dux43eux435-ux440ux430ux441ux43fux440ux435ux434ux435ux43bux435ux43dux438ux435-ux43fux440ux43eux446ux435ux441ux441ux430}

Пусть есть процесс \((X_t)\) и выбраны \(s,t\). Тогда

\[
\varphi_{s,t}(u,v)
=
\mathbb E e^{i(uX_s+vX_t)}.
\]

Если существуют вторые моменты, то

\[
\mathbb EX_s=-i\partial_u\varphi_{s,t}(0,0),
\]

\[
\mathbb E[X_sX_t]
=
-\partial_{uv}\varphi_{s,t}(0,0).
\]

Это и есть связь конечномерной характеристической функции с
корреляционной функцией процесса.

\subsection{6. Схемы доказательств /
обоснований}\label{ux441ux445ux435ux43cux44b-ux434ux43eux43aux430ux437ux430ux442ux435ux43bux44cux441ux442ux432-ux43eux431ux43eux441ux43dux43eux432ux430ux43dux438ux439-5}

\textbf{Почему производная дает математическое ожидание.}

Пусть \(\mathbb E|X|<\infty\). Для фиксированного \(t\) формально:

\[
\frac{d}{dt}e^{itX}=iXe^{itX}.
\]

Чтобы вынести производную под знак ожидания, используем доминирование.
Например, разностное отношение:

\[
\frac{e^{i(t+h)X}-e^{itX}}{h}
=
e^{itX}\frac{e^{ihX}-1}{h}.
\]

Поточечно при \(h\to0\) оно стремится к

\[
iXe^{itX}.
\]

Кроме того,

\[
\left|\frac{e^{ihX}-1}{h}\right|
\le |X|
\]

из оценки \(|e^{iy}-1|\le |y|\). Так как \(\mathbb E|X|<\infty\), можно
применить теорему Лебега о доминированной сходимости. Получаем

\[
\varphi'(t)=\mathbb E[iXe^{itX}].
\]

При \(t=0\):

\[
\varphi'(0)=i\mathbb EX.
\]

\textbf{Почему вторая производная дает второй момент.}

Если \(\mathbb EX^2<\infty\), то аналогично дифференцируем еще раз:

\[
\varphi''(t)
=
\mathbb E[(iX)^2e^{itX}]
=
-\mathbb E[X^2e^{itX}].
\]

При \(t=0\):

\[
\varphi''(0)=-\mathbb EX^2.
\]

Знак минус появляется из-за \(i^2=-1\).

\textbf{Доказательная схема для вектора.}

Для \(X=(X_1,\ldots,X_n)\):

\[
\varphi(t)=\mathbb E e^{i(t_1X_1+\cdots+t_nX_n)}.
\]

Дифференцирование по \(t_j\) выносит множитель \(iX_j\):

\[
\partial_j\varphi(t)
=
\mathbb E[iX_je^{i\langle t,X\rangle}].
\]

Дифференцирование по \(t_i\) и \(t_j\) выносит множитель

\[
(iX_i)(iX_j)=-X_iX_j.
\]

Отсюда:

\[
\partial_{ij}\varphi(0)=-\mathbb E[X_iX_j].
\]

\textbf{Почему ковариационная матрица неотрицательно определена.}

Для любого \(a\in\mathbb R^n\):

\[
a^T\Sigma a
=
\sum_{i,j}a_i a_j\operatorname{Cov}(X_i,X_j)
=
\operatorname{Var}\left(\sum_i a_iX_i\right)
\ge0.
\]

Это также следует из геометрии \(L^2\): ковариационная матрица - матрица
Грама центрированных величин.

\textbf{Почему через логарифм получается ковариация.}

Пусть

\[
\psi=\log\varphi.
\]

Тогда

\[
\partial_i\psi=\frac{\partial_i\varphi}{\varphi}.
\]

В нуле \(\varphi(0)=1\), поэтому

\[
\partial_i\psi(0)=\partial_i\varphi(0)=im_i.
\]

Для второй производной:

\[
\partial_{ij}\psi(0)
=
\partial_{ij}\varphi(0)
-\partial_i\varphi(0)\partial_j\varphi(0).
\]

Подставляем:

\[
\partial_{ij}\varphi(0)=-\mathbb E[X_iX_j],
\]

\[
\partial_i\varphi(0)\partial_j\varphi(0)
=
(im_i)(im_j)
=
-m_im_j.
\]

Получаем

\[
\partial_{ij}\psi(0)
=
-\mathbb E[X_iX_j]+m_im_j
=
-\operatorname{Cov}(X_i,X_j).
\]

\subsection{7. Что написать на
доске}\label{ux447ux442ux43e-ux43dux430ux43fux438ux441ux430ux442ux44c-ux43dux430-ux434ux43eux441ux43aux435-12}

\begin{enumerate}
\def\labelenumi{\arabic{enumi}.}
\tightlist
\item
  Характеристическая функция:
\end{enumerate}

\[
\varphi_X(t)=\mathbb E e^{i\langle t,X\rangle}.
\]

\begin{enumerate}
\def\labelenumi{\arabic{enumi}.}
\setcounter{enumi}{1}
\tightlist
\item
  Одномерные моменты:
\end{enumerate}

\[
\varphi_X^{(k)}(0)=i^k\mathbb EX^k.
\]

\begin{enumerate}
\def\labelenumi{\arabic{enumi}.}
\setcounter{enumi}{2}
\tightlist
\item
  Средний вектор:
\end{enumerate}

\[
m_j=\mathbb EX_j=-i\partial_j\varphi(0).
\]

\begin{enumerate}
\def\labelenumi{\arabic{enumi}.}
\setcounter{enumi}{3}
\tightlist
\item
  Матрица вторых моментов:
\end{enumerate}

\[
R_{ij}=\mathbb E[X_iX_j]=-\partial_{ij}\varphi(0).
\]

\begin{enumerate}
\def\labelenumi{\arabic{enumi}.}
\setcounter{enumi}{4}
\tightlist
\item
  Ковариационная матрица:
\end{enumerate}

\[
\Sigma_{ij}=R_{ij}-m_im_j.
\]

\begin{enumerate}
\def\labelenumi{\arabic{enumi}.}
\setcounter{enumi}{5}
\tightlist
\item
  Через производные \(\varphi\):
\end{enumerate}

\[
\Sigma_{ij}
=
-\partial_{ij}\varphi(0)
+\partial_i\varphi(0)\partial_j\varphi(0).
\]

\begin{enumerate}
\def\labelenumi{\arabic{enumi}.}
\setcounter{enumi}{6}
\tightlist
\item
  Через логарифм:
\end{enumerate}

\[
\partial_{ij}\log\varphi(0)=-\Sigma_{ij}.
\]

\begin{enumerate}
\def\labelenumi{\arabic{enumi}.}
\setcounter{enumi}{7}
\tightlist
\item
  Для процесса:
\end{enumerate}

\[
R(s,t)=\mathbb E[X_sX_t]
=
-\partial_{uv}\varphi_{s,t}(0,0).
\]

\subsection{8. Типичные дополнительные
вопросы}\label{ux442ux438ux43fux438ux447ux43dux44bux435-ux434ux43eux43fux43eux43bux43dux438ux442ux435ux43bux44cux43dux44bux435-ux432ux43eux43fux440ux43eux441ux44b-12}

\textbf{Вопрос. Характеристическая функция всегда существует?}

Да. Для любого \(t\):

\[
|e^{i\langle t,X\rangle}|=1,
\]

поэтому математическое ожидание существует как ожидание ограниченной
комплексной случайной величины.

\textbf{Вопрос. Если характеристическая функция существует всегда,
значит ли это, что все моменты существуют?}

Нет. Существование \(\varphi\) не означает существования моментов.
Например, у распределения Коши характеристическая функция существует, но
математическое ожидание не существует.

\textbf{Вопрос. Что нужно для формулы \(\varphi'(0)=i\mathbb EX\)?}

Достаточно \(\mathbb E|X|<\infty\). Тогда можно дифференцировать под
знаком ожидания.

\textbf{Вопрос. Что нужно для формулы \(\varphi''(0)=-\mathbb EX^2\)?}

Достаточно \(\mathbb EX^2<\infty\).

\textbf{Вопрос. Почему во второй производной стоит минус?}

Потому что

\[
\frac{d^2}{dt^2}e^{itX}
=(iX)^2e^{itX}
=-X^2e^{itX}.
\]

\textbf{Вопрос. Как получить математическое ожидание вектора?}

По компонентам:

\[
\mathbb EX_j=-i\partial_j\varphi(0).
\]

\textbf{Вопрос. Как получить смешанный момент \(\mathbb E[X_iX_j]\)?}

Через смешанную вторую производную:

\[
\mathbb E[X_iX_j]=-\partial_{ij}\varphi(0).
\]

\textbf{Вопрос. Чем отличается матрица вторых моментов от ковариационной
матрицы?}

Матрица вторых моментов:

\[
R_{ij}=\mathbb E[X_iX_j].
\]

Ковариационная матрица:

\[
\Sigma_{ij}=\mathbb E[(X_i-m_i)(X_j-m_j)].
\]

Связь:

\[
\Sigma=R-mm^T.
\]

\textbf{Вопрос. Когда корреляционная матрица совпадает с
ковариационной?}

Когда вектор центрирован:

\[
\mathbb EX=0.
\]

Тогда

\[
R=\Sigma.
\]

\textbf{Вопрос. Что дает логарифм характеристической функции?}

Он дает кумулянты. Первая производная в нуле дает среднее, а вторая -
ковариационную матрицу:

\[
\partial_j\log\varphi(0)=im_j,
\qquad
\partial_{ij}\log\varphi(0)=-\Sigma_{ij}.
\]

\textbf{Вопрос. Как это связано с гауссовским распределением?}

Для гауссовского вектора

\[
\varphi(t)=\exp\left(i\langle t,m\rangle-\frac12t^T\Sigma t\right).
\]

Из логарифма сразу видны \(m\) и \(\Sigma\).

\textbf{Вопрос. Как это связано со случайными процессами?}

Если взять конечномерное распределение

\[
(X_{t_1},\ldots,X_{t_n}),
\]

то производные его характеристической функции в нуле дают средние
\(m(t_i)\) и корреляции \(R(t_i,t_j)\).

\subsection{9. Типичные
ошибки}\label{ux442ux438ux43fux438ux447ux43dux44bux435-ux43eux448ux438ux431ux43aux438-12}

\begin{enumerate}
\def\labelenumi{\arabic{enumi}.}
\tightlist
\item
  Писать \(\varphi''(0)=\mathbb EX^2\) без минуса.
\item
  Забывать множитель \(i\) в первой производной.
\item
  Считать, что характеристическая функция существует только при
  существовании моментов.
\item
  Считать, что существование характеристической функции автоматически
  дает существование всех моментов.
\item
  Путать матрицу вторых моментов \(R=\mathbb E[XX^T]\) и ковариационную
  матрицу \(\Sigma=R-mm^T\).
\item
  Не центрировать величины при переходе от вторых моментов к
  ковариациям.
\item
  Путать корреляционную матрицу в смысле \(R_{ij}=\mathbb E[X_iX_j]\) с
  нормированной матрицей коэффициентов корреляции \(\rho_{ij}\).
\item
  Забывать, что формулы с производными в нуле требуют существования
  соответствующих моментов.
\item
  Дифференцировать по одному аргументу, когда нужен смешанный момент
  \(\mathbb E[X_iX_j]\).
\item
  Делать вывод о независимости из нулевой ковариации.
\end{enumerate}

\subsection{10. Связи с другими
билетами}\label{ux441ux432ux44fux437ux438-ux441-ux434ux440ux443ux433ux438ux43cux438-ux431ux438ux43bux435ux442ux430ux43cux438-12}

\textbf{Билет 08.} Там вводятся характеристические функции конечномерных
распределений и их основные свойства; здесь используются производные
этих функций.

\textbf{Билет 11.} Моменты, дисперсии и ковариации выражаются через
интегралы по распределению; характеристическая функция является другим
способом кодировать те же моменты.

\textbf{Билет 12.} Ковариационная матрица - матрица Грама центрированных
случайных величин в \(L^2\).

\textbf{Билет 14.} Коэффициенты корреляции получаются из ковариационной
матрицы нормировкой на стандартные отклонения.

\textbf{Билет 15.} Согласованные конечномерные распределения можно
задавать через согласованные конечномерные характеристические функции.

\textbf{Билет 18.} Гауссовское распределение полностью задается средним
и ковариационной матрицей, которые видны в характеристической функции.

\textbf{Билет 24.} Корреляционная функция процесса - это второй момент
или ковариация, получаемая из конечномерных характеристических функций.

\textbf{Билет 25.} Спектральные характеристики связаны с корреляционной
функцией как с коэффициентами Фурье.

\subsection{11. Минимум на
удовлетворительно}\label{ux43cux438ux43dux438ux43cux443ux43c-ux43dux430-ux443ux434ux43eux432ux43bux435ux442ux432ux43eux440ux438ux442ux435ux43bux44cux43dux43e-12}

Нужно уметь сказать:

\begin{enumerate}
\def\labelenumi{\arabic{enumi}.}
\tightlist
\item
  Характеристическая функция:
\end{enumerate}

\[
\varphi_X(t)=\mathbb E e^{i\langle t,X\rangle}.
\]

\begin{enumerate}
\def\labelenumi{\arabic{enumi}.}
\setcounter{enumi}{1}
\item
  Если существуют моменты, то производные в нуле дают эти моменты.
\item
  Для случайной величины:
\end{enumerate}

\[
\varphi'(0)=i\mathbb EX,
\qquad
\varphi''(0)=-\mathbb EX^2.
\]

\begin{enumerate}
\def\labelenumi{\arabic{enumi}.}
\setcounter{enumi}{3}
\tightlist
\item
  Для вектора:
\end{enumerate}

\[
\partial_j\varphi(0)=i\mathbb EX_j,
\qquad
\partial_{ij}\varphi(0)=-\mathbb E[X_iX_j].
\]

\begin{enumerate}
\def\labelenumi{\arabic{enumi}.}
\setcounter{enumi}{4}
\tightlist
\item
  Ковариационная матрица:
\end{enumerate}

\[
\Sigma_{ij}=\mathbb E[X_iX_j]-\mathbb EX_i\,\mathbb EX_j.
\]

Главная ловушка: во второй производной стоит минус.

\subsection{12. Минимум на
хорошо}\label{ux43cux438ux43dux438ux43cux443ux43c-ux43dux430-ux445ux43eux440ux43eux448ux43e-12}

Помимо минимума, нужно:

\begin{enumerate}
\def\labelenumi{\arabic{enumi}.}
\tightlist
\item
  назвать условие существования соответствующих моментов;
\item
  объяснить, почему дифференцирование выносит множитель \(iX_j\);
\item
  различить матрицу вторых моментов \(R\) и ковариационную матрицу
  \(\Sigma\);
\item
  записать формулы
\end{enumerate}

\[
m_j=-i\partial_j\varphi(0),
\]

\[
R_{ij}=-\partial_{ij}\varphi(0),
\]

\[
\Sigma=R-mm^T;
\]

\begin{enumerate}
\def\labelenumi{\arabic{enumi}.}
\setcounter{enumi}{4}
\tightlist
\item
  привести пример Бернулли или нормального распределения.
\end{enumerate}

\subsection{13. Что добавить для
отлично}\label{ux447ux442ux43e-ux434ux43eux431ux430ux432ux438ux442ux44c-ux434ux43bux44f-ux43eux442ux43bux438ux447ux43dux43e-12}

Для сильного ответа стоит добавить:

\begin{enumerate}
\def\labelenumi{\arabic{enumi}.}
\tightlist
\item
  мультииндексную формулу
\end{enumerate}

\[
\partial^\alpha\varphi(0)
=
i^{|\alpha|}
\mathbb E[X^\alpha];
\]

\begin{enumerate}
\def\labelenumi{\arabic{enumi}.}
\setcounter{enumi}{1}
\tightlist
\item
  доказательную схему через теорему Лебега о доминированной сходимости;
\item
  формулу через логарифм характеристической функции:
\end{enumerate}

\[
\partial_{ij}\log\varphi(0)=-\Sigma_{ij};
\]

\begin{enumerate}
\def\labelenumi{\arabic{enumi}.}
\setcounter{enumi}{3}
\tightlist
\item
  связь с конечномерными распределениями процесса:
\end{enumerate}

\[
R(s,t)=-\partial_{uv}\varphi_{s,t}(0,0);
\]

\begin{enumerate}
\def\labelenumi{\arabic{enumi}.}
\setcounter{enumi}{4}
\tightlist
\item
  неотрицательную определенность матриц \(R\) и \(\Sigma\);
\item
  связь с гауссовским вектором:
\end{enumerate}

\[
\varphi(t)=
\exp\left(i\langle t,m\rangle-\frac12t^T\Sigma t\right).
\]

\subsection{14. Устная версия ответа на 5
минут}\label{ux443ux441ux442ux43dux430ux44f-ux432ux435ux440ux441ux438ux44f-ux43eux442ux432ux435ux442ux430-ux43dux430-5-ux43cux438ux43dux443ux442-12}

Сначала я напоминаю определение характеристической функции. Для
случайного вектора

\[
X=(X_1,\ldots,X_n)
\]

она равна

\[
\varphi_X(t)=\mathbb E e^{i\langle t,X\rangle}.
\]

Она существует всегда, потому что экспонента имеет модуль \(1\). Но
производные в нуле связаны с моментами только при существовании
соответствующих моментов.

В одномерном случае, если \(\mathbb E|X|^k<\infty\), то

\[
\varphi_X^{(r)}(0)=i^r\mathbb EX^r,
\qquad r\le k.
\]

В частности,

\[
\varphi'(0)=i\mathbb EX,
\qquad
\varphi''(0)=-\mathbb EX^2.
\]

Значит

\[
\mathbb EX=-i\varphi'(0),
\qquad
\operatorname{Var}X=-\varphi''(0)-(\mathbb EX)^2.
\]

Для вектора дифференцируем по координатам:

\[
\partial_j\varphi(0)=i\mathbb EX_j,
\]

\[
\partial_{ij}\varphi(0)=-\mathbb E[X_iX_j].
\]

Поэтому средний вектор и матрица вторых моментов находятся как

\[
m_j=-i\partial_j\varphi(0),
\]

\[
R_{ij}=-\partial_{ij}\varphi(0).
\]

Дальше ковариационная матрица:

\[
\Sigma_{ij}
=
\operatorname{Cov}(X_i,X_j)
=
R_{ij}-m_im_j.
\]

Если вектор центрирован, то \(m=0\), и корреляционная матрица вторых
моментов совпадает с ковариационной.

Доказательная идея простая: при дифференцировании

\[
e^{i\langle t,X\rangle}
\]

по \(t_j\) появляется множитель \(iX_j\), а при втором дифференцировании
по \(t_i,t_j\) появляется \((iX_i)(iX_j)=-X_iX_j\). При существовании
нужных моментов дифференцирование под знаком ожидания оправдывается
доминированной сходимостью.

Завершаю примером. Для Бернулли

\[
\varphi(t)=q+pe^{it}.
\]

Тогда

\[
\varphi'(0)=ip,
\qquad
\varphi''(0)=-p,
\]

поэтому

\[
\mathbb EX=p,
\qquad
\operatorname{Var}X=p-p^2.
\]

Главные ошибки: потерять минус во второй производной, забыть множитель
\(i\), перепутать вторые моменты и ковариации.

\subsection{15. Если осталось 2 минуты перед
экзаменом}\label{ux435ux441ux43bux438-ux43eux441ux442ux430ux43bux43eux441ux44c-2-ux43cux438ux43dux443ux442ux44b-ux43fux435ux440ux435ux434-ux44dux43aux437ux430ux43cux435ux43dux43eux43c-12}

Запомнить четыре строки.

\textbf{Характеристическая функция:}

\[
\varphi_X(t)=\mathbb E e^{i\langle t,X\rangle}.
\]

\textbf{Среднее:}

\[
m_j=\mathbb EX_j=-i\partial_j\varphi(0).
\]

\textbf{Вторые моменты:}

\[
R_{ij}=\mathbb E[X_iX_j]=-\partial_{ij}\varphi(0).
\]

\textbf{Ковариации:}

\[
\Sigma_{ij}=R_{ij}-m_im_j.
\]

Главная ловушка: характеристическая функция существует всегда, но
производные-моменты требуют существования соответствующих моментов; во
второй производной стоит минус.

\newpage

\section{Билет 14. Коэффициенты корреляции и евклидова геометрия
центрированных величин со вторым
моментом}\label{ux431ux438ux43bux435ux442-14.-ux43aux43eux44dux444ux444ux438ux446ux438ux435ux43dux442ux44b-ux43aux43eux440ux440ux435ux43bux44fux446ux438ux438-ux438-ux435ux432ux43aux43bux438ux434ux43eux432ux430-ux433ux435ux43eux43cux435ux442ux440ux438ux44f-ux446ux435ux43dux442ux440ux438ux440ux43eux432ux430ux43dux43dux44bux445-ux432ux435ux43bux438ux447ux438ux43d-ux441ux43e-ux432ux442ux43eux440ux44bux43c-ux43cux43eux43cux435ux43dux442ux43eux43c}

Источники: Сердобольская, лекция 6; Ширяев 1, §11.

\subsection{1. Короткий
ответ}\label{ux43aux43eux440ux43eux442ux43aux438ux439-ux43eux442ux432ux435ux442-13}

Если случайные величины имеют конечные вторые моменты, то их можно
рассматривать как векторы в гильбертовом пространстве
\(L^2(\Omega,\mathcal F,P)\) со скалярным произведением

\[
\langle X,Y\rangle=\mathbb E[XY].
\]

Для центрированных величин

\[
X^0=X-\mathbb EX,\qquad Y^0=Y-\mathbb EY
\]

это скалярное произведение равно ковариации:

\[
\langle X^0,Y^0\rangle
=
\mathbb E[(X-\mathbb EX)(Y-\mathbb EY)]
=
\operatorname{Cov}(X,Y).
\]

Коэффициент корреляции при \(\operatorname{Var}X>0\) и
\(\operatorname{Var}Y>0\):

\[
\rho(X,Y)
=
\frac{\operatorname{Cov}(X,Y)}
{\sqrt{\operatorname{Var}X\,\operatorname{Var}Y}}
=
\frac{\langle X^0,Y^0\rangle}{\|X^0\|_2\|Y^0\|_2}.
\]

Поэтому \(\rho(X,Y)\) - это косинус угла между центрированными векторами
\(X^0\) и \(Y^0\):

\[
\rho(X,Y)=\cos\angle(X^0,Y^0),
\qquad
|\rho(X,Y)|\le1.
\]

Равенство \(\rho=0\) означает некоррелированность, то есть
ортогональность центрированных величин в \(L^2\). Это не означает
независимость в общем случае.

\subsection{2. Полный
ответ}\label{ux43fux43eux43bux43dux44bux439-ux43eux442ux432ux435ux442-13}

\subsubsection{2.1. Введение и геометрическая
интерпретация}\label{ux432ux432ux435ux434ux435ux43dux438ux435-ux438-ux433ux435ux43eux43cux435ux442ux440ux438ux447ux435ux441ux43aux430ux44f-ux438ux43dux442ux435ux440ux43fux440ux435ux442ux430ux446ux438ux44f}

Этот билет продолжает Билет 12 о гильбертовых пространствах и связывает
его с вероятностными понятиями ковариации и корреляции. До этого
случайные величины рассматривались как измеримые функции, распределения
и интегралы. Здесь важна новая точка зрения: случайные величины с
конечным вторым моментом можно складывать, умножать на числа, измерять
их длины и углы между ними.

\subsubsection{\texorpdfstring{2.2. Пространство \(L^2\) и скалярное
произведение}{2.2. Пространство L\^{}2 и скалярное произведение}}\label{ux43fux440ux43eux441ux442ux440ux430ux43dux441ux442ux432ux43e-l2-ux438-ux441ux43aux430ux43bux44fux440ux43dux43eux435-ux43fux440ux43eux438ux437ux432ux435ux434ux435ux43dux438ux435}

Пусть задано вероятностное пространство

\[
(\Omega,\mathcal F,P).
\]

Пространство

\[
L^2(\Omega,\mathcal F,P)
=
\{X:\mathbb E X^2<\infty\}/\sim
\]

состоит из случайных величин с конечным вторым моментом, где величины,
равные почти наверное, считаются одним элементом. В вещественном случае
скалярное произведение задается формулой

\[
\langle X,Y\rangle=\mathbb E[XY].
\]

Оно корректно определено, потому что из неравенства Коши-Буняковского
следует

\[
\mathbb E|XY|
\le
\sqrt{\mathbb EX^2}\sqrt{\mathbb EY^2}
<\infty.
\]

Норма в \(L^2\):

\[
\|X\|_2=\sqrt{\mathbb EX^2}.
\]

Расстояние между двумя случайными величинами:

\[
\|X-Y\|_2
=
\sqrt{\mathbb E[(X-Y)^2]}.
\]

Это среднеквадратичное расстояние. Именно поэтому геометрия \(L^2\)
естественна для дисперсий, ковариаций, регрессии, прогнозирования и
корреляционных характеристик случайных процессов.

\subsubsection{2.3. Центрирование и
ковариация}\label{ux446ux435ux43dux442ux440ux438ux440ux43eux432ux430ux43dux438ux435-ux438-ux43aux43eux432ux430ux440ux438ux430ux446ux438ux44f}

Теперь выделим центрированные величины. Для любой \(X\in L^2\) положим

\[
X^0=X-\mathbb EX.
\]

Тогда

\[
\mathbb EX^0=0,
\qquad
\|X^0\|_2^2
=
\mathbb E[(X-\mathbb EX)^2]
=
\operatorname{Var}X.
\]

Для двух величин \(X,Y\in L^2\) ковариация:

\[
\operatorname{Cov}(X,Y)
=
\mathbb E[(X-\mathbb EX)(Y-\mathbb EY)].
\]

В геометрических обозначениях:

\[
\operatorname{Cov}(X,Y)=\langle X^0,Y^0\rangle.
\]

То есть ковариация - это скалярное произведение центрированных версий
\(X\) и \(Y\).

\subsubsection{2.4. Коэффициент корреляции как косинус
угла}\label{ux43aux43eux44dux444ux444ux438ux446ux438ux435ux43dux442-ux43aux43eux440ux440ux435ux43bux44fux446ux438ux438-ux43aux430ux43a-ux43aux43eux441ux438ux43dux443ux441-ux443ux433ux43bux430}

Коэффициент корреляции определяется, если обе дисперсии положительны:

\[
\operatorname{Var}X>0,
\qquad
\operatorname{Var}Y>0.
\]

Тогда

\[
\rho(X,Y)
=
\frac{\operatorname{Cov}(X,Y)}
{\sqrt{\operatorname{Var}X\,\operatorname{Var}Y}}.
\]

Через геометрию \(L^2\) это

\[
\rho(X,Y)
=
\frac{\langle X^0,Y^0\rangle}{\|X^0\|_2\|Y^0\|_2}.
\]

В обычной евклидовой геометрии для ненулевых векторов \(u,v\):

\[
\cos\angle(u,v)
=
\frac{\langle u,v\rangle}{\|u\|\|v\|}.
\]

Значит

\[
\rho(X,Y)=\cos\angle(X^0,Y^0).
\]

Именно это нужно проговорить как главный смысл билета: корреляция -
нормированная ковариация, то есть косинус угла между центрированными
случайными величинами.

\subsubsection{2.5. Неравенство Коши-Буняковского и диапазон
значений}\label{ux43dux435ux440ux430ux432ux435ux43dux441ux442ux432ux43e-ux43aux43eux448ux438-ux431ux443ux43dux44fux43aux43eux432ux441ux43aux43eux433ux43e-ux438-ux434ux438ux430ux43fux430ux437ux43eux43d-ux437ux43dux430ux447ux435ux43dux438ux439}

Из этой формулы сразу следует

\[
|\rho(X,Y)|\le1.
\]

Доказательство - это неравенство Коши-Буняковского:

\[
|\langle X^0,Y^0\rangle|
\le
\|X^0\|_2\|Y^0\|_2.
\]

Подставляя

\[
\langle X^0,Y^0\rangle=\operatorname{Cov}(X,Y),
\qquad
\|X^0\|_2^2=\operatorname{Var}X,
\qquad
\|Y^0\|_2^2=\operatorname{Var}Y,
\]

получаем

\[
|\operatorname{Cov}(X,Y)|
\le
\sqrt{\operatorname{Var}X\,\operatorname{Var}Y},
\]

а после деления на положительный знаменатель:

\[
|\rho(X,Y)|\le1.
\]

\subsubsection{2.6. Особые случаи: некоррелированность и
независимость}\label{ux43eux441ux43eux431ux44bux435-ux441ux43bux443ux447ux430ux438-ux43dux435ux43aux43eux440ux440ux435ux43bux438ux440ux43eux432ux430ux43dux43dux43eux441ux442ux44c-ux438-ux43dux435ux437ux430ux432ux438ux441ux438ux43cux43eux441ux442ux44c}

Разберем смысл особых случаев.

Если

\[
\rho(X,Y)=0,
\]

то

\[
\operatorname{Cov}(X,Y)=0
\]

и

\[
X^0\perp Y^0
\]

в \(L^2\). Такие величины называются некоррелированными. Важно:
некоррелированность в общем случае слабее независимости. Из
независимости при существовании вторых моментов действительно следует

\[
\operatorname{Cov}(X,Y)=0,
\]

потому что

\[
\mathbb E[XY]=\mathbb EX\,\mathbb EY.
\]

Но обратное неверно.

Классический пример: пусть \(X\) равномерна на \([-1,1]\), а

\[
Y=X^2.
\]

Тогда \(Y\) полностью определяется через \(X\), то есть независимости
нет. Но

\[
\operatorname{Cov}(X,Y)
=
\mathbb E[X^3]-\mathbb EX\,\mathbb EX^2
=
0-0\cdot\mathbb EX^2
=0,
\]

поскольку распределение \(X\) симметрично и \(\mathbb E X^3=0\). Значит
зависимые величины могут быть некоррелированными.

\subsubsection{2.7. Линейная зависимость и вырожденная
дисперсия}\label{ux43bux438ux43dux435ux439ux43dux430ux44f-ux437ux430ux432ux438ux441ux438ux43cux43eux441ux442ux44c-ux438-ux432ux44bux440ux43eux436ux434ux435ux43dux43dux430ux44f-ux434ux438ux441ux43fux435ux440ux441ux438ux44f}

Если

\[
|\rho(X,Y)|=1,
\]

то в неравенстве Коши-Буняковского достигается равенство. Это означает
линейную зависимость центрированных величин в \(L^2\):

\[
Y^0=aX^0
\quad\text{почти наверное}
\]

для некоторого \(a\ne0\). Эквивалентно,

\[
Y=aX+b
\quad\text{почти наверное}.
\]

Если \(a>0\), то \(\rho=1\); если \(a<0\), то \(\rho=-1\).

Если одна из дисперсий равна нулю, например

\[
\operatorname{Var}X=0,
\]

то \(X=\mathbb EX\) почти наверное, то есть \(X\) является константой в
\(L^2\). Тогда угол с ней как с центрированным вектором не определен,
потому что

\[
X^0=0.
\]

Поэтому коэффициент корреляции обычно не определяют при нулевой
дисперсии.

\subsubsection{2.8. Ковариационная и корреляционная
матрицы}\label{ux43aux43eux432ux430ux440ux438ux430ux446ux438ux43eux43dux43dux430ux44f-ux438-ux43aux43eux440ux440ux435ux43bux44fux446ux438ux43eux43dux43dux430ux44f-ux43cux430ux442ux440ux438ux446ux44b}

Для случайного вектора

\[
X=(X_1,\ldots,X_n)^T
\]

с конечными вторыми моментами ковариационная матрица

\[
\Sigma=(\Sigma_{ij}),
\qquad
\Sigma_{ij}=\operatorname{Cov}(X_i,X_j)
\]

является матрицей Грама для центрированных векторов

\[
X_1^0,\ldots,X_n^0
\]

в \(L^2\):

\[
\Sigma_{ij}=\langle X_i^0,X_j^0\rangle.
\]

Отсюда сразу следует ее симметричность и неотрицательная определенность.
Для любого \(a=(a_1,\ldots,a_n)^T\):

\[
a^T\Sigma a
=
\sum_{i,j}a_i a_j\operatorname{Cov}(X_i,X_j)
=
\operatorname{Var}\left(\sum_{i=1}^n a_iX_i\right)
\ge0.
\]

Корреляционная матрица получается нормировкой ковариаций:

\[
\rho_{ij}
=
\rho(X_i,X_j)
=
\frac{\Sigma_{ij}}{\sqrt{\Sigma_{ii}\Sigma_{jj}}},
\]

если соответствующие дисперсии положительны. На диагонали такой матрицы
стоят единицы:

\[
R_{ii}=1.
\]

Геометрически \(R_{ij}\) - косинус угла между \(X_i^0\) и \(X_j^0\).

\subsubsection{2.9. Связь со случайными
процессами}\label{ux441ux432ux44fux437ux44c-ux441ux43e-ux441ux43bux443ux447ux430ux439ux43dux44bux43cux438-ux43fux440ux43eux446ux435ux441ux441ux430ux43cux438}

В конце ответа полезно связать билет с процессами. Для случайной
последовательности или процесса ковариации между значениями в разные
моменты времени образуют корреляционную функцию:

\[
R(s,t)=\operatorname{Cov}(X_s,X_t).
\]

В стационарном широком смысле она зависит только от лага:

\[
R(k)=\operatorname{Cov}(X_0,X_k).
\]

Это уже связь с Билет 17, Билет 24 и Билет 25: корреляционная функция и
спектральные характеристики являются развитием той же геометрии второго
порядка.

\subsection{3. Основные
определения}\label{ux43eux441ux43dux43eux432ux43dux44bux435-ux43eux43fux440ux435ux434ux435ux43bux435ux43dux438ux44f-13}

\begin{itemize}
\tightlist
\item
  \textbf{Центрированная случайная величина}:
\end{itemize}

\[
X^0=X-\mathbb EX.
\]

\begin{itemize}
\tightlist
\item
  \textbf{Дисперсия}:
\end{itemize}

\[
\operatorname{Var}X=\mathbb E[(X-\mathbb EX)^2]=\|X^0\|_2^2.
\]

\begin{itemize}
\tightlist
\item
  \textbf{Ковариация}:
\end{itemize}

\[
\operatorname{Cov}(X,Y)=\mathbb E[(X-\mathbb EX)(Y-\mathbb EY)].
\]

\begin{itemize}
\tightlist
\item
  \textbf{Коэффициент корреляции}:
\end{itemize}

\[
\rho(X,Y)
=
\frac{\operatorname{Cov}(X,Y)}
{\sqrt{\operatorname{Var}X\,\operatorname{Var}Y}}.
\]

\begin{itemize}
\tightlist
\item
  \textbf{Ортогональность в \(L^2\)}:
\end{itemize}

\[
X\perp Y
\quad\Longleftrightarrow\quad
\mathbb E[XY]=0.
\]

\begin{itemize}
\tightlist
\item
  \textbf{Некоррелированность}:
\end{itemize}

\[
X^0\perp Y^0
\quad\Longleftrightarrow\quad
\operatorname{Cov}(X,Y)=0.
\]

\begin{itemize}
\tightlist
\item
  \textbf{Матрица Грама}: матрица скалярных произведений векторов:
\end{itemize}

\[
G_{ij}=\langle u_i,u_j\rangle.
\]

Ковариационная матрица - матрица Грама для центрированных величин.

\subsection{4. Основные теоремы и
свойства}\label{ux43eux441ux43dux43eux432ux43dux44bux435-ux442ux435ux43eux440ux435ux43cux44b-ux438-ux441ux432ux43eux439ux441ux442ux432ux430-13}

\begin{itemize}
\tightlist
\item
  \(\operatorname{Cov}(X,Y)=\langle X^0,Y^0\rangle\).
\item
  \(\operatorname{Var}X=\|X^0\|_2^2\).
\item
  \(\rho(X,Y)=\cos\angle(X^0,Y^0)\) при положительных дисперсиях.
\item
  \(|\rho(X,Y)|\le1\).
\item
  \(\rho(X,Y)=0\) означает некоррелированность, то есть ортогональность
  центрированных величин.
\item
  Независимость при наличии вторых моментов влечет некоррелированность,
  но не наоборот.
\item
  \(|\rho(X,Y)|=1\) тогда и только тогда, когда \(X\) и \(Y\) линейно
  связаны почти наверное:
\end{itemize}

\[
Y=aX+b.
\]

\begin{itemize}
\tightlist
\item
  Ковариационная матрица симметрична и неотрицательно определена.
\item
  Ковариационная матрица является матрицей Грама центрированных
  случайных величин.
\end{itemize}

\subsection{5. Примеры}\label{ux43fux440ux438ux43cux435ux440ux44b-13}

\subsubsection{Независимые
величины}\label{ux43dux435ux437ux430ux432ux438ux441ux438ux43cux44bux435-ux432ux435ux43bux438ux447ux438ux43dux44b}

Если \(X\) и \(Y\) независимы и имеют конечные вторые моменты, то

\[
\mathbb E[XY]=\mathbb EX\,\mathbb EY.
\]

Тогда

\[
\operatorname{Cov}(X,Y)
=
\mathbb E[XY]-\mathbb EX\,\mathbb EY
=0.
\]

Значит независимые величины некоррелированы. Если дисперсии
положительны, то \(\rho(X,Y)=0\).

\subsubsection{Линейная
зависимость}\label{ux43bux438ux43dux435ux439ux43dux430ux44f-ux437ux430ux432ux438ux441ux438ux43cux43eux441ux442ux44c}

Пусть

\[
Y=aX+b,
\qquad a\ne0.
\]

Тогда

\[
Y-\mathbb EY
=
a(X-\mathbb EX).
\]

Поэтому

\[
\operatorname{Cov}(X,Y)=a\operatorname{Var}X,
\qquad
\operatorname{Var}Y=a^2\operatorname{Var}X.
\]

Если \(\operatorname{Var}X>0\), то

\[
\rho(X,Y)
=
\frac{a\operatorname{Var}X}{|a|\operatorname{Var}X}
=
\operatorname{sign}a.
\]

То есть при \(a>0\) корреляция равна \(1\), а при \(a<0\) равна \(-1\).

\subsubsection{Некоррелированные, но зависимые
величины}\label{ux43dux435ux43aux43eux440ux440ux435ux43bux438ux440ux43eux432ux430ux43dux43dux44bux435-ux43dux43e-ux437ux430ux432ux438ux441ux438ux43cux44bux435-ux432ux435ux43bux438ux447ux438ux43dux44b}

Пусть \(X\sim U[-1,1]\), а \(Y=X^2\). Тогда \(Y\) зависит от \(X\), но

\[
\operatorname{Cov}(X,Y)
=
\operatorname{Cov}(X,X^2)
=
\mathbb E X^3-\mathbb EX\,\mathbb EX^2
=0.
\]

Значит \(\rho(X,Y)=0\), если \(\operatorname{Var}Y>0\). Этот пример
показывает главную ловушку: нулевая корреляция не равна независимости.

\subsubsection{Ортогональные векторы в конечномерном
случае}\label{ux43eux440ux442ux43eux433ux43eux43dux430ux43bux44cux43dux44bux435-ux432ux435ux43aux442ux43eux440ux44b-ux432-ux43aux43eux43dux435ux447ux43dux43eux43cux435ux440ux43dux43eux43c-ux441ux43bux443ux447ux430ux435}

Если \(X=(X_1,X_2)\) - центрированный случайный вектор и

\[
\operatorname{Cov}(X_1,X_2)=0,
\]

то координаты \(X_1\) и \(X_2\) ортогональны как элементы \(L^2\). Их
ковариационная матрица диагональна:

\[
\Sigma
=
\begin{pmatrix}
\operatorname{Var}X_1 & 0\\
0 & \operatorname{Var}X_2
\end{pmatrix}.
\]

Но диагональность ковариационной матрицы сама по себе не означает
независимость координат, кроме специальных классов распределений,
например совместно гауссовского случая.

\subsection{6. Схемы доказательств /
обоснований}\label{ux441ux445ux435ux43cux44b-ux434ux43eux43aux430ux437ux430ux442ux435ux43bux44cux441ux442ux432-ux43eux431ux43eux441ux43dux43eux432ux430ux43dux438ux439-6}

\subsubsection{Почему ковариация - скалярное
произведение}\label{ux43fux43eux447ux435ux43cux443-ux43aux43eux432ux430ux440ux438ux430ux446ux438ux44f---ux441ux43aux430ux43bux44fux440ux43dux43eux435-ux43fux440ux43eux438ux437ux432ux435ux434ux435ux43dux438ux435}

Берем центрированные величины

\[
X^0=X-\mathbb EX,
\qquad
Y^0=Y-\mathbb EY.
\]

Тогда по определению скалярного произведения в \(L^2\):

\[
\langle X^0,Y^0\rangle
=
\mathbb E[X^0Y^0]
=
\mathbb E[(X-\mathbb EX)(Y-\mathbb EY)]
=
\operatorname{Cov}(X,Y).
\]

\subsubsection{\texorpdfstring{Почему
\(|\rho|\le1\)}{Почему \textbar\textbackslash rho\textbar\textbackslash le1}}\label{ux43fux43eux447ux435ux43cux443-rhole1}

Применяем Коши-Буняковского в \(L^2\):

\[
|\langle X^0,Y^0\rangle|
\le
\|X^0\|_2\|Y^0\|_2.
\]

После подстановки получаем

\[
|\operatorname{Cov}(X,Y)|
\le
\sqrt{\operatorname{Var}X\,\operatorname{Var}Y}.
\]

Если дисперсии положительны, делим на знаменатель и получаем

\[
|\rho(X,Y)|\le1.
\]

\subsubsection{\texorpdfstring{Почему \(|\rho|=1\) означает линейную
зависимость}{Почему \textbar\textbackslash rho\textbar=1 означает линейную зависимость}}\label{ux43fux43eux447ux435ux43cux443-rho1-ux43eux437ux43dux430ux447ux430ux435ux442-ux43bux438ux43dux435ux439ux43dux443ux44e-ux437ux430ux432ux438ux441ux438ux43cux43eux441ux442ux44c}

В неравенстве Коши-Буняковского равенство достигается тогда и только
тогда, когда векторы линейно зависимы. Поэтому

\[
|\rho(X,Y)|=1
\]

эквивалентно существованию числа \(a\ne0\) такого, что

\[
Y^0=aX^0
\quad\text{почти наверное}.
\]

Раскрывая центрирование:

\[
Y-\mathbb EY=a(X-\mathbb EX),
\]

то есть

\[
Y=aX+(\mathbb EY-a\mathbb EX)
\quad\text{почти наверное}.
\]

\subsubsection{Почему ковариационная матрица неотрицательно
определена}\label{ux43fux43eux447ux435ux43cux443-ux43aux43eux432ux430ux440ux438ux430ux446ux438ux43eux43dux43dux430ux44f-ux43cux430ux442ux440ux438ux446ux430-ux43dux435ux43eux442ux440ux438ux446ux430ux442ux435ux43bux44cux43dux43e-ux43eux43fux440ux435ux434ux435ux43bux435ux43dux430}

Пусть

\[
\Sigma_{ij}=\operatorname{Cov}(X_i,X_j).
\]

Тогда для любого \(a\in\mathbb R^n\):

\[
a^T\Sigma a
=
\sum_{i,j}a_i a_j\operatorname{Cov}(X_i,X_j)
=
\operatorname{Var}\left(\sum_i a_iX_i\right)
\ge0.
\]

Значит \(\Sigma\) неотрицательно определена.

\subsection{7. Что написать на
доске}\label{ux447ux442ux43e-ux43dux430ux43fux438ux441ux430ux442ux44c-ux43dux430-ux434ux43eux441ux43aux435-13}

\begin{enumerate}
\def\labelenumi{\arabic{enumi}.}
\tightlist
\item
  Центрирование:
\end{enumerate}

\[
X^0=X-\mathbb EX.
\]

\begin{enumerate}
\def\labelenumi{\arabic{enumi}.}
\setcounter{enumi}{1}
\tightlist
\item
  Ковариация как скалярное произведение:
\end{enumerate}

\[
\operatorname{Cov}(X,Y)=\langle X^0,Y^0\rangle.
\]

\begin{enumerate}
\def\labelenumi{\arabic{enumi}.}
\setcounter{enumi}{2}
\tightlist
\item
  Дисперсия как квадрат длины:
\end{enumerate}

\[
\operatorname{Var}X=\|X^0\|_2^2.
\]

\begin{enumerate}
\def\labelenumi{\arabic{enumi}.}
\setcounter{enumi}{3}
\tightlist
\item
  Коэффициент корреляции:
\end{enumerate}

\[
\rho(X,Y)
=
\frac{\operatorname{Cov}(X,Y)}
{\sqrt{\operatorname{Var}X\,\operatorname{Var}Y}}
=
\frac{\langle X^0,Y^0\rangle}{\|X^0\|_2\|Y^0\|_2}
=
\cos\angle(X^0,Y^0).
\]

\begin{enumerate}
\def\labelenumi{\arabic{enumi}.}
\setcounter{enumi}{4}
\tightlist
\item
  Главное следствие:
\end{enumerate}

\[
|\rho(X,Y)|\le1.
\]

\begin{enumerate}
\def\labelenumi{\arabic{enumi}.}
\setcounter{enumi}{5}
\tightlist
\item
  Крайние случаи:
\end{enumerate}

\[
\rho=0 \Longleftrightarrow X^0\perp Y^0,
\]

\[
|\rho|=1 \Longleftrightarrow Y=aX+b\quad\text{п.н.}
\]

\subsection{8. Типичные дополнительные
вопросы}\label{ux442ux438ux43fux438ux447ux43dux44bux435-ux434ux43eux43fux43eux43bux43dux438ux442ux435ux43bux44cux43dux44bux435-ux432ux43eux43fux440ux43eux441ux44b-13}

\textbf{Почему нужно центрировать величины?}

Потому что корреляция должна измерять совместные отклонения от средних,
а не сами уровни величин. Без центрирования скалярное произведение
\(\mathbb E[XY]\) зависит от сдвигов. Например, если заменить \(X\) на
\(X+c\), то \(\mathbb E[(X+c)Y]\) обычно изменится, хотя форма
случайного отклонения \(X-\mathbb EX\) не изменилась. Центрирование
убирает постоянную часть:

\[
X^0=X-\mathbb EX.
\]

\textbf{Почему коэффициент корреляции не определяют при нулевой
дисперсии?}

Если \(\operatorname{Var}X=0\), то \(X=\mathbb EX\) почти наверное и

\[
X^0=0.
\]

У нулевого вектора нет направления и угла с другим вектором. В формуле
для \(\rho\) появляется деление на ноль:

\[
\rho(X,Y)=
\frac{\operatorname{Cov}(X,Y)}
{\sqrt{\operatorname{Var}X\,\operatorname{Var}Y}}.
\]

\textbf{В чем разница между ковариацией и корреляцией?}

Ковариация - это ненормированное скалярное произведение центрированных
величин:

\[
\operatorname{Cov}(X,Y)=\langle X^0,Y^0\rangle.
\]

Она зависит от масштаба измерения. Если \(X\) умножить на 10, ковариация
тоже умножится на 10. Корреляция - нормированная ковариация:

\[
\rho(X,Y)=
\frac{\operatorname{Cov}(X,Y)}
{\sqrt{\operatorname{Var}X\,\operatorname{Var}Y}},
\]

поэтому всегда лежит в \([-1,1]\) и показывает только направление
линейной связи.

\textbf{Почему \(|\rho|\le1\)?}

Потому что в \(L^2\) выполняется неравенство Коши-Буняковского:

\[
|\langle X^0,Y^0\rangle|
\le
\|X^0\|_2\|Y^0\|_2.
\]

После подстановки \(\langle X^0,Y^0\rangle=\operatorname{Cov}(X,Y)\) и
\(\|X^0\|_2^2=\operatorname{Var}X\) получаем

\[
|\operatorname{Cov}(X,Y)|
\le
\sqrt{\operatorname{Var}X\,\operatorname{Var}Y},
\]

а значит \(|\rho|\le1\).

\textbf{Что означает \(\rho=0\) геометрически?}

Это означает, что центрированные величины ортогональны в \(L^2\):

\[
X^0\perp Y^0.
\]

Вероятностно это называется некоррелированностью:

\[
\operatorname{Cov}(X,Y)=0.
\]

\textbf{Следует ли независимость из \(\rho=0\)?}

В общем случае нет. Нулевая корреляция означает только отсутствие
линейной связи второго порядка. Например, если \(X\sim U[-1,1]\) и
\(Y=X^2\), то \(Y\) зависит от \(X\), но

\[
\operatorname{Cov}(X,Y)=\operatorname{Cov}(X,X^2)=0.
\]

\textbf{Когда из некоррелированности все-таки следует независимость?}

Важный случай - совместно гауссовские величины. Если вектор \((X,Y)\)
гауссовский и

\[
\operatorname{Cov}(X,Y)=0,
\]

то \(X\) и \(Y\) независимы. Это специальное свойство гауссовских
распределений, а не общее правило.

\textbf{Что означает \(|\rho|=1\)?}

Это означает, что в неравенстве Коши-Буняковского достигнуто равенство,
поэтому центрированные величины линейно зависимы:

\[
Y^0=aX^0
\quad\text{п.н.}
\]

Эквивалентно,

\[
Y=aX+b
\quad\text{п.н.}
\]

При \(a>0\) получаем \(\rho=1\), при \(a<0\) получаем \(\rho=-1\).

\textbf{Почему ковариационная матрица неотрицательно определена?}

Пусть \(\Sigma_{ij}=\operatorname{Cov}(X_i,X_j)\). Тогда для любого
\(a\in\mathbb R^n\):

\[
a^T\Sigma a
=
\operatorname{Var}\left(\sum_i a_iX_i\right)
\ge0.
\]

Значит \(\Sigma\) неотрицательно определена. Геометрически это матрица
Грама центрированных векторов \(X_1^0,\ldots,X_n^0\).

\textbf{Как билет связан с Билет 12, Билет 13, Билет 18, Билет 24 и
Билет 25?}

С Билетом 12 связь через гильбертово пространство \(L^2\). С Билетом 13
- через получение моментов и ковариационной матрицы из
характеристической функции. С Билетом 18 - через гауссовские векторы,
где некоррелированность равносильна независимости. С Билетом 24 и
Билетом 25 - через корреляционные функции и спектральные характеристики
случайных последовательностей.

\subsection{9. Типичные
ошибки}\label{ux442ux438ux43fux438ux447ux43dux44bux435-ux43eux448ux438ux431ux43aux438-13}

\begin{itemize}
\tightlist
\item
  Делать вывод о независимости из \(\rho=0\).
\item
  Забывать центрирование и писать \(\mathbb E[XY]\) вместо
  \(\operatorname{Cov}(X,Y)\) без условия \(\mathbb EX=\mathbb EY=0\).
\item
  Не проверять существование вторых моментов.
\item
  Не проверять, что дисперсии положительны.
\item
  Называть ковариационную матрицу корреляционной без пояснения
  нормировки.
\item
  Путать нулевую ковариацию и ортогональность самих \(X,Y\): на самом
  деле ортогональны центрированные величины \(X^0,Y^0\).
\item
  Забывать, что элементы \(L^2\) рассматриваются с точностью до
  равенства почти наверное.
\end{itemize}

\subsection{10. Связи с другими
билетами}\label{ux441ux432ux44fux437ux438-ux441-ux434ux440ux443ux433ux438ux43cux438-ux431ux438ux43bux435ux442ux430ux43cux438-13}

\begin{itemize}
\tightlist
\item
  Билет 11: математические ожидания, дисперсии и ковариации выражаются
  через интегралы.
\item
  Билет 12: \(L^2\) - гильбертово пространство; корреляция использует
  его скалярное произведение.
\item
  Билет 13: моменты и ковариационная матрица могут быть получены из
  производных характеристической функции.
\item
  Билет 18: для совместно гауссовских величин некоррелированность
  эквивалентна независимости.
\item
  Билет 24: корреляционная функция процесса - ковариация значений
  процесса в разные моменты времени.
\item
  Билет 25: спектральные характеристики строятся из корреляционной
  функции стационарной последовательности.
\end{itemize}

\subsection{11. Минимум на
удовлетворительно}\label{ux43cux438ux43dux438ux43cux443ux43c-ux43dux430-ux443ux434ux43eux432ux43bux435ux442ux432ux43eux440ux438ux442ux435ux43bux44cux43dux43e-13}

Нужно сказать:

\begin{enumerate}
\def\labelenumi{\arabic{enumi}.}
\tightlist
\item
  Для \(X,Y\in L^2\) определены центрированные величины:
\end{enumerate}

\[
X^0=X-\mathbb EX,
\qquad
Y^0=Y-\mathbb EY.
\]

\begin{enumerate}
\def\labelenumi{\arabic{enumi}.}
\setcounter{enumi}{1}
\tightlist
\item
  Ковариация - скалярное произведение центрированных величин:
\end{enumerate}

\[
\operatorname{Cov}(X,Y)=\langle X^0,Y^0\rangle.
\]

\begin{enumerate}
\def\labelenumi{\arabic{enumi}.}
\setcounter{enumi}{2}
\tightlist
\item
  Коэффициент корреляции:
\end{enumerate}

\[
\rho(X,Y)=
\frac{\operatorname{Cov}(X,Y)}
{\sqrt{\operatorname{Var}X\,\operatorname{Var}Y}}.
\]

\begin{enumerate}
\def\labelenumi{\arabic{enumi}.}
\setcounter{enumi}{3}
\tightlist
\item
  Геометрический смысл:
\end{enumerate}

\[
\rho(X,Y)=\cos\angle(X^0,Y^0),
\qquad
|\rho|\le1.
\]

\begin{enumerate}
\def\labelenumi{\arabic{enumi}.}
\setcounter{enumi}{4}
\tightlist
\item
  Главная ловушка: \(\rho=0\) не означает независимость.
\end{enumerate}

\subsection{12. Минимум на
хорошо}\label{ux43cux438ux43dux438ux43cux443ux43c-ux43dux430-ux445ux43eux440ux43eux448ux43e-13}

Помимо минимума, нужно объяснить:

\begin{itemize}
\tightlist
\item
  почему нужно условие \(\mathbb EX^2,\mathbb EY^2<\infty\);
\item
  почему нужны положительные дисперсии для \(\rho\);
\item
  как из Коши-Буняковского получается \(|\rho|\le1\);
\item
  что \(\rho=0\) означает ортогональность центрированных величин;
\item
  что \(|\rho|=1\) означает линейную зависимость почти наверное;
\item
  привести пример зависимых, но некоррелированных величин:
  \(X\sim U[-1,1]\), \(Y=X^2\).
\end{itemize}

\subsection{13. Что добавить для
отлично}\label{ux447ux442ux43e-ux434ux43eux431ux430ux432ux438ux442ux44c-ux434ux43bux44f-ux43eux442ux43bux438ux447ux43dux43e-13}

Для сильного ответа стоит добавить:

\begin{itemize}
\tightlist
\item
  формулировку через гильбертово пространство
  \(L^2(\Omega,\mathcal F,P)\);
\item
  замечание, что элементы \(L^2\) - классы эквивалентности по равенству
  почти наверное;
\item
  доказательство \(|\rho|\le1\) через Коши-Буняковского;
\item
  критерий равенства \(|\rho|=1\);
\item
  объяснение ковариационной матрицы как матрицы Грама;
\item
  доказательство неотрицательной определенности ковариационной матрицы;
\item
  связь с гауссовским случаем: для совместно гауссовских величин
  некоррелированность дает независимость;
\item
  связь с корреляционной функцией случайных процессов и спектральными
  характеристиками.
\end{itemize}

\subsection{14. Устная версия ответа на 5
минут}\label{ux443ux441ux442ux43dux430ux44f-ux432ux435ux440ux441ux438ux44f-ux43eux442ux432ux435ux442ux430-ux43dux430-5-ux43cux438ux43dux443ux442-13}

\begin{enumerate}
\def\labelenumi{\arabic{enumi}.}
\tightlist
\item
  В этом билете случайные величины с конечным вторым моментом
  рассматриваются как векторы в \(L^2\).
\item
  В \(L^2\) скалярное произведение:
\end{enumerate}

\[
\langle X,Y\rangle=\mathbb E[XY].
\]

\begin{enumerate}
\def\labelenumi{\arabic{enumi}.}
\setcounter{enumi}{2}
\tightlist
\item
  Для геометрии корреляции надо перейти к центрированным величинам:
\end{enumerate}

\[
X^0=X-\mathbb EX,
\qquad
Y^0=Y-\mathbb EY.
\]

\begin{enumerate}
\def\labelenumi{\arabic{enumi}.}
\setcounter{enumi}{3}
\tightlist
\item
  Тогда
\end{enumerate}

\[
\operatorname{Cov}(X,Y)=\langle X^0,Y^0\rangle,
\qquad
\operatorname{Var}X=\|X^0\|_2^2.
\]

\begin{enumerate}
\def\labelenumi{\arabic{enumi}.}
\setcounter{enumi}{4}
\tightlist
\item
  При положительных дисперсиях коэффициент корреляции:
\end{enumerate}

\[
\rho(X,Y)=
\frac{\operatorname{Cov}(X,Y)}
{\sqrt{\operatorname{Var}X\,\operatorname{Var}Y}}
=
\frac{\langle X^0,Y^0\rangle}{\|X^0\|_2\|Y^0\|_2}.
\]

\begin{enumerate}
\def\labelenumi{\arabic{enumi}.}
\setcounter{enumi}{5}
\tightlist
\item
  Значит \(\rho\) - косинус угла между \(X^0\) и \(Y^0\):
\end{enumerate}

\[
\rho(X,Y)=\cos\angle(X^0,Y^0).
\]

\begin{enumerate}
\def\labelenumi{\arabic{enumi}.}
\setcounter{enumi}{6}
\tightlist
\item
  По Коши-Буняковскому:
\end{enumerate}

\[
|\rho(X,Y)|\le1.
\]

\begin{enumerate}
\def\labelenumi{\arabic{enumi}.}
\setcounter{enumi}{7}
\tightlist
\item
  Если \(\rho=0\), то \(X^0\) и \(Y^0\) ортогональны, то есть величины
  некоррелированы. Это не значит, что они независимы.
\item
  Если \(|\rho|=1\), то центрированные величины линейно зависимы, то
  есть \(Y=aX+b\) почти наверное.
\item
  Пример ловушки: \(X\sim U[-1,1]\), \(Y=X^2\). Они зависимы, но
  \(\operatorname{Cov}(X,Y)=0\).
\item
  Для случайного вектора ковариационная матрица - матрица Грама
  центрированных координат, поэтому она симметрична и неотрицательно
  определена.
\item
  Завершить можно связью: эта геометрия лежит в основе корреляционных
  функций процессов и спектральных характеристик.
\end{enumerate}

\subsection{15. Если осталось 2 минуты перед
экзаменом}\label{ux435ux441ux43bux438-ux43eux441ux442ux430ux43bux43eux441ux44c-2-ux43cux438ux43dux443ux442ux44b-ux43fux435ux440ux435ux434-ux44dux43aux437ux430ux43cux435ux43dux43eux43c-13}

Запомнить три формулы:

\[
X^0=X-\mathbb EX,
\]

\[
\operatorname{Cov}(X,Y)=\langle X^0,Y^0\rangle,
\]

\[
\rho(X,Y)=
\frac{\operatorname{Cov}(X,Y)}
{\sqrt{\operatorname{Var}X\,\operatorname{Var}Y}}
=
\cos\angle(X^0,Y^0).
\]

Главное следствие:

\[
|\rho|\le1.
\]

Главная ловушка: \(\rho=0\) означает только некоррелированность, а не
независимость. Независимость из некоррелированности следует не всегда;
важное исключение - совместно гауссовские величины.

\newpage

\section{Билет 15. Теорема Колмогорова о согласованных конечномерных
распределениях}\label{ux431ux438ux43bux435ux442-15.-ux442ux435ux43eux440ux435ux43cux430-ux43aux43eux43bux43cux43eux433ux43eux440ux43eux432ux430-ux43e-ux441ux43eux433ux43bux430ux441ux43eux432ux430ux43dux43dux44bux445-ux43aux43eux43dux435ux447ux43dux43eux43cux435ux440ux43dux44bux445-ux440ux430ux441ux43fux440ux435ux434ux435ux43bux435ux43dux438ux44fux445}

Источники: Ширяев 1, §9; Сердобольская, лекция 1.

\subsection{1. Короткий
ответ}\label{ux43aux43eux440ux43eux442ux43aux438ux439-ux43eux442ux432ux435ux442-14}

Теорема Колмогорова отвечает на главный вопрос теории процессов: когда
набор конечномерных распределений действительно является набором
распределений некоторого случайного процесса.

Пусть \(T\) - множество моментов времени, \(E\) - пространство
состояний, например \(E=\mathbb R\), и для каждого конечного набора

\[
t_1,\ldots,t_n\in T
\]

задана вероятностная мера

\[
\mu_{t_1,\ldots,t_n}
\]

на \(E^n\). Эти меры должны быть согласованы:

\begin{enumerate}
\def\labelenumi{\arabic{enumi}.}
\tightlist
\item
  перестановка моментов времени переставляет координаты распределения;
\item
  удаление лишних координат дает соответствующее маргинальное
  распределение.
\end{enumerate}

Тогда существует вероятностная мера \(P\) на пространстве траекторий
\(E^T\) такая, что для координатного процесса

\[
X_t(x)=x(t),\qquad x\in E^T,
\]

его конечномерные распределения равны заданным:

\[
P\{(X_{t_1},\ldots,X_{t_n})\in B\}
=
\mu_{t_1,\ldots,t_n}(B).
\]

Иначе говоря, согласованное семейство конечномерных распределений задает
закон процесса на пространстве траекторий.

\subsection{2. Полный
ответ}\label{ux43fux43eux43bux43dux44bux439-ux43eux442ux432ux435ux442-14}

\subsubsection{2.1. Введение и
мотивация}\label{ux432ux432ux435ux434ux435ux43dux438ux435-ux438-ux43cux43eux442ux438ux432ux430ux446ux438ux44f}

В Билете 06 конечномерные распределения процесса были введены как
распределения случайных векторов

\[
(X_{t_1},\ldots,X_{t_n}).
\]

Если процесс уже задан, то его конечномерные распределения автоматически
существуют и согласованы между собой. В билете 15 вопрос обратный: можно
ли сначала задать все конечномерные распределения, а потом получить
случайный процесс, для которого они действительно являются
распределениями конечных наборов координат.

Именно это делает теорема Колмогорова о согласованных конечномерных
распределениях. Она является базовой конструкцией существования
случайных процессов. Без нее определения вроде ``пусть есть процесс с
такими конечномерными распределениями'' были бы неполными: нужно
доказать, что такой процесс вообще существует.

\subsubsection{2.2. Основные обозначения: время и фазовое
пространство}\label{ux43eux441ux43dux43eux432ux43dux44bux435-ux43eux431ux43eux437ux43dux430ux447ux435ux43dux438ux44f-ux432ux440ux435ux43cux44f-ux438-ux444ux430ux437ux43eux432ux43eux435-ux43fux440ux43eux441ux442ux440ux430ux43dux441ux442ux432ux43e}

Пусть \(T\) - множество индексов. В дискретном времени обычно

\[
T=\mathbb N
\]

или

\[
T=\mathbb Z.
\]

В непрерывном времени может быть

\[
T=\mathbb R_+
\]

или

\[
T=[0,\infty).
\]

Пусть \(E\) - пространство значений процесса. Для экзамена достаточно
уверенно рассматривать случай

\[
E=\mathbb R
\]

с борелевской \(\sigma\)-алгеброй \(\mathcal B(\mathbb R)\), но
формулировка работает и для стандартных борелевских пространств.

\subsubsection{2.3. Пространство траекторий и цилиндрические
множества}\label{ux43fux440ux43eux441ux442ux440ux430ux43dux441ux442ux432ux43e-ux442ux440ux430ux435ux43aux442ux43eux440ux438ux439-ux438-ux446ux438ux43bux438ux43dux434ux440ux438ux447ux435ux441ux43aux438ux435-ux43cux43dux43eux436ux435ux441ux442ux432ux430}

Пространство всех траекторий:

\[
E^T=\{x:T\to E\}.
\]

Для каждого \(t\in T\) есть координатное отображение

\[
X_t:E^T\to E,
\qquad
X_t(x)=x(t).
\]

Если выбрать конечный набор индексов

\[
t_1,\ldots,t_n,
\]

то есть конечномерная проекция

\[
\pi_{t_1,\ldots,t_n}:E^T\to E^n,
\qquad
\pi_{t_1,\ldots,t_n}(x)=(x(t_1),\ldots,x(t_n)).
\]

Цилиндрическим множеством называется прообраз борелевского множества при
такой конечномерной проекции:

\[
C=\pi_{t_1,\ldots,t_n}^{-1}(B),
\qquad
B\in\mathcal B(E^n).
\]

То есть

\[
C=\{x\in E^T:(x(t_1),\ldots,x(t_n))\in B\}.
\]

Интуитивно цилиндрическое событие смотрит только на конечное число
координат траектории и не ограничивает остальные координаты.

\subsubsection{2.4. Цилиндрическая
сигма-алгебра}\label{ux446ux438ux43bux438ux43dux434ux440ux438ux447ux435ux441ux43aux430ux44f-ux441ux438ux433ux43cux430-ux430ux43bux433ux435ux431ux440ux430}

На \(E^T\) берут цилиндрическую, или произведенную, \(\sigma\)-алгебру:

\[
\mathcal E^T
=
\sigma\{\pi_{t_1,\ldots,t_n}^{-1}(B):
n\ge1,\ t_i\in T,\ B\in\mathcal B(E^n)\}.
\]

Это минимальная \(\sigma\)-алгебра, относительно которой все
координатные отображения \(X_t\) измеримы.

\subsubsection{2.5. Условия согласованности
распределений}\label{ux443ux441ux43bux43eux432ux438ux44f-ux441ux43eux433ux43bux430ux441ux43eux432ux430ux43dux43dux43eux441ux442ux438-ux440ux430ux441ux43fux440ux435ux434ux435ux43bux435ux43dux438ux439}

Теперь предположим, что для каждого конечного набора попарно различных
индексов

\[
t_1,\ldots,t_n\in T
\]

задана вероятностная мера

\[
\mu_{t_1,\ldots,t_n}
\]

на \(E^n\). Она должна играть роль распределения вектора

\[
(X_{t_1},\ldots,X_{t_n}).
\]

Но произвольный набор мер не подойдет. Нужны условия согласованности.

Первое условие - согласованность по перестановкам. Если \(\sigma\) -
перестановка множества \(\{1,\ldots,n\}\), то распределение координат в
новом порядке должно быть тем же самым распределением, только с
переставленными координатами. Для любого \(B\in\mathcal B(E^n)\):

\[
\mu_{t_{\sigma(1)},\ldots,t_{\sigma(n)}}(B)
=
\mu_{t_1,\ldots,t_n}
\{(x_1,\ldots,x_n):(x_{\sigma(1)},\ldots,x_{\sigma(n)})\in B\}.
\]

Проще: если поменяли порядок индексов, нужно так же поменять порядок
координат.

Второе условие - маргинальная согласованность. Если к набору индексов
добавили лишние моменты времени, а потом не смотрим на соответствующие
координаты, то должна получиться старая мера. Для \(m\ge1\):

\[
\mu_{t_1,\ldots,t_n,t_{n+1},\ldots,t_{n+m}}
(B\times E^m)
=
\mu_{t_1,\ldots,t_n}(B),
\qquad
B\in\mathcal B(E^n).
\]

Особенно важный частный случай:

\[
\mu_{t_1,\ldots,t_n,t_{n+1}}(B\times E)
=
\mu_{t_1,\ldots,t_n}(B).
\]

Это означает, что удаление координаты соответствует взятию маргинального
распределения.

\subsubsection{2.6. Теорема
Колмогорова}\label{ux442ux435ux43eux440ux435ux43cux430-ux43aux43eux43bux43cux43eux433ux43eux440ux43eux432ux430}

Теперь формулировка теоремы.

\textbf{Теорема Колмогорова о согласованных конечномерных
распределениях.} Пусть для всех конечных наборов индексов
\(t_1,\ldots,t_n\in T\) задано семейство вероятностных мер

\[
\mu_{t_1,\ldots,t_n}
\]

на \(E^n\), где \(E\) - стандартное борелевское пространство, например
\(E=\mathbb R\). Если это семейство согласовано по перестановкам и
маргинализации, то существует единственная вероятностная мера \(P\) на
цилиндрической \(\sigma\)-алгебре пространства \(E^T\) такая, что для
всех \(n\), всех \(t_1,\ldots,t_n\) и всех \(B\in\mathcal B(E^n)\):

\[
P\{\pi_{t_1,\ldots,t_n}\in B\}
=
P\{x:(x(t_1),\ldots,x(t_n))\in B\}
=
\mu_{t_1,\ldots,t_n}(B).
\]

\subsubsection{2.7. Канонический координатный
процесс}\label{ux43aux430ux43dux43eux43dux438ux447ux435ux441ux43aux438ux439-ux43aux43eux43eux440ux434ux438ux43dux430ux442ux43dux44bux439-ux43fux440ux43eux446ux435ux441ux441}

Если положить

\[
X_t(x)=x(t),
\]

то семейство

\[
\{X_t:t\in T\}
\]

является случайным процессом на вероятностном пространстве

\[
(E^T,\mathcal E^T,P),
\]

и его конечномерные распределения совпадают с заданными мерами
\(\mu_{t_1,\ldots,t_n}\).

Такой процесс называется каноническим координатным процессом.
``Канонический'' означает, что исход \(\omega\) сам является траекторией
\(x(\cdot)\), а случайная величина \(X_t\) просто берет у этой
траектории значение в момент \(t\).

\subsubsection{2.8. Идея доказательства и роль
согласованности}\label{ux438ux434ux435ux44f-ux434ux43eux43aux430ux437ux430ux442ux435ux43bux44cux441ux442ux432ux430-ux438-ux440ux43eux43bux44c-ux441ux43eux433ux43bux430ux441ux43eux432ux430ux43dux43dux43eux441ux442ux438}

Идея доказательства такая. Сначала хотим задать вероятность на цилиндрах
формулой

\[
P\left(\pi_{t_1,\ldots,t_n}^{-1}(B)\right)
=
\mu_{t_1,\ldots,t_n}(B).
\]

Здесь сразу возникает вопрос корректности: одно и то же цилиндрическое
множество можно записать разными способами. Например,

\[
\pi_t^{-1}(A)
\]

можно также записать как

\[
\pi_{t,s}^{-1}(A\times E).
\]

Чтобы обе записи давали одну и ту же вероятность, нужна маргинальная
согласованность:

\[
\mu_{t,s}(A\times E)=\mu_t(A).
\]

Если поменяли порядок координат, нужна согласованность по перестановкам.
Поэтому условия теоремы ровно гарантируют, что вероятность цилиндра
определена однозначно.

Дальше на классе цилиндров получается предмера. По теореме продолжения
меры, идея которой была в Билете 03, эту предмеру продолжают на
\(\sigma\)-алгебру, порожденную цилиндрами. Полученная мера \(P\) и есть
закон процесса на \(E^T\).

Уникальность тоже естественна: цилиндрические множества порождают
\(\mathcal E^T\), а значения меры на порождающем классе уже заданы
конечномерными распределениями. Поэтому другая мера с теми же
конечномерными распределениями совпадет с \(P\) на всей цилиндрической
\(\sigma\)-алгебре.

\subsubsection{2.9. Ограничения теоремы и свойства
траекторий}\label{ux43eux433ux440ux430ux43dux438ux447ux435ux43dux438ux44f-ux442ux435ux43eux440ux435ux43cux44b-ux438-ux441ux432ux43eux439ux441ux442ux432ux430-ux442ux440ux430ux435ux43aux442ux43eux440ux438ux439}

Важно понимать ограничение теоремы. Она строит процесс на пространстве
всех функций \(E^T\) с цилиндрической \(\sigma\)-алгеброй. Она не
утверждает автоматически, что траектории будут непрерывными,
ограниченными, монотонными или какими-то регулярными. Например, чтобы
получить процесс с непрерывными траекториями, нужны дополнительные
условия регулярности. Теорема Колмогорова дает существование процесса с
заданными конечномерными законами, но не гарантирует хороших свойств
траекторий.

\subsubsection{2.10. Дискретное
время}\label{ux434ux438ux441ux43aux440ux435ux442ux43dux43eux435-ux432ux440ux435ux43cux44f}

В дискретном времени, когда \(T=\mathbb N\), картина особенно проста.
Пространство траекторий:

\[
E^{\mathbb N}.
\]

Цилиндрическое событие задает условия на конечное число координат:

\[
\{x:x_0\in A_0,\ldots,x_n\in A_n\}.
\]

Если все конечномерные распределения согласованы, то существует мера на
пространстве последовательностей, то есть закон случайной
последовательности. Это напрямую ведет к Билету 16.

\subsection{3. Основные
определения}\label{ux43eux441ux43dux43eux432ux43dux44bux435-ux43eux43fux440ux435ux434ux435ux43bux435ux43dux438ux44f-14}

\subsubsection{Конечномерное распределение
процесса}\label{ux43aux43eux43dux435ux447ux43dux43eux43cux435ux440ux43dux43eux435-ux440ux430ux441ux43fux440ux435ux434ux435ux43bux435ux43dux438ux435-ux43fux440ux43eux446ux435ux441ux441ux430-2}

Для процесса \(\{X_t:t\in T\}\) и конечного набора \(t_1,\ldots,t_n\)
это распределение случайного вектора

\[
(X_{t_1},\ldots,X_{t_n})
\]

на \(E^n\):

\[
\mu_{t_1,\ldots,t_n}(B)
=
P\{(X_{t_1},\ldots,X_{t_n})\in B\}.
\]

\subsubsection{Пространство
траекторий}\label{ux43fux440ux43eux441ux442ux440ux430ux43dux441ux442ux432ux43e-ux442ux440ux430ux435ux43aux442ux43eux440ux438ux439}

Множество всех функций из \(T\) в \(E\):

\[
E^T=\{x:T\to E\}.
\]

Для дискретной последовательности это \(E^{\mathbb N}\).

\subsubsection{Координатное
отображение}\label{ux43aux43eux43eux440ux434ux438ux43dux430ux442ux43dux43eux435-ux43eux442ux43eux431ux440ux430ux436ux435ux43dux438ux435}

\[
X_t:E^T\to E,
\qquad
X_t(x)=x(t).
\]

\subsubsection{Конечномерная
проекция}\label{ux43aux43eux43dux435ux447ux43dux43eux43cux435ux440ux43dux430ux44f-ux43fux440ux43eux435ux43aux446ux438ux44f}

\[
\pi_{t_1,\ldots,t_n}:E^T\to E^n,
\qquad
\pi_{t_1,\ldots,t_n}(x)=(x(t_1),\ldots,x(t_n)).
\]

\subsubsection{Цилиндрическое
множество}\label{ux446ux438ux43bux438ux43dux434ux440ux438ux447ux435ux441ux43aux43eux435-ux43cux43dux43eux436ux435ux441ux442ux432ux43e}

\[
\pi_{t_1,\ldots,t_n}^{-1}(B),
\qquad
B\in\mathcal B(E^n).
\]

Это событие, зависящее только от конечного числа координат.

\subsubsection{\texorpdfstring{Цилиндрическая
\(\sigma\)-алгебра}{Цилиндрическая \textbackslash sigma-алгебра}}\label{ux446ux438ux43bux438ux43dux434ux440ux438ux447ux435ux441ux43aux430ux44f-sigma-ux430ux43bux433ux435ux431ux440ux430}

\[
\mathcal E^T
=
\sigma\{\pi_{t_1,\ldots,t_n}^{-1}(B):
B\in\mathcal B(E^n)\}.
\]

\subsubsection{Согласованное семейство
мер}\label{ux441ux43eux433ux43bux430ux441ux43eux432ux430ux43dux43dux43eux435-ux441ux435ux43cux435ux439ux441ux442ux432ux43e-ux43cux435ux440}

Семейство конечномерных мер \(\mu_{t_1,\ldots,t_n}\), удовлетворяющее
двум условиям: согласованность по перестановкам и согласованность по
маргинализации.

\subsubsection{Канонический
процесс}\label{ux43aux430ux43dux43eux43dux438ux447ux435ux441ux43aux438ux439-ux43fux440ux43eux446ux435ux441ux441}

Процесс на пространстве \(E^T\), заданный координатами:

\[
X_t(x)=x(t).
\]

\subsection{4. Основные теоремы и
свойства}\label{ux43eux441ux43dux43eux432ux43dux44bux435-ux442ux435ux43eux440ux435ux43cux44b-ux438-ux441ux432ux43eux439ux441ux442ux432ux430-14}

\subsubsection{Согласованность по
перестановкам}\label{ux441ux43eux433ux43bux430ux441ux43eux432ux430ux43dux43dux43eux441ux442ux44c-ux43fux43e-ux43fux435ux440ux435ux441ux442ux430ux43dux43eux432ux43aux430ux43c}

Для любой перестановки \(\sigma\) координаты можно переставлять только
вместе с индексами:

\[
\mu_{t_{\sigma(1)},\ldots,t_{\sigma(n)}}(B)
=
\mu_{t_1,\ldots,t_n}
\{x:(x_{\sigma(1)},\ldots,x_{\sigma(n)})\in B\}.
\]

\subsubsection{Маргинальная
согласованность}\label{ux43cux430ux440ux433ux438ux43dux430ux43bux44cux43dux430ux44f-ux441ux43eux433ux43bux430ux441ux43eux432ux430ux43dux43dux43eux441ux442ux44c}

Удаление координат соответствует взятию маргинального распределения:

\[
\mu_{t_1,\ldots,t_n,t_{n+1},\ldots,t_{n+m}}
(B\times E^m)
=
\mu_{t_1,\ldots,t_n}(B).
\]

\subsubsection{Теорема
Колмогорова}\label{ux442ux435ux43eux440ux435ux43cux430-ux43aux43eux43bux43cux43eux433ux43eux440ux43eux432ux430-1}

В стандартном борелевском случае всякое согласованное семейство
конечномерных распределений задает единственную вероятностную меру на
цилиндрической \(\sigma\)-алгебре \(E^T\).

\subsubsection{Координатный процесс имеет заданные конечномерные
распределения}\label{ux43aux43eux43eux440ux434ux438ux43dux430ux442ux43dux44bux439-ux43fux440ux43eux446ux435ux441ux441-ux438ux43cux435ux435ux442-ux437ux430ux434ux430ux43dux43dux44bux435-ux43aux43eux43dux435ux447ux43dux43eux43cux435ux440ux43dux44bux435-ux440ux430ux441ux43fux440ux435ux434ux435ux43bux435ux43dux438ux44f}

\[
P\{(X_{t_1},\ldots,X_{t_n})\in B\}
=
\mu_{t_1,\ldots,t_n}(B).
\]

\subsubsection{\texorpdfstring{Единственность закона на цилиндрической
\(\sigma\)-алгебре}{Единственность закона на цилиндрической \textbackslash sigma-алгебре}}\label{ux435ux434ux438ux43dux441ux442ux432ux435ux43dux43dux43eux441ux442ux44c-ux437ux430ux43aux43eux43dux430-ux43dux430-ux446ux438ux43bux438ux43dux434ux440ux438ux447ux435ux441ux43aux43eux439-sigma-ux430ux43bux433ux435ux431ux440ux435}

Если две меры на \(E^T\) совпадают на всех цилиндрах, то они совпадают
на всей \(\sigma\)-алгебре, порожденной цилиндрами.

\subsubsection{Теорема не дает регулярность
траекторий}\label{ux442ux435ux43eux440ux435ux43cux430-ux43dux435-ux434ux430ux435ux442-ux440ux435ux433ux443ux43bux44fux440ux43dux43eux441ux442ux44c-ux442ux440ux430ux435ux43aux442ux43eux440ux438ux439}

Непрерывность, правосторонняя непрерывность, ограниченность и другие
свойства траекторий требуют отдельных условий.

\subsection{5. Примеры}\label{ux43fux440ux438ux43cux435ux440ux44b-14}

\subsubsection{Независимая одинаково распределенная
последовательность}\label{ux43dux435ux437ux430ux432ux438ux441ux438ux43cux430ux44f-ux43eux434ux438ux43dux430ux43aux43eux432ux43e-ux440ux430ux441ux43fux440ux435ux434ux435ux43bux435ux43dux43dux430ux44f-ux43fux43eux441ux43bux435ux434ux43eux432ux430ux442ux435ux43bux44cux43dux43eux441ux442ux44c}

Пусть \(\nu\) - вероятностная мера на \(E\). Для конечного набора разных
индексов \(t_1,\ldots,t_n\in\mathbb N\) зададим

\[
\mu_{t_1,\ldots,t_n}
=
\nu^{\otimes n}.
\]

То есть все координаты имеют одно и то же распределение \(\nu\) и
независимы. Это семейство согласовано: перестановка координат не меняет
произведение одинаковых мер, а маргинализация произведения дает
произведение меньшего числа множителей:

\[
\nu^{\otimes(n+1)}(B\times E)=\nu^{\otimes n}(B).
\]

По теореме Колмогорова существует случайная последовательность
независимых одинаково распределенных величин с законом \(\nu\).

\subsubsection{Последовательность независимых, но не одинаково
распределенных
величин}\label{ux43fux43eux441ux43bux435ux434ux43eux432ux430ux442ux435ux43bux44cux43dux43eux441ux442ux44c-ux43dux435ux437ux430ux432ux438ux441ux438ux43cux44bux445-ux43dux43e-ux43dux435-ux43eux434ux438ux43dux430ux43aux43eux432ux43e-ux440ux430ux441ux43fux440ux435ux434ux435ux43bux435ux43dux43dux44bux445-ux432ux435ux43bux438ux447ux438ux43d}

Пусть для каждого \(k\in\mathbb N\) задана мера \(\nu_k\) на \(E\). Для
индексов \(k_1,\ldots,k_n\) задаем

\[
\mu_{k_1,\ldots,k_n}
=
\nu_{k_1}\otimes\cdots\otimes\nu_{k_n}.
\]

Это семейство тоже согласовано. Значит можно построить
последовательность независимых величин с заранее заданными разными
одномерными распределениями.

\subsubsection{Марковская цепь по начальному распределению и переходной
матрице}\label{ux43cux430ux440ux43aux43eux432ux441ux43aux430ux44f-ux446ux435ux43fux44c-ux43fux43e-ux43dux430ux447ux430ux43bux44cux43dux43eux43cux443-ux440ux430ux441ux43fux440ux435ux434ux435ux43bux435ux43dux438ux44e-ux438-ux43fux435ux440ux435ux445ux43eux434ux43dux43eux439-ux43cux430ux442ux440ux438ux446ux435}

Пусть пространство состояний счетное, начальное распределение равно
\(\pi\), а переходная матрица равна \(P=(p_{ij})\). Для
\(i_0,\ldots,i_n\) задаем конечномерные вероятности:

\[
\mu_{0,\ldots,n}(i_0,\ldots,i_n)
=
\pi_{i_0}p_{i_0i_1}p_{i_1i_2}\cdots p_{i_{n-1}i_n}.
\]

Эти распределения согласованы: если просуммировать по последней
координате, то из условия

\[
\sum_j p_{ij}=1
\]

получается распределение первых \(n\) координат. Поэтому теорема
Колмогорова дает существование марковской цепи с заданными \(\pi\) и
\(P\). Это связь с Билетом 22 и Билетом 23.

\subsubsection{Гауссовский
процесс}\label{ux433ux430ux443ux441ux441ux43eux432ux441ux43aux438ux439-ux43fux440ux43eux446ux435ux441ux441}

Чтобы задать гауссовский процесс, задают среднюю функцию

\[
m(t)
\]

и ковариационную функцию

\[
K(s,t).
\]

Для каждого конечного набора \(t_1,\ldots,t_n\) берут гауссовское
распределение в \(\mathbb R^n\) со средним вектором

\[
(m(t_1),\ldots,m(t_n))
\]

и ковариационной матрицей

\[
(K(t_i,t_j))_{i,j=1}^n.
\]

Чтобы это было корректно, все такие матрицы должны быть неотрицательно
определены. Тогда гауссовские конечномерные распределения согласованы и
по теореме Колмогорова существует гауссовский процесс. Это связь с
Билетом 18.

\subsubsection{Пример
несогласованности}\label{ux43fux440ux438ux43cux435ux440-ux43dux435ux441ux43eux433ux43bux430ux441ux43eux432ux430ux43dux43dux43eux441ux442ux438}

Пусть пытаются задать

\[
P\{X_1=0\}=1
\]

и одновременно

\[
P\{(X_1,X_2)=(1,1)\}=1.
\]

Такого процесса не существует, потому что маргинал двумерного
распределения по первой координате дает

\[
P\{X_1=1\}=1,
\]

что противоречит заданному одномерному распределению. Это показывает,
зачем нужна согласованность.

\subsection{6. Схемы доказательств /
обоснований}\label{ux441ux445ux435ux43cux44b-ux434ux43eux43aux430ux437ux430ux442ux435ux43bux44cux441ux442ux432-ux43eux431ux43eux441ux43dux43eux432ux430ux43dux438ux439-7}

\subsubsection{Почему условия согласованности
необходимы}\label{ux43fux43eux447ux435ux43cux443-ux443ux441ux43bux43eux432ux438ux44f-ux441ux43eux433ux43bux430ux441ux43eux432ux430ux43dux43dux43eux441ux442ux438-ux43dux435ux43eux431ux445ux43eux434ux438ux43cux44b}

Если процесс уже существует, то для любого конечного набора

\[
(X_{t_1},\ldots,X_{t_n})
\]

его распределение не должно зависеть от того, как мы записали одно и то
же событие.

Если поменяли порядок координат, событие не изменилось по смыслу,
поменялась только запись. Значит нужна согласованность по перестановкам.

Если из вектора

\[
(X_{t_1},\ldots,X_{t_n},X_s)
\]

не смотреть на последнюю координату, то должно получиться распределение

\[
(X_{t_1},\ldots,X_{t_n}).
\]

Значит нужна маргинальная согласованность.

\subsubsection{Идея построения
меры}\label{ux438ux434ux435ux44f-ux43fux43eux441ux442ux440ux43eux435ux43dux438ux44f-ux43cux435ux440ux44b}

\begin{enumerate}
\def\labelenumi{\arabic{enumi}.}
\tightlist
\item
  На цилиндрах задаем:
\end{enumerate}

\[
P(\pi_{t_1,\ldots,t_n}^{-1}(B))
=
\mu_{t_1,\ldots,t_n}(B).
\]

\begin{enumerate}
\def\labelenumi{\arabic{enumi}.}
\setcounter{enumi}{1}
\item
  Согласованность гарантирует, что это определение не зависит от
  представления цилиндра.
\item
  Получаем предмеру на алгебре цилиндров.
\item
  По теореме продолжения меры продолжаем ее на \(\sigma\)-алгебру,
  порожденную цилиндрами.
\item
  На полученном вероятностном пространстве берем координатный процесс:
\end{enumerate}

\[
X_t(x)=x(t).
\]

\subsubsection{Почему закон
единственен}\label{ux43fux43eux447ux435ux43cux443-ux437ux430ux43aux43eux43d-ux435ux434ux438ux43dux441ux442ux432ux435ux43dux435ux43d}

Если две вероятностные меры на \(E^T\) имеют одинаковые конечномерные
распределения, то они совпадают на всех цилиндрах. Цилиндры порождают
\(\mathcal E^T\), значит меры совпадают на всей \(\mathcal E^T\).

\subsection{7. Что написать на
доске}\label{ux447ux442ux43e-ux43dux430ux43fux438ux441ux430ux442ux44c-ux43dux430-ux434ux43eux441ux43aux435-14}

\begin{enumerate}
\def\labelenumi{\arabic{enumi}.}
\tightlist
\item
  Пространство траекторий:
\end{enumerate}

\[
E^T=\{x:T\to E\}.
\]

\begin{enumerate}
\def\labelenumi{\arabic{enumi}.}
\setcounter{enumi}{1}
\tightlist
\item
  Проекция:
\end{enumerate}

\[
\pi_{t_1,\ldots,t_n}(x)=(x(t_1),\ldots,x(t_n)).
\]

\begin{enumerate}
\def\labelenumi{\arabic{enumi}.}
\setcounter{enumi}{2}
\tightlist
\item
  Цилиндры:
\end{enumerate}

\[
\pi_{t_1,\ldots,t_n}^{-1}(B),
\qquad
B\in\mathcal B(E^n).
\]

\begin{enumerate}
\def\labelenumi{\arabic{enumi}.}
\setcounter{enumi}{3}
\tightlist
\item
  Согласованность по маргинализации:
\end{enumerate}

\[
\mu_{t_1,\ldots,t_n,s}(B\times E)
=
\mu_{t_1,\ldots,t_n}(B).
\]

\begin{enumerate}
\def\labelenumi{\arabic{enumi}.}
\setcounter{enumi}{4}
\tightlist
\item
  Согласованность по перестановкам:
\end{enumerate}

\[
\mu_{t_{\sigma(1)},\ldots,t_{\sigma(n)}}
=
\mu_{t_1,\ldots,t_n}\circ p_\sigma^{-1},
\]

где

\[
p_\sigma(x_1,\ldots,x_n)=(x_{\sigma(1)},\ldots,x_{\sigma(n)}).
\]

\begin{enumerate}
\def\labelenumi{\arabic{enumi}.}
\setcounter{enumi}{5}
\tightlist
\item
  Вывод теоремы:
\end{enumerate}

\[
\exists!P\ \text{на}\ (E^T,\mathcal E^T):
\quad
P\{\pi_{t_1,\ldots,t_n}\in B\}
=
\mu_{t_1,\ldots,t_n}(B).
\]

Здесь \(\mathcal E^T\) - цилиндрическая \(\sigma\)-алгебра.

\begin{enumerate}
\def\labelenumi{\arabic{enumi}.}
\setcounter{enumi}{6}
\tightlist
\item
  Канонический процесс:
\end{enumerate}

\[
X_t(x)=x(t).
\]

\subsection{8. Типичные дополнительные
вопросы}\label{ux442ux438ux43fux438ux447ux43dux44bux435-ux434ux43eux43fux43eux43bux43dux438ux442ux435ux43bux44cux43dux44bux435-ux432ux43eux43fux440ux43eux441ux44b-14}

\textbf{Что является главным объектом теоремы?}

Семейство конечномерных распределений

\[
\{\mu_{t_1,\ldots,t_n}\}.
\]

Теорема говорит, когда оно является семейством конечномерных
распределений некоторого процесса.

\textbf{Почему одних одномерных распределений недостаточно?}

Потому что они не задают зависимость между разными моментами времени.
Например, две величины могут иметь одинаковые одномерные распределения,
но быть независимыми, полностью зависимыми или иметь любую допустимую
совместную структуру. Нужны все совместные распределения конечных
наборов координат.

\textbf{Что такое цилиндрическое событие?}

Это событие, зависящее только от конечного числа координат траектории:

\[
\{x:(x(t_1),\ldots,x(t_n))\in B\}.
\]

Остальные координаты не ограничиваются.

\textbf{Зачем нужна согласованность по перестановкам?}

Чтобы распределение не зависело от порядка записи одного и того же
конечного набора координат. Вектор \((X_s,X_t)\) и вектор \((X_t,X_s)\)
имеют меры, связанные перестановкой координат.

\textbf{Зачем нужна маргинальная согласованность?}

Чтобы распределение меньшего набора координат совпадало с маргиналом
распределения большего набора. Например:

\[
P\{(X_s,X_t)\in B\times E\}
=
P\{X_s\in B\}.
\]

\textbf{Что именно строит теорема Колмогорова?}

Она строит вероятностную меру \(P\) на пространстве траекторий \(E^T\) с
цилиндрической \(\sigma\)-алгеброй. После этого координаты
\(X_t(x)=x(t)\) становятся случайным процессом.

\textbf{Гарантирует ли теорема непрерывные траектории?}

Нет. Она гарантирует только существование процесса с заданными
конечномерными распределениями. Регулярность траекторий требует
дополнительных условий.

\textbf{В чем разница между законом процесса и конечномерными
распределениями?}

Закон процесса - это мера на всем пространстве траекторий \(E^T\).
Конечномерные распределения - это ее маргиналы при конечномерных
проекциях.

\textbf{Где здесь используется теорема продолжения меры?}

Сначала мера задается на цилиндрах, то есть на событиях конечномерного
вида. Затем ее продолжают на \(\sigma\)-алгебру, порожденную цилиндрами.

\textbf{Как связана теорема с Билетом 16?}

В Билете 16 дискретный процесс рассматривается как случайная
последовательность, то есть как случайный элемент пространства
\(E^{\mathbb N}\). Теорема Колмогорова как раз дает меру на этом
пространстве по согласованным конечномерным распределениям.

\subsection{9. Типичные
ошибки}\label{ux442ux438ux43fux438ux447ux43dux44bux435-ux43eux448ux438ux431ux43aux438-14}

\begin{itemize}
\tightlist
\item
  Задавать только одномерные распределения и считать, что процесс уже
  определен.
\item
  Забывать маргинальную согласованность.
\item
  Забывать согласованность по перестановкам.
\item
  Путать конечномерное распределение и траекторию.
\item
  Путать закон процесса на \(E^T\) и отдельные распределения на \(E^n\).
\item
  Считать, что теорема автоматически дает непрерывные или регулярные
  траектории.
\item
  Не указывать, на какой \(\sigma\)-алгебре строится мера.
\item
  Формулировать теорему без условий на пространство состояний; для
  аккуратного ответа безопасно говорить о \(E=\mathbb R\) или
  стандартном борелевском пространстве.
\end{itemize}

\subsection{10. Связи с другими
билетами}\label{ux441ux432ux44fux437ux438-ux441-ux434ux440ux443ux433ux438ux43cux438-ux431ux438ux43bux435ux442ux430ux43cux438-14}

\begin{itemize}
\tightlist
\item
  Билет 03: идея продолжения меры с алгебры цилиндров на порожденную
  \(\sigma\)-алгебру.
\item
  Билет 05: координатный процесс состоит из измеримых отображений.
\item
  Билет 06: конечномерные распределения процесса - распределения
  случайных векторов.
\item
  Билет 16: дискретный процесс как случайный элемент пространства
  последовательностей \(E^{\mathbb N}\).
\item
  Билет 17: стационарность в узком смысле формулируется через
  конечномерные распределения.
\item
  Билет 18: гауссовский процесс задается согласованными гауссовскими
  конечномерными распределениями.
\item
  Билет 22 и Билет 23: марковские цепи строятся по начальному
  распределению и переходным вероятностям.
\end{itemize}

\subsection{11. Минимум на
удовлетворительно}\label{ux43cux438ux43dux438ux43cux443ux43c-ux43dux430-ux443ux434ux43eux432ux43bux435ux442ux432ux43eux440ux438ux442ux435ux43bux44cux43dux43e-14}

Нужно сказать:

\begin{enumerate}
\def\labelenumi{\arabic{enumi}.}
\tightlist
\item
  Конечномерное распределение процесса - это закон вектора
\end{enumerate}

\[
(X_{t_1},\ldots,X_{t_n}).
\]

\begin{enumerate}
\def\labelenumi{\arabic{enumi}.}
\setcounter{enumi}{1}
\item
  Чтобы набор таких распределений задавал процесс, он должен быть
  согласован.
\item
  Два условия согласованности:
\end{enumerate}

\[
\text{перестановка индексов} \Longleftrightarrow \text{перестановка координат},
\]

и

\[
\mu_{t_1,\ldots,t_n,s}(B\times E)
=
\mu_{t_1,\ldots,t_n}(B).
\]

\begin{enumerate}
\def\labelenumi{\arabic{enumi}.}
\setcounter{enumi}{3}
\item
  Тогда существует процесс на \(E^T\) с этими конечномерными
  распределениями.
\item
  Главная ловушка: теорема не дает регулярность траекторий.
\end{enumerate}

\subsection{12. Минимум на
хорошо}\label{ux43cux438ux43dux438ux43cux443ux43c-ux43dux430-ux445ux43eux440ux43eux448ux43e-14}

Помимо минимума, нужно объяснить:

\begin{itemize}
\tightlist
\item
  что такое пространство траекторий \(E^T\);
\item
  что такое цилиндрическое событие;
\item
  что такое координатный процесс \(X_t(x)=x(t)\);
\item
  почему согласованность нужна для корректного задания меры на
  цилиндрах;
\item
  как идея доказательства связана с продолжением меры;
\item
  привести пример iid-последовательности или марковской цепи.
\end{itemize}

\subsection{13. Что добавить для
отлично}\label{ux447ux442ux43e-ux434ux43eux431ux430ux432ux438ux442ux44c-ux434ux43bux44f-ux43eux442ux43bux438ux447ux43dux43e-14}

Для сильного ответа стоит добавить:

\begin{itemize}
\tightlist
\item
  аккуратную формулировку через конечномерные проекции
  \(\pi_{t_1,\ldots,t_n}\);
\item
  единственность меры на цилиндрической \(\sigma\)-алгебре;
\item
  различие между законом процесса и конечномерными распределениями;
\item
  предупреждение про отсутствие автоматической регулярности траекторий;
\item
  пример несогласованного семейства;
\item
  пример гауссовского процесса через среднюю и ковариационную функции;
\item
  связь с узкой стационарностью, марковскими цепями и случайными
  последовательностями.
\end{itemize}

\subsection{14. Устная версия ответа на 5
минут}\label{ux443ux441ux442ux43dux430ux44f-ux432ux435ux440ux441ux438ux44f-ux43eux442ux432ux435ux442ux430-ux43dux430-5-ux43cux438ux43dux443ux442-14}

\begin{enumerate}
\def\labelenumi{\arabic{enumi}.}
\tightlist
\item
  Теорема Колмогорова нужна, чтобы по конечномерным распределениям
  построить случайный процесс.
\item
  Пусть для каждого конечного набора \(t_1,\ldots,t_n\) задана мера
  \(\mu_{t_1,\ldots,t_n}\) на \(E^n\).
\item
  Эти меры должны быть согласованы по перестановкам: если меняем порядок
  индексов, меняется порядок координат.
\item
  Они должны быть согласованы по маргинализации:
\end{enumerate}

\[
\mu_{t_1,\ldots,t_n,s}(B\times E)
=
\mu_{t_1,\ldots,t_n}(B).
\]

\begin{enumerate}
\def\labelenumi{\arabic{enumi}.}
\setcounter{enumi}{4}
\tightlist
\item
  Рассматриваем пространство траекторий
\end{enumerate}

\[
E^T=\{x:T\to E\}.
\]

\begin{enumerate}
\def\labelenumi{\arabic{enumi}.}
\setcounter{enumi}{5}
\tightlist
\item
  На нем есть координатные отображения
\end{enumerate}

\[
X_t(x)=x(t).
\]

\begin{enumerate}
\def\labelenumi{\arabic{enumi}.}
\setcounter{enumi}{6}
\tightlist
\item
  Цилиндры имеют вид
\end{enumerate}

\[
\{x:(x(t_1),\ldots,x(t_n))\in B\}.
\]

\begin{enumerate}
\def\labelenumi{\arabic{enumi}.}
\setcounter{enumi}{7}
\tightlist
\item
  По заданным конечномерным мерам задаем вероятность цилиндра:
\end{enumerate}

\[
P(\pi_{t_1,\ldots,t_n}^{-1}(B))
=
\mu_{t_1,\ldots,t_n}(B).
\]

\begin{enumerate}
\def\labelenumi{\arabic{enumi}.}
\setcounter{enumi}{8}
\tightlist
\item
  Согласованность делает это определение корректным.
\item
  Потом мера продолжается на \(\sigma\)-алгебру, порожденную цилиндрами.
\item
  Полученный координатный процесс имеет ровно заданные конечномерные
  распределения.
\item
  Важно: теорема дает существование закона процесса, но не гарантирует
  непрерывность или другие свойства траекторий.
\end{enumerate}

\subsection{15. Если осталось 2 минуты перед
экзаменом}\label{ux435ux441ux43bux438-ux43eux441ux442ux430ux43bux43eux441ux44c-2-ux43cux438ux43dux443ux442ux44b-ux43fux435ux440ux435ux434-ux44dux43aux437ux430ux43cux435ux43dux43eux43c-14}

Запомнить:

\[
E^T=\{x:T\to E\},
\qquad
X_t(x)=x(t).
\]

Цилиндр:

\[
\pi_{t_1,\ldots,t_n}^{-1}(B).
\]

Главное условие:

\[
\mu_{t_1,\ldots,t_n,s}(B\times E)
=
\mu_{t_1,\ldots,t_n}(B),
\]

плюс согласованность по перестановкам.

Главная формулировка: согласованное семейство конечномерных
распределений задает единственную меру на пространстве траекторий с
цилиндрической \(\sigma\)-алгеброй, а координаты становятся процессом.

Главная ловушка: теорема не утверждает регулярность траекторий; она
только строит закон процесса.

\newpage

\section{Билет 16. Дискретный случайный процесс как последовательность
случайных величин и как случайная
последовательность}\label{ux431ux438ux43bux435ux442-16.-ux434ux438ux441ux43aux440ux435ux442ux43dux44bux439-ux441ux43bux443ux447ux430ux439ux43dux44bux439-ux43fux440ux43eux446ux435ux441ux441-ux43aux430ux43a-ux43fux43eux441ux43bux435ux434ux43eux432ux430ux442ux435ux43bux44cux43dux43eux441ux442ux44c-ux441ux43bux443ux447ux430ux439ux43dux44bux445-ux432ux435ux43bux438ux447ux438ux43d-ux438-ux43aux430ux43a-ux441ux43bux443ux447ux430ux439ux43dux430ux44f-ux43fux43eux441ux43bux435ux434ux43eux432ux430ux442ux435ux43bux44cux43dux43eux441ux442ux44c}

Источники: Ширяев 2; Сердобольская, лекция 1.

\subsection{1. Короткий
ответ}\label{ux43aux43eux440ux43eux442ux43aux438ux439-ux43eux442ux432ux435ux442-15}

Дискретный случайный процесс можно понимать двумя эквивалентными
способами.

Первый способ: это семейство случайных величин

\[
\{X_n:n\in\mathbb N_0\}
\]

на одном вероятностном пространстве

\[
(\Omega,\mathcal F,P),
\]

где каждое

\[
X_n:\Omega\to E
\]

измеримо.

Второй способ: это одна случайная последовательность, то есть одно
измеримое отображение

\[
X:\Omega\to E^{\mathbb N_0},
\qquad
X(\omega)=(X_0(\omega),X_1(\omega),\ldots).
\]

Закон процесса - это распределение этого отображения:

\[
P_X(A)=P\{\omega:X(\omega)\in A\},
\qquad
A\in\mathcal E^{\mathbb N_0}.
\]

Конечномерные распределения процесса - это маргиналы закона \(P_X\) на
конечные наборы координат:

\[
P\{(X_{n_1},\ldots,X_{n_k})\in B\}
=
P_X\{x:(x_{n_1},\ldots,x_{n_k})\in B\}.
\]

\subsection{2. Полный
ответ}\label{ux43fux43eux43bux43dux44bux439-ux43eux442ux432ux435ux442-15}

\subsubsection{2.1. Введение и дискретное
время}\label{ux432ux432ux435ux434ux435ux43dux438ux435-ux438-ux434ux438ux441ux43aux440ux435ux442ux43dux43eux435-ux432ux440ux435ux43cux44f}

Этот билет продолжает Билет 15. Там теорема Колмогорова объясняла, как
по согласованным конечномерным распределениям построить меру на
пространстве траекторий. Здесь нужно разобрать самый важный для курса
частный случай: дискретное время, то есть случайные последовательности.

Пусть время дискретно:

\[
n=0,1,2,\ldots
\]

или

\[
n\in\mathbb Z.
\]

Обычно в этом билете достаточно рассматривать

\[
\mathbb N_0=\{0,1,2,\ldots\}.
\]

Пусть пространство состояний равно \(E\), например \(E=\mathbb R\),
\(E=\mathbb R^d\) или счетное множество состояний. На \(E\) задана
\(\sigma\)-алгебра \(\mathcal E\), обычно борелевская \(\mathcal B(E)\).

\subsubsection{2.2. Первая точка зрения: последовательность случайных
величин}\label{ux43fux435ux440ux432ux430ux44f-ux442ux43eux447ux43aux430-ux437ux440ux435ux43dux438ux44f-ux43fux43eux441ux43bux435ux434ux43eux432ux430ux442ux435ux43bux44cux43dux43eux441ux442ux44c-ux441ux43bux443ux447ux430ux439ux43dux44bux445-ux432ux435ux43bux438ux447ux438ux43d}

Первая точка зрения: дискретный случайный процесс - это
последовательность случайных элементов

\[
X_0,X_1,X_2,\ldots
\]

на одном вероятностном пространстве

\[
(\Omega,\mathcal F,P),
\]

где

\[
X_n:(\Omega,\mathcal F)\to(E,\mathcal E)
\]

измеримо для каждого \(n\).

Здесь \(n\) - время, а \(\omega\) - исход. Для фиксированного \(n\)
величина

\[
X_n(\omega)
\]

показывает состояние процесса в момент времени \(n\).

\subsubsection{2.3. Вторая точка зрения: пространство
траекторий}\label{ux432ux442ux43eux440ux430ux44f-ux442ux43eux447ux43aux430-ux437ux440ux435ux43dux438ux44f-ux43fux440ux43eux441ux442ux440ux430ux43dux441ux442ux432ux43e-ux442ux440ux430ux435ux43aux442ux43eux440ux438ux439}

Вторая точка зрения: весь процесс можно рассматривать как одно
отображение в пространство последовательностей:

\[
X:\Omega\to E^{\mathbb N_0},
\]

где

\[
X(\omega)=(X_0(\omega),X_1(\omega),X_2(\omega),\ldots).
\]

То есть каждому исходу \(\omega\) сопоставляется целая траектория
процесса.

Пространство последовательностей:

\[
E^{\mathbb N_0}
=
\{x=(x_0,x_1,x_2,\ldots):x_n\in E\}.
\]

На нем берут произведенную \(\sigma\)-алгебру

\[
\mathcal E^{\mathbb N_0}
=
\sigma\{\pi_n^{-1}(A):n\ge0,\ A\in\mathcal E\},
\]

где

\[
\pi_n:E^{\mathbb N_0}\to E,
\qquad
\pi_n(x)=x_n
\]

\begin{itemize}
\tightlist
\item
  координатная проекция.
\end{itemize}

\subsubsection{2.4. Конечномерные проекции и цилиндрические
множества}\label{ux43aux43eux43dux435ux447ux43dux43eux43cux435ux440ux43dux44bux435-ux43fux440ux43eux435ux43aux446ux438ux438-ux438-ux446ux438ux43bux438ux43dux434ux440ux438ux447ux435ux441ux43aux438ux435-ux43cux43dux43eux436ux435ux441ux442ux432ux430}

Более общие конечномерные проекции имеют вид

\[
\pi_{n_1,\ldots,n_k}:E^{\mathbb N_0}\to E^k,
\qquad
\pi_{n_1,\ldots,n_k}(x)=(x_{n_1},\ldots,x_{n_k}).
\]

Цилиндрическое множество - это событие, зависящее только от конечного
числа координат:

\[
C=\pi_{n_1,\ldots,n_k}^{-1}(B),
\qquad
B\in\mathcal E^k.
\]

Например,

\[
C=\{x:x_0\in A_0,\ x_1\in A_1,\ x_5\in A_5\}
\]

является цилиндром: он накладывает условия только на координаты
\(0,1,5\), а остальные координаты свободны.

\subsubsection{2.5. Доказательство эквивалентности точек
зрения}\label{ux434ux43eux43aux430ux437ux430ux442ux435ux43bux44cux441ux442ux432ux43e-ux44dux43aux432ux438ux432ux430ux43bux435ux43dux442ux43dux43eux441ux442ux438-ux442ux43eux447ux435ux43a-ux437ux440ux435ux43dux438ux44f}

Теперь важно показать эквивалентность двух точек зрения.

Если задана последовательность случайных величин \(X_n\), то определяем

\[
X(\omega)=(X_0(\omega),X_1(\omega),\ldots).
\]

Нужно проверить, что \(X\) измеримо как отображение в
\(E^{\mathbb N_0}\). Поскольку произведенная \(\sigma\)-алгебра
порождена цилиндрами, достаточно проверить прообразы цилиндров. Для
одного координатного цилиндра:

\[
X^{-1}(\pi_n^{-1}(A))
=
\{\omega:X_n(\omega)\in A\}
=
X_n^{-1}(A)\in\mathcal F.
\]

Значит \(X\) измеримо. Для конечномерного цилиндра:

\[
X^{-1}(\pi_{n_1,\ldots,n_k}^{-1}(B))
=
\{\omega:(X_{n_1}(\omega),\ldots,X_{n_k}(\omega))\in B\}\in\mathcal F,
\]

потому что вектор

\[
(X_{n_1},\ldots,X_{n_k})
\]

измерим.

Обратно, если задана случайная последовательность

\[
X:\Omega\to E^{\mathbb N_0},
\]

то отдельные координаты получаются как композиции:

\[
X_n=\pi_n\circ X.
\]

Так как \(\pi_n\) измерима, а \(X\) измеримо, то каждое \(X_n\) -
случайная величина. Значит одно случайное отображение в пространство
последовательностей задает семейство случайных величин.

Эта эквивалентность - центральная мысль билета:

\[
\{X_n\}_{n\ge0}
\quad\Longleftrightarrow\quad
X=(X_0,X_1,\ldots):\Omega\to E^{\mathbb N_0}.
\]

\subsubsection{2.6. Закон процесса и конечномерные
распределения}\label{ux437ux430ux43aux43eux43d-ux43fux440ux43eux446ux435ux441ux441ux430-ux438-ux43aux43eux43dux435ux447ux43dux43eux43cux435ux440ux43dux44bux435-ux440ux430ux441ux43fux440ux435ux434ux435ux43bux435ux43dux438ux44f}

Теперь закон процесса. Как и для любого случайного элемента из Билета
06, распределение случайной последовательности - это образ вероятностной
меры:

\[
P_X=P\circ X^{-1}.
\]

То есть для \(A\in\mathcal E^{\mathbb N_0}\):

\[
P_X(A)=P\{\omega:X(\omega)\in A\}.
\]

Это мера на пространстве последовательностей. Именно она называется
законом процесса.

Конечномерные распределения получаются из закона процесса проекциями.
Для \(n_1,\ldots,n_k\):

\[
P_{n_1,\ldots,n_k}(B)
=
P\{(X_{n_1},\ldots,X_{n_k})\in B\}.
\]

Через закон \(P_X\) это записывается так:

\[
P_{n_1,\ldots,n_k}(B)
=
P_X(\pi_{n_1,\ldots,n_k}^{-1}(B)).
\]

Иными словами, конечномерные распределения - это маргиналы меры \(P_X\).

\subsubsection{2.7. Связь с теоремой Колмогорова и канонический
процесс}\label{ux441ux432ux44fux437ux44c-ux441-ux442ux435ux43eux440ux435ux43cux43eux439-ux43aux43eux43bux43cux43eux433ux43eux440ux43eux432ux430-ux438-ux43aux430ux43dux43eux43dux438ux447ux435ux441ux43aux438ux439-ux43fux440ux43eux446ux435ux441ux441}

Наоборот, если заданы все конечномерные распределения и они согласованы,
то по Билету 15 существует мера \(P_X\) на \(E^{\mathbb N_0}\), для
которой эти распределения являются маргиналами. Тогда на каноническом
пространстве можно взять координатный процесс

\[
X_n(x)=x_n,
\qquad
x=(x_0,x_1,\ldots).
\]

Так получается случайная последовательность с заданными конечномерными
распределениями.

\subsubsection{2.8. Различие между состоянием, траекторией и
законом}\label{ux440ux430ux437ux43bux438ux447ux438ux435-ux43cux435ux436ux434ux443-ux441ux43eux441ux442ux43eux44fux43dux438ux435ux43c-ux442ux440ux430ux435ux43aux442ux43eux440ux438ux435ux439-ux438-ux437ux430ux43aux43eux43dux43eux43c}

Очень важно различать три объекта:

\begin{enumerate}
\def\labelenumi{\arabic{enumi}.}
\tightlist
\item
  \(X_n\) - состояние процесса в момент времени \(n\) как случайная
  величина;
\item
  \(X(\omega)\) - траектория, то есть последовательность значений при
  фиксированном исходе \(\omega\);
\item
  \(P_X\) - закон процесса, то есть вероятность на пространстве всех
  последовательностей.
\end{enumerate}

Траектория - не распределение. Распределение - мера на множестве
возможных траекторий.

\subsubsection{2.9. Определение закона через
цилиндры}\label{ux43eux43fux440ux435ux434ux435ux43bux435ux43dux438ux435-ux437ux430ux43aux43eux43dux430-ux447ux435ux440ux435ux437-ux446ux438ux43bux438ux43dux434ux440ux44b}

Для дискретного времени цилиндры особенно удобны. Например, если
\(E=\mathbb R\), то событие

\[
\{X_0\le a_0,\ X_1\le a_1,\ldots,X_n\le a_n\}
\]

соответствует цилиндру

\[
(-\infty,a_0]\times\cdots\times(-\infty,a_n]\times E\times E\times\cdots.
\]

Вероятность такого цилиндра равна конечномерной функции распределения:

\[
F_{0,\ldots,n}(a_0,\ldots,a_n)
=
P\{X_0\le a_0,\ldots,X_n\le a_n\}.
\]

Поэтому закон случайной последовательности полностью определяется
вероятностями цилиндров, а значит всеми конечномерными распределениями.

\subsection{3. Основные
определения}\label{ux43eux441ux43dux43eux432ux43dux44bux435-ux43eux43fux440ux435ux434ux435ux43bux435ux43dux438ux44f-15}

\subsubsection{Дискретный случайный
процесс}\label{ux434ux438ux441ux43aux440ux435ux442ux43dux44bux439-ux441ux43bux443ux447ux430ux439ux43dux44bux439-ux43fux440ux43eux446ux435ux441ux441}

Семейство случайных элементов

\[
\{X_n:n\in\mathbb N_0\}
\]

на одном вероятностном пространстве, где каждое

\[
X_n:\Omega\to E
\]

измеримо.

\subsubsection{Случайная
последовательность}\label{ux441ux43bux443ux447ux430ux439ux43dux430ux44f-ux43fux43eux441ux43bux435ux434ux43eux432ux430ux442ux435ux43bux44cux43dux43eux441ux442ux44c-1}

Случайный элемент со значениями в пространстве последовательностей:

\[
X:\Omega\to E^{\mathbb N_0},
\qquad
X(\omega)=(X_0(\omega),X_1(\omega),\ldots).
\]

\subsubsection{Пространство
последовательностей}\label{ux43fux440ux43eux441ux442ux440ux430ux43dux441ux442ux432ux43e-ux43fux43eux441ux43bux435ux434ux43eux432ux430ux442ux435ux43bux44cux43dux43eux441ux442ux435ux439}

\[
E^{\mathbb N_0}
=
\{(x_0,x_1,\ldots):x_n\in E\}.
\]

\subsubsection{\texorpdfstring{Произведенная
\(\sigma\)-алгебра}{Произведенная \textbackslash sigma-алгебра}}\label{ux43fux440ux43eux438ux437ux432ux435ux434ux435ux43dux43dux430ux44f-sigma-ux430ux43bux433ux435ux431ux440ux430-1}

\[
\mathcal E^{\mathbb N_0}
=
\sigma\{\pi_n^{-1}(A):n\ge0,\ A\in\mathcal E\}.
\]

\subsubsection{Координатная
проекция}\label{ux43aux43eux43eux440ux434ux438ux43dux430ux442ux43dux430ux44f-ux43fux440ux43eux435ux43aux446ux438ux44f}

\[
\pi_n(x)=x_n.
\]

\subsubsection{Цилиндрическое
событие}\label{ux446ux438ux43bux438ux43dux434ux440ux438ux447ux435ux441ux43aux43eux435-ux441ux43eux431ux44bux442ux438ux435}

\[
\pi_{n_1,\ldots,n_k}^{-1}(B),
\qquad
B\in\mathcal E^k.
\]

Это событие, зависящее от конечного числа координат.

\subsubsection{Траектория
процесса}\label{ux442ux440ux430ux435ux43aux442ux43eux440ux438ux44f-ux43fux440ux43eux446ux435ux441ux441ux430-1}

Для фиксированного исхода \(\omega\) последовательность

\[
X(\omega)=(X_0(\omega),X_1(\omega),\ldots).
\]

\subsubsection{Закон
процесса}\label{ux437ux430ux43aux43eux43d-ux43fux440ux43eux446ux435ux441ux441ux430}

\[
P_X=P\circ X^{-1}
\]

\begin{itemize}
\tightlist
\item
  распределение случайного элемента \(X\) на пространстве
  последовательностей.
\end{itemize}

\subsubsection{Конечномерное
распределение}\label{ux43aux43eux43dux435ux447ux43dux43eux43cux435ux440ux43dux43eux435-ux440ux430ux441ux43fux440ux435ux434ux435ux43bux435ux43dux438ux435}

\[
P_{n_1,\ldots,n_k}(B)
=
P\{(X_{n_1},\ldots,X_{n_k})\in B\}.
\]

\subsection{4. Основные теоремы и
свойства}\label{ux43eux441ux43dux43eux432ux43dux44bux435-ux442ux435ux43eux440ux435ux43cux44b-ux438-ux441ux432ux43eux439ux441ux442ux432ux430-15}

\subsubsection{Эквивалентность двух
описаний}\label{ux44dux43aux432ux438ux432ux430ux43bux435ux43dux442ux43dux43eux441ux442ux44c-ux434ux432ux443ux445-ux43eux43fux438ux441ux430ux43dux438ux439}

Последовательность случайных величин

\[
X_0,X_1,\ldots
\]

эквивалентна одному случайному отображению

\[
X:\Omega\to E^{\mathbb N_0},
\qquad
X(\omega)=(X_0(\omega),X_1(\omega),\ldots).
\]

\subsubsection{Критерий измеримости в пространство
последовательностей}\label{ux43aux440ux438ux442ux435ux440ux438ux439-ux438ux437ux43cux435ux440ux438ux43cux43eux441ux442ux438-ux432-ux43fux440ux43eux441ux442ux440ux430ux43dux441ux442ux432ux43e-ux43fux43eux441ux43bux435ux434ux43eux432ux430ux442ux435ux43bux44cux43dux43eux441ux442ux435ux439}

Отображение

\[
X:\Omega\to E^{\mathbb N_0}
\]

измеримо относительно произведенной \(\sigma\)-алгебры тогда и только
тогда, когда все координаты

\[
X_n=\pi_n\circ X
\]

измеримы.

\subsubsection{Закон процесса задается на пространстве
последовательностей}\label{ux437ux430ux43aux43eux43d-ux43fux440ux43eux446ux435ux441ux441ux430-ux437ux430ux434ux430ux435ux442ux441ux44f-ux43dux430-ux43fux440ux43eux441ux442ux440ux430ux43dux441ux442ux432ux435-ux43fux43eux441ux43bux435ux434ux43eux432ux430ux442ux435ux43bux44cux43dux43eux441ux442ux435ux439}

\[
P_X(A)=P\{X\in A\},
\qquad
A\in\mathcal E^{\mathbb N_0}.
\]

\subsubsection{Конечномерные распределения являются маргиналами
закона}\label{ux43aux43eux43dux435ux447ux43dux43eux43cux435ux440ux43dux44bux435-ux440ux430ux441ux43fux440ux435ux434ux435ux43bux435ux43dux438ux44f-ux44fux432ux43bux44fux44eux442ux441ux44f-ux43cux430ux440ux433ux438ux43dux430ux43bux430ux43cux438-ux437ux430ux43aux43eux43dux430}

\[
P_{n_1,\ldots,n_k}(B)
=
P_X(\pi_{n_1,\ldots,n_k}^{-1}(B)).
\]

\subsubsection{Закон процесса определяется
цилиндрами}\label{ux437ux430ux43aux43eux43d-ux43fux440ux43eux446ux435ux441ux441ux430-ux43eux43fux440ux435ux434ux435ux43bux44fux435ux442ux441ux44f-ux446ux438ux43bux438ux43dux434ux440ux430ux43cux438}

Если две вероятностные меры на \(E^{\mathbb N_0}\) совпадают на всех
цилиндрах, то они совпадают на всей произведенной \(\sigma\)-алгебре.

\subsubsection{Согласованные конечномерные распределения задают закон
последовательности}\label{ux441ux43eux433ux43bux430ux441ux43eux432ux430ux43dux43dux44bux435-ux43aux43eux43dux435ux447ux43dux43eux43cux435ux440ux43dux44bux435-ux440ux430ux441ux43fux440ux435ux434ux435ux43bux435ux43dux438ux44f-ux437ux430ux434ux430ux44eux442-ux437ux430ux43aux43eux43d-ux43fux43eux441ux43bux435ux434ux43eux432ux430ux442ux435ux43bux44cux43dux43eux441ux442ux438}

По теореме Колмогорова из Билета 15 согласованное семейство
конечномерных распределений задает меру на \(E^{\mathbb N_0}\), то есть
закон случайной последовательности.

\subsection{5. Примеры}\label{ux43fux440ux438ux43cux435ux440ux44b-15}

\subsubsection{Независимые броски
монеты}\label{ux43dux435ux437ux430ux432ux438ux441ux438ux43cux44bux435-ux431ux440ux43eux441ux43aux438-ux43cux43eux43dux435ux442ux44b}

Пусть

\[
E=\{0,1\},
\]

где \(1\) означает успех, а \(0\) - неуспех. Последовательность

\[
X_0,X_1,\ldots
\]

независимых бернуллиевских величин с параметром \(p\) имеет
конечномерные вероятности

\[
P\{X_0=x_0,\ldots,X_n=x_n\}
=
\prod_{k=0}^n p^{x_k}(1-p)^{1-x_k},
\qquad
x_k\in\{0,1\}.
\]

Ее закон - мера на пространстве

\[
\{0,1\}^{\mathbb N_0}.
\]

Траектория - бесконечная последовательность нулей и единиц.

\subsubsection{Случайное
блуждание}\label{ux441ux43bux443ux447ux430ux439ux43dux43eux435-ux431ux43bux443ux436ux434ux430ux43dux438ux435}

Пусть \(\xi_1,\xi_2,\ldots\) - независимые одинаково распределенные
приращения, например

\[
P\{\xi_k=1\}=P\{\xi_k=-1\}=\frac12.
\]

Определим

\[
S_0=0,
\qquad
S_n=\sum_{k=1}^n \xi_k.
\]

Тогда

\[
S=(S_0,S_1,S_2,\ldots)
\]

\begin{itemize}
\tightlist
\item
  случайная последовательность со значениями в
  \(\mathbb Z^{\mathbb N_0}\). Каждая траектория - путь блуждания.
\end{itemize}

\subsubsection{Марковская
цепь}\label{ux43cux430ux440ux43aux43eux432ux441ux43aux430ux44f-ux446ux435ux43fux44c}

Пусть \(E\) - счетное множество состояний, \(\pi\) - начальное
распределение, \(P=(p_{ij})\) - переходная матрица. Тогда конечномерные
вероятности имеют вид

\[
P\{X_0=i_0,\ldots,X_n=i_n\}
=
\pi_{i_0}p_{i_0i_1}\cdots p_{i_{n-1}i_n}.
\]

Это задает закон на \(E^{\mathbb N_0}\). Координатная последовательность
с этим законом является марковской цепью. Это связь с Билетом 22 и
Билетом 23.

\subsubsection{Стационарная
последовательность}\label{ux441ux442ux430ux446ux438ux43eux43dux430ux440ux43dux430ux44f-ux43fux43eux441ux43bux435ux434ux43eux432ux430ux442ux435ux43bux44cux43dux43eux441ux442ux44c}

Случайная последовательность называется стационарной в узком смысле,
если ее конечномерные распределения не меняются при сдвиге времени:

\[
(X_{n_1+h},\ldots,X_{n_k+h})
\stackrel d=
(X_{n_1},\ldots,X_{n_k})
\]

для всех допустимых \(h\). В терминах закона на \(E^{\mathbb N_0}\) это
означает инвариантность закона относительно сдвига последовательности.
Это связь с Билетом 17.

\subsubsection{Пример различия траектории и
закона}\label{ux43fux440ux438ux43cux435ux440-ux440ux430ux437ux43bux438ux447ux438ux44f-ux442ux440ux430ux435ux43aux442ux43eux440ux438ux438-ux438-ux437ux430ux43aux43eux43dux430}

Если бросать монету бесконечно много раз, то конкретная траектория может
быть

\[
(1,0,0,1,1,0,\ldots).
\]

Но закон процесса - не эта последовательность, а мера на множестве всех
бесконечных последовательностей из нулей и единиц. Эта мера говорит, с
какими вероятностями возникают цилиндры, например событие

\[
\{x_0=1,\ x_1=0,\ x_2=0\}.
\]

\subsection{6. Схемы доказательств /
обоснований}\label{ux441ux445ux435ux43cux44b-ux434ux43eux43aux430ux437ux430ux442ux435ux43bux44cux441ux442ux432-ux43eux431ux43eux441ux43dux43eux432ux430ux43dux438ux439-8}

\subsubsection{Из последовательности случайных величин получается
случайная
последовательность}\label{ux438ux437-ux43fux43eux441ux43bux435ux434ux43eux432ux430ux442ux435ux43bux44cux43dux43eux441ux442ux438-ux441ux43bux443ux447ux430ux439ux43dux44bux445-ux432ux435ux43bux438ux447ux438ux43d-ux43fux43eux43bux443ux447ux430ux435ux442ux441ux44f-ux441ux43bux443ux447ux430ux439ux43dux430ux44f-ux43fux43eux441ux43bux435ux434ux43eux432ux430ux442ux435ux43bux44cux43dux43eux441ux442ux44c}

Пусть все \(X_n\) измеримы. Определим

\[
X(\omega)=(X_0(\omega),X_1(\omega),\ldots).
\]

Чтобы доказать измеримость \(X\), достаточно проверить прообразы
образующих цилиндров:

\[
X^{-1}(\pi_n^{-1}(A))
=
\{\omega:X_n(\omega)\in A\}
\in\mathcal F.
\]

Значит \(X\) измеримо как отображение в \(E^{\mathbb N_0}\).

\subsubsection{Из случайной последовательности получаются
координаты}\label{ux438ux437-ux441ux43bux443ux447ux430ux439ux43dux43eux439-ux43fux43eux441ux43bux435ux434ux43eux432ux430ux442ux435ux43bux44cux43dux43eux441ux442ux438-ux43fux43eux43bux443ux447ux430ux44eux442ux441ux44f-ux43aux43eux43eux440ux434ux438ux43dux430ux442ux44b}

Пусть \(X:\Omega\to E^{\mathbb N_0}\) измеримо. Тогда

\[
X_n=\pi_n\circ X.
\]

Поскольку координатные проекции \(\pi_n\) измеримы, каждая \(X_n\)
измерима. Значит координаты случайной последовательности являются
случайными величинами.

\subsubsection{Почему закон определяется конечномерными
распределениями}\label{ux43fux43eux447ux435ux43cux443-ux437ux430ux43aux43eux43d-ux43eux43fux440ux435ux434ux435ux43bux44fux435ux442ux441ux44f-ux43aux43eux43dux435ux447ux43dux43eux43cux435ux440ux43dux44bux43cux438-ux440ux430ux441ux43fux440ux435ux434ux435ux43bux435ux43dux438ux44fux43cux438}

Конечномерные распределения задают вероятности всех цилиндров:

\[
P_X(\pi_{n_1,\ldots,n_k}^{-1}(B))
=
P\{(X_{n_1},\ldots,X_{n_k})\in B\}.
\]

Цилиндры порождают произведенную \(\sigma\)-алгебру. Поэтому мера на
пространстве последовательностей определяется своими значениями на
цилиндрах.

\subsubsection{Как используется теорема
Колмогорова}\label{ux43aux430ux43a-ux438ux441ux43fux43eux43bux44cux437ux443ux435ux442ux441ux44f-ux442ux435ux43eux440ux435ux43cux430-ux43aux43eux43bux43cux43eux433ux43eux440ux43eux432ux430}

Если процесс еще не задан, но есть согласованные конечномерные
распределения, то теорема Колмогорова строит меру \(P\) на
\(E^{\mathbb N_0}\). После этого можно взять канонический процесс

\[
X_n(x)=x_n.
\]

Его закон равен \(P\), а конечномерные распределения совпадают с
заданными.

\subsection{7. Что написать на
доске}\label{ux447ux442ux43e-ux43dux430ux43fux438ux441ux430ux442ux44c-ux43dux430-ux434ux43eux441ux43aux435-15}

\begin{enumerate}
\def\labelenumi{\arabic{enumi}.}
\tightlist
\item
  Две точки зрения:
\end{enumerate}

\[
\{X_n\}_{n\ge0}
\quad\Longleftrightarrow\quad
X:\Omega\to E^{\mathbb N_0}.
\]

\begin{enumerate}
\def\labelenumi{\arabic{enumi}.}
\setcounter{enumi}{1}
\tightlist
\item
  Случайная последовательность:
\end{enumerate}

\[
X(\omega)=(X_0(\omega),X_1(\omega),\ldots).
\]

\begin{enumerate}
\def\labelenumi{\arabic{enumi}.}
\setcounter{enumi}{2}
\tightlist
\item
  Пространство последовательностей:
\end{enumerate}

\[
E^{\mathbb N_0}.
\]

\begin{enumerate}
\def\labelenumi{\arabic{enumi}.}
\setcounter{enumi}{3}
\tightlist
\item
  Координатные проекции:
\end{enumerate}

\[
\pi_n(x)=x_n,
\qquad
X_n=\pi_n\circ X.
\]

\begin{enumerate}
\def\labelenumi{\arabic{enumi}.}
\setcounter{enumi}{4}
\tightlist
\item
  Произведенная \(\sigma\)-алгебра:
\end{enumerate}

\[
\mathcal E^{\mathbb N_0}
=
\sigma\{\pi_n^{-1}(A):A\in\mathcal E,\ n\ge0\}.
\]

\begin{enumerate}
\def\labelenumi{\arabic{enumi}.}
\setcounter{enumi}{5}
\tightlist
\item
  Закон процесса:
\end{enumerate}

\[
P_X=P\circ X^{-1}.
\]

\begin{enumerate}
\def\labelenumi{\arabic{enumi}.}
\setcounter{enumi}{6}
\tightlist
\item
  Конечномерные распределения как маргиналы:
\end{enumerate}

\[
P_{n_1,\ldots,n_k}(B)
=
P_X(\pi_{n_1,\ldots,n_k}^{-1}(B)).
\]

\subsection{8. Типичные дополнительные
вопросы}\label{ux442ux438ux43fux438ux447ux43dux44bux435-ux434ux43eux43fux43eux43bux43dux438ux442ux435ux43bux44cux43dux44bux435-ux432ux43eux43fux440ux43eux441ux44b-15}

\textbf{Что является главным объектом билета?}

Дискретный случайный процесс, который можно рассматривать либо как
последовательность случайных величин \(X_0,X_1,\ldots\), либо как один
случайный элемент \(X\) со значениями в пространстве последовательностей
\(E^{\mathbb N_0}\).

\textbf{В чем разница между \(X_n\) и \(X(\omega)\)?}

\(X_n\) - случайная величина, отвечающая моменту времени \(n\). А
\(X(\omega)\) - вся траектория при фиксированном исходе \(\omega\):

\[
X(\omega)=(X_0(\omega),X_1(\omega),\ldots).
\]

\textbf{Что такое пространство последовательностей?}

Это множество всех возможных траекторий:

\[
E^{\mathbb N_0}=\{(x_0,x_1,\ldots):x_n\in E\}.
\]

На нем используется произведенная \(\sigma\)-алгебра, порожденная
координатными цилиндрами.

\textbf{Что такое цилиндрическое событие?}

Это событие, которое задает условия только на конечное число координат:

\[
\{x:(x_{n_1},\ldots,x_{n_k})\in B\}.
\]

Например, \(\{x_0\in A_0,\ x_3\in A_3\}\).

\textbf{Почему отображение \(X:\Omega\to E^{\mathbb N_0}\) измеримо,
если все координаты \(X_n\) измеримы?}

Потому что произведенная \(\sigma\)-алгебра порождена цилиндрами, а
прообраз координатного цилиндра равен

\[
X^{-1}(\pi_n^{-1}(A))=X_n^{-1}(A)\in\mathcal F.
\]

\textbf{Что такое закон процесса?}

Это распределение случайной последовательности:

\[
P_X(A)=P\{X\in A\},
\qquad
A\in\mathcal E^{\mathbb N_0}.
\]

Это мера на множестве всех возможных траекторий.

\textbf{Как конечномерные распределения связаны с законом процесса?}

Они являются маргиналами закона:

\[
P_{n_1,\ldots,n_k}(B)
=
P_X(\pi_{n_1,\ldots,n_k}^{-1}(B)).
\]

\textbf{Почему нельзя путать траекторию и распределение?}

Траектория - это один элемент \(E^{\mathbb N_0}\), то есть конкретная
последовательность значений. Распределение - это вероятностная мера на
множестве всех таких последовательностей.

\textbf{Как билет связан с Билетом 15?}

В Билете 15 доказывается, что согласованные конечномерные распределения
задают меру на пространстве траекторий. В этом билете для дискретного
времени это пространство имеет вид \(E^{\mathbb N_0}\).

\textbf{Как билет связан с марковскими цепями?}

Марковская цепь - это частный случай случайной последовательности. Ее
закон на \(E^{\mathbb N_0}\) задается начальным распределением и
переходными вероятностями:

\[
P\{X_0=i_0,\ldots,X_n=i_n\}
=
\pi_{i_0}p_{i_0i_1}\cdots p_{i_{n-1}i_n}.
\]

\subsection{9. Типичные
ошибки}\label{ux442ux438ux43fux438ux447ux43dux44bux435-ux43eux448ux438ux431ux43aux438-15}

\begin{itemize}
\tightlist
\item
  Путать время \(n\) и исход \(\omega\).
\item
  Путать случайную величину \(X_n\) и значение \(X_n(\omega)\).
\item
  Путать траекторию \(X(\omega)\) и закон процесса \(P_X\).
\item
  Забывать, что все \(X_n\) должны быть заданы на одном вероятностном
  пространстве.
\item
  Забывать произведенную \(\sigma\)-алгебру на \(E^{\mathbb N_0}\).
\item
  Считать, что одномерные распределения \(P_{X_n}\) задают закон
  процесса. Нужны все конечномерные распределения или мера на
  пространстве последовательностей.
\item
  Не различать семейство случайных величин и канонический координатный
  процесс на \(E^{\mathbb N_0}\).
\end{itemize}

\subsection{10. Связи с другими
билетами}\label{ux441ux432ux44fux437ux438-ux441-ux434ux440ux443ux433ux438ux43cux438-ux431ux438ux43bux435ux442ux430ux43cux438-15}

\begin{itemize}
\tightlist
\item
  Билет 05: случайные элементы и измеримые отображения.
\item
  Билет 06: распределение случайного элемента и конечномерные
  распределения.
\item
  Билет 15: согласованные конечномерные распределения задают меру на
  пространстве траекторий.
\item
  Билет 17: стационарность и эргодичность формулируются для случайных
  последовательностей.
\item
  Билет 19: случайные блуждания - примеры случайных последовательностей
  с независимыми приращениями.
\item
  Билет 22 и Билет 23: марковские цепи - случайные последовательности со
  специальной зависимостью.
\item
  Билет 24 и Билет 25: корреляционная функция и спектральные
  характеристики относятся к последовательностям второго порядка.
\end{itemize}

\subsection{11. Минимум на
удовлетворительно}\label{ux43cux438ux43dux438ux43cux443ux43c-ux43dux430-ux443ux434ux43eux432ux43bux435ux442ux432ux43eux440ux438ux442ux435ux43bux44cux43dux43e-15}

Нужно сказать:

\begin{enumerate}
\def\labelenumi{\arabic{enumi}.}
\tightlist
\item
  Дискретный процесс - это последовательность случайных величин:
\end{enumerate}

\[
X_0,X_1,\ldots
\]

\begin{enumerate}
\def\labelenumi{\arabic{enumi}.}
\setcounter{enumi}{1}
\tightlist
\item
  Его можно рассматривать как одну случайную последовательность:
\end{enumerate}

\[
X:\Omega\to E^{\mathbb N_0},
\qquad
X(\omega)=(X_0(\omega),X_1(\omega),\ldots).
\]

\begin{enumerate}
\def\labelenumi{\arabic{enumi}.}
\setcounter{enumi}{2}
\tightlist
\item
  Закон процесса:
\end{enumerate}

\[
P_X=P\circ X^{-1}
\]

на пространстве \(E^{\mathbb N_0}\).

\begin{enumerate}
\def\labelenumi{\arabic{enumi}.}
\setcounter{enumi}{3}
\item
  Конечномерные распределения - маргиналы этого закона.
\item
  Главная ловушка: траектория - не распределение.
\end{enumerate}

\subsection{12. Минимум на
хорошо}\label{ux43cux438ux43dux438ux43cux443ux43c-ux43dux430-ux445ux43eux440ux43eux448ux43e-15}

Помимо минимума, нужно объяснить:

\begin{itemize}
\tightlist
\item
  что такое произведенная \(\sigma\)-алгебра;
\item
  что такое цилиндрическое событие;
\item
  почему измеримость отображения в \(E^{\mathbb N_0}\) проверяется по
  координатам;
\item
  как из закона процесса получить конечномерные распределения;
\item
  как по согласованным конечномерным распределениям построить закон
  последовательности через теорему Колмогорова;
\item
  привести пример: iid-последовательность, случайное блуждание или
  марковская цепь.
\end{itemize}

\subsection{13. Что добавить для
отлично}\label{ux447ux442ux43e-ux434ux43eux431ux430ux432ux438ux442ux44c-ux434ux43bux44f-ux43eux442ux43bux438ux447ux43dux43e-15}

Для сильного ответа стоит добавить:

\begin{itemize}
\tightlist
\item
  четкое различие между \(X_n\), \(X_n(\omega)\), \(X(\omega)\) и
  \(P_X\);
\item
  формулу для цилиндров и конечномерных проекций;
\item
  доказательство эквивалентности двух точек зрения через измеримость
  координат;
\item
  объяснение, почему закон на \(E^{\mathbb N_0}\) определяется
  цилиндрами;
\item
  пример марковской цепи через \(\pi\) и \(P\);
\item
  связь со стационарностью как инвариантностью закона относительно
  сдвига.
\end{itemize}

\subsection{14. Устная версия ответа на 5
минут}\label{ux443ux441ux442ux43dux430ux44f-ux432ux435ux440ux441ux438ux44f-ux43eux442ux432ux435ux442ux430-ux43dux430-5-ux43cux438ux43dux443ux442-15}

\begin{enumerate}
\def\labelenumi{\arabic{enumi}.}
\tightlist
\item
  Дискретный случайный процесс - это семейство случайных величин
  \(X_0,X_1,\ldots\) на одном вероятностном пространстве.
\item
  Его можно рассматривать как одно отображение
\end{enumerate}

\[
X:\Omega\to E^{\mathbb N_0},
\qquad
X(\omega)=(X_0(\omega),X_1(\omega),\ldots).
\]

\begin{enumerate}
\def\labelenumi{\arabic{enumi}.}
\setcounter{enumi}{2}
\tightlist
\item
  Пространство \(E^{\mathbb N_0}\) - это пространство всех траекторий.
\item
  На нем берется произведенная \(\sigma\)-алгебра, порожденная
  координатными проекциями
\end{enumerate}

\[
\pi_n(x)=x_n.
\]

\begin{enumerate}
\def\labelenumi{\arabic{enumi}.}
\setcounter{enumi}{4}
\tightlist
\item
  Цилиндры - события, зависящие от конечного числа координат:
\end{enumerate}

\[
\{x:(x_{n_1},\ldots,x_{n_k})\in B\}.
\]

\begin{enumerate}
\def\labelenumi{\arabic{enumi}.}
\setcounter{enumi}{5}
\tightlist
\item
  Если все \(X_n\) измеримы, то отображение \(X=(X_0,X_1,\ldots)\)
  измеримо, потому что прообразы цилиндров измеримы.
\item
  Обратно, если \(X\) - случайный элемент пространства
  последовательностей, то его координаты
\end{enumerate}

\[
X_n=\pi_n\circ X
\]

являются случайными величинами. 8. Закон процесса - это распределение
\(X\):

\[
P_X=P\circ X^{-1}.
\]

\begin{enumerate}
\def\labelenumi{\arabic{enumi}.}
\setcounter{enumi}{8}
\tightlist
\item
  Конечномерные распределения - маргиналы этого закона:
\end{enumerate}

\[
P_{n_1,\ldots,n_k}(B)
=
P_X(\pi_{n_1,\ldots,n_k}^{-1}(B)).
\]

\begin{enumerate}
\def\labelenumi{\arabic{enumi}.}
\setcounter{enumi}{9}
\tightlist
\item
  По теореме Колмогорова согласованные конечномерные распределения
  задают такой закон на \(E^{\mathbb N_0}\).
\item
  Примеры: независимые броски монеты, случайное блуждание, марковская
  цепь.
\item
  Главная ошибка - путать конкретную траекторию с распределением на
  пространстве траекторий.
\end{enumerate}

\subsection{15. Если осталось 2 минуты перед
экзаменом}\label{ux435ux441ux43bux438-ux43eux441ux442ux430ux43bux43eux441ux44c-2-ux43cux438ux43dux443ux442ux44b-ux43fux435ux440ux435ux434-ux44dux43aux437ux430ux43cux435ux43dux43eux43c-15}

Запомнить главную эквивалентность:

\[
\{X_n\}_{n\ge0}
\quad\Longleftrightarrow\quad
X:\Omega\to E^{\mathbb N_0},
\qquad
X(\omega)=(X_0(\omega),X_1(\omega),\ldots).
\]

Координаты:

\[
X_n=\pi_n\circ X,
\qquad
\pi_n(x)=x_n.
\]

Закон процесса:

\[
P_X=P\circ X^{-1}.
\]

Конечномерные распределения:

\[
P_{n_1,\ldots,n_k}(B)
=
P_X(\pi_{n_1,\ldots,n_k}^{-1}(B)).
\]

Главная ловушка: \(X(\omega)\) - это одна траектория, а \(P_X\) - мера
на множестве всех траекторий.

\newpage

\section{Билет 17. Стационарные и эргодические случайные
последовательности в широком и узком
смыслах}\label{ux431ux438ux43bux435ux442-17.-ux441ux442ux430ux446ux438ux43eux43dux430ux440ux43dux44bux435-ux438-ux44dux440ux433ux43eux434ux438ux447ux435ux441ux43aux438ux435-ux441ux43bux443ux447ux430ux439ux43dux44bux435-ux43fux43eux441ux43bux435ux434ux43eux432ux430ux442ux435ux43bux44cux43dux43eux441ux442ux438-ux432-ux448ux438ux440ux43eux43aux43eux43c-ux438-ux443ux437ux43aux43eux43c-ux441ux43cux44bux441ux43bux430ux445}

Источники: Ширяев 2, гл. V-VI; лекции ДСП.

\subsection{1. Короткий
ответ}\label{ux43aux43eux440ux43eux442ux43aux438ux439-ux43eux442ux432ux435ux442-16}

Случайная последовательность

\[
\{X_n:n\in\mathbb Z\}
\]

называется стационарной в узком смысле, если ее конечномерные
распределения не меняются при сдвиге времени:

\[
(X_{n_1+h},\ldots,X_{n_k+h})
\stackrel d=
(X_{n_1},\ldots,X_{n_k})
\]

для всех \(k\), \(n_1,\ldots,n_k\) и целых \(h\).

Стационарность в широком смысле требует только существования вторых
моментов, постоянного среднего и ковариации, зависящей только от лага:

\[
\mathbb E X_n=m,
\qquad
\operatorname{Cov}(X_n,X_{n+k})=R(k).
\]

Эргодичность в узком смысле означает, что сдвиг-инвариантные события
имеют вероятность только \(0\) или \(1\). Интуитивно: одна достаточно
длинная траектория отражает вероятностные свойства процесса.

Эргодичность в широком смысле обычно понимают как сходимость временного
среднего к математическому ожиданию в среднем квадратичном:

\[
\frac1N\sum_{n=0}^{N-1}X_n
\xrightarrow[L^2]{N\to\infty}
m.
\]

\subsection{2. Полный
ответ}\label{ux43fux43eux43bux43dux44bux439-ux43eux442ux432ux435ux442-16}

\subsubsection{17.1. Введение и основная
идея}\label{ux432ux432ux435ux434ux435ux43dux438ux435-ux438-ux43eux441ux43dux43eux432ux43dux430ux44f-ux438ux434ux435ux44f}

Этот билет продолжает Билет 16, где дискретный процесс рассматривался
как случайная последовательность и как случайный элемент пространства
последовательностей. Теперь нас интересуют свойства последовательности,
которые описывают поведение при сдвиге времени.

Пусть

\[
X=(X_n)_{n\in\mathbb Z}
\]

или

\[
X=(X_n)_{n\ge0}
\]

\begin{itemize}
\tightlist
\item
  случайная последовательность. Для стационарности удобнее работать с
  двусторонним временем \(n\in\mathbb Z\), потому что сдвиг можно
  применять в обе стороны. Для экзамена можно формулировать и для
  \(n\ge0\), но тогда сдвиги берут только такие, чтобы индексы
  оставались неотрицательными.
\end{itemize}

Главная идея стационарности: статистические свойства процесса не
меняются при переносе начала отсчета времени. Если сегодня процесс
выглядит так же, как завтра, в смысле распределений, то он стационарен.

\subsubsection{17.2. Стационарность в узком
смысле}\label{ux441ux442ux430ux446ux438ux43eux43dux430ux440ux43dux43eux441ux442ux44c-ux432-ux443ux437ux43aux43eux43c-ux441ux43cux44bux441ux43bux435}

Сначала определим стационарность в узком смысле. Последовательность
называется стационарной в узком смысле, или строго стационарной, если
для любых

\[
k\ge1,\qquad n_1,\ldots,n_k\in\mathbb Z,\qquad h\in\mathbb Z
\]

выполнено равенство распределений:

\[
(X_{n_1+h},\ldots,X_{n_k+h})
\stackrel d=
(X_{n_1},\ldots,X_{n_k}).
\]

Это означает, что все конечномерные распределения инвариантны
относительно сдвига времени. Например,

\[
X_0\stackrel d=X_{10},
\]

но этого мало: также должно быть

\[
(X_0,X_1,X_5)\stackrel d=(X_{10},X_{11},X_{15})
\]

и аналогично для любого конечного набора моментов.

На языке пространства последовательностей из Билета 16 стационарность в
узком смысле можно выразить через сдвиг. Пусть

\[
S:E^{\mathbb Z}\to E^{\mathbb Z}
\]

задан формулой

\[
(Sx)_n=x_{n+1}.
\]

Если \(P_X\) - закон последовательности на \(E^{\mathbb Z}\), то узкая
стационарность означает

\[
P_X\circ S^{-1}=P_X.
\]

То есть закон процесса как мера на пространстве траекторий инвариантен
относительно сдвига.

\subsubsection{17.3. Стационарность в широком
смысле}\label{ux441ux442ux430ux446ux438ux43eux43dux430ux440ux43dux43eux441ux442ux44c-ux432-ux448ux438ux440ux43eux43aux43eux43c-ux441ux43cux44bux441ux43bux435}

Теперь стационарность в широком смысле. Она относится только к
характеристикам второго порядка, поэтому требует

\[
\mathbb E X_n^2<\infty
\]

для всех \(n\). Последовательность стационарна в широком смысле, если:

\begin{enumerate}
\def\labelenumi{\arabic{enumi}.}
\tightlist
\item
  среднее не зависит от времени:
\end{enumerate}

\[
\mathbb E X_n=m;
\]

\begin{enumerate}
\def\labelenumi{\arabic{enumi}.}
\setcounter{enumi}{1}
\tightlist
\item
  ковариация зависит только от разности индексов:
\end{enumerate}

\[
\operatorname{Cov}(X_n,X_m)=R(m-n).
\]

Часто записывают

\[
R(k)=\operatorname{Cov}(X_n,X_{n+k}).
\]

Если последовательность центрирована, то есть \(m=0\), то

\[
R(k)=\mathbb E[X_nX_{n+k}].
\]

Иногда в курсах случайных процессов корреляционной функцией называют
именно

\[
\mathbb E[X_nX_{n+k}],
\]

а ковариационной - центрированную величину. Важно проговорить, какая
версия используется. В этих билетах удобнее держать в голове:

\[
\operatorname{Cov}(X_n,X_{n+k})
=
\mathbb E[(X_n-m)(X_{n+k}-m)].
\]

\subsubsection{17.4. Связь между типами
стационарности}\label{ux441ux432ux44fux437ux44c-ux43cux435ux436ux434ux443-ux442ux438ux43fux430ux43cux438-ux441ux442ux430ux446ux438ux43eux43dux430ux440ux43dux43eux441ux442ux438}

Узкая стационарность при существовании вторых моментов влечет широкую
стационарность. Действительно, из равенства одномерных распределений
следует одинаковое среднее:

\[
\mathbb E X_n=m.
\]

Из равенства двумерных распределений

\[
(X_n,X_{n+k})\stackrel d=(X_0,X_k)
\]

следует

\[
\operatorname{Cov}(X_n,X_{n+k})
=
\operatorname{Cov}(X_0,X_k).
\]

Обратное в общем случае неверно: широкая стационарность контролирует
только первые два момента, но не все распределение.

Есть важное исключение: для гауссовских последовательностей широкая
стационарность влечет узкую. Причина в том, что конечномерные
распределения гауссовского процесса полностью задаются средним вектором
и ковариационной матрицей. Если среднее постоянно, а ковариация зависит
только от лага, то все конечномерные гауссовские распределения
инвариантны при сдвиге. Это связь с Билетом 18.

\subsubsection{17.5. Понятие эргодичности и инвариантные
события}\label{ux43fux43eux43dux44fux442ux438ux435-ux44dux440ux433ux43eux434ux438ux447ux43dux43eux441ux442ux438-ux438-ux438ux43dux432ux430ux440ux438ux430ux43dux442ux43dux44bux435-ux441ux43eux431ux44bux442ux438ux44f}

Теперь эргодичность. Стационарность говорит, что закон не меняется при
сдвиге. Эргодичность говорит сильнее: в стационарном процессе нет
скрытой случайной характеристики, которая сохраняется при сдвиге и
разбивает траектории на разные режимы.

На пространстве траекторий событие \(A\) называется сдвиг-инвариантным,
если

\[
S^{-1}A=A
\]

с точностью до множеств вероятности ноль. Интуитивно такое событие не
меняется, если сдвинуть всю последовательность на один шаг.

Стационарная последовательность называется эргодической в узком смысле,
если любое сдвиг-инвариантное событие имеет вероятность \(0\) или \(1\):

\[
A=S^{-1}A
\quad\Longrightarrow\quad
P_X(A)\in\{0,1\}.
\]

Это определение часто называют метрической эргодичностью. Его смысл:
сдвиг не оставляет нетривиальных случайных ``меток'' траектории.

\subsubsection{17.6. Практическая интерпретация и временные
средние}\label{ux43fux440ux430ux43aux442ux438ux447ux435ux441ux43aux430ux44f-ux438ux43dux442ux435ux440ux43fux440ux435ux442ux430ux446ux438ux44f-ux438-ux432ux440ux435ux43cux435ux43dux43dux44bux435-ux441ux440ux435ux434ux43dux438ux435}

Практическая интерпретация эргодичности связана с временными средними.
Если процесс стационарен и эргодичен в узком смысле, то по эргодической
теореме для любой измеримой функции \(f\) с
\(\mathbb E|f(X_0)|<\infty\):

\[
\frac1N\sum_{n=0}^{N-1} f(X_n)
\to
\mathbb E f(X_0)
\]

почти наверное. В частности, для \(f(x)=x\), если
\(\mathbb E|X_0|<\infty\):

\[
\frac1N\sum_{n=0}^{N-1}X_n
\to
\mathbb E X_0
\]

почти наверное.

Это и есть главный прикладной смысл: среднее по одной длинной реализации
можно использовать вместо среднего по ансамблю.

\subsubsection{17.7. Эргодичность в широком
смысле}\label{ux44dux440ux433ux43eux434ux438ux447ux43dux43eux441ux442ux44c-ux432-ux448ux438ux440ux43eux43aux43eux43c-ux441ux43cux44bux441ux43bux435}

Эргодичность в широком смысле - более слабое понятие второго порядка.
Для стационарной в широком смысле последовательности с конечными вторыми
моментами говорят, что она эргодична по среднему, или эргодична в
широком смысле относительно среднего, если

\[
\overline X_N
=
\frac1N\sum_{n=0}^{N-1}X_n
\]

сходится к \(m=\mathbb EX_0\) в среднем квадратичном:

\[
\mathbb E|\overline X_N-m|^2\to0.
\]

Эквивалентно,

\[
\operatorname{Var}(\overline X_N)\to0.
\]

Для стационарной в широком смысле последовательности

\[
\operatorname{Var}(\overline X_N)
=
\frac1{N^2}
\sum_{i=0}^{N-1}\sum_{j=0}^{N-1}R(j-i),
\]

где

\[
R(k)=\operatorname{Cov}(X_0,X_k).
\]

Поэтому условия затухания корреляций дают эргодичность среднего. Это
подробно развивается в Билете 24.

\subsubsection{17.8. Сравнение понятий и
примеры}\label{ux441ux440ux430ux432ux43dux435ux43dux438ux435-ux43fux43eux43dux44fux442ux438ux439-ux438-ux43fux440ux438ux43cux435ux440ux44b}

Важно различать:

\begin{enumerate}
\def\labelenumi{\arabic{enumi}.}
\tightlist
\item
  стационарность - инвариантность распределений или моментов при сдвиге;
\item
  эргодичность - возможность восстанавливать вероятностные
  характеристики по одной длинной реализации;
\item
  независимость - совсем другое свойство, гораздо сильнее во многих
  примерах.
\end{enumerate}

Например, iid-последовательность стационарна и эргодична, но
стационарность сама по себе не означает независимость.
Последовательность \(X_n=Y\), где \(Y\) - одна случайная величина,
стационарна, но не эргодична: траектория хранит неизменный случайный
уровень \(Y\).

\subsection{3. Основные
определения}\label{ux43eux441ux43dux43eux432ux43dux44bux435-ux43eux43fux440ux435ux434ux435ux43bux435ux43dux438ux44f-16}

\subsubsection{Сдвиг на пространстве
последовательностей}\label{ux441ux434ux432ux438ux433-ux43dux430-ux43fux440ux43eux441ux442ux440ux430ux43dux441ux442ux432ux435-ux43fux43eux441ux43bux435ux434ux43eux432ux430ux442ux435ux43bux44cux43dux43eux441ux442ux435ux439}

Для двусторонних последовательностей \(x=(x_n)_{n\in\mathbb Z}\):

\[
(Sx)_n=x_{n+1}.
\]

\subsubsection{Стационарность в узком
смысле}\label{ux441ux442ux430ux446ux438ux43eux43dux430ux440ux43dux43eux441ux442ux44c-ux432-ux443ux437ux43aux43eux43c-ux441ux43cux44bux441ux43bux435-1}

Последовательность \((X_n)\) стационарна в узком смысле, если для любых
\(k\), \(n_1,\ldots,n_k\) и \(h\):

\[
(X_{n_1+h},\ldots,X_{n_k+h})
\stackrel d=
(X_{n_1},\ldots,X_{n_k}).
\]

\subsubsection{Стационарность в широком
смысле}\label{ux441ux442ux430ux446ux438ux43eux43dux430ux440ux43dux43eux441ux442ux44c-ux432-ux448ux438ux440ux43eux43aux43eux43c-ux441ux43cux44bux441ux43bux435-1}

Последовательность с конечными вторыми моментами стационарна в широком
смысле, если

\[
\mathbb E X_n=m
\]

не зависит от \(n\), а

\[
\operatorname{Cov}(X_n,X_m)=R(m-n)
\]

зависит только от лага.

\subsubsection{Корреляционная или ковариационная функция стационарной
последовательности}\label{ux43aux43eux440ux440ux435ux43bux44fux446ux438ux43eux43dux43dux430ux44f-ux438ux43bux438-ux43aux43eux432ux430ux440ux438ux430ux446ux438ux43eux43dux43dux430ux44f-ux444ux443ux43dux43aux446ux438ux44f-ux441ux442ux430ux446ux438ux43eux43dux430ux440ux43dux43eux439-ux43fux43eux441ux43bux435ux434ux43eux432ux430ux442ux435ux43bux44cux43dux43eux441ux442ux438}

Для центрированной последовательности:

\[
R(k)=\mathbb E[X_0X_k].
\]

В общем случае:

\[
R(k)=\operatorname{Cov}(X_0,X_k).
\]

\subsubsection{Сдвиг-инвариантное
событие}\label{ux441ux434ux432ux438ux433-ux438ux43dux432ux430ux440ux438ux430ux43dux442ux43dux43eux435-ux441ux43eux431ux44bux442ux438ux435}

Событие \(A\) на пространстве траекторий такое, что

\[
S^{-1}A=A
\]

с точностью до множеств вероятности ноль.

\subsubsection{Эргодичность в узком
смысле}\label{ux44dux440ux433ux43eux434ux438ux447ux43dux43eux441ux442ux44c-ux432-ux443ux437ux43aux43eux43c-ux441ux43cux44bux441ux43bux435}

Стационарная последовательность эргодична, если каждое
сдвиг-инвариантное событие имеет вероятность \(0\) или \(1\):

\[
S^{-1}A=A
\quad\Rightarrow\quad
P(A)\in\{0,1\}.
\]

\subsubsection{Эргодичность в широком смысле по
среднему}\label{ux44dux440ux433ux43eux434ux438ux447ux43dux43eux441ux442ux44c-ux432-ux448ux438ux440ux43eux43aux43eux43c-ux441ux43cux44bux441ux43bux435-ux43fux43e-ux441ux440ux435ux434ux43dux435ux43cux443}

Стационарная в широком смысле последовательность эргодична по среднему,
если

\[
\frac1N\sum_{n=0}^{N-1}X_n
\xrightarrow[L^2]{}
\mathbb E X_0.
\]

\subsection{4. Основные теоремы и
свойства}\label{ux43eux441ux43dux43eux432ux43dux44bux435-ux442ux435ux43eux440ux435ux43cux44b-ux438-ux441ux432ux43eux439ux441ux442ux432ux430-16}

\begin{itemize}
\tightlist
\item
  Узкая стационарность означает инвариантность закона процесса
  относительно сдвига:
\end{itemize}

\[
P_X\circ S^{-1}=P_X.
\]

\begin{itemize}
\tightlist
\item
  Узкая стационарность и конечные вторые моменты влекут широкую
  стационарность.
\item
  Широкая стационарность в общем случае не влечет узкую.
\item
  Для гауссовской последовательности широкая стационарность влечет
  узкую.
\item
  Эргодичность в узком смысле равносильна отсутствию нетривиальных
  сдвиг-инвариантных событий.
\item
  Для стационарной эргодической последовательности временные средние
  интегрируемых функций совпадают с математическими ожиданиями:
\end{itemize}

\[
\frac1N\sum_{n=0}^{N-1}f(X_n)\to\mathbb E f(X_0)
\]

почти наверное.

\begin{itemize}
\tightlist
\item
  Для стационарной в широком смысле последовательности эргодичность
  среднего связана с условием:
\end{itemize}

\[
\operatorname{Var}\left(\frac1N\sum_{n=0}^{N-1}X_n\right)\to0.
\]

\begin{itemize}
\tightlist
\item
  Достаточно, например, чтобы
\end{itemize}

\[
\frac1N\sum_{k=1}^{N}|R(k)|\to0.
\]

В частности, если \(R(k)\to0\), то временное среднее сходится к
математическому ожиданию в \(L^2\). Общий критерий второго порядка:

\[
\operatorname{Var}(\overline X_N)=
\frac1{N^2}\sum_{i,j=0}^{N-1}R(j-i)\to0.
\]

\subsection{5. Примеры}\label{ux43fux440ux438ux43cux435ux440ux44b-16}

\subsubsection{iid-последовательность}\label{iid-ux43fux43eux441ux43bux435ux434ux43eux432ux430ux442ux435ux43bux44cux43dux43eux441ux442ux44c}

Пусть

\[
X_0,X_1,\ldots
\]

независимы и одинаково распределены.

Тогда последовательность стационарна в узком смысле: любой конечный
вектор имеет произведение одних и тех же одномерных законов, и сдвиг
времени не меняет это произведение.

Если \(\mathbb E X_0^2<\infty\), то она стационарна и в широком смысле:

\[
\mathbb E X_n=m,
\]

и

\[
\operatorname{Cov}(X_n,X_{n+k})=
\begin{cases}
\operatorname{Var}X_0, & k=0,\\
0, & k\ne0.
\end{cases}
\]

Такая последовательность эргодична: временные средние сходятся к
математическому ожиданию по закону больших чисел.

\subsubsection{Случайная
константа}\label{ux441ux43bux443ux447ux430ux439ux43dux430ux44f-ux43aux43eux43dux441ux442ux430ux43dux442ux430}

Пусть

\[
X_n=Y
\]

для всех \(n\), где \(Y\) - неконстантная случайная величина.

Эта последовательность стационарна в узком смысле: сдвиг времени вообще
ничего не меняет. Если \(\mathbb E Y^2<\infty\), она стационарна и в
широком смысле:

\[
R(k)=\operatorname{Var}Y
\]

для всех \(k\).

Но она не эргодична. Временное среднее равно

\[
\frac1N\sum_{n=0}^{N-1}X_n=Y,
\]

а не константе \(\mathbb EY\). Одна траектория показывает только свой
случайный уровень.

\subsubsection{Широкая стационарность без
узкой}\label{ux448ux438ux440ux43eux43aux430ux44f-ux441ux442ux430ux446ux438ux43eux43dux430ux440ux43dux43eux441ux442ux44c-ux431ux435ux437-ux443ux437ux43aux43eux439}

Пусть величины \(X_n\) независимы, имеют нулевое среднее и единичную
дисперсию, но распределения чередуются:

\[
X_{2k}\sim N(0,1),
\qquad
X_{2k+1}\sim
\begin{cases}
1, & 1/2,\\
-1, & 1/2.
\end{cases}
\]

Тогда

\[
\mathbb E X_n=0,
\qquad
\operatorname{Var}X_n=1,
\]

а при \(n\ne m\):

\[
\operatorname{Cov}(X_n,X_m)=0
\]

из-за независимости. Значит последовательность стационарна в широком
смысле:

\[
R(0)=1,\qquad R(k)=0,\ k\ne0.
\]

Но она не стационарна в узком смысле, потому что одномерное
распределение зависит от четности \(n\).

\subsubsection{Гауссовская стационарная
последовательность}\label{ux433ux430ux443ux441ux441ux43eux432ux441ux43aux430ux44f-ux441ux442ux430ux446ux438ux43eux43dux430ux440ux43dux430ux44f-ux43fux43eux441ux43bux435ux434ux43eux432ux430ux442ux435ux43bux44cux43dux43eux441ux442ux44c}

Если последовательность гауссовская, то ее конечномерные распределения
полностью задаются средними и ковариациями. Поэтому для гауссовской
последовательности условия

\[
\mathbb E X_n=m,
\qquad
\operatorname{Cov}(X_n,X_{n+k})=R(k)
\]

уже дают узкую стационарность.

\subsubsection{Гармоника со случайной
фазой}\label{ux433ux430ux440ux43cux43eux43dux438ux43aux430-ux441ux43e-ux441ux43bux443ux447ux430ux439ux43dux43eux439-ux444ux430ux437ux43eux439}

Пусть

\[
X_n=\cos(\lambda n+\Phi),
\]

где \(\Phi\) равномерна на \([0,2\pi]\). Тогда сдвиг \(n\mapsto n+h\)
заменяет фазу \(\Phi\) на \(\Phi+\lambda h\), а равномерное
распределение фазы не меняется. Поэтому последовательность стационарна в
узком смысле.

Этот пример полезен тем, что зависимость во времени сильная, но
распределения при сдвиге не меняются.

\subsection{6. Схемы доказательств /
обоснований}\label{ux441ux445ux435ux43cux44b-ux434ux43eux43aux430ux437ux430ux442ux435ux43bux44cux441ux442ux432-ux43eux431ux43eux441ux43dux43eux432ux430ux43dux438ux439-9}

\subsubsection{Почему узкая стационарность влечет
широкую}\label{ux43fux43eux447ux435ux43cux443-ux443ux437ux43aux430ux44f-ux441ux442ux430ux446ux438ux43eux43dux430ux440ux43dux43eux441ux442ux44c-ux432ux43bux435ux447ux435ux442-ux448ux438ux440ux43eux43aux443ux44e}

Пусть последовательность стационарна в узком смысле и

\[
\mathbb E X_n^2<\infty.
\]

Из

\[
X_n\stackrel d=X_0
\]

следует

\[
\mathbb E X_n=\mathbb E X_0=m.
\]

Из

\[
(X_n,X_{n+k})\stackrel d=(X_0,X_k)
\]

следует равенство вторых смешанных моментов:

\[
\mathbb E[X_nX_{n+k}]=\mathbb E[X_0X_k].
\]

Значит

\[
\operatorname{Cov}(X_n,X_{n+k})
=
\operatorname{Cov}(X_0,X_k),
\]

то есть ковариация зависит только от лага.

\subsubsection{Почему широкая стационарность не влечет
узкую}\label{ux43fux43eux447ux435ux43cux443-ux448ux438ux440ux43eux43aux430ux44f-ux441ux442ux430ux446ux438ux43eux43dux430ux440ux43dux43eux441ux442ux44c-ux43dux435-ux432ux43bux435ux447ux435ux442-ux443ux437ux43aux443ux44e}

Широкая стационарность фиксирует только

\[
\mathbb E X_n
\]

и

\[
\operatorname{Cov}(X_n,X_m).
\]

Она не видит третьи, четвертые и вообще все более высокие
распределительные характеристики. Поэтому можно сделать
последовательность с одинаковыми средними и ковариациями, но с разными
одномерными распределениями в разные моменты времени. Пример с
чередованием нормального и радемахеровского распределения показывает это
явно.

\subsubsection{Почему для гауссовских широкая стационарность дает
узкую}\label{ux43fux43eux447ux435ux43cux443-ux434ux43bux44f-ux433ux430ux443ux441ux441ux43eux432ux441ux43aux438ux445-ux448ux438ux440ux43eux43aux430ux44f-ux441ux442ux430ux446ux438ux43eux43dux430ux440ux43dux43eux441ux442ux44c-ux434ux430ux435ux442-ux443ux437ux43aux443ux44e}

Гауссовский вектор полностью определяется средним вектором и
ковариационной матрицей. Для любого набора моментов

\[
n_1,\ldots,n_k
\]

средний вектор равен

\[
(m,\ldots,m),
\]

а ковариационная матрица имеет элементы

\[
\operatorname{Cov}(X_{n_i},X_{n_j})
=
R(n_j-n_i).
\]

После сдвига всех индексов на \(h\) разности не меняются:

\[
(n_j+h)-(n_i+h)=n_j-n_i.
\]

Значит средний вектор и ковариационная матрица не меняются,
следовательно не меняется и гауссовское конечномерное распределение.

\subsubsection{Почему случайная константа не
эргодична}\label{ux43fux43eux447ux435ux43cux443-ux441ux43bux443ux447ux430ux439ux43dux430ux44f-ux43aux43eux43dux441ux442ux430ux43dux442ux430-ux43dux435-ux44dux440ux433ux43eux434ux438ux447ux43dux430}

Если \(X_n=Y\), то событие

\[
A=\{Y>0\}
\]

не меняется при сдвиге траектории. Если

\[
0<P\{Y>0\}<1,
\]

то это нетривиальное сдвиг-инвариантное событие. Значит процесс не
эргодичен.

\subsection{7. Что написать на
доске}\label{ux447ux442ux43e-ux43dux430ux43fux438ux441ux430ux442ux44c-ux43dux430-ux434ux43eux441ux43aux435-16}

\begin{enumerate}
\def\labelenumi{\arabic{enumi}.}
\tightlist
\item
  Узкая стационарность:
\end{enumerate}

\[
(X_{n_1+h},\ldots,X_{n_k+h})
\stackrel d=
(X_{n_1},\ldots,X_{n_k}).
\]

\begin{enumerate}
\def\labelenumi{\arabic{enumi}.}
\setcounter{enumi}{1}
\tightlist
\item
  Сдвиг:
\end{enumerate}

\[
(Sx)_n=x_{n+1},
\qquad
P_X\circ S^{-1}=P_X.
\]

\begin{enumerate}
\def\labelenumi{\arabic{enumi}.}
\setcounter{enumi}{2}
\tightlist
\item
  Широкая стационарность:
\end{enumerate}

\[
\mathbb E X_n=m,
\qquad
\operatorname{Cov}(X_n,X_{n+k})=R(k).
\]

\begin{enumerate}
\def\labelenumi{\arabic{enumi}.}
\setcounter{enumi}{3}
\tightlist
\item
  Узкая плюс вторые моменты:
\end{enumerate}

\[
\text{узкая стационарность}+\mathbb E X_n^2<\infty
\Longrightarrow
\text{широкая стационарность}.
\]

\begin{enumerate}
\def\labelenumi{\arabic{enumi}.}
\setcounter{enumi}{4}
\tightlist
\item
  Эргодичность:
\end{enumerate}

\[
S^{-1}A=A
\Longrightarrow
P(A)\in\{0,1\}.
\]

\begin{enumerate}
\def\labelenumi{\arabic{enumi}.}
\setcounter{enumi}{5}
\tightlist
\item
  Эргодичность среднего:
\end{enumerate}

\[
\frac1N\sum_{n=0}^{N-1}X_n
\to
\mathbb E X_0.
\]

Для широкого смысла обычно:

\[
\frac1N\sum_{n=0}^{N-1}X_n
\xrightarrow{L^2}
\mathbb E X_0.
\]

\subsection{8. Типичные дополнительные
вопросы}\label{ux442ux438ux43fux438ux447ux43dux44bux435-ux434ux43eux43fux43eux43bux43dux438ux442ux435ux43bux44cux43dux44bux435-ux432ux43eux43fux440ux43eux441ux44b-16}

\textbf{Что является главным объектом билета?}

Случайная последовательность \((X_n)\) и ее поведение при сдвиге
времени. Нужно различать инвариантность всех конечномерных
распределений, инвариантность первых двух моментов и эргодичность.

\textbf{В чем разница между узкой и широкой стационарностью?}

Узкая стационарность требует равенства всех конечномерных распределений
после сдвига:

\[
(X_{n_1+h},\ldots,X_{n_k+h})
\stackrel d=
(X_{n_1},\ldots,X_{n_k}).
\]

Широкая стационарность требует только постоянного среднего и ковариации,
зависящей от лага:

\[
\mathbb E X_n=m,
\qquad
\operatorname{Cov}(X_n,X_{n+k})=R(k).
\]

\textbf{Почему узкая стационарность сильнее широкой?}

Потому что она фиксирует все конечномерные распределения. Если есть
вторые моменты, из этих распределений следуют и средние, и ковариации.
Широкая стационарность фиксирует только первые два момента.

\textbf{Когда широкая стационарность влечет узкую?}

Для гауссовских последовательностей. Гауссовские конечномерные
распределения полностью задаются средним вектором и ковариационной
матрицей.

\textbf{Что такое эргодичность в узком смысле?}

Это отсутствие нетривиальных сдвиг-инвариантных событий:

\[
S^{-1}A=A
\Rightarrow
P(A)\in\{0,1\}.
\]

Интуитивно одна длинная траектория достаточно хорошо представляет весь
ансамбль.

\textbf{Что такое эргодичность в широком смысле?}

Обычно в этом курсе это эргодичность среднего: временное среднее
сходится к математическому ожиданию в \(L^2\):

\[
\frac1N\sum_{n=0}^{N-1}X_n
\xrightarrow{L^2}
\mathbb E X_0.
\]

\textbf{Стационарность означает независимость?}

Нет. Стационарность говорит об инвариантности распределений при сдвиге,
но не запрещает зависимость во времени. Например, \(X_n=Y\) для всех
\(n\) стационарна, но координаты полностью зависимы.

\textbf{Эргодичность означает независимость?}

Нет. Независимые одинаково распределенные последовательности обычно
эргодичны, но эргодичность не равна независимости. Эргодичность
запрещает сдвиг-инвариантные случайные режимы, но допускает зависимость.

\textbf{Почему случайная константа стационарна, но не эргодична?}

Потому что сдвиг не меняет траекторию, значит стационарность есть. Но
временное среднее равно случайной величине \(Y\), а не константе
\(\mathbb EY\). Кроме того, событие \(\{Y>0\}\) сдвиг-инвариантно и
может иметь вероятность между \(0\) и \(1\).

\textbf{Как билет связан с Билетом 24?}

В Билете 24 изучается корреляционная функция

\[
R(k)=\operatorname{Cov}(X_0,X_k)
\]

и достаточные условия, при которых временное среднее сходится к
математическому ожиданию в \(L^2\).

\subsection{9. Типичные
ошибки}\label{ux442ux438ux43fux438ux447ux43dux44bux435-ux43eux448ux438ux431ux43aux438-16}

\begin{itemize}
\tightlist
\item
  Смешивать стационарность и эргодичность.
\item
  Считать, что стационарность означает независимость.
\item
  Считать, что широкая стационарность всегда равносильна узкой.
\item
  Забывать условие существования вторых моментов для широкой
  стационарности.
\item
  Проверять только одномерные распределения и объявлять узкую
  стационарность.
\item
  Называть процесс эргодическим без указания смысла: узкий смысл,
  среднее, \(L^2\).
\item
  Путать ковариационную функцию \(R(k)\) и коэффициент корреляции.
\item
  Забывать, что случайная константа является стационарной, но не
  эргодической.
\end{itemize}

\subsection{10. Связи с другими
билетами}\label{ux441ux432ux44fux437ux438-ux441-ux434ux440ux443ux433ux438ux43cux438-ux431ux438ux43bux435ux442ux430ux43cux438-16}

\begin{itemize}
\tightlist
\item
  Билет 16: стационарность формулируется для закона на пространстве
  последовательностей.
\item
  Билет 15: узкая стационарность - свойство согласованных конечномерных
  распределений.
\item
  Билет 14: широкая стационарность использует ковариации и геометрию
  второго порядка.
\item
  Билет 18: для гауссовских последовательностей широкая стационарность
  влечет узкую.
\item
  Билет 21: временные средние и закон больших чисел связаны с
  эргодической идеей.
\item
  Билет 24: корреляционная функция и достаточные условия эргодичности
  среднего.
\item
  Билет 25: спектральные характеристики строятся для стационарных
  последовательностей второго порядка.
\end{itemize}

\subsection{11. Минимум на
удовлетворительно}\label{ux43cux438ux43dux438ux43cux443ux43c-ux43dux430-ux443ux434ux43eux432ux43bux435ux442ux432ux43eux440ux438ux442ux435ux43bux44cux43dux43e-16}

Нужно сказать:

\begin{enumerate}
\def\labelenumi{\arabic{enumi}.}
\tightlist
\item
  Узкая стационарность:
\end{enumerate}

\[
(X_{n_1+h},\ldots,X_{n_k+h})
\stackrel d=
(X_{n_1},\ldots,X_{n_k}).
\]

\begin{enumerate}
\def\labelenumi{\arabic{enumi}.}
\setcounter{enumi}{1}
\tightlist
\item
  Широкая стационарность:
\end{enumerate}

\[
\mathbb E X_n=m,
\qquad
\operatorname{Cov}(X_n,X_{n+k})=R(k).
\]

\begin{enumerate}
\def\labelenumi{\arabic{enumi}.}
\setcounter{enumi}{2}
\item
  Узкая плюс вторые моменты влечет широкую.
\item
  Эргодичность: временные средние отражают вероятностные средние; строго
  - нет нетривиальных сдвиг-инвариантных событий.
\item
  Главная ловушка: стационарность не равна эргодичности и не равна
  независимости.
\end{enumerate}

\subsection{12. Минимум на
хорошо}\label{ux43cux438ux43dux438ux43cux443ux43c-ux43dux430-ux445ux43eux440ux43eux448ux43e-16}

Помимо минимума, нужно объяснить:

\begin{itemize}
\tightlist
\item
  что такое сдвиг на пространстве последовательностей;
\item
  почему узкая стационарность означает инвариантность закона
  относительно сдвига;
\item
  почему широкая стационарность требует вторых моментов;
\item
  почему широкая не всегда влечет узкую;
\item
  почему для гауссовских последовательностей широкая стационарность
  влечет узкую;
\item
  разобрать пример iid-последовательности и случайной константы.
\end{itemize}

\subsection{13. Что добавить для
отлично}\label{ux447ux442ux43e-ux434ux43eux431ux430ux432ux438ux442ux44c-ux434ux43bux44f-ux43eux442ux43bux438ux447ux43dux43e-16}

Для сильного ответа стоит добавить:

\begin{itemize}
\tightlist
\item
  формулировку эргодичности через сдвиг-инвариантные события;
\item
  интерпретацию через эргодическую теорему:
\end{itemize}

\[
\frac1N\sum_{n=0}^{N-1}f(X_n)\to\mathbb E f(X_0);
\]

\begin{itemize}
\tightlist
\item
  определение эргодичности в широком смысле как \(L^2\)-сходимости
  среднего;
\item
  формулу для дисперсии среднего:
\end{itemize}

\[
\operatorname{Var}(\overline X_N)
=
\frac1{N^2}\sum_{i,j=0}^{N-1}R(j-i);
\]

\begin{itemize}
\tightlist
\item
  пример широкой стационарности без узкой;
\item
  пример стационарной неэргодической последовательности;
\item
  связь со спектральной теорией и корреляционной функцией.
\end{itemize}

\subsection{14. Устная версия ответа на 5
минут}\label{ux443ux441ux442ux43dux430ux44f-ux432ux435ux440ux441ux438ux44f-ux43eux442ux432ux435ux442ux430-ux43dux430-5-ux43cux438ux43dux443ux442-16}

\begin{enumerate}
\def\labelenumi{\arabic{enumi}.}
\tightlist
\item
  В этом билете изучаются свойства случайной последовательности при
  сдвиге времени.
\item
  Узкая стационарность означает, что все конечномерные распределения не
  меняются при сдвиге:
\end{enumerate}

\[
(X_{n_1+h},\ldots,X_{n_k+h})
\stackrel d=
(X_{n_1},\ldots,X_{n_k}).
\]

\begin{enumerate}
\def\labelenumi{\arabic{enumi}.}
\setcounter{enumi}{2}
\tightlist
\item
  На пространстве траекторий это означает инвариантность закона
  относительно сдвига:
\end{enumerate}

\[
P_X\circ S^{-1}=P_X.
\]

\begin{enumerate}
\def\labelenumi{\arabic{enumi}.}
\setcounter{enumi}{3}
\tightlist
\item
  Широкая стационарность требует вторых моментов, постоянного среднего и
  ковариации только от лага:
\end{enumerate}

\[
\mathbb E X_n=m,
\qquad
\operatorname{Cov}(X_n,X_{n+k})=R(k).
\]

\begin{enumerate}
\def\labelenumi{\arabic{enumi}.}
\setcounter{enumi}{4}
\tightlist
\item
  Узкая стационарность при наличии вторых моментов влечет широкую.
\item
  Обратное неверно, потому что широкая стационарность контролирует
  только первые два момента.
\item
  Для гауссовских последовательностей обратное верно: среднее и
  ковариация задают все конечномерные распределения.
\item
  Эргодичность в узком смысле означает, что сдвиг-инвариантные события
  имеют вероятность \(0\) или \(1\).
\item
  Практически это означает, что временные средние совпадают с
  вероятностными:
\end{enumerate}

\[
\frac1N\sum_{n=0}^{N-1}f(X_n)\to\mathbb E f(X_0).
\]

\begin{enumerate}
\def\labelenumi{\arabic{enumi}.}
\setcounter{enumi}{9}
\tightlist
\item
  В широком смысле обычно говорят об эргодичности среднего:
\end{enumerate}

\[
\frac1N\sum_{n=0}^{N-1}X_n
\xrightarrow{L^2}
\mathbb E X_0.
\]

\begin{enumerate}
\def\labelenumi{\arabic{enumi}.}
\setcounter{enumi}{10}
\tightlist
\item
  Примеры: iid-последовательность стационарна и эргодична; случайная
  константа стационарна, но не эргодична.
\item
  Завершить можно связью с корреляционной функцией: условия на \(R(k)\)
  дают эргодичность среднего.
\end{enumerate}

\subsection{15. Если осталось 2 минуты перед
экзаменом}\label{ux435ux441ux43bux438-ux43eux441ux442ux430ux43bux43eux441ux44c-2-ux43cux438ux43dux443ux442ux44b-ux43fux435ux440ux435ux434-ux44dux43aux437ux430ux43cux435ux43dux43eux43c-16}

Запомнить три различия.

Узкая стационарность:

\[
(X_{n_1+h},\ldots,X_{n_k+h})
\stackrel d=
(X_{n_1},\ldots,X_{n_k}).
\]

Широкая стационарность:

\[
\mathbb E X_n=m,
\qquad
\operatorname{Cov}(X_n,X_{n+k})=R(k).
\]

Эргодичность:

\[
S^{-1}A=A
\Rightarrow
P(A)\in\{0,1\},
\]

или для среднего во втором порядке:

\[
\frac1N\sum_{n=0}^{N-1}X_n
\xrightarrow{L^2}
\mathbb E X_0.
\]

Главная ловушка: стационарность, эргодичность и независимость - разные
свойства.

\newpage

\section{Билет 18. Гауссовы случайные величины, векторы и
последовательности}\label{ux431ux438ux43bux435ux442-18.-ux433ux430ux443ux441ux441ux43eux432ux44b-ux441ux43bux443ux447ux430ux439ux43dux44bux435-ux432ux435ux43bux438ux447ux438ux43dux44b-ux432ux435ux43aux442ux43eux440ux44b-ux438-ux43fux43eux441ux43bux435ux434ux43eux432ux430ux442ux435ux43bux44cux43dux43eux441ux442ux438}

Источники: Ширяев 1, §13; Ширяев 2, гл. VI.

\subsection{1. Короткий
ответ}\label{ux43aux43eux440ux43eux442ux43aux438ux439-ux43eux442ux432ux435ux442-17}

Нормальная случайная величина с параметрами \(m\in\mathbb R\) и
\(\sigma^2\ge0\) имеет характеристическую функцию

\[
\varphi_X(t)
=
\mathbb E e^{itX}
=
\exp\left\{itm-\frac12\sigma^2t^2\right\}.
\]

Если \(\sigma^2>0\), то ее плотность:

\[
p(x)=\frac1{\sqrt{2\pi\sigma^2}}
\exp\left\{-\frac{(x-m)^2}{2\sigma^2}\right\}.
\]

Случайный вектор

\[
X=(X_1,\ldots,X_n)^T
\]

называется гауссовским, если любая его линейная комбинация

\[
\sum_{j=1}^n u_jX_j
\]

имеет нормальное, возможно вырожденное, распределение. Закон
гауссовского вектора полностью задается средним вектором

\[
m=\mathbb EX
\]

и ковариационной матрицей

\[
\Sigma=\operatorname{Cov}(X).
\]

Его характеристическая функция:

\[
\varphi_X(u)
=
\exp\left\{i\langle u,m\rangle-\frac12u^T\Sigma u\right\}.
\]

Случайная последовательность называется гауссовской, если любой ее
конечный вектор

\[
(X_{n_1},\ldots,X_{n_k})
\]

является гауссовским.

Главные свойства: линейные преобразования гауссовских векторов гауссовы;
для совместно гауссовских величин некоррелированность эквивалентна
независимости; для гауссовской последовательности широкая стационарность
влечет узкую.

\subsection{2. Полный
ответ}\label{ux43fux43eux43bux43dux44bux439-ux43eux442ux432ux435ux442-17}

\subsubsection{2.1. Введение в гауссовские распределения и
процессы}\label{ux432ux432ux435ux434ux435ux43dux438ux435-ux432-ux433ux430ux443ux441ux441ux43eux432ux441ux43aux438ux435-ux440ux430ux441ux43fux440ux435ux434ux435ux43bux435ux43dux438ux44f-ux438-ux43fux440ux43eux446ux435ux441ux441ux44b}

Гауссовы случайные величины и процессы важны потому, что они являются
основным классом распределений второго порядка. Для произвольного
процесса среднее и ковариация обычно не задают весь закон. Для
гауссовского процесса задают: все конечномерные распределения полностью
определяются средними и ковариациями.

\subsubsection{2.2. Одномерное нормальное
распределение}\label{ux43eux434ux43dux43eux43cux435ux440ux43dux43eux435-ux43dux43eux440ux43cux430ux43bux44cux43dux43eux435-ux440ux430ux441ux43fux440ux435ux434ux435ux43bux435ux43dux438ux435}

Начнем с одномерного случая. Случайная величина \(X\) имеет нормальное
распределение с параметрами \(m\) и \(\sigma^2>0\), если ее плотность
равна

\[
p(x)
=
\frac1{\sqrt{2\pi\sigma^2}}
\exp\left\{-\frac{(x-m)^2}{2\sigma^2}\right\}.
\]

Пишут:

\[
X\sim N(m,\sigma^2).
\]

Здесь

\[
\mathbb EX=m,
\qquad
\operatorname{Var}X=\sigma^2.
\]

Стандартная нормальная величина имеет распределение

\[
N(0,1).
\]

Если \(Z\sim N(0,1)\), то

\[
X=m+\sigma Z
\]

имеет распределение \(N(m,\sigma^2)\).

\subsubsection{2.3. Вырожденный случай и характеристическая
функция}\label{ux432ux44bux440ux43eux436ux434ux435ux43dux43dux44bux439-ux441ux43bux443ux447ux430ux439-ux438-ux445ux430ux440ux430ux43aux442ux435ux440ux438ux441ux442ux438ux447ux435ux441ux43aux430ux44f-ux444ux443ux43dux43aux446ux438ux44f}

Разрешают и вырожденный случай \(\sigma^2=0\). Тогда

\[
X=m
\]

почти наверное. Это тоже удобно считать нормальным распределением: оно
не имеет обычной плотности относительно меры Лебега, но его
характеристическая функция получается из той же формулы.

Характеристическая функция нормальной величины:

\[
\varphi_X(t)
=
\mathbb E e^{itX}
=
\exp\left\{itm-\frac12\sigma^2t^2\right\}.
\]

Эта формула важнее плотности, потому что она сразу обобщается на
гауссовские векторы.

\subsubsection{2.4. Определение гауссовского
вектора}\label{ux43eux43fux440ux435ux434ux435ux43bux435ux43dux438ux435-ux433ux430ux443ux441ux441ux43eux432ux441ux43aux43eux433ux43e-ux432ux435ux43aux442ux43eux440ux430}

Теперь случайный вектор. Вектор

\[
X=(X_1,\ldots,X_n)^T
\]

называется гауссовским, если для любого вектора коэффициентов

\[
u=(u_1,\ldots,u_n)^T\in\mathbb R^n
\]

линейная комбинация

\[
\langle u,X\rangle
=
\sum_{j=1}^n u_jX_j
\]

имеет нормальное распределение, возможно вырожденное.

Важно: недостаточно, чтобы каждая координата \(X_j\) отдельно была
нормально распределена. Нужно, чтобы нормальной была любая линейная
комбинация координат. Именно это условие контролирует совместное
распределение.

\subsubsection{2.5. Параметры и распределение гауссовского
вектора}\label{ux43fux430ux440ux430ux43cux435ux442ux440ux44b-ux438-ux440ux430ux441ux43fux440ux435ux434ux435ux43bux435ux43dux438ux435-ux433ux430ux443ux441ux441ux43eux432ux441ux43aux43eux433ux43e-ux432ux435ux43aux442ux43eux440ux430}

Для гауссовского вектора определены:

\[
m=\mathbb EX
=
(\mathbb EX_1,\ldots,\mathbb EX_n)^T
\]

и ковариационная матрица

\[
\Sigma=(\Sigma_{ij}),
\qquad
\Sigma_{ij}
=
\operatorname{Cov}(X_i,X_j).
\]

Матрица \(\Sigma\) симметрична и неотрицательно определена:

\[
u^T\Sigma u
=
\operatorname{Var}\langle u,X\rangle
\ge0.
\]

Характеристическая функция гауссовского вектора имеет вид

\[
\varphi_X(u)
=
\mathbb E e^{i\langle u,X\rangle}
=
\exp\left\{
i\langle u,m\rangle
-\frac12u^T\Sigma u
\right\}.
\]

Эта формула показывает, что закон гауссовского вектора полностью
задается парой

\[
(m,\Sigma).
\]

Обратно, если \(m\in\mathbb R^n\), а \(\Sigma\) - симметричная
неотрицательно определенная матрица, то существует гауссовский вектор с
такими средним и ковариацией. Если \(\Sigma\) положительно определена,
то распределение имеет плотность:

\[
p(x)
=
\frac1{(2\pi)^{n/2}(\det\Sigma)^{1/2}}
\exp\left\{
-\frac12(x-m)^T\Sigma^{-1}(x-m)
\right\}.
\]

Если \(\Sigma\) вырождена, то плотности относительно \(n\)-мерной меры
Лебега нет: распределение сосредоточено на некотором аффинном
подпространстве. Но вектор все равно называется гауссовским, потому что
все линейные комбинации нормальны.

\subsubsection{2.6. Линейные преобразования гауссовских
векторов}\label{ux43bux438ux43dux435ux439ux43dux44bux435-ux43fux440ux435ux43eux431ux440ux430ux437ux43eux432ux430ux43dux438ux44f-ux433ux430ux443ux441ux441ux43eux432ux441ux43aux438ux445-ux432ux435ux43aux442ux43eux440ux43eux432}

Линейные преобразования гауссовских векторов снова гауссовы. Если

\[
Y=AX+b,
\]

где \(A\) - матрица, \(b\) - вектор, а \(X\) гауссовский, то \(Y\)
гауссовский. Его параметры:

\[
\mathbb EY=Am+b,
\qquad
\operatorname{Cov}(Y)=A\Sigma A^T.
\]

Доказательство простое: любая линейная комбинация координат \(Y\)
является линейной комбинацией координат \(X\) плюс константа, а значит
нормальна.

\subsubsection{2.7. Свойства независимости и
некоррелированности}\label{ux441ux432ux43eux439ux441ux442ux432ux430-ux43dux435ux437ux430ux432ux438ux441ux438ux43cux43eux441ux442ux438-ux438-ux43dux435ux43aux43eux440ux440ux435ux43bux438ux440ux43eux432ux430ux43dux43dux43eux441ux442ux438}

Главное специальное свойство гауссовских векторов: для совместно
гауссовских случайных величин некоррелированность равносильна
независимости. В общем случае из

\[
\operatorname{Cov}(X,Y)=0
\]

независимость не следует, как обсуждалось в Билете 14. Но если \((X,Y)\)
- гауссовский вектор, то из нулевой ковариации следует независимость.

Более общая формулировка: если гауссовский вектор разбит на две группы
координат

\[
X=(X^{(1)},X^{(2)}),
\]

и все ковариации между первой и второй группой равны нулю, то эти два
случайных вектора независимы. В терминах матрицы ковариаций это означает
блочно-диагональную матрицу:

\[
\Sigma=
\begin{pmatrix}
\Sigma_1 & 0\\
0 & \Sigma_2
\end{pmatrix}.
\]

Тогда характеристическая функция факторизуется:

\[
\varphi_X(u_1,u_2)
=
\varphi_{X^{(1)}}(u_1)\varphi_{X^{(2)}}(u_2),
\]

а факторизация характеристической функции означает независимость.

\subsubsection{2.8. Гауссовские
последовательности}\label{ux433ux430ux443ux441ux441ux43eux432ux441ux43aux438ux435-ux43fux43eux441ux43bux435ux434ux43eux432ux430ux442ux435ux43bux44cux43dux43eux441ux442ux438}

Теперь гауссовские последовательности. Случайная последовательность

\[
\{X_n:n\in T\}
\]

называется гауссовской, если для любого конечного набора индексов

\[
n_1,\ldots,n_k
\]

случайный вектор

\[
(X_{n_1},\ldots,X_{n_k})
\]

является гауссовским. Эквивалентно: любая конечная линейная комбинация

\[
\sum_{j=1}^k u_jX_{n_j}
\]

имеет нормальное распределение.

Гауссовская последовательность задается средней функцией

\[
m_n=\mathbb EX_n
\]

и ковариационной функцией

\[
K(n,m)=\operatorname{Cov}(X_n,X_m).
\]

Для любого конечного набора индексов \(n_1,\ldots,n_k\) средний вектор
равен

\[
(m_{n_1},\ldots,m_{n_k}),
\]

а ковариационная матрица равна

\[
(K(n_i,n_j))_{i,j=1}^k.
\]

Чтобы такая последовательность существовала, функция \(K\) должна быть
симметричной и неотрицательно определенной:

\[
\sum_{i,j=1}^k a_i a_jK(n_i,n_j)\ge0
\]

для любых \(k\), индексов \(n_i\) и коэффициентов \(a_i\). Тогда
согласованные гауссовские конечномерные распределения задают процесс по
теореме Колмогорова из Билета 15.

\subsubsection{2.9. Стационарность гауссовских
последовательностей}\label{ux441ux442ux430ux446ux438ux43eux43dux430ux440ux43dux43eux441ux442ux44c-ux433ux430ux443ux441ux441ux43eux432ux441ux43aux438ux445-ux43fux43eux441ux43bux435ux434ux43eux432ux430ux442ux435ux43bux44cux43dux43eux441ux442ux435ux439}

Для стационарности гауссовских последовательностей есть важное
упрощение. В Билете 17 различались узкая и широкая стационарность. Для
произвольных процессов широкая стационарность не влечет узкую. Но для
гауссовской последовательности влечет. Если

\[
\mathbb EX_n=m
\]

и

\[
\operatorname{Cov}(X_n,X_{n+k})=R(k),
\]

то все конечномерные распределения не меняются при сдвиге, потому что у
них не меняются средние векторы и ковариационные матрицы. Значит
гауссовская последовательность стационарна в узком смысле.

\subsection{3. Основные
определения}\label{ux43eux441ux43dux43eux432ux43dux44bux435-ux43eux43fux440ux435ux434ux435ux43bux435ux43dux438ux44f-17}

\subsubsection{Нормальная случайная
величина}\label{ux43dux43eux440ux43cux430ux43bux44cux43dux430ux44f-ux441ux43bux443ux447ux430ux439ux43dux430ux44f-ux432ux435ux43bux438ux447ux438ux43dux430}

Случайная величина \(X\) имеет распределение \(N(m,\sigma^2)\), если при
\(\sigma^2>0\) ее плотность равна

\[
p(x)=\frac1{\sqrt{2\pi\sigma^2}}
\exp\left\{-\frac{(x-m)^2}{2\sigma^2}\right\}.
\]

При \(\sigma^2=0\) понимают вырожденное распределение \(X=m\) почти
наверное.

\subsubsection{Характеристическая функция нормальной
величины}\label{ux445ux430ux440ux430ux43aux442ux435ux440ux438ux441ux442ux438ux447ux435ux441ux43aux430ux44f-ux444ux443ux43dux43aux446ux438ux44f-ux43dux43eux440ux43cux430ux43bux44cux43dux43eux439-ux432ux435ux43bux438ux447ux438ux43dux44b}

\[
\varphi_X(t)=\exp\left\{itm-\frac12\sigma^2t^2\right\}.
\]

\subsubsection{Гауссовский
вектор}\label{ux433ux430ux443ux441ux441ux43eux432ux441ux43aux438ux439-ux432ux435ux43aux442ux43eux440}

Вектор \(X=(X_1,\ldots,X_n)^T\) называется гауссовским, если для любого
\(u\in\mathbb R^n\) величина \(\langle u,X\rangle\) нормально
распределена, возможно вырожденно.

\subsubsection{Средний
вектор}\label{ux441ux440ux435ux434ux43dux438ux439-ux432ux435ux43aux442ux43eux440}

\[
m=\mathbb EX.
\]

\subsubsection{Ковариационная
матрица}\label{ux43aux43eux432ux430ux440ux438ux430ux446ux438ux43eux43dux43dux430ux44f-ux43cux430ux442ux440ux438ux446ux430-2}

\[
\Sigma_{ij}=\operatorname{Cov}(X_i,X_j).
\]

\subsubsection{Характеристическая функция гауссовского
вектора}\label{ux445ux430ux440ux430ux43aux442ux435ux440ux438ux441ux442ux438ux447ux435ux441ux43aux430ux44f-ux444ux443ux43dux43aux446ux438ux44f-ux433ux430ux443ux441ux441ux43eux432ux441ux43aux43eux433ux43e-ux432ux435ux43aux442ux43eux440ux430}

\[
\varphi_X(u)
=
\exp\left\{
i\langle u,m\rangle-\frac12u^T\Sigma u
\right\}.
\]

\subsubsection{Гауссовская
последовательность}\label{ux433ux430ux443ux441ux441ux43eux432ux441ux43aux430ux44f-ux43fux43eux441ux43bux435ux434ux43eux432ux430ux442ux435ux43bux44cux43dux43eux441ux442ux44c}

Последовательность \((X_n)\) называется гауссовской, если каждый
конечный вектор

\[
(X_{n_1},\ldots,X_{n_k})
\]

гауссовский.

\subsubsection{Белый гауссовский
шум}\label{ux431ux435ux43bux44bux439-ux433ux430ux443ux441ux441ux43eux432ux441ux43aux438ux439-ux448ux443ux43c}

Последовательность независимых одинаково распределенных гауссовских
величин с нулевым средним:

\[
X_n\sim N(0,\sigma^2),
\qquad
\operatorname{Cov}(X_n,X_m)=0\quad(n\ne m).
\]

\subsection{4. Основные теоремы и
свойства}\label{ux43eux441ux43dux43eux432ux43dux44bux435-ux442ux435ux43eux440ux435ux43cux44b-ux438-ux441ux432ux43eux439ux441ux442ux432ux430-17}

\begin{itemize}
\tightlist
\item
  Нормальный закон задается средним и дисперсией:
\end{itemize}

\[
X\sim N(m,\sigma^2).
\]

\begin{itemize}
\tightlist
\item
  Гауссовский вектор задается средним вектором и ковариационной
  матрицей:
\end{itemize}

\[
\varphi_X(u)
=
\exp\left\{i\langle u,m\rangle-\frac12u^T\Sigma u\right\}.
\]

\begin{itemize}
\tightlist
\item
  Матрица \(\Sigma\) должна быть симметричной и неотрицательно
  определенной.
\item
  Если \(\Sigma>0\), то гауссовский вектор имеет плотность:
\end{itemize}

\[
p(x)
=
\frac1{(2\pi)^{n/2}(\det\Sigma)^{1/2}}
\exp\left\{-\frac12(x-m)^T\Sigma^{-1}(x-m)\right\}.
\]

\begin{itemize}
\tightlist
\item
  Линейное преобразование гауссовского вектора гауссово:
\end{itemize}

\[
Y=AX+b
\quad\Rightarrow\quad
Y\ \text{гауссовский}.
\]

\begin{itemize}
\tightlist
\item
  Для совместно гауссовских величин некоррелированность равносильна
  независимости.
\item
  Нормальность отдельных координат не гарантирует совместную
  гауссовость.
\item
  Гауссовская последовательность задается средней функцией \(m_n\) и
  ковариационной функцией \(K(n,m)\).
\item
  Для гауссовской последовательности широкая стационарность влечет
  узкую.
\end{itemize}

\subsection{5. Примеры}\label{ux43fux440ux438ux43cux435ux440ux44b-17}

\subsubsection{Стандартная нормальная
величина}\label{ux441ux442ux430ux43dux434ux430ux440ux442ux43dux430ux44f-ux43dux43eux440ux43cux430ux43bux44cux43dux430ux44f-ux432ux435ux43bux438ux447ux438ux43dux430}

Если

\[
Z\sim N(0,1),
\]

то

\[
\varphi_Z(t)=e^{-t^2/2}.
\]

Если

\[
X=m+\sigma Z,
\]

то

\[
X\sim N(m,\sigma^2).
\]

\subsubsection{Двумерный гауссовский
вектор}\label{ux434ux432ux443ux43cux435ux440ux43dux44bux439-ux433ux430ux443ux441ux441ux43eux432ux441ux43aux438ux439-ux432ux435ux43aux442ux43eux440}

Пусть

\[
X=(X_1,X_2)^T
\]

имеет средний вектор

\[
m=(m_1,m_2)^T
\]

и ковариационную матрицу

\[
\Sigma=
\begin{pmatrix}
\sigma_1^2 & \rho\sigma_1\sigma_2\\
\rho\sigma_1\sigma_2 & \sigma_2^2
\end{pmatrix},
\qquad
|\rho|\le1.
\]

Если вектор гауссовский, то при \(\rho=0\) координаты независимы. Это
специальное гауссовское свойство.

\subsubsection{Нормальные координаты без совместной
гауссовости}\label{ux43dux43eux440ux43cux430ux43bux44cux43dux44bux435-ux43aux43eux43eux440ux434ux438ux43dux430ux442ux44b-ux431ux435ux437-ux441ux43eux432ux43cux435ux441ux442ux43dux43eux439-ux433ux430ux443ux441ux441ux43eux432ux43eux441ux442ux438}

Пусть \(Z\sim N(0,1)\), а \(\varepsilon\) независимо принимает значения
\(1\) и \(-1\) с вероятностями \(1/2\). Положим

\[
X=Z,
\qquad
Y=\varepsilon Z.
\]

Тогда \(X\) и \(Y\) по отдельности имеют распределение \(N(0,1)\). Но
вектор \((X,Y)\) не является гауссовским: например, сумма

\[
X+Y=(1+\varepsilon)Z
\]

с вероятностью \(1/2\) равна \(0\), а с вероятностью \(1/2\) равна
\(2Z\), поэтому не имеет нормального распределения. Этот пример
показывает, почему в определении нужна нормальность всех линейных
комбинаций, а не только координат.

\subsubsection{Белый гауссовский
шум}\label{ux431ux435ux43bux44bux439-ux433ux430ux443ux441ux441ux43eux432ux441ux43aux438ux439-ux448ux443ux43c-1}

Пусть

\[
X_n
\]

независимы и

\[
X_n\sim N(0,\sigma^2).
\]

Тогда \((X_n)\) - гауссовская последовательность,

\[
\mathbb EX_n=0,
\]

и

\[
R(k)=\operatorname{Cov}(X_0,X_k)
=
\begin{cases}
\sigma^2, & k=0,\\
0, & k\ne0.
\end{cases}
\]

Это базовый пример стационарной гауссовской последовательности.

\subsubsection{Гауссовский
AR(1)}\label{ux433ux430ux443ux441ux441ux43eux432ux441ux43aux438ux439-ar1}

Пусть

\[
X_n=aX_{n-1}+\xi_n,
\qquad |a|<1,
\]

где \(\xi_n\) - белый гауссовский шум с дисперсией \(\sigma_\xi^2\). В
стационарном режиме

\[
\mathbb EX_n=0,
\qquad
R(k)=\frac{\sigma_\xi^2}{1-a^2}a^{|k|}.
\]

Такой процесс гауссовский, потому что строится линейными операциями из
гауссовских величин.

\subsection{6. Схемы доказательств /
обоснований}\label{ux441ux445ux435ux43cux44b-ux434ux43eux43aux430ux437ux430ux442ux435ux43bux44cux441ux442ux432-ux43eux431ux43eux441ux43dux43eux432ux430ux43dux438ux439-10}

\subsubsection{\texorpdfstring{Почему закон гауссовского вектора
задается \(m\) и
\(\Sigma\)}{Почему закон гауссовского вектора задается m и \textbackslash Sigma}}\label{ux43fux43eux447ux435ux43cux443-ux437ux430ux43aux43eux43d-ux433ux430ux443ux441ux441ux43eux432ux441ux43aux43eux433ux43e-ux432ux435ux43aux442ux43eux440ux430-ux437ux430ux434ux430ux435ux442ux441ux44f-m-ux438-sigma}

Для любого \(u\) величина

\[
\langle u,X\rangle
\]

нормальна. Ее среднее:

\[
\mathbb E\langle u,X\rangle=\langle u,m\rangle.
\]

Ее дисперсия:

\[
\operatorname{Var}\langle u,X\rangle=u^T\Sigma u.
\]

Значит

\[
\mathbb E e^{i\langle u,X\rangle}
=
\exp\left\{i\langle u,m\rangle-\frac12u^T\Sigma u\right\}.
\]

Характеристическая функция однозначно задает распределение, поэтому
закон определяется \(m\) и \(\Sigma\).

\subsubsection{Почему линейное преобразование гауссовского вектора
гауссово}\label{ux43fux43eux447ux435ux43cux443-ux43bux438ux43dux435ux439ux43dux43eux435-ux43fux440ux435ux43eux431ux440ux430ux437ux43eux432ux430ux43dux438ux435-ux433ux430ux443ux441ux441ux43eux432ux441ux43aux43eux433ux43e-ux432ux435ux43aux442ux43eux440ux430-ux433ux430ux443ux441ux441ux43eux432ux43e}

Пусть

\[
Y=AX+b.
\]

Для любого \(v\):

\[
\langle v,Y\rangle
=
\langle v,AX+b\rangle
=
\langle A^Tv,X\rangle+\langle v,b\rangle.
\]

Это линейная комбинация координат \(X\) плюс константа, значит
нормальная величина. Следовательно, \(Y\) гауссовский.

\subsubsection{Почему некоррелированность дает независимость для
Гаусса}\label{ux43fux43eux447ux435ux43cux443-ux43dux435ux43aux43eux440ux440ux435ux43bux438ux440ux43eux432ux430ux43dux43dux43eux441ux442ux44c-ux434ux430ux435ux442-ux43dux435ux437ux430ux432ux438ux441ux438ux43cux43eux441ux442ux44c-ux434ux43bux44f-ux433ux430ux443ux441ux441ux430}

Пусть \(X=(X^{(1)},X^{(2)})\) - гауссовский вектор, и ковариации между
группами равны нулю. Тогда

\[
\Sigma=
\begin{pmatrix}
\Sigma_1 & 0\\
0 & \Sigma_2
\end{pmatrix}.
\]

Характеристическая функция:

\[
\varphi(u_1,u_2)
=
\exp\left\{
i\langle u_1,m_1\rangle+i\langle u_2,m_2\rangle
-\frac12u_1^T\Sigma_1u_1
-\frac12u_2^T\Sigma_2u_2
\right\}.
\]

Она раскладывается в произведение:

\[
\varphi(u_1,u_2)
=
\varphi_1(u_1)\varphi_2(u_2).
\]

Факторизация характеристической функции означает независимость.

\subsubsection{Почему широкая стационарность гауссовской
последовательности влечет
узкую}\label{ux43fux43eux447ux435ux43cux443-ux448ux438ux440ux43eux43aux430ux44f-ux441ux442ux430ux446ux438ux43eux43dux430ux440ux43dux43eux441ux442ux44c-ux433ux430ux443ux441ux441ux43eux432ux441ux43aux43eux439-ux43fux43eux441ux43bux435ux434ux43eux432ux430ux442ux435ux43bux44cux43dux43eux441ux442ux438-ux432ux43bux435ux447ux435ux442-ux443ux437ux43aux443ux44e}

Для любого конечного набора индексов \(n_1,\ldots,n_k\) распределение
вектора

\[
(X_{n_1},\ldots,X_{n_k})
\]

гауссовское и задается средним вектором и ковариационной матрицей. Если

\[
\mathbb EX_n=m,
\qquad
\operatorname{Cov}(X_n,X_{n+k})=R(k),
\]

то после сдвига всех индексов на \(h\) средний вектор не меняется, а
ковариационные элементы зависят от тех же разностей:

\[
(n_j+h)-(n_i+h)=n_j-n_i.
\]

Значит не меняется все конечномерное распределение.

\subsection{7. Что написать на
доске}\label{ux447ux442ux43e-ux43dux430ux43fux438ux441ux430ux442ux44c-ux43dux430-ux434ux43eux441ux43aux435-17}

\begin{enumerate}
\def\labelenumi{\arabic{enumi}.}
\tightlist
\item
  Нормальная величина:
\end{enumerate}

\[
X\sim N(m,\sigma^2),
\qquad
\varphi_X(t)=\exp\left\{itm-\frac12\sigma^2t^2\right\}.
\]

\begin{enumerate}
\def\labelenumi{\arabic{enumi}.}
\setcounter{enumi}{1}
\tightlist
\item
  Гауссовский вектор:
\end{enumerate}

\[
X\ \text{гауссовский}
\Longleftrightarrow
\langle u,X\rangle\ \text{нормальна для всех}\ u.
\]

\begin{enumerate}
\def\labelenumi{\arabic{enumi}.}
\setcounter{enumi}{2}
\tightlist
\item
  Характеристическая функция:
\end{enumerate}

\[
\varphi_X(u)
=
\exp\left\{i\langle u,m\rangle-\frac12u^T\Sigma u\right\}.
\]

\begin{enumerate}
\def\labelenumi{\arabic{enumi}.}
\setcounter{enumi}{3}
\tightlist
\item
  Ковариационная матрица:
\end{enumerate}

\[
\Sigma_{ij}=\operatorname{Cov}(X_i,X_j),
\qquad
u^T\Sigma u\ge0.
\]

\begin{enumerate}
\def\labelenumi{\arabic{enumi}.}
\setcounter{enumi}{4}
\tightlist
\item
  Некоррелированность и независимость:
\end{enumerate}

\[
\text{для совместно гауссовских: }
\operatorname{Cov}(X,Y)=0
\Longleftrightarrow
X,Y\ \text{независимы}.
\]

\begin{enumerate}
\def\labelenumi{\arabic{enumi}.}
\setcounter{enumi}{5}
\tightlist
\item
  Гауссовская последовательность:
\end{enumerate}

\[
(X_{n_1},\ldots,X_{n_k})\ \text{гауссовский для всех конечных наборов}.
\]

\begin{enumerate}
\def\labelenumi{\arabic{enumi}.}
\setcounter{enumi}{6}
\tightlist
\item
  Для Гаусса:
\end{enumerate}

\[
\text{широкая стационарность}
\Longrightarrow
\text{узкая стационарность}.
\]

\subsection{8. Типичные дополнительные
вопросы}\label{ux442ux438ux43fux438ux447ux43dux44bux435-ux434ux43eux43fux43eux43bux43dux438ux442ux435ux43bux44cux43dux44bux435-ux432ux43eux43fux440ux43eux441ux44b-17}

\textbf{Что является главным объектом билета?}

Нормальные случайные величины, гауссовские векторы и гауссовские
последовательности. Главное отличие гауссовского объекта: все
конечномерные распределения задаются средними и ковариациями.

\textbf{Почему в определении гауссовского вектора нужны все линейные
комбинации?}

Потому что нормальность отдельных координат не определяет совместный
закон. Нужно требовать, чтобы

\[
\sum_j u_jX_j
\]

была нормальной для любых коэффициентов \(u_j\).

\textbf{Что задает закон гауссовского вектора?}

Средний вектор \(m\) и ковариационная матрица \(\Sigma\).
Характеристическая функция:

\[
\varphi_X(u)
=
\exp\left\{i\langle u,m\rangle-\frac12u^T\Sigma u\right\}.
\]

\textbf{Какие условия нужны на \(\Sigma\)?}

Матрица должна быть симметричной и неотрицательно определенной:

\[
u^T\Sigma u\ge0.
\]

Если \(\Sigma\) положительно определена, есть плотность в
\(\mathbb R^n\); если вырождена, распределение сосредоточено на
подпространстве.

\textbf{Когда некоррелированность означает независимость?}

Для совместно гауссовских величин или групп координат. В общем случае
это неверно.

\textbf{Почему для гауссовских некоррелированность дает независимость?}

Потому что при нулевых межгрупповых ковариациях ковариационная матрица
становится блочно-диагональной, и характеристическая функция
раскладывается в произведение характеристических функций групп.

\textbf{Что такое гауссовская последовательность?}

Это последовательность, у которой любой конечный вектор

\[
(X_{n_1},\ldots,X_{n_k})
\]

гауссовский.

\textbf{Как задать гауссовскую последовательность?}

Нужно задать среднюю функцию \(m_n\) и неотрицательно определенную
ковариационную функцию \(K(n,m)\). Тогда конечномерные гауссовские
распределения согласованы и задают процесс.

\textbf{Почему широкая стационарность гауссовской последовательности
влечет узкую?}

Потому что гауссовские конечномерные распределения полностью
определяются средними и ковариациями. При широкой стационарности они не
меняются при сдвиге.

\textbf{Какая главная ошибка в этом билете?}

Считать, что если каждая координата нормальна, то весь вектор
автоматически гауссовский. Это неверно: нужна нормальность всех линейных
комбинаций.

\subsection{9. Типичные
ошибки}\label{ux442ux438ux43fux438ux447ux43dux44bux435-ux43eux448ux438ux431ux43aux438-17}

\begin{itemize}
\tightlist
\item
  Считать нормальность координат достаточной для совместной гауссовости.
\item
  Забывать вырожденные нормальные распределения.
\item
  Писать плотность многомерного нормального закона при вырожденной
  \(\Sigma\).
\item
  Забывать условие \(\Sigma\ge0\).
\item
  Переносить свойство ``некоррелированность равна независимости'' на
  негауссовские распределения.
\item
  Путать ковариационную матрицу и корреляционную матрицу.
\item
  Считать, что широкая стационарность всегда влечет узкую; это верно не
  вообще, а для гауссовских последовательностей.
\end{itemize}

\subsection{10. Связи с другими
билетами}\label{ux441ux432ux44fux437ux438-ux441-ux434ux440ux443ux433ux438ux43cux438-ux431ux438ux43bux435ux442ux430ux43cux438-17}

\begin{itemize}
\tightlist
\item
  Билет 08: характеристические функции задают распределения.
\item
  Билет 13: средние и ковариации получаются из производных
  характеристической функции.
\item
  Билет 14: ковариация, корреляция и некоррелированность; для Гаусса
  некоррелированность дает независимость.
\item
  Билет 15: согласованные гауссовские конечномерные распределения задают
  гауссовскую последовательность.
\item
  Билет 17: для гауссовских последовательностей широкая стационарность
  влечет узкую.
\item
  Билет 24: корреляционная функция стационарной гауссовской
  последовательности.
\item
  Билет 25: спектральные характеристики стационарных последовательностей
  второго порядка.
\end{itemize}

\subsection{11. Минимум на
удовлетворительно}\label{ux43cux438ux43dux438ux43cux443ux43c-ux43dux430-ux443ux434ux43eux432ux43bux435ux442ux432ux43eux440ux438ux442ux435ux43bux44cux43dux43e-17}

Нужно сказать:

\begin{enumerate}
\def\labelenumi{\arabic{enumi}.}
\tightlist
\item
  Одномерный нормальный закон:
\end{enumerate}

\[
X\sim N(m,\sigma^2),
\qquad
\varphi_X(t)=\exp\left\{itm-\frac12\sigma^2t^2\right\}.
\]

\begin{enumerate}
\def\labelenumi{\arabic{enumi}.}
\setcounter{enumi}{1}
\tightlist
\item
  Гауссовский вектор определяется нормальностью всех линейных
  комбинаций:
\end{enumerate}

\[
\langle u,X\rangle\ \text{нормальна для всех}\ u.
\]

\begin{enumerate}
\def\labelenumi{\arabic{enumi}.}
\setcounter{enumi}{2}
\tightlist
\item
  Его характеристическая функция:
\end{enumerate}

\[
\varphi_X(u)=
\exp\left\{i\langle u,m\rangle-\frac12u^T\Sigma u\right\}.
\]

\begin{enumerate}
\def\labelenumi{\arabic{enumi}.}
\setcounter{enumi}{3}
\item
  Гауссовская последовательность: каждый конечный вектор гауссовский.
\item
  Главная ловушка: нормальность координат не равна совместной
  гауссовости.
\end{enumerate}

\subsection{12. Минимум на
хорошо}\label{ux43cux438ux43dux438ux43cux443ux43c-ux43dux430-ux445ux43eux440ux43eux448ux43e-17}

Помимо минимума, нужно объяснить:

\begin{itemize}
\tightlist
\item
  что такое вырожденный нормальный закон;
\item
  почему \(\Sigma\) должна быть неотрицательно определенной;
\item
  почему линейное преобразование гауссовского вектора гауссово;
\item
  почему для совместно гауссовских некоррелированность равна
  независимости;
\item
  как задается гауссовская последовательность через \(m_n\) и
  \(K(n,m)\);
\item
  почему для гауссовских широкая стационарность влечет узкую.
\end{itemize}

\subsection{13. Что добавить для
отлично}\label{ux447ux442ux43e-ux434ux43eux431ux430ux432ux438ux442ux44c-ux434ux43bux44f-ux43eux442ux43bux438ux447ux43dux43e-17}

Для сильного ответа стоит добавить:

\begin{itemize}
\tightlist
\item
  плотность многомерного нормального закона при \(\Sigma>0\);
\item
  предупреждение про вырожденный случай \(\det\Sigma=0\);
\item
  доказательство через характеристические функции;
\item
  пример нормальных координат без совместной гауссовости;
\item
  блочно-диагональную ковариационную матрицу как критерий независимости
  групп;
\item
  связь с теоремой Колмогорова для построения гауссовской
  последовательности;
\item
  связь со стационарностью и спектральными характеристиками.
\end{itemize}

\subsection{14. Устная версия ответа на 5
минут}\label{ux443ux441ux442ux43dux430ux44f-ux432ux435ux440ux441ux438ux44f-ux43eux442ux432ux435ux442ux430-ux43dux430-5-ux43cux438ux43dux443ux442-17}

\begin{enumerate}
\def\labelenumi{\arabic{enumi}.}
\tightlist
\item
  Начинаю с одномерного нормального закона:
\end{enumerate}

\[
X\sim N(m,\sigma^2),
\qquad
\varphi_X(t)=\exp\left\{itm-\frac12\sigma^2t^2\right\}.
\]

\begin{enumerate}
\def\labelenumi{\arabic{enumi}.}
\setcounter{enumi}{1}
\tightlist
\item
  Гауссовский вектор \(X=(X_1,\ldots,X_n)\) определяется так: любая
  линейная комбинация координат нормальна.
\item
  Это сильнее, чем нормальность каждой координаты отдельно.
\item
  Средний вектор и ковариационная матрица:
\end{enumerate}

\[
m=\mathbb EX,
\qquad
\Sigma_{ij}=\operatorname{Cov}(X_i,X_j).
\]

\begin{enumerate}
\def\labelenumi{\arabic{enumi}.}
\setcounter{enumi}{4}
\tightlist
\item
  Характеристическая функция:
\end{enumerate}

\[
\varphi_X(u)
=
\exp\left\{i\langle u,m\rangle-\frac12u^T\Sigma u\right\}.
\]

\begin{enumerate}
\def\labelenumi{\arabic{enumi}.}
\setcounter{enumi}{5}
\tightlist
\item
  Поэтому гауссовский закон полностью задается \(m\) и \(\Sigma\).
\item
  Линейное преобразование гауссовского вектора снова гауссово:
\end{enumerate}

\[
Y=AX+b.
\]

\begin{enumerate}
\def\labelenumi{\arabic{enumi}.}
\setcounter{enumi}{7}
\tightlist
\item
  Для совместно гауссовских величин нулевая ковариация означает
  независимость. В общем случае это неверно.
\item
  Гауссовская последовательность - это последовательность, у которой
  любой конечный вектор гауссовский.
\item
  Она задается средней функцией \(m_n\) и ковариационной функцией
  \(K(n,m)\).
\item
  Для гауссовской последовательности широкая стационарность влечет
  узкую, потому что конечномерные распределения задаются средними и
  ковариациями.
\item
  Примеры: стандартная нормальная величина, белый гауссовский шум,
  гауссовский \(AR(1)\).
\end{enumerate}

\subsection{15. Если осталось 2 минуты перед
экзаменом}\label{ux435ux441ux43bux438-ux43eux441ux442ux430ux43bux43eux441ux44c-2-ux43cux438ux43dux443ux442ux44b-ux43fux435ux440ux435ux434-ux44dux43aux437ux430ux43cux435ux43dux43eux43c-17}

Запомнить главное.

Одномерная нормальная:

\[
\varphi_X(t)=\exp\left\{itm-\frac12\sigma^2t^2\right\}.
\]

Гауссовский вектор:

\[
\langle u,X\rangle\ \text{нормальна для всех}\ u.
\]

Характеристическая функция:

\[
\varphi_X(u)
=
\exp\left\{i\langle u,m\rangle-\frac12u^T\Sigma u\right\}.
\]

Гауссовская последовательность:

\[
(X_{n_1},\ldots,X_{n_k})\ \text{гауссовский для всех конечных наборов}.
\]

Главные ловушки:

\begin{itemize}
\tightlist
\item
  нормальность координат не означает совместную гауссовость;
\item
  некоррелированность равна независимости только в совместно гауссовском
  случае;
\item
  плотность с \(\Sigma^{-1}\) пишется только при невырожденной
  \(\Sigma\).
\end{itemize}

\newpage

\section{Билет 19. Случайные последовательности с независимыми
приращениями.
Блуждания}\label{ux431ux438ux43bux435ux442-19.-ux441ux43bux443ux447ux430ux439ux43dux44bux435-ux43fux43eux441ux43bux435ux434ux43eux432ux430ux442ux435ux43bux44cux43dux43eux441ux442ux438-ux441-ux43dux435ux437ux430ux432ux438ux441ux438ux43cux44bux43cux438-ux43fux440ux438ux440ux430ux449ux435ux43dux438ux44fux43cux438.-ux431ux43bux443ux436ux434ux430ux43dux438ux44f}

Источники: Ширяев 2; Сердобольская, лекции 2-3.

\subsection{1. Короткий
ответ}\label{ux43aux43eux440ux43eux442ux43aux438ux439-ux43eux442ux432ux435ux442-18}

Случайная последовательность \((S_n)_{n\ge0}\) имеет независимые
приращения, если для любых

\[
0\le n_0<n_1<\cdots<n_k
\]

случайные величины

\[
S_{n_1}-S_{n_0},\quad S_{n_2}-S_{n_1},\quad\ldots,\quad S_{n_k}-S_{n_{k-1}}
\]

независимы.

Главный дискретный пример - случайное блуждание:

\[
S_0=0,
\qquad
S_n=\xi_1+\cdots+\xi_n,
\]

где \(\xi_1,\xi_2,\ldots\) - независимые случайные величины. Если шаги
одинаково распределены, то приращения не только независимы, но и
стационарны:

\[
S_{n+k}-S_n\stackrel d=S_k.
\]

Для iid-шагов характеристическая функция суммы:

\[
\varphi_{S_n}(t)=\left(\varphi_\xi(t)\right)^n.
\]

Если \(\mathbb E\xi_1=a\) и \(\operatorname{Var}\xi_1=\sigma^2\), то

\[
\mathbb ES_n=na,
\qquad
\operatorname{Var}S_n=n\sigma^2.
\]

\subsection{2. Полный
ответ}\label{ux43fux43eux43bux43dux44bux439-ux43eux442ux432ux435ux442-18}

\subsubsection{19.1. Определение независимых
приращений}\label{ux43eux43fux440ux435ux434ux435ux43bux435ux43dux438ux435-ux43dux435ux437ux430ux432ux438ux441ux438ux43cux44bux445-ux43fux440ux438ux440ux430ux449ux435ux43dux438ux439}

Этот билет продолжает Билет 16 о случайных последовательностях и связан
с Билетом 09 о независимости. Здесь важный объект не просто
последовательность независимых величин, а процесс, у которого независимы
изменения на непересекающихся промежутках времени.

Пусть

\[
(S_n)_{n\ge0}
\]

\begin{itemize}
\tightlist
\item
  случайная последовательность со значениями в \(\mathbb R\),
  \(\mathbb Z\) или \(\mathbb R^d\). Приращением за промежуток от \(m\)
  до \(n\) называется
\end{itemize}

\[
S_n-S_m,
\qquad
0\le m<n.
\]

Последовательность имеет независимые приращения, если для любых моментов
времени

\[
0\le n_0<n_1<\cdots<n_k
\]

приращения по непересекающимся промежуткам

\[
S_{n_1}-S_{n_0},\quad
S_{n_2}-S_{n_1},\quad\ldots,\quad
S_{n_k}-S_{n_{k-1}}
\]

независимы.

Важно: независимые приращения не означают, что сами значения
\(S_0,S_1,S_2,\ldots\) независимы. Наоборот, значения обычно сильно
зависимы, потому что

\[
S_{n+1}=S_n+(S_{n+1}-S_n).
\]

Например, в случайном блуждании \(S_{10}\) содержит в себе все шаги,
которые входят в \(S_9\).

\subsubsection{19.2. Стационарные
приращения}\label{ux441ux442ux430ux446ux438ux43eux43dux430ux440ux43dux44bux435-ux43fux440ux438ux440ux430ux449ux435ux43dux438ux44f}

Если распределение приращения зависит только от длины промежутка, то
говорят о стационарных приращениях:

\[
S_{n+k}-S_n\stackrel d=S_k-S_0.
\]

Если обычно берут \(S_0=0\), то:

\[
S_{n+k}-S_n\stackrel d=S_k.
\]

\subsubsection{19.3. Конструкция процесса через суммы независимых
величин}\label{ux43aux43eux43dux441ux442ux440ux443ux43aux446ux438ux44f-ux43fux440ux43eux446ux435ux441ux441ux430-ux447ux435ux440ux435ux437-ux441ux443ux43cux43cux44b-ux43dux435ux437ux430ux432ux438ux441ux438ux43cux44bux445-ux432ux435ux43bux438ux447ux438ux43d}

Главная конструкция дискретного процесса с независимыми приращениями -
сумма независимых шагов. Пусть

\[
\xi_1,\xi_2,\ldots
\]

\begin{itemize}
\tightlist
\item
  независимые случайные величины. Определим
\end{itemize}

\[
S_0=0,
\qquad
S_n=\sum_{j=1}^n\xi_j.
\]

Тогда

\[
S_n-S_m=\sum_{j=m+1}^n\xi_j.
\]

Если промежутки времени не пересекаются, то соответствующие суммы
состоят из непересекающихся наборов независимых величин \(\xi_j\).
Поэтому такие приращения независимы.

Если \(\xi_j\) одинаково распределены, то приращение

\[
S_{n+k}-S_n=\xi_{n+1}+\cdots+\xi_{n+k}
\]

имеет то же распределение, что

\[
S_k=\xi_1+\cdots+\xi_k.
\]

Значит при iid-шагах процесс имеет независимые и стационарные
приращения.

\subsubsection{19.4. Характеристические функции и
моменты}\label{ux445ux430ux440ux430ux43aux442ux435ux440ux438ux441ux442ux438ux447ux435ux441ux43aux438ux435-ux444ux443ux43dux43aux446ux438ux438-ux438-ux43cux43eux43cux435ux43dux442ux44b}

Распределение суммы выражается через свертку. Если шаги iid с законом
\(\mu\), то

\[
S_n\sim \mu^{*n},
\]

где \(\mu^{*n}\) - \(n\)-кратная свертка меры \(\mu\). Через
характеристические функции это особенно просто. Пусть

\[
\varphi_\xi(t)=\mathbb E e^{it\xi_1}.
\]

Тогда из независимости:

\[
\varphi_{S_n}(t)
=
\mathbb E e^{it(\xi_1+\cdots+\xi_n)}
=
\prod_{j=1}^n\mathbb E e^{it\xi_j}
=
\left(\varphi_\xi(t)\right)^n.
\]

Для не одинаково распределенных независимых шагов:

\[
\varphi_{S_n}(t)=\prod_{j=1}^n\varphi_{\xi_j}(t).
\]

Если существуют первые два момента, то они складываются:

\[
\mathbb ES_n
=
\sum_{j=1}^n\mathbb E\xi_j,
\]

и при независимости

\[
\operatorname{Var}S_n
=
\sum_{j=1}^n\operatorname{Var}\xi_j.
\]

В iid-случае, если

\[
\mathbb E\xi_1=a,
\qquad
\operatorname{Var}\xi_1=\sigma^2,
\]

получаем

\[
\mathbb ES_n=na,
\qquad
\operatorname{Var}S_n=n\sigma^2.
\]

\subsubsection{19.5. Случайное
блуждание}\label{ux441ux43bux443ux447ux430ux439ux43dux43eux435-ux431ux43bux443ux436ux434ux430ux43dux438ux435-1}

Случайное блуждание - это процесс вида

\[
S_n=S_0+\xi_1+\cdots+\xi_n,
\]

обычно с \(S_0=0\) или заданным начальным состоянием. Если шаги
принимают значения в \(\mathbb Z\), то блуждание идет по целочисленной
решетке. Самый классический пример - простое блуждание на \(\mathbb Z\):

\[
P\{\xi_k=1\}=p,
\qquad
P\{\xi_k=-1\}=q=1-p.
\]

Если

\[
p=q=\frac12,
\]

блуждание называется простым симметричным.

Для простого блуждания

\[
S_n=\xi_1+\cdots+\xi_n.
\]

Если за \(n\) шагов было \(r\) шагов вправо и \(n-r\) шагов влево, то

\[
S_n=r-(n-r)=2r-n.
\]

Поэтому

\[
P\{S_n=x\}
=
\binom n{(n+x)/2}p^{(n+x)/2}q^{(n-x)/2},
\]

если \(n+x\) четно и \(|x|\le n\); иначе вероятность равна \(0\).

Для простого блуждания:

\[
\mathbb E\xi_1=p-q=2p-1,
\]

\[
\operatorname{Var}\xi_1=1-(p-q)^2=4pq.
\]

Следовательно,

\[
\mathbb ES_n=n(p-q),
\qquad
\operatorname{Var}S_n=4npq.
\]

В симметричном случае:

\[
\mathbb ES_n=0,
\qquad
\operatorname{Var}S_n=n.
\]

\subsubsection{19.6. Марковское свойство и
асимптотика}\label{ux43cux430ux440ux43aux43eux432ux441ux43aux43eux435-ux441ux432ux43eux439ux441ux442ux432ux43e-ux438-ux430ux441ux438ux43cux43fux442ux43eux442ux438ux43aux430}

Случайное блуждание является марковской цепью. Если текущее состояние
равно \(i\), то следующее состояние равно

\[
i+\xi_{n+1}.
\]

Будущий шаг независим от прошлого, поэтому условное распределение
будущего зависит только от текущего состояния. Для простого блуждания на
\(\mathbb Z\) переходные вероятности:

\[
p_{i,i+1}=p,
\qquad
p_{i,i-1}=q,
\]

а остальные переходы равны нулю. Это связывает билет с Билетом 22 и
Билетом 23.

Случайное блуждание обычно не является стационарным по значениям:
распределение \(S_n\) меняется с \(n\), дисперсия растет как \(n\). Но
при iid-шагах оно имеет стационарные приращения. Это типичная ловушка:
стационарные приращения не равны стационарности самого процесса.

Асимптотически блуждание связано с законом больших чисел и центральной
предельной теоремой. Если шаги iid и имеют конечное среднее \(a\), то

\[
\frac{S_n}{n}\xrightarrow{P} a
\]

по закону больших чисел. Если дисперсия \(\sigma^2>0\) конечна, то

\[
\frac{S_n-na}{\sigma\sqrt n}
\Rightarrow
N(0,1).
\]

Это связь с Билетом 21.

\subsection{3. Основные
определения}\label{ux43eux441ux43dux43eux432ux43dux44bux435-ux43eux43fux440ux435ux434ux435ux43bux435ux43dux438ux44f-18}

\subsubsection{Приращение}\label{ux43fux440ux438ux440ux430ux449ux435ux43dux438ux435}

Для последовательности \((S_n)\) приращение на промежутке \((m,n]\):

\[
S_n-S_m.
\]

\subsubsection{Независимые
приращения}\label{ux43dux435ux437ux430ux432ux438ux441ux438ux43cux44bux435-ux43fux440ux438ux440ux430ux449ux435ux43dux438ux44f}

Последовательность имеет независимые приращения, если для любых

\[
0\le n_0<n_1<\cdots<n_k
\]

величины

\[
S_{n_1}-S_{n_0},\ldots,S_{n_k}-S_{n_{k-1}}
\]

независимы.

\subsubsection{Стационарные
приращения}\label{ux441ux442ux430ux446ux438ux43eux43dux430ux440ux43dux44bux435-ux43fux440ux438ux440ux430ux449ux435ux43dux438ux44f-1}

Распределение приращения зависит только от длины промежутка:

\[
S_{n+k}-S_n\stackrel d=S_k-S_0.
\]

Если \(S_0=0\), то

\[
S_{n+k}-S_n\stackrel d=S_k.
\]

\subsubsection{Случайное
блуждание}\label{ux441ux43bux443ux447ux430ux439ux43dux43eux435-ux431ux43bux443ux436ux434ux430ux43dux438ux435-2}

Процесс, задаваемый суммами шагов:

\[
S_0=x,
\qquad
S_n=x+\sum_{j=1}^n\xi_j,
\]

где \(\xi_j\) - независимые шаги.

\subsubsection{\texorpdfstring{Простое блуждание на
\(\mathbb Z\)}{Простое блуждание на \textbackslash mathbb Z}}\label{ux43fux440ux43eux441ux442ux43eux435-ux431ux43bux443ux436ux434ux430ux43dux438ux435-ux43dux430-mathbb-z}

Блуждание с шагами

\[
\xi_j\in\{-1,1\}.
\]

Если

\[
P\{\xi_j=1\}=P\{\xi_j=-1\}=\frac12,
\]

оно называется простым симметричным блужданием.

\subsubsection{Блуждание с одинаково распределенными
шагами}\label{ux431ux43bux443ux436ux434ux430ux43dux438ux435-ux441-ux43eux434ux438ux43dux430ux43aux43eux432ux43e-ux440ux430ux441ux43fux440ux435ux434ux435ux43bux435ux43dux43dux44bux43cux438-ux448ux430ux433ux430ux43cux438}

Блуждание с одинаково распределенными шагами. В этом случае приращения
стационарны.

\subsection{4. Основные теоремы и
свойства}\label{ux43eux441ux43dux43eux432ux43dux44bux435-ux442ux435ux43eux440ux435ux43cux44b-ux438-ux441ux432ux43eux439ux441ux442ux432ux430-18}

\begin{itemize}
\tightlist
\item
  Если
\end{itemize}

\[
S_n=S_0+\sum_{j=1}^n\xi_j
\]

и \(\xi_j\) независимы, то \((S_n)\) имеет независимые приращения.

\begin{itemize}
\tightlist
\item
  Если шаги iid, то приращения независимы и стационарны:
\end{itemize}

\[
S_{n+k}-S_n\stackrel d=S_k-S_0.
\]

\begin{itemize}
\tightlist
\item
  Распределение суммы независимых шагов - свертка:
\end{itemize}

\[
P_{S_n}=\mu_1*\cdots*\mu_n.
\]

В iid-случае:

\[
P_{S_n}=\mu^{*n}.
\]

\begin{itemize}
\tightlist
\item
  Характеристические функции независимых шагов перемножаются:
\end{itemize}

\[
\varphi_{S_n}(t)=\prod_{j=1}^n\varphi_{\xi_j}(t).
\]

В iid-случае:

\[
\varphi_{S_n}(t)=\left(\varphi_\xi(t)\right)^n.
\]

\begin{itemize}
\tightlist
\item
  При существовании моментов:
\end{itemize}

\[
\mathbb ES_n=S_0+\sum_{j=1}^n\mathbb E\xi_j,
\]

\[
\operatorname{Var}S_n=\sum_{j=1}^n\operatorname{Var}\xi_j.
\]

\begin{itemize}
\tightlist
\item
  Случайное блуждание с независимыми шагами является марковской цепью.
\item
  Блуждание обычно не стационарно по значениям, но при iid-шагах имеет
  стационарные приращения.
\item
  Для iid-шагов применимы ЗБЧ и ЦПТ:
\end{itemize}

\[
\frac{S_n}{n}\xrightarrow{P}\mathbb E\xi_1,
\qquad
\frac{S_n-n\mathbb E\xi_1}{\sqrt{n\operatorname{Var}\xi_1}}
\Rightarrow N(0,1).
\]

\subsection{5. Примеры}\label{ux43fux440ux438ux43cux435ux440ux44b-18}

\subsubsection{Простое симметричное
блуждание}\label{ux43fux440ux43eux441ux442ux43eux435-ux441ux438ux43cux43cux435ux442ux440ux438ux447ux43dux43eux435-ux431ux43bux443ux436ux434ux430ux43dux438ux435}

Пусть

\[
P\{\xi_k=1\}=P\{\xi_k=-1\}=\frac12.
\]

Тогда

\[
S_n=\xi_1+\cdots+\xi_n
\]

и

\[
P\{S_n=x\}
=
\binom n{(n+x)/2}2^{-n},
\]

если \(n+x\) четно и \(|x|\le n\). Иначе вероятность равна \(0\).

Среднее и дисперсия:

\[
\mathbb ES_n=0,
\qquad
\operatorname{Var}S_n=n.
\]

\subsubsection{Несимметричное
блуждание}\label{ux43dux435ux441ux438ux43cux43cux435ux442ux440ux438ux447ux43dux43eux435-ux431ux43bux443ux436ux434ux430ux43dux438ux435}

Пусть

\[
P\{\xi_k=1\}=p,
\qquad
P\{\xi_k=-1\}=q=1-p.
\]

Тогда

\[
\mathbb ES_n=n(p-q),
\qquad
\operatorname{Var}S_n=4npq.
\]

Если \(p>q\), блуждание имеет положительный дрейф; если \(p<q\),
отрицательный.

\subsubsection{Сумма
Бернулли}\label{ux441ux443ux43cux43cux430-ux431ux435ux440ux43dux443ux43bux43bux438}

Пусть

\[
\xi_k\sim \operatorname{Bernoulli}(p),
\qquad
\xi_k\in\{0,1\}.
\]

Тогда

\[
S_n=\sum_{k=1}^n\xi_k
\]

считает число успехов за \(n\) испытаний:

\[
S_n\sim \operatorname{Bin}(n,p).
\]

Это процесс с независимыми стационарными приращениями, но он не является
блужданием на всей \(\mathbb Z\), а является неубывающим счетчиком.

\subsubsection{\texorpdfstring{Блуждание в
\(\mathbb R^d\)}{Блуждание в \textbackslash mathbb R\^{}d}}\label{ux431ux43bux443ux436ux434ux430ux43dux438ux435-ux432-mathbb-rd}

Пусть шаги \(\xi_k\) принимают значения в \(\mathbb R^d\). Тогда

\[
S_n=S_0+\xi_1+\cdots+\xi_n
\]

задает блуждание в многомерном пространстве. Например, на решетке
\(\mathbb Z^2\) можно на каждом шаге переходить к одному из соседей.

\subsubsection{Пуассоновский процесс как
аналог}\label{ux43fux443ux430ux441ux441ux43eux43dux43eux432ux441ux43aux438ux439-ux43fux440ux43eux446ux435ux441ux441-ux43aux430ux43a-ux430ux43dux430ux43bux43eux433}

Пуассоновский процесс \(N(t)\) в непрерывном времени имеет независимые и
стационарные приращения:

\[
N(t+s)-N(t)\sim \operatorname{Poisson}(\lambda s).
\]

Это не дискретная случайная последовательность, но полезный аналог: он
показывает, что идея независимых приращений используется и в непрерывном
времени.

\subsection{6. Схемы доказательств /
обоснований}\label{ux441ux445ux435ux43cux44b-ux434ux43eux43aux430ux437ux430ux442ux435ux43bux44cux441ux442ux432-ux43eux431ux43eux441ux43dux43eux432ux430ux43dux438ux439-11}

\subsubsection{Почему сумма независимых шагов имеет независимые
приращения}\label{ux43fux43eux447ux435ux43cux443-ux441ux443ux43cux43cux430-ux43dux435ux437ux430ux432ux438ux441ux438ux43cux44bux445-ux448ux430ux433ux43eux432-ux438ux43cux435ux435ux442-ux43dux435ux437ux430ux432ux438ux441ux438ux43cux44bux435-ux43fux440ux438ux440ux430ux449ux435ux43dux438ux44f}

Пусть

\[
S_n=\sum_{j=1}^n\xi_j.
\]

Для моментов

\[
0\le n_0<n_1<\cdots<n_k
\]

приращение

\[
S_{n_r}-S_{n_{r-1}}
=
\sum_{j=n_{r-1}+1}^{n_r}\xi_j.
\]

Разные приращения зависят от непересекающихся наборов независимых шагов.
Поэтому они независимы.

\subsubsection{Почему характеристическая функция суммы
факторизуется}\label{ux43fux43eux447ux435ux43cux443-ux445ux430ux440ux430ux43aux442ux435ux440ux438ux441ux442ux438ux447ux435ux441ux43aux430ux44f-ux444ux443ux43dux43aux446ux438ux44f-ux441ux443ux43cux43cux44b-ux444ux430ux43aux442ux43eux440ux438ux437ux443ux435ux442ux441ux44f}

Если \(\xi_j\) независимы, то

\[
\varphi_{S_n}(t)
=
\mathbb E\exp\left(it\sum_{j=1}^n\xi_j\right)
=
\mathbb E\prod_{j=1}^n e^{it\xi_j}.
\]

По независимости математическое ожидание произведения равно произведению
математических ожиданий:

\[
\varphi_{S_n}(t)
=
\prod_{j=1}^n\mathbb E e^{it\xi_j}
=
\prod_{j=1}^n\varphi_{\xi_j}(t).
\]

\subsubsection{Почему блуждание является марковской
цепью}\label{ux43fux43eux447ux435ux43cux443-ux431ux43bux443ux436ux434ux430ux43dux438ux435-ux44fux432ux43bux44fux435ux442ux441ux44f-ux43cux430ux440ux43aux43eux432ux441ux43aux43eux439-ux446ux435ux43fux44cux44e}

Для блуждания

\[
S_{n+1}=S_n+\xi_{n+1}.
\]

Шаг \(\xi_{n+1}\) независим от прошлого

\[
S_0,\ldots,S_n.
\]

Поэтому условное распределение \(S_{n+1}\) при известном прошлом зависит
только от \(S_n\). В счетном случае:

\[
P\{S_{n+1}=j\mid S_0=i_0,\ldots,S_n=i\}
=
P\{\xi_{n+1}=j-i\}.
\]

Это марковское свойство.

\subsubsection{Почему блуждание не стационарно по
значениям}\label{ux43fux43eux447ux435ux43cux443-ux431ux43bux443ux436ux434ux430ux43dux438ux435-ux43dux435-ux441ux442ux430ux446ux438ux43eux43dux430ux440ux43dux43e-ux43fux43e-ux437ux43dux430ux447ux435ux43dux438ux44fux43c}

Для симметричного блуждания

\[
\operatorname{Var}S_n=n.
\]

Значит распределение \(S_n\) зависит от \(n\). Например, \(S_0=0\) почти
наверное, а \(S_1\) принимает значения \(\pm1\). Поэтому значения
процесса не стационарны, хотя приращения стационарны.

\subsection{7. Что написать на
доске}\label{ux447ux442ux43e-ux43dux430ux43fux438ux441ux430ux442ux44c-ux43dux430-ux434ux43eux441ux43aux435-18}

\begin{enumerate}
\def\labelenumi{\arabic{enumi}.}
\tightlist
\item
  Приращение:
\end{enumerate}

\[
S_n-S_m.
\]

\begin{enumerate}
\def\labelenumi{\arabic{enumi}.}
\setcounter{enumi}{1}
\tightlist
\item
  Независимые приращения:
\end{enumerate}

\[
S_{n_1}-S_{n_0},\ldots,S_{n_k}-S_{n_{k-1}}
\quad\text{независимы}.
\]

\begin{enumerate}
\def\labelenumi{\arabic{enumi}.}
\setcounter{enumi}{2}
\tightlist
\item
  Блуждание:
\end{enumerate}

\[
S_0=0,
\qquad
S_n=\sum_{j=1}^n\xi_j.
\]

\begin{enumerate}
\def\labelenumi{\arabic{enumi}.}
\setcounter{enumi}{3}
\tightlist
\item
  Стационарные приращения при iid-шагах:
\end{enumerate}

\[
S_{n+k}-S_n\stackrel d=S_k.
\]

\begin{enumerate}
\def\labelenumi{\arabic{enumi}.}
\setcounter{enumi}{4}
\tightlist
\item
  Характеристическая функция:
\end{enumerate}

\[
\varphi_{S_n}(t)=\left(\varphi_\xi(t)\right)^n.
\]

\begin{enumerate}
\def\labelenumi{\arabic{enumi}.}
\setcounter{enumi}{5}
\tightlist
\item
  Моменты:
\end{enumerate}

\[
\mathbb ES_n=n\mathbb E\xi_1,
\qquad
\operatorname{Var}S_n=n\operatorname{Var}\xi_1.
\]

\begin{enumerate}
\def\labelenumi{\arabic{enumi}.}
\setcounter{enumi}{6}
\tightlist
\item
  Простое блуждание:
\end{enumerate}

\[
P\{\xi=1\}=p,
\qquad
P\{\xi=-1\}=q.
\]

\subsection{8. Типичные дополнительные
вопросы}\label{ux442ux438ux43fux438ux447ux43dux44bux435-ux434ux43eux43fux43eux43bux43dux438ux442ux435ux43bux44cux43dux44bux435-ux432ux43eux43fux440ux43eux441ux44b-18}

\textbf{Что является главным объектом билета?}

Случайная последовательность, у которой независимы приращения на
непересекающихся промежутках. Главный пример - случайное блуждание, то
есть сумма независимых шагов.

\textbf{В чем разница между независимыми значениями и независимыми
приращениями?}

Независимые значения означали бы независимость \(S_0,S_1,\ldots\).
Независимые приращения означают независимость разностей \(S_n-S_m\) на
непересекающихся промежутках. У блуждания значения обычно зависимы, но
приращения независимы.

\textbf{Что значит стационарность приращений?}

Это значит, что закон приращения зависит только от длины промежутка:

\[
S_{n+k}-S_n\stackrel d=S_k-S_0.
\]

При \(S_0=0\):

\[
S_{n+k}-S_n\stackrel d=S_k.
\]

\textbf{Почему iid-шаги дают независимые стационарные приращения?}

Независимость приращений следует из независимости непересекающихся
наборов шагов. Стационарность следует из одинакового распределения
шагов: сумма любых \(k\) последовательных шагов имеет один и тот же
закон.

\textbf{Как найти распределение \(S_n\)?}

В общем случае это свертка законов шагов:

\[
P_{S_n}=\mu_1*\cdots*\mu_n.
\]

В iid-случае:

\[
P_{S_n}=\mu^{*n}.
\]

Для простого блуждания получается биномиальная формула.

\textbf{Почему характеристическая функция равна степени?}

Потому что характеристическая функция суммы независимых величин равна
произведению характеристических функций:

\[
\varphi_{S_n}(t)
=
\prod_{j=1}^n\varphi_{\xi_j}(t).
\]

Если шаги одинаково распределены:

\[
\varphi_{S_n}(t)=\varphi_\xi(t)^n.
\]

\textbf{Почему блуждание является марковской цепью?}

Потому что

\[
S_{n+1}=S_n+\xi_{n+1},
\]

а новый шаг независим от прошлого. Поэтому при известном \(S_n\) прошлые
состояния не дают дополнительной информации о \(S_{n+1}\).

\textbf{Случайное блуждание стационарно?}

По значениям обычно нет: распределение \(S_n\) меняется с \(n\). Но при
iid-шагах у него стационарные приращения.

\textbf{Как связаны блуждания с ЗБЧ и ЦПТ?}

Если шаги iid, то

\[
\frac{S_n}{n}\xrightarrow{P}\mathbb E\xi_1
\]

по закону больших чисел, а при конечной ненулевой дисперсии

\[
\frac{S_n-n\mathbb E\xi_1}{\sqrt{n\operatorname{Var}\xi_1}}
\Rightarrow N(0,1).
\]

\textbf{Какая главная ошибка в этом билете?}

Путать независимость приращений с независимостью самих значений процесса
и путать стационарность приращений со стационарностью процесса.

\subsection{9. Типичные
ошибки}\label{ux442ux438ux43fux438ux447ux43dux44bux435-ux43eux448ux438ux431ux43aux438-18}

\begin{itemize}
\tightlist
\item
  Путать независимые значения и независимые приращения.
\item
  Считать, что случайное блуждание стационарно по значениям.
\item
  Забывать начальное состояние \(S_0\).
\item
  Забывать, что стационарность приращений требует одинакового
  распределения шагов.
\item
  Писать \(\varphi_{S_n}=\varphi_\xi^n\) без условия одинакового
  распределения шагов.
\item
  Складывать дисперсии без независимости шагов.
\item
  Путать блуждание на \(\mathbb Z\) и счетчик Бернулли: во втором случае
  шаги \(0/1\), а не \(\pm1\).
\item
  Считать марковское свойство полной независимостью состояний.
\end{itemize}

\subsection{10. Связи с другими
билетами}\label{ux441ux432ux44fux437ux438-ux441-ux434ux440ux443ux433ux438ux43cux438-ux431ux438ux43bux435ux442ux430ux43cux438-18}

\begin{itemize}
\tightlist
\item
  Билет 09: независимость и факторизация математических ожиданий
  используются для приращений и характеристических функций.
\item
  Билет 16: блуждание - случайная последовательность и случайный элемент
  пространства последовательностей.
\item
  Билет 17: у блуждания при iid-шагах стационарные приращения, но обычно
  не стационарные значения.
\item
  Билет 21: ЗБЧ и ЦПТ описывают асимптотику \(S_n\).
\item
  Билет 22: блуждание задает переходную матрицу \(p_{i,i+1}=p\),
  \(p_{i,i-1}=q\).
\item
  Билет 23: случайное блуждание - пример марковской цепи.
\end{itemize}

\subsection{11. Минимум на
удовлетворительно}\label{ux43cux438ux43dux438ux43cux443ux43c-ux43dux430-ux443ux434ux43eux432ux43bux435ux442ux432ux43eux440ux438ux442ux435ux43bux44cux43dux43e-18}

Нужно сказать:

\begin{enumerate}
\def\labelenumi{\arabic{enumi}.}
\tightlist
\item
  Приращение процесса:
\end{enumerate}

\[
S_n-S_m.
\]

\begin{enumerate}
\def\labelenumi{\arabic{enumi}.}
\setcounter{enumi}{1}
\item
  Независимые приращения: приращения на непересекающихся промежутках
  независимы.
\item
  Случайное блуждание:
\end{enumerate}

\[
S_0=0,
\qquad
S_n=\xi_1+\cdots+\xi_n,
\]

где \(\xi_k\) независимы.

\begin{enumerate}
\def\labelenumi{\arabic{enumi}.}
\setcounter{enumi}{3}
\tightlist
\item
  Если шаги iid:
\end{enumerate}

\[
\varphi_{S_n}(t)=\varphi_\xi(t)^n.
\]

\begin{enumerate}
\def\labelenumi{\arabic{enumi}.}
\setcounter{enumi}{4}
\tightlist
\item
  Главная ловушка: независимые приращения не означают независимые
  значения.
\end{enumerate}

\subsection{12. Минимум на
хорошо}\label{ux43cux438ux43dux438ux43cux443ux43c-ux43dux430-ux445ux43eux440ux43eux448ux43e-18}

Помимо минимума, нужно объяснить:

\begin{itemize}
\tightlist
\item
  что такое стационарные приращения;
\item
  почему сумма независимых шагов имеет независимые приращения;
\item
  как распределение суммы связано со сверткой;
\item
  почему характеристические функции перемножаются;
\item
  как считаются \(\mathbb ES_n\) и \(\operatorname{Var}S_n\);
\item
  почему случайное блуждание является марковской цепью;
\item
  разобрать простое симметричное блуждание.
\end{itemize}

\subsection{13. Что добавить для
отлично}\label{ux447ux442ux43e-ux434ux43eux431ux430ux432ux438ux442ux44c-ux434ux43bux44f-ux43eux442ux43bux438ux447ux43dux43e-18}

Для сильного ответа стоит добавить:

\begin{itemize}
\tightlist
\item
  формулу распределения простого блуждания:
\end{itemize}

\[
P\{S_n=x\}
=
\binom n{(n+x)/2}p^{(n+x)/2}q^{(n-x)/2};
\]

\begin{itemize}
\tightlist
\item
  условие четности \(n+x\);
\item
  различие между стационарностью процесса и стационарностью приращений;
\item
  переходные вероятности блуждания как марковской цепи;
\item
  связь с ЗБЧ и ЦПТ;
\item
  пример счетчика Бернулли и пуассоновского процесса как аналога.
\end{itemize}

\subsection{14. Устная версия ответа на 5
минут}\label{ux443ux441ux442ux43dux430ux44f-ux432ux435ux440ux441ux438ux44f-ux43eux442ux432ux435ux442ux430-ux43dux430-5-ux43cux438ux43dux443ux442-18}

\begin{enumerate}
\def\labelenumi{\arabic{enumi}.}
\tightlist
\item
  Приращение последовательности - это разность \(S_n-S_m\).
\item
  Процесс имеет независимые приращения, если приращения на
  непересекающихся промежутках независимы.
\item
  Главный пример - случайное блуждание:
\end{enumerate}

\[
S_0=0,
\qquad
S_n=\xi_1+\cdots+\xi_n.
\]

\begin{enumerate}
\def\labelenumi{\arabic{enumi}.}
\setcounter{enumi}{3}
\tightlist
\item
  Если шаги \(\xi_k\) независимы, то приращения независимы, потому что
  каждое приращение является суммой своего набора шагов.
\item
  Если шаги одинаково распределены, то приращения стационарны:
\end{enumerate}

\[
S_{n+k}-S_n\stackrel d=S_k.
\]

\begin{enumerate}
\def\labelenumi{\arabic{enumi}.}
\setcounter{enumi}{5}
\tightlist
\item
  Распределение \(S_n\) - свертка законов шагов.
\item
  Через характеристические функции:
\end{enumerate}

\[
\varphi_{S_n}(t)=\prod_{j=1}^n\varphi_{\xi_j}(t),
\]

а в iid-случае:

\[
\varphi_{S_n}(t)=\varphi_\xi(t)^n.
\]

\begin{enumerate}
\def\labelenumi{\arabic{enumi}.}
\setcounter{enumi}{7}
\tightlist
\item
  При существовании моментов:
\end{enumerate}

\[
\mathbb ES_n=n\mathbb E\xi_1,
\qquad
\operatorname{Var}S_n=n\operatorname{Var}\xi_1.
\]

\begin{enumerate}
\def\labelenumi{\arabic{enumi}.}
\setcounter{enumi}{8}
\tightlist
\item
  Для простого блуждания \(\xi_k=\pm1\), вероятности шагов \(p\) и
  \(q\).
\item
  Блуждание является марковской цепью, потому что следующий шаг
  независим от прошлого.
\item
  По значениям блуждание обычно не стационарно, но при iid-шагах имеет
  стационарные приращения.
\item
  Асимптотика описывается ЗБЧ и ЦПТ.
\end{enumerate}

\subsection{15. Если осталось 2 минуты перед
экзаменом}\label{ux435ux441ux43bux438-ux43eux441ux442ux430ux43bux43eux441ux44c-2-ux43cux438ux43dux443ux442ux44b-ux43fux435ux440ux435ux434-ux44dux43aux437ux430ux43cux435ux43dux43eux43c-18}

Запомнить:

\[
S_n-S_m
\]

\begin{itemize}
\tightlist
\item
  приращение.
\end{itemize}

Независимые приращения:

\[
S_{n_1}-S_{n_0},\ldots,S_{n_k}-S_{n_{k-1}}
\quad\text{независимы}.
\]

Блуждание:

\[
S_0=0,
\qquad
S_n=\sum_{j=1}^n\xi_j.
\]

При iid-шагах:

\[
S_{n+k}-S_n\stackrel d=S_k,
\qquad
\varphi_{S_n}(t)=\varphi_\xi(t)^n.
\]

Главные ловушки:

\begin{itemize}
\tightlist
\item
  независимые приращения не означают независимые значения;
\item
  стационарные приращения не означают стационарный процесс;
\item
  формула \(\varphi_{S_n}=\varphi_\xi^n\) требует iid-шагов.
\end{itemize}

\newpage

\section{Билет 20. Конечномерные распределения как обобщенные функции.
Дельта-функции Дирака и преобразования
Фурье}\label{ux431ux438ux43bux435ux442-20.-ux43aux43eux43dux435ux447ux43dux43eux43cux435ux440ux43dux44bux435-ux440ux430ux441ux43fux440ux435ux434ux435ux43bux435ux43dux438ux44f-ux43aux430ux43a-ux43eux431ux43eux431ux449ux435ux43dux43dux44bux435-ux444ux443ux43dux43aux446ux438ux438.-ux434ux435ux43bux44cux442ux430-ux444ux443ux43dux43aux446ux438ux438-ux434ux438ux440ux430ux43aux430-ux438-ux43fux440ux435ux43eux431ux440ux430ux437ux43eux432ux430ux43dux438ux44f-ux444ux443ux440ux44cux435}

Источники: Ширяев 1, §12; лекции ДСП.

\subsection{1. Короткий
ответ}\label{ux43aux43eux440ux43eux442ux43aux438ux439-ux43eux442ux432ux435ux442-19}

Конечномерное распределение процесса можно рассматривать не только как
меру на \(\mathbb R^n\), но и как линейный функционал на функциях. Как
мера оно действует на ограниченные борелевские или непрерывные функции;
если говорить строго на языке обобщенных функций, его ограничивают на
тестовые функции из \(C_c^\infty\) или на функции Шварца. Если

\[
\mu=P_{t_1,\ldots,t_n}
\]

\begin{itemize}
\tightlist
\item
  распределение вектора
\end{itemize}

\[
(X_{t_1},\ldots,X_{t_n}),
\]

то каждой ограниченной непрерывной функции \(f\in C_b(\mathbb R^n)\) оно
сопоставляет число

\[
\langle \mu,f\rangle
=
\int_{\mathbb R^n} f(x)\,\mu(dx)
=
\mathbb E f(X_{t_1},\ldots,X_{t_n}).
\]

В этом смысле мера действует как обобщенная функция.

Дельта-мера Дирака в точке \(a\):

\[
\delta_a(B)=\mathbf 1_B(a),
\]

а как функционал:

\[
\langle\delta_a,f\rangle=f(a).
\]

Преобразование Фурье вероятностной меры:

\[
\widehat\mu(u)
=
\int_{\mathbb R^n} e^{i\langle u,x\rangle}\,\mu(dx).
\]

Для распределения случайного вектора это его характеристическая функция:

\[
\varphi_X(u)=\mathbb E e^{i\langle u,X\rangle}.
\]

\subsection{2. Полный
ответ}\label{ux43fux43eux43bux43dux44bux439-ux43eux442ux432ux435ux442-19}

\subsubsection{2.1. Вероятностная мера и конечномерные
распределения}\label{ux432ux435ux440ux43eux44fux442ux43dux43eux441ux442ux43dux430ux44f-ux43cux435ux440ux430-ux438-ux43aux43eux43dux435ux447ux43dux43eux43cux435ux440ux43dux44bux435-ux440ux430ux441ux43fux440ux435ux434ux435ux43bux435ux43dux438ux44f}

Этот билет связывает несколько уже изученных тем: распределения
случайных векторов из Билета 06, характеристические функции из Билета
08, моменты через характеристические функции из Билета 13 и меру Дирака
из Билета 01.

Обычно распределение случайного вектора определяется как вероятностная
мера. Пусть

\[
X=(X_1,\ldots,X_n)
\]

\begin{itemize}
\tightlist
\item
  случайный вектор. Его распределение
\end{itemize}

\[
\mu=P_X
\]

на \(\mathbb R^n\) задается формулой

\[
\mu(B)=P\{X\in B\},
\qquad
B\in\mathcal B(\mathbb R^n).
\]

Для случайного процесса конечномерное распределение - это распределение
вектора

\[
(X_{t_1},\ldots,X_{t_n}).
\]

Его обычно обозначают

\[
P_{t_1,\ldots,t_n}
\]

или

\[
\mu_{t_1,\ldots,t_n}.
\]

\subsubsection{2.2. Распределение как функционал на пространстве
тестовых
функций}\label{ux440ux430ux441ux43fux440ux435ux434ux435ux43bux435ux43dux438ux435-ux43aux430ux43a-ux444ux443ux43dux43aux446ux438ux43eux43dux430ux43b-ux43dux430-ux43fux440ux43eux441ux442ux440ux430ux43dux441ux442ux432ux435-ux442ux435ux441ux442ux43eux432ux44bux445-ux444ux443ux43dux43aux446ux438ux439}

Но меру можно задавать не только значениями на множествах, а и ее
действием на функции. Если \(f\) - ограниченная борелевская или
ограниченная непрерывная функция, то можно взять интеграл

\[
\int_{\mathbb R^n} f(x)\,\mu(dx).
\]

Для вероятностных мер особенно удобен класс ограниченных непрерывных
тестовых функций

\[
C_b(\mathbb R^n).
\]

Ограниченность гарантирует, что интеграл конечен:

\[
\left|\int f\,d\mu\right|
\le
\|f\|_\infty\mu(\mathbb R^n)
=
\|f\|_\infty.
\]

Непрерывность важна для слабой сходимости распределений.

Таким образом, распределение \(\mu\) задает функционал

\[
f\mapsto \langle \mu,f\rangle,
\]

где

\[
\langle \mu,f\rangle
=
\int f(x)\,\mu(dx).
\]

\subsubsection{2.3. Свойства функционала и интерпретация как обобщенной
функции}\label{ux441ux432ux43eux439ux441ux442ux432ux430-ux444ux443ux43dux43aux446ux438ux43eux43dux430ux43bux430-ux438-ux438ux43dux442ux435ux440ux43fux440ux435ux442ux430ux446ux438ux44f-ux43aux430ux43a-ux43eux431ux43eux431ux449ux435ux43dux43dux43eux439-ux444ux443ux43dux43aux446ux438ux438}

Этот функционал линейный:

\[
\langle \mu,af+bg\rangle
=
a\langle\mu,f\rangle+b\langle\mu,g\rangle.
\]

Он положительный:

\[
f\ge0
\quad\Rightarrow\quad
\langle\mu,f\rangle\ge0.
\]

Он нормирован:

\[
\langle\mu,1\rangle=1.
\]

Именно поэтому можно говорить, что распределение рассматривается как
обобщенная функция: оно действует на тестовую функцию и возвращает
число. Здесь слово ``обобщенная'' не означает, что это обязательно
обычная функция плотности. Например, атомарное распределение не имеет
плотности относительно меры Лебега, но как функционал на тестовых
функциях оно определено без проблем.

Для случайного вектора \(X\) действие распределения на тестовую функцию
можно записать через математическое ожидание:

\[
\langle P_X,f\rangle
=
\int_{\mathbb R^n}f(x)\,P_X(dx)
=
\mathbb E f(X).
\]

Для конечномерного распределения процесса:

\[
\langle P_{t_1,\ldots,t_n},f\rangle
=
\mathbb E f(X_{t_1},\ldots,X_{t_n}).
\]

Это особенно удобно: вместо работы с вероятностями всех борелевских
множеств можно проверять интегралы от тестовых функций.

\subsubsection{2.4. Мера Дирака и способы представления
распределений}\label{ux43cux435ux440ux430-ux434ux438ux440ux430ux43aux430-ux438-ux441ux43fux43eux441ux43eux431ux44b-ux43fux440ux435ux434ux441ux442ux430ux432ux43bux435ux43dux438ux44f-ux440ux430ux441ux43fux440ux435ux434ux435ux43bux435ux43dux438ux439}

Теперь дельта-функция Дирака. Строго в теории меры это не обычная
функция, а мера, сосредоточенная в одной точке. Для \(a\in\mathbb R^n\):

\[
\delta_a(B)=
\begin{cases}
1, & a\in B,\\
0, & a\notin B.
\end{cases}
\]

Эквивалентно:

\[
\delta_a(B)=\mathbf 1_B(a).
\]

Как функционал дельта действует подстановкой:

\[
\langle\delta_a,f\rangle
=
\int f(x)\,\delta_a(dx)
=
f(a).
\]

Физическая запись

\[
\int f(x)\delta(x-a)\,dx=f(a)
\]

понимается именно в этом смысле. Строго \(\delta(x-a)\) не является
обычной функцией плотности по Лебегу.

Если случайная величина постоянна:

\[
X=a
\quad\text{п.н.},
\]

то ее распределение равно

\[
P_X=\delta_a.
\]

Если распределение дискретное:

\[
P\{X=x_k\}=p_k,
\qquad
\sum_k p_k=1,
\]

то его можно записать как сумму дельта-мер:

\[
P_X=\sum_k p_k\delta_{x_k}.
\]

Действие на тестовую функцию:

\[
\langle P_X,f\rangle
=
\sum_k p_k f(x_k).
\]

Если распределение абсолютно непрерывно с плотностью \(p(x)\), то

\[
\mu(dx)=p(x)\,dx
\]

и

\[
\langle\mu,f\rangle
=
\int_{\mathbb R^n}f(x)p(x)\,dx.
\]

Так в одной записи объединяются дискретные, непрерывные и смешанные
распределения.

\subsubsection{2.5. Слабая сходимость
распределений}\label{ux441ux43bux430ux431ux430ux44f-ux441ux445ux43eux434ux438ux43cux43eux441ux442ux44c-ux440ux430ux441ux43fux440ux435ux434ux435ux43bux435ux43dux438ux439}

Следующая важная идея - слабая сходимость. Последовательность
вероятностных мер \(\mu_k\) слабо сходится к \(\mu\), если для всех
ограниченных непрерывных функций \(f\):

\[
\int f\,d\mu_k\to\int f\,d\mu.
\]

Пишут:

\[
\mu_k\Rightarrow\mu.
\]

Для случайных величин это означает сходимость по распределению:

\[
X_k\Rightarrow X.
\]

Иными словами, слабая сходимость - это сходимость распределений как
функционалов на тестовых функциях.

\subsubsection{2.6. Преобразование Фурье меры и характеристические
функции}\label{ux43fux440ux435ux43eux431ux440ux430ux437ux43eux432ux430ux43dux438ux435-ux444ux443ux440ux44cux435-ux43cux435ux440ux44b-ux438-ux445ux430ux440ux430ux43aux442ux435ux440ux438ux441ux442ux438ux447ux435ux441ux43aux438ux435-ux444ux443ux43dux43aux446ux438ux438}

Теперь преобразование Фурье меры. Для конечной меры \(\mu\) на
\(\mathbb R^n\) ее преобразование Фурье определяется формулой

\[
\widehat\mu(u)
=
\int_{\mathbb R^n}e^{i\langle u,x\rangle}\,\mu(dx),
\qquad
u\in\mathbb R^n.
\]

Так как

\[
|e^{i\langle u,x\rangle}|=1,
\]

интеграл всегда существует для вероятностной меры.

Если \(\mu=P_X\) - распределение случайного вектора \(X\), то

\[
\widehat\mu(u)
=
\mathbb E e^{i\langle u,X\rangle}
=
\varphi_X(u).
\]

То есть характеристическая функция - это просто преобразование Фурье
распределения.

Для дельта-мер:

\[
\widehat{\delta_a}(u)
=
\int e^{i\langle u,x\rangle}\delta_a(dx)
=
e^{i\langle u,a\rangle}.
\]

Для дискретного распределения:

\[
\mu=\sum_k p_k\delta_{x_k}
\]

получаем

\[
\widehat\mu(u)=\sum_k p_k e^{i\langle u,x_k\rangle}.
\]

Для распределения с плотностью:

\[
\widehat\mu(u)=\int_{\mathbb R^n}e^{i\langle u,x\rangle}p(x)\,dx.
\]

\subsubsection{2.7. Свертка и теорема
единственности}\label{ux441ux432ux435ux440ux442ux43aux430-ux438-ux442ux435ux43eux440ux435ux43cux430-ux435ux434ux438ux43dux441ux442ux432ux435ux43dux43dux43eux441ux442ux438}

Преобразование Фурье удобно тем, что свертка мер превращается в
произведение. Если \(X\) и \(Y\) независимы, то распределение суммы
\(X+Y\) равно свертке распределений:

\[
P_{X+Y}=P_X*P_Y,
\]

а характеристические функции перемножаются:

\[
\varphi_{X+Y}(u)=\varphi_X(u)\varphi_Y(u).
\]

Это использовалось в Билете 08 и Билете 19.

Еще одно ключевое свойство: характеристическая функция однозначно задает
распределение. Если

\[
\widehat\mu(u)=\widehat\nu(u)
\]

для всех \(u\), то

\[
\mu=\nu.
\]

Поэтому преобразование Фурье - не просто удобный числовой образ, а
полноценный способ кодировать распределение.

\subsection{3. Основные
определения}\label{ux43eux441ux43dux43eux432ux43dux44bux435-ux43eux43fux440ux435ux434ux435ux43bux435ux43dux438ux44f-19}

\subsubsection{Конечномерное распределение
процесса}\label{ux43aux43eux43dux435ux447ux43dux43eux43cux435ux440ux43dux43eux435-ux440ux430ux441ux43fux440ux435ux434ux435ux43bux435ux43dux438ux435-ux43fux440ux43eux446ux435ux441ux441ux430-3}

Для процесса \((X_t)\) и набора \(t_1,\ldots,t_n\) это распределение
случайного вектора

\[
(X_{t_1},\ldots,X_{t_n})
\]

на \(\mathbb R^n\):

\[
P_{t_1,\ldots,t_n}(B)
=
P\{(X_{t_1},\ldots,X_{t_n})\in B\}.
\]

\subsubsection{Тестовая
функция}\label{ux442ux435ux441ux442ux43eux432ux430ux44f-ux444ux443ux43dux43aux446ux438ux44f}

В этом билете удобно брать

\[
f\in C_b(\mathbb R^n),
\]

то есть \(f\) непрерывна и ограничена.

\subsubsection{Мера как
функционал}\label{ux43cux435ux440ux430-ux43aux430ux43a-ux444ux443ux43dux43aux446ux438ux43eux43dux430ux43b}

Вероятностная мера \(\mu\) действует на тестовую функцию так:

\[
\langle\mu,f\rangle=\int f\,d\mu.
\]

\subsubsection{Дельта-мера
Дирака}\label{ux434ux435ux43bux44cux442ux430-ux43cux435ux440ux430-ux434ux438ux440ux430ux43aux430}

Мера в точке \(a\):

\[
\delta_a(B)=\mathbf 1_B(a).
\]

Как функционал:

\[
\langle\delta_a,f\rangle=f(a).
\]

\subsubsection{Дискретное распределение через
дельта-меры}\label{ux434ux438ux441ux43aux440ux435ux442ux43dux43eux435-ux440ux430ux441ux43fux440ux435ux434ux435ux43bux435ux43dux438ux435-ux447ux435ux440ux435ux437-ux434ux435ux43bux44cux442ux430-ux43cux435ux440ux44b}

\[
\mu=\sum_k p_k\delta_{x_k}.
\]

\subsubsection{Абсолютно непрерывное
распределение}\label{ux430ux431ux441ux43eux43bux44eux442ux43dux43e-ux43dux435ux43fux440ux435ux440ux44bux432ux43dux43eux435-ux440ux430ux441ux43fux440ux435ux434ux435ux43bux435ux43dux438ux435}

Если есть плотность \(p\), то

\[
\mu(dx)=p(x)\,dx.
\]

\subsubsection{Слабая сходимость
мер}\label{ux441ux43bux430ux431ux430ux44f-ux441ux445ux43eux434ux438ux43cux43eux441ux442ux44c-ux43cux435ux440}

\[
\mu_n\Rightarrow\mu
\quad\Longleftrightarrow\quad
\int f\,d\mu_n\to\int f\,d\mu
\]

для всех \(f\in C_b\).

\subsubsection{Преобразование Фурье
меры}\label{ux43fux440ux435ux43eux431ux440ux430ux437ux43eux432ux430ux43dux438ux435-ux444ux443ux440ux44cux435-ux43cux435ux440ux44b}

\[
\widehat\mu(u)=\int e^{i\langle u,x\rangle}\,\mu(dx).
\]

\subsubsection{Характеристическая
функция}\label{ux445ux430ux440ux430ux43aux442ux435ux440ux438ux441ux442ux438ux447ux435ux441ux43aux430ux44f-ux444ux443ux43dux43aux446ux438ux44f}

Для \(X\sim\mu\):

\[
\varphi_X(u)=\mathbb E e^{i\langle u,X\rangle}=\widehat\mu(u).
\]

\subsection{4. Основные теоремы и
свойства}\label{ux43eux441ux43dux43eux432ux43dux44bux435-ux442ux435ux43eux440ux435ux43cux44b-ux438-ux441ux432ux43eux439ux441ux442ux432ux430-19}

\begin{itemize}
\tightlist
\item
  Вероятностная мера задает положительный нормированный линейный
  функционал:
\end{itemize}

\[
f\mapsto \int f\,d\mu.
\]

\begin{itemize}
\tightlist
\item
  Для конечномерного распределения:
\end{itemize}

\[
\int f\,dP_{t_1,\ldots,t_n}
=
\mathbb E f(X_{t_1},\ldots,X_{t_n}).
\]

\begin{itemize}
\tightlist
\item
  Дельта-мера действует как подстановка:
\end{itemize}

\[
\int f\,d\delta_a=f(a).
\]

\begin{itemize}
\tightlist
\item
  Дискретное распределение записывается как сумма атомов:
\end{itemize}

\[
\mu=\sum_kp_k\delta_{x_k}.
\]

\begin{itemize}
\tightlist
\item
  Характеристическая функция - преобразование Фурье распределения:
\end{itemize}

\[
\varphi_X(u)=\widehat{P_X}(u).
\]

\begin{itemize}
\tightlist
\item
  Преобразование Фурье дельта-меры:
\end{itemize}

\[
\widehat{\delta_a}(u)=e^{i\langle u,a\rangle}.
\]

\begin{itemize}
\tightlist
\item
  Свертка переходит в произведение:
\end{itemize}

\[
\widehat{\mu*\nu}(u)=\widehat\mu(u)\widehat\nu(u).
\]

\begin{itemize}
\tightlist
\item
  Характеристическая функция однозначно задает распределение.
\item
  Слабая сходимость распределений проверяется на ограниченных
  непрерывных тестовых функциях.
\end{itemize}

\subsection{5. Примеры}\label{ux43fux440ux438ux43cux435ux440ux44b-19}

\subsubsection{Постоянная случайная
величина}\label{ux43fux43eux441ux442ux43eux44fux43dux43dux430ux44f-ux441ux43bux443ux447ux430ux439ux43dux430ux44f-ux432ux435ux43bux438ux447ux438ux43dux430}

Если

\[
X=a
\quad\text{п.н.},
\]

то

\[
P_X=\delta_a.
\]

Для тестовой функции:

\[
\mathbb E f(X)=f(a).
\]

Характеристическая функция:

\[
\varphi_X(t)=e^{ita}.
\]

\subsubsection{Дискретное
распределение}\label{ux434ux438ux441ux43aux440ux435ux442ux43dux43eux435-ux440ux430ux441ux43fux440ux435ux434ux435ux43bux435ux43dux438ux435-1}

Пусть

\[
P\{X=x_k\}=p_k.
\]

Тогда

\[
P_X=\sum_kp_k\delta_{x_k}.
\]

Действие на тестовую функцию:

\[
\langle P_X,f\rangle=\sum_kp_kf(x_k).
\]

Характеристическая функция:

\[
\varphi_X(t)=\sum_kp_ke^{itx_k}.
\]

\subsubsection{Бернуллиевская
величина}\label{ux431ux435ux440ux43dux443ux43bux43bux438ux435ux432ux441ux43aux430ux44f-ux432ux435ux43bux438ux447ux438ux43dux430}

Если

\[
P\{X=1\}=p,
\qquad
P\{X=0\}=1-p,
\]

то

\[
P_X=(1-p)\delta_0+p\delta_1.
\]

Характеристическая функция:

\[
\varphi_X(t)=1-p+pe^{it}.
\]

\subsubsection{Распределение с
плотностью}\label{ux440ux430ux441ux43fux440ux435ux434ux435ux43bux435ux43dux438ux435-ux441-ux43fux43bux43eux442ux43dux43eux441ux442ux44cux44e}

Если

\[
\mu(dx)=p(x)\,dx,
\]

то

\[
\langle\mu,f\rangle=\int f(x)p(x)\,dx.
\]

Преобразование Фурье:

\[
\widehat\mu(t)=\int e^{itx}p(x)\,dx.
\]

\subsubsection{Конечномерное распределение
процесса}\label{ux43aux43eux43dux435ux447ux43dux43eux43cux435ux440ux43dux43eux435-ux440ux430ux441ux43fux440ux435ux434ux435ux43bux435ux43dux438ux435-ux43fux440ux43eux446ux435ux441ux441ux430-4}

Для процесса \((X_n)\) двумерное распределение пары \((X_0,X_1)\)
действует на функцию \(f(x,y)\) так:

\[
\langle P_{0,1},f\rangle
=
\mathbb E f(X_0,X_1).
\]

Если процесс детерминированный и

\[
X_0=a,\qquad X_1=b
\]

почти наверное, то

\[
P_{0,1}=\delta_{(a,b)}.
\]

\subsection{6. Схемы доказательств /
обоснований}\label{ux441ux445ux435ux43cux44b-ux434ux43eux43aux430ux437ux430ux442ux435ux43bux44cux441ux442ux432-ux43eux431ux43eux441ux43dux43eux432ux430ux43dux438ux439-12}

\subsubsection{Почему мера задает линейный положительный
функционал}\label{ux43fux43eux447ux435ux43cux443-ux43cux435ux440ux430-ux437ux430ux434ux430ux435ux442-ux43bux438ux43dux435ux439ux43dux44bux439-ux43fux43eux43bux43eux436ux438ux442ux435ux43bux44cux43dux44bux439-ux444ux443ux43dux43aux446ux438ux43eux43dux430ux43b}

Для \(f,g\in C_b\) и чисел \(a,b\):

\[
\int(af+bg)\,d\mu
=
a\int f\,d\mu+b\int g\,d\mu.
\]

Если \(f\ge0\), то

\[
\int f\,d\mu\ge0.
\]

Если \(\mu\) - вероятность, то

\[
\int 1\,d\mu=\mu(\mathbb R^n)=1.
\]

\subsubsection{Почему дельта действует как
подстановка}\label{ux43fux43eux447ux435ux43cux443-ux434ux435ux43bux44cux442ux430-ux434ux435ux439ux441ux442ux432ux443ux435ux442-ux43aux430ux43a-ux43fux43eux434ux441ux442ux430ux43dux43eux432ux43aux430}

Для простой функции это сразу видно:

\[
\int \sum_j c_j\mathbf 1_{A_j}(x)\,\delta_a(dx)
=
\sum_j c_j\mathbf 1_{A_j}(a).
\]

Правая часть равна значению простой функции в точке \(a\). Переходом к
пределу получаем для ограниченной измеримой функции:

\[
\int f\,d\delta_a=f(a).
\]

\subsubsection{Почему характеристическая функция -
Фурье-образ}\label{ux43fux43eux447ux435ux43cux443-ux445ux430ux440ux430ux43aux442ux435ux440ux438ux441ux442ux438ux447ux435ux441ux43aux430ux44f-ux444ux443ux43dux43aux446ux438ux44f---ux444ux443ux440ux44cux435-ux43eux431ux440ux430ux437}

Если \(X\sim\mu\), то по формуле замены переменной:

\[
\mathbb E e^{i\langle u,X\rangle}
=
\int e^{i\langle u,x\rangle}\,\mu(dx).
\]

Левая часть - характеристическая функция, правая - преобразование Фурье
меры.

\subsubsection{Почему при независимости характеристики
перемножаются}\label{ux43fux43eux447ux435ux43cux443-ux43fux440ux438-ux43dux435ux437ux430ux432ux438ux441ux438ux43cux43eux441ux442ux438-ux445ux430ux440ux430ux43aux442ux435ux440ux438ux441ux442ux438ux43aux438-ux43fux435ux440ux435ux43cux43dux43eux436ux430ux44eux442ux441ux44f}

Если \(X\) и \(Y\) независимы, то

\[
\varphi_{X+Y}(u)
=
\mathbb E e^{i\langle u,X+Y\rangle}
=
\mathbb E\left(e^{i\langle u,X\rangle}e^{i\langle u,Y\rangle}\right).
\]

По независимости:

\[
\varphi_{X+Y}(u)
=
\mathbb E e^{i\langle u,X\rangle}
\mathbb E e^{i\langle u,Y\rangle}
=
\varphi_X(u)\varphi_Y(u).
\]

\subsection{7. Что написать на
доске}\label{ux447ux442ux43e-ux43dux430ux43fux438ux441ux430ux442ux44c-ux43dux430-ux434ux43eux441ux43aux435-19}

\begin{enumerate}
\def\labelenumi{\arabic{enumi}.}
\tightlist
\item
  Мера как функционал:
\end{enumerate}

\[
\langle\mu,f\rangle=\int f(x)\,\mu(dx).
\]

\begin{enumerate}
\def\labelenumi{\arabic{enumi}.}
\setcounter{enumi}{1}
\tightlist
\item
  Конечномерное распределение:
\end{enumerate}

\[
\langle P_{t_1,\ldots,t_n},f\rangle
=
\mathbb E f(X_{t_1},\ldots,X_{t_n}).
\]

\begin{enumerate}
\def\labelenumi{\arabic{enumi}.}
\setcounter{enumi}{2}
\tightlist
\item
  Дельта Дирака:
\end{enumerate}

\[
\delta_a(B)=\mathbf 1_B(a),
\qquad
\langle\delta_a,f\rangle=f(a).
\]

\begin{enumerate}
\def\labelenumi{\arabic{enumi}.}
\setcounter{enumi}{3}
\tightlist
\item
  Дискретный закон:
\end{enumerate}

\[
\mu=\sum_kp_k\delta_{x_k}.
\]

\begin{enumerate}
\def\labelenumi{\arabic{enumi}.}
\setcounter{enumi}{4}
\tightlist
\item
  Фурье-образ меры:
\end{enumerate}

\[
\widehat\mu(u)=\int e^{i\langle u,x\rangle}\,\mu(dx).
\]

\begin{enumerate}
\def\labelenumi{\arabic{enumi}.}
\setcounter{enumi}{5}
\tightlist
\item
  Характеристическая функция:
\end{enumerate}

\[
\varphi_X(u)=\widehat{P_X}(u).
\]

\begin{enumerate}
\def\labelenumi{\arabic{enumi}.}
\setcounter{enumi}{6}
\tightlist
\item
  Слабая сходимость:
\end{enumerate}

\[
\mu_n\Rightarrow\mu
\Longleftrightarrow
\int f\,d\mu_n\to\int f\,d\mu
\quad(f\in C_b).
\]

\subsection{8. Типичные дополнительные
вопросы}\label{ux442ux438ux43fux438ux447ux43dux44bux435-ux434ux43eux43fux43eux43bux43dux438ux442ux435ux43bux44cux43dux44bux435-ux432ux43eux43fux440ux43eux441ux44b-19}

\textbf{Что значит рассматривать распределение как обобщенную функцию?}

Это значит рассматривать меру \(\mu\) по ее действию на тестовые
функции:

\[
f\mapsto\langle\mu,f\rangle=\int f\,d\mu.
\]

Даже если у распределения нет плотности, такой функционал определен.

\textbf{Почему берут ограниченные непрерывные тестовые функции?}

Ограниченность гарантирует конечность интеграла по вероятностной мере.
Непрерывность нужна для слабой сходимости: распределения считаются
близкими, если интегралы от всех \(f\in C_b\) близки.

\textbf{Чем дельта Дирака отличается от обычной функции?}

Строго \(\delta_a\) - это мера или функционал, а не обычная функция. Ее
смысл:

\[
\int f\,d\delta_a=f(a).
\]

Нельзя обращаться с ней как с обычной плотностью относительно меры
Лебега.

\textbf{Как записать дискретное распределение через дельта-меры?}

Если

\[
P\{X=x_k\}=p_k,
\]

то

\[
P_X=\sum_kp_k\delta_{x_k}.
\]

\textbf{Что такое преобразование Фурье меры?}

Это функция

\[
\widehat\mu(u)=\int e^{i\langle u,x\rangle}\,\mu(dx).
\]

Для вероятностной меры она всегда определена, потому что модуль
экспоненты равен \(1\).

\textbf{Почему характеристическая функция - это Фурье-образ
распределения?}

Если \(X\sim\mu\), то

\[
\varphi_X(u)
=
\mathbb E e^{i\langle u,X\rangle}
=
\int e^{i\langle u,x\rangle}\,\mu(dx)
=
\widehat\mu(u).
\]

\textbf{Что такое слабая сходимость в этом языке?}

Это сходимость мер как функционалов:

\[
\mu_n\Rightarrow\mu
\]

если для всех \(f\in C_b\):

\[
\langle\mu_n,f\rangle\to\langle\mu,f\rangle.
\]

\textbf{Какой Фурье-образ у дельта-меры?}

\[
\widehat{\delta_a}(u)=e^{i\langle u,a\rangle}.
\]

\textbf{Как связаны свертка и преобразование Фурье?}

Свертка мер переходит в произведение Фурье-образов:

\[
\widehat{\mu*\nu}(u)=\widehat\mu(u)\widehat\nu(u).
\]

Вероятностно это означает, что характеристическая функция суммы
независимых величин равна произведению характеристических функций.

\textbf{Какая главная ошибка в этом билете?}

Путать дельта-функцию с обычной функцией и пытаться использовать
плотность там, где распределение атомарное.

\subsection{9. Типичные
ошибки}\label{ux442ux438ux43fux438ux447ux43dux44bux435-ux43eux448ux438ux431ux43aux438-19}

\begin{itemize}
\tightlist
\item
  Путать дельта-меру с обычной функцией.
\item
  Писать плотность для атомарного распределения без оговорки о
  дельта-мерах.
\item
  Забывать, что \(\delta_a\) действует на функцию как подстановка.
\item
  Путать функцию распределения \(F(x)\) и распределение как меру
  \(\mu\).
\item
  Путать преобразование Фурье плотности и преобразование Фурье меры:
  второе работает и без плотности.
\item
  Забывать знак в соглашении Фурье; в вероятности обычно используют
  \(e^{i\langle u,x\rangle}\).
\item
  Считать, что слабая сходимость проверяется по индикаторам всех
  множеств; стандартная формулировка здесь через \(C_b\).
\end{itemize}

\subsection{10. Связи с другими
билетами}\label{ux441ux432ux44fux437ux438-ux441-ux434ux440ux443ux433ux438ux43cux438-ux431ux438ux43bux435ux442ux430ux43cux438-19}

\begin{itemize}
\tightlist
\item
  Билет 01: мера Дирака как пример меры.
\item
  Билет 06: распределения случайных величин и конечномерные
  распределения.
\item
  Билет 08: характеристические функции как преобразования Фурье
  распределений.
\item
  Билет 09: независимость и произведение характеристических функций.
\item
  Билет 13: моменты и ковариации через производные характеристической
  функции.
\item
  Билет 15: согласованные конечномерные распределения задают процесс.
\item
  Билет 18: гауссовские распределения удобно задаются
  характеристическими функциями.
\item
  Билет 19: для сумм независимых шагов характеристические функции
  перемножаются.
\end{itemize}

\subsection{11. Минимум на
удовлетворительно}\label{ux43cux438ux43dux438ux43cux443ux43c-ux43dux430-ux443ux434ux43eux432ux43bux435ux442ux432ux43eux440ux438ux442ux435ux43bux44cux43dux43e-19}

Нужно сказать:

\begin{enumerate}
\def\labelenumi{\arabic{enumi}.}
\tightlist
\item
  Распределение \(\mu\) можно рассматривать как функционал:
\end{enumerate}

\[
\langle\mu,f\rangle=\int f\,d\mu.
\]

\begin{enumerate}
\def\labelenumi{\arabic{enumi}.}
\setcounter{enumi}{1}
\tightlist
\item
  Для конечномерного распределения процесса:
\end{enumerate}

\[
\langle P_{t_1,\ldots,t_n},f\rangle
=
\mathbb E f(X_{t_1},\ldots,X_{t_n}).
\]

\begin{enumerate}
\def\labelenumi{\arabic{enumi}.}
\setcounter{enumi}{2}
\tightlist
\item
  Дельта Дирака:
\end{enumerate}

\[
\langle\delta_a,f\rangle=f(a).
\]

\begin{enumerate}
\def\labelenumi{\arabic{enumi}.}
\setcounter{enumi}{3}
\tightlist
\item
  Преобразование Фурье меры:
\end{enumerate}

\[
\widehat\mu(u)=\int e^{i\langle u,x\rangle}\,\mu(dx).
\]

\begin{enumerate}
\def\labelenumi{\arabic{enumi}.}
\setcounter{enumi}{4}
\tightlist
\item
  Для распределения случайного вектора это характеристическая функция.
\end{enumerate}

\subsection{12. Минимум на
хорошо}\label{ux43cux438ux43dux438ux43cux443ux43c-ux43dux430-ux445ux43eux440ux43eux448ux43e-19}

Помимо минимума, нужно объяснить:

\begin{itemize}
\tightlist
\item
  почему ограниченность тестовых функций гарантирует конечность
  интеграла;
\item
  почему дельта Дирака не является обычной плотностью;
\item
  как записать дискретное распределение через сумму дельта-мер;
\item
  как записать абсолютно непрерывное распределение через плотность;
\item
  что слабая сходимость означает сходимость интегралов по \(C_b\);
\item
  почему характеристическая функция однозначно задает распределение.
\end{itemize}

\subsection{13. Что добавить для
отлично}\label{ux447ux442ux43e-ux434ux43eux431ux430ux432ux438ux442ux44c-ux434ux43bux44f-ux43eux442ux43bux438ux447ux43dux43e-19}

Для сильного ответа стоит добавить:

\begin{itemize}
\tightlist
\item
  связь с конечномерными распределениями процесса;
\item
  положительность и нормированность функционала
  \(\mu:f\mapsto\int f\,d\mu\);
\item
  примеры: константа, Бернулли, плотность;
\item
  преобразование Фурье дельта-меры;
\item
  свертку и произведение характеристических функций;
\item
  предупреждение про соглашение знака в преобразовании Фурье;
\item
  связь со слабой сходимостью и ЦПТ.
\end{itemize}

\subsection{14. Устная версия ответа на 5
минут}\label{ux443ux441ux442ux43dux430ux44f-ux432ux435ux440ux441ux438ux44f-ux43eux442ux432ux435ux442ux430-ux43dux430-5-ux43cux438ux43dux443ux442-19}

\begin{enumerate}
\def\labelenumi{\arabic{enumi}.}
\tightlist
\item
  Конечномерное распределение процесса - это мера на \(\mathbb R^n\).
\item
  Эту меру можно рассматривать как функционал на тестовых функциях:
\end{enumerate}

\[
\langle\mu,f\rangle=\int f\,d\mu.
\]

\begin{enumerate}
\def\labelenumi{\arabic{enumi}.}
\setcounter{enumi}{2}
\tightlist
\item
  Для процесса:
\end{enumerate}

\[
\langle P_{t_1,\ldots,t_n},f\rangle
=
\mathbb E f(X_{t_1},\ldots,X_{t_n}).
\]

\begin{enumerate}
\def\labelenumi{\arabic{enumi}.}
\setcounter{enumi}{3}
\tightlist
\item
  Такой функционал линейный, положительный и нормированный.
\item
  Дельта Дирака в точке \(a\):
\end{enumerate}

\[
\delta_a(B)=\mathbf 1_B(a),
\qquad
\langle\delta_a,f\rangle=f(a).
\]

\begin{enumerate}
\def\labelenumi{\arabic{enumi}.}
\setcounter{enumi}{5}
\tightlist
\item
  Дискретное распределение записывается как сумма дельт:
\end{enumerate}

\[
\mu=\sum_kp_k\delta_{x_k}.
\]

\begin{enumerate}
\def\labelenumi{\arabic{enumi}.}
\setcounter{enumi}{6}
\tightlist
\item
  Если есть плотность \(p\), то
\end{enumerate}

\[
\langle\mu,f\rangle=\int f(x)p(x)\,dx.
\]

\begin{enumerate}
\def\labelenumi{\arabic{enumi}.}
\setcounter{enumi}{7}
\tightlist
\item
  Преобразование Фурье меры:
\end{enumerate}

\[
\widehat\mu(u)=\int e^{i\langle u,x\rangle}\,\mu(dx).
\]

\begin{enumerate}
\def\labelenumi{\arabic{enumi}.}
\setcounter{enumi}{8}
\tightlist
\item
  Для вероятностной меры это характеристическая функция:
\end{enumerate}

\[
\varphi_X(u)=\widehat{P_X}(u).
\]

\begin{enumerate}
\def\labelenumi{\arabic{enumi}.}
\setcounter{enumi}{9}
\tightlist
\item
  У дельта-меры:
\end{enumerate}

\[
\widehat{\delta_a}(u)=e^{i\langle u,a\rangle}.
\]

\begin{enumerate}
\def\labelenumi{\arabic{enumi}.}
\setcounter{enumi}{10}
\tightlist
\item
  Слабая сходимость означает сходимость интегралов от всех ограниченных
  непрерывных тестовых функций.
\item
  Главная ошибка - считать \(\delta\) обычной функцией и пытаться всюду
  писать плотности.
\end{enumerate}

\subsection{15. Если осталось 2 минуты перед
экзаменом}\label{ux435ux441ux43bux438-ux43eux441ux442ux430ux43bux43eux441ux44c-2-ux43cux438ux43dux443ux442ux44b-ux43fux435ux440ux435ux434-ux44dux43aux437ux430ux43cux435ux43dux43eux43c-19}

Запомнить три формулы.

Мера как функционал:

\[
\langle\mu,f\rangle=\int f\,d\mu.
\]

Дельта Дирака:

\[
\langle\delta_a,f\rangle=f(a).
\]

Фурье-образ меры:

\[
\widehat\mu(u)=\int e^{i\langle u,x\rangle}\,\mu(dx).
\]

Для распределения \(X\):

\[
\widehat{P_X}(u)=\varphi_X(u).
\]

Главная ловушка: \(\delta_a\) - мера/функционал, а не обычная функция
плотности.

\newpage

\section{Билет 21. Закон больших чисел как частный случай центральной
предельной теоремы и как обоснование
аксиоматики}\label{ux431ux438ux43bux435ux442-21.-ux437ux430ux43aux43eux43d-ux431ux43eux43bux44cux448ux438ux445-ux447ux438ux441ux435ux43b-ux43aux430ux43a-ux447ux430ux441ux442ux43dux44bux439-ux441ux43bux443ux447ux430ux439-ux446ux435ux43dux442ux440ux430ux43bux44cux43dux43eux439-ux43fux440ux435ux434ux435ux43bux44cux43dux43eux439-ux442ux435ux43eux440ux435ux43cux44b-ux438-ux43aux430ux43a-ux43eux431ux43eux441ux43dux43eux432ux430ux43dux438ux435-ux430ux43aux441ux438ux43eux43cux430ux442ux438ux43aux438}

Источники: Ширяев 1, гл. I, §5-6; гл. III, §3; Ширяев 2, исторические
комментарии об аксиоматике Колмогорова.

\subsection{1. Короткий
ответ}\label{ux43aux43eux440ux43eux442ux43aux438ux439-ux43eux442ux432ux435ux442-20}

Закон больших чисел говорит, что среднее большого числа независимых
одинаково распределенных наблюдений сходится к математическому ожиданию:

\[
\overline X_n=\frac{X_1+\cdots+X_n}{n}\xrightarrow{P}m,
\qquad
m=\mathbb EX_1.
\]

Центральная предельная теорема уточняет масштаб случайных отклонений.
Если

\[
\mathbb EX_1=m,
\qquad
\operatorname{Var}X_1=\sigma^2\in(0,\infty),
\]

то

\[
\frac{X_1+\cdots+X_n-nm}{\sigma\sqrt n}
\Rightarrow
N(0,1).
\]

Из этой формулы видно, что

\[
X_1+\cdots+X_n-nm
\]

имеет типичный порядок \(\sqrt n\), поэтому после деления на \(n\)
отклонение среднего от \(m\) исчезает:

\[
\overline X_n-m
=
\frac{\sigma}{\sqrt n}
\frac{X_1+\cdots+X_n-nm}{\sigma\sqrt n}
\xrightarrow{P}0.
\]

Для индикаторов событий \(X_k=\mathbf 1_A^{(k)}\) это дает частотную
формулу

\[
\frac{N_n(A)}{n}\xrightarrow{P}P(A).
\]

Поэтому вероятность в аксиоматике Колмогорова не определяется как
частота, но закон больших чисел объясняет, почему вероятностная мера
согласуется с наблюдаемыми частотами в длинных сериях независимых
испытаний.

\subsection{2. Полный
ответ}\label{ux43fux43eux43bux43dux44bux439-ux43eux442ux432ux435ux442-20}

\subsubsection{2.1. Введение и основные
идеи}\label{ux432ux432ux435ux434ux435ux43dux438ux435-ux438-ux43eux441ux43dux43eux432ux43dux44bux435-ux438ux434ux435ux438}

В этом билете нужно связать три идеи:

\begin{itemize}
\tightlist
\item
  закон больших чисел;
\item
  центральную предельную теорему;
\item
  частотное обоснование вероятности как меры.
\end{itemize}

\subsubsection{2.2. Закон больших
чисел}\label{ux437ux430ux43aux43eux43d-ux431ux43eux43bux44cux448ux438ux445-ux447ux438ux441ux435ux43b}

Пусть

\[
X_1,X_2,\ldots
\]

\begin{itemize}
\tightlist
\item
  независимые одинаково распределенные случайные величины. Обозначим
\end{itemize}

\[
S_n=X_1+\cdots+X_n,
\qquad
\overline X_n=\frac{S_n}{n}.
\]

Если существует математическое ожидание

\[
\mathbb EX_1=m,
\]

то закон больших чисел утверждает, что выборочное среднее
\(\overline X_n\) при больших \(n\) близко к \(m\). В слабой форме это
записывается так:

\[
\overline X_n\xrightarrow{P}m,
\]

то есть для каждого \(\varepsilon>0\)

\[
P\{|\overline X_n-m|>\varepsilon\}\to0.
\]

Интуитивно: отдельное наблюдение может сильно колебаться, но при
усреднении независимые случайные отклонения взаимно компенсируются.

\subsubsection{2.3. Центральная предельная
теорема}\label{ux446ux435ux43dux442ux440ux430ux43bux44cux43dux430ux44f-ux43fux440ux435ux434ux435ux43bux44cux43dux430ux44f-ux442ux435ux43eux440ux435ux43cux430}

Центральная предельная теорема дает более точное утверждение. Если,
кроме среднего, существует конечная ненулевая дисперсия

\[
\operatorname{Var}X_1=\sigma^2\in(0,\infty),
\]

то

\[
\frac{S_n-nm}{\sigma\sqrt n}
\Rightarrow
N(0,1).
\]

Здесь \(\Rightarrow\) означает сходимость по распределению.
Эквивалентно, для всех точек непрерывности функции распределения
стандартного нормального закона:

\[
P\left\{
\frac{S_n-nm}{\sigma\sqrt n}\le x
\right\}
\to
\Phi(x),
\]

где

\[
\Phi(x)=\frac{1}{\sqrt{2\pi}}\int_{-\infty}^x e^{-u^2/2}\,du.
\]

\subsubsection{2.4. Вывод ЗБЧ из
ЦПТ}\label{ux432ux44bux432ux43eux434-ux437ux431ux447-ux438ux437-ux446ux43fux442}

Теперь покажем, почему ЗБЧ можно получить из ЦПТ. Обозначим

\[
Z_n=
\frac{S_n-nm}{\sigma\sqrt n}.
\]

По ЦПТ

\[
Z_n\Rightarrow Z,
\qquad
Z\sim N(0,1).
\]

Последовательность, сходящаяся по распределению к настоящей случайной
величине, ограничена по вероятности: ее значения не уходят на
бесконечность с положительной вероятностью. Поэтому

\[
Z_n=O_P(1).
\]

Тогда

\[
\overline X_n-m
=
\frac{S_n-nm}{n}
=
\frac{\sigma}{\sqrt n}Z_n
\xrightarrow{P}0.
\]

Более подробно: для любого \(\varepsilon>0\)

\[
P\{|\overline X_n-m|>\varepsilon\}
=
P\left\{|Z_n|>\frac{\varepsilon\sqrt n}{\sigma}\right\}.
\]

Правая граница уходит в бесконечность, а \(Z_n\) остается ограниченной
по вероятности, значит эта вероятность стремится к нулю. Это и есть
слабый закон больших чисел.

\subsubsection{2.5. Условия применимости
теорем}\label{ux443ux441ux43bux43eux432ux438ux44f-ux43fux440ux438ux43cux435ux43dux438ux43cux43eux441ux442ux438-ux442ux435ux43eux440ux435ux43c}

Важно сказать на экзамене: не всякий закон больших чисел является
буквально частным случаем ЦПТ. ЦПТ требует конечной ненулевой дисперсии,
а слабый закон больших чисел может быть верен при более слабых условиях,
например для iid-величин с конечным первым моментом. Поэтому корректная
фраза такая: \textbf{при условиях ЦПТ закон больших чисел следует из нее
как более грубое утверждение}.

\subsubsection{2.6. Схема Бернулли и частотное
обоснование}\label{ux441ux445ux435ux43cux430-ux431ux435ux440ux43dux443ux43bux43bux438-ux438-ux447ux430ux441ux442ux43eux442ux43dux43eux435-ux43eux431ux43eux441ux43dux43eux432ux430ux43dux438ux435}

Исторически и методически главный пример - схема Бернулли. Пусть
проводится последовательность независимых испытаний, и событие \(A\)
происходит в каждом испытании с вероятностью

\[
p=P(A).
\]

Положим

\[
X_k=\mathbf 1_{\{A\text{ произошло в }k\text{-м испытании}\}}.
\]

Тогда

\[
P\{X_k=1\}=p,
\qquad
P\{X_k=0\}=q=1-p,
\]

и

\[
\mathbb EX_k=p,
\qquad
\operatorname{Var}X_k=pq.
\]

Число появлений события \(A\) за \(n\) испытаний равно

\[
N_n(A)=X_1+\cdots+X_n.
\]

Относительная частота равна

\[
\nu_n(A)=\frac{N_n(A)}{n}.
\]

Закон больших чисел Бернулли утверждает:

\[
\frac{N_n(A)}{n}\xrightarrow{P}p.
\]

То есть для каждого \(\varepsilon>0\)

\[
P\left\{
\left|\frac{N_n(A)}{n}-p\right|>\varepsilon
\right\}
\to0.
\]

Если \(0<p<1\), то из ЦПТ для индикаторов получаем более точную формулу:

\[
\frac{N_n(A)-np}{\sqrt{npq}}
\Rightarrow
N(0,1).
\]

Отсюда снова следует

\[
\frac{N_n(A)}{n}-p
=
\sqrt{\frac{pq}{n}}\,
\frac{N_n(A)-np}{\sqrt{npq}}
\xrightarrow{P}0.
\]

\subsubsection{2.7. Связь с аксиоматикой
Колмогорова}\label{ux441ux432ux44fux437ux44c-ux441-ux430ux43aux441ux438ux43eux43cux430ux442ux438ux43aux43eux439-ux43aux43eux43bux43cux43eux433ux43eux440ux43eux432ux430}

Именно этот пример объясняет связь с аксиоматикой. В аксиоматике
Колмогорова вероятность - это не предел частоты по определению, а мера

\[
P:\mathcal F\to[0,1]
\]

на \(\sigma\)-алгебре событий, такая что

\[
P(\Omega)=1
\]

и для попарно несовместных событий \(A_1,A_2,\ldots\)

\[
P\left(\bigcup_{k=1}^{\infty}A_k\right)
=
\sum_{k=1}^{\infty}P(A_k).
\]

Такой подход математически удобен: он позволяет строить случайные
величины, распределения, интегралы, условные ожидания и случайные
процессы. Но возникает естественный вопрос: почему это число \(P(A)\)
имеет отношение к эксперименту?

Ответ дает закон больших чисел: если событие \(A\) наблюдается в длинной
серии независимых повторений одного и того же опыта, то относительная
частота события будет близка к \(P(A)\) с вероятностью, стремящейся к
единице. Поэтому мера \(P\) получает эмпирическую интерпретацию через
устойчивость частот.

\subsubsection{2.8. Эмпирическая интерпретация и применимость
модели}\label{ux44dux43cux43fux438ux440ux438ux447ux435ux441ux43aux430ux44f-ux438ux43dux442ux435ux440ux43fux440ux435ux442ux430ux446ux438ux44f-ux438-ux43fux440ux438ux43cux435ux43dux438ux43cux43eux441ux442ux44c-ux43cux43eux434ux435ux43bux438}

Тонкость: закон больших чисел не доказывает аксиомы вероятности из
физики. Он показывает, что если мы приняли колмогоровскую модель и
независимость повторных испытаний, то модель предсказывает частотную
устойчивость. Поэтому ЗБЧ является не логическим выводом аксиом из
частот, а математическим обоснованием применимости аксиоматики к
повторяемым экспериментам.

\subsection{3. Основные
определения}\label{ux43eux441ux43dux43eux432ux43dux44bux435-ux43eux43fux440ux435ux434ux435ux43bux435ux43dux438ux44f-20}

\subsubsection{Сходимость по
вероятности}\label{ux441ux445ux43eux434ux438ux43cux43eux441ux442ux44c-ux43fux43e-ux432ux435ux440ux43eux44fux442ux43dux43eux441ux442ux438}

Последовательность случайных величин \(Y_n\) сходится по вероятности к
\(Y\), если для каждого \(\varepsilon>0\)

\[
P\{|Y_n-Y|>\varepsilon\}\to0.
\]

Обозначение:

\[
Y_n\xrightarrow{P}Y.
\]

\subsubsection{Сходимость по
распределению}\label{ux441ux445ux43eux434ux438ux43cux43eux441ux442ux44c-ux43fux43e-ux440ux430ux441ux43fux440ux435ux434ux435ux43bux435ux43dux438ux44e}

Последовательность \(Y_n\) сходится по распределению к \(Y\), если

\[
F_{Y_n}(x)\to F_Y(x)
\]

во всех точках непрерывности \(F_Y\). Обозначение:

\[
Y_n\Rightarrow Y.
\]

Эквивалентно, для всех ограниченных непрерывных функций \(f\):

\[
\mathbb Ef(Y_n)\to\mathbb Ef(Y).
\]

Эта формулировка связана с Билетом 20, где распределения рассматриваются
как функционалы на тестовых функциях.

\subsubsection{Ограниченность по
вероятности}\label{ux43eux433ux440ux430ux43dux438ux447ux435ux43dux43dux43eux441ux442ux44c-ux43fux43e-ux432ux435ux440ux43eux44fux442ux43dux43eux441ux442ux438}

Последовательность \(Y_n\) называется ограниченной по вероятности, если
для каждого \(\delta>0\) существует \(M>0\) такое, что при всех
достаточно больших \(n\)

\[
P\{|Y_n|>M\}<\delta.
\]

Пишут:

\[
Y_n=O_P(1).
\]

Если \(Y_n\Rightarrow Y\), то \(Y_n=O_P(1)\).

\subsubsection{Слабый закон больших
чисел}\label{ux441ux43bux430ux431ux44bux439-ux437ux430ux43aux43eux43d-ux431ux43eux43bux44cux448ux438ux445-ux447ux438ux441ux435ux43b}

Для независимых одинаково распределенных величин \(X_1,X_2,\ldots\) с
\(\mathbb EX_1=m\) слабый ЗБЧ утверждает

\[
\frac{X_1+\cdots+X_n}{n}\xrightarrow{P}m
\]

при соответствующих условиях. Для iid-величин достаточно
\(\mathbb E|X_1|<\infty\).

В частном случае конечной дисперсии это легко доказывается неравенством
Чебышева.

\subsubsection{Усиленный закон больших
чисел}\label{ux443ux441ux438ux43bux435ux43dux43dux44bux439-ux437ux430ux43aux43eux43d-ux431ux43eux43bux44cux448ux438ux445-ux447ux438ux441ux435ux43b}

Усиленный ЗБЧ утверждает почти наверное сходимость:

\[
\frac{X_1+\cdots+X_n}{n}\to m
\quad\text{п.н.}
\]

Это более сильное утверждение, чем сходимость по вероятности.

\subsubsection{Центральная предельная
теорема}\label{ux446ux435ux43dux442ux440ux430ux43bux44cux43dux430ux44f-ux43fux440ux435ux434ux435ux43bux44cux43dux430ux44f-ux442ux435ux43eux440ux435ux43cux430-1}

Если \(X_1,X_2,\ldots\) - iid-величины,

\[
\mathbb EX_1=m,
\qquad
\operatorname{Var}X_1=\sigma^2\in(0,\infty),
\]

то

\[
\frac{S_n-nm}{\sigma\sqrt n}
\Rightarrow
N(0,1).
\]

\subsubsection{Схема
Бернулли}\label{ux441ux445ux435ux43cux430-ux431ux435ux440ux43dux443ux43bux43bux438}

Это последовательность независимых испытаний с двумя исходами: успех с
вероятностью \(p\) и неудача с вероятностью \(q=1-p\). Если \(X_k\) -
индикатор успеха, то

\[
S_n=X_1+\cdots+X_n
\]

имеет биномиальное распределение:

\[
P\{S_n=k\}=C_n^k p^kq^{n-k}.
\]

\subsubsection{Относительная
частота}\label{ux43eux442ux43dux43eux441ux438ux442ux435ux43bux44cux43dux430ux44f-ux447ux430ux441ux442ux43eux442ux430}

Если событие \(A\) произошло \(N_n(A)\) раз в \(n\) испытаниях, то

\[
\nu_n(A)=\frac{N_n(A)}{n}
\]

называется относительной частотой события \(A\).

\subsection{4. Основные теоремы и
свойства}\label{ux43eux441ux43dux43eux432ux43dux44bux435-ux442ux435ux43eux440ux435ux43cux44b-ux438-ux441ux432ux43eux439ux441ux442ux432ux430-20}

\subsubsection{1. Неравенство
Чебышева}\label{ux43dux435ux440ux430ux432ux435ux43dux441ux442ux432ux43e-ux447ux435ux431ux44bux448ux435ux432ux430}

Если случайная величина \(Y\) имеет конечную дисперсию, то для любого
\(\varepsilon>0\)

\[
P\{|Y-\mathbb EY|\ge\varepsilon\}
\le
\frac{\operatorname{Var}Y}{\varepsilon^2}.
\]

Это базовый инструмент доказательства слабого ЗБЧ при конечной
дисперсии.

\subsubsection{2. Слабый ЗБЧ при конечной
дисперсии}\label{ux441ux43bux430ux431ux44bux439-ux437ux431ux447-ux43fux440ux438-ux43aux43eux43dux435ux447ux43dux43eux439-ux434ux438ux441ux43fux435ux440ux441ux438ux438}

Пусть \(X_1,X_2,\ldots\) независимы, одинаково распределены,

\[
\mathbb EX_1=m,
\qquad
\operatorname{Var}X_1=\sigma^2<\infty.
\]

Тогда

\[
\overline X_n\xrightarrow{P}m.
\]

Доказательство:

\[
\mathbb E\overline X_n=m,
\]

а по независимости

\[
\operatorname{Var}\overline X_n
=
\operatorname{Var}\left(\frac{S_n}{n}\right)
=
\frac{1}{n^2}\operatorname{Var}S_n
=
\frac{1}{n^2}n\sigma^2
=
\frac{\sigma^2}{n}.
\]

По неравенству Чебышева:

\[
P\{|\overline X_n-m|>\varepsilon\}
\le
\frac{\sigma^2}{n\varepsilon^2}
\to0.
\]

\subsubsection{3. ЦПТ для
iid-величин}\label{ux446ux43fux442-ux434ux43bux44f-iid-ux432ux435ux43bux438ux447ux438ux43d}

При условиях

\[
\mathbb EX_1=m,
\qquad
\operatorname{Var}X_1=\sigma^2\in(0,\infty)
\]

верна сходимость

\[
\frac{S_n-nm}{\sigma\sqrt n}
\Rightarrow
N(0,1).
\]

Смысл: сумма \(S_n\) сосредоточена около \(nm\), а ее случайные
флуктуации имеют масштаб \(\sqrt n\).

\subsubsection{4. ЗБЧ как следствие
ЦПТ}\label{ux437ux431ux447-ux43aux430ux43a-ux441ux43bux435ux434ux441ux442ux432ux438ux435-ux446ux43fux442}

Если выполнены условия ЦПТ, то

\[
\overline X_n-m
=
\frac{\sigma}{\sqrt n}
\frac{S_n-nm}{\sigma\sqrt n}.
\]

Второй множитель сходится по распределению к \(N(0,1)\), значит он
ограничен по вероятности. Первый множитель стремится к нулю.
Следовательно,

\[
\overline X_n-m\xrightarrow{P}0.
\]

Это и есть слабый ЗБЧ.

\subsubsection{5. Закон больших чисел
Бернулли}\label{ux437ux430ux43aux43eux43d-ux431ux43eux43bux44cux448ux438ux445-ux447ux438ux441ux435ux43b-ux431ux435ux440ux43dux443ux43bux43bux438}

В схеме Бернулли

\[
S_n\sim\operatorname{Bin}(n,p).
\]

Тогда

\[
\frac{S_n}{n}\xrightarrow{P}p.
\]

Через неравенство Чебышева:

\[
\operatorname{Var}\frac{S_n}{n}
=
\frac{pq}{n},
\]

поэтому

\[
P\left\{
\left|\frac{S_n}{n}-p\right|>\varepsilon
\right\}
\le
\frac{pq}{n\varepsilon^2}
\to0.
\]

Через ЦПТ:

\[
\frac{S_n-np}{\sqrt{npq}}
\Rightarrow
N(0,1),
\]

откуда

\[
\frac{S_n}{n}-p
=
\sqrt{\frac{pq}{n}}
\frac{S_n-np}{\sqrt{npq}}
\xrightarrow{P}0.
\]

\subsubsection{6. Частотная интерпретация
вероятности}\label{ux447ux430ux441ux442ux43eux442ux43dux430ux44f-ux438ux43dux442ux435ux440ux43fux440ux435ux442ux430ux446ux438ux44f-ux432ux435ux440ux43eux44fux442ux43dux43eux441ux442ux438}

Если \(A\) - событие с вероятностью \(P(A)=p\), а \(\nu_n(A)\) - частота
его появления в \(n\) независимых повторениях опыта, то

\[
\nu_n(A)\xrightarrow{P}P(A).
\]

Это объясняет, почему вероятность можно практически оценивать частотой:

\[
P(A)\approx \frac{N_n(A)}{n}
\]

при большом \(n\).

\subsection{5. Примеры}\label{ux43fux440ux438ux43cux435ux440ux44b-20}

\subsubsection{Пример 1. Подбрасывание
монеты}\label{ux43fux440ux438ux43cux435ux440-1.-ux43fux43eux434ux431ux440ux430ux441ux44bux432ux430ux43dux438ux435-ux43cux43eux43dux435ux442ux44b}

Пусть \(A\) - выпадение орла, \(P(A)=p\). В \(n\) независимых
подбрасываниях

\[
N_n(A)\sim\operatorname{Bin}(n,p).
\]

Тогда

\[
\frac{N_n(A)}{n}\xrightarrow{P}p.
\]

Если монета правильная, \(p=1/2\), и

\[
\frac{N_n(A)-n/2}{\sqrt{n}/2}
\Rightarrow
N(0,1).
\]

ЗБЧ говорит, что доля орлов близка к \(1/2\), а ЦПТ говорит, какого
порядка типичное отклонение:

\[
N_n(A)-\frac n2
\]

обычно имеет порядок \(\sqrt n\), а не порядок \(n\).

\subsubsection{Пример 2. Среднее
измерений}\label{ux43fux440ux438ux43cux435ux440-2.-ux441ux440ux435ux434ux43dux435ux435-ux438ux437ux43cux435ux440ux435ux43dux438ux439}

Пусть измеряется величина \(m\), а результат \(k\)-го измерения:

\[
X_k=m+\xi_k,
\]

где ошибки \(\xi_k\) независимы,

\[
\mathbb E\xi_k=0,
\qquad
\operatorname{Var}\xi_k=\sigma^2.
\]

Тогда среднее измерений

\[
\overline X_n=m+\frac{\xi_1+\cdots+\xi_n}{n}
\]

сходится к \(m\) по вероятности. ЦПТ уточняет:

\[
\frac{\sqrt n(\overline X_n-m)}{\sigma}
\Rightarrow
N(0,1).
\]

Поэтому ошибка среднего имеет порядок

\[
\frac{\sigma}{\sqrt n}.
\]

\subsubsection{Пример 3. Оценка вероятности
события}\label{ux43fux440ux438ux43cux435ux440-3.-ux43eux446ux435ux43dux43aux430-ux432ux435ux440ux43eux44fux442ux43dux43eux441ux442ux438-ux441ux43eux431ux44bux442ux438ux44f}

Пусть нужно оценить вероятность отказа устройства \(A\). Проводим \(n\)
независимых испытаний и ставим

\[
X_k=\mathbf 1_{\{A\text{ произошло в }k\text{-м испытании}\}}.
\]

Тогда

\[
\widehat p_n=\frac{1}{n}\sum_{k=1}^nX_k
\]

\begin{itemize}
\tightlist
\item
  естественная оценка \(p=P(A)\). По ЗБЧ
\end{itemize}

\[
\widehat p_n\xrightarrow{P}p.
\]

Если \(0<p<1\), то по ЦПТ

\[
\frac{\widehat p_n-p}{\sqrt{p(1-p)/n}}
\Rightarrow
N(0,1).
\]

Это основа нормальных доверительных интервалов для вероятности.

\subsubsection{Пример 4. Случайное
блуждание}\label{ux43fux440ux438ux43cux435ux440-4.-ux441ux43bux443ux447ux430ux439ux43dux43eux435-ux431ux43bux443ux436ux434ux430ux43dux438ux435}

Пусть

\[
S_n=\xi_1+\cdots+\xi_n,
\]

где \(\xi_k\) - iid-шаги,

\[
\mathbb E\xi_k=a,
\qquad
\operatorname{Var}\xi_k=\sigma^2.
\]

Тогда

\[
\frac{S_n}{n}\xrightarrow{P}a,
\]

а при \(\sigma^2>0\)

\[
\frac{S_n-na}{\sigma\sqrt n}
\Rightarrow
N(0,1).
\]

Это связывает билет с Билетом 19: траектория блуждания в среднем идет с
линейной скоростью \(a\), а случайные отклонения от прямой имеют масштаб
\(\sqrt n\).

\subsection{6. Схемы доказательств /
обоснований}\label{ux441ux445ux435ux43cux44b-ux434ux43eux43aux430ux437ux430ux442ux435ux43bux44cux441ux442ux432-ux43eux431ux43eux441ux43dux43eux432ux430ux43dux438ux439-13}

\textbf{Доказательство слабого ЗБЧ через Чебышева.}

\begin{enumerate}
\def\labelenumi{\arabic{enumi}.}
\tightlist
\item
  Берем
\end{enumerate}

\[
\overline X_n=\frac{S_n}{n}.
\]

\begin{enumerate}
\def\labelenumi{\arabic{enumi}.}
\setcounter{enumi}{1}
\tightlist
\item
  Считаем матожидание:
\end{enumerate}

\[
\mathbb E\overline X_n=m.
\]

\begin{enumerate}
\def\labelenumi{\arabic{enumi}.}
\setcounter{enumi}{2}
\tightlist
\item
  Считаем дисперсию, используя независимость:
\end{enumerate}

\[
\operatorname{Var}\overline X_n
=
\frac{\sigma^2}{n}.
\]

\begin{enumerate}
\def\labelenumi{\arabic{enumi}.}
\setcounter{enumi}{3}
\tightlist
\item
  Применяем Чебышева:
\end{enumerate}

\[
P\{|\overline X_n-m|>\varepsilon\}
\le
\frac{\sigma^2}{n\varepsilon^2}.
\]

\begin{enumerate}
\def\labelenumi{\arabic{enumi}.}
\setcounter{enumi}{4}
\tightlist
\item
  Правая часть стремится к нулю, значит
\end{enumerate}

\[
\overline X_n\xrightarrow{P}m.
\]

\textbf{Доказательство ЗБЧ как следствия ЦПТ.}

\begin{enumerate}
\def\labelenumi{\arabic{enumi}.}
\tightlist
\item
  По ЦПТ:
\end{enumerate}

\[
Z_n=\frac{S_n-nm}{\sigma\sqrt n}
\Rightarrow
Z\sim N(0,1).
\]

\begin{enumerate}
\def\labelenumi{\arabic{enumi}.}
\setcounter{enumi}{1}
\item
  Значит \(Z_n=O_P(1)\).
\item
  Переписываем отклонение среднего:
\end{enumerate}

\[
\overline X_n-m
=
\frac{\sigma}{\sqrt n}Z_n.
\]

\begin{enumerate}
\def\labelenumi{\arabic{enumi}.}
\setcounter{enumi}{3}
\tightlist
\item
  Произведение ограниченной по вероятности последовательности на
  числовую последовательность, стремящуюся к нулю, сходится к нулю по
  вероятности:
\end{enumerate}

\[
\frac{\sigma}{\sqrt n}Z_n\xrightarrow{P}0.
\]

\begin{enumerate}
\def\labelenumi{\arabic{enumi}.}
\setcounter{enumi}{4}
\tightlist
\item
  Следовательно,
\end{enumerate}

\[
\overline X_n\xrightarrow{P}m.
\]

\textbf{Доказательство частотной формулы.}

\begin{enumerate}
\def\labelenumi{\arabic{enumi}.}
\tightlist
\item
  Для события \(A\) вводим индикаторы:
\end{enumerate}

\[
X_k=\mathbf 1_{\{A\text{ произошло в }k\text{-м испытании}\}}.
\]

\begin{enumerate}
\def\labelenumi{\arabic{enumi}.}
\setcounter{enumi}{1}
\tightlist
\item
  Тогда
\end{enumerate}

\[
\mathbb EX_k=P(A)=p.
\]

\begin{enumerate}
\def\labelenumi{\arabic{enumi}.}
\setcounter{enumi}{2}
\tightlist
\item
  Частота равна среднему индикаторов:
\end{enumerate}

\[
\nu_n(A)=\frac{X_1+\cdots+X_n}{n}.
\]

\begin{enumerate}
\def\labelenumi{\arabic{enumi}.}
\setcounter{enumi}{3}
\tightlist
\item
  По ЗБЧ:
\end{enumerate}

\[
\nu_n(A)\xrightarrow{P}p=P(A).
\]

Это и есть математическая форма утверждения: ``вероятность проявляется
как устойчивая частота в длинной серии независимых испытаний''.

\subsection{7. Что написать на
доске}\label{ux447ux442ux43e-ux43dux430ux43fux438ux441ux430ux442ux44c-ux43dux430-ux434ux43eux441ux43aux435-20}

Минимальный набор формул:

\[
S_n=X_1+\cdots+X_n,
\qquad
\overline X_n=\frac{S_n}{n}.
\]

ЗБЧ:

\[
\overline X_n\xrightarrow{P}m,
\qquad
m=\mathbb EX_1.
\]

ЦПТ:

\[
\frac{S_n-nm}{\sigma\sqrt n}
\Rightarrow
N(0,1).
\]

Вывод ЗБЧ из ЦПТ:

\[
\overline X_n-m
=
\frac{\sigma}{\sqrt n}
\frac{S_n-nm}{\sigma\sqrt n}
\xrightarrow{P}0.
\]

Частоты:

\[
X_k=\mathbf 1_A^{(k)},
\qquad
\nu_n(A)=\frac{1}{n}\sum_{k=1}^nX_k,
\qquad
\nu_n(A)\xrightarrow{P}P(A).
\]

Фраза про аксиоматику:

\[
P(A)\ \text{не определяется как частота, но ЗБЧ объясняет частотный смысл }P(A).
\]

\subsection{8. Типичные дополнительные
вопросы}\label{ux442ux438ux43fux438ux447ux43dux44bux435-ux434ux43eux43fux43eux43bux43dux438ux442ux435ux43bux44cux43dux44bux435-ux432ux43eux43fux440ux43eux441ux44b-20}

\textbf{В каком смысле ЗБЧ является частным случаем ЦПТ?}

При условиях ЦПТ нормированное отклонение

\[
\frac{S_n-nm}{\sigma\sqrt n}
\]

имеет невырожденный предельный закон, значит отклонение \(S_n-nm\) имеет
порядок \(\sqrt n\). После деления на \(n\) получаем отклонение среднего
порядка \(1/\sqrt n\), которое стремится к нулю по вероятности.

\textbf{Почему нельзя сказать, что любой ЗБЧ следует из ЦПТ?}

Потому что ЦПТ требует конечной ненулевой дисперсии, а некоторые формы
ЗБЧ верны при более слабых условиях. Например, для iid-величин слабый
ЗБЧ верен при \(\mathbb E|X_1|<\infty\), даже если дисперсия бесконечна.

\textbf{Что сильнее: сходимость по вероятности или по распределению?}

Сходимость по вероятности к константе влечет сходимость по распределению
к этой константе. Обратно в общем случае неверно, но если

\[
Y_n\Rightarrow c
\]

к константе \(c\), то

\[
Y_n\xrightarrow{P}c.
\]

В выводе ЗБЧ из ЦПТ мы используем не это напрямую, а факт, что
нормированная сумма ограничена по вероятности.

\textbf{Зачем в ЦПТ нормировка именно на \(\sqrt n\)?}

Для независимых одинаково распределенных величин дисперсии суммируются:

\[
\operatorname{Var}S_n=n\sigma^2.
\]

Стандартное отклонение суммы равно

\[
\sigma\sqrt n.
\]

Поэтому именно деление на \(\sigma\sqrt n\) дает невырожденный
предельный закон.

\textbf{Как связаны ЗБЧ и оценивание вероятности?}

Если \(X_k=\mathbf 1_A^{(k)}\), то

\[
\mathbb EX_k=P(A).
\]

Среднее индикаторов равно частоте события. По ЗБЧ частота сходится к
\(P(A)\), поэтому частота является состоятельной оценкой вероятности.

\textbf{Что именно обосновывает ЗБЧ в аксиоматике Колмогорова?}

Он обосновывает интерпретацию вероятности как идеального значения, около
которого стабилизируются частоты в независимых повторениях. Но сами
аксиомы - \(\sigma\)-алгебра, мера, счетная аддитивность - остаются
математическими предположениями модели.

\textbf{Почему независимость важна?}

Без независимости среднее может не стабилизироваться около
математического ожидания. Например, если все \(X_k\) равны одной и той
же случайной величине \(X\), то

\[
\overline X_n=X
\]

для всех \(n\), и среднее не обязано сходиться к \(\mathbb EX\) по
вероятности.

\textbf{Что будет при \(\sigma^2=0\)?}

Тогда \(X_1=m\) почти наверное, и

\[
\overline X_n=m
\]

почти наверное для всех \(n\). ЗБЧ тривиален, а стандартная формула ЦПТ
с делением на \(\sigma\) неприменима.

\subsection{9. Типичные
ошибки}\label{ux442ux438ux43fux438ux447ux43dux44bux435-ux43eux448ux438ux431ux43aux438-20}

\begin{itemize}
\tightlist
\item
  Говорить, что частота ``равна'' вероятности при большом \(n\).
  Правильно: частота близка к вероятности с большой вероятностью.
\item
  Забывать вид сходимости: в слабом ЗБЧ это сходимость по вероятности,
  не поточечная сходимость для каждого исхода.
\item
  Называть ЗБЧ частным случаем ЦПТ без оговорки о конечной дисперсии.
\item
  Путать масштаб ЦПТ и ЗБЧ: сумма отклоняется на \(\sqrt n\), среднее
  отклоняется на \(1/\sqrt n\).
\item
  Забывать независимость или хотя бы условия слабой зависимости.
\item
  Для схемы Бернулли писать дисперсию частоты как \(npq\) вместо
  \(pq/n\).
\item
  Считать, что ЗБЧ доказывает существование вероятности как предела
  частот. В колмогоровской теории вероятность задается мерой, а ЗБЧ
  объясняет частотную устойчивость внутри модели.
\end{itemize}

\subsection{10. Связи с другими
билетами}\label{ux441ux432ux44fux437ux438-ux441-ux434ux440ux443ux433ux438ux43cux438-ux431ux438ux43bux435ux442ux430ux43cux438-20}

\begin{itemize}
\tightlist
\item
  Билет 02: вероятностное пространство Колмогорова и счетная
  аддитивность меры.
\item
  Билет 08: характеристические функции используются в доказательстве
  ЦПТ.
\item
  Билет 09: независимость нужна для произведения характеристических
  функций и сложения дисперсий.
\item
  Билет 11: математическое ожидание и дисперсия выражаются как
  интегралы.
\item
  Билет 13: моменты связаны с производными характеристической функции.
\item
  Билет 16: последовательность наблюдений можно рассматривать как
  дискретный случайный процесс.
\item
  Билет 17: эргодичность обобщает идею сходимости временных средних.
\item
  Билет 19: для случайных блужданий ЗБЧ и ЦПТ описывают асимптотику
  суммы шагов.
\item
  Билет 20: слабая сходимость распределений формулируется через тестовые
  функции.
\item
  Билет 24: достаточные условия эргодичности дают аналоги ЗБЧ для
  зависимых стационарных последовательностей.
\end{itemize}

\subsection{11. Минимум на
удовлетворительно}\label{ux43cux438ux43dux438ux43cux443ux43c-ux43dux430-ux443ux434ux43eux432ux43bux435ux442ux432ux43eux440ux438ux442ux435ux43bux44cux43dux43e-20}

Нужно сказать:

\begin{enumerate}
\def\labelenumi{\arabic{enumi}.}
\tightlist
\item
  ЗБЧ:
\end{enumerate}

\[
\frac{X_1+\cdots+X_n}{n}\xrightarrow{P}\mathbb EX_1.
\]

\begin{enumerate}
\def\labelenumi{\arabic{enumi}.}
\setcounter{enumi}{1}
\tightlist
\item
  ЦПТ:
\end{enumerate}

\[
\frac{X_1+\cdots+X_n-nm}{\sigma\sqrt n}
\Rightarrow
N(0,1).
\]

\begin{enumerate}
\def\labelenumi{\arabic{enumi}.}
\setcounter{enumi}{2}
\item
  Из ЦПТ следует ЗБЧ, потому что отклонение суммы от \(nm\) имеет
  порядок \(\sqrt n\), а после деления на \(n\) становится малым.
\item
  Для индикаторов событий:
\end{enumerate}

\[
\frac{N_n(A)}{n}\xrightarrow{P}P(A).
\]

\begin{enumerate}
\def\labelenumi{\arabic{enumi}.}
\setcounter{enumi}{4}
\tightlist
\item
  Это объясняет частотный смысл вероятности.
\end{enumerate}

\subsection{12. Минимум на
хорошо}\label{ux43cux438ux43dux438ux43cux443ux43c-ux43dux430-ux445ux43eux440ux43eux448ux43e-20}

Нужно дополнительно:

\begin{itemize}
\tightlist
\item
  строго определить сходимость по вероятности и сходимость по
  распределению;
\item
  назвать условия ЦПТ: iid, конечное среднее, конечная ненулевая
  дисперсия;
\item
  аккуратно вывести
\end{itemize}

\[
\overline X_n-m
=
\frac{\sigma}{\sqrt n}Z_n;
\]

\begin{itemize}
\tightlist
\item
  разобрать схему Бернулли;
\item
  сказать, что ЗБЧ шире ЦПТ и не всегда требует конечной дисперсии.
\end{itemize}

\subsection{13. Что добавить для
отлично}\label{ux447ux442ux43e-ux434ux43eux431ux430ux432ux438ux442ux44c-ux434ux43bux44f-ux43eux442ux43bux438ux447ux43dux43e-20}

Для сильного ответа стоит добавить:

\begin{itemize}
\tightlist
\item
  доказательство слабого ЗБЧ через Чебышева;
\item
  объяснение ограниченности по вероятности при сходимости по
  распределению;
\item
  различие между слабым и усиленным ЗБЧ;
\item
  фразу о том, что аксиоматика не выводится из частот, а получает через
  ЗБЧ экспериментальную интерпретацию;
\item
  связь с эргодичностью: в более общих зависимых процессах роль ЗБЧ
  играют эргодические теоремы.
\end{itemize}

\subsection{14. Устная версия ответа на 5
минут}\label{ux443ux441ux442ux43dux430ux44f-ux432ux435ux440ux441ux438ux44f-ux43eux442ux432ux435ux442ux430-ux43dux430-5-ux43cux438ux43dux443ux442-20}

Закон больших чисел и центральная предельная теорема описывают
асимптотику сумм независимых одинаково распределенных случайных величин.
Пусть

\[
S_n=X_1+\cdots+X_n,
\qquad
\mathbb EX_1=m.
\]

Слабый закон больших чисел утверждает:

\[
\frac{S_n}{n}\xrightarrow{P}m.
\]

То есть для каждого \(\varepsilon>0\)

\[
P\left\{\left|\frac{S_n}{n}-m\right|>\varepsilon\right\}\to0.
\]

Центральная предельная теорема требует дополнительно конечной ненулевой
дисперсии

\[
\operatorname{Var}X_1=\sigma^2\in(0,\infty)
\]

и утверждает:

\[
\frac{S_n-nm}{\sigma\sqrt n}
\Rightarrow
N(0,1).
\]

Из этой формулы сразу видно, что отклонение суммы от \(nm\) имеет
масштаб \(\sqrt n\). Поэтому

\[
\frac{S_n}{n}-m
=
\frac{\sigma}{\sqrt n}
\frac{S_n-nm}{\sigma\sqrt n}
\xrightarrow{P}0.
\]

Значит при условиях ЦПТ закон больших чисел получается как более грубое
следствие: ЦПТ описывает форму флуктуаций, а ЗБЧ говорит только, что
среднее стабилизируется.

Главный пример - схема Бернулли. Если \(A\) происходит с вероятностью
\(p\), а

\[
X_k=\mathbf 1_A^{(k)},
\]

то

\[
\mathbb EX_k=p.
\]

Частота события равна

\[
\nu_n(A)=\frac{X_1+\cdots+X_n}{n}.
\]

По ЗБЧ:

\[
\nu_n(A)\xrightarrow{P}p=P(A).
\]

Это объясняет связь с аксиоматикой Колмогорова. В аксиоматике
вероятность - это мера на \(\sigma\)-алгебре событий, а не частота по
определению. Но закон больших чисел показывает, что в длинных сериях
независимых повторений относительные частоты стабилизируются около этой
меры. Поэтому аксиоматическая вероятность имеет частотную интерпретацию
и применима к экспериментам.

В конце важно добавить оговорку: не всякий ЗБЧ является частным случаем
ЦПТ, потому что ЗБЧ может быть верен при более слабых условиях. Но при
конечной дисперсии ЦПТ действительно влечет слабый закон больших чисел.

\subsection{15. Если осталось 2 минуты перед
экзаменом}\label{ux435ux441ux43bux438-ux43eux441ux442ux430ux43bux43eux441ux44c-2-ux43cux438ux43dux443ux442ux44b-ux43fux435ux440ux435ux434-ux44dux43aux437ux430ux43cux435ux43dux43eux43c-20}

Запомнить три формулы.

ЗБЧ:

\[
\overline X_n=\frac{S_n}{n}\xrightarrow{P}m.
\]

ЦПТ:

\[
\frac{S_n-nm}{\sigma\sqrt n}
\Rightarrow
N(0,1).
\]

Вывод:

\[
\overline X_n-m
=
\frac{\sigma}{\sqrt n}
\frac{S_n-nm}{\sigma\sqrt n}
\xrightarrow{P}0.
\]

Для события \(A\):

\[
\frac{N_n(A)}{n}\xrightarrow{P}P(A).
\]

Главная фраза: вероятность в аксиоматике Колмогорова задается мерой, а
закон больших чисел объясняет, почему эта мера наблюдается как
устойчивая частота в длинной серии независимых испытаний.

\newpage

\section{Билет 22. Марковские переходные функции и
матрицы}\label{ux431ux438ux43bux435ux442-22.-ux43cux430ux440ux43aux43eux432ux441ux43aux438ux435-ux43fux435ux440ux435ux445ux43eux434ux43dux44bux435-ux444ux443ux43dux43aux446ux438ux438-ux438-ux43cux430ux442ux440ux438ux446ux44b}

Источники: Ширяев 2, гл. VIII, §1, §3; Сердобольская, лекция 5.

\subsection{1. Короткий
ответ}\label{ux43aux43eux440ux43eux442ux43aux438ux439-ux43eux442ux432ux435ux442-21}

Переходная функция, или марковское ядро, задает условный закон
следующего состояния при известном текущем состоянии. Если пространство
состояний \((E,\mathcal E)\) произвольное, то переходная функция имеет
вид

\[
P(x,A)=P\{X_{n+1}\in A\mid X_n=x\},
\qquad
x\in E,\ A\in\mathcal E.
\]

В общем пространстве состояний это понимают как регулярную версию
условного распределения \(P(X_{n+1}\in A\mid X_n)\). В счетном случае
запись через условие \(X_n=x\) безопасна буквально. По \(A\) это
вероятностная мера, по \(x\) - измеримая функция.

Если пространство состояний счетное или конечное, то переходная функция
задается матрицей

\[
P=(p_{ij}),
\qquad
p_{ij}=P\{X_{n+1}=j\mid X_n=i\}.
\]

Матрица переходов стохастическая:

\[
p_{ij}\ge0,
\qquad
\sum_j p_{ij}=1.
\]

При соглашении, что распределения записываются строками, эволюция
распределений задается формулой

\[
\pi_{n+1}=\pi_nP.
\]

Переходы за несколько шагов задаются степенями матрицы:

\[
P^{(n)}=P^n.
\]

Главное свойство - уравнение Чепмена-Колмогорова:

\[
P^{(m+n)}(x,A)
=
\int_E P^{(n)}(x,dy)P^{(m)}(y,A).
\]

В матричном случае:

\[
p_{ij}^{(m+n)}
=
\sum_k p_{ik}^{(n)}p_{kj}^{(m)}.
\]

\subsection{2. Полный
ответ}\label{ux43fux43eux43bux43dux44bux439-ux43eux442ux432ux435ux442-21}

\subsubsection{22.1. Предмет билета и
контекст}\label{ux43fux440ux435ux434ux43cux435ux442-ux431ux438ux43bux435ux442ux430-ux438-ux43aux43eux43dux442ux435ux43aux441ux442}

Этот билет является технической основой для Билета 23 о марковских
цепях. В билете 23 главный объект - сам процесс с марковским свойством.
В билете 22 главный объект - переходный механизм: как задается условное
распределение будущего состояния через настоящее.

\subsubsection{22.2. Определение переходной
функции}\label{ux43eux43fux440ux435ux434ux435ux43bux435ux43dux438ux435-ux43fux435ux440ux435ux445ux43eux434ux43dux43eux439-ux444ux443ux43dux43aux446ux438ux438}

Пусть

\[
X=(X_0,X_1,X_2,\ldots)
\]

\begin{itemize}
\tightlist
\item
  случайная последовательность со значениями в измеримом пространстве
  состояний
\end{itemize}

\[
(E,\mathcal E).
\]

Марковская идея состоит в том, что для предсказания следующего состояния
достаточно знать текущее. Если процесс однороден по времени, то условный
закон перехода из состояния \(x\) во множество \(A\) не зависит от
номера шага \(n\):

\[
P\{X_{n+1}\in A\mid X_n=x\}=P(x,A).
\]

\subsubsection{22.3. Свойства переходного
ядра}\label{ux441ux432ux43eux439ux441ux442ux432ux430-ux43fux435ux440ux435ux445ux43eux434ux43dux43eux433ux43e-ux44fux434ux440ux430}

Функция

\[
P:E\times\mathcal E\to[0,1]
\]

называется переходной функцией, переходной вероятностью или марковским
ядром. Она должна удовлетворять двум условиям:

\begin{enumerate}
\def\labelenumi{\arabic{enumi}.}
\tightlist
\item
  При фиксированном \(x\in E\) функция
\end{enumerate}

\[
A\mapsto P(x,A)
\]

является вероятностной мерой на \((E,\mathcal E)\).

\begin{enumerate}
\def\labelenumi{\arabic{enumi}.}
\setcounter{enumi}{1}
\tightlist
\item
  При фиксированном \(A\in\mathcal E\) функция
\end{enumerate}

\[
x\mapsto P(x,A)
\]

является \(\mathcal E\)-измеримой.

Первое условие говорит, что из каждого состояния задано полноценное
распределение следующего состояния. Второе условие нужно для
корректности интегралов вида

\[
\int_E P(x,A)\,\mu(dx).
\]

\subsubsection{22.4. Преобразование
распределений}\label{ux43fux440ux435ux43eux431ux440ux430ux437ux43eux432ux430ux43dux438ux435-ux440ux430ux441ux43fux440ux435ux434ux435ux43bux435ux43dux438ux439}

Если \(\mu\) - распределение \(X_n\), то распределение \(X_{n+1}\)
находится так:

\[
\mu P(A)=\int_E P(x,A)\,\mu(dx).
\]

Здесь \(\mu P\) - новая вероятностная мера. Это просто формула полной
вероятности: чтобы попасть в \(A\) на следующем шаге, нужно
просуммировать или проинтегрировать по всем возможным текущим состояниям
\(x\).

\subsubsection{22.5. Матрица переходных вероятностей в счетном
случае}\label{ux43cux430ux442ux440ux438ux446ux430-ux43fux435ux440ux435ux445ux43eux434ux43dux44bux445-ux432ux435ux440ux43eux44fux442ux43dux43eux441ux442ux435ux439-ux432-ux441ux447ux435ux442ux43dux43eux43c-ux441ux43bux443ux447ux430ux435}

В счетном пространстве состояний

\[
E=\{1,2,\ldots\}
\]

переходная функция полностью задается числами

\[
p_{ij}=P(i,\{j\}).
\]

Тогда

\[
P(i,A)=\sum_{j\in A}p_{ij}.
\]

Матрица

\[
P=(p_{ij})_{i,j\in E}
\]

называется матрицей переходных вероятностей. Она является стохастической
по строкам:

\[
p_{ij}\ge0,
\qquad
\sum_{j\in E}p_{ij}=1.
\]

\subsubsection{22.6. Соглашение о векторно-матричном
умножении}\label{ux441ux43eux433ux43bux430ux448ux435ux43dux438ux435-ux43e-ux432ux435ux43aux442ux43eux440ux43dux43e-ux43cux430ux442ux440ux438ux447ux43dux43eux43c-ux443ux43cux43dux43eux436ux435ux43dux438ux438}

Важно: в этом конспекте используется строковое соглашение. Распределение
записывается как строка

\[
\pi_n=(\pi_n(i))_{i\in E},
\]

и тогда

\[
\pi_{n+1}=\pi_nP.
\]

Покомпонентно:

\[
\pi_{n+1}(j)
=
\sum_i \pi_n(i)p_{ij}.
\]

Если кто-то записывает распределения столбцами, то формула меняется на

\[
\pi_{n+1}=P^T\pi_n.
\]

Поэтому нужно всегда фиксировать соглашение. Типичная ошибка - объявить
строки стохастическими, а потом умножать так, как будто стохастические
столбцы.

\subsubsection{22.7. Вероятности перехода за несколько
шагов}\label{ux432ux435ux440ux43eux44fux442ux43dux43eux441ux442ux438-ux43fux435ux440ux435ux445ux43eux434ux430-ux437ux430-ux43dux435ux441ux43aux43eux43bux44cux43aux43e-ux448ux430ux433ux43eux432}

Теперь рассмотрим переходы за несколько шагов. Для однородной цепи
одношаговое ядро \(P\) не зависит от времени. Двухшаговое ядро
определяется формулой

\[
P^{(2)}(x,A)
=
\int_E P(x,dy)P(y,A).
\]

Смысл: сначала система из \(x\) попадает в промежуточное состояние
\(y\), потом из \(y\) попадает в \(A\).

Рекурсивно:

\[
P^{(1)}(x,A)=P(x,A),
\]

а для \(n\ge1\)

\[
P^{(n+1)}(x,A)
=
\int_E P^{(n)}(x,dy)P(y,A).
\]

Для матриц это ровно матричное умножение:

\[
p_{ij}^{(2)}
=
\sum_k p_{ik}p_{kj},
\]

а вообще

\[
P^{(n)}=P^n.
\]

\subsubsection{22.8. Уравнение Чепмена ---
Колмогорова}\label{ux443ux440ux430ux432ux43dux435ux43dux438ux435-ux447ux435ux43fux43cux435ux43dux430-ux43aux43eux43bux43cux43eux433ux43eux440ux43eux432ux430}

Главная формула билета - уравнение Чепмена-Колмогорова. Для однородной
цепи:

\[
P^{(m+n)}(x,A)
=
\int_E P^{(n)}(x,dy)P^{(m)}(y,A).
\]

Оно говорит, что переход за \(m+n\) шагов можно разложить через
промежуточное состояние после \(n\) шагов. В счетном случае интеграл
заменяется суммой:

\[
p_{ij}^{(m+n)}
=
\sum_{k\in E}p_{ik}^{(n)}p_{kj}^{(m)}.
\]

В матричной записи:

\[
P^{(m+n)}=P^{(n)}P^{(m)}.
\]

Так как для однородной цепи \(P^{(n)}=P^n\), это то же самое, что
обычное тождество

\[
P^{n+m}=P^nP^m.
\]

\subsubsection{22.9. Неоднородные по времени
цепи}\label{ux43dux435ux43eux434ux43dux43eux440ux43eux434ux43dux44bux435-ux43fux43e-ux432ux440ux435ux43cux435ux43dux438-ux446ux435ux43fux438}

В неоднородном случае переходный механизм может зависеть от времени.
Тогда вместо одного ядра берут семейство

\[
P_n(x,A)=P\{X_{n+1}\in A\mid X_n=x\}.
\]

Для счетного пространства это семейство матриц

\[
P_n=(p_{ij}(n)).
\]

Тогда распределения эволюционируют так:

\[
\pi_{n+1}=\pi_nP_n,
\]

а переход от времени \(r\) к времени \(s\) задается произведением
соответствующих ядер:

\[
P_{r,s}=P_rP_{r+1}\cdots P_{s-1},
\qquad
r<s.
\]

В однородном случае все \(P_n=P\), поэтому

\[
P_{r,s}=P^{s-r}.
\]

\subsubsection{22.10. Обобщение на непрерывное
время}\label{ux43eux431ux43eux431ux449ux435ux43dux438ux435-ux43dux430-ux43dux435ux43fux440ux435ux440ux44bux432ux43dux43eux435-ux432ux440ux435ux43cux44f}

Можно также встретить непрерывное время. Тогда переходные вероятности
записывают как

\[
p_{ij}(t)=P\{X_{s+t}=j\mid X_s=i\}
\]

в однородном случае, а матрицы перехода образуют семейство

\[
P(t)=(p_{ij}(t)),
\qquad
t\ge0.
\]

Тогда

\[
P(0)=I,
\]

строки \(P(t)\) суммируются в единицу, и уравнение Чепмена-Колмогорова
принимает вид

\[
P(t+s)=P(t)P(s).
\]

Для текущего курса дискретных случайных процессов обычно достаточно
дискретного времени, но полезно понимать, что это та же идея.

\subsubsection{22.11. Начальное распределение и конечномерные
распределения}\label{ux43dux430ux447ux430ux43bux44cux43dux43eux435-ux440ux430ux441ux43fux440ux435ux434ux435ux43bux435ux43dux438ux435-ux438-ux43aux43eux43dux435ux447ux43dux43eux43cux435ux440ux43dux44bux435-ux440ux430ux441ux43fux440ux435ux434ux435ux43bux435ux43dux438ux44f}

Переходная функция не является начальным распределением. Начальное
распределение отвечает на вопрос, где процесс находится в момент \(0\):

\[
\pi_0(i)=P\{X_0=i\}.
\]

Переходная матрица отвечает на вопрос, как процесс переходит из
состояния в состояние:

\[
p_{ij}=P\{X_{n+1}=j\mid X_n=i\}.
\]

Только вместе начальное распределение и переходные вероятности задают
конечномерные распределения цепи. В счетном однородном случае:

\[
P\{X_0=i_0,\ldots,X_n=i_n\}
=
\pi_0(i_0)p_{i_0i_1}p_{i_1i_2}\cdots p_{i_{n-1}i_n}.
\]

Именно эта формула связывает билет с Билетом 15: такие конечномерные
распределения согласованы, потому что строки \(P\) суммируются в
единицу. Значит по теореме Колмогорова можно построить процесс с
заданными \(\pi_0\) и \(P\).

\subsection{3. Основные
определения}\label{ux43eux441ux43dux43eux432ux43dux44bux435-ux43eux43fux440ux435ux434ux435ux43bux435ux43dux438ux44f-21}

\subsubsection{Пространство
состояний}\label{ux43fux440ux43eux441ux442ux440ux430ux43dux441ux442ux432ux43e-ux441ux43eux441ux442ux43eux44fux43dux438ux439}

Измеримое пространство

\[
(E,\mathcal E),
\]

значения в котором принимает процесс \(X_n\). В конечном или счетном
случае обычно берут

\[
E=\{1,2,\ldots,r\}
\]

или

\[
E=\mathbb Z,\ \mathbb N.
\]

\subsubsection{Переходная
функция}\label{ux43fux435ux440ux435ux445ux43eux434ux43dux430ux44f-ux444ux443ux43dux43aux446ux438ux44f}

Функция

\[
P(x,A),
\qquad
x\in E,\ A\in\mathcal E,
\]

которая задает вероятность перехода из состояния \(x\) в множество \(A\)
за один шаг.

\subsubsection{Марковское
ядро}\label{ux43cux430ux440ux43aux43eux432ux441ux43aux43eux435-ux44fux434ux440ux43e}

То же самое, что переходная функция, но с акцентом на измеримую
структуру:

\begin{itemize}
\tightlist
\item
  \(P(x,\cdot)\) - вероятностная мера;
\item
  \(P(\cdot,A)\) - измеримая функция.
\end{itemize}

\subsubsection{Переходная
вероятность}\label{ux43fux435ux440ux435ux445ux43eux434ux43dux430ux44f-ux432ux435ux440ux43eux44fux442ux43dux43eux441ux442ux44c}

В счетном случае:

\[
p_{ij}=P\{X_{n+1}=j\mid X_n=i\}.
\]

Это вероятность перехода из состояния \(i\) в состояние \(j\) за один
шаг.

\subsubsection{Стохастическая
матрица}\label{ux441ux442ux43eux445ux430ux441ux442ux438ux447ux435ux441ux43aux430ux44f-ux43cux430ux442ux440ux438ux446ux430}

Матрица

\[
P=(p_{ij})
\]

называется стохастической по строкам, если

\[
p_{ij}\ge0,
\qquad
\sum_jp_{ij}=1
\]

для каждой строки \(i\).

\subsubsection{Однородность по
времени}\label{ux43eux434ux43dux43eux440ux43eux434ux43dux43eux441ux442ux44c-ux43fux43e-ux432ux440ux435ux43cux435ux43dux438}

Цепь или переходный механизм называется однородным, если условный закон
перехода не зависит от времени:

\[
P\{X_{n+1}\in A\mid X_n=x\}=P(x,A)
\]

для всех \(n\).

В счетном случае:

\[
P\{X_{n+1}=j\mid X_n=i\}=p_{ij}.
\]

\subsubsection{Неоднородные
переходы}\label{ux43dux435ux43eux434ux43dux43eux440ux43eux434ux43dux44bux435-ux43fux435ux440ux435ux445ux43eux434ux44b}

Если переходы зависят от времени, пишут

\[
P_n(x,A)=P\{X_{n+1}\in A\mid X_n=x\}.
\]

В счетном случае:

\[
p_{ij}(n)=P\{X_{n+1}=j\mid X_n=i\}.
\]

\subsubsection{\texorpdfstring{\(n\)-шаговая переходная
функция}{n-шаговая переходная функция}}\label{n-ux448ux430ux433ux43eux432ux430ux44f-ux43fux435ux440ux435ux445ux43eux434ux43dux430ux44f-ux444ux443ux43dux43aux446ux438ux44f}

Для однородной цепи:

\[
P^{(n)}(x,A)=P\{X_n\in A\mid X_0=x\}.
\]

В счетном случае:

\[
p_{ij}^{(n)}=P\{X_n=j\mid X_0=i\}.
\]

\subsubsection{Начальное
распределение}\label{ux43dux430ux447ux430ux43bux44cux43dux43eux435-ux440ux430ux441ux43fux440ux435ux434ux435ux43bux435ux43dux438ux435}

Распределение состояния в момент \(0\):

\[
\pi_0(A)=P\{X_0\in A\}.
\]

В счетном случае:

\[
\pi_0(i)=P\{X_0=i\}.
\]

\subsection{4. Основные теоремы и
свойства}\label{ux43eux441ux43dux43eux432ux43dux44bux435-ux442ux435ux43eux440ux435ux43cux44b-ux438-ux441ux432ux43eux439ux441ux442ux432ux430-21}

\subsubsection{1. Стохастичность переходной
матрицы}\label{ux441ux442ux43eux445ux430ux441ux442ux438ux447ux43dux43eux441ux442ux44c-ux43fux435ux440ux435ux445ux43eux434ux43dux43eux439-ux43cux430ux442ux440ux438ux446ux44b}

Для любой переходной матрицы

\[
p_{ij}\ge0,
\qquad
\sum_jp_{ij}=1.
\]

Доказательство: при фиксированном \(i\) числа \(p_{ij}\) являются
вероятностями несовместных событий \(\{X_{n+1}=j\}\), которые образуют
полную группу исходов по следующему состоянию.

\subsubsection{2. Переходная функция переводит распределения в
распределения}\label{ux43fux435ux440ux435ux445ux43eux434ux43dux430ux44f-ux444ux443ux43dux43aux446ux438ux44f-ux43fux435ux440ux435ux432ux43eux434ux438ux442-ux440ux430ux441ux43fux440ux435ux434ux435ux43bux435ux43dux438ux44f-ux432-ux440ux430ux441ux43fux440ux435ux434ux435ux43bux435ux43dux438ux44f}

Если \(\mu\) - вероятностная мера на \((E,\mathcal E)\), то

\[
\mu P(A)=\int_E P(x,A)\,\mu(dx)
\]

тоже вероятностная мера.

Проверка:

\[
\mu P(A)\ge0,
\]

счетная аддитивность следует из счетной аддитивности \(P(x,\cdot)\) и
теоремы Леви или монотонной сходимости, а нормировка:

\[
\mu P(E)=\int_E P(x,E)\,\mu(dx)=\int_E1\,\mu(dx)=1.
\]

\subsubsection{3. Эволюция распределений в счетном
случае}\label{ux44dux432ux43eux43bux44eux446ux438ux44f-ux440ux430ux441ux43fux440ux435ux434ux435ux43bux435ux43dux438ux439-ux432-ux441ux447ux435ux442ux43dux43eux43c-ux441ux43bux443ux447ux430ux435}

Если

\[
\pi_n(i)=P\{X_n=i\},
\]

то

\[
\pi_{n+1}(j)
=
\sum_i\pi_n(i)p_{ij}.
\]

В матричной записи:

\[
\pi_{n+1}=\pi_nP.
\]

Доказательство - формула полной вероятности:

\[
P\{X_{n+1}=j\}
=
\sum_iP\{X_n=i\}P\{X_{n+1}=j\mid X_n=i\}.
\]

\subsubsection{4. Двухшаговые
переходы}\label{ux434ux432ux443ux445ux448ux430ux433ux43eux432ux44bux435-ux43fux435ux440ux435ux445ux43eux434ux44b}

В общем пространстве:

\[
P^{(2)}(x,A)
=
\int_E P(x,dy)P(y,A).
\]

В счетном пространстве:

\[
p_{ij}^{(2)}
=
\sum_k p_{ik}p_{kj}.
\]

Это обычное умножение матриц:

\[
P^{(2)}=P^2.
\]

\subsubsection{\texorpdfstring{5. \(n\)-шаговые
переходы}{5. n-шаговые переходы}}\label{n-ux448ux430ux433ux43eux432ux44bux435-ux43fux435ux440ux435ux445ux43eux434ux44b}

Для однородной цепи:

\[
P^{(n)}=P^n.
\]

В счетном случае:

\[
p_{ij}^{(n)}
=
(P^n)_{ij}.
\]

При этом

\[
\pi_n=\pi_0P^n.
\]

\subsubsection{6. Уравнение
Чепмена-Колмогорова}\label{ux443ux440ux430ux432ux43dux435ux43dux438ux435-ux447ux435ux43fux43cux435ux43dux430-ux43aux43eux43bux43cux43eux433ux43eux440ux43eux432ux430-1}

Для всех \(m,n\ge0\):

\[
P^{(m+n)}(x,A)
=
\int_E P^{(n)}(x,dy)P^{(m)}(y,A).
\]

В счетном случае:

\[
p_{ij}^{(m+n)}
=
\sum_kp_{ik}^{(n)}p_{kj}^{(m)}.
\]

В матричном виде:

\[
P^{(m+n)}=P^{(n)}P^{(m)}.
\]

Смысл: переход за \(m+n\) шагов раскладывается по промежуточному
состоянию через \(n\) шагов.

\subsubsection{7. Формула конечномерных
распределений}\label{ux444ux43eux440ux43cux443ux43bux430-ux43aux43eux43dux435ux447ux43dux43eux43cux435ux440ux43dux44bux445-ux440ux430ux441ux43fux440ux435ux434ux435ux43bux435ux43dux438ux439}

Для однородной цепи со счетным пространством состояний:

\[
P\{X_0=i_0,\ldots,X_n=i_n\}
=
\pi_0(i_0)p_{i_0i_1}\cdots p_{i_{n-1}i_n}.
\]

Эта формула показывает, что начальное распределение и переходная матрица
полностью задают закон цепи.

\subsubsection{8. Непрерывное время: матрицы
перехода}\label{ux43dux435ux43fux440ux435ux440ux44bux432ux43dux43eux435-ux432ux440ux435ux43cux44f-ux43cux430ux442ux440ux438ux446ux44b-ux43fux435ux440ux435ux445ux43eux434ux430}

Если рассматривается однородный марковский процесс с конечным числом
состояний и непрерывным временем, то

\[
P(t)=(p_{ij}(t)),
\qquad
p_{ij}(t)=P\{X_{s+t}=j\mid X_s=i\}.
\]

Тогда

\[
P(0)=I,
\qquad
p_{ij}(t)\ge0,
\qquad
\sum_jp_{ij}(t)=1,
\]

и

\[
P(t+s)=P(t)P(s).
\]

Это непрерывновременной аналог формулы \(P^{m+n}=P^mP^n\).

\subsection{5. Примеры}\label{ux43fux440ux438ux43cux435ux440ux44b-21}

\subsubsection{Пример 1. Двухсостоянийная
цепь}\label{ux43fux440ux438ux43cux435ux440-1.-ux434ux432ux443ux445ux441ux43eux441ux442ux43eux44fux43dux438ux439ux43dux430ux44f-ux446ux435ux43fux44c}

Пусть

\[
E=\{0,1\}.
\]

Матрица переходов:

\[
P=
\begin{pmatrix}
1-a & a\\
b & 1-b
\end{pmatrix},
\qquad
0\le a,b\le1.
\]

Из состояния \(0\) цепь переходит в \(1\) с вероятностью \(a\), а
остается в \(0\) с вероятностью \(1-a\). Из состояния \(1\) цепь
переходит в \(0\) с вероятностью \(b\), а остается в \(1\) с
вероятностью \(1-b\).

Строки суммируются в единицу:

\[
(1-a)+a=1,
\qquad
b+(1-b)=1.
\]

Если начальное распределение

\[
\pi_0=(\alpha,1-\alpha),
\]

то

\[
\pi_1=\pi_0P.
\]

\subsubsection{\texorpdfstring{Пример 2. Простое случайное блуждание на
\(\mathbb Z\)}{Пример 2. Простое случайное блуждание на \textbackslash mathbb Z}}\label{ux43fux440ux438ux43cux435ux440-2.-ux43fux440ux43eux441ux442ux43eux435-ux441ux43bux443ux447ux430ux439ux43dux43eux435-ux431ux43bux443ux436ux434ux430ux43dux438ux435-ux43dux430-mathbb-z}

Пусть

\[
E=\mathbb Z,
\]

и из состояния \(i\) цепь переходит в \(i+1\) с вероятностью \(p\), а в
\(i-1\) с вероятностью \(q=1-p\):

\[
p_{i,i+1}=p,
\qquad
p_{i,i-1}=q.
\]

Все остальные переходы равны нулю:

\[
p_{ij}=0,
\qquad
j\ne i+1,\ i-1.
\]

Это переходная матрица случайного блуждания из Билета 19. Она
бесконечная, но каждая строка имеет сумму

\[
p+q=1.
\]

\subsubsection{Пример 3. Поглощающее
состояние}\label{ux43fux440ux438ux43cux435ux440-3.-ux43fux43eux433ux43bux43eux449ux430ux44eux449ux435ux435-ux441ux43eux441ux442ux43eux44fux43dux438ux435}

Пусть

\[
E=\{0,1,2\}.
\]

Состояние \(0\) поглощающее:

\[
p_{00}=1,
\qquad
p_{01}=p_{02}=0.
\]

Например,

\[
P=
\begin{pmatrix}
1 & 0 & 0\\
1/3 & 1/3 & 1/3\\
1/2 & 0 & 1/2
\end{pmatrix}.
\]

Если цепь попала в \(0\), то она уже не выходит из \(0\). Это видно по
первой строке матрицы.

\subsubsection{Пример 4. Детерминированная
динамика}\label{ux43fux440ux438ux43cux435ux440-4.-ux434ux435ux442ux435ux440ux43cux438ux43dux438ux440ux43eux432ux430ux43dux43dux430ux44f-ux434ux438ux43dux430ux43cux438ux43aux430}

Пусть есть отображение

\[
T:E\to E.
\]

Если из состояния \(x\) система всегда переходит в \(T(x)\), то

\[
P(x,A)=\mathbf 1_A(Tx).
\]

В счетном случае:

\[
p_{ij}=
\begin{cases}
1, & j=T(i),\\
0, & j\ne T(i).
\end{cases}
\]

Это тоже марковское ядро. Случайность здесь может быть только в
начальном состоянии.

\subsubsection{Пример 5. Ядро с
плотностью}\label{ux43fux440ux438ux43cux435ux440-5.-ux44fux434ux440ux43e-ux441-ux43fux43bux43eux442ux43dux43eux441ux442ux44cux44e}

Пусть

\[
E=\mathbb R.
\]

Если при текущем состоянии \(x\) следующее состояние имеет плотность
\(p(x,y)\) по \(y\), то

\[
P(x,A)=\int_A p(x,y)\,dy.
\]

Для корректности нужно:

\[
p(x,y)\ge0,
\qquad
\int_{\mathbb R}p(x,y)\,dy=1
\]

для каждого \(x\), и измеримость по \(x\). Это непрерывный аналог
переходной матрицы.

\subsection{6. Схемы доказательств /
обоснований}\label{ux441ux445ux435ux43cux44b-ux434ux43eux43aux430ux437ux430ux442ux435ux43bux44cux441ux442ux432-ux43eux431ux43eux441ux43dux43eux432ux430ux43dux438ux439-14}

\textbf{Почему строки матрицы суммируются в единицу.}

Фиксируем текущее состояние \(i\). Тогда события

\[
\{X_{n+1}=j\},
\qquad
j\in E,
\]

образуют полную группу возможных значений следующего состояния. Поэтому

\[
\sum_jp_{ij}
=
\sum_jP\{X_{n+1}=j\mid X_n=i\}
=
1.
\]

\textbf{Почему \(\pi_{n+1}=\pi_nP\).}

Для каждого \(j\):

\[
\pi_{n+1}(j)
=
P\{X_{n+1}=j\}.
\]

Разбиваем по значениям \(X_n\):

\[
P\{X_{n+1}=j\}
=
\sum_iP\{X_n=i,\ X_{n+1}=j\}.
\]

Далее:

\[
P\{X_n=i,\ X_{n+1}=j\}
=
P\{X_n=i\}P\{X_{n+1}=j\mid X_n=i\}.
\]

Значит

\[
\pi_{n+1}(j)=\sum_i\pi_n(i)p_{ij}.
\]

Это и есть \(j\)-я координата строки \(\pi_nP\).

\textbf{Доказательство Чепмена-Колмогорова в счетном случае.}

Нужно найти вероятность перехода из \(i\) в \(j\) за \(m+n\) шагов:

\[
p_{ij}^{(m+n)}
=
P\{X_{m+n}=j\mid X_0=i\}.
\]

Разбиваем по промежуточному состоянию \(k\) в момент \(n\):

\[
p_{ij}^{(m+n)}
=
\sum_k
P\{X_n=k,\ X_{m+n}=j\mid X_0=i\}.
\]

По правилу умножения условных вероятностей:

\[
P\{X_n=k,\ X_{m+n}=j\mid X_0=i\}
=
P\{X_n=k\mid X_0=i\}
P\{X_{m+n}=j\mid X_n=k,\ X_0=i\}.
\]

По марковскому свойству и однородности будущее после момента \(n\)
зависит от прошлого только через \(X_n=k\):

\[
P\{X_{m+n}=j\mid X_n=k,\ X_0=i\}
=
P\{X_m=j\mid X_0=k\}
=
p_{kj}^{(m)}.
\]

Поэтому

\[
p_{ij}^{(m+n)}
=
\sum_k p_{ik}^{(n)}p_{kj}^{(m)}.
\]

\textbf{Доказательство Чепмена-Колмогорова для ядра.}

В общем пространстве состояний сумма по \(k\) заменяется интегралом по
промежуточному состоянию \(y\):

\[
P^{(m+n)}(x,A)
=
\int_E P^{(n)}(x,dy)P^{(m)}(y,A).
\]

Это та же формула полной вероятности, записанная для условных
распределений.

\textbf{Почему \(\pi_0\) и \(P\) задают конечномерные распределения.}

Для траектории

\[
i_0,i_1,\ldots,i_n
\]

используем цепное правило:

\[
P\{X_0=i_0,\ldots,X_n=i_n\}
=
P\{X_0=i_0\}
\prod_{k=1}^n
P\{X_k=i_k\mid X_0=i_0,\ldots,X_{k-1}=i_{k-1}\}.
\]

По марковскому свойству:

\[
P\{X_k=i_k\mid X_0=i_0,\ldots,X_{k-1}=i_{k-1}\}
=
p_{i_{k-1}i_k}.
\]

Значит

\[
P\{X_0=i_0,\ldots,X_n=i_n\}
=
\pi_0(i_0)p_{i_0i_1}\cdots p_{i_{n-1}i_n}.
\]

\subsection{7. Что написать на
доске}\label{ux447ux442ux43e-ux43dux430ux43fux438ux441ux430ux442ux44c-ux43dux430-ux434ux43eux441ux43aux435-21}

Переходное ядро:

\[
P(x,A)=P\{X_{n+1}\in A\mid X_n=x\}.
\]

Условия ядра:

\[
P(x,\cdot)\ \text{-- вероятность},
\qquad
P(\cdot,A)\ \text{-- измерима}.
\]

Матрица переходов:

\[
p_{ij}=P\{X_{n+1}=j\mid X_n=i\},
\qquad
p_{ij}\ge0,
\qquad
\sum_jp_{ij}=1.
\]

Эволюция распределения:

\[
\pi_{n+1}=\pi_nP,
\qquad
\pi_n=\pi_0P^n.
\]

Чепмен-Колмогоров:

\[
P^{(m+n)}(x,A)
=
\int_E P^{(n)}(x,dy)P^{(m)}(y,A).
\]

В счетном случае:

\[
p_{ij}^{(m+n)}
=
\sum_k p_{ik}^{(n)}p_{kj}^{(m)}.
\]

Формула траектории:

\[
P\{X_0=i_0,\ldots,X_n=i_n\}
=
\pi_0(i_0)p_{i_0i_1}\cdots p_{i_{n-1}i_n}.
\]

\subsection{8. Типичные дополнительные
вопросы}\label{ux442ux438ux43fux438ux447ux43dux44bux435-ux434ux43eux43fux43eux43bux43dux438ux442ux435ux43bux44cux43dux44bux435-ux432ux43eux43fux440ux43eux441ux44b-21}

\textbf{Что такое марковское ядро?}

Это функция \(P(x,A)\), которая при фиксированном \(x\) является
вероятностной мерой по \(A\), а при фиксированном \(A\) является
измеримой функцией по \(x\). Она задает условный закон следующего
состояния при текущем состоянии \(x\).

\textbf{Чем переходная матрица отличается от распределения цепи?}

Распределение \(\pi_n\) говорит, где цепь находится в момент \(n\).
Переходная матрица \(P\) говорит, как цепь переходит из одного состояния
в другое. Формула связи:

\[
\pi_{n+1}=\pi_nP.
\]

\textbf{Почему матрица переходов стохастическая именно по строкам?}

Потому что строка с номером \(i\) - это условное распределение
следующего состояния при условии \(X_n=i\):

\[
(p_{ij})_j.
\]

Сумма вероятностей всех возможных следующих состояний равна единице:

\[
\sum_jp_{ij}=1.
\]

\textbf{Что изменится при столбцовом соглашении?}

Если распределения записывать столбцами, то формула эволюции станет

\[
\pi_{n+1}=P^T\pi_n
\]

при той же строковой стохастичности. Либо можно использовать
столбцово-стохастические матрицы, но тогда нужно явно менять соглашение.
В этом конспекте используется строковое соглашение.

\textbf{Как получить переход за два шага?}

Нужно просуммировать по промежуточному состоянию:

\[
p_{ij}^{(2)}=\sum_kp_{ik}p_{kj}.
\]

Это элемент матрицы \(P^2\).

\textbf{Как формулируется уравнение Чепмена-Колмогорова?}

Для ядер:

\[
P^{(m+n)}(x,A)
=
\int_E P^{(n)}(x,dy)P^{(m)}(y,A).
\]

Для матриц:

\[
p_{ij}^{(m+n)}
=
\sum_kp_{ik}^{(n)}p_{kj}^{(m)}.
\]

\textbf{Почему Чепмен-Колмогоров следует из полной вероятности?}

Переход за \(m+n\) шагов можно разбить по состоянию \(k\) в
промежуточный момент \(n\). Суммируем вероятность попасть из \(i\) в
\(k\) за \(n\) шагов и затем из \(k\) в \(j\) за \(m\) шагов.

\textbf{Что значит однородность?}

Однородность означает, что переходный механизм не зависит от времени:

\[
P\{X_{n+1}\in A\mid X_n=x\}=P(x,A)
\]

для всех \(n\). В неоднородном случае нужно писать \(P_n(x,A)\).

\textbf{Можно ли задать марковскую цепь только матрицей переходов?}

Не полностью. Матрица задает переходный механизм, но еще нужно начальное
распределение \(\pi_0\). Вместе \(\pi_0\) и \(P\) задают все
конечномерные распределения.

\textbf{Почему случайное блуждание является примером переходной
матрицы?}

Потому что при известном текущем положении \(i\) следующее положение
зависит только от нового шага:

\[
p_{i,i+1}=p,
\qquad
p_{i,i-1}=q.
\]

Прошлая траектория не нужна для описания следующего перехода.

\subsection{9. Типичные
ошибки}\label{ux442ux438ux43fux438ux447ux43dux44bux435-ux43eux448ux438ux431ux43aux438-21}

\begin{itemize}
\tightlist
\item
  Делать стохастическими столбцы без оговорки, а потом использовать
  строковую формулу \(\pi_{n+1}=\pi_nP\).
\item
  Путать переходную матрицу \(P\) и распределение \(\pi_n\).
\item
  Забывать начальное распределение: одна матрица \(P\) не задает закон
  цепи полностью.
\item
  Писать \(p_{ij}^{(2)}=p_{ij}^2\). Правильно:
\end{itemize}

\[
p_{ij}^{(2)}=\sum_kp_{ik}p_{kj}.
\]

\begin{itemize}
\tightlist
\item
  В уравнении Чепмена-Колмогорова забывать сумму или интеграл по
  промежуточному состоянию.
\item
  Путать одношаговую вероятность \(p_{ij}\) и \(n\)-шаговую вероятность
  \(p_{ij}^{(n)}\).
\item
  Называть любую неотрицательную матрицу переходной. Нужна нормировка
  строк.
\item
  Путать однородность цепи с одинаковостью распределений \(X_n\).
  Однородность говорит о переходах, а не о том, что \(\pi_n\) постоянно.
\item
  Для общего пространства состояний забывать измеримость
  \(x\mapsto P(x,A)\).
\end{itemize}

\subsection{10. Связи с другими
билетами}\label{ux441ux432ux44fux437ux438-ux441-ux434ux440ux443ux433ux438ux43cux438-ux431ux438ux43bux435ux442ux430ux43cux438-21}

\begin{itemize}
\tightlist
\item
  Билет 05: переходное ядро требует измеримости по состоянию.
\item
  Билет 06: переходные функции задают условные распределения будущего
  состояния.
\item
  Билет 09: формулы с ядрами и распределениями опираются на
  интегрирование по произведениям и полную вероятность.
\item
  Билет 10: переходная функция является регулярным условным
  распределением.
\item
  Билет 15: по начальному распределению и переходным ядрам строятся
  согласованные конечномерные распределения.
\item
  Билет 16: марковская цепь - частный случай случайной
  последовательности.
\item
  Билет 19: случайное блуждание задается простой переходной матрицей.
\item
  Билет 23: переходные матрицы используются для определения и анализа
  марковских цепей.
\item
  Билет 24: для стационарных марковских цепей переходы влияют на
  корреляции и эргодичность.
\end{itemize}

\subsection{11. Минимум на
удовлетворительно}\label{ux43cux438ux43dux438ux43cux443ux43c-ux43dux430-ux443ux434ux43eux432ux43bux435ux442ux432ux43eux440ux438ux442ux435ux43bux44cux43dux43e-21}

Нужно сказать:

\begin{enumerate}
\def\labelenumi{\arabic{enumi}.}
\tightlist
\item
  Переходная функция:
\end{enumerate}

\[
P(x,A)=P\{X_{n+1}\in A\mid X_n=x\}.
\]

\begin{enumerate}
\def\labelenumi{\arabic{enumi}.}
\setcounter{enumi}{1}
\tightlist
\item
  В счетном случае:
\end{enumerate}

\[
p_{ij}=P\{X_{n+1}=j\mid X_n=i\}.
\]

\begin{enumerate}
\def\labelenumi{\arabic{enumi}.}
\setcounter{enumi}{2}
\tightlist
\item
  Матрица переходов стохастическая:
\end{enumerate}

\[
p_{ij}\ge0,
\qquad
\sum_jp_{ij}=1.
\]

\begin{enumerate}
\def\labelenumi{\arabic{enumi}.}
\setcounter{enumi}{3}
\tightlist
\item
  Распределения обновляются по формуле:
\end{enumerate}

\[
\pi_{n+1}=\pi_nP.
\]

\begin{enumerate}
\def\labelenumi{\arabic{enumi}.}
\setcounter{enumi}{4}
\tightlist
\item
  Чепмен-Колмогоров:
\end{enumerate}

\[
p_{ij}^{(m+n)}
=
\sum_kp_{ik}^{(n)}p_{kj}^{(m)}.
\]

\subsection{12. Минимум на
хорошо}\label{ux43cux438ux43dux438ux43cux443ux43c-ux43dux430-ux445ux43eux440ux43eux448ux43e-21}

Помимо минимума, нужно:

\begin{itemize}
\tightlist
\item
  объяснить, что \(P(x,\cdot)\) - мера, а \(P(\cdot,A)\) - измеримая
  функция;
\item
  различать однородные и неоднородные переходы;
\item
  вывести \(\pi_{n+1}=\pi_nP\) через полную вероятность;
\item
  вывести Чепмена-Колмогорова через промежуточное состояние;
\item
  привести пример: двухсостоянийная цепь, блуждание или поглощающее
  состояние.
\end{itemize}

\subsection{13. Что добавить для
отлично}\label{ux447ux442ux43e-ux434ux43eux431ux430ux432ux438ux442ux44c-ux434ux43bux44f-ux43eux442ux43bux438ux447ux43dux43e-21}

Для сильного ответа стоит добавить:

\begin{itemize}
\tightlist
\item
  формулу для общего ядра:
\end{itemize}

\[
\mu P(A)=\int_E P(x,A)\,\mu(dx);
\]

\begin{itemize}
\tightlist
\item
  формулу конечномерных распределений:
\end{itemize}

\[
P\{X_0=i_0,\ldots,X_n=i_n\}
=
\pi_0(i_0)p_{i_0i_1}\cdots p_{i_{n-1}i_n};
\]

\begin{itemize}
\tightlist
\item
  связь с теоремой Колмогорова: строки суммируются в единицу, поэтому
  конечномерные распределения согласованы;
\item
  замечание про непрерывное время:
\end{itemize}

\[
P(t+s)=P(t)P(s);
\]

\begin{itemize}
\tightlist
\item
  предупреждение про строковое и столбцовое соглашения.
\end{itemize}

\subsection{14. Устная версия ответа на 5
минут}\label{ux443ux441ux442ux43dux430ux44f-ux432ux435ux440ux441ux438ux44f-ux43eux442ux432ux435ux442ux430-ux43dux430-5-ux43cux438ux43dux443ux442-21}

Переходные функции и матрицы описывают вероятностный механизм движения
марковской системы. Пусть процесс \(X_n\) принимает значения в
пространстве состояний \((E,\mathcal E)\). Переходная функция задается
формулой

\[
P(x,A)=P\{X_{n+1}\in A\mid X_n=x\}.
\]

При фиксированном \(x\) это вероятностная мера по \(A\), а при
фиксированном \(A\) - измеримая функция по \(x\). Поэтому переходное
ядро можно применять к распределению \(\mu\):

\[
\mu P(A)=\int_E P(x,A)\,\mu(dx).
\]

Это распределение следующего состояния, если текущее состояние имеет
распределение \(\mu\).

В счетном пространстве состояний ядро задается матрицей

\[
p_{ij}=P\{X_{n+1}=j\mid X_n=i\}.
\]

Матрица стохастическая:

\[
p_{ij}\ge0,
\qquad
\sum_jp_{ij}=1.
\]

Здесь строки соответствуют текущему состоянию, а столбцы - следующему.
Если распределения записывать строками, то

\[
\pi_{n+1}=\pi_nP.
\]

Покомпонентно:

\[
\pi_{n+1}(j)=\sum_i\pi_n(i)p_{ij}.
\]

Переход за несколько шагов получается умножением матриц:

\[
P^{(n)}=P^n.
\]

Главная формула - Чепмен-Колмогоров:

\[
p_{ij}^{(m+n)}
=
\sum_kp_{ik}^{(n)}p_{kj}^{(m)}.
\]

Она получается из формулы полной вероятности: чтобы попасть из \(i\) в
\(j\) за \(m+n\) шагов, разбиваем траектории по промежуточному состоянию
\(k\) через \(n\) шагов.

Для общего пространства состояний сумма заменяется интегралом:

\[
P^{(m+n)}(x,A)
=
\int_E P^{(n)}(x,dy)P^{(m)}(y,A).
\]

Пример - случайное блуждание на \(\mathbb Z\):

\[
p_{i,i+1}=p,
\qquad
p_{i,i-1}=q,
\qquad
p+q=1.
\]

Все остальные переходы равны нулю. Это марковский переходный механизм:
чтобы узнать следующее положение, достаточно знать текущее.

В конце важно сказать, что переходная матрица не заменяет начальное
распределение. Закон цепи задается парой \(\pi_0\) и \(P\), а
конечномерные вероятности имеют вид

\[
P\{X_0=i_0,\ldots,X_n=i_n\}
=
\pi_0(i_0)p_{i_0i_1}\cdots p_{i_{n-1}i_n}.
\]

\subsection{15. Если осталось 2 минуты перед
экзаменом}\label{ux435ux441ux43bux438-ux43eux441ux442ux430ux43bux43eux441ux44c-2-ux43cux438ux43dux443ux442ux44b-ux43fux435ux440ux435ux434-ux44dux43aux437ux430ux43cux435ux43dux43eux43c-21}

Запомнить:

\[
P(x,A)=P\{X_{n+1}\in A\mid X_n=x\}.
\]

В счетном случае:

\[
p_{ij}=P\{X_{n+1}=j\mid X_n=i\},
\qquad
p_{ij}\ge0,
\qquad
\sum_jp_{ij}=1.
\]

Распределение:

\[
\pi_{n+1}=\pi_nP.
\]

\(n\)-шаговые переходы:

\[
P^{(n)}=P^n.
\]

Чепмен-Колмогоров:

\[
p_{ij}^{(m+n)}
=
\sum_kp_{ik}^{(n)}p_{kj}^{(m)}.
\]

Главная ловушка: строки, а не столбцы, являются условными
распределениями следующего состояния при фиксированном текущем
состоянии.

\newpage

\section{Билет 23. Марковские случайные цепи и их
характеристики}\label{ux431ux438ux43bux435ux442-23.-ux43cux430ux440ux43aux43eux432ux441ux43aux438ux435-ux441ux43bux443ux447ux430ux439ux43dux44bux435-ux446ux435ux43fux438-ux438-ux438ux445-ux445ux430ux440ux430ux43aux442ux435ux440ux438ux441ux442ux438ux43aux438}

Источники: Ширяев 2, гл. VIII, §1, §4-7; Сердобольская, лекция 5.

\subsection{1. Короткий
ответ}\label{ux43aux43eux440ux43eux442ux43aux438ux439-ux43eux442ux432ux435ux442-22}

Марковская цепь - это случайная последовательность

\[
X_0,X_1,X_2,\ldots
\]

со значениями в конечном или счетном пространстве состояний \(E\), у
которой условное распределение будущего при известном настоящем не
зависит от прошлого:

\[
P\{X_{n+1}=j\mid X_0=i_0,\ldots,X_n=i\}
=
P\{X_{n+1}=j\mid X_n=i\}.
\]

Если цепь однородна по времени, то

\[
P\{X_{n+1}=j\mid X_n=i\}=p_{ij},
\]

где

\[
P=(p_{ij})
\]

\begin{itemize}
\tightlist
\item
  стохастическая матрица переходов:
\end{itemize}

\[
p_{ij}\ge0,
\qquad
\sum_jp_{ij}=1.
\]

Закон цепи задается начальным распределением \(\pi_0\) и матрицей \(P\).
Для траектории:

\[
P\{X_0=i_0,\ldots,X_n=i_n\}
=
\pi_0(i_0)p_{i_0i_1}\cdots p_{i_{n-1}i_n}.
\]

Распределения развиваются по формуле

\[
\pi_{n+1}=\pi_nP,
\qquad
\pi_n=\pi_0P^n.
\]

Главные характеристики состояний: достижимость, сообщаемость, классы
сообщающихся состояний, неразложимость, период,
возвратность/невозвратность, положительная возвратность. Стационарное
распределение \(\pi\) удовлетворяет

\[
\pi=\pi P.
\]

Для конечной неразложимой цепи существует единственное стационарное
распределение \(\pi\). Если цепь дополнительно апериодична, то

\[
p_{ij}^{(n)}\to\pi_j.
\]

\subsection{2. Полный
ответ}\label{ux43fux43eux43bux43dux44bux439-ux43eux442ux432ux435ux442-22}

\subsubsection{2.1. Понятие марковской цепи и марковское
свойство}\label{ux43fux43eux43dux44fux442ux438ux435-ux43cux430ux440ux43aux43eux432ux441ux43aux43eux439-ux446ux435ux43fux438-ux438-ux43cux430ux440ux43aux43eux432ux441ux43aux43eux435-ux441ux432ux43eux439ux441ux442ux432ux43e}

Марковские цепи - основной класс случайных последовательностей с
зависимостью. В iid-последовательности все наблюдения независимы. В
марковской цепи зависимость есть, но она локальная: для прогноза
следующего состояния достаточно знать текущее состояние.

Пусть пространство состояний

\[
E=\{1,2,\ldots\}
\]

конечно или счетно. Последовательность случайных величин

\[
X=(X_n)_{n\ge0}
\]

со значениями в \(E\) называется марковской цепью, если для любых
допустимых состояний

\[
i_0,\ldots,i_n,j
\]

при условии, что условная вероятность определена, выполняется марковское
свойство:

\[
P\{X_{n+1}=j\mid X_0=i_0,\ldots,X_{n-1}=i_{n-1},X_n=i_n\}
=
P\{X_{n+1}=j\mid X_n=i_n\}.
\]

Интуиция: прошлое влияет на будущее только через настоящее состояние.
Если известно \(X_n\), то дополнительная информация о
\(X_0,\ldots,X_{n-1}\) не меняет условный закон \(X_{n+1}\).

\subsubsection{2.2. Однородность и матрица переходных
вероятностей}\label{ux43eux434ux43dux43eux440ux43eux434ux43dux43eux441ux442ux44c-ux438-ux43cux430ux442ux440ux438ux446ux430-ux43fux435ux440ux435ux445ux43eux434ux43dux44bux445-ux432ux435ux440ux43eux44fux442ux43dux43eux441ux442ux435ux439}

Цепь называется однородной по времени, если эти условные вероятности не
зависят от \(n\):

\[
P\{X_{n+1}=j\mid X_n=i\}=p_{ij}.
\]

Матрица

\[
P=(p_{ij})_{i,j\in E}
\]

называется матрицей переходных вероятностей. Она уже подробно
обсуждалась в Билете 22. Здесь важно помнить:

\[
p_{ij}\ge0,
\qquad
\sum_jp_{ij}=1.
\]

В строковом соглашении распределение в момент \(n\) записывается строкой

\[
\pi_n=(\pi_n(i))_{i\in E},
\qquad
\pi_n(i)=P\{X_n=i\}.
\]

Тогда

\[
\pi_{n+1}=\pi_nP,
\qquad
\pi_n=\pi_0P^n.
\]

Покомпонентно:

\[
\pi_{n+1}(j)=\sum_i\pi_n(i)p_{ij}.
\]

\subsubsection{2.3. Описание цепи и формула распределения
траекторий}\label{ux43eux43fux438ux441ux430ux43dux438ux435-ux446ux435ux43fux438-ux438-ux444ux43eux440ux43cux443ux43bux430-ux440ux430ux441ux43fux440ux435ux434ux435ux43bux435ux43dux438ux44f-ux442ux440ux430ux435ux43aux442ux43eux440ux438ux439}

Марковская цепь полностью задается двумя объектами:

\begin{enumerate}
\def\labelenumi{\arabic{enumi}.}
\tightlist
\item
  начальным распределением
\end{enumerate}

\[
\pi_0(i)=P\{X_0=i\};
\]

\begin{enumerate}
\def\labelenumi{\arabic{enumi}.}
\setcounter{enumi}{1}
\tightlist
\item
  матрицей переходов \(P=(p_{ij})\).
\end{enumerate}

Действительно, для любой конечной траектории

\[
i_0,i_1,\ldots,i_n
\]

имеет место формула:

\[
P\{X_0=i_0,\ldots,X_n=i_n\}
=
\pi_0(i_0)p_{i_0i_1}p_{i_1i_2}\cdots p_{i_{n-1}i_n}.
\]

Это ключевая формула билета. Она получается последовательным применением
правила умножения вероятностей и марковского свойства:

\[
P\{X_k=i_k\mid X_0=i_0,\ldots,X_{k-1}=i_{k-1}\}
=
p_{i_{k-1}i_k}.
\]

Важно не путать марковость с независимостью. В марковской цепи состояний
обычно зависимы. Например, если \(p_{ii}\) велико, то знание \(X_n=i\)
сильно влияет на распределение \(X_{n+1}\). Марковость говорит не об
отсутствии зависимости, а о специальной форме зависимости: будущее
условно независимо от прошлого при фиксированном настоящем.

\subsubsection{2.4. Переходные вероятности за n шагов и уравнение
Чепмена-Колмогорова}\label{ux43fux435ux440ux435ux445ux43eux434ux43dux44bux435-ux432ux435ux440ux43eux44fux442ux43dux43eux441ux442ux438-ux437ux430-n-ux448ux430ux433ux43eux432-ux438-ux443ux440ux430ux432ux43dux435ux43dux438ux435-ux447ux435ux43fux43cux435ux43dux430-ux43aux43eux43bux43cux43eux433ux43eux440ux43eux432ux430}

Для анализа цепи изучают не только одношаговые переходы \(p_{ij}\), но и
\(n\)-шаговые переходы:

\[
p_{ij}^{(n)}=P\{X_n=j\mid X_0=i\}.
\]

Для однородной цепи

\[
P^{(n)}=P^n,
\]

и выполнено уравнение Чепмена-Колмогорова:

\[
p_{ij}^{(m+n)}
=
\sum_kp_{ik}^{(n)}p_{kj}^{(m)}.
\]

\subsubsection{2.5. Классификация состояний: сообщаемость и
неразложимость}\label{ux43aux43bux430ux441ux441ux438ux444ux438ux43aux430ux446ux438ux44f-ux441ux43eux441ux442ux43eux44fux43dux438ux439-ux441ux43eux43eux431ux449ux430ux435ux43cux43eux441ux442ux44c-ux438-ux43dux435ux440ux430ux437ux43bux43eux436ux438ux43cux43eux441ux442ux44c}

Теперь перейдем к характеристикам состояний.

Состояние \(j\) называется достижимым из состояния \(i\), если
существует \(n\ge0\) такое, что

\[
p_{ij}^{(n)}>0.
\]

Пишут

\[
i\to j.
\]

Состояния \(i\) и \(j\) называются сообщающимися, если

\[
i\to j
\quad\text{и}\quad
j\to i.
\]

Пишут

\[
i\leftrightarrow j.
\]

Сообщаемость является отношением эквивалентности: она рефлексивна,
симметрична и транзитивна. Поэтому пространство состояний разбивается на
классы сообщающихся состояний.

Цепь называется неразложимой, если любые два состояния сообщаются:

\[
\forall i,j\in E
\quad
i\leftrightarrow j.
\]

Иными словами, из любого состояния можно попасть в любое другое с
положительной вероятностью за некоторое число шагов.

\subsubsection{2.6. Замкнутые классы и периодичность
состояний}\label{ux437ux430ux43cux43aux43dux443ux442ux44bux435-ux43aux43bux430ux441ux441ux44b-ux438-ux43fux435ux440ux438ux43eux434ux438ux447ux43dux43eux441ux442ux44c-ux441ux43eux441ux442ux43eux44fux43dux438ux439}

Класс состояний называется замкнутым, если из него нельзя выйти:

\[
i\in C,\ i\to j
\quad\Rightarrow\quad
j\in C.
\]

Если цепь попала в замкнутый класс, она остается в нем навсегда.
Одноточечный замкнутый класс называется поглощающим состоянием:

\[
p_{ii}=1.
\]

Еще одна важная характеристика - период. Период состояния \(i\):

\[
d(i)=\gcd\{n\ge1:p_{ii}^{(n)}>0\}.
\]

Если

\[
d(i)=1,
\]

состояние называется апериодическим. Если \(d(i)>1\), возвращения в
\(i\) возможны только в моменты времени, кратные некоторому числу или
лежащие в одном циклическом классе. В одном классе сообщающихся
состояний все состояния имеют один и тот же период.

\subsubsection{2.7. Возвратность и невозвратность: определения и
критерии}\label{ux432ux43eux437ux432ux440ux430ux442ux43dux43eux441ux442ux44c-ux438-ux43dux435ux432ux43eux437ux432ux440ux430ux442ux43dux43eux441ux442ux44c-ux43eux43fux440ux435ux434ux435ux43bux435ux43dux438ux44f-ux438-ux43aux440ux438ux442ux435ux440ux438ux438}

Теперь возвратность. Пусть

\[
\tau_i=\inf\{n\ge1:X_n=i\}
\]

\begin{itemize}
\tightlist
\item
  время первого возвращения в \(i\) после выхода из \(i\); если
  возвращения нет, считают \(\tau_i=\infty\). Если цепь стартует из
  \(i\), то вероятность когда-нибудь вернуться в \(i\) равна
\end{itemize}

\[
f_{ii}=P_i\{\tau_i<\infty\}.
\]

Состояние \(i\) называется возвратным, если

\[
f_{ii}=1.
\]

Состояние \(i\) называется невозвратным, или транзиентным, если

\[
f_{ii}<1.
\]

Интуитивно: возвратное состояние почти наверное посещается снова, если
из него стартовать. Невозвратное состояние с положительной вероятностью
покидается навсегда.

Критерий через \(n\)-шаговые вероятности:

\[
i\ \text{возвратно}
\quad\Longleftrightarrow\quad
\sum_{n=0}^{\infty}p_{ii}^{(n)}=\infty.
\]

И

\[
i\ \text{невозвратно}
\quad\Longleftrightarrow\quad
\sum_{n=0}^{\infty}p_{ii}^{(n)}<\infty.
\]

\subsubsection{2.8. Среднее время возвращения и типы
возвратности}\label{ux441ux440ux435ux434ux43dux435ux435-ux432ux440ux435ux43cux44f-ux432ux43eux437ux432ux440ux430ux449ux435ux43dux438ux44f-ux438-ux442ux438ux43fux44b-ux432ux43eux437ux432ux440ux430ux442ux43dux43eux441ux442ux438}

Для возвратного состояния вводят среднее время возвращения:

\[
\mu_i=E_i\tau_i.
\]

Если

\[
\mu_i<\infty,
\]

то состояние называется положительно возвратным. Если

\[
\mu_i=\infty,
\]

то оно называется нулевым возвратным.

Для конечных неразложимых цепей все состояния положительно возвратны.
Для счетных цепей это уже не всегда верно: например, симметричное
случайное блуждание на \(\mathbb Z\) возвратно, но не положительно
возвратно.

\subsubsection{2.9. Стационарное распределение и равновесный
режим}\label{ux441ux442ux430ux446ux438ux43eux43dux430ux440ux43dux43eux435-ux440ux430ux441ux43fux440ux435ux434ux435ux43bux435ux43dux438ux435-ux438-ux440ux430ux432ux43dux43eux432ux435ux441ux43dux44bux439-ux440ux435ux436ux438ux43c}

Стационарное распределение - это вероятностный вектор

\[
\pi=(\pi_i)_{i\in E}
\]

такой, что

\[
\pi_i\ge0,
\qquad
\sum_i\pi_i=1,
\]

и

\[
\pi=\pi P.
\]

Покомпонентно:

\[
\pi_j=\sum_i\pi_i p_{ij}.
\]

Если \(X_0\) имеет распределение \(\pi\), то

\[
\pi_1=\pi P=\pi.
\]

Значит все \(X_n\) имеют одно и то же распределение \(\pi\). Поэтому
стационарное распределение описывает равновесный режим цепи. Важно:
стационарное распределение цепи - это не то же самое, что стационарность
любой конкретной реализации. Это распределение, которое сохраняется
переходом.

\subsubsection{2.10. Предельное поведение и условия
эргодичности}\label{ux43fux440ux435ux434ux435ux43bux44cux43dux43eux435-ux43fux43eux432ux435ux434ux435ux43dux438ux435-ux438-ux443ux441ux43bux43eux432ux438ux44f-ux44dux440ux433ux43eux434ux438ux447ux43dux43eux441ux442ux438}

Основная предельная картина такая:

\begin{itemize}
\tightlist
\item
  если цепь конечна, неразложима и апериодична, то существует
  единственное стационарное распределение \(\pi\);
\item
  для любых начальных состояний \(i\) и любых состояний \(j\)
\end{itemize}

\[
p_{ij}^{(n)}\to\pi_j;
\]

\begin{itemize}
\tightlist
\item
  для любого начального распределения \(\pi_0\)
\end{itemize}

\[
\pi_0P^n\to\pi.
\]

В счетном случае нужны более тонкие условия. По Ширяеву: существование
единственного стационарного распределения связано с наличием
единственного неразложимого положительно возвратного класса, а
существование эргодического предельного распределения требует
неразложимости, положительной возвратности и апериодичности.

\subsection{3. Основные
определения}\label{ux43eux441ux43dux43eux432ux43dux44bux435-ux43eux43fux440ux435ux434ux435ux43bux435ux43dux438ux44f-22}

\subsubsection{Марковская
цепь}\label{ux43cux430ux440ux43aux43eux432ux441ux43aux430ux44f-ux446ux435ux43fux44c-1}

Последовательность \((X_n)_{n\ge0}\) со значениями в \(E\) называется
марковской цепью, если

\[
P\{X_{n+1}=j\mid X_0=i_0,\ldots,X_n=i_n\}
=
P\{X_{n+1}=j\mid X_n=i_n\}.
\]

\subsubsection{Однородная марковская
цепь}\label{ux43eux434ux43dux43eux440ux43eux434ux43dux430ux44f-ux43cux430ux440ux43aux43eux432ux441ux43aux430ux44f-ux446ux435ux43fux44c}

Цепь называется однородной, если

\[
P\{X_{n+1}=j\mid X_n=i\}=p_{ij}
\]

не зависит от \(n\).

\subsubsection{Матрица
переходов}\label{ux43cux430ux442ux440ux438ux446ux430-ux43fux435ux440ux435ux445ux43eux434ux43eux432}

Матрица

\[
P=(p_{ij}),
\qquad
p_{ij}=P\{X_{n+1}=j\mid X_n=i\}.
\]

Она стохастическая по строкам:

\[
p_{ij}\ge0,
\qquad
\sum_jp_{ij}=1.
\]

\subsubsection{Начальное
распределение}\label{ux43dux430ux447ux430ux43bux44cux43dux43eux435-ux440ux430ux441ux43fux440ux435ux434ux435ux43bux435ux43dux438ux435-1}

Вероятностный вектор

\[
\pi_0(i)=P\{X_0=i\}.
\]

\subsubsection{\texorpdfstring{Распределение в момент
\(n\)}{Распределение в момент n}}\label{ux440ux430ux441ux43fux440ux435ux434ux435ux43bux435ux43dux438ux435-ux432-ux43cux43eux43cux435ux43dux442-n}

\[
\pi_n(i)=P\{X_n=i\}.
\]

Для строкового соглашения:

\[
\pi_n=\pi_0P^n.
\]

\subsubsection{\texorpdfstring{\(n\)-шаговая переходная
вероятность}{n-шаговая переходная вероятность}}\label{n-ux448ux430ux433ux43eux432ux430ux44f-ux43fux435ux440ux435ux445ux43eux434ux43dux430ux44f-ux432ux435ux440ux43eux44fux442ux43dux43eux441ux442ux44c}

\[
p_{ij}^{(n)}=P\{X_n=j\mid X_0=i\}.
\]

\subsubsection{Достижимость}\label{ux434ux43eux441ux442ux438ux436ux438ux43cux43eux441ux442ux44c}

Состояние \(j\) достижимо из \(i\), если

\[
\exists n\ge0:\ p_{ij}^{(n)}>0.
\]

Обозначение:

\[
i\to j.
\]

\subsubsection{Сообщаемость}\label{ux441ux43eux43eux431ux449ux430ux435ux43cux43eux441ux442ux44c}

Состояния \(i\) и \(j\) сообщаются, если

\[
i\to j
\quad\text{и}\quad
j\to i.
\]

Обозначение:

\[
i\leftrightarrow j.
\]

\subsubsection{Класс сообщающихся
состояний}\label{ux43aux43bux430ux441ux441-ux441ux43eux43eux431ux449ux430ux44eux449ux438ux445ux441ux44f-ux441ux43eux441ux442ux43eux44fux43dux438ux439}

Класс эквивалентности по отношению сообщаемости.

\subsubsection{Неразложимая
цепь}\label{ux43dux435ux440ux430ux437ux43bux43eux436ux438ux43cux430ux44f-ux446ux435ux43fux44c}

Цепь называется неразложимой, если все состояния принадлежат одному
классу сообщаемости:

\[
\forall i,j\in E:\ i\leftrightarrow j.
\]

\subsubsection{Замкнутый
класс}\label{ux437ux430ux43cux43aux43dux443ux442ux44bux439-ux43aux43bux430ux441ux441}

Класс \(C\) замкнут, если из него нельзя выйти:

\[
i\in C,\ p_{ij}^{(n)}>0
\quad\Rightarrow\quad
j\in C.
\]

Для одношаговых переходов достаточно проверять:

\[
i\in C,\ p_{ij}>0
\quad\Rightarrow\quad
j\in C.
\]

\subsubsection{Поглощающее
состояние}\label{ux43fux43eux433ux43bux43eux449ux430ux44eux449ux435ux435-ux441ux43eux441ux442ux43eux44fux43dux438ux435}

Состояние \(i\) поглощающее, если

\[
p_{ii}=1.
\]

\subsubsection{Период
состояния}\label{ux43fux435ux440ux438ux43eux434-ux441ux43eux441ux442ux43eux44fux43dux438ux44f}

\[
d(i)=\gcd\{n\ge1:p_{ii}^{(n)}>0\}.
\]

Если \(d(i)=1\), состояние апериодично. Если \(d(i)>1\), состояние
периодично.

\subsubsection{Время первого
возвращения}\label{ux432ux440ux435ux43cux44f-ux43fux435ux440ux432ux43eux433ux43e-ux432ux43eux437ux432ux440ux430ux449ux435ux43dux438ux44f}

При старте из \(i\):

\[
\tau_i=\inf\{n\ge1:X_n=i\}.
\]

\subsubsection{Возвратное
состояние}\label{ux432ux43eux437ux432ux440ux430ux442ux43dux43eux435-ux441ux43eux441ux442ux43eux44fux43dux438ux435}

Состояние \(i\) возвратно, если

\[
P_i\{\tau_i<\infty\}=1.
\]

\subsubsection{Невозвратное
состояние}\label{ux43dux435ux432ux43eux437ux432ux440ux430ux442ux43dux43eux435-ux441ux43eux441ux442ux43eux44fux43dux438ux435}

Состояние \(i\) невозвратно, если

\[
P_i\{\tau_i<\infty\}<1.
\]

\subsubsection{Положительно возвратное
состояние}\label{ux43fux43eux43bux43eux436ux438ux442ux435ux43bux44cux43dux43e-ux432ux43eux437ux432ux440ux430ux442ux43dux43eux435-ux441ux43eux441ux442ux43eux44fux43dux438ux435}

Возвратное состояние \(i\) положительно возвратно, если

\[
E_i\tau_i<\infty.
\]

\subsubsection{Нулевое возвратное
состояние}\label{ux43dux443ux43bux435ux432ux43eux435-ux432ux43eux437ux432ux440ux430ux442ux43dux43eux435-ux441ux43eux441ux442ux43eux44fux43dux438ux435}

Возвратное состояние \(i\) нулевое возвратное, если

\[
E_i\tau_i=\infty.
\]

\subsubsection{Стационарное
распределение}\label{ux441ux442ux430ux446ux438ux43eux43dux430ux440ux43dux43eux435-ux440ux430ux441ux43fux440ux435ux434ux435ux43bux435ux43dux438ux435}

Вероятностный вектор \(\pi\) такой, что

\[
\pi=\pi P.
\]

Покомпонентно:

\[
\pi_j=\sum_i\pi_i p_{ij}.
\]

\subsubsection{Эргодическая цепь в конечном
случае}\label{ux44dux440ux433ux43eux434ux438ux447ux435ux441ux43aux430ux44f-ux446ux435ux43fux44c-ux432-ux43aux43eux43dux435ux447ux43dux43eux43c-ux441ux43bux443ux447ux430ux435}

Часто конечную цепь называют эргодической, если она неразложима и
апериодична. Тогда существует единственное стационарное распределение и

\[
p_{ij}^{(n)}\to\pi_j.
\]

\subsection{4. Основные теоремы и
свойства}\label{ux43eux441ux43dux43eux432ux43dux44bux435-ux442ux435ux43eux440ux435ux43cux44b-ux438-ux441ux432ux43eux439ux441ux442ux432ux430-22}

\subsubsection{1. Закон
траектории}\label{ux437ux430ux43aux43eux43d-ux442ux440ux430ux435ux43aux442ux43eux440ux438ux438}

Если цепь однородна, имеет начальное распределение \(\pi_0\) и матрицу
переходов \(P\), то

\[
P\{X_0=i_0,\ldots,X_n=i_n\}
=
\pi_0(i_0)\prod_{k=1}^n p_{i_{k-1}i_k}.
\]

\subsubsection{2. Эволюция
распределений}\label{ux44dux432ux43eux43bux44eux446ux438ux44f-ux440ux430ux441ux43fux440ux435ux434ux435ux43bux435ux43dux438ux439}

Для строковых распределений:

\[
\pi_{n+1}=\pi_nP,
\qquad
\pi_n=\pi_0P^n.
\]

\subsubsection{3.
Чепмен-Колмогоров}\label{ux447ux435ux43fux43cux435ux43d-ux43aux43eux43bux43cux43eux433ux43eux440ux43eux432}

Для \(n\)-шаговых переходов:

\[
p_{ij}^{(m+n)}
=
\sum_kp_{ik}^{(n)}p_{kj}^{(m)}.
\]

Эта формула доказывает, что

\[
P^{(n)}=P^n.
\]

\subsubsection{4. Сообщаемость - отношение
эквивалентности}\label{ux441ux43eux43eux431ux449ux430ux435ux43cux43eux441ux442ux44c---ux43eux442ux43dux43eux448ux435ux43dux438ux435-ux44dux43aux432ux438ux432ux430ux43bux435ux43dux442ux43dux43eux441ux442ux438}

Свойство

\[
i\leftrightarrow j
\]

рефлексивно, симметрично и транзитивно. Транзитивность следует из
Чепмена-Колмогорова: если \(p_{ij}^{(n)}>0\) и \(p_{jk}^{(m)}>0\), то

\[
p_{ik}^{(n+m)}
\ge
p_{ij}^{(n)}p_{jk}^{(m)}
>0.
\]

\subsubsection{5. В одном классе сообщаемости свойства
совпадают}\label{ux432-ux43eux434ux43dux43eux43c-ux43aux43bux430ux441ux441ux435-ux441ux43eux43eux431ux449ux430ux435ux43cux43eux441ux442ux438-ux441ux432ux43eux439ux441ux442ux432ux430-ux441ux43eux432ux43fux430ux434ux430ux44eux442}

Если \(i\leftrightarrow j\), то состояния \(i\) и \(j\) имеют один и тот
же период. Также возвратность и невозвратность являются свойствами
класса: в одном классе либо все состояния возвратны, либо все
невозвратны.

\subsubsection{6. Критерий
возвратности}\label{ux43aux440ux438ux442ux435ux440ux438ux439-ux432ux43eux437ux432ux440ux430ux442ux43dux43eux441ux442ux438}

Состояние \(i\) возвратно тогда и только тогда, когда

\[
\sum_{n=0}^{\infty}p_{ii}^{(n)}=\infty.
\]

Состояние \(i\) невозвратно тогда и только тогда, когда

\[
\sum_{n=0}^{\infty}p_{ii}^{(n)}<\infty.
\]

\subsubsection{7. Стационарное распределение сохраняется
переходом}\label{ux441ux442ux430ux446ux438ux43eux43dux430ux440ux43dux43eux435-ux440ux430ux441ux43fux440ux435ux434ux435ux43bux435ux43dux438ux435-ux441ux43eux445ux440ux430ux43dux44fux435ux442ux441ux44f-ux43fux435ux440ux435ux445ux43eux434ux43eux43c}

Если

\[
\pi=\pi P,
\]

и

\[
\pi_0=\pi,
\]

то

\[
\pi_n=\pi
\]

для всех \(n\).

\subsubsection{8. Конечная неразложимая
цепь}\label{ux43aux43eux43dux435ux447ux43dux430ux44f-ux43dux435ux440ux430ux437ux43bux43eux436ux438ux43cux430ux44f-ux446ux435ux43fux44c}

Если \(E\) конечно и цепь неразложима, то существует единственное
стационарное распределение \(\pi\), причем

\[
\pi_i>0
\]

для всех \(i\). Все состояния положительно возвратны.

\subsubsection{9. Конечная неразложимая апериодическая
цепь}\label{ux43aux43eux43dux435ux447ux43dux430ux44f-ux43dux435ux440ux430ux437ux43bux43eux436ux438ux43cux430ux44f-ux430ux43fux435ux440ux438ux43eux434ux438ux447ux435ux441ux43aux430ux44f-ux446ux435ux43fux44c}

Если \(E\) конечно, цепь неразложима и апериодична, то

\[
p_{ij}^{(n)}\to\pi_j
\]

для всех \(i,j\). То есть предельное распределение не зависит от
начального состояния.

\subsubsection{10. Счетный
случай}\label{ux441ux447ux435ux442ux43dux44bux439-ux441ux43bux443ux447ux430ux439}

Для счетной цепи существование единственного стационарного распределения
требует положительной возвратности соответствующего замкнутого
неразложимого класса. Для эргодического предела дополнительно нужна
апериодичность.

\subsection{5. Примеры}\label{ux43fux440ux438ux43cux435ux440ux44b-22}

\subsubsection{Пример 1. Двухсостоянийная
цепь}\label{ux43fux440ux438ux43cux435ux440-1.-ux434ux432ux443ux445ux441ux43eux441ux442ux43eux44fux43dux438ux439ux43dux430ux44f-ux446ux435ux43fux44c-1}

Пусть

\[
E=\{0,1\},
\]

и

\[
P=
\begin{pmatrix}
1-a & a\\
b & 1-b
\end{pmatrix},
\qquad
0<a,b<1.
\]

Из любого состояния можно попасть в любое другое:

\[
0\leftrightarrow1.
\]

Цепь неразложима. Так как

\[
p_{00}=1-a>0,
\qquad
p_{11}=1-b>0,
\]

то возврат в каждое состояние возможен за один шаг, и период равен
\(1\). Цепь апериодична.

Стационарное распределение \(\pi=(\pi_0,\pi_1)\) находится из

\[
\pi=\pi P,
\qquad
\pi_0+\pi_1=1.
\]

Получаем

\[
\pi_0a=\pi_1b,
\]

поэтому

\[
\pi_0=\frac{b}{a+b},
\qquad
\pi_1=\frac{a}{a+b}.
\]

\subsubsection{Пример 2. Поглощающее
состояние}\label{ux43fux440ux438ux43cux435ux440-2.-ux43fux43eux433ux43bux43eux449ux430ux44eux449ux435ux435-ux441ux43eux441ux442ux43eux44fux43dux438ux435}

Пусть

\[
P=
\begin{pmatrix}
1 & 0\\
q & 1-q
\end{pmatrix},
\qquad
0<q\le1.
\]

Состояние \(0\) поглощающее:

\[
p_{00}=1.
\]

Из состояния \(1\) можно попасть в \(0\), но из \(0\) нельзя вернуться в
\(1\). Поэтому цепь разложима. Стационарное распределение:

\[
\pi=(1,0).
\]

Состояние \(1\) невозвратно: если цепь ушла из \(1\) в \(0\), она уже не
вернется.

\subsubsection{\texorpdfstring{Пример 3. Простое случайное блуждание на
\(\mathbb Z\)}{Пример 3. Простое случайное блуждание на \textbackslash mathbb Z}}\label{ux43fux440ux438ux43cux435ux440-3.-ux43fux440ux43eux441ux442ux43eux435-ux441ux43bux443ux447ux430ux439ux43dux43eux435-ux431ux43bux443ux436ux434ux430ux43dux438ux435-ux43dux430-mathbb-z}

Пусть

\[
p_{i,i+1}=p,
\qquad
p_{i,i-1}=q=1-p.
\]

Если \(0<p<1\), все состояния сообщаются, цепь неразложима. Возвраты в
исходное состояние возможны только за четное число шагов, поэтому период
равен

\[
d=2.
\]

При \(p=q=1/2\) блуждание на \(\mathbb Z\) возвратно, но среднее время
возвращения бесконечно, поэтому оно нулевое возвратное и не имеет
стационарного вероятностного распределения.

При \(p\ne q\) блуждание на \(\mathbb Z\) невозвратно.

\subsubsection{Пример 4. Случайное блуждание на конечном
графе}\label{ux43fux440ux438ux43cux435ux440-4.-ux441ux43bux443ux447ux430ux439ux43dux43eux435-ux431ux43bux443ux436ux434ux430ux43dux438ux435-ux43dux430-ux43aux43eux43dux435ux447ux43dux43eux43c-ux433ux440ux430ux444ux435}

Пусть есть конечный связный неориентированный граф. Из вершины \(i\)
цепь переходит равновероятно в одну из соседних вершин:

\[
p_{ij}=
\begin{cases}
\frac{1}{\deg(i)}, & i\sim j,\\
0, & \text{иначе}.
\end{cases}
\]

Если граф связный, цепь неразложима. Стационарное распределение:

\[
\pi_i=\frac{\deg(i)}{\sum_k\deg(k)}.
\]

Если граф двудольный, цепь имеет период \(2\). Если в графе есть петля
или нарушена двудольность, часто получается апериодичность.

\subsubsection{Пример 5. Цепь
рождения-гибели}\label{ux43fux440ux438ux43cux435ux440-5.-ux446ux435ux43fux44c-ux440ux43eux436ux434ux435ux43dux438ux44f-ux433ux438ux431ux435ux43bux438}

На состояниях

\[
E=\{0,1,2,\ldots\}
\]

переходы возможны только в соседние состояния:

\[
p_{i,i+1}=p_i,
\qquad
p_{i,i-1}=q_i,
\qquad
p_{ii}=r_i,
\]

где

\[
p_i+q_i+r_i=1.
\]

Такие цепи моделируют численность очереди, популяции или запаса. Их
характеристики зависят от соотношения вероятностей рождения и гибели:
цепь может иметь стационарное распределение, а может уходить на
бесконечность.

\subsection{6. Схемы доказательств /
обоснований}\label{ux441ux445ux435ux43cux44b-ux434ux43eux43aux430ux437ux430ux442ux435ux43bux44cux441ux442ux432-ux43eux431ux43eux441ux43dux43eux432ux430ux43dux438ux439-15}

\textbf{Доказательство формулы траектории.}

По правилу умножения вероятностей:

\[
P\{X_0=i_0,\ldots,X_n=i_n\}
=
P\{X_0=i_0\}
\prod_{k=1}^n
P\{X_k=i_k\mid X_0=i_0,\ldots,X_{k-1}=i_{k-1}\}.
\]

По марковскому свойству:

\[
P\{X_k=i_k\mid X_0=i_0,\ldots,X_{k-1}=i_{k-1}\}
=
P\{X_k=i_k\mid X_{k-1}=i_{k-1}\}
=
p_{i_{k-1}i_k}.
\]

Значит

\[
P\{X_0=i_0,\ldots,X_n=i_n\}
=
\pi_0(i_0)\prod_{k=1}^n p_{i_{k-1}i_k}.
\]

\textbf{Доказательство \(\pi_{n+1}=\pi_nP\).}

Для каждого \(j\):

\[
\pi_{n+1}(j)
=
P\{X_{n+1}=j\}
=
\sum_iP\{X_n=i,\ X_{n+1}=j\}
\]

\[
=
\sum_iP\{X_n=i\}P\{X_{n+1}=j\mid X_n=i\}
=
\sum_i\pi_n(i)p_{ij}.
\]

Это ровно \(j\)-я координата строки \(\pi_nP\).

\textbf{Доказательство транзитивности сообщаемости.}

Пусть

\[
i\to j,
\qquad
j\to k.
\]

Тогда существуют \(n,m\) такие, что

\[
p_{ij}^{(n)}>0,
\qquad
p_{jk}^{(m)}>0.
\]

По Чепмену-Колмогорову:

\[
p_{ik}^{(n+m)}
=
\sum_l p_{il}^{(n)}p_{lk}^{(m)}
\ge
p_{ij}^{(n)}p_{jk}^{(m)}
>0.
\]

Значит

\[
i\to k.
\]

Аналогично доказывается обратная достижимость, если есть сообщаемость.

\textbf{Идея критерия возвратности.}

Величина

\[
\sum_{n=0}^{\infty}p_{ii}^{(n)}
\]

равна математическому ожиданию числа посещений состояния \(i\) при
старте из \(i\):

\[
E_i\sum_{n=0}^{\infty}\mathbf 1_{\{X_n=i\}}
=
\sum_{n=0}^{\infty}P_i\{X_n=i\}
=
\sum_{n=0}^{\infty}p_{ii}^{(n)}.
\]

Если состояние невозвратно, число возвращений имеет конечное среднее.
Если состояние возвратно, посещений бесконечно много почти наверное, и
сумма расходится.

\textbf{Почему стационарное распределение сохраняется.}

Если

\[
\pi=\pi P,
\]

и

\[
\pi_0=\pi,
\]

то

\[
\pi_1=\pi_0P=\pi P=\pi.
\]

По индукции:

\[
\pi_n=\pi
\]

для всех \(n\).

\textbf{Идея предельной теоремы для конечной эргодической цепи.}

Неразложимость означает, что цепь не распадается на независимые
замкнутые части: из любого состояния можно добраться до любого другого.
Положительная возвратность в конечном случае гарантирует существование
равновесных долей времени. Апериодичность убирает осцилляции по
четным/нечетным моментам. Поэтому матрица \(P^n\) стабилизируется по
строкам:

\[
(P^n)_{ij}\to\pi_j.
\]

Все строки предельной матрицы одинаковы и равны стационарному
распределению.

\subsection{7. Что написать на
доске}\label{ux447ux442ux43e-ux43dux430ux43fux438ux441ux430ux442ux44c-ux43dux430-ux434ux43eux441ux43aux435-22}

Марковское свойство:

\[
P\{X_{n+1}=j\mid X_0=i_0,\ldots,X_n=i\}
=
P\{X_{n+1}=j\mid X_n=i\}
=p_{ij}.
\]

Матрица переходов:

\[
p_{ij}\ge0,
\qquad
\sum_jp_{ij}=1.
\]

Закон траектории:

\[
P\{X_0=i_0,\ldots,X_n=i_n\}
=
\pi_0(i_0)p_{i_0i_1}\cdots p_{i_{n-1}i_n}.
\]

Эволюция распределений:

\[
\pi_{n+1}=\pi_nP,
\qquad
\pi_n=\pi_0P^n.
\]

Достижимость и сообщаемость:

\[
i\to j
\Longleftrightarrow
\exists n:\ p_{ij}^{(n)}>0,
\qquad
i\leftrightarrow j
\Longleftrightarrow
i\to j,\ j\to i.
\]

Период:

\[
d(i)=\gcd\{n\ge1:p_{ii}^{(n)}>0\}.
\]

Возвратность:

\[
P_i\{\tau_i<\infty\}=1,
\qquad
\tau_i=\inf\{n\ge1:X_n=i\}.
\]

Стационарное распределение:

\[
\pi=\pi P,
\qquad
\sum_i\pi_i=1.
\]

Предельная формула для конечной эргодической цепи:

\[
p_{ij}^{(n)}\to\pi_j.
\]

\subsection{8. Типичные дополнительные
вопросы}\label{ux442ux438ux43fux438ux447ux43dux44bux435-ux434ux43eux43fux43eux43bux43dux438ux442ux435ux43bux44cux43dux44bux435-ux432ux43eux43fux440ux43eux441ux44b-22}

\textbf{Марковская цепь - это независимая последовательность?}

Нет. Марковская цепь обычно зависима. Марковское свойство означает
только условную независимость будущего от прошлого при фиксированном
настоящем.

\textbf{Что нужно, чтобы задать закон марковской цепи?}

Нужно начальное распределение \(\pi_0\) и переходная матрица \(P\).
Тогда все конечномерные распределения задаются формулой

\[
P\{X_0=i_0,\ldots,X_n=i_n\}
=
\pi_0(i_0)\prod_{k=1}^n p_{i_{k-1}i_k}.
\]

\textbf{Чем отличается стационарное распределение от стационарной цепи?}

Стационарное распределение \(\pi\) удовлетворяет \(\pi=\pi P\). Если
начальное распределение равно \(\pi\), то процесс становится
стационарным в узком смысле по сдвигам времени. Но сама матрица \(P\)
может иметь стационарное распределение, даже если цепь стартовала не из
него.

\textbf{Что такое неразложимость?}

Это возможность добраться из любого состояния в любое другое:

\[
\forall i,j\quad \exists n:\ p_{ij}^{(n)}>0.
\]

\textbf{Что такое период и зачем он нужен?}

Период состояния \(i\):

\[
d(i)=\gcd\{n\ge1:p_{ii}^{(n)}>0\}.
\]

Если период больше \(1\), вероятности переходов могут не сходиться, а
колебаться. Например, простое блуждание на двудольном графе возвращается
в исходную долю только в четные моменты.

\textbf{Что такое возвратность?}

Состояние возвратно, если при старте из него цепь почти наверное
когда-нибудь вернется:

\[
P_i\{\tau_i<\infty\}=1.
\]

\textbf{Что такое положительная возвратность?}

Это возвратность с конечным средним временем возвращения:

\[
E_i\tau_i<\infty.
\]

Именно положительная возвратность связана с существованием стационарного
распределения.

\textbf{Почему для конечной неразложимой цепи есть стационарное
распределение?}

В конечном случае вероятностный симплекс компактен, а переход
\(\pi\mapsto\pi P\) переводит его в себя. При неразложимости
стационарное распределение единственно и имеет все положительные
координаты.

\textbf{Зачем нужна апериодичность в предельной теореме?}

Без апериодичности цепь может колебаться между циклическими подклассами.
Например, если цепь каждый шаг переходит из \(0\) в \(1\) и из \(1\) в
\(0\), то стационарное распределение есть, но \(P^n\) не сходится.

\textbf{Как связана марковская цепь со случайным блужданием?}

Случайное блуждание

\[
S_{n+1}=S_n+\xi_{n+1}
\]

является марковской цепью, потому что следующий шаг независим от
прошлого, а условное распределение \(S_{n+1}\) зависит только от
\(S_n\).

\subsection{9. Типичные
ошибки}\label{ux442ux438ux43fux438ux447ux43dux44bux435-ux43eux448ux438ux431ux43aux438-22}

\begin{itemize}
\tightlist
\item
  Считать марковость независимостью состояний.
\item
  Забывать начальное распределение \(\pi_0\).
\item
  Путать матрицу переходов \(P\) и распределение \(\pi_n\).
\item
  Путать стационарность цепи и стационарное распределение.
\item
  Писать \(\pi P=\pi\) без условия \(\sum_i\pi_i=1\) и \(\pi_i\ge0\).
\item
  Называть цепь эргодической только из-за существования стационарного
  распределения. Для сходимости \(p_{ij}^{(n)}\to\pi_j\) нужна еще
  апериодичность и отсутствие разложения.
\item
  Забывать, что в счетном случае не всякая неразложимая цепь имеет
  стационарное вероятностное распределение.
\item
  Путать возвратность и положительную возвратность.
\item
  Считать, что если состояние достижимо из другого, то обратная
  достижимость автоматическая. Это неверно.
\item
  Забывать, что период является свойством класса сообщающихся состояний.
\end{itemize}

\subsection{10. Связи с другими
билетами}\label{ux441ux432ux44fux437ux438-ux441-ux434ux440ux443ux433ux438ux43cux438-ux431ux438ux43bux435ux442ux430ux43cux438-22}

\begin{itemize}
\tightlist
\item
  Билет 10: марковское свойство формулируется через условные вероятности
  и условные распределения.
\item
  Билет 15: начальное распределение и матрица переходов задают
  согласованные конечномерные распределения.
\item
  Билет 16: марковская цепь - частный случай дискретного случайного
  процесса.
\item
  Билет 17: если цепь стартует из стационарного распределения, она
  стационарна в узком смысле.
\item
  Билет 19: случайное блуждание - важный пример марковской цепи.
\item
  Билет 21: эргодические и предельные свойства цепей являются аналогами
  закона больших чисел для зависимых последовательностей.
\item
  Билет 22: переходные функции и матрицы - техническая основа марковских
  цепей.
\item
  Билет 24: эргодичность стационарных последовательностей и
  корреляционные условия связаны с предельным поведением марковских
  цепей.
\end{itemize}

\subsection{11. Минимум на
удовлетворительно}\label{ux43cux438ux43dux438ux43cux443ux43c-ux43dux430-ux443ux434ux43eux432ux43bux435ux442ux432ux43eux440ux438ux442ux435ux43bux44cux43dux43e-22}

Нужно сказать:

\begin{enumerate}
\def\labelenumi{\arabic{enumi}.}
\tightlist
\item
  Марковское свойство:
\end{enumerate}

\[
P\{X_{n+1}=j\mid X_0=i_0,\ldots,X_n=i\}
=
P\{X_{n+1}=j\mid X_n=i\}.
\]

\begin{enumerate}
\def\labelenumi{\arabic{enumi}.}
\setcounter{enumi}{1}
\tightlist
\item
  В однородном случае:
\end{enumerate}

\[
P\{X_{n+1}=j\mid X_n=i\}=p_{ij}.
\]

\begin{enumerate}
\def\labelenumi{\arabic{enumi}.}
\setcounter{enumi}{2}
\tightlist
\item
  Закон траектории:
\end{enumerate}

\[
P\{X_0=i_0,\ldots,X_n=i_n\}
=
\pi_0(i_0)p_{i_0i_1}\cdots p_{i_{n-1}i_n}.
\]

\begin{enumerate}
\def\labelenumi{\arabic{enumi}.}
\setcounter{enumi}{3}
\tightlist
\item
  Эволюция распределений:
\end{enumerate}

\[
\pi_{n+1}=\pi_nP.
\]

\begin{enumerate}
\def\labelenumi{\arabic{enumi}.}
\setcounter{enumi}{4}
\tightlist
\item
  Стационарное распределение:
\end{enumerate}

\[
\pi=\pi P.
\]

\subsection{12. Минимум на
хорошо}\label{ux43cux438ux43dux438ux43cux443ux43c-ux43dux430-ux445ux43eux440ux43eux448ux43e-22}

Помимо минимума, нужно:

\begin{itemize}
\tightlist
\item
  объяснить, почему марковость не равна независимости;
\item
  определить достижимость и сообщаемость;
\item
  сказать, что классы сообщаемости разбивают пространство состояний;
\item
  определить неразложимость и период;
\item
  привести пример поглощающей цепи или случайного блуждания;
\item
  назвать условие сходимости для конечной цепи: неразложимость плюс
  апериодичность.
\end{itemize}

\subsection{13. Что добавить для
отлично}\label{ux447ux442ux43e-ux434ux43eux431ux430ux432ux438ux442ux44c-ux434ux43bux44f-ux43eux442ux43bux438ux447ux43dux43e-22}

Для сильного ответа стоит добавить:

\begin{itemize}
\tightlist
\item
  возвратность и невозвратность через \(P_i\{\tau_i<\infty\}\);
\item
  критерий возвратности:
\end{itemize}

\[
\sum_{n=0}^{\infty}p_{ii}^{(n)}=\infty;
\]

\begin{itemize}
\tightlist
\item
  положительную и нулевую возвратность;
\item
  точную роль стационарного распределения;
\item
  формулировку предельной теоремы:
\end{itemize}

\[
p_{ij}^{(n)}\to\pi_j;
\]

\begin{itemize}
\tightlist
\item
  предупреждение, что в счетном случае неразложимости и апериодичности
  недостаточно без положительной возвратности.
\end{itemize}

\subsection{14. Устная версия ответа на 5
минут}\label{ux443ux441ux442ux43dux430ux44f-ux432ux435ux440ux441ux438ux44f-ux43eux442ux432ux435ux442ux430-ux43dux430-5-ux43cux438ux43dux443ux442-22}

Марковская цепь - это случайная последовательность \(X_0,X_1,\ldots\) со
значениями в конечном или счетном множестве состояний \(E\), для которой
будущее при известном настоящем не зависит от прошлого:

\[
P\{X_{n+1}=j\mid X_0=i_0,\ldots,X_n=i\}
=
P\{X_{n+1}=j\mid X_n=i\}.
\]

Если цепь однородна, то правая часть не зависит от \(n\):

\[
P\{X_{n+1}=j\mid X_n=i\}=p_{ij}.
\]

Матрица \(P=(p_{ij})\) стохастическая:

\[
p_{ij}\ge0,
\qquad
\sum_jp_{ij}=1.
\]

Цепь задается начальным распределением \(\pi_0\) и матрицей \(P\). Для
любой траектории:

\[
P\{X_0=i_0,\ldots,X_n=i_n\}
=
\pi_0(i_0)p_{i_0i_1}\cdots p_{i_{n-1}i_n}.
\]

Распределение в момент \(n\) считается по формуле

\[
\pi_n=\pi_0P^n.
\]

Дальше основные характеристики. Состояние \(j\) достижимо из \(i\), если
\(p_{ij}^{(n)}>0\) для некоторого \(n\). Если \(i\) достижимо из \(j\) и
\(j\) достижимо из \(i\), состояния сообщаются. Сообщаемость разбивает
пространство состояний на классы. Если все состояния сообщаются, цепь
неразложима.

Период состояния:

\[
d(i)=\gcd\{n\ge1:p_{ii}^{(n)}>0\}.
\]

Если \(d(i)=1\), состояние апериодично. Период важен для сходимости: при
периоде \(2\) цепь может бесконечно колебаться между двумя частями
пространства.

Возвратность: при старте из \(i\) вводим

\[
\tau_i=\inf\{n\ge1:X_n=i\}.
\]

Состояние возвратно, если

\[
P_i\{\tau_i<\infty\}=1,
\]

и невозвратно, если эта вероятность меньше единицы. Если возвратное
состояние имеет конечное среднее время возвращения, оно положительно
возвратно.

Стационарное распределение - это вероятностный вектор \(\pi\), который
сохраняется переходом:

\[
\pi=\pi P.
\]

Если цепь стартует из \(\pi\), то все \(X_n\) имеют распределение
\(\pi\). Для конечной неразложимой апериодической цепи стационарное
распределение единственно и

\[
p_{ij}^{(n)}\to\pi_j.
\]

Главные ошибки: считать марковость независимостью, забывать начальное
распределение и путать стационарное распределение со стационарностью
цепи.

\subsection{15. Если осталось 2 минуты перед
экзаменом}\label{ux435ux441ux43bux438-ux43eux441ux442ux430ux43bux43eux441ux44c-2-ux43cux438ux43dux443ux442ux44b-ux43fux435ux440ux435ux434-ux44dux43aux437ux430ux43cux435ux43dux43eux43c-22}

Запомнить:

\[
P\{X_{n+1}=j\mid X_0=i_0,\ldots,X_n=i\}
=p_{ij}.
\]

Закон траектории:

\[
P\{X_0=i_0,\ldots,X_n=i_n\}
=
\pi_0(i_0)p_{i_0i_1}\cdots p_{i_{n-1}i_n}.
\]

Распределения:

\[
\pi_{n+1}=\pi_nP,
\qquad
\pi_n=\pi_0P^n.
\]

Стационарное распределение:

\[
\pi=\pi P.
\]

Характеристики:

\[
i\to j \Leftrightarrow \exists n:\ p_{ij}^{(n)}>0,
\qquad
d(i)=\gcd\{n:p_{ii}^{(n)}>0\}.
\]

Возвратность:

\[
P_i\{\tau_i<\infty\}=1.
\]

Для конечной неразложимой апериодической цепи:

\[
p_{ij}^{(n)}\to\pi_j.
\]

Главная ловушка: марковская цепь - не независимая последовательность;
зависимость есть, но она проходит через текущее состояние.

\newpage

\section{Билет 24. Корреляционная функция. Достаточные условия
эргодичности
последовательности}\label{ux431ux438ux43bux435ux442-24.-ux43aux43eux440ux440ux435ux43bux44fux446ux438ux43eux43dux43dux430ux44f-ux444ux443ux43dux43aux446ux438ux44f.-ux434ux43eux441ux442ux430ux442ux43eux447ux43dux44bux435-ux443ux441ux43bux43eux432ux438ux44f-ux44dux440ux433ux43eux434ux438ux447ux43dux43eux441ux442ux438-ux43fux43eux441ux43bux435ux434ux43eux432ux430ux442ux435ux43bux44cux43dux43eux441ux442ux438}

Источники: Ширяев 2, гл. V-VI; рукописные лекции ДСП, лекции 9-10;
Сердобольская, разделы о ковариационной функции.

\subsection{1. Короткий
ответ}\label{ux43aux43eux440ux43eux442ux43aux438ux439-ux43eux442ux432ux435ux442-23}

Для случайной последовательности со вторыми моментами ковариационная
функция описывает линейную зависимость значений в разные моменты
времени. Если последовательность стационарна в широком смысле и

\[
\mathbb E X_n=m,
\]

то ее ковариационная функция зависит только от лага:

\[
R(k)=\operatorname{Cov}(X_{n+k},X_n)
=
\mathbb E[(X_{n+k}-m)\overline{(X_n-m)}].
\]

Если \(R(0)>0\), нормированная корреляционная функция:

\[
\rho(k)=\frac{R(k)}{R(0)}.
\]

Эргодичность по среднему в широком смысле означает

\[
\overline X_N=\frac1N\sum_{n=0}^{N-1}X_n
\xrightarrow[L^2]{N\to\infty}
m.
\]

Главная формула:

\[
\mathbb E|\overline X_N-m|^2
=
\frac1{N^2}\sum_{i=0}^{N-1}\sum_{j=0}^{N-1}R(i-j)
=
\frac1N\sum_{|k|<N}\left(1-\frac{|k|}{N}\right)R(k).
\]

Поэтому достаточное условие эргодичности:

\[
\frac1N\sum_{|k|<N}\left(1-\frac{|k|}{N}\right)R(k)\to0.
\]

Более простые достаточные условия:

\[
R(k)\to0
\]

или

\[
\sum_{k\in\mathbb Z}|R(k)|<\infty.
\]

\subsection{2. Полный
ответ}\label{ux43fux43eux43bux43dux44bux439-ux43eux442ux432ux435ux442-23}

\subsubsection{2.1. Введение и связь с предыдущими
темами}\label{ux432ux432ux435ux434ux435ux43dux438ux435-ux438-ux441ux432ux44fux437ux44c-ux441-ux43fux440ux435ux434ux44bux434ux443ux449ux438ux43cux438-ux442ux435ux43cux430ux43cux438}

Этот билет продолжает Билет 17 о стационарности и эргодичности. Там были
даны определения стационарности в узком и широком смыслах. Здесь нужно
разобрать, как эргодичность среднего проверяется через корреляционную,
точнее ковариационную, функцию.

\subsubsection{2.2. Определение ковариации и стационарность в широком
смысле}\label{ux43eux43fux440ux435ux434ux435ux43bux435ux43dux438ux435-ux43aux43eux432ux430ux440ux438ux430ux446ux438ux438-ux438-ux441ux442ux430ux446ux438ux43eux43dux430ux440ux43dux43eux441ux442ux44c-ux432-ux448ux438ux440ux43eux43aux43eux43c-ux441ux43cux44bux441ux43bux435}

Пусть

\[
X=(X_n)_{n\in\mathbb Z}
\]

\begin{itemize}
\tightlist
\item
  случайная последовательность с конечными вторыми моментами:
\end{itemize}

\[
\mathbb E|X_n|^2<\infty.
\]

Для двух случайных величин \(X_n\) и \(X_m\) ковариация равна

\[
\operatorname{Cov}(X_n,X_m)
=
\mathbb E[(X_n-\mathbb EX_n)\overline{(X_m-\mathbb EX_m)}].
\]

Для действительных величин комплексное сопряжение можно не писать.

Если последовательность стационарна в широком смысле, то:

\[
\mathbb EX_n=m
\]

не зависит от \(n\), а ковариация зависит только от разности индексов:

\[
\operatorname{Cov}(X_n,X_m)=R(n-m).
\]

Обычно фиксируют

\[
R(k)=\operatorname{Cov}(X_k,X_0).
\]

Если процесс центрирован, то есть \(m=0\), то

\[
R(k)=\mathbb E X_k\overline{X_0}.
\]

В рукописных лекциях это записывается как скалярное произведение в
гильбертовом пространстве \(L_2\):

\[
R(k)=(X_k,X_0)_{L_2}.
\]

\subsubsection{2.3. Различие между ковариационной и корреляционной
функциями}\label{ux440ux430ux437ux43bux438ux447ux438ux435-ux43cux435ux436ux434ux443-ux43aux43eux432ux430ux440ux438ux430ux446ux438ux43eux43dux43dux43eux439-ux438-ux43aux43eux440ux440ux435ux43bux44fux446ux438ux43eux43dux43dux43eux439-ux444ux443ux43dux43aux446ux438ux44fux43cux438}

Важно различать две функции:

\begin{enumerate}
\def\labelenumi{\arabic{enumi}.}
\tightlist
\item
  ковариационная функция:
\end{enumerate}

\[
R(k)=\operatorname{Cov}(X_k,X_0);
\]

\begin{enumerate}
\def\labelenumi{\arabic{enumi}.}
\setcounter{enumi}{1}
\tightlist
\item
  корреляционная функция в нормированном смысле:
\end{enumerate}

\[
\rho(k)=\frac{R(k)}{R(0)}
\]

при \(R(0)>0\).

В разных курсах слово ``корреляционная'' иногда используют для \(R(k)\),
а иногда для \(\rho(k)\). На экзамене лучше сразу сказать: ``я буду
называть \(R(k)\) ковариационной функцией, а нормированную величину
\(\rho(k)\) - корреляционной''.

\subsubsection{2.4. Основные свойства ковариационной
функции}\label{ux43eux441ux43dux43eux432ux43dux44bux435-ux441ux432ux43eux439ux441ux442ux432ux430-ux43aux43eux432ux430ux440ux438ux430ux446ux438ux43eux43dux43dux43eux439-ux444ux443ux43dux43aux446ux438ux438}

Ковариационная функция имеет несколько базовых свойств.

Во-первых,

\[
R(0)=\operatorname{Var}X_0\ge0.
\]

Если

\[
R(0)=0,
\]

то

\[
X_n=m
\]

почти наверное для всех \(n\), то есть процесс тривиален в смысле
второго порядка.

Во-вторых, для комплексного случая

\[
R(-k)=\overline{R(k)}.
\]

Для действительного процесса:

\[
R(-k)=R(k).
\]

В-третьих, из неравенства Коши-Буняковского:

\[
|R(k)|
=
|\operatorname{Cov}(X_k,X_0)|
\le
\sqrt{\operatorname{Var}X_k\,\operatorname{Var}X_0}
=
R(0).
\]

В-четвертых, \(R\) положительно определена. Для любых комплексных чисел
\(a_1,\ldots,a_r\) и целых индексов \(n_1,\ldots,n_r\):

\[
\sum_{\alpha=1}^r\sum_{\beta=1}^r
a_\alpha \overline{a_\beta}R(n_\alpha-n_\beta)
\ge0.
\]

Это просто дисперсия линейной комбинации:

\[
\mathbb E\left|
\sum_{\alpha=1}^r a_\alpha (X_{n_\alpha}-m)
\right|^2
\ge0.
\]

\subsubsection{2.5. Понятие эргодичности по среднему в широком
смысле}\label{ux43fux43eux43dux44fux442ux438ux435-ux44dux440ux433ux43eux434ux438ux447ux43dux43eux441ux442ux438-ux43fux43e-ux441ux440ux435ux434ux43dux435ux43cux443-ux432-ux448ux438ux440ux43eux43aux43eux43c-ux441ux43cux44bux441ux43bux435}

Теперь эргодичность. В широком смысле чаще всего рассматривают
эргодичность по среднему: можно ли среднее по одной длинной реализации
заменить математическим ожиданием?

Определим временное среднее:

\[
\overline X_N
=
\frac1N\sum_{n=0}^{N-1}X_n.
\]

Последовательность называется эргодичной по среднему в широком смысле,
если

\[
\overline X_N
\xrightarrow[L^2]{N\to\infty}
m.
\]

То есть

\[
\mathbb E|\overline X_N-m|^2\to0.
\]

Так как

\[
\mathbb E\overline X_N=m,
\]

это равносильно

\[
\operatorname{Var}(\overline X_N)\to0.
\]

\subsubsection{2.6. Вывод формулы для дисперсии временного
среднего}\label{ux432ux44bux432ux43eux434-ux444ux43eux440ux43cux443ux43bux44b-ux434ux43bux44f-ux434ux438ux441ux43fux435ux440ux441ux438ux438-ux432ux440ux435ux43cux435ux43dux43dux43eux433ux43e-ux441ux440ux435ux434ux43dux435ux433ux43e}

Вычислим эту дисперсию. Для центрированной последовательности
\(Y_n=X_n-m\):

\[
\overline X_N-m
=
\frac1N\sum_{n=0}^{N-1}Y_n.
\]

Тогда

\[
\mathbb E|\overline X_N-m|^2
=
\frac1{N^2}
\sum_{i=0}^{N-1}\sum_{j=0}^{N-1}
\mathbb E Y_i\overline{Y_j}.
\]

По стационарности в широком смысле:

\[
\mathbb E Y_i\overline{Y_j}=R(i-j).
\]

Значит

\[
\mathbb E|\overline X_N-m|^2
=
\frac1{N^2}
\sum_{i=0}^{N-1}\sum_{j=0}^{N-1}R(i-j).
\]

Если собрать пары с одинаковым лагом \(k=i-j\), то число таких пара
равно \(N-|k|\). Поэтому

\[
\mathbb E|\overline X_N-m|^2
=
\frac1N
\sum_{|k|<N}
\left(1-\frac{|k|}{N}\right)R(k).
\]

Для действительной последовательности это можно записать как

\[
\operatorname{Var}(\overline X_N)
=
\frac{R(0)}{N}
+
\frac{2}{N}
\sum_{k=1}^{N-1}
\left(1-\frac{k}{N}\right)R(k).
\]

\subsubsection{2.7. Основной критерий эргодичности и его спектральная
интерпретация}\label{ux43eux441ux43dux43eux432ux43dux43eux439-ux43aux440ux438ux442ux435ux440ux438ux439-ux44dux440ux433ux43eux434ux438ux447ux43dux43eux441ux442ux438-ux438-ux435ux433ux43e-ux441ux43fux435ux43aux442ux440ux430ux43bux44cux43dux430ux44f-ux438ux43dux442ux435ux440ux43fux440ux435ux442ux430ux446ux438ux44f}

Отсюда сразу получаем главный критерий второго порядка:

\[
\overline X_N\to m\ \text{в }L^2
\quad\Longleftrightarrow\quad
\frac1{N^2}
\sum_{i=0}^{N-1}\sum_{j=0}^{N-1}R(i-j)\to0.
\]

Для центрированной стационарной \(L^2\)-последовательности это можно
формулировать через среднее Чезаро ковариаций:

\[
\frac1N\sum_{k=0}^{N-1}R(k)\to0
\]

как необходимое и достаточное условие сходимости средних арифметических
в среднеквадратическом смысле. В спектральном языке это эквивалентно
отсутствию атома спектральной меры в нуле:

\[
F(\{0\})=0.
\]

Но для билета достаточно знать простые достаточные условия.

\subsubsection{2.8. Достаточные условия
эргодичности}\label{ux434ux43eux441ux442ux430ux442ux43eux447ux43dux44bux435-ux443ux441ux43bux43eux432ux438ux44f-ux44dux440ux433ux43eux434ux438ux447ux43dux43eux441ux442ux438}

\textbf{Первое достаточное условие:} корреляции затухают:

\[
R(k)\to0,
\qquad
|k|\to\infty.
\]

Тогда взвешенное среднее ковариаций стремится к нулю, и

\[
\overline X_N\xrightarrow{L^2}m.
\]

\textbf{Второе достаточное условие:} абсолютная суммируемость
ковариаций:

\[
\sum_{k\in\mathbb Z}|R(k)|<\infty.
\]

Тогда

\[
\operatorname{Var}(\overline X_N)
\le
\frac1N\sum_{|k|<N}|R(k)|
\le
\frac1N\sum_{k\in\mathbb Z}|R(k)|
\to0.
\]

Это очень удобное условие: если корреляции убывают достаточно быстро,
среднее эргодично.

\textbf{Третье достаточное условие:} среднее абсолютных корреляций
уходит в ноль:

\[
\frac1N\sum_{k=1}^{N-1}|R(k)|\to0.
\]

Тогда из формулы для дисперсии среднего также следует

\[
\operatorname{Var}(\overline X_N)\to0.
\]

\subsubsection{2.9. Интуитивное понимание и границы
применимости}\label{ux438ux43dux442ux443ux438ux442ux438ux432ux43dux43eux435-ux43fux43eux43dux438ux43cux430ux43dux438ux435-ux438-ux433ux440ux430ux43dux438ux446ux44b-ux43fux440ux438ux43cux435ux43dux438ux43cux43eux441ux442ux438}

Интуитивно все эти условия говорят одно: далекие значения процесса почти
не связаны линейно, поэтому усреднение по времени уничтожает случайные
колебания.

Важно понимать ограничения. Эргодичность по среднему в широком смысле -
это не полная эргодичность в узком смысле. Она касается только средних
самих \(X_n\) и только сходимости в \(L^2\). Узкая эргодичность сильнее
и относится ко всем сдвиг-инвариантным событиям или временным средним
функций \(f(X_n)\).

\subsection{3. Основные
определения}\label{ux43eux441ux43dux43eux432ux43dux44bux435-ux43eux43fux440ux435ux434ux435ux43bux435ux43dux438ux44f-23}

\subsubsection{Стационарность в широком
смысле}\label{ux441ux442ux430ux446ux438ux43eux43dux430ux440ux43dux43eux441ux442ux44c-ux432-ux448ux438ux440ux43eux43aux43eux43c-ux441ux43cux44bux441ux43bux435-2}

Последовательность \((X_n)_{n\in\mathbb Z}\) стационарна в широком
смысле, если

\[
\mathbb E|X_n|^2<\infty,
\]

\[
\mathbb EX_n=m
\]

не зависит от \(n\), и

\[
\operatorname{Cov}(X_n,X_m)
\]

зависит только от \(n-m\).

\subsubsection{Ковариационная
функция}\label{ux43aux43eux432ux430ux440ux438ux430ux446ux438ux43eux43dux43dux430ux44f-ux444ux443ux43dux43aux446ux438ux44f}

Для стационарной в широком смысле последовательности:

\[
R(k)=\operatorname{Cov}(X_k,X_0).
\]

В общем нестационарном случае:

\[
R(n,m)=\operatorname{Cov}(X_n,X_m).
\]

\subsubsection{Корреляционная
функция}\label{ux43aux43eux440ux440ux435ux43bux44fux446ux438ux43eux43dux43dux430ux44f-ux444ux443ux43dux43aux446ux438ux44f}

Если \(R(0)>0\), то

\[
\rho(k)=\frac{R(k)}{R(0)}.
\]

Она нормирована:

\[
\rho(0)=1,
\qquad
|\rho(k)|\le1.
\]

\subsubsection{Центрированная
последовательность}\label{ux446ux435ux43dux442ux440ux438ux440ux43eux432ux430ux43dux43dux430ux44f-ux43fux43eux441ux43bux435ux434ux43eux432ux430ux442ux435ux43bux44cux43dux43eux441ux442ux44c}

\[
Y_n=X_n-\mathbb EX_n.
\]

Для стационарной последовательности:

\[
Y_n=X_n-m.
\]

\subsubsection{Эргодичность по среднему в широком
смысле}\label{ux44dux440ux433ux43eux434ux438ux447ux43dux43eux441ux442ux44c-ux43fux43e-ux441ux440ux435ux434ux43dux435ux43cux443-ux432-ux448ux438ux440ux43eux43aux43eux43c-ux441ux43cux44bux441ux43bux435}

Стационарная в широком смысле последовательность эргодична по среднему,
если

\[
\frac1N\sum_{n=0}^{N-1}X_n
\xrightarrow{L^2}
\mathbb EX_0.
\]

\subsubsection{Среднеквадратичная
сходимость}\label{ux441ux440ux435ux434ux43dux435ux43aux432ux430ux434ux440ux430ux442ux438ux447ux43dux430ux44f-ux441ux445ux43eux434ux438ux43cux43eux441ux442ux44c}

\[
Z_N\xrightarrow{L^2}Z
\]

означает

\[
\mathbb E|Z_N-Z|^2\to0.
\]

\subsubsection{Положительная определенность ковариационной
функции}\label{ux43fux43eux43bux43eux436ux438ux442ux435ux43bux44cux43dux430ux44f-ux43eux43fux440ux435ux434ux435ux43bux435ux43dux43dux43eux441ux442ux44c-ux43aux43eux432ux430ux440ux438ux430ux446ux438ux43eux43dux43dux43eux439-ux444ux443ux43dux43aux446ux438ux438}

Функция \(R:\mathbb Z\to\mathbb C\) положительно определена, если

\[
\sum_{\alpha,\beta=1}^r
a_\alpha\overline{a_\beta}R(n_\alpha-n_\beta)\ge0
\]

для любых \(r\), \(a_1,\ldots,a_r\in\mathbb C\) и
\(n_1,\ldots,n_r\in\mathbb Z\).

\subsection{4. Основные теоремы и
свойства}\label{ux43eux441ux43dux43eux432ux43dux44bux435-ux442ux435ux43eux440ux435ux43cux44b-ux438-ux441ux432ux43eux439ux441ux442ux432ux430-23}

\subsubsection{1. Свойства ковариационной
функции}\label{ux441ux432ux43eux439ux441ux442ux432ux430-ux43aux43eux432ux430ux440ux438ux430ux446ux438ux43eux43dux43dux43eux439-ux444ux443ux43dux43aux446ux438ux438}

Для стационарной в широком смысле последовательности:

\[
R(0)\ge0,
\]

\[
R(-k)=\overline{R(k)},
\]

\[
|R(k)|\le R(0).
\]

Если процесс действительный, то

\[
R(-k)=R(k).
\]

\subsubsection{2. Положительная
определенность}\label{ux43fux43eux43bux43eux436ux438ux442ux435ux43bux44cux43dux430ux44f-ux43eux43fux440ux435ux434ux435ux43bux435ux43dux43dux43eux441ux442ux44c}

Ковариационная функция стационарной последовательности положительно
определена:

\[
\sum_{\alpha,\beta}
a_\alpha\overline{a_\beta}R(n_\alpha-n_\beta)\ge0.
\]

Доказательство:

\[
\sum_{\alpha,\beta}
a_\alpha\overline{a_\beta}R(n_\alpha-n_\beta)
=
\mathbb E\left|
\sum_\alpha a_\alpha(X_{n_\alpha}-m)
\right|^2
\ge0.
\]

\subsubsection{3. Дисперсия временного
среднего}\label{ux434ux438ux441ux43fux435ux440ux441ux438ux44f-ux432ux440ux435ux43cux435ux43dux43dux43eux433ux43e-ux441ux440ux435ux434ux43dux435ux433ux43e}

Пусть

\[
\overline X_N=\frac1N\sum_{n=0}^{N-1}X_n.
\]

Тогда

\[
\operatorname{Var}(\overline X_N)
=
\frac1{N^2}
\sum_{i=0}^{N-1}\sum_{j=0}^{N-1}R(i-j).
\]

Эквивалентно:

\[
\operatorname{Var}(\overline X_N)
=
\frac1N
\sum_{|k|<N}
\left(1-\frac{|k|}{N}\right)R(k).
\]

В действительном случае:

\[
\operatorname{Var}(\overline X_N)
=
\frac{R(0)}{N}
+
\frac{2}{N}\sum_{k=1}^{N-1}
\left(1-\frac{k}{N}\right)R(k).
\]

\subsubsection{4. Критерий эргодичности по
среднему}\label{ux43aux440ux438ux442ux435ux440ux438ux439-ux44dux440ux433ux43eux434ux438ux447ux43dux43eux441ux442ux438-ux43fux43e-ux441ux440ux435ux434ux43dux435ux43cux443}

Последовательность эргодична по среднему в широком смысле тогда и только
тогда, когда

\[
\operatorname{Var}(\overline X_N)\to0.
\]

По формуле выше это равносильно

\[
\frac1{N^2}
\sum_{i,j=0}^{N-1}R(i-j)\to0.
\]

\subsubsection{5. Достаточное условие через затухание
корреляций}\label{ux434ux43eux441ux442ux430ux442ux43eux447ux43dux43eux435-ux443ux441ux43bux43eux432ux438ux435-ux447ux435ux440ux435ux437-ux437ux430ux442ux443ux445ux430ux43dux438ux435-ux43aux43eux440ux440ux435ux43bux44fux446ux438ux439}

Если

\[
R(k)\to0,
\qquad
|k|\to\infty,
\]

то

\[
\overline X_N\xrightarrow{L^2}m.
\]

\subsubsection{6. Достаточное условие через абсолютную
суммируемость}\label{ux434ux43eux441ux442ux430ux442ux43eux447ux43dux43eux435-ux443ux441ux43bux43eux432ux438ux435-ux447ux435ux440ux435ux437-ux430ux431ux441ux43eux43bux44eux442ux43dux443ux44e-ux441ux443ux43cux43cux438ux440ux443ux435ux43cux43eux441ux442ux44c}

Если

\[
\sum_{k\in\mathbb Z}|R(k)|<\infty,
\]

то

\[
\operatorname{Var}(\overline X_N)=O\left(\frac1N\right)
\]

и последовательность эргодична по среднему.

\subsubsection{7. Спектральная
формулировка}\label{ux441ux43fux435ux43aux442ux440ux430ux43bux44cux43dux430ux44f-ux444ux43eux440ux43cux443ux43bux438ux440ux43eux432ux43aux430}

По теореме Герглотца положительно определенная ковариационная функция
имеет представление

\[
R(k)=\int_{-\pi}^{\pi}e^{ik\lambda}F(d\lambda),
\]

где \(F\) - конечная неотрицательная спектральная мера. Тогда средние
сходятся к \(m\) в \(L^2\) тогда и только тогда, когда у спектральной
меры нет атома в нуле:

\[
F(\{0\})=0.
\]

Это уже мост к Билету 25.

\subsection{5. Примеры}\label{ux43fux440ux438ux43cux435ux440ux44b-23}

\subsubsection{Пример 1. Белый
шум}\label{ux43fux440ux438ux43cux435ux440-1.-ux431ux435ux43bux44bux439-ux448ux443ux43c}

Пусть

\[
\mathbb E X_n=0,
\qquad
\operatorname{Var}X_n=\sigma^2,
\]

и величины некоррелированы:

\[
\operatorname{Cov}(X_n,X_m)=0,
\qquad
n\ne m.
\]

Тогда

\[
R(k)=
\begin{cases}
\sigma^2, & k=0,\\
0, & k\ne0.
\end{cases}
\]

Значит

\[
\operatorname{Var}(\overline X_N)=\frac{\sigma^2}{N}\to0.
\]

Белый шум эргодичен по среднему.

\subsubsection{Пример 2. Геометрически затухающие
корреляции}\label{ux43fux440ux438ux43cux435ux440-2.-ux433ux435ux43eux43cux435ux442ux440ux438ux447ux435ux441ux43aux438-ux437ux430ux442ux443ux445ux430ux44eux449ux438ux435-ux43aux43eux440ux440ux435ux43bux44fux446ux438ux438}

Пусть

\[
R(k)=\sigma^2 a^{|k|},
\qquad
|a|<1.
\]

Тогда

\[
\sum_{k\in\mathbb Z}|R(k)|<\infty.
\]

Следовательно,

\[
\overline X_N\xrightarrow{L^2}m.
\]

Такой вид ковариаций возникает, например, у стационарной авторегрессии
первого порядка \(AR(1)\):

\[
X_n=aX_{n-1}+\varepsilon_n,
\qquad
|a|<1.
\]

\subsubsection{Пример 3. Случайная
константа}\label{ux43fux440ux438ux43cux435ux440-3.-ux441ux43bux443ux447ux430ux439ux43dux430ux44f-ux43aux43eux43dux441ux442ux430ux43dux442ux430}

Пусть

\[
X_n=Z
\]

для всех \(n\), где \(Z\) - нетривиальная случайная величина,

\[
\mathbb EZ=m,
\qquad
\operatorname{Var}Z>0.
\]

Тогда последовательность стационарна, но

\[
\overline X_N=Z
\]

для всех \(N\). Поэтому

\[
\overline X_N\nrightarrow m
\]

в \(L^2\), если \(Z\) не константа. Ковариационная функция:

\[
R(k)=\operatorname{Var}Z
\]

для всех \(k\). Она не затухает, и эргодичности среднего нет.

\subsubsection{Пример 4. Общий случай с суммируемыми
ковариациями}\label{ux43fux440ux438ux43cux435ux440-4.-ux43eux431ux449ux438ux439-ux441ux43bux443ux447ux430ux439-ux441-ux441ux443ux43cux43cux438ux440ux443ux435ux43cux44bux43cux438-ux43aux43eux432ux430ux440ux438ux430ux446ux438ux44fux43cux438}

Если

\[
|R(k)|\le \frac{C}{(1+|k|)^{1+\delta}},
\qquad
\delta>0,
\]

то

\[
\sum_{k\in\mathbb Z}|R(k)|<\infty.
\]

Значит среднее эргодично. Это типичный случай ``быстрого забывания''
прошлого.

\subsubsection{Пример 5. Некоррелированность не равна
независимости}\label{ux43fux440ux438ux43cux435ux440-5.-ux43dux435ux43aux43eux440ux440ux435ux43bux438ux440ux43eux432ux430ux43dux43dux43eux441ux442ux44c-ux43dux435-ux440ux430ux432ux43dux430-ux43dux435ux437ux430ux432ux438ux441ux438ux43cux43eux441ux442ux438}

Можно иметь

\[
R(k)=0
\]

при \(k\ne0\), но величины не будут независимыми. Для эргодичности по
среднему в широком смысле достаточно контроля вторых моментов, но для
узкой эргодичности этого уже мало.

\subsection{6. Схемы доказательств /
обоснований}\label{ux441ux445ux435ux43cux44b-ux434ux43eux43aux430ux437ux430ux442ux435ux43bux44cux441ux442ux432-ux43eux431ux43eux441ux43dux43eux432ux430ux43dux438ux439-16}

\textbf{Доказательство свойств \(R(k)\).}

Если процесс центрирован, то

\[
R(k)=\mathbb E X_k\overline{X_0}.
\]

Тогда

\[
R(-k)=\mathbb E X_{-k}\overline{X_0}.
\]

По стационарности пара \((X_{-k},X_0)\) имеет те же характеристики
второго порядка, что \((X_0,X_k)\), поэтому

\[
R(-k)=\overline{R(k)}.
\]

Неравенство

\[
|R(k)|\le R(0)
\]

следует из Коши-Буняковского:

\[
|\mathbb E X_k\overline{X_0}|
\le
\sqrt{\mathbb E|X_k|^2\,\mathbb E|X_0|^2}
=R(0).
\]

\textbf{Вывод формулы для дисперсии среднего.}

Центрируем:

\[
Y_n=X_n-m.
\]

Тогда

\[
\overline X_N-m=\frac1N\sum_{n=0}^{N-1}Y_n.
\]

Следовательно,

\[
\mathbb E|\overline X_N-m|^2
=
\frac1{N^2}
\mathbb E
\left|
\sum_{n=0}^{N-1}Y_n
\right|^2.
\]

Раскрываем квадрат:

\[
\mathbb E
\left|
\sum_{n=0}^{N-1}Y_n
\right|^2
=
\sum_{i=0}^{N-1}\sum_{j=0}^{N-1}
\mathbb E Y_i\overline{Y_j}.
\]

По стационарности:

\[
\mathbb E Y_i\overline{Y_j}=R(i-j).
\]

Получаем:

\[
\mathbb E|\overline X_N-m|^2
=
\frac1{N^2}
\sum_{i,j=0}^{N-1}R(i-j).
\]

Если собрать одинаковые лаги:

\[
\mathbb E|\overline X_N-m|^2
=
\frac1N
\sum_{|k|<N}
\left(1-\frac{|k|}{N}\right)R(k).
\]

\textbf{Почему абсолютная суммируемость достаточна.}

Если

\[
\sum_{k\in\mathbb Z}|R(k)|<\infty,
\]

то

\[
\operatorname{Var}(\overline X_N)
\le
\frac1N
\sum_{|k|<N}|R(k)|
\le
\frac1N
\sum_{k\in\mathbb Z}|R(k)|
\to0.
\]

Значит

\[
\overline X_N\to m
\]

в \(L^2\).

\textbf{Почему \(R(k)\to0\) достаточно.}

Из \(|R(k)|\le R(0)\) следует ограниченность ковариаций. Если
\(R(k)\to0\), то средние по Чезаро тоже стремятся к нулю:

\[
\frac1N\sum_{k=0}^{N-1}R(k)\to0.
\]

С учетом симметрии \(R(-k)=\overline{R(k)}\) это дает исчезновение
взвешенного среднего ковариаций в формуле для

\[
\operatorname{Var}(\overline X_N).
\]

Следовательно,

\[
\operatorname{Var}(\overline X_N)\to0.
\]

\subsection{7. Что написать на
доске}\label{ux447ux442ux43e-ux43dux430ux43fux438ux441ux430ux442ux44c-ux43dux430-ux434ux43eux441ux43aux435-23}

Ковариационная функция:

\[
R(k)=\operatorname{Cov}(X_k,X_0)
=
\mathbb E[(X_k-m)\overline{(X_0-m)}].
\]

Нормированная корреляция:

\[
\rho(k)=\frac{R(k)}{R(0)}.
\]

Свойства:

\[
R(0)\ge0,
\qquad
R(-k)=\overline{R(k)},
\qquad
|R(k)|\le R(0).
\]

Временное среднее:

\[
\overline X_N=\frac1N\sum_{n=0}^{N-1}X_n.
\]

Дисперсия среднего:

\[
\mathbb E|\overline X_N-m|^2
=
\frac1{N^2}
\sum_{i,j=0}^{N-1}R(i-j)
=
\frac1N
\sum_{|k|<N}
\left(1-\frac{|k|}{N}\right)R(k).
\]

Достаточные условия:

\[
R(k)\to0
\quad\Rightarrow\quad
\overline X_N\xrightarrow{L^2}m,
\]

и

\[
\sum_{k\in\mathbb Z}|R(k)|<\infty
\quad\Rightarrow\quad
\overline X_N\xrightarrow{L^2}m.
\]

Спектральная связка:

\[
R(k)=\int_{-\pi}^{\pi}e^{ik\lambda}F(d\lambda),
\qquad
F(\{0\})=0.
\]

\subsection{8. Типичные дополнительные
вопросы}\label{ux442ux438ux43fux438ux447ux43dux44bux435-ux434ux43eux43fux43eux43bux43dux438ux442ux435ux43bux44cux43dux44bux435-ux432ux43eux43fux440ux43eux441ux44b-23}

\textbf{Чем ковариационная функция отличается от корреляционной?}

Ковариационная функция:

\[
R(k)=\operatorname{Cov}(X_k,X_0).
\]

Нормированная корреляционная функция:

\[
\rho(k)=\frac{R(k)}{R(0)}.
\]

У \(\rho(0)=1\), а у \(R(0)=\operatorname{Var}X_0\).

\textbf{Почему нужно центрирование?}

Корреляция должна измерять совместные отклонения от среднего, поэтому
нужно брать

\[
X_n-m.
\]

Если забыть центрирование, то постоянная ненулевая средняя компонента
даст ложный вклад в зависимость.

\textbf{Почему \(\operatorname{Var}(\overline X_N)\to0\) означает
эргодичность по среднему?}

Потому что

\[
\mathbb E\overline X_N=m.
\]

Значит

\[
\mathbb E|\overline X_N-m|^2
=
\operatorname{Var}(\overline X_N).
\]

Сходимость этой величины к нулю и есть сходимость в \(L^2\).

\textbf{Достаточно ли условия \(R(k)\to0\)?}

Да, для эргодичности среднего в широком смысле это достаточное условие.
Оно не обязательно для всех возможных вариантов, но удобное и часто
используемое.

\textbf{Почему абсолютная суммируемость сильнее?}

Если

\[
\sum_k|R(k)|<\infty,
\]

то ковариации не просто стремятся к нулю, а делают это достаточно
быстро. Тогда дисперсия среднего убывает как \(O(1/N)\).

\textbf{Пример неэргодичной последовательности?}

Случайная константа:

\[
X_n=Z.
\]

Тогда

\[
\overline X_N=Z
\]

для всех \(N\), и среднее не сходится к \(\mathbb EZ\), если \(Z\) не
константа.

\textbf{Связь со спектральной мерой?}

Если

\[
R(k)=\int e^{ik\lambda}F(d\lambda),
\]

то атом спектральной меры в нуле соответствует случайной постоянной
компоненте, которая не исчезает при усреднении. Поэтому условие
эргодичности среднего:

\[
F(\{0\})=0.
\]

\textbf{В чем отличие от узкой эргодичности?}

Узкая эргодичность говорит о сдвиг-инвариантных событиях и временных
средних функций. Эргодичность в широком смысле по среднему говорит
только о сходимости

\[
\frac1N\sum X_n
\]

в \(L^2\).

\subsection{9. Типичные
ошибки}\label{ux442ux438ux43fux438ux447ux43dux44bux435-ux43eux448ux438ux431ux43aux438-23}

\begin{itemize}
\tightlist
\item
  Путать ковариационную функцию \(R(k)\) и нормированный коэффициент
  корреляции \(\rho(k)\).
\item
  Забывать центрирование: писать \(\mathbb EX_kX_0\) вместо
  \(\mathbb E[(X_k-m)(X_0-m)]\) без условия \(m=0\).
\item
  Говорить ``эргодична'' без уточнения: в узком смысле, по среднему, в
  \(L^2\).
\item
  Считать, что \(R(k)\to0\) означает независимость. Это только
  некоррелированность в пределе.
\item
  Забывать множитель \(N-|k|\) при переходе от двойной суммы к сумме по
  лагам.
\item
  Думать, что стационарность сама по себе дает эргодичность. Случайная
  константа - контрпример.
\item
  Считать абсолютную суммируемость необходимой. Она достаточна, но не
  необходима.
\item
  В комплексном случае забывать сопряжение и свойство
  \(R(-k)=\overline{R(k)}\).
\end{itemize}

\subsection{10. Связи с другими
билетами}\label{ux441ux432ux44fux437ux438-ux441-ux434ux440ux443ux433ux438ux43cux438-ux431ux438ux43bux435ux442ux430ux43cux438-23}

\begin{itemize}
\tightlist
\item
  Билет 12: ковариация как скалярное произведение в \(L_2\) и
  положительная определенность.
\item
  Билет 14: коэффициент корреляции и геометрия центрированных величин.
\item
  Билет 17: стационарность и эргодичность в широком и узком смыслах.
\item
  Билет 18: для гауссовских последовательностей ковариация определяет
  все конечномерные распределения.
\item
  Билет 21: эргодичность среднего - аналог закона больших чисел для
  зависимых стационарных последовательностей.
\item
  Билет 23: стационарные марковские цепи дают примеры зависимых
  последовательностей, где корреляции часто затухают.
\item
  Билет 25: спектральная мера является преобразованием Фурье
  ковариационной функции.
\end{itemize}

\subsection{11. Минимум на
удовлетворительно}\label{ux43cux438ux43dux438ux43cux443ux43c-ux43dux430-ux443ux434ux43eux432ux43bux435ux442ux432ux43eux440ux438ux442ux435ux43bux44cux43dux43e-23}

Нужно сказать:

\begin{enumerate}
\def\labelenumi{\arabic{enumi}.}
\tightlist
\item
  Для стационарной в широком смысле последовательности
\end{enumerate}

\[
R(k)=\operatorname{Cov}(X_k,X_0).
\]

\begin{enumerate}
\def\labelenumi{\arabic{enumi}.}
\setcounter{enumi}{1}
\tightlist
\item
  Временное среднее:
\end{enumerate}

\[
\overline X_N=\frac1N\sum_{n=0}^{N-1}X_n.
\]

\begin{enumerate}
\def\labelenumi{\arabic{enumi}.}
\setcounter{enumi}{2}
\tightlist
\item
  Эргодичность по среднему:
\end{enumerate}

\[
\overline X_N\xrightarrow{L^2}\mathbb EX_0.
\]

\begin{enumerate}
\def\labelenumi{\arabic{enumi}.}
\setcounter{enumi}{3}
\tightlist
\item
  Главная формула:
\end{enumerate}

\[
\operatorname{Var}(\overline X_N)
=
\frac1{N^2}\sum_{i,j=0}^{N-1}R(i-j).
\]

\begin{enumerate}
\def\labelenumi{\arabic{enumi}.}
\setcounter{enumi}{4}
\tightlist
\item
  Если корреляции затухают достаточно быстро, например
  \(\sum_k|R(k)|<\infty\), то среднее эргодично.
\end{enumerate}

\subsection{12. Минимум на
хорошо}\label{ux43cux438ux43dux438ux43cux443ux43c-ux43dux430-ux445ux43eux440ux43eux448ux43e-23}

Помимо минимума, нужно:

\begin{itemize}
\tightlist
\item
  различить \(R(k)\) и \(\rho(k)\);
\item
  назвать свойства \(R(0)\ge0\), \(R(-k)=\overline{R(k)}\),
  \(|R(k)|\le R(0)\);
\item
  вывести формулу
\end{itemize}

\[
\operatorname{Var}(\overline X_N)
=
\frac1N\sum_{|k|<N}\left(1-\frac{|k|}{N}\right)R(k);
\]

\begin{itemize}
\tightlist
\item
  привести пример белого шума и случайной константы;
\item
  объяснить, почему случайная константа стационарна, но не эргодична по
  среднему.
\end{itemize}

\subsection{13. Что добавить для
отлично}\label{ux447ux442ux43e-ux434ux43eux431ux430ux432ux438ux442ux44c-ux434ux43bux44f-ux43eux442ux43bux438ux447ux43dux43e-23}

Для сильного ответа стоит добавить:

\begin{itemize}
\tightlist
\item
  положительную определенность ковариационной функции;
\item
  условие Чезаро:
\end{itemize}

\[
\frac1N\sum_{k=0}^{N-1}R(k)\to0;
\]

\begin{itemize}
\tightlist
\item
  спектральную интерпретацию:
\end{itemize}

\[
F(\{0\})=0;
\]

\begin{itemize}
\tightlist
\item
  связь с теоремой Герглотца:
\end{itemize}

\[
R(k)=\int_{-\pi}^{\pi}e^{ik\lambda}F(d\lambda);
\]

\begin{itemize}
\tightlist
\item
  четкое различие между эргодичностью в широком смысле и узкой
  эргодичностью.
\end{itemize}

\subsection{14. Устная версия ответа на 5
минут}\label{ux443ux441ux442ux43dux430ux44f-ux432ux435ux440ux441ux438ux44f-ux43eux442ux432ux435ux442ux430-ux43dux430-5-ux43cux438ux43dux443ux442-23}

Пусть \(X=(X_n)_{n\in\mathbb Z}\) - стационарная в широком смысле
случайная последовательность:

\[
\mathbb EX_n=m,
\qquad
\mathbb E|X_n|^2<\infty,
\]

а ковариация зависит только от лага. Тогда ковариационная функция:

\[
R(k)=\operatorname{Cov}(X_k,X_0)
=
\mathbb E[(X_k-m)\overline{(X_0-m)}].
\]

Если \(R(0)>0\), нормированная корреляционная функция:

\[
\rho(k)=\frac{R(k)}{R(0)}.
\]

Основные свойства:

\[
R(0)\ge0,
\qquad
R(-k)=\overline{R(k)},
\qquad
|R(k)|\le R(0).
\]

Кроме того, \(R\) положительно определена, потому что для любой линейной
комбинации

\[
\mathbb E\left|\sum_\alpha a_\alpha(X_{n_\alpha}-m)\right|^2\ge0.
\]

Эргодичность по среднему в широком смысле означает:

\[
\overline X_N=\frac1N\sum_{n=0}^{N-1}X_n
\xrightarrow{L^2}
m.
\]

Так как \(\mathbb E\overline X_N=m\), это равносильно

\[
\operatorname{Var}(\overline X_N)\to0.
\]

Вычисляем дисперсию:

\[
\operatorname{Var}(\overline X_N)
=
\frac1{N^2}\sum_{i=0}^{N-1}\sum_{j=0}^{N-1}R(i-j).
\]

Собирая одинаковые лаги:

\[
\operatorname{Var}(\overline X_N)
=
\frac1N\sum_{|k|<N}
\left(1-\frac{|k|}{N}\right)R(k).
\]

Отсюда достаточные условия эргодичности: если

\[
R(k)\to0
\]

или тем более

\[
\sum_{k\in\mathbb Z}|R(k)|<\infty,
\]

то

\[
\operatorname{Var}(\overline X_N)\to0,
\]

а значит

\[
\overline X_N\to m
\]

в среднем квадратичном.

Пример эргодичного среднего - белый шум: \(R(k)=0\) при \(k\ne0\), и
дисперсия среднего равна \(\sigma^2/N\). Контрпример - случайная
константа \(X_n=Z\): процесс стационарен, но \(\overline X_N=Z\),
поэтому среднее не сходится к \(\mathbb EZ\), если \(Z\) не вырождена.

В конце можно добавить связь со спектрами: по теореме Герглотца

\[
R(k)=\int_{-\pi}^{\pi}e^{ik\lambda}F(d\lambda).
\]

Эргодичность среднего соответствует отсутствию атома спектральной меры в
нуле:

\[
F(\{0\})=0.
\]

\subsection{15. Если осталось 2 минуты перед
экзаменом}\label{ux435ux441ux43bux438-ux43eux441ux442ux430ux43bux43eux441ux44c-2-ux43cux438ux43dux443ux442ux44b-ux43fux435ux440ux435ux434-ux44dux43aux437ux430ux43cux435ux43dux43eux43c-23}

Запомнить:

\[
R(k)=\operatorname{Cov}(X_k,X_0).
\]

Свойства:

\[
R(0)\ge0,
\qquad
R(-k)=\overline{R(k)},
\qquad
|R(k)|\le R(0).
\]

Эргодичность среднего:

\[
\frac1N\sum_{n=0}^{N-1}X_n
\xrightarrow{L^2}
\mathbb EX_0.
\]

Главная формула:

\[
\operatorname{Var}(\overline X_N)
=
\frac1{N^2}\sum_{i,j=0}^{N-1}R(i-j).
\]

Достаточные условия:

\[
R(k)\to0
\]

или

\[
\sum_k|R(k)|<\infty.
\]

Главная ловушка: стационарность не означает эргодичность; случайная
константа стационарна, но не эргодична по среднему.

\newpage

\section{Билет 25. Спектральные
характеристики}\label{ux431ux438ux43bux435ux442-25.-ux441ux43fux435ux43aux442ux440ux430ux43bux44cux43dux44bux435-ux445ux430ux440ux430ux43aux442ux435ux440ux438ux441ux442ux438ux43aux438}

Источники: Ширяев 2, гл. VI (Спектральная теория); лекции ДСП;
Сердобольская.

\subsection{1. Короткий
ответ}\label{ux43aux43eux440ux43eux442ux43aux438ux439-ux43eux442ux432ux435ux442-24}

Любая стационарная в широком смысле случайная последовательность (с
дискретным временем)

\[
X = (X_n)_{n \in \mathbb{Z}}
\]

и её ковариационная функция

\[
R(k) = \operatorname{Cov}(X_{n+k}, X_n)
\]

допускают спектральное представление. Это значит, что процесс можно
представить как суперпозицию некоррелированных гармонических колебаний с
различными частотами

\[
\lambda \in [-\pi, \pi].
\]

Главный результат для ковариационной функции --- \textbf{теорема
Герглотца}: любая ковариационная функция \(R(k)\) представляется как
преобразование Фурье некоторой конечной меры на \([-\pi, \pi]\)
(спектральной меры \(F\)):

\[
R(k) = \int_{-\pi}^{\pi} e^{i k \lambda} F(d\lambda).
\]

Если мера \(F\) имеет плотность \(f(\lambda)\) по мере Лебега, то она
называется \textbf{спектральной плотностью}, и \(R(k)\) связана с ней
взаимно обратными преобразованиями Фурье.

Сам процесс \(X_n\) также представляется интегралом:

\[
X_n = \int_{-\pi}^{\pi} e^{i n \lambda} Z(d\lambda),
\]

где \(Z\) --- случайная мера с ортогональными приращениями. Это
представление изометрично связывает гильбертово пространство величин
процесса \(H_X\) (см. Билет 12) и пространство функций \(L_2(F)\).

\subsection{2. Полный
ответ}\label{ux43fux43eux43bux43dux44bux439-ux43eux442ux432ux435ux442-24}

\subsubsection{2.1. Введение: ковариационная функция и её
свойства}\label{ux432ux432ux435ux434ux435ux43dux438ux435-ux43aux43eux432ux430ux440ux438ux430ux446ux438ux43eux43dux43dux430ux44f-ux444ux443ux43dux43aux446ux438ux44f-ux438-ux435ux451-ux441ux432ux43eux439ux441ux442ux432ux430}

Рассмотрим комплекснозначную случайную последовательность со вторыми
моментами

\[
\mathbb E|X_n|^2 < \infty.
\]

Как было показано в Билете 17, процесс стационарен в широком смысле,
если его математическое ожидание постоянно:

\[
\mathbb E X_n = m,
\]

а ковариация зависит только от сдвига:

\[
\operatorname{Cov}(X_{n+k}, X_n) = R(k).
\]

Без ограничения общности далее будем считать, что процесс центрирован:

\[
m = 0, \quad R(k) = \mathbb E[X_k \overline{X_0}].
\]

В Билете 24 было доказано, что ковариационная функция неотрицательно
определена: для любых комплексных чисел \(a_1, \ldots, a_r\) и моментов
времени \(t_1, \ldots, t_r\) выполнено

\[
\sum_{j=1}^r \sum_{l=1}^r a_j \overline{a_l} R(t_j - t_l) = \mathbb E \left| \sum_{j=1}^r a_j X_{t_j} \right|^2 \ge 0.
\]

Это свойство роднит ковариационную функцию с характеристическими
функциями распределений (см. Билет 08). Для последних работает теорема
Бохнера---Хинчина, утверждающая, что непрерывная положительно
определенная функция является преобразованием Фурье вероятностной меры.
Аналогичный результат для дискретного времени называется теоремой
Герглотца.

\subsubsection{2.2. Теорема Герглотца о
спектре}\label{ux442ux435ux43eux440ux435ux43cux430-ux433ux435ux440ux433ux43bux43eux442ux446ux430-ux43e-ux441ux43fux435ux43aux442ux440ux435}

Так как время \(t \in \mathbb{Z}\) дискретно, экспоненты
\(e^{i t \lambda}\) периодичны по частоте \(\lambda\) с периодом
\(2\pi\). Поэтому ``весь спектр'' частот, доступный для процесса с
дискретным временем, помещается на одном периоде, например на
полуинтервале

\[
[-\pi, \pi).
\]

\textbf{Теорема Герглотца} утверждает, что \(R(k)\) неотрицательно
определена тогда и только тогда, когда она является преобразованием
Фурье-Стилтьеса:

\[
R(k) = \int_{-\pi}^{\pi} e^{i k \lambda} F(d\lambda),
\]

где \(F\) --- конечная мера на борелевских множествах \([-\pi, \pi)\),
называемая \textbf{спектральной мерой}. Выбор \([-\pi,\pi)\) вместо
\([-\pi,\pi]\) нужен только для того, чтобы не считать одну и ту же
частоту дважды: точки \(-\pi\) и \(\pi\) дают одинаковые значения
\(e^{ik\lambda}\) при целых \(k\).

Физический смысл спектральной меры: дисперсия процесса (полная мощность)
равна

\[
\operatorname{Var}(X_n) = R(0) = \int_{-\pi}^{\pi} F(d\lambda) = F([-\pi, \pi)).
\]

Мера интервала \(F([\lambda_1, \lambda_2])\) показывает вклад этих
частот в полную мощность процесса. Доля мощности равна
\(F([\lambda_1,\lambda_2])/R(0)\) при \(R(0)>0\).

Часто вместо меры пишут \textbf{спектральную функцию}

\[
\mathcal F(\lambda)=F([-\pi,\lambda]), \qquad -\pi\le \lambda<\pi.
\]

Это не новая сущность, а способ записать ту же меру через неубывающую
функцию. Тогда интеграл выше читается как интеграл Лебега-Стилтьеса:

\[
R(k)=\int_{-\pi}^{\pi}e^{ik\lambda}\,d\mathcal F(\lambda).
\]

Спектральная мера единственна: если две конечные меры имеют одинаковые
коэффициенты Фурье \(\int e^{ik\lambda}F(d\lambda)\) для всех
\(k\in\mathbb Z\), то они совпадают. Идея доказательства единственности:
тригонометрические полиномы равномерно приближают непрерывные функции на
окружности, значит совпадают интегралы от всех непрерывных функций, а
значит совпадают и меры.

\subsubsection{2.3. Спектральная
плотность}\label{ux441ux43fux435ux43aux442ux440ux430ux43bux44cux43dux430ux44f-ux43fux43bux43eux442ux43dux43eux441ux442ux44c}

В общем случае по теореме Лебега о разложении (связано с Билетом 10),
спектральная мера распадается на три компоненты: абсолютно непрерывную,
дискретную (атомарную) и сингулярную:

\[
F = F_{ac} + F_d + F_{sing}.
\]

Если мера абсолютно непрерывна относительно меры Лебега, существует
функция \(f(\lambda) \ge 0\), называемая \textbf{спектральной
плотностью}, такая что:

\[
F(d\lambda) = f(\lambda) d\lambda.
\]

В этом случае теорема Герглотца принимает вид:

\[
R(k) = \int_{-\pi}^{\pi} e^{i k \lambda} f(\lambda) d\lambda.
\]

Если ряд \(\sum_{k=-\infty}^{\infty} |R(k)|\) сходится, то спектральная
плотность существует, непрерывна, и её можно найти по формуле обращения:

\[
f(\lambda) = \frac{1}{2\pi} \sum_{k=-\infty}^{\infty} R(k) e^{-i k \lambda}.
\]

В более слабой форме, если \(\sum_k |R(k)|^2<\infty\), то спектральная
функция тоже имеет плотность, но уже в смысле \(L_2\):

\[
f(\lambda)=\frac1{2\pi}\sum_{k=-\infty}^{\infty}R(k)e^{-ik\lambda}
\quad \text{в } L_2[-\pi,\pi].
\]

Для экзамена удобно помнить иерархию: абсолютная суммируемость
ковариаций дает хорошую непрерывную плотность; квадрат-суммируемость
дает плотность как \(L_2\)-функцию; без этих условий остается общая
спектральная мера, у которой могут быть атомы и сингулярная часть.

Для вещественного процесса \(R(k)\) вещественна и четна
(\(R(k) = R(-k)\)). Поэтому спектральная плотность также является четной
функцией:

\[
f(\lambda) = f(-\lambda),
\]

а синусы в формуле Эйлера взаимно уничтожаются:

\[
f(\lambda) = \frac{1}{2\pi} R(0) + \frac{1}{\pi} \sum_{k=1}^{\infty} R(k) \cos(k \lambda).
\]

\subsubsection{2.4. Спектральное представление стационарной
последовательности}\label{ux441ux43fux435ux43aux442ux440ux430ux43bux44cux43dux43eux435-ux43fux440ux435ux434ux441ux442ux430ux432ux43bux435ux43dux438ux435-ux441ux442ux430ux446ux438ux43eux43dux430ux440ux43dux43eux439-ux43fux43eux441ux43bux435ux434ux43eux432ux430ux442ux435ux43bux44cux43dux43eux441ux442ux438}

Мы выразили через частоты функцию ковариации. Можно ли выразить через
гармоники сам случайный процесс?

Да. Любую центрированную стационарную в широком смысле
последовательность можно записать в виде \textbf{стохастического
интеграла}:

\[
X_n = \int_{-\pi}^{\pi} e^{i n \lambda} Z(d\lambda),
\]

где \(Z\) --- случайная мера с ортогональными приращениями.

Что значит ``случайная мера с ортогональными приращениями''? Каждому
борелевскому множеству \(\Delta \subset [-\pi, \pi)\) ставится в
соответствие комплекснозначная случайная величина
\(Z(\Delta) \in L_2(\Omega)\). 1. Она центрирована:
\(\mathbb E Z(\Delta) = 0\). 2. Ее дисперсия задается спектральной мерой
\(F\):

\[
\mathbb E|Z(\Delta)|^2 = F(\Delta).
\]

\begin{enumerate}
\def\labelenumi{\arabic{enumi}.}
\setcounter{enumi}{2}
\tightlist
\item
  Для непересекающихся \(\Delta_1\) и \(\Delta_2\) величины
  \(Z(\Delta_1)\) и \(Z(\Delta_2)\) некоррелированы (ортогональны в
  пространстве \(L_2(\Omega)\)):
\end{enumerate}

\[
\mathbb E[Z(\Delta_1) \overline{Z(\Delta_2)}] = 0.
\]

\subsubsection{2.5. Изометрия гильбертовых
пространств}\label{ux438ux437ux43eux43cux435ux442ux440ux438ux44f-ux433ux438ux43bux44cux431ux435ux440ux442ux43eux432ux44bux445-ux43fux440ux43eux441ux442ux440ux430ux43dux441ux442ux432}

Понятие ортогональной случайной меры и интеграла по ней напрямую
опирается на теорию гильбертовых пространств (см. Билет 12).

Пусть \(H_X\) --- замкнутое линейное подпространство в \(L_2(\Omega)\),
порожденное величинами \(\{X_n\}_{n \in \mathbb{Z}}\). А \(L_2(F)\) ---
пространство детерминированных (не случайных!) комплекснозначных функций
на \([-\pi, \pi]\), интегрируемых с квадратом по мере \(F\).

Интеграл \(\int g(\lambda) Z(d\lambda)\) строится как предел линейных
комбинаций: мы берем ступенчатую функцию
\(g(\lambda) = \sum c_j \mathbf{1}_{\Delta_j}(\lambda)\) и определяем
для нее интеграл как \(\sum c_j Z(\Delta_j)\).

Из свойства ортогональности приращений \(Z\) следует фундаментальное
свойство: скалярное произведение случайных величин (в \(H_X\)) в
точности равно скалярному произведению функций (в \(L_2(F)\)).

Действительно, для ступенчатых функций:

\[
\mathbb E \left[ \left( \sum_j c_j Z(\Delta_j) \right) \overline{\left( \sum_k d_k Z(\Delta_k) \right)} \right]
=
\sum_j c_j \overline{d_j} F(\Delta_j)
=
\int_{-\pi}^{\pi} g(\lambda) \overline{h(\lambda)} F(d\lambda).
\]

Таким образом, отображение

\[
H_X \ni \int_{-\pi}^{\pi} g(\lambda) Z(d\lambda) \quad \longleftrightarrow \quad g(\lambda) \in L_2(F)
\]

является \textbf{изометрическим изоморфизмом} между гильбертовым
пространством случайных величин и гильбертовым пространством функций.

Элементу \(X_n \in H_X\) соответствует функция
\(e^{i n \lambda} \in L_2(F)\). Скалярное произведение
\(\langle X_n, X_0 \rangle_{H_X} = R(n)\) переходит в:

\[
\int_{-\pi}^{\pi} e^{i n \lambda} \cdot 1 \, F(d\lambda),
\]

что в точности возвращает нас к теореме Герглотца.

\subsubsection{2.6. Связь спектра с усреднением и
эргодичностью}\label{ux441ux432ux44fux437ux44c-ux441ux43fux435ux43aux442ux440ux430-ux441-ux443ux441ux440ux435ux434ux43dux435ux43dux438ux435ux43c-ux438-ux44dux440ux433ux43eux434ux438ux447ux43dux43eux441ux442ux44cux44e}

Спектр сразу объясняет результат из Билета 24 про временные средние. Для
центрированного процесса

\[
\overline X_N=\frac1N\sum_{n=0}^{N-1}X_n
\]

по спектральному представлению:

\[
\overline X_N
=
\int_{-\pi}^{\pi}
\left(\frac1N\sum_{n=0}^{N-1}e^{in\lambda}\right)
Z(d\lambda).
\]

По изометрии:

\[
\mathbb E|\overline X_N|^2
=
\int_{-\pi}^{\pi}
\left|
\frac1N\sum_{n=0}^{N-1}e^{in\lambda}
\right|^2
F(d\lambda).
\]

Множитель под интегралом является нормированным ядром Фейера: при
\(\lambda\ne0\) он стремится к нулю, а в точке \(\lambda=0\) равен
\(1\). Поэтому предел дисперсии среднего равен массе спектральной меры в
нуле:

\[
\mathbb E|\overline X_N|^2 \to F(\{0\}).
\]

Значит, эргодичность по среднему в широком смысле эквивалентна
отсутствию атома в нулевой частоте:

\[
F(\{0\})=0.
\]

Интуитивно атом в нуле --- это случайная постоянная компонента процесса.
Она не исчезает при усреднении по времени.

\subsection{3. Основные
определения}\label{ux43eux441ux43dux43eux432ux43dux44bux435-ux43eux43fux440ux435ux434ux435ux43bux435ux43dux438ux44f-24}

\textbf{Спектральная мера} \(F\) --- конечная мера на борелевской
\(\sigma\)-алгебре одного периода частот, обычно \([-\pi, \pi)\), такая
что ковариационная функция \(R(k)\) является ее преобразованием
Фурье-Стилтьеса.

\textbf{Спектральная функция} \(\mathcal F(\lambda)\) --- функция
распределения спектральной меры:

\[
\mathcal F(\lambda)=F([-\pi,\lambda]).
\]

Она неубывающая и правая непрерывная; скачок
\(\mathcal F(\lambda)-\mathcal F(\lambda-)\) равен массе атома в частоте
\(\lambda\).

\textbf{Спектральная плотность} \(f(\lambda)\) --- производная
Радона-Никодима (см. Билет 10) спектральной меры относительно меры
Лебега, существующая, если мера \(F\) абсолютно непрерывна:

\[
F(d\lambda) = f(\lambda) d\lambda.
\]

\textbf{Случайная мера с ортогональными приращениями} \(Z\) --- функция
множества, заданная на борелевской \(\sigma\)-алгебре \([-\pi, \pi)\),
значениями которой являются случайные величины из \(L_2(\Omega)\),
удовлетворяющая: 1. \(\mathbb E Z(\Delta) = 0\). 2.
\(\Delta_1 \cap \Delta_2 = \emptyset \implies Z(\Delta_1 \cup \Delta_2) = Z(\Delta_1) + Z(\Delta_2)\)
почти наверное. 3.
\(\Delta_1 \cap \Delta_2 = \emptyset \implies \mathbb E[Z(\Delta_1) \overline{Z(\Delta_2)}] = 0\).

\textbf{Структурная мера} \(\mu_Z\) для меры \(Z\) --- детерминированная
мера \(\mu_Z(\Delta) = \mathbb E|Z(\Delta)|^2\). (В контексте
стационарных процессов структурная мера совпадает со спектральной мерой
\(F\)).

\textbf{Частотная характеристика фильтра} с импульсной характеристикой
\(h_k\) --- функция

\[
H(\lambda)=\sum_{k=-\infty}^{\infty}h_k e^{-ik\lambda},
\]

которая показывает, как линейный фильтр усиливает или подавляет разные
частоты.

\subsection{4. Основные теоремы и
свойства}\label{ux43eux441ux43dux43eux432ux43dux44bux435-ux442ux435ux43eux440ux435ux43cux44b-ux438-ux441ux432ux43eux439ux441ux442ux432ux430-24}

\subsubsection{Теорема
Герглотца}\label{ux442ux435ux43eux440ux435ux43cux430-ux433ux435ux440ux433ux43bux43eux442ux446ux430}

\textbf{Условие:} Последовательность чисел \(R(k)\),
\(k \in \mathbb{Z}\), является ковариационной функцией стационарной в
широком смысле последовательности. \textbf{Утверждение:} Это выполнено
тогда и только тогда, когда существует конечная мера \(F\) на
\([-\pi, \pi)\), такая что:

\[
R(k) = \int_{-\pi}^{\pi} e^{i k \lambda} F(d\lambda).
\]

Мера \(F\) определяется функцией \(R\) единственным образом. Кроме того,

\[
F([-\pi,\pi))=R(0),
\]

то есть полная масса спектральной меры равна дисперсии центрированного
процесса.

\textbf{Схема доказательства:} 1. Необходимость. Так как \(R(k)\)
положительно определена на решетке \(\mathbb{Z}\), строят сглаженные
неотрицательные тригонометрические многочлены:

\[
f_N(\lambda) = \frac{1}{2\pi N} \sum_{n=0}^{N-1} \sum_{m=0}^{N-1} R(n-m) e^{-i(n-m)\lambda} \ge 0.
\]

Положительность следует из неотрицательной определенности \(R\).
Интеграл от \(f_N(\lambda)\) равен \(R(0)\). Меры
\(F_N(d\lambda) = f_N(\lambda) d\lambda\) образуют слабо компактное
семейство (по теореме Прохорова). Их слабый предел \(F\) и будет искомой
мерой.

\begin{enumerate}
\def\labelenumi{\arabic{enumi}.}
\setcounter{enumi}{1}
\tightlist
\item
  Достаточность. Если интеграл существует, то положительная
  определенность \(R(k)\) очевидна:
\end{enumerate}

\[
\sum_{n,m} a_n \overline{a_m} R(n-m) = \int_{-\pi}^{\pi} \left| \sum_n a_n e^{i n \lambda} \right|^2 F(d\lambda) \ge 0.
\]

По теореме Колмогорова (см. Билет 15) существование неотрицательно
определенной матрицы всегда гарантирует существование гауссовского
стационарного процесса с такой ковариацией.

\subsubsection{Теорема Карунена о спектральном
представлении}\label{ux442ux435ux43eux440ux435ux43cux430-ux43aux430ux440ux443ux43dux435ux43dux430-ux43e-ux441ux43fux435ux43aux442ux440ux430ux43bux44cux43dux43eux43c-ux43fux440ux435ux434ux441ux442ux430ux432ux43bux435ux43dux438ux438}

\textbf{Условие:} Пусть \(X_n\) --- центрированная стационарная в
широком смысле последовательность со спектральной мерой \(F\).
\textbf{Утверждение:} Существует случайная мера \(Z\) с ортогональными
приращениями, такая что \(\mathbb E|Z(\Delta)|^2 = F(\Delta)\) и

\[
X_n = \int_{-\pi}^{\pi} e^{i n \lambda} Z(d\lambda).
\]

\textbf{Схема доказательства:} Зададим отображение
\(I: e^{i n \lambda} \mapsto X_n\). Оно сохраняет скалярное
произведение, так как
\(\langle e^{i n \lambda}, e^{i m \lambda} \rangle_{L_2(F)} = \int e^{i(n-m)\lambda}F(d\lambda) = R(n-m) = \langle X_n, X_m \rangle_{H_X}\).
По линейности продолжим \(I\) на тригонометрические полиномы. Множество
тригонометрических полиномов плотно в \(L_2(F)\). Значит, мы можем по
непрерывности продолжить \(I\) на все пространство \(L_2(F)\). Определим
меру \(Z\) как образы индикаторов:

\[
Z(\Delta) = I(\mathbf{1}_\Delta(\lambda)).
\]

Тогда \(Z(\Delta)\) ортогональны для непересекающихся множеств (т.к.
ортогональны их индикаторы в \(L_2(F)\)), а \(X_n\) получается
интегрированием функции \(e^{in\lambda}\) по этой мере.

\subsubsection{Влияние линейных
фильтров}\label{ux432ux43bux438ux44fux43dux438ux435-ux43bux438ux43dux435ux439ux43dux44bux445-ux444ux438ux43bux44cux442ux440ux43eux432}

Если пропустить процесс \(X_n\) через устойчивый линейный фильтр с
импульсной характеристикой \(h_k\) (где \(\sum |h_k| < \infty\)):

\[
Y_n = \sum_{k=-\infty}^{\infty} h_k X_{n-k},
\]

то спектральная плотность преобразуется путем умножения на квадрат
модуля передаточной функции (частотной характеристики) фильтра:

\[
f_Y(\lambda) = |H(\lambda)|^2 f_X(\lambda), \quad \text{где } H(\lambda) = \sum_{k=-\infty}^{\infty} h_k e^{-i k \lambda}.
\]

\subsection{5. Примеры}\label{ux43fux440ux438ux43cux435ux440ux44b-24}

\subsubsection{Пример 1. Белый шум (процесс без
памяти)}\label{ux43fux440ux438ux43cux435ux440-1.-ux431ux435ux43bux44bux439-ux448ux443ux43c-ux43fux440ux43eux446ux435ux441ux441-ux431ux435ux437-ux43fux430ux43cux44fux442ux438}

Пусть \(X_n = \varepsilon_n\), где \(\varepsilon_n\) ---
некоррелированные величины:

\[
\mathbb E \varepsilon_n = 0, \quad \operatorname{Var}(\varepsilon_n) = \sigma^2, \quad \operatorname{Cov}(\varepsilon_n, \varepsilon_m) = 0 \text{ при } n \ne m.
\]

Ковариационная функция:

\[
R(0) = \sigma^2, \quad R(k) = 0 \text{ для } k \ne 0.
\]

Ряд Фурье для плотности:

\[
f(\lambda) = \frac{1}{2\pi} \sum_{k=-\infty}^{\infty} R(k) e^{-i k \lambda} = \frac{\sigma^2}{2\pi}.
\]

\textbf{Вывод:} Спектральная плотность белого шума константна. В
процессе присутствуют все частоты от \(-\pi\) до \(\pi\) в равной
пропорции. Отсюда аналогия с белым светом в оптике.

\subsubsection{Пример 2. Гармоническое колебание со случайной фазой
(линейчатый
спектр)}\label{ux43fux440ux438ux43cux435ux440-2.-ux433ux430ux440ux43cux43eux43dux438ux447ux435ux441ux43aux43eux435-ux43aux43eux43bux435ux431ux430ux43dux438ux435-ux441ux43e-ux441ux43bux443ux447ux430ux439ux43dux43eux439-ux444ux430ux437ux43eux439-ux43bux438ux43dux435ux439ux447ux430ux442ux44bux439-ux441ux43fux435ux43aux442ux440}

Пусть

\[
X_n = A \cos(\omega n + \Phi),
\]

где \(\omega \in (0, \pi)\) и \(A > 0\) --- детерминированные числа, а
\(\Phi \sim U[0, 2\pi]\) --- случайная величина.

Считаем ковариацию:

\[
R(k) = \mathbb E[A \cos(\omega(n+k) + \Phi) A \cos(\omega n + \Phi)].
\]

Используя формулу произведения косинусов, получаем:

\[
R(k) = \frac{A^2}{2} \cos(\omega k).
\]

Представим косинус через экспоненты:

\[
R(k) = \frac{A^2}{4} e^{i \omega k} + \frac{A^2}{4} e^{-i \omega k}.
\]

\textbf{Вывод:} Это уже преобразование Фурье для меры \(F\), которая
состоит ровно из двух атомов (дельта-функций Дирака, см. Билет 20) в
точках \(\lambda_1 = \omega\) и \(\lambda_2 = -\omega\), масса каждого
из которых равна \(\frac{A^2}{4}\). Плотности \(f(\lambda)\) по мере
Лебега у такого процесса \textbf{не существует}. Весь ``спектр''
сосредоточен в двух конкретных частотах.

\subsubsection{Пример 3. Авторегрессия первого порядка
AR(1)}\label{ux43fux440ux438ux43cux435ux440-3.-ux430ux432ux442ux43eux440ux435ux433ux440ux435ux441ux441ux438ux44f-ux43fux435ux440ux432ux43eux433ux43e-ux43fux43eux440ux44fux434ux43aux430-ar1}

Пусть процесс задан разностным уравнением:

\[
X_n = a X_{n-1} + \varepsilon_n, \quad |a| < 1,
\]

где \(\varepsilon_n\) --- белый шум с дисперсией \(\sigma^2\).

Как мы знаем, его ковариационная функция равна:

\[
R(k) = \frac{\sigma^2}{1-a^2} a^{|k|}.
\]

Найдем спектральную плотность. Так как процесс получается из белого шума
фильтром \(H(\lambda) = \frac{1}{1 - a e^{-i\lambda}}\) (это следует из
\(X_n - a X_{n-1} = \varepsilon_n\), т.е.
\(X(1 - a e^{-i\lambda}) = \varepsilon\)), то по свойству фильтра:

\[
f_X(\lambda) = |H(\lambda)|^2 f_\varepsilon(\lambda) = \left| \frac{1}{1 - a e^{-i\lambda}} \right|^2 \frac{\sigma^2}{2\pi}.
\]

Распишем знаменатель:

\[
|1 - a \cos\lambda + i a \sin\lambda|^2 = (1 - a \cos\lambda)^2 + (a \sin\lambda)^2 = 1 - 2a \cos\lambda + a^2.
\]

\textbf{Вывод:}

\[
f(\lambda) = \frac{\sigma^2}{2\pi (1 - 2a \cos\lambda + a^2)}.
\]

Если \(a > 0\), спектр имеет пик на низких частотах (возле
\(\lambda = 0\)), процесс имеет длинные ``волны''. Если \(a < 0\), пик
на высоких частотах (возле \(\lambda = \pi\)), процесс ``мелко дрожит''.

\subsubsection{Пример 4. Скользящее среднее
MA(1)}\label{ux43fux440ux438ux43cux435ux440-4.-ux441ux43aux43eux43bux44cux437ux44fux449ux435ux435-ux441ux440ux435ux434ux43dux435ux435-ma1}

\[
X_n = \varepsilon_n + \theta \varepsilon_{n-1}.
\]

Ковариация: \(R(0) = \sigma^2(1+\theta^2)\),
\(R(\pm 1) = \sigma^2\theta\), \(R(|k|>1) = 0\).

Спектральная плотность:

\[
f(\lambda) = \frac{1}{2\pi} \left( \sigma^2(1+\theta^2) + 2\sigma^2\theta \cos\lambda \right) = \frac{\sigma^2}{2\pi} |1 + \theta e^{-i\lambda}|^2.
\]

\subsection{6. Контрпримеры и
предупреждения}\label{ux43aux43eux43dux442ux440ux43fux440ux438ux43cux435ux440ux44b-ux438-ux43fux440ux435ux434ux443ux43fux440ux435ux436ux434ux435ux43dux438ux44f}

\textbf{Не каждый процесс имеет плотность.} Ошибка думать, что
спектральная плотность существует всегда. Если в процессе есть строго
периодические неразрушаемые компоненты (как в Примере 2), их спектр
дискретный. Формула
\(f(\lambda) = \frac{1}{2\pi}\sum R(k) e^{-ik\lambda}\) в этом случае
просто не будет сходиться.

\textbf{Спектральная мера не является вероятностной.} В теореме
Бохнера-Хинчина (для характеристических функций) интеграл от меры равен
\(\varphi(0) = 1\), то есть мера вероятностная. В теореме Герглотца
интеграл от меры равен \(R(0) = \operatorname{Var}(X)\), что обычно не
равно 1.

\textbf{Случайная мера \(Z\) комплекснозначна даже для вещественного
процесса.} Ортогональная мера \(Z\) в спектральном разложении
\(X_n = \int e^{in\lambda} Z(d\lambda)\) принимает комплексные значения,
чтобы компенсировать экспоненты \(e^{in\lambda}\). Чтобы сам процесс
\(X_n\) был вещественным, необходимо выполнение условия симметрии:

\[
Z(\Delta) = \overline{Z(-\Delta)}.
\]

\subsection{7. Главные формулы для
доски}\label{ux433ux43bux430ux432ux43dux44bux435-ux444ux43eux440ux43cux443ux43bux44b-ux434ux43bux44f-ux434ux43eux441ux43aux438}

\begin{enumerate}
\def\labelenumi{\arabic{enumi}.}
\tightlist
\item
  \textbf{Теорема Герглотца:}
\end{enumerate}

\[
R(k) = \int_{-\pi}^{\pi} e^{i k \lambda} F(d\lambda).
\]

\begin{enumerate}
\def\labelenumi{\arabic{enumi}.}
\setcounter{enumi}{1}
\tightlist
\item
  \textbf{Прямое и обратное преобразования Фурье для плотности:}
\end{enumerate}

\[
f(\lambda) = \frac{1}{2\pi} \sum_{k=-\infty}^{\infty} R(k) e^{-i k \lambda}, \quad
R(k) = \int_{-\pi}^{\pi} e^{i k \lambda} f(\lambda) d\lambda.
\]

\begin{enumerate}
\def\labelenumi{\arabic{enumi}.}
\setcounter{enumi}{2}
\tightlist
\item
  \textbf{Спектральное представление процесса:}
\end{enumerate}

\[
X_n = \int_{-\pi}^{\pi} e^{i n \lambda} Z(d\lambda).
\]

\begin{enumerate}
\def\labelenumi{\arabic{enumi}.}
\setcounter{enumi}{3}
\tightlist
\item
  \textbf{Ортогональность приращений спектральной случайной меры:}
\end{enumerate}

\[
\mathbb E[Z(\Delta_1) \overline{Z(\Delta_2)}] = F(\Delta_1 \cap \Delta_2).
\]

\begin{enumerate}
\def\labelenumi{\arabic{enumi}.}
\setcounter{enumi}{4}
\tightlist
\item
  \textbf{Белый шум:}
\end{enumerate}

\[
f(\lambda) = \frac{\sigma^2}{2\pi}.
\]

\begin{enumerate}
\def\labelenumi{\arabic{enumi}.}
\setcounter{enumi}{5}
\tightlist
\item
  \textbf{Эргодичность по среднему через спектр:}
\end{enumerate}

\[
\mathbb E|\overline X_N|^2
=
\int_{-\pi}^{\pi}
\left|
\frac1N\sum_{n=0}^{N-1}e^{in\lambda}
\right|^2F(d\lambda),
\qquad
\overline X_N\xrightarrow{L^2}0 \Longleftrightarrow F(\{0\})=0.
\]

\subsection{8. Типичные дополнительные
вопросы}\label{ux442ux438ux43fux438ux447ux43dux44bux435-ux434ux43eux43fux43eux43bux43dux438ux442ux435ux43bux44cux43dux44bux435-ux432ux43eux43fux440ux43eux441ux44b-24}

\textbf{Вопрос:} Почему для дискретного времени спектр ограничен
отрезком \([-\pi, \pi]\), а для непрерывного --- всей прямой
\(\mathbb{R}\)? \textbf{Ответ:} Для дискретного времени
\(n \in \mathbb{Z}\) базисные функции \(e^{i n \lambda}\) периодичны по
\(\lambda\) с периодом \(2\pi\). То есть
\(e^{i n (\lambda + 2\pi)} = e^{i n \lambda}\). Мы физически не можем
отличить частоту \(\lambda\) от \(\lambda + 2\pi\) (эффект наложения
частот, или алиасинг). Для непрерывного времени \(t \in \mathbb{R}\)
функции \(e^{i t \lambda}\) при разных \(\lambda\) линейно независимы на
всей оси.

\textbf{Вопрос:} Как по спектральной плотности найти дисперсию процесса?
\textbf{Ответ:} Проинтегрировать ее по всем частотам:
\(\operatorname{Var}(X) = R(0) = \int_{-\pi}^{\pi} f(\lambda) d\lambda\).

\textbf{Вопрос:} Может ли спектральная плотность быть отрицательной в
каких-то точках? \textbf{Ответ:} Нет, \(f(\lambda) \ge 0\). Это прямое
следствие свойства неотрицательной определенности ковариационной
функции.

\textbf{Вопрос:} Чем отличаются спектральная мера, спектральная функция
и спектральная плотность? \textbf{Ответ:} Спектральная мера \(F\) ---
основной объект. Спектральная функция
\(\mathcal F(\lambda)=F([-\pi,\lambda])\) просто записывает эту меру как
функцию распределения по частоте. Спектральная плотность \(f\)
существует только в абсолютно непрерывном случае, когда
\(F(d\lambda)=f(\lambda)d\lambda\).

\textbf{Вопрос:} Почему спектральная мера не обязана быть вероятностной?
\textbf{Ответ:} Ее полная масса равна не \(1\), а
\(R(0)=\operatorname{Var}X_0\) для центрированного процесса. Если
хочется вероятностную меру, можно нормировать \(F/R(0)\) при \(R(0)>0\),
но в спектральной теории обычно важна именно ненормированная мощность.

\textbf{Вопрос:} Что делает линейный фильтр со спектром? \textbf{Ответ:}
Если \(Y_n=\sum_k h_kX_{n-k}\) и \(H(\lambda)=\sum_kh_ke^{-ik\lambda}\),
то спектральная мера умножается на \(|H(\lambda)|^2\). При наличии
плотности это записывается как
\(f_Y(\lambda)=|H(\lambda)|^2f_X(\lambda)\).

\textbf{Вопрос:} Связь спектра и эргодичности по среднему?
\textbf{Ответ:} В Билете 24 мы вывели, что среднее
\(\frac{1}{N}\sum_{n=0}^{N-1} X_n\) сходится в \(L_2\) к нулю (если
\(m=0\)). Через спектральную меру дисперсия среднего выражается как
интеграл от ядра Фейера. В пределе среднее сходится к величине
\(Z(\{0\})\). Поэтому эргодичность эквивалентна отсутствию атома в нуле
у спектральной меры: \(F(\{0\}) = 0\).

\subsection{9. Типичные
ошибки}\label{ux442ux438ux43fux438ux447ux43dux44bux435-ux43eux448ux438ux431ux43aux438-24}

\begin{itemize}
\tightlist
\item
  \textbf{Путать спектральную функцию \(\mathcal F(\lambda)\) и
  спектральную плотность \(f(\lambda)\).} Забывать, что у процесса может
  не быть плотности, если мера сингулярна или дискретна.
\item
  \textbf{Ошибки в нормировке.} Часто забывают множитель \(1/2\pi\). В
  нашем курсе он стоит перед суммой для \(f(\lambda)\), а перед
  интегралом для \(R(k)\) его нет.
\item
  \textbf{Путать \(Z(\Delta)\) и \(F(\Delta)\).} \(Z\) --- это
  \emph{случайная} мера (ее значения --- случайные величины), а \(F\)
  --- \emph{детерминированная} мера (ее значения --- действительные
  числа, дисперсии).
\item
  \textbf{Считать интеграл Риманом.} Интеграл
  \(X_n = \int e^{in\lambda} Z(d\lambda)\) --- это не интеграл
  Лебега-Стилтьеса, а стохастический интеграл, понимаемый как предел в
  \(L_2(\Omega)\).
\end{itemize}

\subsection{10. Связи с другими
билетами}\label{ux441ux432ux44fux437ux438-ux441-ux434ux440ux443ux433ux438ux43cux438-ux431ux438ux43bux435ux442ux430ux43cux438-24}

\begin{itemize}
\tightlist
\item
  Билет 08: Характеристические функции. Теорема Герглотца --- родная
  сестра теоремы Бохнера-Хинчина. Ковариационная функция играет роль
  характеристической функции, а спектральная мера --- роль распределения
  вероятностей.
\item
  Билет 12: Гильбертовы пространства. Изометрия между \(H_X\) и
  \(L_2(F)\) --- красивейший пример применения функционального анализа в
  теории вероятностей.
\item
  Билет 17: Стационарность в широком смысле --- основное условие для
  всей спектральной теории.
\item
  Билет 20: Дельта-функции Дирака нужны для описания дискретных спектров
  (как в Примере 2).
\item
  Билет 24: Эргодичность тесно связана со спектром (атом в нуле).
\end{itemize}

\subsection{11. Минимум на
удовлетворительно}\label{ux43cux438ux43dux438ux43cux443ux43c-ux43dux430-ux443ux434ux43eux432ux43bux435ux442ux432ux43eux440ux438ux442ux435ux43bux44cux43dux43e-24}

Нужно знать: 1. Что такое теорема Герглотца (написать формулу
\(R(k) = \int_{-\pi}^{\pi} e^{i k \lambda} dF(\lambda)\)). 2. Что спектр
для дискретного времени ограничен: \(\lambda \in [-\pi, \pi]\). 3. Знать
формулы, связывающие \(R(k)\) и \(f(\lambda)\) (преобразование Фурье с
множителем \(1/2\pi\)). 4. Знать, что если спектральная плотность
существует, то она неотрицательна: \(f(\lambda) \ge 0\); в общем случае
есть спектральная мера. 5. Уметь привести пример белого шума: его
ковариация равна 0 везде кроме \(k=0\), а его спектральная плотность ---
константа \(f(\lambda) = \frac{\sigma^2}{2\pi}\).

\subsection{12. Минимум на
хорошо}\label{ux43cux438ux43dux438ux43cux443ux43c-ux43dux430-ux445ux43eux440ux43eux448ux43e-24}

К минимуму добавляется: 1. Спектральное представление самого процесса
\(X_n = \int_{-\pi}^{\pi} e^{i n \lambda} Z(d\lambda)\). 2. Понимание
того, что такое ортогональная случайная мера \(Z(\Delta)\) и ее связь с
\(F(\Delta)\) через уравнение дисперсий:
\(\mathbb E|Z(\Delta)|^2 = F(\Delta)\). 3. Вывод спектральной плотности
для авторегрессии AR(1) или скользящего среднего MA(1).

\subsection{13. Что добавить для
отлично}\label{ux447ux442ux43e-ux434ux43eux431ux430ux432ux438ux442ux44c-ux434ux43bux44f-ux43eux442ux43bux438ux447ux43dux43e-24}

Для отличного ответа нужно уверенно владеть математическим аппаратом: 1.
Рассказать про изометрию гильбертовых пространств \(H_X\) и \(L_2(F)\).
2. Описать схему доказательства теоремы Герглотца (например, через
приближение функциями \(f_N(\lambda)\) и теорему Хелли/Прохорова). 3.
Привести пример с дискретным спектром (гармоническое колебание) и
объяснить, почему там нет плотности. 4. Объяснить эффект действия
линейного фильтра на спектральную плотность.

\subsection{14. Устная версия ответа на 5
минут}\label{ux443ux441ux442ux43dux430ux44f-ux432ux435ux440ux441ux438ux44f-ux43eux442ux432ux435ux442ux430-ux43dux430-5-ux43cux438ux43dux443ux442-24}

Спектральная теория позволяет посмотреть на стационарный в широком
смысле процесс как на сумму гармоник с разными частотами и случайными
амплитудами.

Так как время \(n\) у нас целое, все доступные частоты умещаются на
отрезке \([-\pi, \pi]\) --- вне его синусы и косинусы начнут
повторяться. Главный результат теории --- теорема Герглотца. Она
говорит, что ковариационная функция \(R(k)\) может быть представлена как
интеграл Фурье от некоторой конечной меры \(F\) на отрезке
\([-\pi, \pi]\): \emph{(пишем на доске формулу Герглотца)}. Мера \(F\)
называется спектральной мерой. Если у нее есть плотность \(f(\lambda)\),
то она неотрицательна, и \(R(k)\) с \(f(\lambda)\) связаны парой взаимно
обратных преобразований Фурье.

Кроме того, по теореме Карунена мы можем разложить в такой же интеграл и
сам процесс: \emph{(пишем интеграл
\(X_n = \int e^{in\lambda} Z(d\lambda)\))}. Здесь \(Z\) --- это уже не
обычная мера, а случайная функция множества с ортогональными
приращениями. Её дисперсия для любого интервала \(\Delta\) в точности
равна значению спектральной меры \(F(\Delta)\). Это представление
порождает идеальный словарь перевода: случайные величины из \(L_2\)
переводятся в детерминированные функции из \(L_2(F)\). Скалярные
произведения сохраняются.

Простейший пример --- белый шум. У него ковариация --- это ``иголочка''
в нуле. Значит, его спектральная плотность --- это константа. Все
частоты представлены равномерно. Если же процесс --- это просто синус со
случайной фазой, то его спектр дискретный --- это две дельта-функции на
частоте синуса.

\subsection{15. Если осталось 2 минуты перед
экзаменом}\label{ux435ux441ux43bux438-ux43eux441ux442ux430ux43bux43eux441ux44c-2-ux43cux438ux43dux443ux442ux44b-ux43fux435ux440ux435ux434-ux44dux43aux437ux430ux43cux435ux43dux43eux43c-24}

\begin{itemize}
\tightlist
\item
  \textbf{Время:} дискретно (\(n \in \mathbb{Z}\)). Значит, частоты
  ограничены (\(\lambda \in [-\pi, \pi]\)).
\item
  \textbf{Теорема Герглотца:}
  \(R(k) = \int_{-\pi}^{\pi} e^{i k \lambda} F(d\lambda)\). Связывает
  ковариацию и спектральную меру (дисперсию на частотах).
\item
  \textbf{Плотность:}
  \(f(\lambda) = \frac{1}{2\pi} \sum R(k) e^{-i k \lambda}\). Всегда
  \(\ge 0\).
\item
  \textbf{Сам процесс:}
  \(X_n = \int_{-\pi}^{\pi} e^{i n \lambda} Z(d\lambda)\). \(Z\) ---
  случайная ортогональная мера.
\item
  \textbf{Связь мер:} \(\mathbb E|Z(\Delta)|^2 = F(\Delta)\).
\item
  \textbf{Белый шум:}
  \(R(k) = \delta_{k0} \sigma^2 \implies f(\lambda) = \sigma^2 / 2\pi\)
  (константа).
\item
  \textbf{Ошибка:} путать \(F\) (обычная мера, задает распределение
  дисперсии) и \(Z\) (случайная мера).
\end{itemize}

\newpage

\end{document}
